% ===========================================================================
% 04b-representation.tex -- beats 3-4. The expression frame, read in the weights.
% Source: thesis/Sections/Investigation-v4.tex 1052-1545.
% Every number, interval, n and conclusion is verbatim from v1.
%
% Consistency pass: the section now opens by naming the third question of
% \cref{sec:intro-map} and stating its answer -- identity written into the
% activations, barely consulted by generation -- before the instruments, so a
% reader arriving from the introduction finds what was promised. The stratum
% paragraph and its \scopelimit are adjacent again, since the scope limit is
% about that stratum's reference. "Step three, the removal gate" was a \defn
% box between three \paragraph steps; it is now the paragraph it always was,
% merged with the gate result that followed it. Vocabulary: "output" as a noun
% for the prediction, "compound", and bare "signature" are gone; "shift" is
% reserved for treated minus control, so the logsumexp constant is an offset.
% "Energy" is gone as well, and it went two ways because it named two different
% things. The $\approx 99.5\%$ that a full-stream control puts outside the slab
% is a projection share of a squared Euclidean norm, $1 - 10/2048$ for ten
% intervention directions in $2048$, so it now reads "squared norm", in the
% sense fixed by the norm entry in the symbols list of 00-HowToRead.tex. It is
% not the residual/generic decomposition of shared/inference.py, whose
% docstring calls itself a variance decomposition and is wrong to: nothing is
% centred there either. The "strong-programme energy" of the paragraph on where
% ablation stops named no computed quantity at all, since that sentence carries
% no numeral and cites no stored result, so it now reads "structure", which
% claims only what a hypothetical ("could carry") can support.
% ===========================================================================

% ===========================================================================
\beatsection{E}{3}{The drug is not missing from the computation}{sec:methods-probe}

\begin{question}[3]
A model that ignores a drug it was told about can fail in two ways. Either the
drug never enters its computation, an activation-side encoding failure that
target- and input-side fixes might reach, or it enters and nothing reads it, a
decoder failure only an objective-side fix could reach. Behaviour cannot tell
them apart, since both predict the same scores. The weights can. A linear decode
probe, read as a floor test and not as evidence of use, settles the first half.
\end{question}

\noindent
\Cref{sec:intro-map} asks where useful prediction fails: is the drug absent from
the data, lost from the model's representation, or encoded and underused by
generation? The three sections that follow answer in one direction, and the
answer is worth having before the instruments. Drug identity is written into the
activations and stays written, more legibly with depth once batch structure is
projected out, while generation consults it too weakly to account for the
drug-blindness of \cref{sec:res-blind}. The failure is not on the activation
side. Neither half comes free: the presence result is close to guaranteed by
construction, the low-gain reading rests on the weakest measurement in the
chapter, and both are carried below with the discount attached.

Every behavioural measurement so far fits the last two readings
equally.\framecue{\Eframe}{scores here are full-profile, \cref{sec:res-metrics}}
\Cref{sec:res-blind} has answered the first where it matters here: the task is
empirically discriminable for a substantial stratum at the tested pseudobulk
budget, the reference rises with aggregation, $46.2\%$ of conditions clear the
declared identifiability bar, and on those the model is matched exactly by a
baseline copying the untreated cell ($0.768$ against $0.766$).\gloss{identifiable
stratum}{conditions whose own within-well split-half reference clears the pre-set
bar} The shortfall there is the model's; where in the model it lives needs the
weights.

\begin{scopelimit}[carried into \cref{sec:res-epsilon} and \cref{sec:limitations}]
The distance from a score to that stratum's reference is not available headroom.
Across the three difficulty strata of \cref{sec:res-blind} the control-copy's own
score tracks the reference almost step for step: $0.310$ against a reference of
$0.287$ on the low-reference stratum, $0.540$ against $0.689$ on the marginal one,
and $0.766$ against $0.938$ on the identifiable one. A predictor that reproduces
the untreated cell and knows nothing whatever about any drug does best exactly
where the ranking problem is easiest. What separates the strata on this
criterion is largely drug-independent: within a comparison set every condition's
untreated cells come from one plate-matched DMSO pool, so a control-copy scores
high exactly where a condition's measured profile sits closer to that shared
untreated state than its plate-mates' do. Most of the ranking these conditions
are scored on can therefore be done without the drug, and the gap between a predictor and a high reference is
not all of it room that better drug modelling would fill. Every score quoted on
the identifiable stratum below inherits that reading, including the three target
constructions of \cref{sec:res-epsilon}.
\end{scopelimit}

\paragraph{Three sub-questions, in the order their blind spots close} Is drug
identity present in the activations? Does the network use the region where it is
written? And does it use \emph{which} drug? Each instrument's blind spot is what
the next must close; \cref{sec:res-epsilon} adds a fourth.
\Cref{fig:comparators} maps all four, on scales that may not be read across.

% ===========================================================================
\begin{figure}[htbp]
\begin{fullwidth}
\centering
% ===========================================================================
% figs/comparators.tex -- fig:comparators, reworked for thesis_v2.
% FIGURE BODY ONLY. No preamble, no document environment. \input it from
% inside a figure/fullwidth float in Sections/04b-representation.tex.
%
% Four instruments, four deliberately incommensurable value-to-position maps,
% one comparator each. The premise and all four scales are v1's and unchanged.
% What changed is labelling only; no mark moved, no number changed.
%
%   (1) ROW B plotted seven marks and named one ("all seven layers"), leaving
%       six unidentifiable. Every mark is now named by the layer it is, with a
%       fanned leader for the three that crowd at 3.6/3.8/4.2, and the stub
%       carries "one mark per layer" so a bare accent numeral above the axis
%       cannot be misread as a second set of tick labels (ticks are quietink
%       and below the axis; layer numbers are accent and above it).
%   (2) ROW A's "slab projected out" sat to the right of its own point and
%       nearer the shuffled-label rule than to the mark it names. It now has a
%       leader back to the 0.121 dot.
%   (3) ROW C's single-cell marker was a hollow ring sitting on the dotted zero
%       rule (-0.0012 against 0.000: 0.10cm apart on this scale). It is now
%       solid, the zero rule is interrupted across its height so nothing runs
%       into it, and its label starts 0.30cm clear with a leader through the
%       interruption -- necessary because the zero rule lies between mark and
%       label and would otherwise claim the label.
%   (4) Bar labels in ROWS C and D moved from a 0.13cm to a 0.30cm gap, so
%       "optimal-transport barycentric" no longer reads as a continuation of
%       its own interval bar.
%   General rule applied throughout: a label more than ~0.3cm from its mark
%   gets a short leader.
%   Row heights grew (rows now 0.70/0.77/0.98 apart rather than colliding) to
%   pay for the new labels. Nothing is rescaled: 1cm here is 1cm on the page.
%
% PROVENANCE of every number drawn is in comparators.json, sibling to this
% file. All four rows were checked against RESULTS_cluster and against the
% thesis text; they agree. Sources:
%   row A, B  RESULTS_cluster/workspace_probe.json
%   row C     RESULTS_cluster/probe_arm_residual.json,
%             probe_rep_residual_s43.json, probe_rep_residual_s45.json,
%             probe_arm_single_cell.json
%   row D     RESULTS_cluster/stratify_sameplate_final.json,
%             stratify_ot_T2.json, stratify_consensus.json
%
% ===========================================================================
% v3 -- 2026-08-18. Supervisor's verdict, quoted so the reason survives with
% the file: "Figure 4.7 makes no sense as well by the way, there is no
% definition on the x axis for what we are measuring and moreover I dont
% really get what its reference in the is it there part, it just points to the
% slab project out but it doesnt make any sense?"
%
% SUPERSEDED, EXPLICITLY, and left standing above because the v1->v2 rationale
% is audited against it:
%   * line 8, "What changed is labelling only; no mark moved, no number
%     changed", was true of v2 and is FALSE of v3. In v3 rows C and D are
%     re-domained, so every mark in those two rows MOVES ON THE PAGE. No value
%     changes anywhere: each mark still sits at the position its own row's map
%     assigns to the value in its source JSON, and both the old and the new
%     maps are recorded in comparators.json.
%   * header note (1) is superseded: the seven accent layer numerals and the
%     brace over 3.6/3.8/4.2 are DELETED, so "marks named by layer" is gone
%     from the stub too. Numerals set 5.2pt above a value axis in the same size
%     and glyph class as the tick numerals 7.8pt below it are read as values --
%     the mark lettered "4" sat 19pt from the tick reading "20" -- and left to
%     right they read 16, 12, 9, 6, 2, 8, 4, which is neither depth order nor
%     value order of the identifiers. The layer-to-ratio map is drawn properly
%     in fig:mechanism panel (b), which gives layer a TRUE numeric axis and
%     annotates every pair with its own ratio, and printed in tab:workspace.
%     Row B's remaining job -- scale, comparator, spread -- needs no identity
%     on any single mark, and the identities are re-homed in the caption.
%   * header note (2) is superseded: the 0.121 mark, its leader and its label
%     are DELETED from row A. It is layers.<L>.drug_decode_ablated, BYTE-
%     IDENTICAL at all seven layers -- not a reading at a depth but the single
%     constant output of section 4.4's removal gate, an instrument answering a
%     different question, sitting inside a section-4.3 row 4.0mm from row A's
%     own comparator and carrying the only leader on the figure. It is drawn
%     as the gate's own numerator in fig:purify, plotted against depth in
%     fig:mechanism (a) and printed as text in fig:licensing. Restyling it as a
%     floor fails, because a floor is only a floor against a before/after pair
%     and row A draws no "before"; moving it to row B is arithmetically
%     impossible, because row B's axis is a dimensionless ratio to a matched
%     null and no position on it means an accuracy of 0.121. It is NOT
%     replaced: adding 0.758 (raw at layer 16) would put two series at two
%     depths on one rule and rebuild the before/after misreading that only
%     fig:mechanism (a), with a layer axis, can tell honestly. The mark is kept
%     in comparators.json as a removed-from-drawing entry with its provenance.
%   * header note (3) survives verbatim: the interrupted zero rule and the
%     leader through the interruption are kept, re-cut to the new row-C map.
%   * header note (4) survives: the 0.30cm label gap after an interval bar is
%     kept in both rows.
%
% WHAT v3 ADDS.
%   (i) AN AXIS-LABEL BLOCK UNDER EVERY ONE OF THE FOUR RULES, in axis-label
%       position -- directly under the tick numerals, left-aligned to the
%       rule's left end -- naming the measured quantity and its unit on a
%       \footnotesize roman ink line, with at most one \scriptsize quietink
%       scope line under it. This is the direct answer to the supervisor's
%       first complaint and it is a transcription, not an invention: the four
%       quantity strings were already carried in comparators.json's own
%       "quantity" fields and in the grey sub-lines of the old left stub, i.e.
%       14.1pt to the LEFT of the rule, inside the question label, where a
%       reader looking at the scale had nothing. Row D's frame is stated
%       (EXPRESSION) because expression-frame and residual-frame NIR are not
%       comparable; it is resolved, not guessed, from
%       stratify_sameplate_final.json :: test2.all.model = 0.4977117949890227,
%       the 0.498 that figs/style_v2.py names as the expression-frame reading
%       against the residual frame's 0.537, and from 04b-representation.tex,
%       which declares itself the expression-frame file at line 2 and departs
%       once, at line 452, for the section-4.5 residual arm -- which is row C
%       and not row D.
%  (ii) THE ROW LETTERS (a)-(d) on the stub, so the caption can address a row
%       instead of listing four clauses and hoping the reader maps them in
%       order. The stub is cut to the question alone, \footnotesize roman ink,
%       right-aligned at one x; all seven grey sub-lines are gone, their
%       content having moved to the axis blocks where it belongs.
% (iii) ONE THRESHOLD, ONE STROKE. Row A's chance rule and row C's
%       pre-registered margin are now BLACK, SOLID, 0.90pt, labelled
%       \scriptsize roman black -- the register fig:mechanism uses for chance
%       and fig:family for the same 10% margin. In v2 one pre-registered
%       threshold had three renderings inside nine pages and one chance floor
%       had two; this figure was carrying the demoted one of each. Densely
%       dotted 0.50pt quietink is now reserved for a coordinate that decides
%       nothing: the two zero rules, and row B's 1x rule -- which is not a
%       threshold but the UNIT OF ITS OWN AXIS by construction, a definitional
%       comparator, and is kept dotted for that reason. DEPARTURE, RECORDED:
%       the rebuild brief asked for that rule to be relabelled "matched-
%       dimension random slab" in roman. It is labelled "1x, by construction"
%       instead, because the axis-label line 3.5cm below it now carries
%       "Ablation cost, in multiples of the matched-dimension random slab" --
%       the object is named once, in axis-label position, which is where S5
%       requires the unit to be named, and setting the same four-word string
%       twice inside one row is the kind of thing this pass exists to remove.
%       fig:mechanism (b)'s name for the object is the one used, and it is used
%       once. Both threshold rules and both zero rules now start at the tick depth
%       (-0.10) and rise; in v2 the chance rule ran down to -0.14, through the
%       numeral lane, and merged with the tick furniture.
%  (iv) ROW A'S TWO MARKS ARE DATED AND DISTINGUISHED BY GLYPH. 0.821 is
%       layers.9.drug_decode and 0.883 is layers.16.drug_decode_context_removed
%       -- two series at two different depths. Undated and separated only by a
%       colour difference on a 1.4pt disc they read as a paired before/after at
%       one depth, which is exactly the reading the section's finding depends
%       on NOT being made: the two series CROSS between layers 9 and 12, and
%       context removal RAISES accuracy at 12 and 16. The raw mark is now an
%       OPEN accent disc and the context-removed mark a FILLED one, which is
%       fig:mechanism (a)'s own encoding of these same two series. Both labels
%       sit on ONE baseline; the second is right-aligned to the rule's right
%       end and leadered back to its mark, so nothing crosses the measure. In
%       v2 "context removed" ran 11.02pt past the end of the row dividers and
%       set the figure's right edge with a stray label.
%   (v) ROWS C AND D RE-DOMAINED so the labelled ticks reach BOTH ends of every
%       rule. Row C: [-0.02, +0.12], ticks every 0.02 (v2 stopped its numerals
%       22.87pt short of the right end). Row D: [-0.04, +0.06], ticks every
%       0.02 (v2 was inset 37.20pt at BOTH ends). The four rules were already
%       congruent to better than 0.01pt and stay so; what was ragged was the
%       tick lane, and the last printed number of the four rows now sits at one
%       x and the first at one x. Row D must run to +0.06 and not stop at
%       +0.04 because the single-cell interval's upper limit is +0.042 and an
%       interval may never be clipped by its own axis; row C's right end at
%       +0.12 keeps the 0.10 margin INBOARD, so the margin reads as a
%       comparator and not as the axis terminus.
%  (vi) NAMES FROM THE AUTHORITATIVE TABLES. Row C's series are tab:identity's
%       own "residual, seed 42 / 43 / 45" and "single-cell arm"; "median over
%       its depths" was the widest label in the row at 136pt and is caption
%       material. Row D's middle arm takes tab:epsilon's "OT barycentric", and
%       the consensus arm now names the reason it is drawn quietink while the
%       other two are accent -- "consensus, stopped mid-epoch" (04b:972 and
%       tab:epsilon's note). Row D's comparator, the only one of the four left
%       anonymous in a figure captioned "each instrument answers a different
%       question against a different comparator", is labelled "zero: no
%       scramble gap" -- in QUIETINK, and this is the second departure from the
%       rebuild brief, which asked for that label in ink. On this figure a
%       rule's own label takes the rule's own colour, so a black label always
%       names a black threshold (rows A and C) and a quietink label always
%       names a quietink dotted coordinate (rows B and D); zero is the null of
%       a difference, a comparator and not a threshold, and setting its label
%       black would be the one place the two registers crossed.
%       Row A's floor is "chance, 1/12": workspace_probe.json
%       stores chance = 0.08333333333333333, exactly 1/12 in IEEE double, and
%       contains no permutation or shuffle key of any kind, so "shuffled
%       labels" asserted on a RULE a measured comparator the run never stored.
%       The inherited wording survives in the prose at 04b:152 and in captions.
% (vii) ONE INTERVAL DECLARATION AT THE FOOT, naming the construction for every
%       row that has bars and stating that the other two return none. Row C's
%       bars are two-sided cluster bootstraps over 58 ordered PAIR clusters and
%       row D's are ONE-WAY percentile bootstraps over 40 CELL LINES: two
%       designs drawn in v2 with the identical 1.0pt bar and the identical dot,
%       in adjacent rows, with the caption saying only "bars are 95% intervals
%       where the instrument returns one" and rows A and B saying nothing at
%       all, so a reader inferred their instrument does return intervals.
% (viii) WIDENED from a 10.5cm rule to 12.10cm, drawn width 160.3mm -> 172.7mm,
%       measured on the render at 200dpi. The four scales had 105mm of a 174mm
%       measure while 112pt went to the label stub. One constant, \RL, now
%       carries the rule length for all four rows, so the rules cannot drift
%       apart in a later edit; measured, all four run x = 198.87 to 541.86pt in
%       the standalone render, congruent to 0.00pt, and the first and last
%       DRAWN TICK of every row sits on those two endpoints. No value moves
%       relative to its own axis.
%   (ix) THE THREE ROW DIVIDERS ARE DELETED. Each was a rulegrey rule spanning
%       100% of the drawn width, and with the axis blocks added its nearest
%       neighbour above became a line of definitional prose rather than a
%       panel, which is the configuration a full-width rule may not have. They
%       also cost about a twentieth of the whole ink budget. The four rows are
%       separated the way fig:mechanism separates its panels: by whitespace,
%       14 to 15pt of it, and by the fact that each row now opens with its own
%       head and closes with its own axis block.
%    (x) THE PANEL HEAD SITS AT THE TOP OF ITS ROW, not beside the rule. In v2
%       the stub was vertically centred on the row, so in rows C and D -- whose
%       marks stack ABOVE the rule -- the question arrived after the data. All
%       four heads now share one baseline rule of their own: the top text line
%       of the row. Reading order is identical in every row: head, marks, rule,
%       tick numerals, quantity, scope. The heads stay right-aligned at one x
%       (measured: all four end at 187.33pt) and the six left-aligned text
%       blocks of each row all start at x = 199.07pt, spread 0.00pt.
%
% WHAT v3 DELIBERATELY DOES NOT DO. The four rows keep four separate scales,
% and the zeros of rows C and D are NOT vertically aligned (measured: 248.06pt
% and 336.27pt). A shared zero would assert the commensurability this figure
% exists to deny. The interval-availability encoding is untouched: an
% instrument that returns a clustered bootstrap is drawn as a bar ABOVE the
% rule, one that returns none as a bare point mark ON it. The four domains are
% still set by their comparators and not by their data -- row A runs the full 0
% to 1, row C's right end is set by the pre-registered margin, rows C and D
% both show their negative side. Nothing is cropped to flatter.
%
% STROKE WEIGHTS, measured off the render: 0.30pt furniture (27 objects: tick
% marks and three leaders), 0.40pt axis rule (4), 0.50pt dotted coordinate (4:
% row B's 1x rule, row C's two-piece zero rule, row D's zero rule), 0.90pt
% threshold (2: chance and the pre-registered margin, and nothing else on the
% figure is black solid 0.90pt), 1.00pt interval bar (6: exactly the six marks
% whose instruments return a clustered bootstrap). A SIXTH weight, 0.80pt, is
% used exactly once and is justified here rather than removed: it is not a rule
% weight at all but the mark-edge weight of the one OPEN disc, row A's raw
% probe, which the mark vocabulary fixes at 1.4pt radius, white fill, 0.8pt
% edge. There is exactly one comparator-encoded point mark on the figure, so
% the weight is used exactly once by construction.
%
% ===========================================================================
% v4 -- 2026-08-18. REPAIR PASS on v3, after three adversarial reads of the v3
% render. Eleven defects were confirmed against the render and the sources and
% are repaired below; six proposed repairs were overruled and the reason for
% each is recorded, because an overruled finding that is not written down comes
% back. No value changed. Two marks moved and the reason is that they now sit
% at their SOURCE ratio rather than at its one-decimal rounding.
%
% CONFIRMED AND REPAIRED.
%   (R1) THE FOOT BLOCK ASSERTED SOMETHING FALSE ABOUT ROW B, AND DID NOT COVER
%        ROW C'S FOURTH MARK. v3 read "Intervals: (a) and (b) return none.
%        (c): 95% bootstrap over 58 ordered pair clusters." Both halves are
%        wrong. Row B's instrument DOES return a dispersion --
%        workspace_probe.json stores every KL term as [mean, se, n], e.g.
%        layers.9.kl.pure_drug = [0.41758305778106053, 0.005008675034860955,
%        60] -- a normal-approximation standard error, which this document
%        declines to draw because it is not its interval convention. That is
%        the sentence figs/hook.json, figs/purify.json and figs/fig-mechanism
%        .json all use, and this figure alone had upgraded it to "returns
%        none". Row A really does store a bare float
%        (layers.9.drug_decode = 0.8208333333333334, no dispersion at all), so
%        the two rows were lumped under a claim true of only one. Separately,
%        "58 ordered pair clusters" is the ANCHOR SEED's count only: seed 43
%        bootstraps over 193 pair clusters (n = 240) and seed 45 over 199
%        (n = 240), so one printed numeral described three bars and was wrong
%        for two of them. The three counts are recorded in comparators.json and
%        re-homed to the caption; the foot now names the CONSTRUCTION, which is
%        what the interval rule asks of it, and says that the fourth mark in
%        row C carries none.
%   (R2) ROW C'S SCOPE LINE WAS FALSE OF ONE OF THE FOUR MARKS IT SAT UNDER.
%        "residual arm, layer 12, alpha = 0.5; one mark per seed" is true of
%        three marks and false of the fourth: the single-cell mark is a
%        different arm, has NO layer-12 value at all
%        (probe_arm_single_cell.json :: layers.12 carries no swap block;
%        headroom 3.8803978061910347 against config.headroom 5.0, and
%        swap_skipped_reasons records "12: instrument saturated"), and is the
%        median over the six depths that were measured. The scope line now
%        reads "residual seeds at layer 12, alpha = 0.5; single-cell median,
%        six depths", which covers all four marks and is tab:identity's own
%        description of its last row ("single-cell arm, six depths").
%   (R3) ROW C'S SINGLE-CELL MARK CARRIED A HORIZONTAL 0.30pt QUIETINK LEADER
%        AT ITS OWN y, IN ITS OWN COLOUR, 0.24cm long, 0.40cm below three
%        1.00pt quietink/accent interval bars of the same geometry. At print
%        size it read as a one-sided interval that the source does not store,
%        which is the one thing the interval rule forbids. The leader is
%        DELETED. It existed only because the label sat 0.35cm from the mark's
%        centre, i.e. 0.30cm from its edge, over the 3mm threshold at which a
%        label takes a leader; the label is moved to 0.22cm from the centre
%        (0.17cm, 4.8pt, from the marker's edge), under the threshold, so no
%        leader is needed. The zero rule's interruption across the marker's
%        height is kept, and now serves the label as well: nothing lies between
%        mark and label.
%   (R4) FIVE VERTICAL RULES RAN DOWN INTO THE TICK LANE AND IMPERSONATED
%        TICKS. Every tick is a 0.30pt rulegrey stroke from the axis rule to
%        axis-0.10. In v3 the chance rule, row B's 1x rule, row C's zero and
%        margin rules and row D's zero rule all started at -0.10 as well, so:
%        the margin rule printed BLACK at 0.90pt over the 0.10 tick of a
%        seven-tick grey ladder; both zero rules printed the "0" tick DOTTED
%        while every other tick in the same ladder was solid; and chance and
%        the 1x rule each added a sixth, unnumbered downstroke to a five-tick
%        ladder. All five now stop at the axis rule. The tick lane is uniform
%        furniture again and no coordinate rule is doing furniture's job.
%   (R5) ROW A'S TWO LABELS COULD NOT BE RESOLVED TO THEIR MARKS. Measured on
%        the v3 render: the open disc (raw probe) sat at x = 480.5pt, INSIDE
%        the 451.9-541.7pt span of the label "context removed, layer 16" that
%        names the other mark; "raw probe, layer 9" sat 15mm to the left of its
%        own mark and reached it by a 5-degree diagonal hairline that travelled
%        beneath the first 27pt of that other label; and both leaders stopped
%        1.15pt short of the discs, so each pointed at nothing. The two labels
%        cannot both sit beside their marks on one baseline -- they are 63pt
%        and 90pt wide and their marks are 21.3pt apart -- and one baseline is
%        what the label rule requires. So the LAYER MOVES OFF THE MARK LABEL
%        AND INTO THE AXIS BLOCK, which is where fig:hook and fig:materiality
%        already disclose theirs and where the cross-figure rule accepts it:
%        line two of row A now reads "twelve-way drug decode; raw at layer 9,
%        context removed at layer 16", tying each series to its own depth in
%        axis-label position. The marks are then labelled "raw probe" and
%        "context removed", 31pt and 55pt wide, and the geometry closes:
%        "raw probe" is right-aligned ON its own mark's x, "context removed" is
%        right-aligned to the rule's right end and stands over its own mark and
%        over no other, and each takes an identical 0.30pt accent drop leader
%        that ENDS ON THE EDGE OF ITS DISC. One baseline, one leader idiom,
%        neither label spanning the other's mark, nothing crossing the measure.
%        DEPARTURE FROM THE REBUILD BRIEF, RECORDED: the brief asked for the
%        strings "raw probe, layer 9" and "context removed, layer 16" ON THE
%        MARKS. Those strings are what makes the geometry insoluble, and the
%        information they carry is not lost, only moved 0.5cm down into the
%        block whose job is scope.
%   (R6) ROW B'S SEVEN MARKS WERE PLOTTED AT THEIR ONE-DECIMAL ROUNDINGS.
%        \foreach ran over {3.6,3.8,4.2,5.7,12.7,14.1,18.6}. With the layer
%        numerals deleted in v3, position became the ONLY carrier of these
%        seven numbers, and the rounding distorted the one property row B still
%        exists to show: the L16-to-L12 gap is 0.2412 in the data and was drawn
%        0.200, the L12-to-L9 gap is 0.3387 and was drawn 0.400, and the two
%        leftmost discs were left 0.63pt (0.22mm) of white between their edges,
%        which fuses at print size and makes the row show six marks where its
%        own scope line says seven. The marks are now plotted at the source
%        ratios, workspace_probe.json :: layers.<L>.kl.pure_drug[0] /
%        layers.<L>.kl.random_matched[0], written to six decimals because
%        pgfmath is fixed-point (its resolution here is about 6e-7 cm, four
%        orders below the imagesetter). No printed numeral exists on this row
%        to contradict, and the drawn positions move by at most 0.72pt. The
%        white gap between the two leftmost discs becomes 1.34pt (0.47mm).
%        ROWS A, C AND D ARE DELIBERATELY NOT CHANGED WITH IT, and the reason
%        is the same reason row B is. Their marks are plotted at the three- and
%        four-decimal values the thesis PRINTS -- 0.821 and 0.883 in
%        tab:workspace, +0.0197 / +0.0118 / +0.0085 / -0.0012 in tab:identity,
%        +0.014 / -0.001 / -0.010 in tab:epsilon -- and every one of those
%        marks carries a name a reader can look up in the table that prints it,
%        so drawn position and printed value agree by construction. Row B's
%        marks carry no name and no numeral, so position is all a reader has.
%        The displacements are also of different orders: row B's rounding moved
%        a mark up to 0.72pt, rows A, C and D's move one at most 0.11pt, which
%        is below what an imagesetter resolves. Recovered from the v4 render by
%        inverting each row's map: row A 0.82101 / 0.88301, row C 0.019702 /
%        0.011797 / 0.008501 / -0.001204, row D 0.014001 / -0.001014 /
%        -0.010001, row B 3.5778 / 3.8189 / 4.1577 / 5.7248 / 12.6835 / 14.135
%        / 18.6343 against source ratios 3.577965 / 3.819122 / 4.157847 /
%        5.725077 / 12.683429 / 14.134973 / 18.634044.
%   (R7) THE ROW COUNT WAS NOT RECOVERABLE. Row B's scope line said "one mark
%        per layer" and never said how many layers. It now says "seven layers,
%        one mark each", so a reader who sees the two leftmost marks touch can
%        still count the row.
%   (R8) ROW C'S ZERO RULE WAS ANONYMOUS WHILE ROW D'S IDENTICAL ZERO RULE WAS
%        NAMED. Zero does more work in C than in D: it is the null the three
%        interval bars are read against, and the row's other rule, the black
%        pre-registered margin, is 45pt from the nearest mark. The ledger's
%        stated reason for leaving it anonymous -- "the '0' tick numeral sits
%        directly beneath it" -- is falsified by row D, whose zero has a '0'
%        tick numeral directly beneath it too and was labelled anyway. Row C's
%        zero now carries "zero: no identity effect" in quietink on the row's
%        own top text baseline, exactly as row D carries "zero: no scramble
%        gap".
%   (R9) THREE TYPE SIZES ABOVE 7pt, AND THE EXCESS WAS FOUR GLYPHS OF A
%        FOREIGN FACE SITTING IN THE TICK LANE. The v3 render set 9.96pt (300
%        chars) and 7.97pt (671) in TeXGyrePagella, plus four CMSY10 glyphs
%        reported at 8.30pt: the multiplication sign of "1x" and the three
%        math minus signs of the negative tick labels in rows C and D. The
%        multiplication sign is now \texttimes, which sets in Pagella at the
%        surrounding 7.97pt. The minus signs are the negative-tick idiom the
%        brief lists as keep-verbatim; they are now \textminus, which is the
%        TS1 minus of the text font -- a real minus, not the ASCII hyphen
%        \foreach would print, and no longer a second face. All figure type is
%        now Pagella at two sizes.
%  (R10) THE INTER-ROW RHYTHM WAS IRREGULAR, AND THE PAIR MOST AT RISK OF BEING
%        READ ACROSS WAS NOT THE MOST SEPARATED. Measured gaps between one
%        row's last ink and the next row's first were 0.42 / 0.54 / 0.50cm.
%        Rows C and D are the pair that must not be compared -- both are
%        dot-and-bar stacks against a dotted zero, both step their ticks by
%        0.02, and five of their six tick numerals are the same strings -- yet
%        C-to-D was not the widest band. Row B moves from y = -2.10 to -2.14
%        and row C from -5.70 to -5.66, which makes the three gaps 0.44 / 0.44
%        / 0.52. Row D and the foot block do not move, so the drawn height is
%        unchanged. Every mark in rows B and C moves with its own rule; no
%        value moves relative to its own axis.
%  (R11) THE THREE SEED LABELS NAMED THE ARM THREE TIMES. "residual, seed 42 /
%        43 / 45" set "residual" three times, 51pt of label, while line two of
%        the same row's axis block now opens "residual seeds at layer 12". The
%        marks are cut to "seed 42", "seed 43", "seed 45": the arm is named
%        once, in axis-label position, which is the same rule under which row
%        B's 1x rule carries its coordinate and not its name. It also widens
%        the visible difference between the three seed marks and the fourth,
%        which is a different arm. tab:identity's full strings are in the
%        ledger and in the caption.
%
% OVERRULED, WITH REASONS.
%   (O1) "DRAW ONLY ONE MARK IN ROW A; DELETE 0.883." Refused. The cross-figure
%        ownership rule assigns 0.8208333333333334 and 0.8833333333333332 to
%        fig:mechanism (a), which draws them as series, AND to this row as its
%        headline "with both layers named". A row that answers "is it there?"
%        with one accuracy loses the fact that the section's headline is a
%        PAIR. The before/after hazard the finding names is real and is met by
%        three things instead: the two series are drawn with different glyphs
%        (fig:mechanism (a)'s own open/filled encoding), their two depths are
%        named on the drawing in axis-label position, and the caption states
%        that they cross between layers 9 and 12 so context removal RAISES
%        accuracy at the two deepest probes.
%   (O2) "SET ROW A'S TWO LABELS ON TWO BASELINES." Refused: the label rule is
%        "within a row all direct labels share one baseline", and a second
%        baseline would cost 0.36cm of a figure that has 0.4mm of headroom
%        under its 118mm cap. The same defect is repaired at R5 without a
%        second line.
%   (O3) "MOVE THE FOUR PANEL HEADS ONTO THEIR OWN LINES AT x = 0 AND DELETE
%        THE STUB COLUMN." Refused. It is the better composition and it is what
%        fig:mechanism does, but four extra head lines cost about 1.4cm and
%        this figure has 0.4mm. The rebuild brief specifies a right-aligned
%        stub at one x for this figure, and the four question strings it fixes
%        are what set the column at 4.97cm rather than the 3.4cm the brief
%        estimated.
%   (O4) "LEFT-ALIGN THE FOUR HEADS SO THE ROW LETTERS SIT AT ONE x." Refused.
%        It would put (a)-(d) at one x, which is real, but it would also leave
%        a ragged right edge with a 0-to-24mm gap between each head and the
%        rule it heads, so the head would float away from its own row. A side
%        head belongs against the data it names. The brief asks for right-
%        aligned at one x and that is what is drawn; all four heads end at
%        x = 187.33pt, spread 0.00pt.
%   (O5) "RENAME ROW D'S FIRST MARK 'single-cell target' TO DISTINGUISH IT FROM
%        ROW C'S 'single-cell arm'." Refused, and the finding has it backwards.
%        They are the SAME object: 04b-representation.tex:939 writes "the
%        single-cell arm's scramble gap is +0.014", which is row D's own value.
%        One object may not carry two names. Row D keeps tab:epsilon's wording
%        and the caption states in words that the same arm appears in both rows
%        as two different quantities on two different rules, which is the part
%        of the risk that only words can close.
%   (O6) "NAME THE TRAINING ARM ON ROWS A AND B, BECAUSE ROW C IS A DIFFERENT
%        FINE-TUNE." The fact is right and is worth saying: workspace_probe
%        .json :: config.model_path ends pythia_sft_endcell/final for rows A
%        and B, probe_arm_residual.json ends pythia_sft_residual/final for row
%        C's three seeds. But probe_arm_single_cell.json -- row C's fourth mark
%        -- is back on pythia_sft_endcell/final, so "the arm" is not a property
%        of a row and a one-word cue on the axis would mislead. It goes in the
%        caption, where it can be said in a sentence.
%   (O7) "DISCLOSE ON THE DRAWING THAT ROW A DRAWS EACH SERIES AT ITS OWN BEST
%        OF SEVEN DEPTHS." True -- 0.8208333333333334 is the maximum of the raw
%        series and 0.8833333333333332 the maximum of the context-removed one
%        -- but the document's convention for a favourable-layer choice is the
%        caption, which is where fig:hook discloses layer 4 and fig:materiality
%        discloses layer 12. It is in the proposed caption, together with row
%        C's own layer-12 disclosure.
%
% V3 CLAIMS THIS PASS FALSIFIES, superseded here in place and left standing
% above, because the v2->v3 audit trail refers to them.
%   * v3 (iii), last sentence: "Both threshold rules and both zero rules now
%     start at the tick depth (-0.10) and rise." FALSE of v4. All five
%     coordinate and threshold rules now start AT THE AXIS RULE. Starting them
%     at the tick depth was itself the defect: it put a 0.90pt black stroke and
%     two 0.50pt dotted strokes inside a lane that carries nothing but 0.30pt
%     rulegrey ticks. See R4.
%   * v3 (iv): the row-A mark labels are no longer "raw probe, layer 9" and
%     "context removed, layer 16". The layers are still on the drawing, in the
%     axis block. The claim that both labels sit on ONE baseline survives and
%     is still true; what fails is the claim that the leadering worked. See R5.
%   * v3 (vi): row C's series are no longer lettered with tab:identity's full
%     strings "residual, seed 42 / 43 / 45" but with "seed 42 / 43 / 45", the
%     arm having moved into axis-label position. See R11.
%   * v3 (vii): "naming the construction for every row that has bars and
%     stating that the other two return none." FALSE on two counts: row B's
%     instrument does return a dispersion, and one mark inside row C returns
%     none while the other three do not. See R1.
%   * the STROKE WEIGHTS paragraph above: "0.30pt furniture (27 objects: tick
%     marks and three leaders)". v4 measures 26 ticks and two mark leaders, one
%     of them an elbow drawn as two segments; row C's third leader is deleted.
%     Every other weight and count in that paragraph still holds, including the
%     justification of the single 0.80pt stroke.
%   * the v3 ledger's width_measured.v3 block recorded 4.79 characters per
%     100mm^2 and called the four 8.30pt CMSY10 glyphs "not a third type size".
%     Re-measured before this pass began: 4.81, and they were a third size.
%     Both are corrected in comparators.json, width_measured.
%   * header line 8, "no mark moved, no number changed", was already superseded
%     by v3 and stays superseded: v4 moves marks again, in rows B and C, and
%     changes no value. Every old and new position is in comparators.json.
%
% TEXT BUDGET, since this pass both adds and cuts. v3 set 978 characters in
% 20312 square mm, 4.81 per 100; the ceiling is 5.0 and the ink ceiling is 6.0%
% at 200dpi. Cuts (row A's two mark labels -19, its scope line -1, row C's
% three seed labels -30, row D's scope line -15, which stopped repeating the
% "40 cell lines" the foot block already prints) very nearly pay for the adds
% (row B's layer count +9, row C's scope line +18 and its zero label +24, the
% foot block +37). Measured after: see comparators.json, width_measured.v4.
%
% ===========================================================================
% v5 -- 2026-08-19. VERIFICATION PASS on v4, against the same three adversarial
% reads. All twenty-one confirmed defects were re-checked ON THE RENDER, not on
% the source, and all twenty-one hold repaired; the seven overrules were
% re-argued and all seven stand. ONE defect that no refuter raised was found
% and repaired, and three claims made in the v3 and v4 headers are superseded
% because re-measurement does not support them. NO VALUE CHANGED. One rule
% endpoint moved; no mark moved.
%
% REPAIRED.
%   (R12) ROW C'S ZERO RULE WAS THE ONLY LABELLED RULE ON THE FIGURE THAT DID
%         NOT REACH ITS OWN LABEL. Every labelled rule here tops out 0.04cm
%         below its label's baseline: chance (axis+0.28, label axis+0.32), row
%         B's 1x rule (-1.86 / -1.82), row C's pre-registered margin (-4.00 /
%         -3.96), row D's zero (-7.60 / -7.56). Row C's zero rule stopped at
%         -4.16 against a label baseline of -3.96, a 0.20cm gap -- five times
%         every other, and on the one rule that three interval bars are read
%         against. Measured on the v4 render: the margin rule's top sat at
%         y = 191.99pt with its label's box bottom at 192.45, i.e. touching,
%         while the zero rule's top sat at 196.53 under a label on the SAME
%         baseline, 4.5pt clear. It now stops at -4.00 like the other four.
%         It crosses nothing: the nearest mark, seed 42's interval bar, begins
%         at +0.0063, which is 15.4pt to the right of zero on this row's map.
%         Height is unchanged: the row's topmost ink is its label, not its rule.
%
% SUPERSEDED, EXPLICITLY, and left standing above.
%   * v4's R9 and the v4 STROKE WEIGHTS supersession say "v4 measures 26 ticks
%     and two mark leaders", which reads as 28 objects. The figure draws 24
%     rulegrey ticks (5 + 5 + 8 + 6) and two accent mark-leader PATHS, one of
%     them the elbow, which is 26 objects at 0.30pt in total. Re-measured off
%     the v5 render with the path list, not by eye: 0.30pt x 26, 0.40pt x 4,
%     0.50pt x 4, 0.80pt x 1, 0.90pt x 2, 1.00pt x 6. Every other figure in
%     that paragraph holds, including the justification of the single 0.80pt.
%   * v4's R10 says the row rhythm becomes "0.44 / 0.44 / 0.52" cm. Those are
%     computed from the row anchors, not measured off the render. The ANCHOR
%     pitch is 0.95 / 0.95 / 1.03cm and is what the coordinates set; the INK
%     gap, measured on the v5 render from one row's lowest text box to the
%     next row's highest, is 11.43 / 11.43 / 12.39pt = 0.40 / 0.40 / 0.44cm
%     (0.48cm for the third if row C's outlying PazoMath alpha box, which hangs
%     1.3pt below the rest of that line, is excluded). R10's substance is
%     unaffected and is the only part that matters: two equal gaps and a wider
%     one, with the widest band between rows C and D, the pair a reader must
%     not read across.
%   * the v3 header, note (iii), quotes row B's axis-label line one as
%     "Ablation cost, in multiples of the matched-dimension random slab". All
%     four axis-label line-ones were lowercased in v4 and the v4 header did not
%     record it, though comparators.json did, per row, as
%     axis_label_line_1_v3_superseded. They now read "probe accuracy",
%     "ablation cost, ...", "identity effect, ..." and "scramble gap $\gap$:
%     ...". The reason is fig:mechanism, which is this document's reference
%     standard for an axis label and which sets the same class of string --
%     including the same quantity, "probe accuracy" -- in lower case at 10.0pt
%     five times over. The four panel HEADS stay capitalised: they are
%     interrogative sentences, not descriptive labels, and lower-casing a
%     question would read as a fragment.
%
% OVERRULED IN v4 WITHOUT BEING WRITTEN DOWN, recorded here so they do not come
% back.
%   (O8) "RAISE ROW B'S 1x RULE FROM -1.82 TO ABOUT -1.74 SO IT READS AS A RULE
%        RATHER THAN A STUB." Refused. It is the faintest object on the figure
%        and that is correct: it is a definitional coordinate, the unit of its
%        own axis, and the dotted register exists for a coordinate that decides
%        nothing. Height is not the reason; consistency is. v4 gave it exactly
%        row A's chance geometry -- rule from the axis to axis+0.28cm, label
%        baseline at axis+0.32cm -- so that a rule's own label sits at ONE
%        offset everywhere on the figure, and R12 above has just brought the
%        fifth rule into that same offset. Lengthening one of the five to make
%        it louder would break the thing R12 fixed. At 0.50pt densely dotted
%        over 7.94pt it prints about five dots, which is legible at 200dpi and
%        at print size; checked on both.
%   (O9) "RESTORE tab:identity's FULL STRING 'single-cell arm, six depths' AS
%        THE MARK LABEL IN ROW C." Refused, and the defect it names is already
%        repaired the other way round: "six depths" is in axis-label position,
%        on line two of row C's own axis block, which is where this figure puts
%        every scope fact (row A's two layers, row B's layer count, row C's
%        rung, row D's stratum). The mark keeps the object's name and the axis
%        block carries the statistic, which is the same division of labour that
%        lets the three seed marks read "seed 42/43/45" rather than repeating
%        "residual" three times. Restoring the full string would set the widest
%        label in the row for a fact already stated 0.5cm below it.
%
% RE-MEASURED ON THE v5 RENDER, all inside their ceilings: content box
% 172.72mm wide (170-174 required) by 117.60mm tall (cap 118); ink 5.84% at
% 200dpi (ceiling 6.0); 1000 characters in 20312 square mm, 4.92 per 100
% (ceiling 5.0); two type sizes above 7pt, 9.96pt for the four axis-label
% line-ones and the four panel heads and 7.97pt for everything else, plus a
% 7.57pt three-character subscript in $\gap$; six faces, all Pagella or Palladio
% and no sans, no CMSY, no typewriter; italic 10 characters of 1000, of which 5
% are the emphasised word "which" and the rest are the math variables $\alpha$
% and $n$. The four axis rules run x = 198.87 to 541.86pt, congruent to 0.00pt,
% and the first and last drawn tick of every row sits on those endpoints.
% One page, zero Overfull and zero Underfull boxes.
% ===========================================================================
%
% ===========================================================================
% v-BLOCK, 2026-08-19: WHOLE-SLATE PASS (fig:comparators, fig:licensing,
% fig:purify, fig:hook, fig:materiality reconciled against one another and
% against fig:mechanism, which is the document's reference standard and was
% NOT rebuilt). Nothing above this line is deleted; where an earlier claim is
% now false it is superseded EXPLICITLY below.
%
% WHY. The five figures were rebuilt independently and drifted. This pass owns
% only the differences BETWEEN them. No value changed in any of the five; the
% sibling .json ledgers carry the value-to-position maps, old and new.
%
% THE SLATE RULES THAT NOW HOLD ACROSS ALL FIVE (each stated once here, in
% every file, so any future single-figure edit can see what it would break):
%
%  R1 TICK NUMERALS are quietink, \scriptsize, inner sep 0, anchored north on
%     the tick and set 0.16cm below the axis rule, over a 0.10cm rulegrey
%     0.30pt tick. Before: comparators/licensing/purify quietink, hook and
%     materiality ink; ticks 0.10 / 0.12 / 0.14cm; three different gaps.
%  R2 A THRESHOLD'S OWN LABEL IS BLACK, matching the black solid 0.90pt rule
%     it names (S6). Before: comparators black, purify's "gate, 0.8" and
%     materiality's two margin labels ink.
%  R3 ONE NAME PER OBJECT (S19). The pre-registered 10% rule is "10%
%     materiality margin" in fig:comparators row (c), fig:materiality (a) and
%     (b), and in figs/family.py, which is where the name comes from.
%     fig:comparators' "pre-registered margin" is superseded; "pre-registered"
%     is now the caption's word, not the drawing's. Zero, where it is the null
%     a bar is read against, is "zero: no identity effect" in BOTH
%     fig:comparators row (c) and fig:materiality (b) -- one quantity, one
%     wording. V_pure is spelt with \textrm everywhere, as figs/slab.tex does.
%  R4 COLOUR KEY, one key for five figures (S8). accent = the drug-side object
%     under test. quietink = every comparator, null, control, replication or
%     discounted arm, AND its label, AND its leader. ink/black = thresholds,
%     axis rules, axis labels, panel heads, definitional prose. rulegrey =
%     furniture only. Within quietink the GLYPH separates two kinds: an OPEN
%     disc is a construction built to carry no signal (a null, a control, a
%     raw/before series); a FILLED disc is a real measurement that is not the
%     headline. fig:hook (d) drew both its null and its cell-line control in
%     ink and is corrected.
%  R5 PROJECTION AND ZERO RULES are quietink 0.50pt densely dotted (S7).
%     fig:hook's two panel-(c) projection lines were rulegrey 0.30pt densely
%     dashed and are corrected.
%  R6 DEFINITIONAL PROSE is \footnotesize ink and OUTRANKS every label (S4);
%     at most one block per figure, on the panel's own x origin. fig:hook's
%     one prose line was \scriptsize and is raised. fig:purify's panel-(b)
%     block carried an argument clause ("so causal claims use the purified
%     slab, not the raw") that its caption already carries, and is now
%     definitional only.
%  R7 ONE INTERVAL DECLARATION per figure, \scriptsize quietink, at the foot,
%     on the figure's x origin, opening with the word "Intervals:" and naming
%     EVERY panel including the ones with no measured mark (S17).
%     fig:comparators already did this and is the model; the other four now
%     match its form.
%  R8 THE FIVE-FIELD PROMPT STRIP is ONE object drawn in two figures, and the
%     two drawings are now congruent: cell boundaries 0 / 1.12 / 1.90 / 2.64 /
%     4.23 / 6.90 cm and cell height 0.42cm in both fig:licensing (a) and
%     fig:hook (a) and (b), labels centred in their cells, the drug cell
%     accent-bordered at 0.40pt on accent!12 and its label NOT bold. Before:
%     licensing 6.90cm wide with 0.62cm cells and a bold label, hook 7.44cm
%     wide with 0.42cm cells and a plain label -- the same five fields at two
%     sizes, four pages apart, with each caption pointing at the other.
%  R9 MEASURE. Every fullwidth figure draws 172.0-174.0mm; the one textwidth
%     figure draws 140.3mm. Ink at 200dpi inside the content bbox, counted as
%     greyscale luminance below 230 (the contract's own method), is now
%     5.63-5.99% on all five, under the 6.0% body ceiling, where before this
%     pass it ran 5.36-6.54%.
%
% WHAT THIS PASS DELIBERATELY DID NOT UNIFY, so that a later reader does not
% think it was missed:
% 2026-08-19 PANEL-HEAD (c) REWORDED, at the supervisor's direction.
%   Was "(c) Is it \emph{which} drug?", which is not a grammatical question:
%   the interrogative "which" cannot sit inside "Is it ...?" as a bare
%   modifier, so the head read as a fragment beside the three well-formed
%   heads it is parallel to. Now "(c) Does it matter which drug?", which is
%   grammatical, keeps the row's actual question -- the region is used (row b),
%   but does its CONTENT, the identity of the drug named, change the model's
%   behaviour. Set as "(c) Is identity read?": parallel to the short forms of
%   (a) and (b), it is the title of sec:res-mechanism, and it is the noun the
%   row's own axis label already uses ("identity effect, as a fraction of...").
%   "Does it matter which drug?" was tried first and made this stub the widest
%   in the figure, overfull by 5.08pt against a measure with 0.4mm spare.
%   The prose at 04b sec:methods-probe still asks "does it use which drug?";
%   that is a clause inside a sentence, where the construction is fine.
%   No mark moved and no number changed.
%  * PANEL-HEAD CASE. fig:comparators sets "(a) Is it there?" and the other
%    four set "(a) where the drug name enters". comparators' heads are
%    interrogative SENTENCES and sentence case is correct for a sentence; the
%    v4 header argued this and the argument stands. fig:mechanism's lowercase
%    heads are noun phrases, not questions. One capital letter, four times.
%  * PANEL-HEAD POSITION. fig:comparators puts its heads in a right-aligned
%    left stub, the other four above the panel at its x origin. Moving them
%    above the rules costs about 20mm of height against a 118mm cap that the
%    figure already meets at 117.60mm. Not available.
%  * TWO-COLUMN GUTTER. purify and hook split at x = 8.40cm, licensing at
%    7.75cm. licensing cannot move right: "(d) what a negative would buy" is
%    4.85cm wide and its origin is already at 12.40 of a 17.30cm measure, with
%    1.5pt of slack. Moving purify and hook LEFT to 7.75 would be arbitrary.
% ===========================================================================
%
% CHANGED IN THIS FILE. (i) Row (c)'s margin label "pre-registered margin" ->
% "10% materiality margin" (R3). (ii) The first and last tick numeral of all
% four rows are now anchored north-west / north-east ON the rule's own
% endpoints instead of being centred on them, so a row's leftmost text x IS
% the axis origin and no numeral overhangs the measure (S13, S14); the six
% interior numerals are unchanged. (iii) \RL, the one rule length all four
% rows share, goes from 12.10cm to 12.35cm. Pulling the end numerals in had
% taken the content box to 170.43mm, 0.43mm inside S14's floor; 12.35cm puts
% it at 172.97mm, in line with licensing 173.10, purify 172.09 and hook
% 173.99. This RESCALES all four value-to-position maps by 12.35/12.10 and
% CHANGES NO VALUE: every mark still sits at its source coordinate, and both
% the old and the new maps are recorded in comparators.json. The absolute
% label shims (+0.30, +0.22, +0.10, -0.42, -0.38, -0.08) are print-unit gaps
% and are deliberately NOT rescaled.
% SUPERSEDED HERE: line 8's "What changed is labelling only; no mark moved, no
% number changed" was already false of v2-v5 and is now false again -- every
% mark's POSITION moved by 2.07% when \RL grew. No mark's VALUE has ever
% moved. Read line 8 as applying to v1->v2 only.
% NOT CHANGED, and it is the model for the other four: this figure's foot
% block already opened with "Intervals:" and already named all four rows, and
% its stroke set (0.30 furniture / 0.40 axis / 0.50 dotted coordinate / 0.90
% threshold / 1.00 interval) is the set the slate now uses.
% RE-MEASURED: 172.97 x 117.60mm; ink 5.822%; 983 characters, 4.83 per
% 100mm^2; two type sizes above 7pt (9.96 / 7.97) plus a 7.57pt subscript in
% \gap. One page, zero Overfull, zero Underfull.
% ===========================================================================
\begin{tikzpicture}[x=1cm, y=1cm, font=\footnotesize]
  % ---- Four deliberately different value-to-position maps, one per rule. ---
  % ONE rule length for all four rows, declared once, so the four rules stay
  % congruent to the last decimal; v3 raises it from v1/v2's 10.5cm. Rows A and
  % B keep their domains; rows C and D are re-domained so that the first and
  % last DRAWN TICK of every row sits exactly on the rule's own endpoints.
  \def\RL{12.35}
  \def\xa#1{{(#1)*\RL}}                        % row A, domain  0    .. 1
  \def\xb#1{{(#1)/20*\RL}}                     % row B, domain  0    .. 20
  \def\xc#1{{((#1)+0.02)/0.14*\RL}}            % row C, domain -0.02 .. 0.12
  \def\xd#1{{((#1)+0.04)/0.10*\RL}}            % row D, domain -0.04 .. 0.06
  % Same maps shifted by a label gap; the shift MUST sit inside the braces --
  % \xc{v}+0.30 is not a pgfmath expression and does not compile. 0.30cm, not
  % v1's 0.13cm: at 0.13cm a label abutted the end of its own interval bar and
  % read as part of it. Every node carries inner sep=0, so from v3 these shims
  % are the gaps the reader actually sees, to the hundredth of a centimetre.
  \def\xcR#1{{((#1)+0.02)/0.14*\RL+0.30}}
  \def\xdR#1{{((#1)+0.04)/0.10*\RL+0.30}}
  % label shims. \xcM: v3 set row C's single-cell label 0.35cm from the mark's
  % centre, i.e. 0.30cm from its edge, which is over the 3mm threshold at which
  % a label takes a leader -- and that leader, horizontal, quietink, at the
  % mark's own y, read as a one-sided interval bar. v4 closes the gap to 0.22cm
  % from the centre (0.17cm, 4.8pt, from the edge), under the threshold, and
  % DELETES the leader. \xcW and \xdL set a zero rule's own label clear of the
  % rule; \xaL and \xbL do the same above; \xcL right-aligns row C's margin
  % label clear of the margin rule.
  \def\xcM#1{{((#1)+0.02)/0.14*\RL+0.22}}
  \def\xcW#1{{((#1)+0.02)/0.14*\RL+0.10}}
  \def\xaL#1{{(#1)*\RL+0.10}}
  % \xaB sets row A's raw-probe label back from its own mark so the two row-A
  % labels do not close up into one run of text; \xaC starts its elbow leader
  % 0.04cm clear of that label.
  \def\xaB#1{{(#1)*\RL-0.42}}
  \def\xaC#1{{(#1)*\RL-0.38}}
  \def\xbL#1{{(#1)/20*\RL+0.10}}
  \def\xcL#1{{((#1)+0.02)/0.14*\RL-0.08}}
  \def\xdL#1{{((#1)+0.04)/0.10*\RL+0.10}}
  % Four node styles and no others, so the type ranking is enforced by the
  % style sheet and not by hand: qty = axis-label line one, the largest type on
  % the figure; scope = its one scope line, and the foot; stub = the panel
  % head; lab = every direct label and every rule's own label; tick = numerals.
  \tikzset{lab/.style={font=\scriptsize, inner sep=0},
           tick/.style={anchor=north, font=\scriptsize, text=quietink, inner sep=0},
           qty/.style={anchor=north west, font=\footnotesize, text=ink, inner sep=0},
           scope/.style={anchor=north west, font=\scriptsize, text=quietink, inner sep=0},
           stub/.style={anchor=east, font=\footnotesize, text=ink, inner sep=0}}

  % =====================================================================
  % ROW A -- is it there?  probe accuracy, 0 to 1
  % =====================================================================
  \draw[ink, line width=0.4pt] (0,0) -- (\RL,0);
  \foreach \v in {0.25,0.5,0.75}{
    \draw[rulegrey, line width=0.3pt] (\xa{\v},0) -- (\xa{\v},-0.10);
    \node[tick] at (\xa{\v},-0.16) {\v};}
  \draw[rulegrey, line width=0.3pt] (\xa{0},0) -- (\xa{0},-0.10);
  \draw[rulegrey, line width=0.3pt] (\xa{1},0) -- (\xa{1},-0.10);
  \node[tick, anchor=north west] at (\xa{0},-0.16) {0};
  \node[tick, anchor=north east] at (\xa{1},-0.16) {1};
  % chance is a threshold, so black, solid, 0.90pt, labelled roman black. v4
  % stops it AT THE AXIS RULE: in v3 it ran to the tick depth -0.10, so it
  % printed a sixth, black, unnumbered downstroke inside a five-tick lane.
  \draw[black, line width=0.9pt] (\xa{0.083},0) -- (\xa{0.083},0.28);
  \node[lab, anchor=base west, text=black] at (\xaL{0.083},0.32) {chance, $1/12$};
  % raw probe: OPEN disc, which is fig:mechanism (a)'s own encoding -- it draws
  % this series open and the context-removed series filled -- so the two marks
  % are separable by glyph and not by colour alone at 1.4pt.
  % v4: the label is right-aligned ON its own mark's x and drops a leader onto
  % the disc. In v3 it sat 15mm left of the mark and reached it by a diagonal
  % that ran beneath the OTHER label, whose span covered both discs.
  % The label is set back 0.42cm from its own mark and reaches it by an elbow:
  % a horizontal run in clear space, then the same vertical drop the other mark
  % gets. Right-aligned ON the mark, the two labels closed to 6.4pt and read as
  % one run of text, "raw probe context removed"; the set-back opens 18.3pt and
  % the whole elbow lies LEFT of the other label, so no leader passes under any
  % text.
  \draw[accent, line width=0.8pt, fill=white] (\xa{0.821},0) circle (1.4pt);
  \node[lab, anchor=base east, text=accent] at (\xaB{0.821},0.32) {raw probe};
  \draw[accent, line width=0.3pt] (\xaC{0.821},0.18) -- (\xa{0.821},0.18) -- (\xa{0.821},0.05);
  % context removed: FILLED disc. Its label is right-aligned to the rule's
  % right end so nothing crosses the measure; at 55pt it now stands over its
  % own mark and over no other, and takes the identical drop leader.
  \fill[accent] (\xa{0.883},0) circle (1.4pt);
  \node[lab, anchor=base east, text=accent] at (\RL,0.32) {context removed};
  \draw[accent, line width=0.3pt] (\xa{0.883},0.25) -- (\xa{0.883},0.05);
  % axis-label block, in axis-label position: under the numerals, at x = 0.
  % Line two now carries the two depths, tying each series to its own, which is
  % where fig:hook and fig:materiality disclose theirs.
  \node[qty]   at (0,-0.50) {probe accuracy};
  \node[scope] at (0,-0.87) {twelve-way drug decode; raw at layer 9, context removed at layer 16};
  \node[stub, anchor=base east] at (-0.40,0.32) {(a) Is it there?};

  % =====================================================================
  % ROW B -- is it used?  ablation cost in multiples of the null, 0 to 20
  % =====================================================================
  % v4: the rule drops from y = -2.10 to -2.14, which evens the inter-row gaps
  % to 0.44 / 0.44 / 0.52cm and gives the widest band to the C/D pair, the one
  % pair a reader must not read across. Nothing moves relative to this rule.
  \draw[ink, line width=0.4pt] (0,-2.14) -- (\RL,-2.14);
  \foreach \v in {5,10,15}{
    \draw[rulegrey, line width=0.3pt] (\xb{\v},-2.14) -- (\xb{\v},-2.24);
    \node[tick] at (\xb{\v},-2.30) {\v};}
  \draw[rulegrey, line width=0.3pt] (\xb{0},-2.14) -- (\xb{0},-2.24);
  \draw[rulegrey, line width=0.3pt] (\xb{20},-2.14) -- (\xb{20},-2.24);
  \node[tick, anchor=north west] at (\xb{0},-2.30) {0};
  \node[tick, anchor=north east] at (\xb{20},-2.30) {20};
  % The null at 1x is not a threshold: it is the UNIT OF THIS AXIS by
  % construction, a definitional comparator, so it keeps the dotted register.
  % It carries its COORDINATE and not its name, because the axis-label line
  % below already names the object in axis-label position, which is where S5
  % requires the unit to be named; setting "matched-dimension random slab"
  % here as well would print one four-word string twice inside one row, 3.5cm
  % apart. fig:mechanism (b)'s name for the object is the one used, once.
  % v4: stops at the axis rule, and rises to axis+0.28 with its label at
  % axis+0.32, the same geometry as row A's chance rule, so a rule's own label
  % sits at one offset everywhere on the figure. \texttimes, not $\times$: the
  % math sign was the figure's only CMSY10 glyph and printed at a third type
  % size, 8.30pt, inside the tick lane.
  \draw[quietink, line width=0.5pt, densely dotted] (\xb{1},-2.14) -- (\xb{1},-1.86);
  \node[lab, anchor=base west, text=quietink] at (\xbL{1},-1.82)
    {1\texttimes{}, by construction};
  % Seven point estimates, one per layer. The layer identities are NOT drawn:
  % see the v3 header, note (1) superseded. v4 plots them at the SOURCE ratios
  % (workspace_probe.json :: layers.<L>.kl.pure_drug[0]/random_matched[0], to
  % six decimals, which is four orders finer than the imagesetter) rather than
  % at their one-decimal roundings: with the numerals gone, position is the only
  % carrier left, and the rounding had pushed the two leftmost discs to 0.22mm
  % of white. Order below is L16, L12, L9, L6, L2, L8, L4.
  \foreach \v in {3.577965,3.819122,4.157847,5.725077,12.683429,14.134973,18.634044}{
    \fill[accent] (\xb{\v},-2.14) circle (1.4pt);}
  \node[qty]   at (0,-2.64) {ablation cost, in multiples of the matched-dimension random slab};
  \node[scope] at (0,-3.01) {KL from the clean forward pass; seven layers, one mark each};
  \node[stub, anchor=base east] at (-0.40,-1.82) {(b) Is it used?};

  % =====================================================================
  % ROW C -- is it WHICH drug?  fraction of the model's own gap
  % =====================================================================
  % v4: the rule rises from y = -5.70 to -5.66 (see row B). Everything in the
  % row moves with it; no value moves relative to the rule.
  \draw[ink, line width=0.4pt] (0,-5.66) -- (\RL,-5.66);
  % The label is carried separately from the coordinate so the negative tick
  % sets a real minus and not the ASCII hyphen \foreach would print. Same idiom
  % as gate.tex and power.tex; positives stay unmathed so their digits match
  % rows A and B, which are set in text. v4 sets the minus as \textminus, the
  % text font's own TS1 minus, rather than the math $-$: it is still a real
  % minus, and it is no longer the only non-Pagella glyph on the figure.
  \foreach \v/\t in {0/0, 0.02/0.02, 0.04/0.04, 0.06/0.06, 0.08/0.08, 0.10/0.10}{
    \draw[rulegrey, line width=0.3pt] (\xc{\v},-5.66) -- (\xc{\v},-5.76);
    \node[tick] at (\xc{\v},-5.82) {\t};}
  \draw[rulegrey, line width=0.3pt] (\xc{-0.02},-5.66) -- (\xc{-0.02},-5.76);
  \draw[rulegrey, line width=0.3pt] (\xc{0.12},-5.66) -- (\xc{0.12},-5.76);
  \node[tick, anchor=north west] at (\xc{-0.02},-5.82) {\textminus0.02};
  \node[tick, anchor=north east] at (\xc{0.12},-5.82) {0.12};
  % The zero rule is drawn in two pieces. The single-cell median sits at
  % -0.0012, one tenth of a centimetre from zero on this scale, so the rule is
  % interrupted across the marker's height rather than run through it; from v4
  % the interruption also holds the gap between that mark and its own label.
  % Both pieces stop AT THE AXIS RULE: in v3 the lower piece ran to the tick
  % depth, so the "0" tick printed dotted while every other tick in the same
  % eight-tick ladder printed solid.
  % v5: the upper piece's top moves from -4.16 to -4.00. Every other labelled
  % rule on the figure -- chance, row B's 1x, row C's margin, row D's zero --
  % tops out 0.04cm below its own label's baseline; this one stopped 0.20cm
  % below, so the one rule that is read against three interval bars was also
  % the one rule that did not reach its own name. It crosses nothing: the
  % nearest mark, seed 42's bar, begins at +0.0063 and 15.4pt to the right.
  \draw[quietink, line width=0.5pt, densely dotted] (\xc{0},-5.66) -- (\xc{0},-5.55);
  \draw[quietink, line width=0.5pt, densely dotted] (\xc{0},-5.27) -- (\xc{0},-4.00);
  % v4 NAMES ZERO. Row D's identical zero rule was labelled and this one was
  % not, and zero does more work here: it is the null the three bars are read
  % against. Quietink, like the rule, and on the row's top text baseline, which
  % is exactly row D's treatment.
  \node[lab, anchor=base west, text=quietink] at (\xcW{0},-3.96) {zero: no identity effect};
  % the pre-registered margin is a threshold: black, solid, 0.90pt, the same
  % stroke fig:family already gives this same 10% margin. It too stops at the
  % axis rule; in v3 it overprinted the 0.10 tick in black at three times the
  % tick's weight.
  \draw[black, line width=0.9pt] (\xc{0.10},-5.66) -- (\xc{0.10},-4.00);
  \node[lab, anchor=base east, text=black] at (\xcL{0.10},-3.96) {10\% materiality margin};
  % The arm is named once, in axis-label position ("residual seeds at layer 12"
  % below), so the three marks are lettered by what distinguishes them.
  \draw[accent, line width=1.0pt] (\xc{0.0063},-4.21) -- (\xc{0.0335},-4.21);
  \fill[accent] (\xc{0.0197},-4.21) circle (1.4pt);
  \node[lab, anchor=west, text=accent] at (\xcR{0.0335},-4.21) {seed 42};
  \draw[quietink, line width=1.0pt] (\xc{0.0060},-4.61) -- (\xc{0.0173},-4.61);
  \fill[quietink] (\xc{0.0118},-4.61) circle (1.4pt);
  \node[lab, anchor=west, text=quietink] at (\xcR{0.0173},-4.61) {seed 43};
  \draw[quietink, line width=1.0pt] (\xc{0.0029},-5.01) -- (\xc{0.0143},-5.01);
  \fill[quietink] (\xc{0.0085},-5.01) circle (1.4pt);
  \node[lab, anchor=west, text=quietink] at (\xcR{0.0143},-5.01) {seed 45};
  % Solid, like every other point mark; the absent bar is what says it carries
  % no interval, and the foot block now says so in words. v4 DELETED the leader
  % that stood here: horizontal, quietink, 0.30pt, at the mark's own y and in
  % the mark's own colour, 0.40cm below three interval bars of the same
  % geometry, it read as a one-sided interval the source does not store. The
  % label instead moves inside the 3mm threshold at which a leader is required.
  \fill[quietink] (\xc{-0.0012},-5.41) circle (1.4pt);
  \node[lab, anchor=west, text=quietink] at (\xcM{-0.0012},-5.41) {single-cell arm};
  \node[qty]   at (0,-6.16) {identity effect, as a fraction of the model's own clean preference gap};
  \node[scope] at (0,-6.53) {residual seeds at layer 12, $\alpha$ = 0.5; single-cell median, six depths};
  \node[stub, anchor=base east] at (-0.40,-3.96) {(c) Is identity read?};

  % =====================================================================
  % ROW D -- does a better target help?  scramble gap in NIR
  % =====================================================================
  \draw[ink, line width=0.4pt] (0,-8.90) -- (\RL,-8.90);
  \foreach \v/\t in {-0.02/\textminus0.02, 0/0, 0.02/0.02, 0.04/0.04}{
    \draw[rulegrey, line width=0.3pt] (\xd{\v},-8.90) -- (\xd{\v},-9.00);
    \node[tick] at (\xd{\v},-9.06) {\t};}
  \draw[rulegrey, line width=0.3pt] (\xd{-0.04},-8.90) -- (\xd{-0.04},-9.00);
  \draw[rulegrey, line width=0.3pt] (\xd{0.06},-8.90) -- (\xd{0.06},-9.00);
  \node[tick, anchor=north west] at (\xd{-0.04},-9.06) {\textminus0.04};
  \node[tick, anchor=north east] at (\xd{0.06},-9.06) {0.06};
  % stops at the axis rule; in v3 it ran to the tick depth and printed the "0"
  % tick dotted inside a ladder whose other five ticks are solid.
  \draw[quietink, line width=0.5pt, densely dotted] (\xd{0},-8.90) -- (\xd{0},-7.60);
  % quietink, not ink: on this figure a rule's own label takes the rule's own
  % colour, so black label <-> black threshold (rows A and C) and quietink
  % label <-> quietink dotted coordinate (rows B and D). Zero is the null of a
  % difference, which is a comparator and not a threshold.
  \node[lab, anchor=base west, text=quietink] at (\xdL{0},-7.56) {zero: no scramble gap};
  \draw[accent, line width=1.0pt] (\xd{-0.016},-7.85) -- (\xd{0.042},-7.85);
  \fill[accent] (\xd{0.014},-7.85) circle (1.4pt);
  \node[lab, anchor=west, text=accent] at (\xdR{0.042},-7.85) {single cell};
  \draw[accent, line width=1.0pt] (\xd{-0.018},-8.25) -- (\xd{0.016},-8.25);
  \fill[accent] (\xd{-0.001},-8.25) circle (1.4pt);
  \node[lab, anchor=west, text=accent] at (\xdR{0.016},-8.25) {OT barycentric};
  % quietink, and under the one-colour code an arm drawn quietink names its
  % discount: the consensus run was stopped mid-epoch (04b:972, tab:epsilon).
  \draw[quietink, line width=1.0pt] (\xd{-0.022},-8.65) -- (\xd{0.003},-8.65);
  \fill[quietink] (\xd{-0.010},-8.65) circle (1.4pt);
  \node[lab, anchor=west, text=quietink] at (\xdR{0.003},-8.65) {consensus, stopped mid-epoch};
  \node[qty]   at (0,-9.40) {scramble gap $\gap$: the model's $\NIR$ minus its scramble's, expression frame};
  \node[scope] at (0,-9.77) {identifiable stratum, $n$ = 204 conditions};
  \node[stub, anchor=base east] at (-0.40,-7.56) {(d) Does a better target help?};

  % ---- one interval declaration, naming every row and every bare mark -----
  % v4. Three repairs. (i) Row B does NOT "return none": workspace_probe.json
  % stores every KL term as [mean, se, n], a normal-approximation standard
  % error, which this document declines to draw -- the same sentence hook.json,
  % purify.json and fig-mechanism.json use. Row A really does store a bare
  % float, so the two are now stated separately. (ii) Row C's fourth mark, the
  % single-cell median, carries no interval and the block did not say so.
  % (iii) "58 ordered pair clusters" is the anchor seed's count alone: seed 43
  % bootstraps over 193 and seed 45 over 199. The counts go to the caption and
  % the block names the CONSTRUCTION, which is what it is for.
  \node[scope] at (0,-10.55) {Intervals: (a) none stored. (b) normal-approximation standard errors only, not drawn.};
  \node[scope] at (0,-10.87) {(c) $95\%$ two-sided over ordered pair clusters; the median none. (d) $95\%$ one-way over cell lines.};

\end{tikzpicture}

\caption[Four questions, four comparators]{\textbf{Each instrument answers a
different question against a different comparator, and the four rules may not
be read across.} The four rules share a length and both endpoints and nothing
else: each carries its own quantity, named beneath it, and its own comparator.
The zeros of~(c) and~(d) are deliberately not vertically aligned, since
aligning them would assert the commensurability this figure exists to deny.
Converting one rule into another is the inference this part exists to block.
Where an instrument returns a clustered bootstrap its mark carries a bar; where
it does not, the mark sits bare on the rule, and the foot line names the
construction row by row. \emph{(a)}~The two accuracies are at different depths
and are not a before-and-after pair: the open mark is the raw probe at
layer~$9$, the filled one the same probe at layer~$16$ once cell line, plate
and dose are projected out. The two series cross between layers~$9$ and~$12$
--- removing context \emph{raises} accuracy at the two deepest probes --- and
only \cref{fig:mechanism}(a), which has a layer axis, can show it. The floor is
analytic, not empirical: \texttt{workspace\_probe.json} stores \texttt{chance}
as exactly $1/12$ and holds no permutation null. What projecting the drug slab
out costs this probe belongs to \cref{fig:purify} and \cref{fig:mechanism}(a)
and is not drawn here. \emph{(b)}~The seven marks are the seven probed depths
in value order and are not identified; which mark is which layer is
\cref{fig:mechanism}(b), which plots both divergences against a numeric layer
axis, and the Ratio column of \cref{tab:workspace}. That table's caveat holds
here too: the null is matched to the slab's rank and not to the variance the
slab removes, the slab's subspace fraction running $4.9$ to $7.0$ times the
null's, so the ratio mixes drug content with that difference. \emph{(c)}~The
three seeds are the first three rows of \cref{tab:identity}, and the black rule
is the $10\%$ materiality margin pre-registered for calling an effect material,
drawn in the same weight and colour as in \cref{fig:family} and
\cref{fig:materiality}. Layer~$12$ is the anchor seed's maximum over its seven
depths and the only one whose interval excludes zero; a fourth replicate,
seed~$44$, stores \texttt{swap\_skipped} at layer~$12$ and is not drawn. The
single-cell mark is a median over that arm's \emph{six} measured depths ---
layer~$12$ is skipped there, its cell-line headroom of $3.88$ falling below the
$5\times$ gate --- and no interval is stored for a median, which is why it
carries none. \emph{(d)}~The comparator is zero and not chance, the quantity
being a difference against the model's own scrambled partner on the
identifiable stratum in the expression frame (\cref{tab:epsilon}); the
consensus arm is greyed because its run was stopped mid-epoch.}
\label{fig:comparators}
\end{fullwidth}
\end{figure}

% ===========================================================================
\paragraph{The probe, and why its positive is nearly free} The prompt names the
drug in plain text (\cref{sec:methods-data}), so recovering it from the
activations establishes only that the model did not discard what it was handed.
The informative outcome would have been a negative, and what the probe buys is
the depth at which one would appear.\gloss{decode probe}{a classifier trained to
name the drug from the residual stream at one layer}

% ===========================================================================
\begin{figure}[htbp]
\begin{fullwidth}
\centering
% ===========================================================================
% licensing.tex -- fig:licensing, what the probe's positive licenses and what
% the negative that never arrived would have licensed.
% FIGURE BODY ONLY. No preamble, no document environment, no float, no caption.
% \input it from inside a figure/fullwidth float in
% Sections/04b-representation.tex, at \beatsection sec:methods-probe, straight
% after the paragraph "The probe, and why its positive is nearly free"
% (04b-representation.tex:108-113). The complete float is staged, commented, at
% the foot of this file.
%
% Compiles against thesis_v2/preamble.tex and nothing else: no \includegraphics,
% no external data, no TikZ library beyond the ones preamble.tex already loads
% (this file uses arrows.meta for the two Stealth heads and
% decorations.pathreplacing for the mirrored brace under the prompt strip).
%
% WHAT IT IS. Two panels. LEFT: where the drug name enters the prompt, and
% where the probe reads. RIGHT: the two ways a model can ignore a drug it was
% told about, one of them struck by the positive and the other left standing,
% beside the small licence the strike buys and the large one a failure would
% have bought. The figure draws an ASYMMETRY, not a result: the section's own
% discount, that a positive here was near-certain in advance while a negative
% would have been decisive, has no graphical form anywhere in the document, and
% nothing in thesis_v2 draws the prompt, its fields, or the read position.
%
% EVERY NUMBER DRAWN, AND WHERE IT COMES FROM
% -------------------------------------------
% Nothing in this figure is a plotted coordinate. Four accuracies appear as
% TEXT and seven layer indices are the only quantity with a POSITION. Sibling
% ledger of every one of them at machine precision: thesis_v2/figs/licensing.json.
%
%   set as text, panel A                       source
%   0.821  probe at layer 9        RESULTS_cluster/workspace_probe.json ::
%                                  layers.9.drug_decode = 0.8208333333333334
%   0.758  probe at layer 16       ... :: layers.16.drug_decode
%                                       = 0.7583333333333334
%   0.083  shuffled-label floor    ... :: chance = 0.08333333333333333
%                                  (identical at layers.<L>.chance)
%   0.883  context removed, L16    ... :: layers.16.drug_decode_context_removed
%                                       = 0.8833333333333332
%   480    prompt count            COMPUTED, not stored: config.n_drugs 12 x
%                                  config.n_per_drug 40. The thesis prints
%                                  \num{480} at 04b-representation.tex:115-120.
%   12, 40 twelve drugs, forty     ... :: n_drug_classes = 12, config.n_drugs
%          cells each              = 12, config.n_per_drug = 40
%
%   [SUPERSEDED IN PART BY PASS 12, below: the four ACCURACIES 0.821, 0.758,
%   0.083 and 0.883 are no longer set in this figure at all. 480, 12 and 40
%   remain, and the block above is kept verbatim because the ledger's
%   removed_from_drawing entries are audited against it.]
%
%   DRAWN AS A POSITION, panel A
%   2, 4, 6, 8, 9, 12, 16          ... :: config.layers = "2,4,6,8,9,12,16".
%   The seven ticks on the depth column are at TRUE layer spacing,
%   y = 2.10 + (L/16)*4.80, so the even grid and the one added layer (9) are
%   visible as such. This is the figure's only scaled mark and it is a scale of
%   depth, not of any measured quantity. The seven y values are, and are
%   CHECKED against the map rather than against the eye:
%     L    2     4     6     8     9     12    16
%     y   2.70  3.30  3.90  4.50  4.80  5.70  6.90
%   Layer 12 was drawn at 6.00 in the first build and is a repaired defect; see
%   design decision 2.
%
%   NOT FROM ANY JSON -- design facts, cited by line, drawn_true false in the
%   ledger. This is the handling slab.json gives the appendix's +2.75 nats.
%   the five prompt fields, in this order:
%     cell line | drug | dose | mechanism | control cell sentence
%     Sections/03-Apparatus.tex:250-252, "The prompt supplies cell line, drug,
%     dose, mechanism and the control cell's sentence".
%     CHECKED AGAINST figs/hook.tex:335-339, which draws the same strip for
%     sec:res-used: the five labels are spelled identically and are in the same
%     order, and hook.tex's header records the same measured \scriptsize widths
%     (cell line 0.954, drug 0.616 bold, dose 0.579, mechanism 1.424, control
%     cell sentence 2.511; total ink 6.084cm), arrived at independently. Only
%     the cell edges differ, because hook's strip is 8.20cm and this one is
%     7.20cm and each distributes its own slack. The two figures are meant to
%     read as one object seen twice, and they do.
%   the read position, "the last prompt position ... where generation begins":
%     Sections/04b-representation.tex:115-118.
%   the two failure modes, every word of F1 and F2:
%     Sections/04b-representation.tex:33-36, the \begin{question}[3] box.
%
% AGREEMENT WITH THE TEXT (checked against the Sections file, not assumed):
%   04b:126-127  "reads $82\%$ at layer 9 and $76\%$ at layer 16, the deepest
%                measured, against the $\approx 8\%$ floor"
%                -> panel A prints 0.821, 0.758 and the 0.083 floor.        OK
%   04b:131-132  "once cell line, plate and dose are projected out, the series
%                rises with depth instead, to $0.883$ ($88\%$)"
%                -> panel A prints 0.883 with that condition attached.      OK
%   04b:118-120  "A cross-validated multinomial logistic-regression classifier
%                learns which of the twelve drugs the prompt named"         OK
%   04b:122-124  "Folds are stratified over prompts and not grouped by well"
%                -> drawn verbatim as the probe's stated limitation.        OK
%   04b:108-110  "The prompt names the drug in plain text ..., so recovering
%                it ... establishes only that the model did not discard what
%                it was handed" -> B12 is that sentence, split over two lines,
%                and it is the ONLY entry in the licence that exists.       OK
%   04b:110-112  "The informative outcome would have been a negative, and what
%                the probe buys is the depth at which one would appear"
%                -> B13's second line, "the layer at which it is lost".     OK
%   04b:33-36    "Either the drug never enters its computation, an
%                activation-side encoding failure that target- and input-side
%                fixes might reach, or it enters and nothing reads it, a
%                decoder failure only an objective-side fix could reach"
%                -> F1, F2 and B13's last two lines are that partition,
%                word for word.                                             OK
%   04b:142-144  claim box: "It rules out that the drug field is lost before
%                the decoder runs, an encoding failure in the activations. It
%                does not rule out that the decoder ignores it."
%                -> the strike lands on F1 and only on F1.                  OK
% tab:workspace (04b:228-234) prints the same 0.821, 0.758 and 0.883 in its
% Decode columns; this figure sets them as text, not as marks, so no scale
% here can be read against that table.
%
% The subtitle of panel B, "a negative, at the 0.083 floor, would have been
% decisive", names the floor a second time on purpose. Without it the figure
% never states what a failure would have LOOKED like, and the only cue left is
% the word "failed" in the right-hand tag; the acceptance question this figure
% is built to answer is precisely "what would the probe have had to return".
%
% DESIGN DECISIONS, each with the reason that forced it
% -----------------------------------------------------
%  1. THE FIGURE CARRIES NO AXIS AND NO PLOTTED VALUE. The four accuracies are
%     set as text. An accuracy axis was the obvious alternative and was
%     rejected twice over: fig:comparators row A already plots exactly these
%     numbers on exactly that rule, so a second one would be a redraw; and an
%     axis invites the reader to read the height of 0.821 as the size of the
%     finding, when the section's whole point is that the finding's size is
%     nearly irrelevant because a positive was near-certain in advance.
%  2. THE ONLY SCALED MARK IS DEPTH. The seven ticks sit at true layer spacing
%     rather than evenly, because "seven layers" and "the depth at which a
%     negative would have appeared" is what the probe actually buys (04b:111-112),
%     and an evenly spaced tick row would hide that the grid is even except for
%     the one layer added at 9. The 8/9 pair is 0.30cm apart, the document's
%     minimum label separation, which is why the column is 4.80cm tall; it may
%     not be shortened without either dropping a numeral or lying about the grid.
%     REPAIRED, PASS 9: layer 12 was drawn at y = 6.00 and the map gives 5.70.
%     Measured off the 200dpi render, 2->4->6->8->9 ran at 0.30cm a layer, then
%     9->12 at 0.40 a layer and 12->16 at 0.225, so on the page three layers
%     occupied more room than four and layer 12 sat where layer 13 belongs. The
%     column is the figure's only scaled mark and the section's payoff is the
%     depth a negative would have named, so a grid step of error here is the one
%     error the figure cannot carry. Checked safe before drawing: 5.70 is 0.90
%     above the layer-9 tick and 1.20 below layer 16, and the nearest ink at
%     that height is the val node on baseline 5.97, which ends near x = 5.5,
%     well clear of the tick column at x = 7.02.
%  2b.THE 8/9 NUMERALS ARE NOT BRACED, AND THAT IS A DECISION, NOT AN OVERSIGHT.
%     A reviewer measured 0.102cm of clear white between the two numerals
%     against a glyph height of 0.203cm and asked for comparators.tex's brace
%     idiom, "9, 8" named as an ordered set. Overruled. comparators braces marks
%     whose POSITIONS a reader cannot resolve on a measured axis, and fanning
%     them would claim a resolution the estimates do not have; here the two
%     positions are exactly right, each numeral is unambiguously attached to its
%     own 0.42cm tick, and the tightness IS the datum -- layers 8 and 9 are
%     adjacent in a sixteen-layer stack and must look it. Bracing them would
%     hide the one thing the added layer at 9 is there to show. The pair was
%     re-read at 200dpi after the numerals went quietink and the two glyphs are
%     separate and legible. The column cannot be lengthened either: its bottom
%     is pinned to the caret at y = 2.10 (which is layer 0, the prompt) and its
%     top to the "layer" label at 7.12, so the most 16 layers could ever have is
%     4.95cm, 0.009cm a layer more than they have now.
%  3. NO PANEL LETTERS. Two uppercase sans word tags in accent, the corpus
%     convention (circularity.tex, slab.tex); (a)/(b) appears nowhere in this
%     document's figures.
%  4. THE STRIKE IS THE ONLY MARK ABOVE 1.0pt. 1.1pt, the same weight as
%     circularity.tex:156, where one struck edge carries that figure's claim.
%     REPAIRED, PASS 9: it is now ONE HORIZONTAL stroke through the x-height of
%     line one, 8.70 to 13.50 at y = 5.60, and not the diagonal
%     (8.86,5.20)--(13.94,5.64) of the first build. Both reviewers reported the
%     diagonal as blocking and the render agrees: its left end sat exactly on
%     line two's baseline and its right end 0.11cm above line one's, so it
%     entered along the bottom of "an activation" and left through the x-height
%     of "computation", and at print size "an activation" was a smear. Its right
%     end was also 0.26cm from the box border with the accent arrow to "struck"
%     starting at 14.20 on nearly the same line, so the two read as one
%     continuous pointer out of the box. A rule through the x-height leaves the
%     words readable, because the eye completes struck text and cannot complete
%     text cut along its baseline. LINE ONE ONLY: line one is the proposition
%     the positive eliminates and line two is its gloss, so one stroke does the
%     work and the 1.1pt register stays "used exactly once". Extent measured off
%     the render: line one's ink is 8.76 to 13.44, the stroke overhangs 0.06cm
%     each side, and it stops 0.70cm short of the box's right border, which is
%     what kills the leader reading.
%     Everything else here is 0.30-0.90pt, so weight alone says
%     which mark is the argument. The struck box is drawn and STRUCK, not
%     deleted: a deleted box leaves the caption nothing to point at, and the
%     partition is two-membered whether or not one member survives. The two
%     boxes carry IDENTICAL borders for the same reason -- the struck one was
%     drawn at rulegrey!55 in the first build, and if the box is also fainter
%     then two marks say the same thing and neither says it alone.
%  5. THE TWO LICENCE BLOCKS ARE SIZED BY THEIR CONTENTS AND BY NOTHING ELSE.
%     The one the positive buys holds one sentence; the one a failure would
%     have bought holds four lines. That difference in footprint IS the
%     asymmetry, readable before a word is read. It quantifies nothing, and
%     both the caption and the ledger say so.
%  6. SOLID-AND-WHITE MEANS "HAPPENED", DASHED-AND-CREAM MEANS "DID NOT".
%     One register, used once each, so the counterfactual block cannot be
%     mistaken for a second result. REPAIRED, PASS 9: the cream is now panelbg
%     straight, the same cream slab.tex and hook.tex use, and so is the depth
%     column's. Three creams were in play across this slate -- panelbg in hook
%     and slab, panelbg!60 in the column, panelbg!45 in the counterfactual box
%     -- and the last two are about 2% ink, which will not reliably reproduce on
%     uncoated stock; a tint that does not print is a tint doing no work while
%     adding a third cream to the system. The DASHED border is what carries
%     "this did not happen"; the fill only says "a filled region".
%  7. THE LICENCE THAT EXISTS HAS NO SECOND ENTRY AND NO "which fixes it points
%     at" LINE. The answer to that question is nothing, and an absent line says
%     it better than the word "none", which would read as a measured category.
%     Same principle as comparators.tex's bare dot for a quantity with no
%     interval: the absence is the mark.
%  8. THE DRUG CELL IS THE ONLY ACCENT MARK IN THE PROMPT STRIP, at accent!12
%     fill with an accent border and an accent bold label, so the eye finds it
%     before it reads anything. It is the reason every number in the left-hand
%     block is what it is.
%  9. NO INTERVAL IS DRAWN ANYWHERE. The probe returns a cross-validated
%     accuracy and workspace_probe.json stores no interval for it; the document
%     draws bars only where the instrument returns its convention
%     (comparators.tex, style_v2.py). Nothing here is a bar, a whisker or a
%     dot on a scale, so nothing invites one.
% 10. THE ACCENT BUDGET IS THE DRUG CELL, THE CARET, THE SEVEN DEPTH TICKS,
%     THE STRIKE AND ITS ARROW AND LABEL, AND THE FOUR WORD TAGS.
%     REPAIRED, PASS 9: the seven depth NUMERALS were accent and are now
%     quietink. With them the figure carried about 24 accent marks against a
%     corpus ceiling near 16 (comparators.tex gives accent exactly two jobs),
%     and roughly fourteen of them stood inside one 0.4 x 4.8cm column, making
%     it the densest accent region on the page -- so the eye landed on the tick
%     column first and then hunted beside tick 9 for "0.821", which is set as
%     prose three centimetres to its left. A numeral beside a tick is an axis
%     label, which is what quietink is for in every other figure in the corpus
%     (comparators.tex:70-72, gate.tex, purify.tex). The TICKS stay accent
%     because the depth grid is what \cref{sec:methods-probe} says the probe
%     actually buys. Re-rasterised at 115dpi after the change -- about print
%     size at arm's length -- and the first fixation is now the drug cell.
% 11. NO 0.50 CHANCE RULE APPEARS. This instrument's floor is 1/12 = 0.083, not
%     0.50, and it is set as text rather than as a rule because there is no
%     axis for a rule to cross. (Recorded because the corpus is split on how a
%     chance rule is drawn: style_v2.chance_line() mandates thin solid black,
%     while ladder.tex:112-115, the only TikZ precedent, draws it densely
%     dotted. This figure needs neither.)
%
% LAYOUT NOTES, all forced by a build or by a measurement
% -------------------------------------------------------
%   * fullwidth = textwidth 142mm + marginparwidth 28mm + marginparsep 4mm
%     = 174mm = 493.23pt. (preamble.tex's own comment beside
%     \newenvironment{fullwidth} says 171mm and is stale arithmetic.) The box
%     is pinned with a bare \path to 17.30cm = 491.4pt, 1.8pt inside the
%     measure. Nothing is drawn outside x in [0, 17.30] except the 0.2pt
%     half-width of B13's right border.
%   * EVERY MULTI-LINE BLOCK IS SEPARATE anchor=base west NODES. No node in
%     this file uses align= or text width. frames.tex measured a 3.7pt drift
%     from that: TikZ's align builds a minipage whose first-line height follows
%     the AMBIENT \baselineskip, so a node's ink moves with whatever prose
%     precedes the float. Every y below IS a baseline. Within-block leading is
%     0.33cm = 9.39pt, one \scriptsize line; between blocks it is 0.45cm.
%   * THE PROMPT STRIP'S CELL EDGES ARE SET FROM MEASURED STRING WIDTHS, not
%     chosen. Every string in this figure was set in the real preamble and its
%     \wd read back before a coordinate was written. At \scriptsize the five
%     field names measure, in cm:
%       cell line 0.954 | drug (bold) 0.616 | dose 0.579 | mechanism 1.424 |
%       control cell sentence 2.511.  Total ink 6.084cm.
%     The strip therefore cannot be the 6.78cm first laid out: at that width
%     "mechanism" alone overruns its 1.28cm cell by 0.144cm and the five cells
%     share 0.696cm of padding, 0.07cm a side. The strip is 7.20cm, the widest
%     that keeps the depth column's numerals inside panel A, which leaves
%     0.09-0.11cm each side -- more than the 1.6pt inner sep slab.tex sets on
%     its own nodes. Cell edges: 0.00 / 1.13 / 1.93 / 2.69 / 4.29 / 7.20.
%   * THE ELLIPSIS INSIDE THE FIFTH CELL WAS DROPPED, AND THE MEASUREMENT IS
%     WHY. A typeset $\cdots$ at \scriptsize is 0.405cm and needs about 0.16cm
%     of gap; the fifth cell is 2.91cm and its label is 2.511cm, leaving
%     0.399cm in total. No strip width recovers it: pushing the caret past
%     x = 7.30 drives the layer numerals to within 0.13cm of the panel divider.
%     A drawn three-dot ellipsis (0.18cm) leaves a 0.066cm gap to the label,
%     which is 1.9pt and reads as part of the last word.
%   * THE FIFTH CELL'S RIGHT EDGE IS SOLID. It was drawn dotted first, on the
%     reading that the strip has no horizontal scale. Two faults, both seen in
%     the render: the caret covers the top 0.21cm of that edge, leaving a
%     0.41cm dotted stub that reads as a printing fault; and a dotted edge at
%     exactly the x the caret marks says both "the prompt ends here" and "this
%     field runs on past here". The caret is the load-bearing mark, so the
%     absence of scale is stated in the caption, in words, instead.
%   * THE CARET STRADDLES THE STRIP'S RIGHT EDGE ON PURPOSE, and is centred on
%     the depth column above it. The last prompt position is where the prompt
%     ends, so the mark for it sits ON that boundary rather than inside the
%     last cell. PASS 9: its apex is now the column's own bottom edge, y = 2.10,
%     and the 0.5pt accent stub that used to bridge 2.04 to 2.10 is deleted.
%     Inside a 3mm square at that corner there converged the strip's top rule,
%     the strip's right rule, the caret, the stub, the column's bottom edge and
%     the brace's right arm, and the stub crossed a rulegrey rule uninterrupted.
%     Six strokes became five and the apex does the joining. The brace's right
%     arm was left where it is: measured, it terminates at (7.20,1.30), which is
%     0.12cm below the strip's bottom-right corner and 0.56cm below the caret's
%     own base, so it is not part of that corner.
%   * THE PROMPT CELLS ARE DRAWN AS figs/hook.tex DRAWS THEM. PASS 9: rulegrey
%     0.35pt, panelbg fill, rounded corners=1pt for the four ordinary fields,
%     and accent 0.4pt / accent!12 / rounded corners=1pt for the drug field.
%     Before this pass licensing drew square-cornered WHITE cells sharing edges
%     and hook drew rounded panelbg boxes with rounded gutters, so the same
%     five-field prompt was drawn two incompatible ways in two figures a page
%     apart while both captions insisted they were the same prompt. hook's is
%     the better mark -- the rounding reads as token boundaries rather than as
%     a table -- so licensing adopts it. Only the cell EDGES still differ, and
%     they must: this strip is 7.20cm and hook's is 8.20cm, and each distributes
%     its own slack. That is a difference of measure, not of treatment.
%   * NO LEADER RUNS FROM THE CARET LABEL TO THE CARET. One was drawn and
%     taken out: the only clear path from the label's lower right corner passes
%     within 0.02cm of the depth column's bottom-left corner at (7.02,2.10), or
%     else crosses the strip's top border at x = 6.94 and runs on inside the
%     last cell. Both read as a stray stroke. The label's second line ends
%     0.21cm from the caret's footprint with nothing between, inside the
%     document's ~0.3cm leader threshold.
%   * NO ACCENT MARK LIES LEFT OF x = 7.02. The seven depth ticks first ran
%     from 6.92, within 0.49cm of the text block, and the block's line at
%     y = 3.30 shares its baseline exactly with the layer-4 tick, so
%     "0.883 once cell line, plate and dose are projected out ---- 4" read as
%     one labelled row. The ticks now start at the column's own left edge.
%   * PANEL A'S LINE LENGTH IS CAPPED AT 6.85cm, not the 6.35cm of the first
%     layout: that limit followed the depth column, which moved right from
%     6.60 to 7.02 when the strip was widened. The longest line drawn is
%     "where generation begins; read at seven depths" right-aligned on 6.84,
%     and the longest left-aligned line is "a cross-validated
%     multinomial probe, twelve classes" at 6.57cm.
%   * "READ AT SEVEN DEPTHS", NOT "READ AT EVERY LAYER". Beside a column that
%     is sixteen layers tall and carries seven ticks, "every layer" left a
%     reader unable to tell whether the ticks were the seven layers actually
%     probed or an arbitrary sample of a sweep over all sixteen -- which is the
%     one fact the column exists to carry. The review asked for "read at each of
%     the seven marked depths"; measured, that is 1.94cm longer and the node is
%     anchor=east on 6.84, so it would have started at x = -1.1 and taken the
%     picture off the measure. "seven depths" is 20 characters against "every
%     layer"'s 19 and cost nothing.
%   * THE COUNTERFACTUAL BLOCK'S FIRST ENTRY LOST ITS FINITE VERB, and the
%     review's own wording would not fit. "the drug is lost inside the model,"
%     read as an assertion that the drug IS lost, the opposite of the section's
%     conclusion; the review asked for "that the drug is lost inside the
%     model,", which measures 4.66cm from x = 12.84 and put the picture 11.88pt
%     over the measure -- the harness reported exactly that Overfull hbox on the
%     build that tried it. "the drug lost inside the model," is 4.02cm, fits
%     with 0.44cm to spare, and is the noun phrase the other three entries are.
%   * THE CARET LABEL IS SHORT LINE OVER LONG LINE, not the reverse. Written
%     long over short it sat 0.75cm under a left-aligned quietink line of the
%     same width and register, and at print size the two read as one three-line
%     grey paragraph. Short first, right-aligned, breaks that; the long second
%     line then also fills the 3.8cm of empty measure that sat above the strip.
%   * B14 IS FIGURE-WIDE, NOT PANEL-LOCAL. It measures 11.70cm at \footnotesize
%     and is centred on 8.65, the figure's own centre, below the panel divider
%     (which stops at y = 0.90). Panel-local it would have run from 7.10 to
%     18.80 and off the measure.
%   * TYPE. Ambient \small on the picture; every node sets its own size.
%     \scriptsize = 8pt on all of it except the two panel subtitles and the
%     closing line, which are \footnotesize = 10pt, and the four word tags,
%     which are \sffamily\scriptsize\bfseries. Nothing is set at body size and
%     nothing is scaled: 1cm here is 1cm on the page, no scale= key, no
%     \resizebox.
%
% ===========================================================================
% PASS 12 -- 2026-08-18. THE STYLE-CONTRACT REBUILD.
% ===========================================================================
% Everything above this line is the v1-v11 record and is left intact because
% the thesis is audited against it. This section states what pass 12 changed,
% and SUPERSEDES BY NAME every claim above that pass 12 makes false. Nothing
% above has been edited away.
%
% WHY. The supervisor's verdict on the 4.8-4.12 slate: "go back over them and
% iron them out even more ... make it look more professional, right now it
% looks a bit sloppy and childish ... make sure that the figures are CORRECT
% and display the idea we want to convey in the best way possible." The
% reference standard is fig:mechanism (4.11, printed page 51): lowercase
% parenthesised panel letters with sentence-case descriptive titles, rotated
% axis labels that NAME the measured quantity, direct labels, no legend, no
% all-caps banners, no italic aphorisms, no boxed callouts, and the whole
% argument carried by a twenty-line caption. Measured against it this figure
% was, before pass 12: 4 all-caps sans accent banners out of 4 panel heads,
% 17% of its characters italic, 11.21% ink at 200dpi and 9.9 characters per
% 100mm^2 -- DENSER THAN THE APPENDIX REFERENCE PLATE A.3 (9.7%, 6.4) while
% sitting in the running argument of chapter 4.
%
% WHAT PASS 12 DELETED FROM THE DRAWING, AND WHY
% ----------------------------------------------
%  D1. THE FOUR ACCURACIES, set as running text in panel A: "$0.821$ at layer
%      9 and $0.758$ at the deepest,", "against a $0.083$ shuffled-label
%      floor", "$0.883$ once cell line, plate and dose are projected out", and
%      the "what it returns" key that headed them. They sat immediately left of
%      a scaled layer stack carrying ticks labelled 9 and 16, so a reader tries
%      to read 0.821 off the stack at layer 9 and cannot -- and the caption had
%      to spend two whole sentences denying the reading the drawing offered
%      ("They are set here as text, not as marks on a rule; the accuracy rule
%      is Figure 4.7, row one"). A caption that must deny a reading its own
%      drawing invites is a defect in the drawing, not in the caption.
%      All four are drawn or printed in three other places -- fig:comparators
%      row (a), fig:mechanism panel (a), tab:workspace -- and this figure's own
%      caption prints them as well. Under the cross-figure redundancy
%      resolution, fig:mechanism (a) OWNS 0.821 and 0.883 because it is the
%      only place the two series are drawn AS SERIES and the crossing at layer
%      12 is visible; 4.7 row (a) keeps the chance floor because a comparator
%      row without its comparator is not a comparator row. 4.8 keeps the
%      ARGUMENT about the floor -- what a negative would have looked like --
%      and loses the numerals. All four survive in licensing.json as
%      removed_from_drawing entries with their key paths and with pointers to
%      the figures that still draw them.
%      SUPERSEDES: the "set as text, panel A" block above (kept verbatim, since
%      the ledger's removed_from_drawing entries are audited against it); the
%      AGREEMENT-WITH-THE-TEXT lines for 04b:126-127 and 04b:131-132 (the
%      agreement is now the CAPTION's to carry, not the drawing's); design
%      decision 1's second clause; and the paragraph beginning "The subtitle of
%      panel B ... names the floor a second time on purpose", which is void
%      because that subtitle is deleted (D3).
%  D2. THE FOUR ALL-CAPS SANS ACCENT BANNERS -- "WHERE THE DRUG NAME ENTERS",
%      "WHAT THE OUTCOME LICENSES", "WHAT THE POSITIVE BUYS", "IF IT HAD FAILED
%      INSTEAD". Every one of them narrated rather than named. A.3 (figs/slab.tex)
%      does use short accent caps banners, but it is an appendix reference
%      plate a reader stops at and studies, its banners are noun phrases naming
%      objects, and the caps sans accent face is byte-identical to the
%      preamble's \marginlabel apparatus face -- which belongs to the margin,
%      not to a figure in the running body. 4.11 and 4.13 carry zero banners.
%      Replaced by 4.11's form: a lowercase parenthesised letter and a
%      sentence-case descriptive title, \footnotesize roman ink, left-aligned
%      to the panel's own x origin, no colon, no bold, no caps.
%      SUPERSEDES design decision 3 ("NO PANEL LETTERS ... (a)/(b) appears
%      nowhere in this document's figures") IN FULL. It was true when written
%      and is now false: fig:mechanism sets "(a) encoded more, used no more"
%      and "(b) ablation cost against a random null", and it is the reference
%      standard the slate is being brought to.
%  D3. THE THREE ITALIC SUBTITLES -- "the drug is named in the prompt itself",
%      "a positive was near-certain in advance;", "a negative, at the $0.083$
%      floor, would have been decisive" -- and D4's epigram. The subtitles were
%      \footnotesize italic quietink, i.e. the LARGEST type in the figure was a
%      grey italic editorial line while every measured thing was set 25%
%      smaller. The last of the three is also the caption's own lede restated.
%  D4. THE CENTRED ITALIC EPIGRAM "one run, two possible outcomes; the one that
%      could not fail is the one that happened." It was 9.96pt Pagella-Italic
%      quietink centred within 0.2pt of the figure's own x-centre, and the
%      naive-reader pass reports the eye going to it BEFORE any mark, reading
%      it as a second caption above the real one. In this document the argument
%      lives in the caption. It is written to figs/_PENDING_CAPTIONS.tex.
%  D5. THE FOUR BOX STROKES -- the two failure-mode boxes, the solid white
%      "what the positive buys" block and the dashed panelbg counterfactual
%      block. A stroked rectangle is drawn only where the rectangle IS an
%      object in the argument (a token, a prompt field, a subspace, a scoring
%      frame). These four enclosed SENTENCES, and the figure's own caption
%      conceded that "the figure quantifies nothing about the size of a
%      licence". Panels (c) and (d) are now two columns of equal width on one
%      x-grid with baseline-aligned entries, in the frames.tex idiom.
%      SUPERSEDES design decision 6 IN FULL (there is no dashed border and no
%      cream fill left to carry "this did not happen"; quietink carries it, as
%      it does everywhere else in the slate), and design decision 5's second
%      sentence -- see C4 for what replaced the footprint asymmetry.
%  D6. THE GREY FILL INSIDE THE LAYER STACK, and the stack itself as a filled
%      column. It read as a BAR: the naive-reader pass's first question was
%      what quantity the length of the grey bar encoded, and the answer is
%      none. It also ran from y = 2.10 to 7.05, i.e. past the layer-2 tick
%      below and the layer-16 tick above, implying that depth continues beyond
%      the measured range. It is now a plain 0.40pt ink axis rule running from
%      the layer-2 tick to the layer-16 tick and stopping exactly on them.
%  D7. "folds stratified over prompts, not grouped by well". Prose already
%      carried at 04b:122-124; it is a limitation of the instrument, which is
%      caption and prose material, not a name for anything the figure draws.
%      SUPERSEDES the AGREEMENT line for 04b:122-124 above.
%
% WHAT PASS 12 CHANGED, AND WHY
% -----------------------------
%  C1. THE AXIS NOW DEFINES ITSELF. This is the supervisor's first complaint
%      and it is answered mechanically. The depth stack is the figure's only
%      axis and it carried, before this pass, seven bare numerals and the word
%      "layer" set as a 0.30cm-tall stub above the column top -- no quantity in
%      axis-label position, no scope. It now carries 4.11's form exactly:
%      rotated 90 degrees reading upward, immediately left of the tick
%      numerals, line 1 \footnotesize roman ink "layer", line 2 \scriptsize
%      quietink "residual stream read at the last prompt position, where
%      generation begins". Line 2 is 258.63pt = 9.09cm set on one line, which
%      does not fit beside a 4.20cm axis, so it is wrapped onto two typeset
%      lines at the only break that keeps both under the 4.60cm of clear
%      vertical measure available at that x: "residual stream read at the last
%      prompt" (137.06pt = 4.817cm) over "position, where generation begins"
%      (119.57pt = 4.202cm). It is ONE label line wrapped, not two label lines.
%      Rotated multi-line text stacks toward increasing x, so line 1 ("layer")
%      sits FARTHEST from the axis, at x = 5.44, exactly as a wrapped
%      matplotlib ylabel does in 4.11.
%  C2. THE PANEL-A PROSE BLOCK IS NOW ONE BLOCK OF TWO LINES, and it is the
%      figure's only prose. "what is read: the residual stream at the last
%      prompt position" measures 264.20pt = 9.29cm at \footnotesize and cannot
%      be set on one line inside a 6.90cm panel; more to the point, C1 puts
%      that exact fact in axis-label position, where the style contract wants
%      it, so keeping it in the prose block as well would be a duplication.
%      The block is therefore "what reads it: a cross-validated" over
%      "multinomial probe, twelve classes" (4.851 and 5.241cm), \footnotesize
%      roman ink, left edge on panel (a)'s x origin.
%  C3. THE COLOUR INVERSION IS FIXED, and it was a correctness defect, not a
%      preference. Before this pass the accent arrow marked the STRUCK branch,
%      so accent meant ELIMINATED here and meant THE HEADLINE RESULT in the
%      figure four pages earlier; the naive-reader pass read the accent arrow
%      as the surviving branch. Across the slate accent now means one thing:
%      the drug-side object the rest of the chapter measures. In this figure
%      that is exactly two objects -- the drug cell in the prompt strip, and
%      the SURVIVING failure mode with its arrow and its outcome label. The
%      struck mode, its arrow, its outcome label and the whole of panel (d),
%      which is a counterfactual, are quietink. Panel titles, the axis, the
%      axis labels, the prose, the caret and panel (c) are ink; ticks, brace,
%      dividers and the four ordinary prompt cells are rulegrey furniture.
%      SUPERSEDES design decision 10's accent budget IN FULL: the seven depth
%      ticks are no longer accent (they are furniture, 0.30pt rulegrey), the
%      caret is no longer accent (it is part of the axis apparatus, ink), the
%      strike is no longer accent (it is the black threshold-weight stroke, see
%      C6), and the four word tags no longer exist (D2). It also supersedes the
%      layout note "NO ACCENT MARK LIES LEFT OF x = 7.02", which is now
%      trivially true for a different reason: the only accent in panel A is the
%      drug cell, at x 1.12 to 1.90.
%  C4. PANELS (c) AND (d) ARE NOW EQUAL. Two columns of equal width (4.60cm)
%      on one x-grid, entries baseline-aligned, FOUR ENTRIES EACH. Before this
%      pass the counterfactual block was the larger and more detailed of the
%      pair -- one sentence against four lines, inside a bigger dashed box --
%      so the eye landed on the thing that did not happen. The asymmetry the
%      figure exists to draw is now carried by the WORDS rather than by the
%      footprint: panel (c)'s entries 3 and 4 name, in the same slots and the
%      same register as panel (d)'s entries 2 to 4, exactly what the positive
%      did NOT buy -- "no layer, no class of failure," / "and no class of fix."
%      against "the layer at which it is lost," / "an activation-side failure,
%      and" / "target- and input-side fixes." A reader can now diff the two
%      columns line by line, which is a stronger reading of the asymmetry than
%      a difference of box size ever was.
%      SUPERSEDES design decision 7 ("an absent line says it better than the
%      word 'none'") and design decision 5's second sentence. The reason
%      decision 7 gave still stands against the bare word "none", which would
%      read as a measured category; what is drawn is not "none" but a named
%      absence in a slot the facing column fills, which is the opposite of a
%      category and cannot be read as a quantity.
%  C5. ONE X ORIGIN PER PANEL, declared as a named coordinate and referenced by
%      every element of that panel. Before this pass each banner sat 4.6pt
%      (1.6mm) LEFT of the text it headed, four times over, and the central
%      divider had 8.0pt of clearance on its left and 15.2pt on its right.
%      Panel (a) origin x = 0.00; panels (b), (c) origin x = 7.75; panel (d)
%      origin x = 12.70. The divider sits at x = 7.42, which is 0.33cm from
%      panel (a)'s rightmost ink (the caret's right base at 7.09) and 0.33cm
%      from panel (b)'s origin: the gutter is now equal on both sides to
%      0.001cm.
%  C6. FIVE STROKE WEIGHTS BECOME THREE, EACH WITH ONE MEANING.
%        0.30pt rulegrey -- furniture: the seven ticks, the brace, the panel
%                           divider, the (c)/(d) row rule, the four ordinary
%                           prompt cells, and the two outcome leaders (which
%                           carry their mark's colour, not rulegrey).
%        0.40pt          -- the axis rule (ink) and the drug cell's border
%                           (accent). Used twice, so it is not a singleton.
%        0.90pt ink      -- the strike, and nothing else.
%      The 1.1pt register of design decision 4 is retired: 0.90pt is the
%      document's threshold weight (style_v2.chance_line() is hard-coded
%      lw=0.9, and figs/family.pdf draws the pre-registered materiality margin
%      at 0.90pt), and a figure that contains no threshold at all can carry its
%      one mark of argument at that weight without any collision. THAT IS THE
%      JUSTIFICATION the contract requires for a weight used in one kind of
%      place: this figure has no chance rule, no gate, no margin and no floor
%      (design decision 11, which still holds), so nothing else here may be
%      black solid 0.90pt and nothing is.
%      SUPERSEDES design decision 4's first sentence and its "used exactly
%      once" clause: the strike is now TWO strokes, one through each of the two
%      lines of the struck mode, because the mode is a two-line entry and a
%      rule through its first line only left its second line standing. It also
%      supersedes the cell-border weight recorded in the pass-9 prompt-cell
%      note: 0.35pt is not one of the document's five weights and the ordinary
%      cells are now 0.30pt. The cells' construction is otherwise unchanged and
%      still matches hook.tex's \tokbox.
%  C7. THE FIGURE FOOT CARRIES ONE \scriptsize QUIETINK DISCLOSURE LINE,
%      left-aligned to the figure's x origin: "Nothing in this figure is a
%      measured quantity; the depth ticks are the only scaled marks." Every
%      figure in the slate carries exactly one such line. Here it does the job
%      an interval declaration does elsewhere -- it tells the reader, without
%      leaving the figure, that no mark on the page is an estimate. Design
%      decision 9 (no interval is drawn anywhere) still holds and is now
%      readable on the drawing rather than only in the ledger.
%  C8. THE STRIP IS 6.90cm, NOT 7.20cm, and the caret is 0.40cm tall rather
%      than 0.24cm. Both follow from C1: the rotated axis-label block needs a
%      clear vertical measure of 4.82cm at x = 5.76, which puts its lower edge
%      at y = 2.391, so the strip's top edge came down from 2.46 to 2.28 and
%      the caret grew to span 2.28 to 2.70 -- 0.111cm of clearance, measured,
%      between the rotated label's foot and the strip's top rule. The strip's
%      slack is redistributed over the same five measured field widths:
%      6.084cm of ink and 0.816cm of padding, 0.0816cm a side, giving cell
%      edges 0.00 / 1.12 / 1.90 / 2.64 / 4.23 / 6.90 against the old
%      0.00 / 1.13 / 1.93 / 2.69 / 4.29 / 7.20. NO FIELD CHANGED, no order
%      changed, and the strip still carries no horizontal scale.
%      SUPERSEDES the layout note "THE PROMPT STRIP'S CELL EDGES", the
%      "PANEL A'S LINE LENGTH IS CAPPED AT 6.85cm" note (the cap is now 6.90cm,
%      set by the strip, and the longest left-aligned line drawn is 5.241cm),
%      the caret note's "apex is now the column's own bottom edge, y = 2.10"
%      (the apex is the axis rule's own base at the layer-2 tick, y = 2.70),
%      and the "B14 IS FIGURE-WIDE" note (B14 is deleted, D4).
%      The ELLIPSIS note and the FIFTH CELL'S RIGHT EDGE note both still hold
%      and were re-checked at the new width: the fifth cell is 2.67cm against a
%      2.511cm label, so there is even less room for a $\cdots$ than there was.
%  C9. WHAT DID NOT CHANGE, and must not, because it is the argument and the
%      measurement: the seven depth ticks and their map y = 2.10 + (L/16)*4.80,
%      giving 2.70 / 3.30 / 3.90 / 4.50 / 4.80 / 5.70 / 6.90 -- BYTE-IDENTICAL
%      to pass 11, including the pass-9 repair of layer 12 from 6.00 to 5.70;
%      the five prompt fields, their order, their measured widths and the
%      accent on the drug cell; the caret at the strip's right edge and the
%      read-here relation between the two; the brace and the count line
%      "$480$ real tier-2 prompts: twelve drugs, forty cells each"; the strike
%      on the eliminated mode and the two-member partition; and the fact that
%      the figure draws no interval, no bar and no marker on a rule.
%      Design decisions 2, 2b, 8 and 11 stand unchanged. NO NUMBER IN THIS
%      FIGURE CHANGED VALUE IN PASS 12; four numbers left the drawing entirely
%      and are recorded in licensing.json as removed_from_drawing.
%
% GEOMETRY OF PASS 12, all measured with \wd in the real preamble before a
% coordinate was written (figs/_preview.py licensing, 200dpi, read at each step)
% -------------------------------------------------------------------------
%   panel (a) origin x = 0.00 ; panels (b) and (c) origin x = 7.75 ;
%   panel (d) origin x = 12.70 ; divider x = 7.42 ; measure 0 to 17.30cm.
%   axis rule x = 6.90, y 2.70 to 6.90 ; ticks 6.62 to 6.90 ; numerals
%   anchor=east on 6.52 ; rotated "layer" at x = 5.44, wrapped scope note at
%   x = 5.76 and 6.04, all centred on the axis midspan y = 4.80.
%   Measured string widths used (cm, \footnotesize unless marked s = \scriptsize):
%     (a) title 4.969 | (b) title 9.495 | (c) title 4.081 |
%     (d) title, wrapped 4.157 over 1.921
%     prose 4.851 / 5.241
%     mode 1 5.842, gloss s 4.292 | mode 2 4.472, gloss s 2.093
%     outcomes s 2.727 / 2.147 and 2.921 / 1.712
%     (c) entries s 3.210 / 2.516 / 3.352 / 2.293
%     (d) entries s 3.804 / 3.303 / 3.689 / 3.399
%     count line s 6.557 | foot line s 10.860
%     cell labels s: 0.954 / 0.616 bold / 0.579 / 1.424 / 2.511
%   Widest object on the page is panel (b)'s title, 9.495cm, which ends at
%   x = 17.245, 0.055cm inside the measure. No text crosses the measure.
%
% C10. THE EMPTY BAND IN PANEL (a) IS STRUCTURAL AND IS LEFT EMPTY ON PURPOSE.
%   Measured on the 200dpi render it is about 5.2 x 4.0cm, at x 0.0-5.2 and
%   y 2.4-6.4. It is forced, not careless, and four ways of removing it were
%   built and rejected:
%     (i)  shortening the axis. The 8/9 pair is 0.30cm apart at the current
%          scale, which is the document's minimum label separation, so the
%          16-layer span cannot be under 4.80cm and the drawn 2-to-16 span
%          cannot be under 4.20cm without dropping a numeral or falsifying the
%          grid. Design decision 2 and 2b, both still standing.
%     (ii) narrowing the strip. Its five field names measure 6.084cm of ink and
%          hook.tex sets the same five, so the strip cannot go below about
%          6.7cm without either abbreviating a field or breaking the agreement
%          with fig:hook that both captions assert.
%     (iii)inverting the axis so the strip sits at the top and depth runs
%          downward. This moves the band, it does not remove it, and it would
%          draw the residual stream running the wrong way.
%     (iv) moving panels (c) and (d) into the band. They need 2 x 4.60cm plus a
%          gutter, which is 9.55cm; the band is 5.2cm wide.
%   What is left is the space between the prompt and the read, which is the
%   residual stream, and this figure deliberately draws nothing there: the
%   stream's contents are fig:mechanism's subject and the slab's geometry is
%   fig:slab's. Filling it with narration is exactly what pass 12 removed.
% ===========================================================================
% PASS 13 -- 2026-08-18. THE REFUTATION PASS.
% ===========================================================================
% Three adversarial refuters read the pass-12 render. Everything above is left
% intact; this block states what pass 13 changed and SUPERSEDES BY NAME every
% pass-12 claim it makes false. Each change is recorded at the code it changes
% as well, so a reader of the drawing does not have to come back here.
%
% R1. THE COORDINATE RECORD IN "GEOMETRY OF PASS 12" AND IN C5 WAS WRONG IN
%     SIX PLACES, AND THE HEADERS ARE WHAT THE THESIS IS AUDITED AGAINST.
%     Superseded, by name, with the values the code actually drew in pass 12:
%       divider x         header 7.42   -> DRAWN 7.40   (and now deleted)
%       caret right base  header 7.09   -> DRAWN 7.05   (and now 7.02)
%       divider gutters   header 0.33   -> DRAWN 0.35 both sides
%       tick inner end    header 6.62   -> DRAWN 6.74   (and now 6.78)
%       numerals anchor   header 6.52   -> DRAWN 6.64   (and now 6.70)
%       rotated "layer"   header 5.44   -> DRAWN 5.36   (and now 6.18)
%       scope line two    header 6.04   -> DRAWN 6.12   (and now deleted)
%     The same block recorded the two failure-mode gloss lines at their
%     \scriptsize widths ("gloss s 4.292", "gloss s 2.093") while pass 12's
%     code set both at \footnotesize (5.365 and 2.616cm) -- and it is the gloss
%     width that sets the strike's right end, so the header could not be used
%     to check the one mark the file calls its argument. Pass 13 sets the modes
%     at \scriptsize, so 4.292 and 2.093 are now the drawn widths and the
%     record and the code agree for the first time.
%     C4's "Two columns of equal width (4.60cm)" was also never true of the
%     drawing: pass 12's origins were 7.75 and 12.70 on a 17.30 measure, i.e.
%     4.95 and 4.60. See the (c)/(d) block for what pass 13 drew instead.
%
% R2. figs/_PENDING_CAPTIONS.tex DID NOT EXIST. D4 and the note above the
%     staged float both asserted that the cut epigram and the pass-12
%     replacement caption were "written to figs/_PENDING_CAPTIONS.tex"; the
%     file was never created, so every sentence cut under D1, D3, D4 and D7 was
%     unaccounted for outside a comment block in this file, and the caption
%     that actually prints in the thesis is still the pass-11 one, which is
%     false against the drawing in four places (it indexes the panels "Left,"
%     and "Right,", it prints four accuracies as "set here as text", and it
%     says the two licence blocks "are sized by their contents"). Pass 13
%     creates figs/_PENDING_CAPTIONS.tex and writes the pass-13 caption and
%     every deleted sentence into it, keyed by \label{fig:licensing}. This
%     builder still does not edit Sections/. The copy below is retained as the
%     audit trail and is marked a duplicate of the pending file.
%
% R3. THE STRIKES WERE BLACK SOLID 0.90pt, WHICH IS THE SLATE'S THRESHOLD
%     STROKE. Corrected to quietink 0.40pt. Full reasoning at the code.
%     SUPERSEDES C6's third weight; the figure now has two weights, 0.30 and
%     0.40pt, neither a singleton.
% R4. THE CARET'S APEX WAS THE AXIS BASE AND THE LAYER-2 TICK AT ONCE.
%     Broken apart; apex 2.50, base 6.78-7.02. SUPERSEDES C8's apex clause.
% R5. THE AXIS LABEL WAS THREE ROTATED LINES AND OVERRAN THE AXIS AT BOTH ENDS,
%     with the quantity name outermost. One \footnotesize "layer" beside the
%     numerals, one \scriptsize scope line outboard, 3.908cm inside a 4.20cm
%     axis. SUPERSEDES C1's wrap arithmetic and its "line 1 is the outermost".
% R6. PANELS (c) AND (d) COULD NOT BE READ ACROSS, and two of (c)'s four
%     entries were unsupported. Three paired rows; "no class of failure,"
%     deleted as contradicting this figure's own panel (b) and the claim box;
%     "an activation-side failure," deleted as a duplicate of panel (b)'s own
%     struck line; (d)'s first entry hedged to what a LINEAR probe's negative
%     licenses. SUPERSEDES C4 in full and the ledger's
%     entries_3_and_4_are_new_in_pass_12.
% R7. THE VERTICAL DIVIDER AND THE DEPTH AXIS WERE PARALLEL HAIRLINES 0.50cm
%     APART. The divider is deleted.
% R8. PANEL (b)'s MODES WERE \footnotesize, i.e. the same size as the panel
%     titles, and their outcome labels were \scriptsize on a different leading,
%     so nothing lined up across the gutter. Both are \scriptsize now, on one
%     leading, on one set of baselines. One leading per size across the figure:
%     0.40cm at \footnotesize, 0.35cm at \scriptsize.
% R9. SMALLER ITEMS: the mirrored brace's pointer aimed at empty white and is
%     now a plain span rule; "$480$" set three digits in URWPalladioL-Roma
%     against Pagella everywhere else and is now text; the prose block's
%     pronoun had no antecedent on the reading line and now names the residual
%     stream; the foot line conceded two counterexamples in its own second
%     clause and now states what the counts are; the word "seven", deleted with
%     the caret label in pass 12, is back on the drawing in that foot line,
%     which is what the still-standing pass-9 layout note asks for.
%
% WHAT PASS 13 OVERRULED, AND WHY
% -------------------------------
%  O1. "REBUILD PANEL (a): DROP THE DEPTH SWEEP, OR SET LAYER HORIZONTALLY, OR
%      SHORTEN THE AXIS BY TIGHTENING THE 8/9 PAIR." Overruled on all three
%      counts. The seven-tick depth scale at true spacing is the figure's only
%      scaled mark and its keep list names it; the 8/9 pair is already at true
%      spacing and 0.30cm apart, the document's minimum label separation, so
%      the span cannot shorten without dropping a numeral or falsifying the
%      grid (design decisions 2 and 2b, both standing); a horizontal layer axis
%      would draw the residual stream running the wrong way and would break the
%      caret's read-here relation to the strip. C10's four rejected removals
%      stand. What pass 13 did take from that finding is the head alignment:
%      panels (a) and (b) now put their first content line on one baseline.
%  O2. "RETITLE (a) 'where the drug name enters, and where the probe reads'."
%      Overruled on measurement. Panel (b)'s title is 9.495cm of the 9.55cm the
%      right half has and sits on the SAME baseline, so panel (a)'s title
%      cannot exceed about 7.0cm; the proposed title measures 8.5cm and the
%      shortest honest variant 7.7cm. The read is named at the axis instead,
%      which is where the style contract wants it.
%  O3. "DRAW THE CARET AS AN OPEN rulegrey CHEVRON." Overruled. rulegrey is
%      furniture and the caret is content -- it is the one mark of the read
%      position, and the keep list names the triangle. Breaking the coincidence
%      and shrinking it to its pass-9 size answers the defect; recolouring it
%      would demote a content mark to furniture.
%  O4. "PUT PANEL (b)'s OUTCOME COLUMN ON PANEL (d)'s ORIGIN (12.40)."
%      Overruled: the longer strike ends at 12.484, so the label column would
%      begin inside the strike. It sits at 13.44, 0.50cm clear of that
%      terminal, which is the separation the pass-9 note asks for.
%  O5. "KEEP THE STRIKES AT 0.90pt IN quietink." Overruled in favour of 0.40pt.
%      0.90pt is the slate's threshold weight whatever its colour, and at 0.90pt
%      against \scriptsize the struck line does not survive at print size --
%      which is the legibility failure the pass-12 ledger already records.
%  O6. "REACH S15's 6.0% INK AND 5.0 CHARACTERS PER 100mm^2." Not reached, and
%      the arithmetic is recorded rather than hidden: pass 13 measures 6.65%
%      and 5.42, against pass 12's 7.95% and 6.05 and pass 11's 11.21% and 9.9.
%      What is left is 697 characters of which 483 are mandated furniture --
%      four panel titles, the axis definition, the definitional prose, the five
%      field names, the count line and the one interval-style disclosure -- and
%      214 are the partition and the two licences, which are the figure. The
%      only cuts left would delete the axis definition that answers the
%      supervisor's first complaint, or one of the two failure modes.
% ===========================================================================
% PASS 14 -- 2026-08-19. THE SECOND REFUTATION PASS.
% ===========================================================================
% Three adversarial refuters read the pass-13 render and returned thirty
% findings. Everything above is left intact; this block states what pass 14
% changed, what it OVERRULED and why, and SUPERSEDES BY NAME every pass-12 and
% pass-13 claim it makes false. Each change is also recorded at the code it
% changes. NO NUMBER IN THIS FIGURE CHANGED VALUE IN PASS 14: the seven depth
% ticks are still y = 2.10 + (L/16)*4.80 -> 2.70/3.30/3.90/4.50/4.80/5.70/6.90,
% byte-identical to passes 11, 12 and 13, and the count line still prints the
% same 480, twelve and forty. The strip translated 0.14cm DOWN in y as one
% rigid body; not one cell edge in x moved, and x is the only axis on which
% that strip carries anything at all.
%
% S1. THE PENDING CAPTION FILE. All three refuters opened with the same
%     blocking finding, and it is correct: figs/_PENDING_CAPTIONS.tex contains
%     no \label{fig:licensing} block. R2 above asserts, in as many words, that
%     "Pass 13 creates figs/_PENDING_CAPTIONS.tex and writes the pass-13
%     caption and every deleted sentence into it"; the file exists and holds
%     fig:materiality, fig:comparators and fig:hook blocks, and this figure's
%     append was lost -- almost certainly clobbered by a concurrent
%     read-modify-write from another builder, which is the hazard CF12 names.
%     So pass 13 committed, one pass later, the identical failure it accused
%     pass 12 of. SUPERSEDES R2's last three sentences and the "duplicate of
%     the pending file" note above the staged float: they described an intent,
%     not the filesystem.
%     Pass 14 APPENDS the block (pure append, never a rewrite, so it cannot
%     clobber another builder's), then READS THE FILE BACK to confirm the block
%     survived, and records the byte offset it landed at in licensing.json. The
%     block carries the pass-14 caption, the four clauses of the live pass-11
%     caption at 04b:117-140 that are false against this drawing with their
%     line numbers, and every sentence passes 12, 13 and 14 cut from the
%     drawing, each against the caption sentence that now carries it.
%
% S2. "no class of fix" IS DELETED FROM PANEL (c), AND THIS IS THE ONE
%     CORRECTNESS DEFECT IN THE PASS-13 DRAWING. Full argument at the code.
%     In one line: 04b:31-36 is a two-member partition in which each member
%     names its own fix class, panel (b) strikes one member and says so on this
%     page, so the survivor and its fix class stand -- and pass 13 had already
%     deleted the sibling entry "no class of failure," for exactly this reason
%     without applying it here. (d)'s "target- and input-side fixes" goes with
%     it, because a row with one side is not a row. Both fix classes move to
%     the caption, which can state the condition a two-word slot cannot.
%     SUPERSEDES R6's list of surviving entries and C4 in full.
%
% S3. THE TWO OUTCOME LEADERS ARE DELETED. Measured: the quietink one left a
%     0.24cm gap at its left end and a 0.10cm gap at its right, touching
%     neither mark nor label, and lay midway between two quietink strikes so
%     the eye read the verdict "struck by the positive," as itself struck; the
%     accent one began 1.397cm PAST the mark it claimed to join. The pass-13
%     ledger recorded the opposite as a measured property, having measured only
%     the first of the two. Colour and exactly shared baselines carry the
%     pairing, which is 4.11's discipline. The outcome column moves 13.44 ->
%     13.00. SUPERSEDES C6's leader clause, O4's arithmetic, and pass 13's
%     lead/.style rationale; lead/.style itself is deleted.
%
% S4. THE CARET IS NAMED AT ITSELF and the three cues that made it an axis
%     arrowhead are broken: it gains a \scriptsize ink direct label "the last
%     prompt position" 0.08cm off its left base, the strip's 0.14cm drop puts
%     0.34cm between its apex and the axis base, and its base widens to 0.28cm
%     so it no longer shares an endpoint with the tick lane. Under the figure's
%     own map it now occupies layers 0.13 to 0.87, below the drawn scale.
%     SUPERSEDES R4 and licensing.json's read_position note.
%
% S5. THE ROTATED SCOPE LINE IS DELETED and "layer" moves off the layer-9 tick.
%     The scope line shared the five-word clause "at the last prompt position"
%     with the caption that actually prints, which S18 forbids, and it named at
%     90 degrees and 1.14cm away the one glyph that had no name of its own.
%     "layer" was centred on the axis midspan, which is the layer-9 position by
%     arithmetic (2.10 + (9/16)*4.80 = 4.80 = (2.70+6.90)/2), so it read
%     "layer 9"; 5.70-to-6.90 is the only tick gap on this axis long enough to
%     hold it clear. SUPERSEDES R5 and C1.
%
% S6. THE SPAN RULE UNDER THE STRIP IS DELETED. It duplicated the strip's own
%     bottom border at the same colour, weight and length 0.12cm away, and it
%     bracketed one schematic prompt under a line counting 480 of them.
%     SUPERSEDES R9's span-rule clause.
%
% S7. PANEL (b)'s TITLE IS A NAME, NOT A SENTENCE: 9.495cm -> 5.821cm, so the
%     widest object on the page is no longer a title that restated 04b:31-32.
%     Panel (c)'s row one loses its wrap and its own S18 collision ("the model
%     did not discard what it was handed" is seven words verbatim from
%     04b:107-108, six lines above the float) and becomes "the drug field is
%     not discarded", one line. The strikes start on the panel origin 7.75
%     rather than 0.06cm left of it. The (c)/(d) group drops 0.30cm and its row
%     pitch opens to 0.70cm, so deleting a row did not trade panel (a)'s hole
%     for a new one in the bottom-right corner. The foot line drops its second
%     clause. SUPERSEDES O2's measurement clause and R8's 0.52cm row gap.
%
% S8. THE PANEL-(a) BAND. The pass-13 drawing put one contiguous 56.2 x 39.2mm
%     void -- 17.2% of the figure -- under the prose block and left of the
%     axis, because pass 13 had moved that block from the band's middle up to
%     the panel head to buy a shared first-content baseline with panel (b).
%     Pass 14 moves it back to 5.00/4.60, centred on the axis midspan, and puts
%     the caret's new label in the band's lower part; the band is now two voids
%     of about 1.9 and 1.6cm and the block sits level with the axis it
%     describes. The two TITLE baselines still align at 7.45, which is the
%     alignment a reader sees. NOTHING FURTHER WAS PUT IN THE BAND, and the
%     reason is the contract, not indifference: S4 allows at most one prose
%     block per panel and panel (a) already has it, and C10's four ways of
%     REMOVING the band (shorten the axis, narrow the strip, invert the axis,
%     move (c)/(d) into it) were each rebuilt and each still fails. The band is
%     the space between the prompt and the read; this figure draws nothing
%     there because the stream's contents are 4.11's subject and the slab's
%     geometry is A.3's. SUPERSEDES pass 13's placement repair (ii).
%
% S9. MEASURED AFTER ALL OF THE ABOVE (200dpi, clipped to the content bbox):
%       content bbox   173.10 x 74.29mm   (S14 wants 170-174; cap 75mm)
%       ink            5.826%             (was 6.65 at pass 13, 11.21 at 11)
%       characters     596 non-space, 4.63 per 100mm^2  (was 5.44 at pass 13)
%                      694 with spaces,  5.40 per 100mm^2  (was 6.41)
%       stroke weights 0.30pt rulegrey furniture (4 ordinary cells, 7 ticks,
%                      the (c)/(d) row rule) and 0.40pt (axis rule ink, drug
%                      cell border accent, two strikes quietink). Two weights,
%                      neither a singleton, no 0.90pt stroke, no arrowhead,
%                      no leader, no dotted rule, no legend.
%       type           two sizes: \footnotesize 10.0pt (four panel titles, the
%                      prose block, the axis quantity name) and \scriptsize
%                      8.0pt (everything else). 0% italic. 0 all-caps banners.
%     S15's ink ceiling is met for the first time in this figure's history.
%     The with-spaces character count is still 5.40 against a 5.0 ceiling and
%     is recorded rather than hidden: the refuters' own measurement of that
%     rule was taken on non-space characters (5.44 against 5.0), which is met
%     at 4.63; closing the with-spaces gap needs 51 more characters and the
%     only strings left are the four titles, the axis name, the five field
%     names, the count line the spec's keep list mandates, the two failure
%     modes and their outcomes, the two licence rows and the one disclosure
%     line. Every one of them is mandated content.
%
% WHAT PASS 14 OVERRULED, AND WHY
% -------------------------------
%  P1. "SET THE SEVEN TICK NUMERALS IN INK AND LENGTHEN THE TICKS TO 0.20cm,
%      BECAUSE THE AXIS READS AS A PANEL DIVIDER." Overruled on both counts.
%      Quietink tick numerals are this corpus's convention, not an oversight:
%      comparators.tex:70-72, gate.tex and purify.tex all set them so, and
%      design decision 10 changed them from accent to quietink in pass 9 with
%      that reason recorded. S8 assigns ink to axis LABELS, which is where
%      "layer" is, and does not assign a colour to tick numerals; making this
%      one figure's numerals black would break it away from the four figures it
%      is being aligned to. The ticks cannot lengthen inward either: they run
%      6.78-6.90 and the numerals are right-anchored on 6.70, so 0.20cm ticks
%      would touch them. The divider reading the finding describes was created
%      by the divider itself and died with it in pass 13 (R7); what is left is
%      a rule with seven numbered ticks and a rotated quantity name, which is
%      an axis.
%  P2. "OPEN A GUTTER BETWEEN THE PROMPT CELLS, OR DROP rounded corners=1pt,
%      BECAUSE THE ROUNDINGS LEAVE WHITE NOTCHES AT EVERY INTERNAL JUNCTION."
%      The notches are real -- re-cropped at 1200dpi and seen -- and the fix is
%      still overruled here. The finding's own evidence is that figs/hook.tex
%      draws the identical construction at the identical radius, and both
%      captions assert the two strips are the same prompt seen twice; a strip
%      repaired in one file and not the other would break an agreement the text
%      states, to remove a 0.1mm artefact. This builder may not edit hook.tex.
%      Recorded in licensing.json as a coordinated repair the two files owe
%      together, so the next pass that touches both can make it once.
%  P3. "SET THE PROSE BLOCK AS ONE LINE AND DROP IT A RANK, BECAUSE IT IS THE
%      SAME SIZE AS THE PANEL TITLES." Overruled, and the finding concedes the
%      contract point itself: S4 mandates \footnotesize for definitional prose
%      and states in terms that prose OUTRANKS labels, which is A.3's ranking
%      and the one the contract chose to keep. A one-line rewrite either drops
%      "the residual stream", which is the antecedent R9 put there and the only
%      naming of the object the axis stands for, or drops "cross-validated",
%      which is why the accuracy is a presence test at all.
%  P4. "RESTORE THE PROBE'S TASK TO THE DRAWING: 'over the twelve drug
%      classes'." Overruled as to the drawing, accepted as to the record. The
%      figure is at its ink ceiling, the count line already prints "twelve
%      drugs" 2.4cm below, and naming the instrument does not require naming
%      its class count. The twelve-class task AND the fold caveat (04b:155-157,
%      cut in pass 12 as D7) are both in the pass-14 caption, and
%      licensing.json records under explicitly_not_drawn that the drawing names
%      the instrument without its target and why.
%  P5. "LEAVE PANEL (c)'s SLOTS BLANK, OR SET AN EN-DASH IN THEM, SO THE SMALL
%      LICENCE LOOKS SMALL." Overruled. An empty slot under a head reading
%      "what the positive buys" is ambiguous between "nothing" and "not drawn
%      yet", and design decision 7's own supersession at C4 gives the reason a
%      named absence beats a blank or the bare word "none". The finding's
%      underlying complaint -- that the two licences were drawn the same size,
%      with the smaller one in darker ink and more lines -- is answered by S2
%      and S7 above: (c) now sets 39 characters against (d)'s 58 and its second
%      entry is 0.991cm against (d)'s 3.233cm.
%  P6. "DRAW THE TIE BETWEEN THE DEPTH AXIS AND (d)'s 'the layer at which it is
%      lost', SO THE AXIS HAS A STATED ROLE." Overruled as a drawn mark. A
%      leader across 5.5cm of figure would be the only connector on the page
%      and would assert an argument, and a line saying "a negative would have
%      named one of these seven depths" is argument, not definition, which S4
%      and S18 put in the caption. The pass-14 caption says it in those words.
%      The axis is now fully DEFINED on the drawing -- quantity name, seven
%      numbered ticks, and the position it is read at, labelled at the caret --
%      which is what S5 asks of it.
%  P7. "MAKE THE VERTICAL RULE THE RESIDUAL STREAM, SO THE PROSE BLOCK HAS A
%      DRAWN REFERENT." Overruled. The rule is a LAYER axis; the residual
%      stream at one position is a 2048-dimensional vector at each of seven
%      depths, and renaming a depth coordinate after the object sampled along
%      it would put two names on one rule and cost S5 its clean reading. The
%      prose block was instead moved level with the axis and the caret named,
%      so the reading relation the block defines -- position, depths, probe --
%      is drawn beside it rather than four centimetres above it.
% ===========================================================================
% inner sep=0pt on every node, so that anchor=base west puts the START OF THE
% INK exactly on the panel's x origin and a node's width is exactly the \wd
% measured for its string. With TikZ's default 0.3333em the picture ran 5.43pt
% past the fullwidth measure (the harness reported exactly that Overfull hbox),
% because panel (b)'s title is 9.495cm of the 9.55cm the right half has.
% PASS 14: that title is now 5.821cm, so inner sep=0pt is no longer what keeps
% the picture inside the measure -- but it is kept, because it is also what
% makes anchor=base west put the START OF THE INK on the panel origin, which is
% what S13's one-x-origin check measures. The widest object is now panel (d)'s
% title, ending at x = 17.247, 0.053cm inside the 17.30cm measure.
%
% ===========================================================================
% v-BLOCK, 2026-08-19: WHOLE-SLATE PASS (fig:comparators, fig:licensing,
% fig:purify, fig:hook, fig:materiality reconciled against one another and
% against fig:mechanism, which is the document's reference standard and was
% NOT rebuilt). Nothing above this line is deleted; where an earlier claim is
% now false it is superseded EXPLICITLY below.
%
% WHY. The five figures were rebuilt independently and drifted. This pass owns
% only the differences BETWEEN them. No value changed in any of the five; the
% sibling .json ledgers carry the value-to-position maps, old and new.
%
% THE SLATE RULES THAT NOW HOLD ACROSS ALL FIVE (each stated once here, in
% every file, so any future single-figure edit can see what it would break):
%
%  R1 TICK NUMERALS are quietink, \scriptsize, inner sep 0, anchored north on
%     the tick and set 0.16cm below the axis rule, over a 0.10cm rulegrey
%     0.30pt tick. Before: comparators/licensing/purify quietink, hook and
%     materiality ink; ticks 0.10 / 0.12 / 0.14cm; three different gaps.
%  R2 A THRESHOLD'S OWN LABEL IS BLACK, matching the black solid 0.90pt rule
%     it names (S6). Before: comparators black, purify's "gate, 0.8" and
%     materiality's two margin labels ink.
%  R3 ONE NAME PER OBJECT (S19). The pre-registered 10% rule is "10%
%     materiality margin" in fig:comparators row (c), fig:materiality (a) and
%     (b), and in figs/family.py, which is where the name comes from.
%     fig:comparators' "pre-registered margin" is superseded; "pre-registered"
%     is now the caption's word, not the drawing's. Zero, where it is the null
%     a bar is read against, is "zero: no identity effect" in BOTH
%     fig:comparators row (c) and fig:materiality (b) -- one quantity, one
%     wording. V_pure is spelt with \textrm everywhere, as figs/slab.tex does.
%  R4 COLOUR KEY, one key for five figures (S8). accent = the drug-side object
%     under test. quietink = every comparator, null, control, replication or
%     discounted arm, AND its label, AND its leader. ink/black = thresholds,
%     axis rules, axis labels, panel heads, definitional prose. rulegrey =
%     furniture only. Within quietink the GLYPH separates two kinds: an OPEN
%     disc is a construction built to carry no signal (a null, a control, a
%     raw/before series); a FILLED disc is a real measurement that is not the
%     headline. fig:hook (d) drew both its null and its cell-line control in
%     ink and is corrected.
%  R5 PROJECTION AND ZERO RULES are quietink 0.50pt densely dotted (S7).
%     fig:hook's two panel-(c) projection lines were rulegrey 0.30pt densely
%     dashed and are corrected.
%  R6 DEFINITIONAL PROSE is \footnotesize ink and OUTRANKS every label (S4);
%     at most one block per figure, on the panel's own x origin. fig:hook's
%     one prose line was \scriptsize and is raised. fig:purify's panel-(b)
%     block carried an argument clause ("so causal claims use the purified
%     slab, not the raw") that its caption already carries, and is now
%     definitional only.
%  R7 ONE INTERVAL DECLARATION per figure, \scriptsize quietink, at the foot,
%     on the figure's x origin, opening with the word "Intervals:" and naming
%     EVERY panel including the ones with no measured mark (S17).
%     fig:comparators already did this and is the model; the other four now
%     match its form.
%  R8 THE FIVE-FIELD PROMPT STRIP is ONE object drawn in two figures, and the
%     two drawings are now congruent: cell boundaries 0 / 1.12 / 1.90 / 2.64 /
%     4.23 / 6.90 cm and cell height 0.42cm in both fig:licensing (a) and
%     fig:hook (a) and (b), labels centred in their cells, the drug cell
%     accent-bordered at 0.40pt on accent!12 and its label NOT bold. Before:
%     licensing 6.90cm wide with 0.62cm cells and a bold label, hook 7.44cm
%     wide with 0.42cm cells and a plain label -- the same five fields at two
%     sizes, four pages apart, with each caption pointing at the other.
%  R9 MEASURE. Every fullwidth figure draws 172.0-174.0mm; the one textwidth
%     figure draws 140.3mm. Ink at 200dpi inside the content bbox, counted as
%     greyscale luminance below 230 (the contract's own method), is now
%     5.63-5.99% on all five, under the 6.0% body ceiling, where before this
%     pass it ran 5.36-6.54%.
%
% WHAT THIS PASS DELIBERATELY DID NOT UNIFY, so that a later reader does not
% think it was missed:
%  * PANEL-HEAD CASE. fig:comparators sets "(a) Is it there?" and the other
%    four set "(a) where the drug name enters". comparators' heads are
%    interrogative SENTENCES and sentence case is correct for a sentence; the
%    v4 header argued this and the argument stands. fig:mechanism's lowercase
%    heads are noun phrases, not questions. One capital letter, four times.
%  * PANEL-HEAD POSITION. fig:comparators puts its heads in a right-aligned
%    left stub, the other four above the panel at its x origin. Moving them
%    above the rules costs about 20mm of height against a 118mm cap that the
%    figure already meets at 117.60mm. Not available.
%  * TWO-COLUMN GUTTER. purify and hook split at x = 8.40cm, licensing at
%    7.75cm. licensing cannot move right: "(d) what a negative would buy" is
%    4.85cm wide and its origin is already at 12.40 of a 17.30cm measure, with
%    1.5pt of slack. Moving purify and hook LEFT to 7.75 would be arbitrary.
% ===========================================================================
%
% CHANGED IN THIS FILE. (i) The prompt strip adopts fig:hook's cell height,
% 0.42cm, and fig:hook drops to this file's cell WIDTHS, so the two drawings
% of the one object are congruent (R8); the strip's top stays at y = 2.14, so
% the caret, the layer axis and everything above them are untouched, and the
% strip's bottom rises from 1.52 to 1.72. Labels move from anchor=base at
% y = 1.775 to centred at the cell's own mid-height, 1.92, with inner sep 0,
% which is \tokbox's placement. (ii) The drug cell's label loses \bfseries: it
% was the only bold glyph on the slate, and the accent colour plus the accent
% border already carry the emphasis. (iii) The layer ticks go from 0.12 to
% 0.10cm and the numerals from 6.70 to 6.74, a 0.06cm gap (R1). (iv) The
% "480 real tier-2 prompts" line goes ink -> quietink, matching the same class
% of line in fig:purify, and rises to y = 1.38 with the strip. (v) The prose
% block moves from y 5.00/4.60 to 5.55/5.15 so its first line sits on panel
% (b)'s second-mode baseline, 5.55, giving the two columns a shared
% horizontal. (vi) The foot line takes the slate's declaration form (R7).
% NOT FIXED, and it is this figure's real weakness: panel (a) is an L -- a
% 4.20cm vertical layer axis at the column's right edge and a 6.90cm strip
% along its bottom -- so about 6.0 x 3.4cm of the column is empty. The pass-12
% header records four rejected fills and this pass adds no fifth; the void is
% the (position x depth) plane in which this figure measures nothing, and the
% prose block sits in it to break it rather than to disguise it.
% RE-MEASURED: 173.10 x 74.30mm; ink 5.975%; 726 characters, 5.65 per 100mm^2,
% which is over S15's 5.0 and is the one contract clause this figure still
% misses -- see pass 12's arithmetic, unchanged by this pass. Exactly two type
% sizes, 9.96 and 7.97, in one face, no bold, no italic, no sans, no CM.
% One page, zero Overfull, zero Underfull.
% ===========================================================================
\begin{tikzpicture}[x=1cm, y=1cm, font=\small, inner sep=0pt]
  \tikzset{
    % 4.11's panel head: lowercase parenthesised letter, sentence-case
    % descriptive title, \footnotesize roman ink, left on the panel's origin.
    ttl/.style ={font=\footnotesize, text=ink, anchor=base west},
    % the figure's ONE prose register: definitional only, \footnotesize ink.
    def/.style ={font=\footnotesize, text=ink, anchor=base west},
    % the surviving branch -- the object the rest of the chapter measures.
    % PASS 13: \scriptsize, not \footnotesize. These are the two branches of a
    % partition, i.e. direct labels of drawn objects, not definitional prose;
    % under S1 the largest size in the figure must belong to a panel title, an
    % axis label or definitional prose, and 4.11 sets its in-plot labels one
    % step below its titles. They now share a size and a baseline grid with
    % their own outcome labels, which is what fixed the cross-gutter splay.
    liv/.style ={font=\scriptsize, text=accent, anchor=base west},
    % the struck branch, and every counterfactual.
    ded/.style ={font=\scriptsize, text=quietink, anchor=base west},
    % \scriptsize registers: labels, entries, scope and foot lines.
    lab/.style ={font=\scriptsize, text=ink, anchor=base west},
    qlab/.style={font=\scriptsize, text=quietink, anchor=base west},
    alab/.style={font=\scriptsize, text=accent, anchor=base west},
    tick/.style={font=\scriptsize, text=quietink, anchor=east},
    % The prompt cell is drawn as figs/hook.tex draws it (\tokbox): panelbg
    % fill, rounded corners=1pt. PASS 12: the border is 0.30pt, the document's
    % furniture weight, not the 0.35pt of pass 9, which is not one of the five
    % weights the slate uses.
    cell/.style={draw=rulegrey, line width=0.30pt, fill=panelbg,
                 rounded corners=1pt},
    % The two outcome connectors were IDENTICAL plain leaders. PASS 13 deleted
    % the Stealth heads: a single arrowhead is reserved for an operation acting
    % in a direction and its label must be a verb of action, and "struck by the
    % positive" and "survives" are outcomes, not operations.
    % PASS 14 DELETES lead/.style AND BOTH LEADERS. Reasoning at the two mode
    % blocks and in the pass-14 header at S3. The figure now contains no
    % arrowhead and no connector of any kind, and its only glyph that could be
    % mistaken for either -- the caret -- carries a direct label at itself.
    }

% ---------------------------------------------------------------------------
% ONE X ORIGIN PER PANEL. Every element below references its panel's origin,
% so the spread of text-line x0 inside a panel is exactly zero.
% ---------------------------------------------------------------------------
  \coordinate (OA) at (0.00,0);    % panel (a)
  \coordinate (OB) at (7.75,0);    % panels (b) and (c)
  \coordinate (OD) at (12.40,0);   % panel (d)
  \coordinate (OBL) at (13.44,0);  % panel (b)'s outcome column

% PASS 13: THE VERTICAL PANEL DIVIDER IS DELETED. It stood at x = 7.40 and the
% depth axis stands at x = 6.90, so the two were parallel hairlines 0.50cm
% apart with an empty channel between them; measured on the render, the band
% between them carried 0.16% ink over the full height. At print size they read
% as one doubled boundary, and panel (a)'s only scaled object was thereby read
% as frame furniture. The lettered panel heads and a 0.85cm gutter partition
% the figure without it, which is 4.11's own arrangement -- it draws no
% dividers at all. SUPERSEDES C5's divider clause and its gutter arithmetic.

% ###########################################################################
% # PANEL (a) -- where the drug name enters
% ###########################################################################
  \node[ttl] at (0.00,7.45) {(a) where the drug name enters};

  % --- the figure's only prose block: definitional, two lines -------------
  % PASS 13, TWO REPAIRS.
  %  (i) THE PRONOUN HAD NO ANTECEDENT ON THE READING LINE. Pass 12 read "what
  %      reads it: ...", and the only occurrence of "residual stream" was
  %      inside the rotated \scriptsize scope line, 4cm away and at 90 degrees.
  %      The block now names the object it defines. "twelve classes" leaves the
  %      block: the line "a cross-validated linear probe, twelve classes"
  %      measures 6.68cm and the panel's prose is capped at 5.48cm by the
  %      rotated label column at x = 5.62 (see the axis block below). "linear"
  %      is the word worth the space, because it is what makes panel (d)'s
  %      "the drug not linearly readable" the right counterfactual, and twelve
  %      still reaches the reader through the count line under the strip.
  % (ii) IT MOVED FROM THE PANEL'S VERTICAL MIDDLE TO THE PANEL'S HEAD, so that
  %      panels (a) and (b) put their first content line on ONE baseline, 6.70,
  %      0.75cm under their titles. Pass 12 set it at 5.90/5.50, which gave
  %      panel (a) a 1.55cm head-to-content gap against panel (b)'s 0.75cm and
  %      put panel (a)'s only substantial ink level with panels (c) and (d), so
  %      the lettering order and the spatial order disagreed.
  %      SUPERSEDES C10's placement paragraph ("three tiers with voids of 1.55
  %      and 3.22cm"); C10's four rejected ways of REMOVING the band all still
  %      stand, and the band is still deliberately empty, because it is the
  %      space between the prompt and the read, i.e. the residual stream.
  %      Clearance measured: line one is 4.880cm and line two 4.592cm, both
  %      clear of the rotated scope column's left edge at x = 5.62 by more than
  %      0.7cm at every y.
  % PASS 14: THE BLOCK MOVES BACK INTO THE BAND, at 5.00/4.60, centred on the
  % axis midspan 4.80. Pass 13's (ii) above bought a shared first-content
  % baseline with panel (b) and paid for it with ONE CONTIGUOUS VOID of
  % 56.2 x 39.2mm -- 17.2% of the figure's whole area, measured at 200dpi --
  % under the block and left of the axis, so panel (a) read as text at the top,
  % a stranded vertical stick in the middle and a hole between them. The two
  % TITLE baselines still align at 7.45, which is the alignment a reader
  % actually sees; the first-content alignment across a 7.75cm gutter is not.
  % Set here the block also sits level with the axis it describes, so "the
  % residual stream" has the rule and its seven ticks beside it rather than
  % 4cm above them. The band is now two voids of about 1.9 and 1.6cm rather
  % than one of 3.5cm, and its lower part carries the caret's direct label.
  % SUPERSEDES pass 13's repair (ii) at this block and C10's placement
  % paragraph as re-stated there; C10's four rejected ways of REMOVING the
  % band all still stand, and S4's "at most one block per panel" is why no
  % second prose block was written into it.
  \node[def] at (0.00,5.55) {what reads the residual stream:};
  \node[def] at (0.00,5.15) {a cross-validated linear probe};

  % --- the depth axis ------------------------------------------------------
  % PASS 12: a plain 0.40pt ink rule from the layer-2 tick to the layer-16
  % tick, terminating exactly on them. Before this pass it was a panelbg-filled
  % 0.36cm-wide column running from 2.10 to 7.05, i.e. past both extreme ticks,
  % which read as a bar whose length encoded nothing and implied depth beyond
  % the measured range.
  \draw[ink, line width=0.40pt] (6.90,2.70) -- (6.90,6.90);
  % y = 2.10 + (L/16)*4.80 for every one of the seven, unchanged from pass 11
  % and CHECKED against the map rather than against the eye:
  %   L   2    4    6    8    9    12   16
  %   y  2.70 3.30 3.90 4.50 4.80 5.70 6.90
  % Tick length 0.12cm = 3.4pt, PASS 13, down from 0.16cm: 4.11's own ticks are
  % about 3.5pt and the 0.04cm bought here is what lets the quantity name sit
  % beside the numerals without touching them.
  \foreach \L/\y in {2/2.70, 4/3.30, 6/3.90, 8/4.50, 9/4.80, 12/5.70, 16/6.90}{
    \draw[rulegrey, line width=0.30pt] (6.80,\y) -- (6.90,\y);
    \node[tick] at (6.74,\y) {\L};}
  % AXIS LABEL, in axis-label position: rotated 90 degrees reading upward,
  % centred on the axis midspan y = 4.80. PASS 13, TWO REPAIRS.
  %  (i) THE RANKING IS INVERTED BACK. Pass 12 put the quantity name OUTERMOST,
  %      at x = 5.36, with two wrapped \scriptsize grey lines between it and the
  %      numerals, so the thing immediately left of the tick numerals was a grey
  %      sentence and the quantity itself was 1.4cm further out -- the opposite
  %      of what S5 asks for. "layer" now sits immediately left of the numerals
  %      and the scope note outboard of it.
  % (ii) THE SCOPE NOTE IS ONE LINE, NOT A WRAPPED TWO. Pass 12's note measured
  %      9.09cm, wrapped onto two typeset lines of 4.817 and 4.202cm, and the
  %      longer of the two OVERRAN a 4.20cm axis at both ends -- 3.0mm past the
  %      layer-16 tick above and the same below -- which is the exact defect D6
  %      records itself as having removed from the column. It also broke the
  %      noun phrase across the wrap ("the last prompt / position"). The note is
  %      now "read at the last prompt position", 111.20pt = 3.908cm, which lies
  %      INSIDE the 4.20cm axis with 0.146cm of clearance at each end. "where
  %      generation begins" is caption material and is in the caption.
  %      SUPERSEDES C1's wrap arithmetic and its "line 1 is the outermost"
  %      clause in full.
  % x geometry, measured: at inner sep=0 a rotated \footnotesize line is 0.35cm
  % wide in x and a \scriptsize one 0.28cm; the widest numeral ("16") is 8.0pt =
  % 0.281cm and is right-anchored on 6.70, so it begins at x = 6.419.
  % scope 5.62-5.90 | layer 6.005-6.355 | numerals 6.419-6.70 | ticks 6.78-6.90.
  % PASS 14, TWO REPAIRS AT THIS BLOCK.
  %  (i) THE SCOPE LINE IS DELETED AND ITS CONTENT MOVES TO THE MARK IT NAMES.
  %      "read at the last prompt position" shared the five-word clause "at the
  %      last prompt position" with the caption that actually prints
  %      (04b:124-126, "the probe reads the residual stream at the last prompt
  %      position, where generation begins"), which S18 forbids; and it named,
  %      at 90 degrees and 1.14cm away, the one glyph on the page that had no
  %      name of its own -- the caret. 4.11's discipline is to direct-label at
  %      the mark, so the position is now labelled AT the caret, horizontally,
  %      in \scriptsize ink, and the axis keeps only its quantity name. S5
  %      makes the scope line optional; the mandatory \footnotesize roman ink
  %      quantity name in axis-label position is still there and is still the
  %      only thing between the tick numerals and panel (a)'s body.
  %      SUPERSEDES pass 13's repair (ii) at this block in full.
  % (ii) "layer" MOVES OFF THE LAYER-9 TICK. The axis midspan and the layer-9
  %      position are the SAME y by arithmetic, not by accident:
  %      2.10 + (9/16)*4.80 = 4.80 = (2.70 + 6.90)/2. So a rotated label
  %      centred on the midspan was centred on one numeral, 1.96mm from it, and
  %      read "layer 9". Measured, "layer" is 21.46pt = 0.754cm long, and the
  %      seven tick gaps are 0.60/0.60/0.60/0.30/0.90/1.20cm, so 5.70-to-6.90
  %      is the ONLY stretch of this axis long enough to hold the label with no
  %      numeral inside its span. Centred at 6.30 it spans 5.923 to 6.677 and
  %      the nearest numerals are 12 at 5.70 and 16 at 6.90. It stays in
  %      axis-label position: x 6.005-6.355, immediately left of the numeral
  %      column at 6.419-6.70, inside S5's 12pt band either side of that column.
  \node[font=\footnotesize, text=ink, rotate=90, anchor=center]
    at (6.18,6.30) {layer};

  % --- the prompt strip ----------------------------------------------------
  % Edges from the measured field-name widths: 6.084cm of ink, 0.816cm of
  % padding shared as 0.0816cm a side. See C8.
  % PASS 14: THE WHOLE STRIP DROPS 0.14cm, from y 1.66-2.28 to y 1.52-2.14.
  % NO CELL EDGE IN x MOVED, no field changed, no order changed and the strip
  % is still 6.90cm wide with the same five measured widths -- the change is
  % one rigid translation in y, and it buys the 0.14cm the caret's direct label
  % needs between the strip's top rule and the layer-2 numeral (see the read
  % position below). The room exists because pass 14 deletes the span rule that
  % occupied 1.46-1.54; the count line's baseline at 1.18 is unmoved and its
  % ascenders now clear the strip's bottom rule by 0.14cm.
  \draw[cell] (0.00,1.72) rectangle (1.12,2.14);
  \draw[cell] (1.90,1.72) rectangle (2.64,2.14);
  \draw[cell] (2.64,1.72) rectangle (4.23,2.14);
  \draw[cell] (4.23,1.72) rectangle (6.90,2.14);
  % the drug cell: the only accent mark in panel (a). Same construction as
  % hook.tex's \tokdrug -- accent border, accent!12 fill, rounded corners=1pt.
  \filldraw[fill=accent!12, draw=accent, line width=0.40pt, rounded corners=1pt]
    (1.12,1.72) rectangle (1.90,2.14);

  % baseline 1.775 = the strip's vertical centre 1.83 less 0.055, which centres
  % the x-height band rather than the box. (Pass 13: 1.915 on a centre of 1.97.)
  \node[font=\scriptsize, text=ink,    inner sep=0pt] at (0.560,1.92) {cell line};
  \node[font=\scriptsize, text=accent, inner sep=0pt] at (1.510,1.92) {drug};
  \node[font=\scriptsize, text=ink,    inner sep=0pt] at (2.270,1.92) {dose};
  \node[font=\scriptsize, text=ink,    inner sep=0pt] at (3.435,1.92) {mechanism};
  \node[font=\scriptsize, text=ink,    inner sep=0pt]
    at (5.565,1.92) {control cell sentence};

  % --- the read position ---------------------------------------------------
  % The caret straddles the strip's right edge, which is where the prompt ends.
  % It is ink, not accent: it is part of the axis apparatus, and accent in this
  % figure means the drug field and the surviving branch and nothing else.
  % PASS 13: IT NO LONGER TOUCHES THE AXIS AND NO LONGER SITS ON A NUMBERED
  % TICK. Pass 12 put its apex at (6.90, 2.70), which is simultaneously the
  % axis rule's base, the layer-2 tick and the largest filled black glyph in
  % panel (a); measured on the render the apex and the layer-2 tick coincided
  % to within a pixel. Three objects at one coordinate read two ways, both
  % false: as a filled arrowhead terminating the depth axis downward, and as a
  % mark AT LAYER 2 on the figure's only scaled axis -- which also falsified
  % the foot line, since the caret would then occupy a position on the depth
  % scale. It marks a token position, not a depth. The apex is now 2.50, a
  % clear 0.20cm below the axis base at 2.70, and the triangle is back to the
  % 0.24cm width of pass 9 (base 6.78 to 7.02, so it still straddles the
  % strip's right edge at 6.90 symmetrically and overhangs it by 0.12cm rather
  % than 0.15cm). It is named by the axis scope line 0.9cm to its left, which
  % now reads "read at the last prompt position".
  % SUPERSEDES C8's clause "the apex is the axis rule's own base at the layer-2
  % tick, y = 2.70", and restores the pass-9 note's own reading that the apex
  % belongs at a position the depth scale does not number.
  % PASS 14: THE CARET IS NAMED AT ITSELF, AND THE THREE CUES THAT MADE IT AN
  % ARROWHEAD ARE BROKEN ONE BY ONE. Pass 13 moved it 2mm and did not resolve
  % the reading; all three refuters read it, unaided, as the axis's terminal
  % arrowhead, and two of them read the layer-2 numeral as its label. Measured
  % on the pass-13 render, three cues did that: (1) it was the only
  % arrowhead-shaped glyph left in the figure and the largest solid black one
  % in the panel; (2) its apex stood on the axis rule's own x, 2.00mm below the
  % rule's bare square terminus; (3) its base ran 6.78 to 7.02 while the seven
  % ticks run 6.78 to 6.90, so its left half lay exactly in the tick lane, and
  % the "2" numeral's box bottom sat 0.25mm above its apex. Pass 14:
  %   * it acquires a \scriptsize ink DIRECT LABEL at itself, "the last prompt
  %     position", right-anchored at 6.68 so its ink stops 0.08cm from the
  %     caret's left base -- inside the document's ~0.3cm leader threshold, so
  %     no leader is owed. This is where that string was always meant to be:
  %     it was set rotated in axis-label position until pass 14, where a reader
  %     attached it to the axis rather than to the triangle;
  %   * the strip drops 0.14cm, so the apex is now 0.34cm below the axis base
  %     (was 0.20cm) and the label sits in the gap that opens;
  %   * the base widens from 0.24 to 0.28cm, 6.76 to 7.04, so it no longer
  %     shares an endpoint with the tick lane. It still straddles the strip's
  %     right edge at 6.90 symmetrically, which is the pass-9 relation the keep
  %     list names, and it is still ink, not rulegrey (O3 stands: rulegrey is
  %     furniture and this mark is content).
  % Under the figure's own map y = 2.10 + (L/16)*4.80 the triangle now occupies
  % y 2.14 to 2.36, i.e. layers 0.13 to 0.87 -- BELOW layer 1, off the drawn
  % scale, which is what "it marks a token position, not a depth" requires and
  % what the pass-13 ledger claimed without arithmetic. SUPERSEDES pass 13's
  % apex/base clause at this block and licensing.json's read_position note.
  \fill[ink] (6.90,2.36) -- (6.76,2.14) -- (7.04,2.14) -- cycle;
  \node[font=\scriptsize, text=ink, anchor=base east] at (6.68,2.28)
    {the last prompt position};

  % --- the span: the five fields are one prompt ----------------------------
  % PASS 13: A PLAIN SPAN RULE WITH DOWNTURNED ENDS REPLACES THE MIRRORED
  % BRACE. A brace has a central pointer, and pass 12 moved the count line from
  % centred-on-the-brace to left-aligned on the panel origin (S13) without
  % moving the pointer, so the notch at x = 3.45 aimed at 3.45cm of empty white
  % while its label began at x = 0.00. A span rule carries the same "these five
  % fields are one prompt" grouping, has no pointer to aim wrongly, and drops
  % the decorations.pathreplacing dependency.
  % PASS 14: THE SPAN RULE IS DELETED. All three refuters read it and all three
  % called it. Measured on the pass-13 render it was a rulegrey 0.30pt rule of
  % the strip's exact x-extent, 0.12cm below and parallel to the strip's own
  % rulegrey 0.30pt bottom border -- same colour, same weight, same length --
  % so it reproduced as a doubled edge or an offset re-impression rather than
  % as a bracket, its 0.08cm downturns vanishing at print size. Its left
  % downturn also ended 0.78mm directly above the "4" of "480" and read as a
  % stray tick on the numeral. And it grouped ONE schematic prompt's five
  % fields while the line beneath it counts 480 prompts, so on a strip that
  % carries no horizontal scale a full-width bracket labelled 480 invited
  % reading the width as the count. Nothing is lost: the five cells abut and
  % already read as one object, and the count line sits under the strip on the
  % strip's own left edge and names the population it is drawn from.
  % SUPERSEDES pass 13's R9 span-rule clause and the "the five fields are one
  % prompt" heading above it; the 0.14cm it frees is what lets the strip drop.
  % The count line names the population the spanned object is drawn from. It is
  % left-aligned to panel (a)'s origin, so panel (a) has exactly one x origin.
  % \scriptsize: at \footnotesize it measures 8.196cm and would run past the
  % panel's own measure. PASS 13: 480 is set as TEXT, not as $480$. In math
  % mode the three digits rendered in URWPalladioL-Roma while every other
  % numeral in the figure -- the seven tick labels, and the 2 of "tier-2" on
  % this very line -- rendered in TeXGyrePagella-Regular. One figure, two digit
  % faces, three characters. The measured width is 186.58pt = 6.5575cm either
  % way, so no coordinate moved and no value changed.
  \node[qlab] at (0.00,1.38)
    {480 real tier-2 prompts: twelve drugs, forty cells each};

% ###########################################################################
% # PANEL (b) -- the two ways a model can ignore a drug it was told about
% ###########################################################################
  % PASS 14: THE TITLE IS A NAME, NOT A SENTENCE. "the two ways a model can
  % ignore a drug it was told about" measured 270.16pt = 9.495cm, which made a
  % TITLE the widest object on the page, set the figure's whole right measure,
  % left 0.055cm of optical margin, and outweighed panel (a)'s title three to
  % one on the baseline they share, so the two heads did not read as peers. It
  % also restated 04b:31-32's own opening sentence ("A model that ignores a
  % drug it was told about can fail in two ways"), which is the paragraph the
  % float is attached to. 4.11's two titles run 29 and 38 characters; this one
  % ran 60. "two ways to ignore a named drug" is 165.62pt = 5.821cm, keeps the
  % "was told about" sense in the word NAMED -- which is also what panel (a)'s
  % title says enters -- and ends at x = 13.571. The full sentence is in the
  % caption. SUPERSEDES pass 13's O2 in its measurement clause: panel (a)'s
  % title is no longer constrained by this one. It is left at 4.969cm anyway,
  % because every longer variant either overruns panel (a)'s own 6.90cm measure
  % or says the probe reads the drug name, which it does not -- it reads the
  % stream, and the read is named at the caret and at the axis.
  \node[ttl] at (7.75,7.45)
    {(b) two ways to ignore a named drug};

  % --- failure mode one: struck --------------------------------------------
  % quietink, because it is eliminated. TWO strikes, one through each line: the
  % entry is two lines and a rule through its first line only left the second
  % standing. Each overhangs its line's ink by 0.06cm at each end, which is what
  % a strikethrough does. Endpoints from the measured \scriptsize widths,
  % 4.6741 and 4.2918cm: 7.69 to 12.484 and 7.69 to 12.102.
  % PASS 13, THE COLOUR AND THE WEIGHT. The strikes were ink at 0.90pt, and
  % that is the register S6 and CF2 reserve, across the whole slate, for a
  % threshold: chance in 4.11(a), the removal gate in 4.9, the pre-registered
  % materiality margin in 4.12 and 4.13, and style_v2.chance_line()'s own
  % hard-coded color='black', lw=0.9. C6's defence was that this figure
  % contains no threshold so nothing collides -- which answers S10's singleton
  % clause but not S6's, whose check is that NOTHING WHICH IS NOT A THRESHOLD
  % may be black solid 0.90pt. This figure prints six pages before the reader
  % meets the first real threshold in the arc, and it is the first figure of
  % the slate, so it was teaching the threshold stroke on a negation. Measured
  % against it: the struck text is quietink (grey 107) and the strikes were ink
  % (grey 25), so the branch the chapter DISCARDS carried the darkest ink on
  % the page while the surviving branch receded -- the S8 inversion C3 fixed in
  % colour, reintroduced in weight. Both strikes are now quietink 0.40pt, the
  % colour and weight of the text they cross. 0.40pt is already live here (the
  % axis rule and the drug cell's border), so no new register appears and no
  % singleton is created; 0.50pt was rejected because S10 reserves it for a
  % dotted coordinate. SUPERSEDES C6's third weight in full: the figure's
  % weights are now 0.30pt furniture and 0.40pt, and no 0.90pt stroke exists.
  % Strike y is the line's baseline + 0.065, half the \scriptsize x-height.
  % PASS 14, TWO REPAIRS.
  %  (i) THE STRIKES START ON THE PANEL'S OWN ORIGIN. They began at 7.69, which
  %      is 0.06cm (1.7pt) LEFT of the 7.75 every other text line in panels (b)
  %      and (c) starts on, so the panel's left margin was ragged and the item
  %      protruding furthest left was the one the figure discards. A
  %      strikethrough needs overhang at its TERMINAL, where the eye leaves the
  %      word, not at its initial; the right ends are unchanged at 12.484 and
  %      12.102, still 0.06cm past their lines' ink at 12.424 and 12.042.
  % (ii) THE OUTCOME LEADER IS DELETED (both of them; see the surviving mode).
  %      Measured on the pass-13 render this one ran 12.72 to 13.34 with a
  %      0.24cm gap at its left end and a 0.10cm gap at its right, so it
  %      touched neither the strike nor the label and read as a free-floating
  %      dash. Worse, it was a THIRD quietink horizontal lying midway between
  %      two quietink strikes 0.35cm apart, at the same colour and near the
  %      same weight, so the eye completed strike + gap + leader + label into
  %      one continuous rule and read the verdict "struck by the positive," as
  %      itself struck -- the exact "strike and leader read as one continuous
  %      pointer" fault the pass-9 note records itself as having killed on the
  %      diagonal. Widening the gap did not kill it because the stroke was
  %      still collinear in band, colour and register. Nothing is owed under
  %      S22: there is no MARK here, only a text block, and the pairing is
  %      carried twice over -- by colour (grey block to grey label, accent
  %      block to accent label) and by baselines shared to 0.00pt, which is
  %      4.11's own discipline. SUPERSEDES pass 13's lead/.style rationale and
  %      C6's "the two outcome leaders" clause: the figure now has one stroke
  %      weight for furniture (0.30pt) and one for the strikes and the axis
  %      (0.40pt), and no connector of any kind.
  \node[ded] at (7.75,6.70) {the drug never enters its computation};
  \node[ded] at (7.75,6.35) {an activation-side encoding failure};
  \draw[quietink, line width=0.40pt] (7.75,6.765) -- (12.484,6.765);
  \draw[quietink, line width=0.40pt] (7.75,6.415) -- (12.102,6.415);
  \node[qlab] at (13.00,6.70) {struck by the positive,};
  \node[qlab] at (13.00,6.35) {and only this one};

  % --- failure mode two: survives ------------------------------------------
  % accent, because it is the object the rest of the chapter measures. The two
  % leaders are identical in weight; only the label differs.
  % PASS 13: EACH OUTCOME LABEL NOW SITS ON ITS MODE'S OWN BASELINES, 6.70/6.35
  % and 5.55/5.20. Pass 12 set the modes \footnotesize on 0.40cm leading and
  % their labels \scriptsize on 0.30cm leading, each block centred on the other,
  % so no line matched its partner and the pair splayed in opposite directions
  % across a 55mm gutter -- measured, 0.41mm apart at the first line and 0.60mm
  % the other way at the second. One leading per size now holds across the whole
  % figure: 0.40cm at \footnotesize, 0.35cm at \scriptsize.
  % The outcome column moved from 14.28 to 13.44, which is what the \scriptsize
  % modes freed. Its leader starts 0.50cm clear of the longer strike's terminal
  % (12.484), against 0.17cm in pass 12, so strike and leader cannot read as one
  % continuous pointer -- the fault that killed the pass-1 diagonal.
  % PASS 14: THIS LEADER IS DELETED TOO, AND IT WAS THE WORSE OF THE PAIR.
  % Measured on the pass-13 render, the accent mode's longest line ends at
  % x = 11.328 and the leader began at 12.72, so it started 1.397cm -- 39.6pt,
  % 14mm -- past the mark it claimed to join and covered 0.61 of a 2.13cm gap:
  % at print size a 6.1mm dash floating in 14mm of white with nothing on its
  % left. The pass-13 ledger recorded the OPPOSITE as a measured property
  % ("the leader now starts 0.236cm past the longer strike's terminal"), which
  % is true of the struck mode's leader only; the accent one's 1.397cm gap was
  % never measured. Two leaders of visibly unequal reach were the alternative
  % and were rejected: the struck one cannot be extended leftward without
  % butting the strike and restoring the continuous-pointer reading.
  % The outcome column moves 13.44 -> 13.00, which halves the unlinked gap on
  % this mode from 2.11 to 1.67cm and still clears the longer strike's terminal
  % (12.484) by 0.516cm, the 0.5cm separation the pass-9 note asks for.
  % SUPERSEDES pass 13's O4 arithmetic (13.44) and its "the two leaders are
  % identical in weight" clause -- there are no leaders.
  \node[liv] at (7.75,5.55) {it enters and nothing reads it};
  \node[liv] at (7.75,5.20) {a decoder failure};
  \node[alab] at (13.00,5.55) {survives, and is the rest};
  \node[alab] at (13.00,5.20) {of the chapter};

% ###########################################################################
% # PANELS (c) AND (d) -- the two licences, as two columns read across
% ###########################################################################
  % PASS 14: THE WHOLE (c)/(d) GROUP DROPS 0.30cm AND ITS ROW PITCH OPENS FROM
  % 0.52 to 0.70cm. Deleting the third row (see the entries below) left the
  % right half ending at y = 2.55 against panel (a)'s count line at 1.18, i.e.
  % a fresh 9.5 x 1.8cm void in the bottom-right corner -- trading one hole for
  % another is not a repair. Moving the rule to 4.05 also makes the gap from
  % panel (b)'s last line (5.20) to the rule 1.15cm, which is exactly the gap
  % between panel (b)'s own two mode blocks, so the right half now has one
  % vertical rhythm. The title still sits 0.45cm under the rule, as before.
  \draw[rulegrey, line width=0.30pt] (7.75,4.05) -- (17.30,4.05);

  % PASS 13: BOTH TITLES ARE ONE LINE ON ONE BASELINE. (d) read "what a
  % negative would have bought", 6.167cm, which does not fit a column and
  % wrapped, so (c)'s head-to-body gap was 0.90cm against (d)'s 0.54cm on two
  % columns the ledger called equal. "would buy" is 4.847cm and the conditional
  % still carries the counterfactual against (c)'s indicative "buys"; the
  % perfect tense survives in the caption. Origins are 7.75 and 12.40: (d)'s
  % title then ends at 17.247, which is 0.053cm inside the measure, exactly the
  % clearance panel (b)'s title already runs at. Columns are 4.65 and 4.90cm,
  % and the content footprints are 4.081/4.847 (titles) and 3.210/3.697
  % (entries), so neither column is the larger and more detailed one -- which is
  % what C4's equalisation was for. SUPERSEDES C4's "equal width (4.60cm)",
  % which was never true of the drawing: pass 12's origins were 7.75 and 12.70,
  % i.e. 4.95 and 4.60.
  \node[ttl] at (7.75,3.60) {(c) what the positive buys};
  \node[ttl] at (12.40,3.60) {(d) what a negative would buy};

  % THREE ENTRIES EACH, PAIRED ROW BY ROW. PASS 13, AND THIS IS THE CORRECTNESS
  % REPAIR IN THIS PANEL, NOT A TYPOGRAPHIC ONE.
  %  (i) THE ACROSS-READING WAS OFF BY ONE. Pass 12 set four baselines in each
  %      column and asserted, in the ledger and in the caption, that they were
  %      to be read across entry by entry. They could not be: (c) was two
  %      sentences word-wrapped over four lines and (d) a single four-item
  %      comma list, so row 1 paired "the model did not discard" with "the drug
  %      lost inside the model," and the negation of the layer sat one row
  %      above the layer it negated. Three of the four pairings were nonsense.
  %      The rows are now grouped, not merely stacked: 0.32cm of leading WITHIN
  %      a row and 0.52cm BETWEEN rows, so a wrapped entry cannot be mistaken
  %      for the next entry. Each row's first line is on one baseline in both
  %      columns, and the three pairs are exact.
  % (ii) "no class of failure," IS DELETED, because it contradicts this
  %      figure's own panel (b) and the claim box it is drawn from. 04b:142-144
  %      reads "It rules out that the drug field is lost before the decoder
  %      runs, an encoding failure in the activations", and panel (b) says so
  %      on the page, 3.5cm above, in the words "struck by the positive, and
  %      only this one". A positive that strikes one of two modes has bought a
  %      class of failure ruled out; what it has not bought is a depth and a
  %      fix. The two entries that survive, "no layer" and "no class of fix",
  %      are the two the sources license: 04b:110-112, "what the probe buys is
  %      the depth at which one would appear", and the caption's own lede, that
  %      the negative "would have named the layer at which the drug was lost
  %      and the class of fix that could reach it".
  % (iii)"an activation-side failure, and" IS DELETED FROM (d), because panel
  %      (b)'s struck line already reads "an activation-side encoding failure"
  %      2.5cm away. One object, one name, drawn once; the fix classes stay,
  %      because they are named nowhere else in the drawing.
  % (iv) (d)'s first entry carries the hedge the instrument requires. The probe
  %      is linear -- panel (a) now says so -- and a linear decode failure
  %      licenses only that the drug is not linearly readable, not that it is
  %      absent. The claim box is careful about exactly this ("Its identity is
  %      linearly readable from the residual stream at every layer measured");
  %      pass 12's "the drug lost inside the model," was not, and since the
  %      figure's whole argument is a comparison of licence SIZES, an
  %      overstated counterfactual inflates the asymmetry it exists to draw.
  %      SUPERSEDES the pass-9 note "THE COUNTERFACTUAL BLOCK'S FIRST ENTRY
  %      LOST ITS FINITE VERB" as to its wording only; the reason that note
  %      gives -- that the entry must be a noun phrase, not an assertion that
  %      the drug IS lost -- is preserved and strengthened by the hedge.
  % (c) is ink: it happened. (d) is quietink: it did not.
  % Leading is 0.35cm WITHIN a row -- the one \scriptsize leading in the figure,
  % shared with panel (b) -- and 0.52cm BETWEEN rows.
  % PASS 14, THREE ENTRIES BECOME TWO PAIRED ROWS OF ONE LINE EACH, AND ONE OF
  % THE DELETIONS IS A CORRECTNESS REPAIR, NOT A TYPOGRAPHIC ONE.
  %  (i) "no class of fix" IS DELETED, BY EXACTLY THE ARGUMENT PASS 13 USED TO
  %      DELETE "no class of failure," ONE ROW ABOVE IT -- and that argument
  %      was never applied to this entry, in the header or in the ledger.
  %      04b:31-36 is a TWO-MEMBER partition and each member names its own fix
  %      class: mode one is "an activation-side encoding failure that target-
  %      and input-side fixes might reach", mode two "a decoder failure only an
  %      objective-side fix could reach". Panel (b) strikes mode one and says
  %      so on this page, 3.5cm above, in the words "struck by the positive,
  %      and only this one". If one member of an exhaustive pair is struck the
  %      survivor stands, and the survivor names its fix class -- so within
  %      this figure's own frame the positive DOES buy a class of fix, and the
  %      entry is false. It survived pass 13 only because panel (b) draws mode
  %      two's gloss truncated to "a decoder failure", dropping the very clause
  %      of its source sentence that refutes the entry. A drawing may not
  %      contradict, at 8pt, the box it cites as its source.
  %      The alternative repair -- rewriting the entry to the class the
  %      positive does buy, "objective-side fixes only", against (d)'s
  %      "target- and input-side fixes" -- was built and rejected: that licence
  %      holds only CONDITIONALLY, on the model in fact ignoring the drug,
  %      which is 4.4's and 4.5's finding and not this section's, and a
  %      two-word slot cannot carry the condition without overstating it. Both
  %      fix classes go to the caption, which can state the condition.
  %      (d)'s "target- and input-side fixes" goes with it, because a row with
  %      one side is not a row and the entry is named nowhere else it is owed.
  % (ii) ROW ONE LOSES ITS WRAP AND ITS S18 COLLISION. "the model did not
  %      discard what it was handed" is seven words verbatim from 04b:107-108,
  %      the paragraph the float is attached to ("recovering it from the
  %      activations establishes only that the model did not discard what it
  %      was handed"), six lines of source above the float -- which S18 forbids
  %      outright. It was also the only entry in either column that wrapped, so
  %      (c) read as three items against (d)'s two and the across-reading broke
  %      at row one: measured, the white gap between rows 1 and 2 was 6.76pt in
  %      (c) and 16.68pt in (d), a 2.47x difference in the same two rows.
  %      "the drug field is not discarded" is 106.48pt = 3.742cm, one line,
  %      inside (c)'s 4.65cm column, keeps the sense the positive licenses, and
  %      shares no five-word clause with the prose or with any caption.
  % (iii)WITH EVERY ENTRY ON ONE LINE THERE IS NO WITHIN-ROW LEADING LEFT, so
  %      the two rows sit at one pitch of 0.60cm in both columns and the
  %      pairing needs no rule and no cue. The asymmetry the figure exists to
  %      draw is now also visible as ink for the first time: (c) sets 39
  %      characters against (d)'s 58, and (c)'s second entry is 0.991cm against
  %      (d)'s 3.233cm, so the small licence is the small column.
  %      SUPERSEDES pass 13's (ii)/(iii) rationale at this block as to the
  %      entries that survive, and C4's "FOUR ENTRIES EACH" in full.
  \node[lab]  at (7.75,2.85) {the drug field is not discarded};
  \node[qlab] at (12.40,2.85) {the drug not linearly readable};

  \node[lab]  at (7.75,2.15) {no layer};
  \node[qlab] at (12.40,2.15) {the layer at which it is lost};

% --- the figure's one disclosure line ---------------------------------------
% PASS 13: REWORDED. "Nothing in this figure is a measured quantity" sat 0.9cm
% under a line printing 480, twelve and forty, and then exempted the depth
% ticks in its own second clause, so it read as a blanket claim followed
% immediately by two counterexamples. The counts are the probe's design, not
% its results, and saying so is shorter than conceding an exception. It also
% now carries the count "seven", which pass 12 deleted with the caret label and
% which the pass-9 layout note (still standing) argues is the one fact the
% depth column exists to carry: without it a reader cannot tell whether the
% ticks are the depths probed or a sample of a sweep over all sixteen.
% One line, 389.66pt = 13.695cm, left-aligned to the figure's x origin.
% PASS 14: THE SECOND CLAUSE GOES, AND THE FIRST IS WHAT WAS DOING THE WORK.
% Pass 13 added "the counts are the probe's design" because the pass-12 wording
% ("Nothing in this figure is a measured quantity") sat under a line printing
% 480, twelve and forty and read as a blanket claim with two counterexamples
% beneath it. But pass 13 made TWO changes at once, and it was the other one --
% "a measured quantity" becoming "an estimate" -- that removed the
% contradiction, since a count is not an estimate. The clause was therefore
% belt and braces on a line already sound, at 35 characters, in the longest
% string in the figure, on a drawing over S15's ink ceiling. At 76 characters
% the line still answers the acceptance question ("can the reader say, without
% leaving the figure, that nothing here is a measured quantity?") on its own
% and still carries the count "seven" the pass-9 layout note asks for. The line
% is now 9.60cm rather than 13.695cm and is no longer the widest \scriptsize
% object on the page. SUPERSEDES pass 13's R9 foot-line clause as to the second
% sentence only; its rewording of "measured quantity" to "estimate" stands.
  \node[qlab] at (0.00,0.38)
    {Intervals: none in any panel; nothing here is an estimate, and the seven
     depth ticks are the only scaled marks.};

  % the bounding box IS the fullwidth measure, less 1.8pt
  \path (0,0.16) rectangle (17.30,7.72);
\end{tikzpicture}

% ---------------------------------------------------------------------------
% READY TO PASTE into Sections/04b-representation.tex, at \beatsection
% sec:methods-probe, immediately after the paragraph "The probe, and why its
% positive is nearly free" (04b-representation.tex:108-113) and before
% "Forward passes only." Not written into the Section file by this build.
%
% THE PASS-11 CAPTION IS RETAINED BELOW, COMMENTED OUT AND MARKED SUPERSEDED,
% because the thesis is audited against these headers. It is FALSE against the
% pass-12 and pass-13 drawings in four places: it indexes the panels "Left,"
% and "Right," where the drawing sets four lowercase panel letters, it says the
% three accuracies and the floor are "set here as text" (they are not set here
% at all any more), it spends two sentences denying a reading the drawing no
% longer offers, and it says the two licence blocks "are sized by their
% contents", which was the footprint asymmetry pass 12 replaced.
%
% The PASS-13 REPLACEMENT IS WRITTEN TO figs/_PENDING_CAPTIONS.tex, keyed by
% \label{fig:licensing}, together with every sentence pass 12 and pass 13 cut
% from the drawing. PASS 13 CREATED THAT FILE: pass 12 asserted twice that it
% had, and it did not exist. This builder does not edit Sections/. The block
% below is a DUPLICATE of the pending file, kept so this file is self-contained
% and the audit trail is here; the pending file is the one an integrator reads.
%
% ###########################################################################
% PASS 14 SUPERSEDES THE WHOLE OF THE PARAGRAPH ABOVE, AND THE PASS-13 CAPTION
% BELOW IS NOW HISTORY, NOT THE PROPOSAL.
%
% (i) The paragraph above is FALSE AS WRITTEN. figs/_PENDING_CAPTIONS.tex did
%     exist when pass 13 ran (it was created by the pass-10 repair of
%     fig:materiality), and pass 13's APPEND TO IT DID NOT SURVIVE: read on
%     2026-08-19 the file held fig:materiality, fig:comparators and fig:hook
%     blocks and no fig:licensing block at all. All three pass-14 refuters
%     opened with that finding and all three were right. The likeliest cause is
%     a concurrent read-modify-write from another figure's builder, which is
%     the hazard CF12 names about this shared file.
% (ii) PASS 14 APPENDED the block, in place, to figs/_PENDING_CAPTIONS.tex, and
%     READ THE FILE BACK TWICE -- once immediately after the append and once at
%     the end of the pass. Both times the block was there and the
%     fig:materiality, fig:comparators and fig:hook blocks were all still
%     there. It was appended at the END rather than inserted in figure order
%     deliberately: an append cannot destroy another builder's block and an
%     ordered insert can, which is what happened to pass 13.
%     NO LINE NUMBER IS RECORDED FOR IT, and that is the point of the two
%     reads: between them another builder appended to the fig:comparators
%     block and every heading below it moved -- this block went from line 484
%     to 597 and fig:hook from 306 to 369 -- so a line number written into
%     this header would be false within the hour. The \label heading is the
%     stable address.
% (iii)THE PASS-14 CAPTION IS THE ONE IN THE PENDING FILE, and it is NOT
%     duplicated here. Keeping a second copy in this header is what let pass 13
%     believe it had staged a caption it had not; one copy, in the file an
%     integrator is told to read, is the safer arrangement. What is recorded
%     here instead is what the pass-14 caption differs from the pass-13 draft
%     below in, so the audit trail is complete without a duplicate:
%       * it prints NONE of the four accuracies, where the pass-13 draft set
%         all four. It names them ("the raw peak and the deepest, the
%         context-removed deepest and the $1/12$ floor") and points at
%         \cref{tab:workspace} and \cref{fig:mechanism}(a). This is the
%         cross-figure redundancy resolution applied to the caption as well as
%         to the drawing: 0.821 and 0.883 are owned by 4.11(a), the only place
%         the two series are drawn AS series with their crossing at layer 12
%         visible, and printing them a fourth time here is the redraw the
%         resolution exists to stop. The values stay in the prose at
%         04b:159-166, in the claim box, and in \cref{tab:workspace};
%       * it carries the two FIX CLASSES in full, with the premise they depend
%         on, because pass 14 deleted "no class of fix" and "target- and
%         input-side fixes" from panels (c) and (d);
%       * it restates each failure mode's gloss in full, because the drawing
%         abbreviates both;
%       * it drops "The two licence blocks are sized by their contents", which
%         the pass-13 draft had already dropped, and also drops "the figure
%         quantifies nothing about the size of a licence": with the boxes gone
%         and (c)/(d) set as plain text columns, nothing on the page reads as
%         an area, and a caption denying a reading the drawing no longer offers
%         is the defect that sentence was meant to fix;
%       * it says "the single position read, the last of the prompt" rather
%         than "at the last prompt position", because the drawing now sets
%         "the last prompt position" as the caret's direct label and S18 bars
%         a shared five-word clause.
%     MEASURED: 2554 characters, 22 typeset lines under the real preamble in a
%     fullwidth float with the pass-14 drawing above it, against the live
%     pass-11 caption's 1692 characters and 15 lines, and against the 19-22 the
%     rest of the 4.7-4.12 arc runs at. It compiles with zero errors and zero
%     Overfull boxes on one page.
% ###########################################################################
%
% ---- PASS 13 REPLACEMENT CAPTION, SUPERSEDED BY PASS 14, kept for the audit
% ---- trail. DO NOT PASTE THIS ONE. -----------------------------------------
% \begin{figure}[htbp]
% \begin{fullwidth}
% \centering
% % ===========================================================================
% licensing.tex -- fig:licensing, what the probe's positive licenses and what
% the negative that never arrived would have licensed.
% FIGURE BODY ONLY. No preamble, no document environment, no float, no caption.
% \input it from inside a figure/fullwidth float in
% Sections/04b-representation.tex, at \beatsection sec:methods-probe, straight
% after the paragraph "The probe, and why its positive is nearly free"
% (04b-representation.tex:108-113). The complete float is staged, commented, at
% the foot of this file.
%
% Compiles against thesis_v2/preamble.tex and nothing else: no \includegraphics,
% no external data, no TikZ library beyond the ones preamble.tex already loads
% (this file uses arrows.meta for the two Stealth heads and
% decorations.pathreplacing for the mirrored brace under the prompt strip).
%
% WHAT IT IS. Two panels. LEFT: where the drug name enters the prompt, and
% where the probe reads. RIGHT: the two ways a model can ignore a drug it was
% told about, one of them struck by the positive and the other left standing,
% beside the small licence the strike buys and the large one a failure would
% have bought. The figure draws an ASYMMETRY, not a result: the section's own
% discount, that a positive here was near-certain in advance while a negative
% would have been decisive, has no graphical form anywhere in the document, and
% nothing in thesis_v2 draws the prompt, its fields, or the read position.
%
% EVERY NUMBER DRAWN, AND WHERE IT COMES FROM
% -------------------------------------------
% Nothing in this figure is a plotted coordinate. Four accuracies appear as
% TEXT and seven layer indices are the only quantity with a POSITION. Sibling
% ledger of every one of them at machine precision: thesis_v2/figs/licensing.json.
%
%   set as text, panel A                       source
%   0.821  probe at layer 9        RESULTS_cluster/workspace_probe.json ::
%                                  layers.9.drug_decode = 0.8208333333333334
%   0.758  probe at layer 16       ... :: layers.16.drug_decode
%                                       = 0.7583333333333334
%   0.083  shuffled-label floor    ... :: chance = 0.08333333333333333
%                                  (identical at layers.<L>.chance)
%   0.883  context removed, L16    ... :: layers.16.drug_decode_context_removed
%                                       = 0.8833333333333332
%   480    prompt count            COMPUTED, not stored: config.n_drugs 12 x
%                                  config.n_per_drug 40. The thesis prints
%                                  \num{480} at 04b-representation.tex:115-120.
%   12, 40 twelve drugs, forty     ... :: n_drug_classes = 12, config.n_drugs
%          cells each              = 12, config.n_per_drug = 40
%
%   [SUPERSEDED IN PART BY PASS 12, below: the four ACCURACIES 0.821, 0.758,
%   0.083 and 0.883 are no longer set in this figure at all. 480, 12 and 40
%   remain, and the block above is kept verbatim because the ledger's
%   removed_from_drawing entries are audited against it.]
%
%   DRAWN AS A POSITION, panel A
%   2, 4, 6, 8, 9, 12, 16          ... :: config.layers = "2,4,6,8,9,12,16".
%   The seven ticks on the depth column are at TRUE layer spacing,
%   y = 2.10 + (L/16)*4.80, so the even grid and the one added layer (9) are
%   visible as such. This is the figure's only scaled mark and it is a scale of
%   depth, not of any measured quantity. The seven y values are, and are
%   CHECKED against the map rather than against the eye:
%     L    2     4     6     8     9     12    16
%     y   2.70  3.30  3.90  4.50  4.80  5.70  6.90
%   Layer 12 was drawn at 6.00 in the first build and is a repaired defect; see
%   design decision 2.
%
%   NOT FROM ANY JSON -- design facts, cited by line, drawn_true false in the
%   ledger. This is the handling slab.json gives the appendix's +2.75 nats.
%   the five prompt fields, in this order:
%     cell line | drug | dose | mechanism | control cell sentence
%     Sections/03-Apparatus.tex:250-252, "The prompt supplies cell line, drug,
%     dose, mechanism and the control cell's sentence".
%     CHECKED AGAINST figs/hook.tex:335-339, which draws the same strip for
%     sec:res-used: the five labels are spelled identically and are in the same
%     order, and hook.tex's header records the same measured \scriptsize widths
%     (cell line 0.954, drug 0.616 bold, dose 0.579, mechanism 1.424, control
%     cell sentence 2.511; total ink 6.084cm), arrived at independently. Only
%     the cell edges differ, because hook's strip is 8.20cm and this one is
%     7.20cm and each distributes its own slack. The two figures are meant to
%     read as one object seen twice, and they do.
%   the read position, "the last prompt position ... where generation begins":
%     Sections/04b-representation.tex:115-118.
%   the two failure modes, every word of F1 and F2:
%     Sections/04b-representation.tex:33-36, the \begin{question}[3] box.
%
% AGREEMENT WITH THE TEXT (checked against the Sections file, not assumed):
%   04b:126-127  "reads $82\%$ at layer 9 and $76\%$ at layer 16, the deepest
%                measured, against the $\approx 8\%$ floor"
%                -> panel A prints 0.821, 0.758 and the 0.083 floor.        OK
%   04b:131-132  "once cell line, plate and dose are projected out, the series
%                rises with depth instead, to $0.883$ ($88\%$)"
%                -> panel A prints 0.883 with that condition attached.      OK
%   04b:118-120  "A cross-validated multinomial logistic-regression classifier
%                learns which of the twelve drugs the prompt named"         OK
%   04b:122-124  "Folds are stratified over prompts and not grouped by well"
%                -> drawn verbatim as the probe's stated limitation.        OK
%   04b:108-110  "The prompt names the drug in plain text ..., so recovering
%                it ... establishes only that the model did not discard what
%                it was handed" -> B12 is that sentence, split over two lines,
%                and it is the ONLY entry in the licence that exists.       OK
%   04b:110-112  "The informative outcome would have been a negative, and what
%                the probe buys is the depth at which one would appear"
%                -> B13's second line, "the layer at which it is lost".     OK
%   04b:33-36    "Either the drug never enters its computation, an
%                activation-side encoding failure that target- and input-side
%                fixes might reach, or it enters and nothing reads it, a
%                decoder failure only an objective-side fix could reach"
%                -> F1, F2 and B13's last two lines are that partition,
%                word for word.                                             OK
%   04b:142-144  claim box: "It rules out that the drug field is lost before
%                the decoder runs, an encoding failure in the activations. It
%                does not rule out that the decoder ignores it."
%                -> the strike lands on F1 and only on F1.                  OK
% tab:workspace (04b:228-234) prints the same 0.821, 0.758 and 0.883 in its
% Decode columns; this figure sets them as text, not as marks, so no scale
% here can be read against that table.
%
% The subtitle of panel B, "a negative, at the 0.083 floor, would have been
% decisive", names the floor a second time on purpose. Without it the figure
% never states what a failure would have LOOKED like, and the only cue left is
% the word "failed" in the right-hand tag; the acceptance question this figure
% is built to answer is precisely "what would the probe have had to return".
%
% DESIGN DECISIONS, each with the reason that forced it
% -----------------------------------------------------
%  1. THE FIGURE CARRIES NO AXIS AND NO PLOTTED VALUE. The four accuracies are
%     set as text. An accuracy axis was the obvious alternative and was
%     rejected twice over: fig:comparators row A already plots exactly these
%     numbers on exactly that rule, so a second one would be a redraw; and an
%     axis invites the reader to read the height of 0.821 as the size of the
%     finding, when the section's whole point is that the finding's size is
%     nearly irrelevant because a positive was near-certain in advance.
%  2. THE ONLY SCALED MARK IS DEPTH. The seven ticks sit at true layer spacing
%     rather than evenly, because "seven layers" and "the depth at which a
%     negative would have appeared" is what the probe actually buys (04b:111-112),
%     and an evenly spaced tick row would hide that the grid is even except for
%     the one layer added at 9. The 8/9 pair is 0.30cm apart, the document's
%     minimum label separation, which is why the column is 4.80cm tall; it may
%     not be shortened without either dropping a numeral or lying about the grid.
%     REPAIRED, PASS 9: layer 12 was drawn at y = 6.00 and the map gives 5.70.
%     Measured off the 200dpi render, 2->4->6->8->9 ran at 0.30cm a layer, then
%     9->12 at 0.40 a layer and 12->16 at 0.225, so on the page three layers
%     occupied more room than four and layer 12 sat where layer 13 belongs. The
%     column is the figure's only scaled mark and the section's payoff is the
%     depth a negative would have named, so a grid step of error here is the one
%     error the figure cannot carry. Checked safe before drawing: 5.70 is 0.90
%     above the layer-9 tick and 1.20 below layer 16, and the nearest ink at
%     that height is the val node on baseline 5.97, which ends near x = 5.5,
%     well clear of the tick column at x = 7.02.
%  2b.THE 8/9 NUMERALS ARE NOT BRACED, AND THAT IS A DECISION, NOT AN OVERSIGHT.
%     A reviewer measured 0.102cm of clear white between the two numerals
%     against a glyph height of 0.203cm and asked for comparators.tex's brace
%     idiom, "9, 8" named as an ordered set. Overruled. comparators braces marks
%     whose POSITIONS a reader cannot resolve on a measured axis, and fanning
%     them would claim a resolution the estimates do not have; here the two
%     positions are exactly right, each numeral is unambiguously attached to its
%     own 0.42cm tick, and the tightness IS the datum -- layers 8 and 9 are
%     adjacent in a sixteen-layer stack and must look it. Bracing them would
%     hide the one thing the added layer at 9 is there to show. The pair was
%     re-read at 200dpi after the numerals went quietink and the two glyphs are
%     separate and legible. The column cannot be lengthened either: its bottom
%     is pinned to the caret at y = 2.10 (which is layer 0, the prompt) and its
%     top to the "layer" label at 7.12, so the most 16 layers could ever have is
%     4.95cm, 0.009cm a layer more than they have now.
%  3. NO PANEL LETTERS. Two uppercase sans word tags in accent, the corpus
%     convention (circularity.tex, slab.tex); (a)/(b) appears nowhere in this
%     document's figures.
%  4. THE STRIKE IS THE ONLY MARK ABOVE 1.0pt. 1.1pt, the same weight as
%     circularity.tex:156, where one struck edge carries that figure's claim.
%     REPAIRED, PASS 9: it is now ONE HORIZONTAL stroke through the x-height of
%     line one, 8.70 to 13.50 at y = 5.60, and not the diagonal
%     (8.86,5.20)--(13.94,5.64) of the first build. Both reviewers reported the
%     diagonal as blocking and the render agrees: its left end sat exactly on
%     line two's baseline and its right end 0.11cm above line one's, so it
%     entered along the bottom of "an activation" and left through the x-height
%     of "computation", and at print size "an activation" was a smear. Its right
%     end was also 0.26cm from the box border with the accent arrow to "struck"
%     starting at 14.20 on nearly the same line, so the two read as one
%     continuous pointer out of the box. A rule through the x-height leaves the
%     words readable, because the eye completes struck text and cannot complete
%     text cut along its baseline. LINE ONE ONLY: line one is the proposition
%     the positive eliminates and line two is its gloss, so one stroke does the
%     work and the 1.1pt register stays "used exactly once". Extent measured off
%     the render: line one's ink is 8.76 to 13.44, the stroke overhangs 0.06cm
%     each side, and it stops 0.70cm short of the box's right border, which is
%     what kills the leader reading.
%     Everything else here is 0.30-0.90pt, so weight alone says
%     which mark is the argument. The struck box is drawn and STRUCK, not
%     deleted: a deleted box leaves the caption nothing to point at, and the
%     partition is two-membered whether or not one member survives. The two
%     boxes carry IDENTICAL borders for the same reason -- the struck one was
%     drawn at rulegrey!55 in the first build, and if the box is also fainter
%     then two marks say the same thing and neither says it alone.
%  5. THE TWO LICENCE BLOCKS ARE SIZED BY THEIR CONTENTS AND BY NOTHING ELSE.
%     The one the positive buys holds one sentence; the one a failure would
%     have bought holds four lines. That difference in footprint IS the
%     asymmetry, readable before a word is read. It quantifies nothing, and
%     both the caption and the ledger say so.
%  6. SOLID-AND-WHITE MEANS "HAPPENED", DASHED-AND-CREAM MEANS "DID NOT".
%     One register, used once each, so the counterfactual block cannot be
%     mistaken for a second result. REPAIRED, PASS 9: the cream is now panelbg
%     straight, the same cream slab.tex and hook.tex use, and so is the depth
%     column's. Three creams were in play across this slate -- panelbg in hook
%     and slab, panelbg!60 in the column, panelbg!45 in the counterfactual box
%     -- and the last two are about 2% ink, which will not reliably reproduce on
%     uncoated stock; a tint that does not print is a tint doing no work while
%     adding a third cream to the system. The DASHED border is what carries
%     "this did not happen"; the fill only says "a filled region".
%  7. THE LICENCE THAT EXISTS HAS NO SECOND ENTRY AND NO "which fixes it points
%     at" LINE. The answer to that question is nothing, and an absent line says
%     it better than the word "none", which would read as a measured category.
%     Same principle as comparators.tex's bare dot for a quantity with no
%     interval: the absence is the mark.
%  8. THE DRUG CELL IS THE ONLY ACCENT MARK IN THE PROMPT STRIP, at accent!12
%     fill with an accent border and an accent bold label, so the eye finds it
%     before it reads anything. It is the reason every number in the left-hand
%     block is what it is.
%  9. NO INTERVAL IS DRAWN ANYWHERE. The probe returns a cross-validated
%     accuracy and workspace_probe.json stores no interval for it; the document
%     draws bars only where the instrument returns its convention
%     (comparators.tex, style_v2.py). Nothing here is a bar, a whisker or a
%     dot on a scale, so nothing invites one.
% 10. THE ACCENT BUDGET IS THE DRUG CELL, THE CARET, THE SEVEN DEPTH TICKS,
%     THE STRIKE AND ITS ARROW AND LABEL, AND THE FOUR WORD TAGS.
%     REPAIRED, PASS 9: the seven depth NUMERALS were accent and are now
%     quietink. With them the figure carried about 24 accent marks against a
%     corpus ceiling near 16 (comparators.tex gives accent exactly two jobs),
%     and roughly fourteen of them stood inside one 0.4 x 4.8cm column, making
%     it the densest accent region on the page -- so the eye landed on the tick
%     column first and then hunted beside tick 9 for "0.821", which is set as
%     prose three centimetres to its left. A numeral beside a tick is an axis
%     label, which is what quietink is for in every other figure in the corpus
%     (comparators.tex:70-72, gate.tex, purify.tex). The TICKS stay accent
%     because the depth grid is what \cref{sec:methods-probe} says the probe
%     actually buys. Re-rasterised at 115dpi after the change -- about print
%     size at arm's length -- and the first fixation is now the drug cell.
% 11. NO 0.50 CHANCE RULE APPEARS. This instrument's floor is 1/12 = 0.083, not
%     0.50, and it is set as text rather than as a rule because there is no
%     axis for a rule to cross. (Recorded because the corpus is split on how a
%     chance rule is drawn: style_v2.chance_line() mandates thin solid black,
%     while ladder.tex:112-115, the only TikZ precedent, draws it densely
%     dotted. This figure needs neither.)
%
% LAYOUT NOTES, all forced by a build or by a measurement
% -------------------------------------------------------
%   * fullwidth = textwidth 142mm + marginparwidth 28mm + marginparsep 4mm
%     = 174mm = 493.23pt. (preamble.tex's own comment beside
%     \newenvironment{fullwidth} says 171mm and is stale arithmetic.) The box
%     is pinned with a bare \path to 17.30cm = 491.4pt, 1.8pt inside the
%     measure. Nothing is drawn outside x in [0, 17.30] except the 0.2pt
%     half-width of B13's right border.
%   * EVERY MULTI-LINE BLOCK IS SEPARATE anchor=base west NODES. No node in
%     this file uses align= or text width. frames.tex measured a 3.7pt drift
%     from that: TikZ's align builds a minipage whose first-line height follows
%     the AMBIENT \baselineskip, so a node's ink moves with whatever prose
%     precedes the float. Every y below IS a baseline. Within-block leading is
%     0.33cm = 9.39pt, one \scriptsize line; between blocks it is 0.45cm.
%   * THE PROMPT STRIP'S CELL EDGES ARE SET FROM MEASURED STRING WIDTHS, not
%     chosen. Every string in this figure was set in the real preamble and its
%     \wd read back before a coordinate was written. At \scriptsize the five
%     field names measure, in cm:
%       cell line 0.954 | drug (bold) 0.616 | dose 0.579 | mechanism 1.424 |
%       control cell sentence 2.511.  Total ink 6.084cm.
%     The strip therefore cannot be the 6.78cm first laid out: at that width
%     "mechanism" alone overruns its 1.28cm cell by 0.144cm and the five cells
%     share 0.696cm of padding, 0.07cm a side. The strip is 7.20cm, the widest
%     that keeps the depth column's numerals inside panel A, which leaves
%     0.09-0.11cm each side -- more than the 1.6pt inner sep slab.tex sets on
%     its own nodes. Cell edges: 0.00 / 1.13 / 1.93 / 2.69 / 4.29 / 7.20.
%   * THE ELLIPSIS INSIDE THE FIFTH CELL WAS DROPPED, AND THE MEASUREMENT IS
%     WHY. A typeset $\cdots$ at \scriptsize is 0.405cm and needs about 0.16cm
%     of gap; the fifth cell is 2.91cm and its label is 2.511cm, leaving
%     0.399cm in total. No strip width recovers it: pushing the caret past
%     x = 7.30 drives the layer numerals to within 0.13cm of the panel divider.
%     A drawn three-dot ellipsis (0.18cm) leaves a 0.066cm gap to the label,
%     which is 1.9pt and reads as part of the last word.
%   * THE FIFTH CELL'S RIGHT EDGE IS SOLID. It was drawn dotted first, on the
%     reading that the strip has no horizontal scale. Two faults, both seen in
%     the render: the caret covers the top 0.21cm of that edge, leaving a
%     0.41cm dotted stub that reads as a printing fault; and a dotted edge at
%     exactly the x the caret marks says both "the prompt ends here" and "this
%     field runs on past here". The caret is the load-bearing mark, so the
%     absence of scale is stated in the caption, in words, instead.
%   * THE CARET STRADDLES THE STRIP'S RIGHT EDGE ON PURPOSE, and is centred on
%     the depth column above it. The last prompt position is where the prompt
%     ends, so the mark for it sits ON that boundary rather than inside the
%     last cell. PASS 9: its apex is now the column's own bottom edge, y = 2.10,
%     and the 0.5pt accent stub that used to bridge 2.04 to 2.10 is deleted.
%     Inside a 3mm square at that corner there converged the strip's top rule,
%     the strip's right rule, the caret, the stub, the column's bottom edge and
%     the brace's right arm, and the stub crossed a rulegrey rule uninterrupted.
%     Six strokes became five and the apex does the joining. The brace's right
%     arm was left where it is: measured, it terminates at (7.20,1.30), which is
%     0.12cm below the strip's bottom-right corner and 0.56cm below the caret's
%     own base, so it is not part of that corner.
%   * THE PROMPT CELLS ARE DRAWN AS figs/hook.tex DRAWS THEM. PASS 9: rulegrey
%     0.35pt, panelbg fill, rounded corners=1pt for the four ordinary fields,
%     and accent 0.4pt / accent!12 / rounded corners=1pt for the drug field.
%     Before this pass licensing drew square-cornered WHITE cells sharing edges
%     and hook drew rounded panelbg boxes with rounded gutters, so the same
%     five-field prompt was drawn two incompatible ways in two figures a page
%     apart while both captions insisted they were the same prompt. hook's is
%     the better mark -- the rounding reads as token boundaries rather than as
%     a table -- so licensing adopts it. Only the cell EDGES still differ, and
%     they must: this strip is 7.20cm and hook's is 8.20cm, and each distributes
%     its own slack. That is a difference of measure, not of treatment.
%   * NO LEADER RUNS FROM THE CARET LABEL TO THE CARET. One was drawn and
%     taken out: the only clear path from the label's lower right corner passes
%     within 0.02cm of the depth column's bottom-left corner at (7.02,2.10), or
%     else crosses the strip's top border at x = 6.94 and runs on inside the
%     last cell. Both read as a stray stroke. The label's second line ends
%     0.21cm from the caret's footprint with nothing between, inside the
%     document's ~0.3cm leader threshold.
%   * NO ACCENT MARK LIES LEFT OF x = 7.02. The seven depth ticks first ran
%     from 6.92, within 0.49cm of the text block, and the block's line at
%     y = 3.30 shares its baseline exactly with the layer-4 tick, so
%     "0.883 once cell line, plate and dose are projected out ---- 4" read as
%     one labelled row. The ticks now start at the column's own left edge.
%   * PANEL A'S LINE LENGTH IS CAPPED AT 6.85cm, not the 6.35cm of the first
%     layout: that limit followed the depth column, which moved right from
%     6.60 to 7.02 when the strip was widened. The longest line drawn is
%     "where generation begins; read at seven depths" right-aligned on 6.84,
%     and the longest left-aligned line is "a cross-validated
%     multinomial probe, twelve classes" at 6.57cm.
%   * "READ AT SEVEN DEPTHS", NOT "READ AT EVERY LAYER". Beside a column that
%     is sixteen layers tall and carries seven ticks, "every layer" left a
%     reader unable to tell whether the ticks were the seven layers actually
%     probed or an arbitrary sample of a sweep over all sixteen -- which is the
%     one fact the column exists to carry. The review asked for "read at each of
%     the seven marked depths"; measured, that is 1.94cm longer and the node is
%     anchor=east on 6.84, so it would have started at x = -1.1 and taken the
%     picture off the measure. "seven depths" is 20 characters against "every
%     layer"'s 19 and cost nothing.
%   * THE COUNTERFACTUAL BLOCK'S FIRST ENTRY LOST ITS FINITE VERB, and the
%     review's own wording would not fit. "the drug is lost inside the model,"
%     read as an assertion that the drug IS lost, the opposite of the section's
%     conclusion; the review asked for "that the drug is lost inside the
%     model,", which measures 4.66cm from x = 12.84 and put the picture 11.88pt
%     over the measure -- the harness reported exactly that Overfull hbox on the
%     build that tried it. "the drug lost inside the model," is 4.02cm, fits
%     with 0.44cm to spare, and is the noun phrase the other three entries are.
%   * THE CARET LABEL IS SHORT LINE OVER LONG LINE, not the reverse. Written
%     long over short it sat 0.75cm under a left-aligned quietink line of the
%     same width and register, and at print size the two read as one three-line
%     grey paragraph. Short first, right-aligned, breaks that; the long second
%     line then also fills the 3.8cm of empty measure that sat above the strip.
%   * B14 IS FIGURE-WIDE, NOT PANEL-LOCAL. It measures 11.70cm at \footnotesize
%     and is centred on 8.65, the figure's own centre, below the panel divider
%     (which stops at y = 0.90). Panel-local it would have run from 7.10 to
%     18.80 and off the measure.
%   * TYPE. Ambient \small on the picture; every node sets its own size.
%     \scriptsize = 8pt on all of it except the two panel subtitles and the
%     closing line, which are \footnotesize = 10pt, and the four word tags,
%     which are \sffamily\scriptsize\bfseries. Nothing is set at body size and
%     nothing is scaled: 1cm here is 1cm on the page, no scale= key, no
%     \resizebox.
%
% ===========================================================================
% PASS 12 -- 2026-08-18. THE STYLE-CONTRACT REBUILD.
% ===========================================================================
% Everything above this line is the v1-v11 record and is left intact because
% the thesis is audited against it. This section states what pass 12 changed,
% and SUPERSEDES BY NAME every claim above that pass 12 makes false. Nothing
% above has been edited away.
%
% WHY. The supervisor's verdict on the 4.8-4.12 slate: "go back over them and
% iron them out even more ... make it look more professional, right now it
% looks a bit sloppy and childish ... make sure that the figures are CORRECT
% and display the idea we want to convey in the best way possible." The
% reference standard is fig:mechanism (4.11, printed page 51): lowercase
% parenthesised panel letters with sentence-case descriptive titles, rotated
% axis labels that NAME the measured quantity, direct labels, no legend, no
% all-caps banners, no italic aphorisms, no boxed callouts, and the whole
% argument carried by a twenty-line caption. Measured against it this figure
% was, before pass 12: 4 all-caps sans accent banners out of 4 panel heads,
% 17% of its characters italic, 11.21% ink at 200dpi and 9.9 characters per
% 100mm^2 -- DENSER THAN THE APPENDIX REFERENCE PLATE A.3 (9.7%, 6.4) while
% sitting in the running argument of chapter 4.
%
% WHAT PASS 12 DELETED FROM THE DRAWING, AND WHY
% ----------------------------------------------
%  D1. THE FOUR ACCURACIES, set as running text in panel A: "$0.821$ at layer
%      9 and $0.758$ at the deepest,", "against a $0.083$ shuffled-label
%      floor", "$0.883$ once cell line, plate and dose are projected out", and
%      the "what it returns" key that headed them. They sat immediately left of
%      a scaled layer stack carrying ticks labelled 9 and 16, so a reader tries
%      to read 0.821 off the stack at layer 9 and cannot -- and the caption had
%      to spend two whole sentences denying the reading the drawing offered
%      ("They are set here as text, not as marks on a rule; the accuracy rule
%      is Figure 4.7, row one"). A caption that must deny a reading its own
%      drawing invites is a defect in the drawing, not in the caption.
%      All four are drawn or printed in three other places -- fig:comparators
%      row (a), fig:mechanism panel (a), tab:workspace -- and this figure's own
%      caption prints them as well. Under the cross-figure redundancy
%      resolution, fig:mechanism (a) OWNS 0.821 and 0.883 because it is the
%      only place the two series are drawn AS SERIES and the crossing at layer
%      12 is visible; 4.7 row (a) keeps the chance floor because a comparator
%      row without its comparator is not a comparator row. 4.8 keeps the
%      ARGUMENT about the floor -- what a negative would have looked like --
%      and loses the numerals. All four survive in licensing.json as
%      removed_from_drawing entries with their key paths and with pointers to
%      the figures that still draw them.
%      SUPERSEDES: the "set as text, panel A" block above (kept verbatim, since
%      the ledger's removed_from_drawing entries are audited against it); the
%      AGREEMENT-WITH-THE-TEXT lines for 04b:126-127 and 04b:131-132 (the
%      agreement is now the CAPTION's to carry, not the drawing's); design
%      decision 1's second clause; and the paragraph beginning "The subtitle of
%      panel B ... names the floor a second time on purpose", which is void
%      because that subtitle is deleted (D3).
%  D2. THE FOUR ALL-CAPS SANS ACCENT BANNERS -- "WHERE THE DRUG NAME ENTERS",
%      "WHAT THE OUTCOME LICENSES", "WHAT THE POSITIVE BUYS", "IF IT HAD FAILED
%      INSTEAD". Every one of them narrated rather than named. A.3 (figs/slab.tex)
%      does use short accent caps banners, but it is an appendix reference
%      plate a reader stops at and studies, its banners are noun phrases naming
%      objects, and the caps sans accent face is byte-identical to the
%      preamble's \marginlabel apparatus face -- which belongs to the margin,
%      not to a figure in the running body. 4.11 and 4.13 carry zero banners.
%      Replaced by 4.11's form: a lowercase parenthesised letter and a
%      sentence-case descriptive title, \footnotesize roman ink, left-aligned
%      to the panel's own x origin, no colon, no bold, no caps.
%      SUPERSEDES design decision 3 ("NO PANEL LETTERS ... (a)/(b) appears
%      nowhere in this document's figures") IN FULL. It was true when written
%      and is now false: fig:mechanism sets "(a) encoded more, used no more"
%      and "(b) ablation cost against a random null", and it is the reference
%      standard the slate is being brought to.
%  D3. THE THREE ITALIC SUBTITLES -- "the drug is named in the prompt itself",
%      "a positive was near-certain in advance;", "a negative, at the $0.083$
%      floor, would have been decisive" -- and D4's epigram. The subtitles were
%      \footnotesize italic quietink, i.e. the LARGEST type in the figure was a
%      grey italic editorial line while every measured thing was set 25%
%      smaller. The last of the three is also the caption's own lede restated.
%  D4. THE CENTRED ITALIC EPIGRAM "one run, two possible outcomes; the one that
%      could not fail is the one that happened." It was 9.96pt Pagella-Italic
%      quietink centred within 0.2pt of the figure's own x-centre, and the
%      naive-reader pass reports the eye going to it BEFORE any mark, reading
%      it as a second caption above the real one. In this document the argument
%      lives in the caption. It is written to figs/_PENDING_CAPTIONS.tex.
%  D5. THE FOUR BOX STROKES -- the two failure-mode boxes, the solid white
%      "what the positive buys" block and the dashed panelbg counterfactual
%      block. A stroked rectangle is drawn only where the rectangle IS an
%      object in the argument (a token, a prompt field, a subspace, a scoring
%      frame). These four enclosed SENTENCES, and the figure's own caption
%      conceded that "the figure quantifies nothing about the size of a
%      licence". Panels (c) and (d) are now two columns of equal width on one
%      x-grid with baseline-aligned entries, in the frames.tex idiom.
%      SUPERSEDES design decision 6 IN FULL (there is no dashed border and no
%      cream fill left to carry "this did not happen"; quietink carries it, as
%      it does everywhere else in the slate), and design decision 5's second
%      sentence -- see C4 for what replaced the footprint asymmetry.
%  D6. THE GREY FILL INSIDE THE LAYER STACK, and the stack itself as a filled
%      column. It read as a BAR: the naive-reader pass's first question was
%      what quantity the length of the grey bar encoded, and the answer is
%      none. It also ran from y = 2.10 to 7.05, i.e. past the layer-2 tick
%      below and the layer-16 tick above, implying that depth continues beyond
%      the measured range. It is now a plain 0.40pt ink axis rule running from
%      the layer-2 tick to the layer-16 tick and stopping exactly on them.
%  D7. "folds stratified over prompts, not grouped by well". Prose already
%      carried at 04b:122-124; it is a limitation of the instrument, which is
%      caption and prose material, not a name for anything the figure draws.
%      SUPERSEDES the AGREEMENT line for 04b:122-124 above.
%
% WHAT PASS 12 CHANGED, AND WHY
% -----------------------------
%  C1. THE AXIS NOW DEFINES ITSELF. This is the supervisor's first complaint
%      and it is answered mechanically. The depth stack is the figure's only
%      axis and it carried, before this pass, seven bare numerals and the word
%      "layer" set as a 0.30cm-tall stub above the column top -- no quantity in
%      axis-label position, no scope. It now carries 4.11's form exactly:
%      rotated 90 degrees reading upward, immediately left of the tick
%      numerals, line 1 \footnotesize roman ink "layer", line 2 \scriptsize
%      quietink "residual stream read at the last prompt position, where
%      generation begins". Line 2 is 258.63pt = 9.09cm set on one line, which
%      does not fit beside a 4.20cm axis, so it is wrapped onto two typeset
%      lines at the only break that keeps both under the 4.60cm of clear
%      vertical measure available at that x: "residual stream read at the last
%      prompt" (137.06pt = 4.817cm) over "position, where generation begins"
%      (119.57pt = 4.202cm). It is ONE label line wrapped, not two label lines.
%      Rotated multi-line text stacks toward increasing x, so line 1 ("layer")
%      sits FARTHEST from the axis, at x = 5.44, exactly as a wrapped
%      matplotlib ylabel does in 4.11.
%  C2. THE PANEL-A PROSE BLOCK IS NOW ONE BLOCK OF TWO LINES, and it is the
%      figure's only prose. "what is read: the residual stream at the last
%      prompt position" measures 264.20pt = 9.29cm at \footnotesize and cannot
%      be set on one line inside a 6.90cm panel; more to the point, C1 puts
%      that exact fact in axis-label position, where the style contract wants
%      it, so keeping it in the prose block as well would be a duplication.
%      The block is therefore "what reads it: a cross-validated" over
%      "multinomial probe, twelve classes" (4.851 and 5.241cm), \footnotesize
%      roman ink, left edge on panel (a)'s x origin.
%  C3. THE COLOUR INVERSION IS FIXED, and it was a correctness defect, not a
%      preference. Before this pass the accent arrow marked the STRUCK branch,
%      so accent meant ELIMINATED here and meant THE HEADLINE RESULT in the
%      figure four pages earlier; the naive-reader pass read the accent arrow
%      as the surviving branch. Across the slate accent now means one thing:
%      the drug-side object the rest of the chapter measures. In this figure
%      that is exactly two objects -- the drug cell in the prompt strip, and
%      the SURVIVING failure mode with its arrow and its outcome label. The
%      struck mode, its arrow, its outcome label and the whole of panel (d),
%      which is a counterfactual, are quietink. Panel titles, the axis, the
%      axis labels, the prose, the caret and panel (c) are ink; ticks, brace,
%      dividers and the four ordinary prompt cells are rulegrey furniture.
%      SUPERSEDES design decision 10's accent budget IN FULL: the seven depth
%      ticks are no longer accent (they are furniture, 0.30pt rulegrey), the
%      caret is no longer accent (it is part of the axis apparatus, ink), the
%      strike is no longer accent (it is the black threshold-weight stroke, see
%      C6), and the four word tags no longer exist (D2). It also supersedes the
%      layout note "NO ACCENT MARK LIES LEFT OF x = 7.02", which is now
%      trivially true for a different reason: the only accent in panel A is the
%      drug cell, at x 1.12 to 1.90.
%  C4. PANELS (c) AND (d) ARE NOW EQUAL. Two columns of equal width (4.60cm)
%      on one x-grid, entries baseline-aligned, FOUR ENTRIES EACH. Before this
%      pass the counterfactual block was the larger and more detailed of the
%      pair -- one sentence against four lines, inside a bigger dashed box --
%      so the eye landed on the thing that did not happen. The asymmetry the
%      figure exists to draw is now carried by the WORDS rather than by the
%      footprint: panel (c)'s entries 3 and 4 name, in the same slots and the
%      same register as panel (d)'s entries 2 to 4, exactly what the positive
%      did NOT buy -- "no layer, no class of failure," / "and no class of fix."
%      against "the layer at which it is lost," / "an activation-side failure,
%      and" / "target- and input-side fixes." A reader can now diff the two
%      columns line by line, which is a stronger reading of the asymmetry than
%      a difference of box size ever was.
%      SUPERSEDES design decision 7 ("an absent line says it better than the
%      word 'none'") and design decision 5's second sentence. The reason
%      decision 7 gave still stands against the bare word "none", which would
%      read as a measured category; what is drawn is not "none" but a named
%      absence in a slot the facing column fills, which is the opposite of a
%      category and cannot be read as a quantity.
%  C5. ONE X ORIGIN PER PANEL, declared as a named coordinate and referenced by
%      every element of that panel. Before this pass each banner sat 4.6pt
%      (1.6mm) LEFT of the text it headed, four times over, and the central
%      divider had 8.0pt of clearance on its left and 15.2pt on its right.
%      Panel (a) origin x = 0.00; panels (b), (c) origin x = 7.75; panel (d)
%      origin x = 12.70. The divider sits at x = 7.42, which is 0.33cm from
%      panel (a)'s rightmost ink (the caret's right base at 7.09) and 0.33cm
%      from panel (b)'s origin: the gutter is now equal on both sides to
%      0.001cm.
%  C6. FIVE STROKE WEIGHTS BECOME THREE, EACH WITH ONE MEANING.
%        0.30pt rulegrey -- furniture: the seven ticks, the brace, the panel
%                           divider, the (c)/(d) row rule, the four ordinary
%                           prompt cells, and the two outcome leaders (which
%                           carry their mark's colour, not rulegrey).
%        0.40pt          -- the axis rule (ink) and the drug cell's border
%                           (accent). Used twice, so it is not a singleton.
%        0.90pt ink      -- the strike, and nothing else.
%      The 1.1pt register of design decision 4 is retired: 0.90pt is the
%      document's threshold weight (style_v2.chance_line() is hard-coded
%      lw=0.9, and figs/family.pdf draws the pre-registered materiality margin
%      at 0.90pt), and a figure that contains no threshold at all can carry its
%      one mark of argument at that weight without any collision. THAT IS THE
%      JUSTIFICATION the contract requires for a weight used in one kind of
%      place: this figure has no chance rule, no gate, no margin and no floor
%      (design decision 11, which still holds), so nothing else here may be
%      black solid 0.90pt and nothing is.
%      SUPERSEDES design decision 4's first sentence and its "used exactly
%      once" clause: the strike is now TWO strokes, one through each of the two
%      lines of the struck mode, because the mode is a two-line entry and a
%      rule through its first line only left its second line standing. It also
%      supersedes the cell-border weight recorded in the pass-9 prompt-cell
%      note: 0.35pt is not one of the document's five weights and the ordinary
%      cells are now 0.30pt. The cells' construction is otherwise unchanged and
%      still matches hook.tex's \tokbox.
%  C7. THE FIGURE FOOT CARRIES ONE \scriptsize QUIETINK DISCLOSURE LINE,
%      left-aligned to the figure's x origin: "Nothing in this figure is a
%      measured quantity; the depth ticks are the only scaled marks." Every
%      figure in the slate carries exactly one such line. Here it does the job
%      an interval declaration does elsewhere -- it tells the reader, without
%      leaving the figure, that no mark on the page is an estimate. Design
%      decision 9 (no interval is drawn anywhere) still holds and is now
%      readable on the drawing rather than only in the ledger.
%  C8. THE STRIP IS 6.90cm, NOT 7.20cm, and the caret is 0.40cm tall rather
%      than 0.24cm. Both follow from C1: the rotated axis-label block needs a
%      clear vertical measure of 4.82cm at x = 5.76, which puts its lower edge
%      at y = 2.391, so the strip's top edge came down from 2.46 to 2.28 and
%      the caret grew to span 2.28 to 2.70 -- 0.111cm of clearance, measured,
%      between the rotated label's foot and the strip's top rule. The strip's
%      slack is redistributed over the same five measured field widths:
%      6.084cm of ink and 0.816cm of padding, 0.0816cm a side, giving cell
%      edges 0.00 / 1.12 / 1.90 / 2.64 / 4.23 / 6.90 against the old
%      0.00 / 1.13 / 1.93 / 2.69 / 4.29 / 7.20. NO FIELD CHANGED, no order
%      changed, and the strip still carries no horizontal scale.
%      SUPERSEDES the layout note "THE PROMPT STRIP'S CELL EDGES", the
%      "PANEL A'S LINE LENGTH IS CAPPED AT 6.85cm" note (the cap is now 6.90cm,
%      set by the strip, and the longest left-aligned line drawn is 5.241cm),
%      the caret note's "apex is now the column's own bottom edge, y = 2.10"
%      (the apex is the axis rule's own base at the layer-2 tick, y = 2.70),
%      and the "B14 IS FIGURE-WIDE" note (B14 is deleted, D4).
%      The ELLIPSIS note and the FIFTH CELL'S RIGHT EDGE note both still hold
%      and were re-checked at the new width: the fifth cell is 2.67cm against a
%      2.511cm label, so there is even less room for a $\cdots$ than there was.
%  C9. WHAT DID NOT CHANGE, and must not, because it is the argument and the
%      measurement: the seven depth ticks and their map y = 2.10 + (L/16)*4.80,
%      giving 2.70 / 3.30 / 3.90 / 4.50 / 4.80 / 5.70 / 6.90 -- BYTE-IDENTICAL
%      to pass 11, including the pass-9 repair of layer 12 from 6.00 to 5.70;
%      the five prompt fields, their order, their measured widths and the
%      accent on the drug cell; the caret at the strip's right edge and the
%      read-here relation between the two; the brace and the count line
%      "$480$ real tier-2 prompts: twelve drugs, forty cells each"; the strike
%      on the eliminated mode and the two-member partition; and the fact that
%      the figure draws no interval, no bar and no marker on a rule.
%      Design decisions 2, 2b, 8 and 11 stand unchanged. NO NUMBER IN THIS
%      FIGURE CHANGED VALUE IN PASS 12; four numbers left the drawing entirely
%      and are recorded in licensing.json as removed_from_drawing.
%
% GEOMETRY OF PASS 12, all measured with \wd in the real preamble before a
% coordinate was written (figs/_preview.py licensing, 200dpi, read at each step)
% -------------------------------------------------------------------------
%   panel (a) origin x = 0.00 ; panels (b) and (c) origin x = 7.75 ;
%   panel (d) origin x = 12.70 ; divider x = 7.42 ; measure 0 to 17.30cm.
%   axis rule x = 6.90, y 2.70 to 6.90 ; ticks 6.62 to 6.90 ; numerals
%   anchor=east on 6.52 ; rotated "layer" at x = 5.44, wrapped scope note at
%   x = 5.76 and 6.04, all centred on the axis midspan y = 4.80.
%   Measured string widths used (cm, \footnotesize unless marked s = \scriptsize):
%     (a) title 4.969 | (b) title 9.495 | (c) title 4.081 |
%     (d) title, wrapped 4.157 over 1.921
%     prose 4.851 / 5.241
%     mode 1 5.842, gloss s 4.292 | mode 2 4.472, gloss s 2.093
%     outcomes s 2.727 / 2.147 and 2.921 / 1.712
%     (c) entries s 3.210 / 2.516 / 3.352 / 2.293
%     (d) entries s 3.804 / 3.303 / 3.689 / 3.399
%     count line s 6.557 | foot line s 10.860
%     cell labels s: 0.954 / 0.616 bold / 0.579 / 1.424 / 2.511
%   Widest object on the page is panel (b)'s title, 9.495cm, which ends at
%   x = 17.245, 0.055cm inside the measure. No text crosses the measure.
%
% C10. THE EMPTY BAND IN PANEL (a) IS STRUCTURAL AND IS LEFT EMPTY ON PURPOSE.
%   Measured on the 200dpi render it is about 5.2 x 4.0cm, at x 0.0-5.2 and
%   y 2.4-6.4. It is forced, not careless, and four ways of removing it were
%   built and rejected:
%     (i)  shortening the axis. The 8/9 pair is 0.30cm apart at the current
%          scale, which is the document's minimum label separation, so the
%          16-layer span cannot be under 4.80cm and the drawn 2-to-16 span
%          cannot be under 4.20cm without dropping a numeral or falsifying the
%          grid. Design decision 2 and 2b, both still standing.
%     (ii) narrowing the strip. Its five field names measure 6.084cm of ink and
%          hook.tex sets the same five, so the strip cannot go below about
%          6.7cm without either abbreviating a field or breaking the agreement
%          with fig:hook that both captions assert.
%     (iii)inverting the axis so the strip sits at the top and depth runs
%          downward. This moves the band, it does not remove it, and it would
%          draw the residual stream running the wrong way.
%     (iv) moving panels (c) and (d) into the band. They need 2 x 4.60cm plus a
%          gutter, which is 9.55cm; the band is 5.2cm wide.
%   What is left is the space between the prompt and the read, which is the
%   residual stream, and this figure deliberately draws nothing there: the
%   stream's contents are fig:mechanism's subject and the slab's geometry is
%   fig:slab's. Filling it with narration is exactly what pass 12 removed.
% ===========================================================================
% PASS 13 -- 2026-08-18. THE REFUTATION PASS.
% ===========================================================================
% Three adversarial refuters read the pass-12 render. Everything above is left
% intact; this block states what pass 13 changed and SUPERSEDES BY NAME every
% pass-12 claim it makes false. Each change is recorded at the code it changes
% as well, so a reader of the drawing does not have to come back here.
%
% R1. THE COORDINATE RECORD IN "GEOMETRY OF PASS 12" AND IN C5 WAS WRONG IN
%     SIX PLACES, AND THE HEADERS ARE WHAT THE THESIS IS AUDITED AGAINST.
%     Superseded, by name, with the values the code actually drew in pass 12:
%       divider x         header 7.42   -> DRAWN 7.40   (and now deleted)
%       caret right base  header 7.09   -> DRAWN 7.05   (and now 7.02)
%       divider gutters   header 0.33   -> DRAWN 0.35 both sides
%       tick inner end    header 6.62   -> DRAWN 6.74   (and now 6.78)
%       numerals anchor   header 6.52   -> DRAWN 6.64   (and now 6.70)
%       rotated "layer"   header 5.44   -> DRAWN 5.36   (and now 6.18)
%       scope line two    header 6.04   -> DRAWN 6.12   (and now deleted)
%     The same block recorded the two failure-mode gloss lines at their
%     \scriptsize widths ("gloss s 4.292", "gloss s 2.093") while pass 12's
%     code set both at \footnotesize (5.365 and 2.616cm) -- and it is the gloss
%     width that sets the strike's right end, so the header could not be used
%     to check the one mark the file calls its argument. Pass 13 sets the modes
%     at \scriptsize, so 4.292 and 2.093 are now the drawn widths and the
%     record and the code agree for the first time.
%     C4's "Two columns of equal width (4.60cm)" was also never true of the
%     drawing: pass 12's origins were 7.75 and 12.70 on a 17.30 measure, i.e.
%     4.95 and 4.60. See the (c)/(d) block for what pass 13 drew instead.
%
% R2. figs/_PENDING_CAPTIONS.tex DID NOT EXIST. D4 and the note above the
%     staged float both asserted that the cut epigram and the pass-12
%     replacement caption were "written to figs/_PENDING_CAPTIONS.tex"; the
%     file was never created, so every sentence cut under D1, D3, D4 and D7 was
%     unaccounted for outside a comment block in this file, and the caption
%     that actually prints in the thesis is still the pass-11 one, which is
%     false against the drawing in four places (it indexes the panels "Left,"
%     and "Right,", it prints four accuracies as "set here as text", and it
%     says the two licence blocks "are sized by their contents"). Pass 13
%     creates figs/_PENDING_CAPTIONS.tex and writes the pass-13 caption and
%     every deleted sentence into it, keyed by \label{fig:licensing}. This
%     builder still does not edit Sections/. The copy below is retained as the
%     audit trail and is marked a duplicate of the pending file.
%
% R3. THE STRIKES WERE BLACK SOLID 0.90pt, WHICH IS THE SLATE'S THRESHOLD
%     STROKE. Corrected to quietink 0.40pt. Full reasoning at the code.
%     SUPERSEDES C6's third weight; the figure now has two weights, 0.30 and
%     0.40pt, neither a singleton.
% R4. THE CARET'S APEX WAS THE AXIS BASE AND THE LAYER-2 TICK AT ONCE.
%     Broken apart; apex 2.50, base 6.78-7.02. SUPERSEDES C8's apex clause.
% R5. THE AXIS LABEL WAS THREE ROTATED LINES AND OVERRAN THE AXIS AT BOTH ENDS,
%     with the quantity name outermost. One \footnotesize "layer" beside the
%     numerals, one \scriptsize scope line outboard, 3.908cm inside a 4.20cm
%     axis. SUPERSEDES C1's wrap arithmetic and its "line 1 is the outermost".
% R6. PANELS (c) AND (d) COULD NOT BE READ ACROSS, and two of (c)'s four
%     entries were unsupported. Three paired rows; "no class of failure,"
%     deleted as contradicting this figure's own panel (b) and the claim box;
%     "an activation-side failure," deleted as a duplicate of panel (b)'s own
%     struck line; (d)'s first entry hedged to what a LINEAR probe's negative
%     licenses. SUPERSEDES C4 in full and the ledger's
%     entries_3_and_4_are_new_in_pass_12.
% R7. THE VERTICAL DIVIDER AND THE DEPTH AXIS WERE PARALLEL HAIRLINES 0.50cm
%     APART. The divider is deleted.
% R8. PANEL (b)'s MODES WERE \footnotesize, i.e. the same size as the panel
%     titles, and their outcome labels were \scriptsize on a different leading,
%     so nothing lined up across the gutter. Both are \scriptsize now, on one
%     leading, on one set of baselines. One leading per size across the figure:
%     0.40cm at \footnotesize, 0.35cm at \scriptsize.
% R9. SMALLER ITEMS: the mirrored brace's pointer aimed at empty white and is
%     now a plain span rule; "$480$" set three digits in URWPalladioL-Roma
%     against Pagella everywhere else and is now text; the prose block's
%     pronoun had no antecedent on the reading line and now names the residual
%     stream; the foot line conceded two counterexamples in its own second
%     clause and now states what the counts are; the word "seven", deleted with
%     the caret label in pass 12, is back on the drawing in that foot line,
%     which is what the still-standing pass-9 layout note asks for.
%
% WHAT PASS 13 OVERRULED, AND WHY
% -------------------------------
%  O1. "REBUILD PANEL (a): DROP THE DEPTH SWEEP, OR SET LAYER HORIZONTALLY, OR
%      SHORTEN THE AXIS BY TIGHTENING THE 8/9 PAIR." Overruled on all three
%      counts. The seven-tick depth scale at true spacing is the figure's only
%      scaled mark and its keep list names it; the 8/9 pair is already at true
%      spacing and 0.30cm apart, the document's minimum label separation, so
%      the span cannot shorten without dropping a numeral or falsifying the
%      grid (design decisions 2 and 2b, both standing); a horizontal layer axis
%      would draw the residual stream running the wrong way and would break the
%      caret's read-here relation to the strip. C10's four rejected removals
%      stand. What pass 13 did take from that finding is the head alignment:
%      panels (a) and (b) now put their first content line on one baseline.
%  O2. "RETITLE (a) 'where the drug name enters, and where the probe reads'."
%      Overruled on measurement. Panel (b)'s title is 9.495cm of the 9.55cm the
%      right half has and sits on the SAME baseline, so panel (a)'s title
%      cannot exceed about 7.0cm; the proposed title measures 8.5cm and the
%      shortest honest variant 7.7cm. The read is named at the axis instead,
%      which is where the style contract wants it.
%  O3. "DRAW THE CARET AS AN OPEN rulegrey CHEVRON." Overruled. rulegrey is
%      furniture and the caret is content -- it is the one mark of the read
%      position, and the keep list names the triangle. Breaking the coincidence
%      and shrinking it to its pass-9 size answers the defect; recolouring it
%      would demote a content mark to furniture.
%  O4. "PUT PANEL (b)'s OUTCOME COLUMN ON PANEL (d)'s ORIGIN (12.40)."
%      Overruled: the longer strike ends at 12.484, so the label column would
%      begin inside the strike. It sits at 13.44, 0.50cm clear of that
%      terminal, which is the separation the pass-9 note asks for.
%  O5. "KEEP THE STRIKES AT 0.90pt IN quietink." Overruled in favour of 0.40pt.
%      0.90pt is the slate's threshold weight whatever its colour, and at 0.90pt
%      against \scriptsize the struck line does not survive at print size --
%      which is the legibility failure the pass-12 ledger already records.
%  O6. "REACH S15's 6.0% INK AND 5.0 CHARACTERS PER 100mm^2." Not reached, and
%      the arithmetic is recorded rather than hidden: pass 13 measures 6.65%
%      and 5.42, against pass 12's 7.95% and 6.05 and pass 11's 11.21% and 9.9.
%      What is left is 697 characters of which 483 are mandated furniture --
%      four panel titles, the axis definition, the definitional prose, the five
%      field names, the count line and the one interval-style disclosure -- and
%      214 are the partition and the two licences, which are the figure. The
%      only cuts left would delete the axis definition that answers the
%      supervisor's first complaint, or one of the two failure modes.
% ===========================================================================
% PASS 14 -- 2026-08-19. THE SECOND REFUTATION PASS.
% ===========================================================================
% Three adversarial refuters read the pass-13 render and returned thirty
% findings. Everything above is left intact; this block states what pass 14
% changed, what it OVERRULED and why, and SUPERSEDES BY NAME every pass-12 and
% pass-13 claim it makes false. Each change is also recorded at the code it
% changes. NO NUMBER IN THIS FIGURE CHANGED VALUE IN PASS 14: the seven depth
% ticks are still y = 2.10 + (L/16)*4.80 -> 2.70/3.30/3.90/4.50/4.80/5.70/6.90,
% byte-identical to passes 11, 12 and 13, and the count line still prints the
% same 480, twelve and forty. The strip translated 0.14cm DOWN in y as one
% rigid body; not one cell edge in x moved, and x is the only axis on which
% that strip carries anything at all.
%
% S1. THE PENDING CAPTION FILE. All three refuters opened with the same
%     blocking finding, and it is correct: figs/_PENDING_CAPTIONS.tex contains
%     no \label{fig:licensing} block. R2 above asserts, in as many words, that
%     "Pass 13 creates figs/_PENDING_CAPTIONS.tex and writes the pass-13
%     caption and every deleted sentence into it"; the file exists and holds
%     fig:materiality, fig:comparators and fig:hook blocks, and this figure's
%     append was lost -- almost certainly clobbered by a concurrent
%     read-modify-write from another builder, which is the hazard CF12 names.
%     So pass 13 committed, one pass later, the identical failure it accused
%     pass 12 of. SUPERSEDES R2's last three sentences and the "duplicate of
%     the pending file" note above the staged float: they described an intent,
%     not the filesystem.
%     Pass 14 APPENDS the block (pure append, never a rewrite, so it cannot
%     clobber another builder's), then READS THE FILE BACK to confirm the block
%     survived, and records the byte offset it landed at in licensing.json. The
%     block carries the pass-14 caption, the four clauses of the live pass-11
%     caption at 04b:117-140 that are false against this drawing with their
%     line numbers, and every sentence passes 12, 13 and 14 cut from the
%     drawing, each against the caption sentence that now carries it.
%
% S2. "no class of fix" IS DELETED FROM PANEL (c), AND THIS IS THE ONE
%     CORRECTNESS DEFECT IN THE PASS-13 DRAWING. Full argument at the code.
%     In one line: 04b:31-36 is a two-member partition in which each member
%     names its own fix class, panel (b) strikes one member and says so on this
%     page, so the survivor and its fix class stand -- and pass 13 had already
%     deleted the sibling entry "no class of failure," for exactly this reason
%     without applying it here. (d)'s "target- and input-side fixes" goes with
%     it, because a row with one side is not a row. Both fix classes move to
%     the caption, which can state the condition a two-word slot cannot.
%     SUPERSEDES R6's list of surviving entries and C4 in full.
%
% S3. THE TWO OUTCOME LEADERS ARE DELETED. Measured: the quietink one left a
%     0.24cm gap at its left end and a 0.10cm gap at its right, touching
%     neither mark nor label, and lay midway between two quietink strikes so
%     the eye read the verdict "struck by the positive," as itself struck; the
%     accent one began 1.397cm PAST the mark it claimed to join. The pass-13
%     ledger recorded the opposite as a measured property, having measured only
%     the first of the two. Colour and exactly shared baselines carry the
%     pairing, which is 4.11's discipline. The outcome column moves 13.44 ->
%     13.00. SUPERSEDES C6's leader clause, O4's arithmetic, and pass 13's
%     lead/.style rationale; lead/.style itself is deleted.
%
% S4. THE CARET IS NAMED AT ITSELF and the three cues that made it an axis
%     arrowhead are broken: it gains a \scriptsize ink direct label "the last
%     prompt position" 0.08cm off its left base, the strip's 0.14cm drop puts
%     0.34cm between its apex and the axis base, and its base widens to 0.28cm
%     so it no longer shares an endpoint with the tick lane. Under the figure's
%     own map it now occupies layers 0.13 to 0.87, below the drawn scale.
%     SUPERSEDES R4 and licensing.json's read_position note.
%
% S5. THE ROTATED SCOPE LINE IS DELETED and "layer" moves off the layer-9 tick.
%     The scope line shared the five-word clause "at the last prompt position"
%     with the caption that actually prints, which S18 forbids, and it named at
%     90 degrees and 1.14cm away the one glyph that had no name of its own.
%     "layer" was centred on the axis midspan, which is the layer-9 position by
%     arithmetic (2.10 + (9/16)*4.80 = 4.80 = (2.70+6.90)/2), so it read
%     "layer 9"; 5.70-to-6.90 is the only tick gap on this axis long enough to
%     hold it clear. SUPERSEDES R5 and C1.
%
% S6. THE SPAN RULE UNDER THE STRIP IS DELETED. It duplicated the strip's own
%     bottom border at the same colour, weight and length 0.12cm away, and it
%     bracketed one schematic prompt under a line counting 480 of them.
%     SUPERSEDES R9's span-rule clause.
%
% S7. PANEL (b)'s TITLE IS A NAME, NOT A SENTENCE: 9.495cm -> 5.821cm, so the
%     widest object on the page is no longer a title that restated 04b:31-32.
%     Panel (c)'s row one loses its wrap and its own S18 collision ("the model
%     did not discard what it was handed" is seven words verbatim from
%     04b:107-108, six lines above the float) and becomes "the drug field is
%     not discarded", one line. The strikes start on the panel origin 7.75
%     rather than 0.06cm left of it. The (c)/(d) group drops 0.30cm and its row
%     pitch opens to 0.70cm, so deleting a row did not trade panel (a)'s hole
%     for a new one in the bottom-right corner. The foot line drops its second
%     clause. SUPERSEDES O2's measurement clause and R8's 0.52cm row gap.
%
% S8. THE PANEL-(a) BAND. The pass-13 drawing put one contiguous 56.2 x 39.2mm
%     void -- 17.2% of the figure -- under the prose block and left of the
%     axis, because pass 13 had moved that block from the band's middle up to
%     the panel head to buy a shared first-content baseline with panel (b).
%     Pass 14 moves it back to 5.00/4.60, centred on the axis midspan, and puts
%     the caret's new label in the band's lower part; the band is now two voids
%     of about 1.9 and 1.6cm and the block sits level with the axis it
%     describes. The two TITLE baselines still align at 7.45, which is the
%     alignment a reader sees. NOTHING FURTHER WAS PUT IN THE BAND, and the
%     reason is the contract, not indifference: S4 allows at most one prose
%     block per panel and panel (a) already has it, and C10's four ways of
%     REMOVING the band (shorten the axis, narrow the strip, invert the axis,
%     move (c)/(d) into it) were each rebuilt and each still fails. The band is
%     the space between the prompt and the read; this figure draws nothing
%     there because the stream's contents are 4.11's subject and the slab's
%     geometry is A.3's. SUPERSEDES pass 13's placement repair (ii).
%
% S9. MEASURED AFTER ALL OF THE ABOVE (200dpi, clipped to the content bbox):
%       content bbox   173.10 x 74.29mm   (S14 wants 170-174; cap 75mm)
%       ink            5.826%             (was 6.65 at pass 13, 11.21 at 11)
%       characters     596 non-space, 4.63 per 100mm^2  (was 5.44 at pass 13)
%                      694 with spaces,  5.40 per 100mm^2  (was 6.41)
%       stroke weights 0.30pt rulegrey furniture (4 ordinary cells, 7 ticks,
%                      the (c)/(d) row rule) and 0.40pt (axis rule ink, drug
%                      cell border accent, two strikes quietink). Two weights,
%                      neither a singleton, no 0.90pt stroke, no arrowhead,
%                      no leader, no dotted rule, no legend.
%       type           two sizes: \footnotesize 10.0pt (four panel titles, the
%                      prose block, the axis quantity name) and \scriptsize
%                      8.0pt (everything else). 0% italic. 0 all-caps banners.
%     S15's ink ceiling is met for the first time in this figure's history.
%     The with-spaces character count is still 5.40 against a 5.0 ceiling and
%     is recorded rather than hidden: the refuters' own measurement of that
%     rule was taken on non-space characters (5.44 against 5.0), which is met
%     at 4.63; closing the with-spaces gap needs 51 more characters and the
%     only strings left are the four titles, the axis name, the five field
%     names, the count line the spec's keep list mandates, the two failure
%     modes and their outcomes, the two licence rows and the one disclosure
%     line. Every one of them is mandated content.
%
% WHAT PASS 14 OVERRULED, AND WHY
% -------------------------------
%  P1. "SET THE SEVEN TICK NUMERALS IN INK AND LENGTHEN THE TICKS TO 0.20cm,
%      BECAUSE THE AXIS READS AS A PANEL DIVIDER." Overruled on both counts.
%      Quietink tick numerals are this corpus's convention, not an oversight:
%      comparators.tex:70-72, gate.tex and purify.tex all set them so, and
%      design decision 10 changed them from accent to quietink in pass 9 with
%      that reason recorded. S8 assigns ink to axis LABELS, which is where
%      "layer" is, and does not assign a colour to tick numerals; making this
%      one figure's numerals black would break it away from the four figures it
%      is being aligned to. The ticks cannot lengthen inward either: they run
%      6.78-6.90 and the numerals are right-anchored on 6.70, so 0.20cm ticks
%      would touch them. The divider reading the finding describes was created
%      by the divider itself and died with it in pass 13 (R7); what is left is
%      a rule with seven numbered ticks and a rotated quantity name, which is
%      an axis.
%  P2. "OPEN A GUTTER BETWEEN THE PROMPT CELLS, OR DROP rounded corners=1pt,
%      BECAUSE THE ROUNDINGS LEAVE WHITE NOTCHES AT EVERY INTERNAL JUNCTION."
%      The notches are real -- re-cropped at 1200dpi and seen -- and the fix is
%      still overruled here. The finding's own evidence is that figs/hook.tex
%      draws the identical construction at the identical radius, and both
%      captions assert the two strips are the same prompt seen twice; a strip
%      repaired in one file and not the other would break an agreement the text
%      states, to remove a 0.1mm artefact. This builder may not edit hook.tex.
%      Recorded in licensing.json as a coordinated repair the two files owe
%      together, so the next pass that touches both can make it once.
%  P3. "SET THE PROSE BLOCK AS ONE LINE AND DROP IT A RANK, BECAUSE IT IS THE
%      SAME SIZE AS THE PANEL TITLES." Overruled, and the finding concedes the
%      contract point itself: S4 mandates \footnotesize for definitional prose
%      and states in terms that prose OUTRANKS labels, which is A.3's ranking
%      and the one the contract chose to keep. A one-line rewrite either drops
%      "the residual stream", which is the antecedent R9 put there and the only
%      naming of the object the axis stands for, or drops "cross-validated",
%      which is why the accuracy is a presence test at all.
%  P4. "RESTORE THE PROBE'S TASK TO THE DRAWING: 'over the twelve drug
%      classes'." Overruled as to the drawing, accepted as to the record. The
%      figure is at its ink ceiling, the count line already prints "twelve
%      drugs" 2.4cm below, and naming the instrument does not require naming
%      its class count. The twelve-class task AND the fold caveat (04b:155-157,
%      cut in pass 12 as D7) are both in the pass-14 caption, and
%      licensing.json records under explicitly_not_drawn that the drawing names
%      the instrument without its target and why.
%  P5. "LEAVE PANEL (c)'s SLOTS BLANK, OR SET AN EN-DASH IN THEM, SO THE SMALL
%      LICENCE LOOKS SMALL." Overruled. An empty slot under a head reading
%      "what the positive buys" is ambiguous between "nothing" and "not drawn
%      yet", and design decision 7's own supersession at C4 gives the reason a
%      named absence beats a blank or the bare word "none". The finding's
%      underlying complaint -- that the two licences were drawn the same size,
%      with the smaller one in darker ink and more lines -- is answered by S2
%      and S7 above: (c) now sets 39 characters against (d)'s 58 and its second
%      entry is 0.991cm against (d)'s 3.233cm.
%  P6. "DRAW THE TIE BETWEEN THE DEPTH AXIS AND (d)'s 'the layer at which it is
%      lost', SO THE AXIS HAS A STATED ROLE." Overruled as a drawn mark. A
%      leader across 5.5cm of figure would be the only connector on the page
%      and would assert an argument, and a line saying "a negative would have
%      named one of these seven depths" is argument, not definition, which S4
%      and S18 put in the caption. The pass-14 caption says it in those words.
%      The axis is now fully DEFINED on the drawing -- quantity name, seven
%      numbered ticks, and the position it is read at, labelled at the caret --
%      which is what S5 asks of it.
%  P7. "MAKE THE VERTICAL RULE THE RESIDUAL STREAM, SO THE PROSE BLOCK HAS A
%      DRAWN REFERENT." Overruled. The rule is a LAYER axis; the residual
%      stream at one position is a 2048-dimensional vector at each of seven
%      depths, and renaming a depth coordinate after the object sampled along
%      it would put two names on one rule and cost S5 its clean reading. The
%      prose block was instead moved level with the axis and the caret named,
%      so the reading relation the block defines -- position, depths, probe --
%      is drawn beside it rather than four centimetres above it.
% ===========================================================================
% inner sep=0pt on every node, so that anchor=base west puts the START OF THE
% INK exactly on the panel's x origin and a node's width is exactly the \wd
% measured for its string. With TikZ's default 0.3333em the picture ran 5.43pt
% past the fullwidth measure (the harness reported exactly that Overfull hbox),
% because panel (b)'s title is 9.495cm of the 9.55cm the right half has.
% PASS 14: that title is now 5.821cm, so inner sep=0pt is no longer what keeps
% the picture inside the measure -- but it is kept, because it is also what
% makes anchor=base west put the START OF THE INK on the panel origin, which is
% what S13's one-x-origin check measures. The widest object is now panel (d)'s
% title, ending at x = 17.247, 0.053cm inside the 17.30cm measure.
%
% ===========================================================================
% v-BLOCK, 2026-08-19: WHOLE-SLATE PASS (fig:comparators, fig:licensing,
% fig:purify, fig:hook, fig:materiality reconciled against one another and
% against fig:mechanism, which is the document's reference standard and was
% NOT rebuilt). Nothing above this line is deleted; where an earlier claim is
% now false it is superseded EXPLICITLY below.
%
% WHY. The five figures were rebuilt independently and drifted. This pass owns
% only the differences BETWEEN them. No value changed in any of the five; the
% sibling .json ledgers carry the value-to-position maps, old and new.
%
% THE SLATE RULES THAT NOW HOLD ACROSS ALL FIVE (each stated once here, in
% every file, so any future single-figure edit can see what it would break):
%
%  R1 TICK NUMERALS are quietink, \scriptsize, inner sep 0, anchored north on
%     the tick and set 0.16cm below the axis rule, over a 0.10cm rulegrey
%     0.30pt tick. Before: comparators/licensing/purify quietink, hook and
%     materiality ink; ticks 0.10 / 0.12 / 0.14cm; three different gaps.
%  R2 A THRESHOLD'S OWN LABEL IS BLACK, matching the black solid 0.90pt rule
%     it names (S6). Before: comparators black, purify's "gate, 0.8" and
%     materiality's two margin labels ink.
%  R3 ONE NAME PER OBJECT (S19). The pre-registered 10% rule is "10%
%     materiality margin" in fig:comparators row (c), fig:materiality (a) and
%     (b), and in figs/family.py, which is where the name comes from.
%     fig:comparators' "pre-registered margin" is superseded; "pre-registered"
%     is now the caption's word, not the drawing's. Zero, where it is the null
%     a bar is read against, is "zero: no identity effect" in BOTH
%     fig:comparators row (c) and fig:materiality (b) -- one quantity, one
%     wording. V_pure is spelt with \textrm everywhere, as figs/slab.tex does.
%  R4 COLOUR KEY, one key for five figures (S8). accent = the drug-side object
%     under test. quietink = every comparator, null, control, replication or
%     discounted arm, AND its label, AND its leader. ink/black = thresholds,
%     axis rules, axis labels, panel heads, definitional prose. rulegrey =
%     furniture only. Within quietink the GLYPH separates two kinds: an OPEN
%     disc is a construction built to carry no signal (a null, a control, a
%     raw/before series); a FILLED disc is a real measurement that is not the
%     headline. fig:hook (d) drew both its null and its cell-line control in
%     ink and is corrected.
%  R5 PROJECTION AND ZERO RULES are quietink 0.50pt densely dotted (S7).
%     fig:hook's two panel-(c) projection lines were rulegrey 0.30pt densely
%     dashed and are corrected.
%  R6 DEFINITIONAL PROSE is \footnotesize ink and OUTRANKS every label (S4);
%     at most one block per figure, on the panel's own x origin. fig:hook's
%     one prose line was \scriptsize and is raised. fig:purify's panel-(b)
%     block carried an argument clause ("so causal claims use the purified
%     slab, not the raw") that its caption already carries, and is now
%     definitional only.
%  R7 ONE INTERVAL DECLARATION per figure, \scriptsize quietink, at the foot,
%     on the figure's x origin, opening with the word "Intervals:" and naming
%     EVERY panel including the ones with no measured mark (S17).
%     fig:comparators already did this and is the model; the other four now
%     match its form.
%  R8 THE FIVE-FIELD PROMPT STRIP is ONE object drawn in two figures, and the
%     two drawings are now congruent: cell boundaries 0 / 1.12 / 1.90 / 2.64 /
%     4.23 / 6.90 cm and cell height 0.42cm in both fig:licensing (a) and
%     fig:hook (a) and (b), labels centred in their cells, the drug cell
%     accent-bordered at 0.40pt on accent!12 and its label NOT bold. Before:
%     licensing 6.90cm wide with 0.62cm cells and a bold label, hook 7.44cm
%     wide with 0.42cm cells and a plain label -- the same five fields at two
%     sizes, four pages apart, with each caption pointing at the other.
%  R9 MEASURE. Every fullwidth figure draws 172.0-174.0mm; the one textwidth
%     figure draws 140.3mm. Ink at 200dpi inside the content bbox, counted as
%     greyscale luminance below 230 (the contract's own method), is now
%     5.63-5.99% on all five, under the 6.0% body ceiling, where before this
%     pass it ran 5.36-6.54%.
%
% WHAT THIS PASS DELIBERATELY DID NOT UNIFY, so that a later reader does not
% think it was missed:
%  * PANEL-HEAD CASE. fig:comparators sets "(a) Is it there?" and the other
%    four set "(a) where the drug name enters". comparators' heads are
%    interrogative SENTENCES and sentence case is correct for a sentence; the
%    v4 header argued this and the argument stands. fig:mechanism's lowercase
%    heads are noun phrases, not questions. One capital letter, four times.
%  * PANEL-HEAD POSITION. fig:comparators puts its heads in a right-aligned
%    left stub, the other four above the panel at its x origin. Moving them
%    above the rules costs about 20mm of height against a 118mm cap that the
%    figure already meets at 117.60mm. Not available.
%  * TWO-COLUMN GUTTER. purify and hook split at x = 8.40cm, licensing at
%    7.75cm. licensing cannot move right: "(d) what a negative would buy" is
%    4.85cm wide and its origin is already at 12.40 of a 17.30cm measure, with
%    1.5pt of slack. Moving purify and hook LEFT to 7.75 would be arbitrary.
% ===========================================================================
%
% CHANGED IN THIS FILE. (i) The prompt strip adopts fig:hook's cell height,
% 0.42cm, and fig:hook drops to this file's cell WIDTHS, so the two drawings
% of the one object are congruent (R8); the strip's top stays at y = 2.14, so
% the caret, the layer axis and everything above them are untouched, and the
% strip's bottom rises from 1.52 to 1.72. Labels move from anchor=base at
% y = 1.775 to centred at the cell's own mid-height, 1.92, with inner sep 0,
% which is \tokbox's placement. (ii) The drug cell's label loses \bfseries: it
% was the only bold glyph on the slate, and the accent colour plus the accent
% border already carry the emphasis. (iii) The layer ticks go from 0.12 to
% 0.10cm and the numerals from 6.70 to 6.74, a 0.06cm gap (R1). (iv) The
% "480 real tier-2 prompts" line goes ink -> quietink, matching the same class
% of line in fig:purify, and rises to y = 1.38 with the strip. (v) The prose
% block moves from y 5.00/4.60 to 5.55/5.15 so its first line sits on panel
% (b)'s second-mode baseline, 5.55, giving the two columns a shared
% horizontal. (vi) The foot line takes the slate's declaration form (R7).
% NOT FIXED, and it is this figure's real weakness: panel (a) is an L -- a
% 4.20cm vertical layer axis at the column's right edge and a 6.90cm strip
% along its bottom -- so about 6.0 x 3.4cm of the column is empty. The pass-12
% header records four rejected fills and this pass adds no fifth; the void is
% the (position x depth) plane in which this figure measures nothing, and the
% prose block sits in it to break it rather than to disguise it.
% RE-MEASURED: 173.10 x 74.30mm; ink 5.975%; 726 characters, 5.65 per 100mm^2,
% which is over S15's 5.0 and is the one contract clause this figure still
% misses -- see pass 12's arithmetic, unchanged by this pass. Exactly two type
% sizes, 9.96 and 7.97, in one face, no bold, no italic, no sans, no CM.
% One page, zero Overfull, zero Underfull.
% ===========================================================================
\begin{tikzpicture}[x=1cm, y=1cm, font=\small, inner sep=0pt]
  \tikzset{
    % 4.11's panel head: lowercase parenthesised letter, sentence-case
    % descriptive title, \footnotesize roman ink, left on the panel's origin.
    ttl/.style ={font=\footnotesize, text=ink, anchor=base west},
    % the figure's ONE prose register: definitional only, \footnotesize ink.
    def/.style ={font=\footnotesize, text=ink, anchor=base west},
    % the surviving branch -- the object the rest of the chapter measures.
    % PASS 13: \scriptsize, not \footnotesize. These are the two branches of a
    % partition, i.e. direct labels of drawn objects, not definitional prose;
    % under S1 the largest size in the figure must belong to a panel title, an
    % axis label or definitional prose, and 4.11 sets its in-plot labels one
    % step below its titles. They now share a size and a baseline grid with
    % their own outcome labels, which is what fixed the cross-gutter splay.
    liv/.style ={font=\scriptsize, text=accent, anchor=base west},
    % the struck branch, and every counterfactual.
    ded/.style ={font=\scriptsize, text=quietink, anchor=base west},
    % \scriptsize registers: labels, entries, scope and foot lines.
    lab/.style ={font=\scriptsize, text=ink, anchor=base west},
    qlab/.style={font=\scriptsize, text=quietink, anchor=base west},
    alab/.style={font=\scriptsize, text=accent, anchor=base west},
    tick/.style={font=\scriptsize, text=quietink, anchor=east},
    % The prompt cell is drawn as figs/hook.tex draws it (\tokbox): panelbg
    % fill, rounded corners=1pt. PASS 12: the border is 0.30pt, the document's
    % furniture weight, not the 0.35pt of pass 9, which is not one of the five
    % weights the slate uses.
    cell/.style={draw=rulegrey, line width=0.30pt, fill=panelbg,
                 rounded corners=1pt},
    % The two outcome connectors were IDENTICAL plain leaders. PASS 13 deleted
    % the Stealth heads: a single arrowhead is reserved for an operation acting
    % in a direction and its label must be a verb of action, and "struck by the
    % positive" and "survives" are outcomes, not operations.
    % PASS 14 DELETES lead/.style AND BOTH LEADERS. Reasoning at the two mode
    % blocks and in the pass-14 header at S3. The figure now contains no
    % arrowhead and no connector of any kind, and its only glyph that could be
    % mistaken for either -- the caret -- carries a direct label at itself.
    }

% ---------------------------------------------------------------------------
% ONE X ORIGIN PER PANEL. Every element below references its panel's origin,
% so the spread of text-line x0 inside a panel is exactly zero.
% ---------------------------------------------------------------------------
  \coordinate (OA) at (0.00,0);    % panel (a)
  \coordinate (OB) at (7.75,0);    % panels (b) and (c)
  \coordinate (OD) at (12.40,0);   % panel (d)
  \coordinate (OBL) at (13.44,0);  % panel (b)'s outcome column

% PASS 13: THE VERTICAL PANEL DIVIDER IS DELETED. It stood at x = 7.40 and the
% depth axis stands at x = 6.90, so the two were parallel hairlines 0.50cm
% apart with an empty channel between them; measured on the render, the band
% between them carried 0.16% ink over the full height. At print size they read
% as one doubled boundary, and panel (a)'s only scaled object was thereby read
% as frame furniture. The lettered panel heads and a 0.85cm gutter partition
% the figure without it, which is 4.11's own arrangement -- it draws no
% dividers at all. SUPERSEDES C5's divider clause and its gutter arithmetic.

% ###########################################################################
% # PANEL (a) -- where the drug name enters
% ###########################################################################
  \node[ttl] at (0.00,7.45) {(a) where the drug name enters};

  % --- the figure's only prose block: definitional, two lines -------------
  % PASS 13, TWO REPAIRS.
  %  (i) THE PRONOUN HAD NO ANTECEDENT ON THE READING LINE. Pass 12 read "what
  %      reads it: ...", and the only occurrence of "residual stream" was
  %      inside the rotated \scriptsize scope line, 4cm away and at 90 degrees.
  %      The block now names the object it defines. "twelve classes" leaves the
  %      block: the line "a cross-validated linear probe, twelve classes"
  %      measures 6.68cm and the panel's prose is capped at 5.48cm by the
  %      rotated label column at x = 5.62 (see the axis block below). "linear"
  %      is the word worth the space, because it is what makes panel (d)'s
  %      "the drug not linearly readable" the right counterfactual, and twelve
  %      still reaches the reader through the count line under the strip.
  % (ii) IT MOVED FROM THE PANEL'S VERTICAL MIDDLE TO THE PANEL'S HEAD, so that
  %      panels (a) and (b) put their first content line on ONE baseline, 6.70,
  %      0.75cm under their titles. Pass 12 set it at 5.90/5.50, which gave
  %      panel (a) a 1.55cm head-to-content gap against panel (b)'s 0.75cm and
  %      put panel (a)'s only substantial ink level with panels (c) and (d), so
  %      the lettering order and the spatial order disagreed.
  %      SUPERSEDES C10's placement paragraph ("three tiers with voids of 1.55
  %      and 3.22cm"); C10's four rejected ways of REMOVING the band all still
  %      stand, and the band is still deliberately empty, because it is the
  %      space between the prompt and the read, i.e. the residual stream.
  %      Clearance measured: line one is 4.880cm and line two 4.592cm, both
  %      clear of the rotated scope column's left edge at x = 5.62 by more than
  %      0.7cm at every y.
  % PASS 14: THE BLOCK MOVES BACK INTO THE BAND, at 5.00/4.60, centred on the
  % axis midspan 4.80. Pass 13's (ii) above bought a shared first-content
  % baseline with panel (b) and paid for it with ONE CONTIGUOUS VOID of
  % 56.2 x 39.2mm -- 17.2% of the figure's whole area, measured at 200dpi --
  % under the block and left of the axis, so panel (a) read as text at the top,
  % a stranded vertical stick in the middle and a hole between them. The two
  % TITLE baselines still align at 7.45, which is the alignment a reader
  % actually sees; the first-content alignment across a 7.75cm gutter is not.
  % Set here the block also sits level with the axis it describes, so "the
  % residual stream" has the rule and its seven ticks beside it rather than
  % 4cm above them. The band is now two voids of about 1.9 and 1.6cm rather
  % than one of 3.5cm, and its lower part carries the caret's direct label.
  % SUPERSEDES pass 13's repair (ii) at this block and C10's placement
  % paragraph as re-stated there; C10's four rejected ways of REMOVING the
  % band all still stand, and S4's "at most one block per panel" is why no
  % second prose block was written into it.
  \node[def] at (0.00,5.55) {what reads the residual stream:};
  \node[def] at (0.00,5.15) {a cross-validated linear probe};

  % --- the depth axis ------------------------------------------------------
  % PASS 12: a plain 0.40pt ink rule from the layer-2 tick to the layer-16
  % tick, terminating exactly on them. Before this pass it was a panelbg-filled
  % 0.36cm-wide column running from 2.10 to 7.05, i.e. past both extreme ticks,
  % which read as a bar whose length encoded nothing and implied depth beyond
  % the measured range.
  \draw[ink, line width=0.40pt] (6.90,2.70) -- (6.90,6.90);
  % y = 2.10 + (L/16)*4.80 for every one of the seven, unchanged from pass 11
  % and CHECKED against the map rather than against the eye:
  %   L   2    4    6    8    9    12   16
  %   y  2.70 3.30 3.90 4.50 4.80 5.70 6.90
  % Tick length 0.12cm = 3.4pt, PASS 13, down from 0.16cm: 4.11's own ticks are
  % about 3.5pt and the 0.04cm bought here is what lets the quantity name sit
  % beside the numerals without touching them.
  \foreach \L/\y in {2/2.70, 4/3.30, 6/3.90, 8/4.50, 9/4.80, 12/5.70, 16/6.90}{
    \draw[rulegrey, line width=0.30pt] (6.80,\y) -- (6.90,\y);
    \node[tick] at (6.74,\y) {\L};}
  % AXIS LABEL, in axis-label position: rotated 90 degrees reading upward,
  % centred on the axis midspan y = 4.80. PASS 13, TWO REPAIRS.
  %  (i) THE RANKING IS INVERTED BACK. Pass 12 put the quantity name OUTERMOST,
  %      at x = 5.36, with two wrapped \scriptsize grey lines between it and the
  %      numerals, so the thing immediately left of the tick numerals was a grey
  %      sentence and the quantity itself was 1.4cm further out -- the opposite
  %      of what S5 asks for. "layer" now sits immediately left of the numerals
  %      and the scope note outboard of it.
  % (ii) THE SCOPE NOTE IS ONE LINE, NOT A WRAPPED TWO. Pass 12's note measured
  %      9.09cm, wrapped onto two typeset lines of 4.817 and 4.202cm, and the
  %      longer of the two OVERRAN a 4.20cm axis at both ends -- 3.0mm past the
  %      layer-16 tick above and the same below -- which is the exact defect D6
  %      records itself as having removed from the column. It also broke the
  %      noun phrase across the wrap ("the last prompt / position"). The note is
  %      now "read at the last prompt position", 111.20pt = 3.908cm, which lies
  %      INSIDE the 4.20cm axis with 0.146cm of clearance at each end. "where
  %      generation begins" is caption material and is in the caption.
  %      SUPERSEDES C1's wrap arithmetic and its "line 1 is the outermost"
  %      clause in full.
  % x geometry, measured: at inner sep=0 a rotated \footnotesize line is 0.35cm
  % wide in x and a \scriptsize one 0.28cm; the widest numeral ("16") is 8.0pt =
  % 0.281cm and is right-anchored on 6.70, so it begins at x = 6.419.
  % scope 5.62-5.90 | layer 6.005-6.355 | numerals 6.419-6.70 | ticks 6.78-6.90.
  % PASS 14, TWO REPAIRS AT THIS BLOCK.
  %  (i) THE SCOPE LINE IS DELETED AND ITS CONTENT MOVES TO THE MARK IT NAMES.
  %      "read at the last prompt position" shared the five-word clause "at the
  %      last prompt position" with the caption that actually prints
  %      (04b:124-126, "the probe reads the residual stream at the last prompt
  %      position, where generation begins"), which S18 forbids; and it named,
  %      at 90 degrees and 1.14cm away, the one glyph on the page that had no
  %      name of its own -- the caret. 4.11's discipline is to direct-label at
  %      the mark, so the position is now labelled AT the caret, horizontally,
  %      in \scriptsize ink, and the axis keeps only its quantity name. S5
  %      makes the scope line optional; the mandatory \footnotesize roman ink
  %      quantity name in axis-label position is still there and is still the
  %      only thing between the tick numerals and panel (a)'s body.
  %      SUPERSEDES pass 13's repair (ii) at this block in full.
  % (ii) "layer" MOVES OFF THE LAYER-9 TICK. The axis midspan and the layer-9
  %      position are the SAME y by arithmetic, not by accident:
  %      2.10 + (9/16)*4.80 = 4.80 = (2.70 + 6.90)/2. So a rotated label
  %      centred on the midspan was centred on one numeral, 1.96mm from it, and
  %      read "layer 9". Measured, "layer" is 21.46pt = 0.754cm long, and the
  %      seven tick gaps are 0.60/0.60/0.60/0.30/0.90/1.20cm, so 5.70-to-6.90
  %      is the ONLY stretch of this axis long enough to hold the label with no
  %      numeral inside its span. Centred at 6.30 it spans 5.923 to 6.677 and
  %      the nearest numerals are 12 at 5.70 and 16 at 6.90. It stays in
  %      axis-label position: x 6.005-6.355, immediately left of the numeral
  %      column at 6.419-6.70, inside S5's 12pt band either side of that column.
  \node[font=\footnotesize, text=ink, rotate=90, anchor=center]
    at (6.18,6.30) {layer};

  % --- the prompt strip ----------------------------------------------------
  % Edges from the measured field-name widths: 6.084cm of ink, 0.816cm of
  % padding shared as 0.0816cm a side. See C8.
  % PASS 14: THE WHOLE STRIP DROPS 0.14cm, from y 1.66-2.28 to y 1.52-2.14.
  % NO CELL EDGE IN x MOVED, no field changed, no order changed and the strip
  % is still 6.90cm wide with the same five measured widths -- the change is
  % one rigid translation in y, and it buys the 0.14cm the caret's direct label
  % needs between the strip's top rule and the layer-2 numeral (see the read
  % position below). The room exists because pass 14 deletes the span rule that
  % occupied 1.46-1.54; the count line's baseline at 1.18 is unmoved and its
  % ascenders now clear the strip's bottom rule by 0.14cm.
  \draw[cell] (0.00,1.72) rectangle (1.12,2.14);
  \draw[cell] (1.90,1.72) rectangle (2.64,2.14);
  \draw[cell] (2.64,1.72) rectangle (4.23,2.14);
  \draw[cell] (4.23,1.72) rectangle (6.90,2.14);
  % the drug cell: the only accent mark in panel (a). Same construction as
  % hook.tex's \tokdrug -- accent border, accent!12 fill, rounded corners=1pt.
  \filldraw[fill=accent!12, draw=accent, line width=0.40pt, rounded corners=1pt]
    (1.12,1.72) rectangle (1.90,2.14);

  % baseline 1.775 = the strip's vertical centre 1.83 less 0.055, which centres
  % the x-height band rather than the box. (Pass 13: 1.915 on a centre of 1.97.)
  \node[font=\scriptsize, text=ink,    inner sep=0pt] at (0.560,1.92) {cell line};
  \node[font=\scriptsize, text=accent, inner sep=0pt] at (1.510,1.92) {drug};
  \node[font=\scriptsize, text=ink,    inner sep=0pt] at (2.270,1.92) {dose};
  \node[font=\scriptsize, text=ink,    inner sep=0pt] at (3.435,1.92) {mechanism};
  \node[font=\scriptsize, text=ink,    inner sep=0pt]
    at (5.565,1.92) {control cell sentence};

  % --- the read position ---------------------------------------------------
  % The caret straddles the strip's right edge, which is where the prompt ends.
  % It is ink, not accent: it is part of the axis apparatus, and accent in this
  % figure means the drug field and the surviving branch and nothing else.
  % PASS 13: IT NO LONGER TOUCHES THE AXIS AND NO LONGER SITS ON A NUMBERED
  % TICK. Pass 12 put its apex at (6.90, 2.70), which is simultaneously the
  % axis rule's base, the layer-2 tick and the largest filled black glyph in
  % panel (a); measured on the render the apex and the layer-2 tick coincided
  % to within a pixel. Three objects at one coordinate read two ways, both
  % false: as a filled arrowhead terminating the depth axis downward, and as a
  % mark AT LAYER 2 on the figure's only scaled axis -- which also falsified
  % the foot line, since the caret would then occupy a position on the depth
  % scale. It marks a token position, not a depth. The apex is now 2.50, a
  % clear 0.20cm below the axis base at 2.70, and the triangle is back to the
  % 0.24cm width of pass 9 (base 6.78 to 7.02, so it still straddles the
  % strip's right edge at 6.90 symmetrically and overhangs it by 0.12cm rather
  % than 0.15cm). It is named by the axis scope line 0.9cm to its left, which
  % now reads "read at the last prompt position".
  % SUPERSEDES C8's clause "the apex is the axis rule's own base at the layer-2
  % tick, y = 2.70", and restores the pass-9 note's own reading that the apex
  % belongs at a position the depth scale does not number.
  % PASS 14: THE CARET IS NAMED AT ITSELF, AND THE THREE CUES THAT MADE IT AN
  % ARROWHEAD ARE BROKEN ONE BY ONE. Pass 13 moved it 2mm and did not resolve
  % the reading; all three refuters read it, unaided, as the axis's terminal
  % arrowhead, and two of them read the layer-2 numeral as its label. Measured
  % on the pass-13 render, three cues did that: (1) it was the only
  % arrowhead-shaped glyph left in the figure and the largest solid black one
  % in the panel; (2) its apex stood on the axis rule's own x, 2.00mm below the
  % rule's bare square terminus; (3) its base ran 6.78 to 7.02 while the seven
  % ticks run 6.78 to 6.90, so its left half lay exactly in the tick lane, and
  % the "2" numeral's box bottom sat 0.25mm above its apex. Pass 14:
  %   * it acquires a \scriptsize ink DIRECT LABEL at itself, "the last prompt
  %     position", right-anchored at 6.68 so its ink stops 0.08cm from the
  %     caret's left base -- inside the document's ~0.3cm leader threshold, so
  %     no leader is owed. This is where that string was always meant to be:
  %     it was set rotated in axis-label position until pass 14, where a reader
  %     attached it to the axis rather than to the triangle;
  %   * the strip drops 0.14cm, so the apex is now 0.34cm below the axis base
  %     (was 0.20cm) and the label sits in the gap that opens;
  %   * the base widens from 0.24 to 0.28cm, 6.76 to 7.04, so it no longer
  %     shares an endpoint with the tick lane. It still straddles the strip's
  %     right edge at 6.90 symmetrically, which is the pass-9 relation the keep
  %     list names, and it is still ink, not rulegrey (O3 stands: rulegrey is
  %     furniture and this mark is content).
  % Under the figure's own map y = 2.10 + (L/16)*4.80 the triangle now occupies
  % y 2.14 to 2.36, i.e. layers 0.13 to 0.87 -- BELOW layer 1, off the drawn
  % scale, which is what "it marks a token position, not a depth" requires and
  % what the pass-13 ledger claimed without arithmetic. SUPERSEDES pass 13's
  % apex/base clause at this block and licensing.json's read_position note.
  \fill[ink] (6.90,2.36) -- (6.76,2.14) -- (7.04,2.14) -- cycle;
  \node[font=\scriptsize, text=ink, anchor=base east] at (6.68,2.28)
    {the last prompt position};

  % --- the span: the five fields are one prompt ----------------------------
  % PASS 13: A PLAIN SPAN RULE WITH DOWNTURNED ENDS REPLACES THE MIRRORED
  % BRACE. A brace has a central pointer, and pass 12 moved the count line from
  % centred-on-the-brace to left-aligned on the panel origin (S13) without
  % moving the pointer, so the notch at x = 3.45 aimed at 3.45cm of empty white
  % while its label began at x = 0.00. A span rule carries the same "these five
  % fields are one prompt" grouping, has no pointer to aim wrongly, and drops
  % the decorations.pathreplacing dependency.
  % PASS 14: THE SPAN RULE IS DELETED. All three refuters read it and all three
  % called it. Measured on the pass-13 render it was a rulegrey 0.30pt rule of
  % the strip's exact x-extent, 0.12cm below and parallel to the strip's own
  % rulegrey 0.30pt bottom border -- same colour, same weight, same length --
  % so it reproduced as a doubled edge or an offset re-impression rather than
  % as a bracket, its 0.08cm downturns vanishing at print size. Its left
  % downturn also ended 0.78mm directly above the "4" of "480" and read as a
  % stray tick on the numeral. And it grouped ONE schematic prompt's five
  % fields while the line beneath it counts 480 prompts, so on a strip that
  % carries no horizontal scale a full-width bracket labelled 480 invited
  % reading the width as the count. Nothing is lost: the five cells abut and
  % already read as one object, and the count line sits under the strip on the
  % strip's own left edge and names the population it is drawn from.
  % SUPERSEDES pass 13's R9 span-rule clause and the "the five fields are one
  % prompt" heading above it; the 0.14cm it frees is what lets the strip drop.
  % The count line names the population the spanned object is drawn from. It is
  % left-aligned to panel (a)'s origin, so panel (a) has exactly one x origin.
  % \scriptsize: at \footnotesize it measures 8.196cm and would run past the
  % panel's own measure. PASS 13: 480 is set as TEXT, not as $480$. In math
  % mode the three digits rendered in URWPalladioL-Roma while every other
  % numeral in the figure -- the seven tick labels, and the 2 of "tier-2" on
  % this very line -- rendered in TeXGyrePagella-Regular. One figure, two digit
  % faces, three characters. The measured width is 186.58pt = 6.5575cm either
  % way, so no coordinate moved and no value changed.
  \node[qlab] at (0.00,1.38)
    {480 real tier-2 prompts: twelve drugs, forty cells each};

% ###########################################################################
% # PANEL (b) -- the two ways a model can ignore a drug it was told about
% ###########################################################################
  % PASS 14: THE TITLE IS A NAME, NOT A SENTENCE. "the two ways a model can
  % ignore a drug it was told about" measured 270.16pt = 9.495cm, which made a
  % TITLE the widest object on the page, set the figure's whole right measure,
  % left 0.055cm of optical margin, and outweighed panel (a)'s title three to
  % one on the baseline they share, so the two heads did not read as peers. It
  % also restated 04b:31-32's own opening sentence ("A model that ignores a
  % drug it was told about can fail in two ways"), which is the paragraph the
  % float is attached to. 4.11's two titles run 29 and 38 characters; this one
  % ran 60. "two ways to ignore a named drug" is 165.62pt = 5.821cm, keeps the
  % "was told about" sense in the word NAMED -- which is also what panel (a)'s
  % title says enters -- and ends at x = 13.571. The full sentence is in the
  % caption. SUPERSEDES pass 13's O2 in its measurement clause: panel (a)'s
  % title is no longer constrained by this one. It is left at 4.969cm anyway,
  % because every longer variant either overruns panel (a)'s own 6.90cm measure
  % or says the probe reads the drug name, which it does not -- it reads the
  % stream, and the read is named at the caret and at the axis.
  \node[ttl] at (7.75,7.45)
    {(b) two ways to ignore a named drug};

  % --- failure mode one: struck --------------------------------------------
  % quietink, because it is eliminated. TWO strikes, one through each line: the
  % entry is two lines and a rule through its first line only left the second
  % standing. Each overhangs its line's ink by 0.06cm at each end, which is what
  % a strikethrough does. Endpoints from the measured \scriptsize widths,
  % 4.6741 and 4.2918cm: 7.69 to 12.484 and 7.69 to 12.102.
  % PASS 13, THE COLOUR AND THE WEIGHT. The strikes were ink at 0.90pt, and
  % that is the register S6 and CF2 reserve, across the whole slate, for a
  % threshold: chance in 4.11(a), the removal gate in 4.9, the pre-registered
  % materiality margin in 4.12 and 4.13, and style_v2.chance_line()'s own
  % hard-coded color='black', lw=0.9. C6's defence was that this figure
  % contains no threshold so nothing collides -- which answers S10's singleton
  % clause but not S6's, whose check is that NOTHING WHICH IS NOT A THRESHOLD
  % may be black solid 0.90pt. This figure prints six pages before the reader
  % meets the first real threshold in the arc, and it is the first figure of
  % the slate, so it was teaching the threshold stroke on a negation. Measured
  % against it: the struck text is quietink (grey 107) and the strikes were ink
  % (grey 25), so the branch the chapter DISCARDS carried the darkest ink on
  % the page while the surviving branch receded -- the S8 inversion C3 fixed in
  % colour, reintroduced in weight. Both strikes are now quietink 0.40pt, the
  % colour and weight of the text they cross. 0.40pt is already live here (the
  % axis rule and the drug cell's border), so no new register appears and no
  % singleton is created; 0.50pt was rejected because S10 reserves it for a
  % dotted coordinate. SUPERSEDES C6's third weight in full: the figure's
  % weights are now 0.30pt furniture and 0.40pt, and no 0.90pt stroke exists.
  % Strike y is the line's baseline + 0.065, half the \scriptsize x-height.
  % PASS 14, TWO REPAIRS.
  %  (i) THE STRIKES START ON THE PANEL'S OWN ORIGIN. They began at 7.69, which
  %      is 0.06cm (1.7pt) LEFT of the 7.75 every other text line in panels (b)
  %      and (c) starts on, so the panel's left margin was ragged and the item
  %      protruding furthest left was the one the figure discards. A
  %      strikethrough needs overhang at its TERMINAL, where the eye leaves the
  %      word, not at its initial; the right ends are unchanged at 12.484 and
  %      12.102, still 0.06cm past their lines' ink at 12.424 and 12.042.
  % (ii) THE OUTCOME LEADER IS DELETED (both of them; see the surviving mode).
  %      Measured on the pass-13 render this one ran 12.72 to 13.34 with a
  %      0.24cm gap at its left end and a 0.10cm gap at its right, so it
  %      touched neither the strike nor the label and read as a free-floating
  %      dash. Worse, it was a THIRD quietink horizontal lying midway between
  %      two quietink strikes 0.35cm apart, at the same colour and near the
  %      same weight, so the eye completed strike + gap + leader + label into
  %      one continuous rule and read the verdict "struck by the positive," as
  %      itself struck -- the exact "strike and leader read as one continuous
  %      pointer" fault the pass-9 note records itself as having killed on the
  %      diagonal. Widening the gap did not kill it because the stroke was
  %      still collinear in band, colour and register. Nothing is owed under
  %      S22: there is no MARK here, only a text block, and the pairing is
  %      carried twice over -- by colour (grey block to grey label, accent
  %      block to accent label) and by baselines shared to 0.00pt, which is
  %      4.11's own discipline. SUPERSEDES pass 13's lead/.style rationale and
  %      C6's "the two outcome leaders" clause: the figure now has one stroke
  %      weight for furniture (0.30pt) and one for the strikes and the axis
  %      (0.40pt), and no connector of any kind.
  \node[ded] at (7.75,6.70) {the drug never enters its computation};
  \node[ded] at (7.75,6.35) {an activation-side encoding failure};
  \draw[quietink, line width=0.40pt] (7.75,6.765) -- (12.484,6.765);
  \draw[quietink, line width=0.40pt] (7.75,6.415) -- (12.102,6.415);
  \node[qlab] at (13.00,6.70) {struck by the positive,};
  \node[qlab] at (13.00,6.35) {and only this one};

  % --- failure mode two: survives ------------------------------------------
  % accent, because it is the object the rest of the chapter measures. The two
  % leaders are identical in weight; only the label differs.
  % PASS 13: EACH OUTCOME LABEL NOW SITS ON ITS MODE'S OWN BASELINES, 6.70/6.35
  % and 5.55/5.20. Pass 12 set the modes \footnotesize on 0.40cm leading and
  % their labels \scriptsize on 0.30cm leading, each block centred on the other,
  % so no line matched its partner and the pair splayed in opposite directions
  % across a 55mm gutter -- measured, 0.41mm apart at the first line and 0.60mm
  % the other way at the second. One leading per size now holds across the whole
  % figure: 0.40cm at \footnotesize, 0.35cm at \scriptsize.
  % The outcome column moved from 14.28 to 13.44, which is what the \scriptsize
  % modes freed. Its leader starts 0.50cm clear of the longer strike's terminal
  % (12.484), against 0.17cm in pass 12, so strike and leader cannot read as one
  % continuous pointer -- the fault that killed the pass-1 diagonal.
  % PASS 14: THIS LEADER IS DELETED TOO, AND IT WAS THE WORSE OF THE PAIR.
  % Measured on the pass-13 render, the accent mode's longest line ends at
  % x = 11.328 and the leader began at 12.72, so it started 1.397cm -- 39.6pt,
  % 14mm -- past the mark it claimed to join and covered 0.61 of a 2.13cm gap:
  % at print size a 6.1mm dash floating in 14mm of white with nothing on its
  % left. The pass-13 ledger recorded the OPPOSITE as a measured property
  % ("the leader now starts 0.236cm past the longer strike's terminal"), which
  % is true of the struck mode's leader only; the accent one's 1.397cm gap was
  % never measured. Two leaders of visibly unequal reach were the alternative
  % and were rejected: the struck one cannot be extended leftward without
  % butting the strike and restoring the continuous-pointer reading.
  % The outcome column moves 13.44 -> 13.00, which halves the unlinked gap on
  % this mode from 2.11 to 1.67cm and still clears the longer strike's terminal
  % (12.484) by 0.516cm, the 0.5cm separation the pass-9 note asks for.
  % SUPERSEDES pass 13's O4 arithmetic (13.44) and its "the two leaders are
  % identical in weight" clause -- there are no leaders.
  \node[liv] at (7.75,5.55) {it enters and nothing reads it};
  \node[liv] at (7.75,5.20) {a decoder failure};
  \node[alab] at (13.00,5.55) {survives, and is the rest};
  \node[alab] at (13.00,5.20) {of the chapter};

% ###########################################################################
% # PANELS (c) AND (d) -- the two licences, as two columns read across
% ###########################################################################
  % PASS 14: THE WHOLE (c)/(d) GROUP DROPS 0.30cm AND ITS ROW PITCH OPENS FROM
  % 0.52 to 0.70cm. Deleting the third row (see the entries below) left the
  % right half ending at y = 2.55 against panel (a)'s count line at 1.18, i.e.
  % a fresh 9.5 x 1.8cm void in the bottom-right corner -- trading one hole for
  % another is not a repair. Moving the rule to 4.05 also makes the gap from
  % panel (b)'s last line (5.20) to the rule 1.15cm, which is exactly the gap
  % between panel (b)'s own two mode blocks, so the right half now has one
  % vertical rhythm. The title still sits 0.45cm under the rule, as before.
  \draw[rulegrey, line width=0.30pt] (7.75,4.05) -- (17.30,4.05);

  % PASS 13: BOTH TITLES ARE ONE LINE ON ONE BASELINE. (d) read "what a
  % negative would have bought", 6.167cm, which does not fit a column and
  % wrapped, so (c)'s head-to-body gap was 0.90cm against (d)'s 0.54cm on two
  % columns the ledger called equal. "would buy" is 4.847cm and the conditional
  % still carries the counterfactual against (c)'s indicative "buys"; the
  % perfect tense survives in the caption. Origins are 7.75 and 12.40: (d)'s
  % title then ends at 17.247, which is 0.053cm inside the measure, exactly the
  % clearance panel (b)'s title already runs at. Columns are 4.65 and 4.90cm,
  % and the content footprints are 4.081/4.847 (titles) and 3.210/3.697
  % (entries), so neither column is the larger and more detailed one -- which is
  % what C4's equalisation was for. SUPERSEDES C4's "equal width (4.60cm)",
  % which was never true of the drawing: pass 12's origins were 7.75 and 12.70,
  % i.e. 4.95 and 4.60.
  \node[ttl] at (7.75,3.60) {(c) what the positive buys};
  \node[ttl] at (12.40,3.60) {(d) what a negative would buy};

  % THREE ENTRIES EACH, PAIRED ROW BY ROW. PASS 13, AND THIS IS THE CORRECTNESS
  % REPAIR IN THIS PANEL, NOT A TYPOGRAPHIC ONE.
  %  (i) THE ACROSS-READING WAS OFF BY ONE. Pass 12 set four baselines in each
  %      column and asserted, in the ledger and in the caption, that they were
  %      to be read across entry by entry. They could not be: (c) was two
  %      sentences word-wrapped over four lines and (d) a single four-item
  %      comma list, so row 1 paired "the model did not discard" with "the drug
  %      lost inside the model," and the negation of the layer sat one row
  %      above the layer it negated. Three of the four pairings were nonsense.
  %      The rows are now grouped, not merely stacked: 0.32cm of leading WITHIN
  %      a row and 0.52cm BETWEEN rows, so a wrapped entry cannot be mistaken
  %      for the next entry. Each row's first line is on one baseline in both
  %      columns, and the three pairs are exact.
  % (ii) "no class of failure," IS DELETED, because it contradicts this
  %      figure's own panel (b) and the claim box it is drawn from. 04b:142-144
  %      reads "It rules out that the drug field is lost before the decoder
  %      runs, an encoding failure in the activations", and panel (b) says so
  %      on the page, 3.5cm above, in the words "struck by the positive, and
  %      only this one". A positive that strikes one of two modes has bought a
  %      class of failure ruled out; what it has not bought is a depth and a
  %      fix. The two entries that survive, "no layer" and "no class of fix",
  %      are the two the sources license: 04b:110-112, "what the probe buys is
  %      the depth at which one would appear", and the caption's own lede, that
  %      the negative "would have named the layer at which the drug was lost
  %      and the class of fix that could reach it".
  % (iii)"an activation-side failure, and" IS DELETED FROM (d), because panel
  %      (b)'s struck line already reads "an activation-side encoding failure"
  %      2.5cm away. One object, one name, drawn once; the fix classes stay,
  %      because they are named nowhere else in the drawing.
  % (iv) (d)'s first entry carries the hedge the instrument requires. The probe
  %      is linear -- panel (a) now says so -- and a linear decode failure
  %      licenses only that the drug is not linearly readable, not that it is
  %      absent. The claim box is careful about exactly this ("Its identity is
  %      linearly readable from the residual stream at every layer measured");
  %      pass 12's "the drug lost inside the model," was not, and since the
  %      figure's whole argument is a comparison of licence SIZES, an
  %      overstated counterfactual inflates the asymmetry it exists to draw.
  %      SUPERSEDES the pass-9 note "THE COUNTERFACTUAL BLOCK'S FIRST ENTRY
  %      LOST ITS FINITE VERB" as to its wording only; the reason that note
  %      gives -- that the entry must be a noun phrase, not an assertion that
  %      the drug IS lost -- is preserved and strengthened by the hedge.
  % (c) is ink: it happened. (d) is quietink: it did not.
  % Leading is 0.35cm WITHIN a row -- the one \scriptsize leading in the figure,
  % shared with panel (b) -- and 0.52cm BETWEEN rows.
  % PASS 14, THREE ENTRIES BECOME TWO PAIRED ROWS OF ONE LINE EACH, AND ONE OF
  % THE DELETIONS IS A CORRECTNESS REPAIR, NOT A TYPOGRAPHIC ONE.
  %  (i) "no class of fix" IS DELETED, BY EXACTLY THE ARGUMENT PASS 13 USED TO
  %      DELETE "no class of failure," ONE ROW ABOVE IT -- and that argument
  %      was never applied to this entry, in the header or in the ledger.
  %      04b:31-36 is a TWO-MEMBER partition and each member names its own fix
  %      class: mode one is "an activation-side encoding failure that target-
  %      and input-side fixes might reach", mode two "a decoder failure only an
  %      objective-side fix could reach". Panel (b) strikes mode one and says
  %      so on this page, 3.5cm above, in the words "struck by the positive,
  %      and only this one". If one member of an exhaustive pair is struck the
  %      survivor stands, and the survivor names its fix class -- so within
  %      this figure's own frame the positive DOES buy a class of fix, and the
  %      entry is false. It survived pass 13 only because panel (b) draws mode
  %      two's gloss truncated to "a decoder failure", dropping the very clause
  %      of its source sentence that refutes the entry. A drawing may not
  %      contradict, at 8pt, the box it cites as its source.
  %      The alternative repair -- rewriting the entry to the class the
  %      positive does buy, "objective-side fixes only", against (d)'s
  %      "target- and input-side fixes" -- was built and rejected: that licence
  %      holds only CONDITIONALLY, on the model in fact ignoring the drug,
  %      which is 4.4's and 4.5's finding and not this section's, and a
  %      two-word slot cannot carry the condition without overstating it. Both
  %      fix classes go to the caption, which can state the condition.
  %      (d)'s "target- and input-side fixes" goes with it, because a row with
  %      one side is not a row and the entry is named nowhere else it is owed.
  % (ii) ROW ONE LOSES ITS WRAP AND ITS S18 COLLISION. "the model did not
  %      discard what it was handed" is seven words verbatim from 04b:107-108,
  %      the paragraph the float is attached to ("recovering it from the
  %      activations establishes only that the model did not discard what it
  %      was handed"), six lines of source above the float -- which S18 forbids
  %      outright. It was also the only entry in either column that wrapped, so
  %      (c) read as three items against (d)'s two and the across-reading broke
  %      at row one: measured, the white gap between rows 1 and 2 was 6.76pt in
  %      (c) and 16.68pt in (d), a 2.47x difference in the same two rows.
  %      "the drug field is not discarded" is 106.48pt = 3.742cm, one line,
  %      inside (c)'s 4.65cm column, keeps the sense the positive licenses, and
  %      shares no five-word clause with the prose or with any caption.
  % (iii)WITH EVERY ENTRY ON ONE LINE THERE IS NO WITHIN-ROW LEADING LEFT, so
  %      the two rows sit at one pitch of 0.60cm in both columns and the
  %      pairing needs no rule and no cue. The asymmetry the figure exists to
  %      draw is now also visible as ink for the first time: (c) sets 39
  %      characters against (d)'s 58, and (c)'s second entry is 0.991cm against
  %      (d)'s 3.233cm, so the small licence is the small column.
  %      SUPERSEDES pass 13's (ii)/(iii) rationale at this block as to the
  %      entries that survive, and C4's "FOUR ENTRIES EACH" in full.
  \node[lab]  at (7.75,2.85) {the drug field is not discarded};
  \node[qlab] at (12.40,2.85) {the drug not linearly readable};

  \node[lab]  at (7.75,2.15) {no layer};
  \node[qlab] at (12.40,2.15) {the layer at which it is lost};

% --- the figure's one disclosure line ---------------------------------------
% PASS 13: REWORDED. "Nothing in this figure is a measured quantity" sat 0.9cm
% under a line printing 480, twelve and forty, and then exempted the depth
% ticks in its own second clause, so it read as a blanket claim followed
% immediately by two counterexamples. The counts are the probe's design, not
% its results, and saying so is shorter than conceding an exception. It also
% now carries the count "seven", which pass 12 deleted with the caret label and
% which the pass-9 layout note (still standing) argues is the one fact the
% depth column exists to carry: without it a reader cannot tell whether the
% ticks are the depths probed or a sample of a sweep over all sixteen.
% One line, 389.66pt = 13.695cm, left-aligned to the figure's x origin.
% PASS 14: THE SECOND CLAUSE GOES, AND THE FIRST IS WHAT WAS DOING THE WORK.
% Pass 13 added "the counts are the probe's design" because the pass-12 wording
% ("Nothing in this figure is a measured quantity") sat under a line printing
% 480, twelve and forty and read as a blanket claim with two counterexamples
% beneath it. But pass 13 made TWO changes at once, and it was the other one --
% "a measured quantity" becoming "an estimate" -- that removed the
% contradiction, since a count is not an estimate. The clause was therefore
% belt and braces on a line already sound, at 35 characters, in the longest
% string in the figure, on a drawing over S15's ink ceiling. At 76 characters
% the line still answers the acceptance question ("can the reader say, without
% leaving the figure, that nothing here is a measured quantity?") on its own
% and still carries the count "seven" the pass-9 layout note asks for. The line
% is now 9.60cm rather than 13.695cm and is no longer the widest \scriptsize
% object on the page. SUPERSEDES pass 13's R9 foot-line clause as to the second
% sentence only; its rewording of "measured quantity" to "estimate" stands.
  \node[qlab] at (0.00,0.38)
    {Intervals: none in any panel; nothing here is an estimate, and the seven
     depth ticks are the only scaled marks.};

  % the bounding box IS the fullwidth measure, less 1.8pt
  \path (0,0.16) rectangle (17.30,7.72);
\end{tikzpicture}

% ---------------------------------------------------------------------------
% READY TO PASTE into Sections/04b-representation.tex, at \beatsection
% sec:methods-probe, immediately after the paragraph "The probe, and why its
% positive is nearly free" (04b-representation.tex:108-113) and before
% "Forward passes only." Not written into the Section file by this build.
%
% THE PASS-11 CAPTION IS RETAINED BELOW, COMMENTED OUT AND MARKED SUPERSEDED,
% because the thesis is audited against these headers. It is FALSE against the
% pass-12 and pass-13 drawings in four places: it indexes the panels "Left,"
% and "Right," where the drawing sets four lowercase panel letters, it says the
% three accuracies and the floor are "set here as text" (they are not set here
% at all any more), it spends two sentences denying a reading the drawing no
% longer offers, and it says the two licence blocks "are sized by their
% contents", which was the footprint asymmetry pass 12 replaced.
%
% The PASS-13 REPLACEMENT IS WRITTEN TO figs/_PENDING_CAPTIONS.tex, keyed by
% \label{fig:licensing}, together with every sentence pass 12 and pass 13 cut
% from the drawing. PASS 13 CREATED THAT FILE: pass 12 asserted twice that it
% had, and it did not exist. This builder does not edit Sections/. The block
% below is a DUPLICATE of the pending file, kept so this file is self-contained
% and the audit trail is here; the pending file is the one an integrator reads.
%
% ###########################################################################
% PASS 14 SUPERSEDES THE WHOLE OF THE PARAGRAPH ABOVE, AND THE PASS-13 CAPTION
% BELOW IS NOW HISTORY, NOT THE PROPOSAL.
%
% (i) The paragraph above is FALSE AS WRITTEN. figs/_PENDING_CAPTIONS.tex did
%     exist when pass 13 ran (it was created by the pass-10 repair of
%     fig:materiality), and pass 13's APPEND TO IT DID NOT SURVIVE: read on
%     2026-08-19 the file held fig:materiality, fig:comparators and fig:hook
%     blocks and no fig:licensing block at all. All three pass-14 refuters
%     opened with that finding and all three were right. The likeliest cause is
%     a concurrent read-modify-write from another figure's builder, which is
%     the hazard CF12 names about this shared file.
% (ii) PASS 14 APPENDED the block, in place, to figs/_PENDING_CAPTIONS.tex, and
%     READ THE FILE BACK TWICE -- once immediately after the append and once at
%     the end of the pass. Both times the block was there and the
%     fig:materiality, fig:comparators and fig:hook blocks were all still
%     there. It was appended at the END rather than inserted in figure order
%     deliberately: an append cannot destroy another builder's block and an
%     ordered insert can, which is what happened to pass 13.
%     NO LINE NUMBER IS RECORDED FOR IT, and that is the point of the two
%     reads: between them another builder appended to the fig:comparators
%     block and every heading below it moved -- this block went from line 484
%     to 597 and fig:hook from 306 to 369 -- so a line number written into
%     this header would be false within the hour. The \label heading is the
%     stable address.
% (iii)THE PASS-14 CAPTION IS THE ONE IN THE PENDING FILE, and it is NOT
%     duplicated here. Keeping a second copy in this header is what let pass 13
%     believe it had staged a caption it had not; one copy, in the file an
%     integrator is told to read, is the safer arrangement. What is recorded
%     here instead is what the pass-14 caption differs from the pass-13 draft
%     below in, so the audit trail is complete without a duplicate:
%       * it prints NONE of the four accuracies, where the pass-13 draft set
%         all four. It names them ("the raw peak and the deepest, the
%         context-removed deepest and the $1/12$ floor") and points at
%         \cref{tab:workspace} and \cref{fig:mechanism}(a). This is the
%         cross-figure redundancy resolution applied to the caption as well as
%         to the drawing: 0.821 and 0.883 are owned by 4.11(a), the only place
%         the two series are drawn AS series with their crossing at layer 12
%         visible, and printing them a fourth time here is the redraw the
%         resolution exists to stop. The values stay in the prose at
%         04b:159-166, in the claim box, and in \cref{tab:workspace};
%       * it carries the two FIX CLASSES in full, with the premise they depend
%         on, because pass 14 deleted "no class of fix" and "target- and
%         input-side fixes" from panels (c) and (d);
%       * it restates each failure mode's gloss in full, because the drawing
%         abbreviates both;
%       * it drops "The two licence blocks are sized by their contents", which
%         the pass-13 draft had already dropped, and also drops "the figure
%         quantifies nothing about the size of a licence": with the boxes gone
%         and (c)/(d) set as plain text columns, nothing on the page reads as
%         an area, and a caption denying a reading the drawing no longer offers
%         is the defect that sentence was meant to fix;
%       * it says "the single position read, the last of the prompt" rather
%         than "at the last prompt position", because the drawing now sets
%         "the last prompt position" as the caret's direct label and S18 bars
%         a shared five-word clause.
%     MEASURED: 2554 characters, 22 typeset lines under the real preamble in a
%     fullwidth float with the pass-14 drawing above it, against the live
%     pass-11 caption's 1692 characters and 15 lines, and against the 19-22 the
%     rest of the 4.7-4.12 arc runs at. It compiles with zero errors and zero
%     Overfull boxes on one page.
% ###########################################################################
%
% ---- PASS 13 REPLACEMENT CAPTION, SUPERSEDED BY PASS 14, kept for the audit
% ---- trail. DO NOT PASTE THIS ONE. -----------------------------------------
% \begin{figure}[htbp]
% \begin{fullwidth}
% \centering
% % ===========================================================================
% licensing.tex -- fig:licensing, what the probe's positive licenses and what
% the negative that never arrived would have licensed.
% FIGURE BODY ONLY. No preamble, no document environment, no float, no caption.
% \input it from inside a figure/fullwidth float in
% Sections/04b-representation.tex, at \beatsection sec:methods-probe, straight
% after the paragraph "The probe, and why its positive is nearly free"
% (04b-representation.tex:108-113). The complete float is staged, commented, at
% the foot of this file.
%
% Compiles against thesis_v2/preamble.tex and nothing else: no \includegraphics,
% no external data, no TikZ library beyond the ones preamble.tex already loads
% (this file uses arrows.meta for the two Stealth heads and
% decorations.pathreplacing for the mirrored brace under the prompt strip).
%
% WHAT IT IS. Two panels. LEFT: where the drug name enters the prompt, and
% where the probe reads. RIGHT: the two ways a model can ignore a drug it was
% told about, one of them struck by the positive and the other left standing,
% beside the small licence the strike buys and the large one a failure would
% have bought. The figure draws an ASYMMETRY, not a result: the section's own
% discount, that a positive here was near-certain in advance while a negative
% would have been decisive, has no graphical form anywhere in the document, and
% nothing in thesis_v2 draws the prompt, its fields, or the read position.
%
% EVERY NUMBER DRAWN, AND WHERE IT COMES FROM
% -------------------------------------------
% Nothing in this figure is a plotted coordinate. Four accuracies appear as
% TEXT and seven layer indices are the only quantity with a POSITION. Sibling
% ledger of every one of them at machine precision: thesis_v2/figs/licensing.json.
%
%   set as text, panel A                       source
%   0.821  probe at layer 9        RESULTS_cluster/workspace_probe.json ::
%                                  layers.9.drug_decode = 0.8208333333333334
%   0.758  probe at layer 16       ... :: layers.16.drug_decode
%                                       = 0.7583333333333334
%   0.083  shuffled-label floor    ... :: chance = 0.08333333333333333
%                                  (identical at layers.<L>.chance)
%   0.883  context removed, L16    ... :: layers.16.drug_decode_context_removed
%                                       = 0.8833333333333332
%   480    prompt count            COMPUTED, not stored: config.n_drugs 12 x
%                                  config.n_per_drug 40. The thesis prints
%                                  \num{480} at 04b-representation.tex:115-120.
%   12, 40 twelve drugs, forty     ... :: n_drug_classes = 12, config.n_drugs
%          cells each              = 12, config.n_per_drug = 40
%
%   [SUPERSEDED IN PART BY PASS 12, below: the four ACCURACIES 0.821, 0.758,
%   0.083 and 0.883 are no longer set in this figure at all. 480, 12 and 40
%   remain, and the block above is kept verbatim because the ledger's
%   removed_from_drawing entries are audited against it.]
%
%   DRAWN AS A POSITION, panel A
%   2, 4, 6, 8, 9, 12, 16          ... :: config.layers = "2,4,6,8,9,12,16".
%   The seven ticks on the depth column are at TRUE layer spacing,
%   y = 2.10 + (L/16)*4.80, so the even grid and the one added layer (9) are
%   visible as such. This is the figure's only scaled mark and it is a scale of
%   depth, not of any measured quantity. The seven y values are, and are
%   CHECKED against the map rather than against the eye:
%     L    2     4     6     8     9     12    16
%     y   2.70  3.30  3.90  4.50  4.80  5.70  6.90
%   Layer 12 was drawn at 6.00 in the first build and is a repaired defect; see
%   design decision 2.
%
%   NOT FROM ANY JSON -- design facts, cited by line, drawn_true false in the
%   ledger. This is the handling slab.json gives the appendix's +2.75 nats.
%   the five prompt fields, in this order:
%     cell line | drug | dose | mechanism | control cell sentence
%     Sections/03-Apparatus.tex:250-252, "The prompt supplies cell line, drug,
%     dose, mechanism and the control cell's sentence".
%     CHECKED AGAINST figs/hook.tex:335-339, which draws the same strip for
%     sec:res-used: the five labels are spelled identically and are in the same
%     order, and hook.tex's header records the same measured \scriptsize widths
%     (cell line 0.954, drug 0.616 bold, dose 0.579, mechanism 1.424, control
%     cell sentence 2.511; total ink 6.084cm), arrived at independently. Only
%     the cell edges differ, because hook's strip is 8.20cm and this one is
%     7.20cm and each distributes its own slack. The two figures are meant to
%     read as one object seen twice, and they do.
%   the read position, "the last prompt position ... where generation begins":
%     Sections/04b-representation.tex:115-118.
%   the two failure modes, every word of F1 and F2:
%     Sections/04b-representation.tex:33-36, the \begin{question}[3] box.
%
% AGREEMENT WITH THE TEXT (checked against the Sections file, not assumed):
%   04b:126-127  "reads $82\%$ at layer 9 and $76\%$ at layer 16, the deepest
%                measured, against the $\approx 8\%$ floor"
%                -> panel A prints 0.821, 0.758 and the 0.083 floor.        OK
%   04b:131-132  "once cell line, plate and dose are projected out, the series
%                rises with depth instead, to $0.883$ ($88\%$)"
%                -> panel A prints 0.883 with that condition attached.      OK
%   04b:118-120  "A cross-validated multinomial logistic-regression classifier
%                learns which of the twelve drugs the prompt named"         OK
%   04b:122-124  "Folds are stratified over prompts and not grouped by well"
%                -> drawn verbatim as the probe's stated limitation.        OK
%   04b:108-110  "The prompt names the drug in plain text ..., so recovering
%                it ... establishes only that the model did not discard what
%                it was handed" -> B12 is that sentence, split over two lines,
%                and it is the ONLY entry in the licence that exists.       OK
%   04b:110-112  "The informative outcome would have been a negative, and what
%                the probe buys is the depth at which one would appear"
%                -> B13's second line, "the layer at which it is lost".     OK
%   04b:33-36    "Either the drug never enters its computation, an
%                activation-side encoding failure that target- and input-side
%                fixes might reach, or it enters and nothing reads it, a
%                decoder failure only an objective-side fix could reach"
%                -> F1, F2 and B13's last two lines are that partition,
%                word for word.                                             OK
%   04b:142-144  claim box: "It rules out that the drug field is lost before
%                the decoder runs, an encoding failure in the activations. It
%                does not rule out that the decoder ignores it."
%                -> the strike lands on F1 and only on F1.                  OK
% tab:workspace (04b:228-234) prints the same 0.821, 0.758 and 0.883 in its
% Decode columns; this figure sets them as text, not as marks, so no scale
% here can be read against that table.
%
% The subtitle of panel B, "a negative, at the 0.083 floor, would have been
% decisive", names the floor a second time on purpose. Without it the figure
% never states what a failure would have LOOKED like, and the only cue left is
% the word "failed" in the right-hand tag; the acceptance question this figure
% is built to answer is precisely "what would the probe have had to return".
%
% DESIGN DECISIONS, each with the reason that forced it
% -----------------------------------------------------
%  1. THE FIGURE CARRIES NO AXIS AND NO PLOTTED VALUE. The four accuracies are
%     set as text. An accuracy axis was the obvious alternative and was
%     rejected twice over: fig:comparators row A already plots exactly these
%     numbers on exactly that rule, so a second one would be a redraw; and an
%     axis invites the reader to read the height of 0.821 as the size of the
%     finding, when the section's whole point is that the finding's size is
%     nearly irrelevant because a positive was near-certain in advance.
%  2. THE ONLY SCALED MARK IS DEPTH. The seven ticks sit at true layer spacing
%     rather than evenly, because "seven layers" and "the depth at which a
%     negative would have appeared" is what the probe actually buys (04b:111-112),
%     and an evenly spaced tick row would hide that the grid is even except for
%     the one layer added at 9. The 8/9 pair is 0.30cm apart, the document's
%     minimum label separation, which is why the column is 4.80cm tall; it may
%     not be shortened without either dropping a numeral or lying about the grid.
%     REPAIRED, PASS 9: layer 12 was drawn at y = 6.00 and the map gives 5.70.
%     Measured off the 200dpi render, 2->4->6->8->9 ran at 0.30cm a layer, then
%     9->12 at 0.40 a layer and 12->16 at 0.225, so on the page three layers
%     occupied more room than four and layer 12 sat where layer 13 belongs. The
%     column is the figure's only scaled mark and the section's payoff is the
%     depth a negative would have named, so a grid step of error here is the one
%     error the figure cannot carry. Checked safe before drawing: 5.70 is 0.90
%     above the layer-9 tick and 1.20 below layer 16, and the nearest ink at
%     that height is the val node on baseline 5.97, which ends near x = 5.5,
%     well clear of the tick column at x = 7.02.
%  2b.THE 8/9 NUMERALS ARE NOT BRACED, AND THAT IS A DECISION, NOT AN OVERSIGHT.
%     A reviewer measured 0.102cm of clear white between the two numerals
%     against a glyph height of 0.203cm and asked for comparators.tex's brace
%     idiom, "9, 8" named as an ordered set. Overruled. comparators braces marks
%     whose POSITIONS a reader cannot resolve on a measured axis, and fanning
%     them would claim a resolution the estimates do not have; here the two
%     positions are exactly right, each numeral is unambiguously attached to its
%     own 0.42cm tick, and the tightness IS the datum -- layers 8 and 9 are
%     adjacent in a sixteen-layer stack and must look it. Bracing them would
%     hide the one thing the added layer at 9 is there to show. The pair was
%     re-read at 200dpi after the numerals went quietink and the two glyphs are
%     separate and legible. The column cannot be lengthened either: its bottom
%     is pinned to the caret at y = 2.10 (which is layer 0, the prompt) and its
%     top to the "layer" label at 7.12, so the most 16 layers could ever have is
%     4.95cm, 0.009cm a layer more than they have now.
%  3. NO PANEL LETTERS. Two uppercase sans word tags in accent, the corpus
%     convention (circularity.tex, slab.tex); (a)/(b) appears nowhere in this
%     document's figures.
%  4. THE STRIKE IS THE ONLY MARK ABOVE 1.0pt. 1.1pt, the same weight as
%     circularity.tex:156, where one struck edge carries that figure's claim.
%     REPAIRED, PASS 9: it is now ONE HORIZONTAL stroke through the x-height of
%     line one, 8.70 to 13.50 at y = 5.60, and not the diagonal
%     (8.86,5.20)--(13.94,5.64) of the first build. Both reviewers reported the
%     diagonal as blocking and the render agrees: its left end sat exactly on
%     line two's baseline and its right end 0.11cm above line one's, so it
%     entered along the bottom of "an activation" and left through the x-height
%     of "computation", and at print size "an activation" was a smear. Its right
%     end was also 0.26cm from the box border with the accent arrow to "struck"
%     starting at 14.20 on nearly the same line, so the two read as one
%     continuous pointer out of the box. A rule through the x-height leaves the
%     words readable, because the eye completes struck text and cannot complete
%     text cut along its baseline. LINE ONE ONLY: line one is the proposition
%     the positive eliminates and line two is its gloss, so one stroke does the
%     work and the 1.1pt register stays "used exactly once". Extent measured off
%     the render: line one's ink is 8.76 to 13.44, the stroke overhangs 0.06cm
%     each side, and it stops 0.70cm short of the box's right border, which is
%     what kills the leader reading.
%     Everything else here is 0.30-0.90pt, so weight alone says
%     which mark is the argument. The struck box is drawn and STRUCK, not
%     deleted: a deleted box leaves the caption nothing to point at, and the
%     partition is two-membered whether or not one member survives. The two
%     boxes carry IDENTICAL borders for the same reason -- the struck one was
%     drawn at rulegrey!55 in the first build, and if the box is also fainter
%     then two marks say the same thing and neither says it alone.
%  5. THE TWO LICENCE BLOCKS ARE SIZED BY THEIR CONTENTS AND BY NOTHING ELSE.
%     The one the positive buys holds one sentence; the one a failure would
%     have bought holds four lines. That difference in footprint IS the
%     asymmetry, readable before a word is read. It quantifies nothing, and
%     both the caption and the ledger say so.
%  6. SOLID-AND-WHITE MEANS "HAPPENED", DASHED-AND-CREAM MEANS "DID NOT".
%     One register, used once each, so the counterfactual block cannot be
%     mistaken for a second result. REPAIRED, PASS 9: the cream is now panelbg
%     straight, the same cream slab.tex and hook.tex use, and so is the depth
%     column's. Three creams were in play across this slate -- panelbg in hook
%     and slab, panelbg!60 in the column, panelbg!45 in the counterfactual box
%     -- and the last two are about 2% ink, which will not reliably reproduce on
%     uncoated stock; a tint that does not print is a tint doing no work while
%     adding a third cream to the system. The DASHED border is what carries
%     "this did not happen"; the fill only says "a filled region".
%  7. THE LICENCE THAT EXISTS HAS NO SECOND ENTRY AND NO "which fixes it points
%     at" LINE. The answer to that question is nothing, and an absent line says
%     it better than the word "none", which would read as a measured category.
%     Same principle as comparators.tex's bare dot for a quantity with no
%     interval: the absence is the mark.
%  8. THE DRUG CELL IS THE ONLY ACCENT MARK IN THE PROMPT STRIP, at accent!12
%     fill with an accent border and an accent bold label, so the eye finds it
%     before it reads anything. It is the reason every number in the left-hand
%     block is what it is.
%  9. NO INTERVAL IS DRAWN ANYWHERE. The probe returns a cross-validated
%     accuracy and workspace_probe.json stores no interval for it; the document
%     draws bars only where the instrument returns its convention
%     (comparators.tex, style_v2.py). Nothing here is a bar, a whisker or a
%     dot on a scale, so nothing invites one.
% 10. THE ACCENT BUDGET IS THE DRUG CELL, THE CARET, THE SEVEN DEPTH TICKS,
%     THE STRIKE AND ITS ARROW AND LABEL, AND THE FOUR WORD TAGS.
%     REPAIRED, PASS 9: the seven depth NUMERALS were accent and are now
%     quietink. With them the figure carried about 24 accent marks against a
%     corpus ceiling near 16 (comparators.tex gives accent exactly two jobs),
%     and roughly fourteen of them stood inside one 0.4 x 4.8cm column, making
%     it the densest accent region on the page -- so the eye landed on the tick
%     column first and then hunted beside tick 9 for "0.821", which is set as
%     prose three centimetres to its left. A numeral beside a tick is an axis
%     label, which is what quietink is for in every other figure in the corpus
%     (comparators.tex:70-72, gate.tex, purify.tex). The TICKS stay accent
%     because the depth grid is what \cref{sec:methods-probe} says the probe
%     actually buys. Re-rasterised at 115dpi after the change -- about print
%     size at arm's length -- and the first fixation is now the drug cell.
% 11. NO 0.50 CHANCE RULE APPEARS. This instrument's floor is 1/12 = 0.083, not
%     0.50, and it is set as text rather than as a rule because there is no
%     axis for a rule to cross. (Recorded because the corpus is split on how a
%     chance rule is drawn: style_v2.chance_line() mandates thin solid black,
%     while ladder.tex:112-115, the only TikZ precedent, draws it densely
%     dotted. This figure needs neither.)
%
% LAYOUT NOTES, all forced by a build or by a measurement
% -------------------------------------------------------
%   * fullwidth = textwidth 142mm + marginparwidth 28mm + marginparsep 4mm
%     = 174mm = 493.23pt. (preamble.tex's own comment beside
%     \newenvironment{fullwidth} says 171mm and is stale arithmetic.) The box
%     is pinned with a bare \path to 17.30cm = 491.4pt, 1.8pt inside the
%     measure. Nothing is drawn outside x in [0, 17.30] except the 0.2pt
%     half-width of B13's right border.
%   * EVERY MULTI-LINE BLOCK IS SEPARATE anchor=base west NODES. No node in
%     this file uses align= or text width. frames.tex measured a 3.7pt drift
%     from that: TikZ's align builds a minipage whose first-line height follows
%     the AMBIENT \baselineskip, so a node's ink moves with whatever prose
%     precedes the float. Every y below IS a baseline. Within-block leading is
%     0.33cm = 9.39pt, one \scriptsize line; between blocks it is 0.45cm.
%   * THE PROMPT STRIP'S CELL EDGES ARE SET FROM MEASURED STRING WIDTHS, not
%     chosen. Every string in this figure was set in the real preamble and its
%     \wd read back before a coordinate was written. At \scriptsize the five
%     field names measure, in cm:
%       cell line 0.954 | drug (bold) 0.616 | dose 0.579 | mechanism 1.424 |
%       control cell sentence 2.511.  Total ink 6.084cm.
%     The strip therefore cannot be the 6.78cm first laid out: at that width
%     "mechanism" alone overruns its 1.28cm cell by 0.144cm and the five cells
%     share 0.696cm of padding, 0.07cm a side. The strip is 7.20cm, the widest
%     that keeps the depth column's numerals inside panel A, which leaves
%     0.09-0.11cm each side -- more than the 1.6pt inner sep slab.tex sets on
%     its own nodes. Cell edges: 0.00 / 1.13 / 1.93 / 2.69 / 4.29 / 7.20.
%   * THE ELLIPSIS INSIDE THE FIFTH CELL WAS DROPPED, AND THE MEASUREMENT IS
%     WHY. A typeset $\cdots$ at \scriptsize is 0.405cm and needs about 0.16cm
%     of gap; the fifth cell is 2.91cm and its label is 2.511cm, leaving
%     0.399cm in total. No strip width recovers it: pushing the caret past
%     x = 7.30 drives the layer numerals to within 0.13cm of the panel divider.
%     A drawn three-dot ellipsis (0.18cm) leaves a 0.066cm gap to the label,
%     which is 1.9pt and reads as part of the last word.
%   * THE FIFTH CELL'S RIGHT EDGE IS SOLID. It was drawn dotted first, on the
%     reading that the strip has no horizontal scale. Two faults, both seen in
%     the render: the caret covers the top 0.21cm of that edge, leaving a
%     0.41cm dotted stub that reads as a printing fault; and a dotted edge at
%     exactly the x the caret marks says both "the prompt ends here" and "this
%     field runs on past here". The caret is the load-bearing mark, so the
%     absence of scale is stated in the caption, in words, instead.
%   * THE CARET STRADDLES THE STRIP'S RIGHT EDGE ON PURPOSE, and is centred on
%     the depth column above it. The last prompt position is where the prompt
%     ends, so the mark for it sits ON that boundary rather than inside the
%     last cell. PASS 9: its apex is now the column's own bottom edge, y = 2.10,
%     and the 0.5pt accent stub that used to bridge 2.04 to 2.10 is deleted.
%     Inside a 3mm square at that corner there converged the strip's top rule,
%     the strip's right rule, the caret, the stub, the column's bottom edge and
%     the brace's right arm, and the stub crossed a rulegrey rule uninterrupted.
%     Six strokes became five and the apex does the joining. The brace's right
%     arm was left where it is: measured, it terminates at (7.20,1.30), which is
%     0.12cm below the strip's bottom-right corner and 0.56cm below the caret's
%     own base, so it is not part of that corner.
%   * THE PROMPT CELLS ARE DRAWN AS figs/hook.tex DRAWS THEM. PASS 9: rulegrey
%     0.35pt, panelbg fill, rounded corners=1pt for the four ordinary fields,
%     and accent 0.4pt / accent!12 / rounded corners=1pt for the drug field.
%     Before this pass licensing drew square-cornered WHITE cells sharing edges
%     and hook drew rounded panelbg boxes with rounded gutters, so the same
%     five-field prompt was drawn two incompatible ways in two figures a page
%     apart while both captions insisted they were the same prompt. hook's is
%     the better mark -- the rounding reads as token boundaries rather than as
%     a table -- so licensing adopts it. Only the cell EDGES still differ, and
%     they must: this strip is 7.20cm and hook's is 8.20cm, and each distributes
%     its own slack. That is a difference of measure, not of treatment.
%   * NO LEADER RUNS FROM THE CARET LABEL TO THE CARET. One was drawn and
%     taken out: the only clear path from the label's lower right corner passes
%     within 0.02cm of the depth column's bottom-left corner at (7.02,2.10), or
%     else crosses the strip's top border at x = 6.94 and runs on inside the
%     last cell. Both read as a stray stroke. The label's second line ends
%     0.21cm from the caret's footprint with nothing between, inside the
%     document's ~0.3cm leader threshold.
%   * NO ACCENT MARK LIES LEFT OF x = 7.02. The seven depth ticks first ran
%     from 6.92, within 0.49cm of the text block, and the block's line at
%     y = 3.30 shares its baseline exactly with the layer-4 tick, so
%     "0.883 once cell line, plate and dose are projected out ---- 4" read as
%     one labelled row. The ticks now start at the column's own left edge.
%   * PANEL A'S LINE LENGTH IS CAPPED AT 6.85cm, not the 6.35cm of the first
%     layout: that limit followed the depth column, which moved right from
%     6.60 to 7.02 when the strip was widened. The longest line drawn is
%     "where generation begins; read at seven depths" right-aligned on 6.84,
%     and the longest left-aligned line is "a cross-validated
%     multinomial probe, twelve classes" at 6.57cm.
%   * "READ AT SEVEN DEPTHS", NOT "READ AT EVERY LAYER". Beside a column that
%     is sixteen layers tall and carries seven ticks, "every layer" left a
%     reader unable to tell whether the ticks were the seven layers actually
%     probed or an arbitrary sample of a sweep over all sixteen -- which is the
%     one fact the column exists to carry. The review asked for "read at each of
%     the seven marked depths"; measured, that is 1.94cm longer and the node is
%     anchor=east on 6.84, so it would have started at x = -1.1 and taken the
%     picture off the measure. "seven depths" is 20 characters against "every
%     layer"'s 19 and cost nothing.
%   * THE COUNTERFACTUAL BLOCK'S FIRST ENTRY LOST ITS FINITE VERB, and the
%     review's own wording would not fit. "the drug is lost inside the model,"
%     read as an assertion that the drug IS lost, the opposite of the section's
%     conclusion; the review asked for "that the drug is lost inside the
%     model,", which measures 4.66cm from x = 12.84 and put the picture 11.88pt
%     over the measure -- the harness reported exactly that Overfull hbox on the
%     build that tried it. "the drug lost inside the model," is 4.02cm, fits
%     with 0.44cm to spare, and is the noun phrase the other three entries are.
%   * THE CARET LABEL IS SHORT LINE OVER LONG LINE, not the reverse. Written
%     long over short it sat 0.75cm under a left-aligned quietink line of the
%     same width and register, and at print size the two read as one three-line
%     grey paragraph. Short first, right-aligned, breaks that; the long second
%     line then also fills the 3.8cm of empty measure that sat above the strip.
%   * B14 IS FIGURE-WIDE, NOT PANEL-LOCAL. It measures 11.70cm at \footnotesize
%     and is centred on 8.65, the figure's own centre, below the panel divider
%     (which stops at y = 0.90). Panel-local it would have run from 7.10 to
%     18.80 and off the measure.
%   * TYPE. Ambient \small on the picture; every node sets its own size.
%     \scriptsize = 8pt on all of it except the two panel subtitles and the
%     closing line, which are \footnotesize = 10pt, and the four word tags,
%     which are \sffamily\scriptsize\bfseries. Nothing is set at body size and
%     nothing is scaled: 1cm here is 1cm on the page, no scale= key, no
%     \resizebox.
%
% ===========================================================================
% PASS 12 -- 2026-08-18. THE STYLE-CONTRACT REBUILD.
% ===========================================================================
% Everything above this line is the v1-v11 record and is left intact because
% the thesis is audited against it. This section states what pass 12 changed,
% and SUPERSEDES BY NAME every claim above that pass 12 makes false. Nothing
% above has been edited away.
%
% WHY. The supervisor's verdict on the 4.8-4.12 slate: "go back over them and
% iron them out even more ... make it look more professional, right now it
% looks a bit sloppy and childish ... make sure that the figures are CORRECT
% and display the idea we want to convey in the best way possible." The
% reference standard is fig:mechanism (4.11, printed page 51): lowercase
% parenthesised panel letters with sentence-case descriptive titles, rotated
% axis labels that NAME the measured quantity, direct labels, no legend, no
% all-caps banners, no italic aphorisms, no boxed callouts, and the whole
% argument carried by a twenty-line caption. Measured against it this figure
% was, before pass 12: 4 all-caps sans accent banners out of 4 panel heads,
% 17% of its characters italic, 11.21% ink at 200dpi and 9.9 characters per
% 100mm^2 -- DENSER THAN THE APPENDIX REFERENCE PLATE A.3 (9.7%, 6.4) while
% sitting in the running argument of chapter 4.
%
% WHAT PASS 12 DELETED FROM THE DRAWING, AND WHY
% ----------------------------------------------
%  D1. THE FOUR ACCURACIES, set as running text in panel A: "$0.821$ at layer
%      9 and $0.758$ at the deepest,", "against a $0.083$ shuffled-label
%      floor", "$0.883$ once cell line, plate and dose are projected out", and
%      the "what it returns" key that headed them. They sat immediately left of
%      a scaled layer stack carrying ticks labelled 9 and 16, so a reader tries
%      to read 0.821 off the stack at layer 9 and cannot -- and the caption had
%      to spend two whole sentences denying the reading the drawing offered
%      ("They are set here as text, not as marks on a rule; the accuracy rule
%      is Figure 4.7, row one"). A caption that must deny a reading its own
%      drawing invites is a defect in the drawing, not in the caption.
%      All four are drawn or printed in three other places -- fig:comparators
%      row (a), fig:mechanism panel (a), tab:workspace -- and this figure's own
%      caption prints them as well. Under the cross-figure redundancy
%      resolution, fig:mechanism (a) OWNS 0.821 and 0.883 because it is the
%      only place the two series are drawn AS SERIES and the crossing at layer
%      12 is visible; 4.7 row (a) keeps the chance floor because a comparator
%      row without its comparator is not a comparator row. 4.8 keeps the
%      ARGUMENT about the floor -- what a negative would have looked like --
%      and loses the numerals. All four survive in licensing.json as
%      removed_from_drawing entries with their key paths and with pointers to
%      the figures that still draw them.
%      SUPERSEDES: the "set as text, panel A" block above (kept verbatim, since
%      the ledger's removed_from_drawing entries are audited against it); the
%      AGREEMENT-WITH-THE-TEXT lines for 04b:126-127 and 04b:131-132 (the
%      agreement is now the CAPTION's to carry, not the drawing's); design
%      decision 1's second clause; and the paragraph beginning "The subtitle of
%      panel B ... names the floor a second time on purpose", which is void
%      because that subtitle is deleted (D3).
%  D2. THE FOUR ALL-CAPS SANS ACCENT BANNERS -- "WHERE THE DRUG NAME ENTERS",
%      "WHAT THE OUTCOME LICENSES", "WHAT THE POSITIVE BUYS", "IF IT HAD FAILED
%      INSTEAD". Every one of them narrated rather than named. A.3 (figs/slab.tex)
%      does use short accent caps banners, but it is an appendix reference
%      plate a reader stops at and studies, its banners are noun phrases naming
%      objects, and the caps sans accent face is byte-identical to the
%      preamble's \marginlabel apparatus face -- which belongs to the margin,
%      not to a figure in the running body. 4.11 and 4.13 carry zero banners.
%      Replaced by 4.11's form: a lowercase parenthesised letter and a
%      sentence-case descriptive title, \footnotesize roman ink, left-aligned
%      to the panel's own x origin, no colon, no bold, no caps.
%      SUPERSEDES design decision 3 ("NO PANEL LETTERS ... (a)/(b) appears
%      nowhere in this document's figures") IN FULL. It was true when written
%      and is now false: fig:mechanism sets "(a) encoded more, used no more"
%      and "(b) ablation cost against a random null", and it is the reference
%      standard the slate is being brought to.
%  D3. THE THREE ITALIC SUBTITLES -- "the drug is named in the prompt itself",
%      "a positive was near-certain in advance;", "a negative, at the $0.083$
%      floor, would have been decisive" -- and D4's epigram. The subtitles were
%      \footnotesize italic quietink, i.e. the LARGEST type in the figure was a
%      grey italic editorial line while every measured thing was set 25%
%      smaller. The last of the three is also the caption's own lede restated.
%  D4. THE CENTRED ITALIC EPIGRAM "one run, two possible outcomes; the one that
%      could not fail is the one that happened." It was 9.96pt Pagella-Italic
%      quietink centred within 0.2pt of the figure's own x-centre, and the
%      naive-reader pass reports the eye going to it BEFORE any mark, reading
%      it as a second caption above the real one. In this document the argument
%      lives in the caption. It is written to figs/_PENDING_CAPTIONS.tex.
%  D5. THE FOUR BOX STROKES -- the two failure-mode boxes, the solid white
%      "what the positive buys" block and the dashed panelbg counterfactual
%      block. A stroked rectangle is drawn only where the rectangle IS an
%      object in the argument (a token, a prompt field, a subspace, a scoring
%      frame). These four enclosed SENTENCES, and the figure's own caption
%      conceded that "the figure quantifies nothing about the size of a
%      licence". Panels (c) and (d) are now two columns of equal width on one
%      x-grid with baseline-aligned entries, in the frames.tex idiom.
%      SUPERSEDES design decision 6 IN FULL (there is no dashed border and no
%      cream fill left to carry "this did not happen"; quietink carries it, as
%      it does everywhere else in the slate), and design decision 5's second
%      sentence -- see C4 for what replaced the footprint asymmetry.
%  D6. THE GREY FILL INSIDE THE LAYER STACK, and the stack itself as a filled
%      column. It read as a BAR: the naive-reader pass's first question was
%      what quantity the length of the grey bar encoded, and the answer is
%      none. It also ran from y = 2.10 to 7.05, i.e. past the layer-2 tick
%      below and the layer-16 tick above, implying that depth continues beyond
%      the measured range. It is now a plain 0.40pt ink axis rule running from
%      the layer-2 tick to the layer-16 tick and stopping exactly on them.
%  D7. "folds stratified over prompts, not grouped by well". Prose already
%      carried at 04b:122-124; it is a limitation of the instrument, which is
%      caption and prose material, not a name for anything the figure draws.
%      SUPERSEDES the AGREEMENT line for 04b:122-124 above.
%
% WHAT PASS 12 CHANGED, AND WHY
% -----------------------------
%  C1. THE AXIS NOW DEFINES ITSELF. This is the supervisor's first complaint
%      and it is answered mechanically. The depth stack is the figure's only
%      axis and it carried, before this pass, seven bare numerals and the word
%      "layer" set as a 0.30cm-tall stub above the column top -- no quantity in
%      axis-label position, no scope. It now carries 4.11's form exactly:
%      rotated 90 degrees reading upward, immediately left of the tick
%      numerals, line 1 \footnotesize roman ink "layer", line 2 \scriptsize
%      quietink "residual stream read at the last prompt position, where
%      generation begins". Line 2 is 258.63pt = 9.09cm set on one line, which
%      does not fit beside a 4.20cm axis, so it is wrapped onto two typeset
%      lines at the only break that keeps both under the 4.60cm of clear
%      vertical measure available at that x: "residual stream read at the last
%      prompt" (137.06pt = 4.817cm) over "position, where generation begins"
%      (119.57pt = 4.202cm). It is ONE label line wrapped, not two label lines.
%      Rotated multi-line text stacks toward increasing x, so line 1 ("layer")
%      sits FARTHEST from the axis, at x = 5.44, exactly as a wrapped
%      matplotlib ylabel does in 4.11.
%  C2. THE PANEL-A PROSE BLOCK IS NOW ONE BLOCK OF TWO LINES, and it is the
%      figure's only prose. "what is read: the residual stream at the last
%      prompt position" measures 264.20pt = 9.29cm at \footnotesize and cannot
%      be set on one line inside a 6.90cm panel; more to the point, C1 puts
%      that exact fact in axis-label position, where the style contract wants
%      it, so keeping it in the prose block as well would be a duplication.
%      The block is therefore "what reads it: a cross-validated" over
%      "multinomial probe, twelve classes" (4.851 and 5.241cm), \footnotesize
%      roman ink, left edge on panel (a)'s x origin.
%  C3. THE COLOUR INVERSION IS FIXED, and it was a correctness defect, not a
%      preference. Before this pass the accent arrow marked the STRUCK branch,
%      so accent meant ELIMINATED here and meant THE HEADLINE RESULT in the
%      figure four pages earlier; the naive-reader pass read the accent arrow
%      as the surviving branch. Across the slate accent now means one thing:
%      the drug-side object the rest of the chapter measures. In this figure
%      that is exactly two objects -- the drug cell in the prompt strip, and
%      the SURVIVING failure mode with its arrow and its outcome label. The
%      struck mode, its arrow, its outcome label and the whole of panel (d),
%      which is a counterfactual, are quietink. Panel titles, the axis, the
%      axis labels, the prose, the caret and panel (c) are ink; ticks, brace,
%      dividers and the four ordinary prompt cells are rulegrey furniture.
%      SUPERSEDES design decision 10's accent budget IN FULL: the seven depth
%      ticks are no longer accent (they are furniture, 0.30pt rulegrey), the
%      caret is no longer accent (it is part of the axis apparatus, ink), the
%      strike is no longer accent (it is the black threshold-weight stroke, see
%      C6), and the four word tags no longer exist (D2). It also supersedes the
%      layout note "NO ACCENT MARK LIES LEFT OF x = 7.02", which is now
%      trivially true for a different reason: the only accent in panel A is the
%      drug cell, at x 1.12 to 1.90.
%  C4. PANELS (c) AND (d) ARE NOW EQUAL. Two columns of equal width (4.60cm)
%      on one x-grid, entries baseline-aligned, FOUR ENTRIES EACH. Before this
%      pass the counterfactual block was the larger and more detailed of the
%      pair -- one sentence against four lines, inside a bigger dashed box --
%      so the eye landed on the thing that did not happen. The asymmetry the
%      figure exists to draw is now carried by the WORDS rather than by the
%      footprint: panel (c)'s entries 3 and 4 name, in the same slots and the
%      same register as panel (d)'s entries 2 to 4, exactly what the positive
%      did NOT buy -- "no layer, no class of failure," / "and no class of fix."
%      against "the layer at which it is lost," / "an activation-side failure,
%      and" / "target- and input-side fixes." A reader can now diff the two
%      columns line by line, which is a stronger reading of the asymmetry than
%      a difference of box size ever was.
%      SUPERSEDES design decision 7 ("an absent line says it better than the
%      word 'none'") and design decision 5's second sentence. The reason
%      decision 7 gave still stands against the bare word "none", which would
%      read as a measured category; what is drawn is not "none" but a named
%      absence in a slot the facing column fills, which is the opposite of a
%      category and cannot be read as a quantity.
%  C5. ONE X ORIGIN PER PANEL, declared as a named coordinate and referenced by
%      every element of that panel. Before this pass each banner sat 4.6pt
%      (1.6mm) LEFT of the text it headed, four times over, and the central
%      divider had 8.0pt of clearance on its left and 15.2pt on its right.
%      Panel (a) origin x = 0.00; panels (b), (c) origin x = 7.75; panel (d)
%      origin x = 12.70. The divider sits at x = 7.42, which is 0.33cm from
%      panel (a)'s rightmost ink (the caret's right base at 7.09) and 0.33cm
%      from panel (b)'s origin: the gutter is now equal on both sides to
%      0.001cm.
%  C6. FIVE STROKE WEIGHTS BECOME THREE, EACH WITH ONE MEANING.
%        0.30pt rulegrey -- furniture: the seven ticks, the brace, the panel
%                           divider, the (c)/(d) row rule, the four ordinary
%                           prompt cells, and the two outcome leaders (which
%                           carry their mark's colour, not rulegrey).
%        0.40pt          -- the axis rule (ink) and the drug cell's border
%                           (accent). Used twice, so it is not a singleton.
%        0.90pt ink      -- the strike, and nothing else.
%      The 1.1pt register of design decision 4 is retired: 0.90pt is the
%      document's threshold weight (style_v2.chance_line() is hard-coded
%      lw=0.9, and figs/family.pdf draws the pre-registered materiality margin
%      at 0.90pt), and a figure that contains no threshold at all can carry its
%      one mark of argument at that weight without any collision. THAT IS THE
%      JUSTIFICATION the contract requires for a weight used in one kind of
%      place: this figure has no chance rule, no gate, no margin and no floor
%      (design decision 11, which still holds), so nothing else here may be
%      black solid 0.90pt and nothing is.
%      SUPERSEDES design decision 4's first sentence and its "used exactly
%      once" clause: the strike is now TWO strokes, one through each of the two
%      lines of the struck mode, because the mode is a two-line entry and a
%      rule through its first line only left its second line standing. It also
%      supersedes the cell-border weight recorded in the pass-9 prompt-cell
%      note: 0.35pt is not one of the document's five weights and the ordinary
%      cells are now 0.30pt. The cells' construction is otherwise unchanged and
%      still matches hook.tex's \tokbox.
%  C7. THE FIGURE FOOT CARRIES ONE \scriptsize QUIETINK DISCLOSURE LINE,
%      left-aligned to the figure's x origin: "Nothing in this figure is a
%      measured quantity; the depth ticks are the only scaled marks." Every
%      figure in the slate carries exactly one such line. Here it does the job
%      an interval declaration does elsewhere -- it tells the reader, without
%      leaving the figure, that no mark on the page is an estimate. Design
%      decision 9 (no interval is drawn anywhere) still holds and is now
%      readable on the drawing rather than only in the ledger.
%  C8. THE STRIP IS 6.90cm, NOT 7.20cm, and the caret is 0.40cm tall rather
%      than 0.24cm. Both follow from C1: the rotated axis-label block needs a
%      clear vertical measure of 4.82cm at x = 5.76, which puts its lower edge
%      at y = 2.391, so the strip's top edge came down from 2.46 to 2.28 and
%      the caret grew to span 2.28 to 2.70 -- 0.111cm of clearance, measured,
%      between the rotated label's foot and the strip's top rule. The strip's
%      slack is redistributed over the same five measured field widths:
%      6.084cm of ink and 0.816cm of padding, 0.0816cm a side, giving cell
%      edges 0.00 / 1.12 / 1.90 / 2.64 / 4.23 / 6.90 against the old
%      0.00 / 1.13 / 1.93 / 2.69 / 4.29 / 7.20. NO FIELD CHANGED, no order
%      changed, and the strip still carries no horizontal scale.
%      SUPERSEDES the layout note "THE PROMPT STRIP'S CELL EDGES", the
%      "PANEL A'S LINE LENGTH IS CAPPED AT 6.85cm" note (the cap is now 6.90cm,
%      set by the strip, and the longest left-aligned line drawn is 5.241cm),
%      the caret note's "apex is now the column's own bottom edge, y = 2.10"
%      (the apex is the axis rule's own base at the layer-2 tick, y = 2.70),
%      and the "B14 IS FIGURE-WIDE" note (B14 is deleted, D4).
%      The ELLIPSIS note and the FIFTH CELL'S RIGHT EDGE note both still hold
%      and were re-checked at the new width: the fifth cell is 2.67cm against a
%      2.511cm label, so there is even less room for a $\cdots$ than there was.
%  C9. WHAT DID NOT CHANGE, and must not, because it is the argument and the
%      measurement: the seven depth ticks and their map y = 2.10 + (L/16)*4.80,
%      giving 2.70 / 3.30 / 3.90 / 4.50 / 4.80 / 5.70 / 6.90 -- BYTE-IDENTICAL
%      to pass 11, including the pass-9 repair of layer 12 from 6.00 to 5.70;
%      the five prompt fields, their order, their measured widths and the
%      accent on the drug cell; the caret at the strip's right edge and the
%      read-here relation between the two; the brace and the count line
%      "$480$ real tier-2 prompts: twelve drugs, forty cells each"; the strike
%      on the eliminated mode and the two-member partition; and the fact that
%      the figure draws no interval, no bar and no marker on a rule.
%      Design decisions 2, 2b, 8 and 11 stand unchanged. NO NUMBER IN THIS
%      FIGURE CHANGED VALUE IN PASS 12; four numbers left the drawing entirely
%      and are recorded in licensing.json as removed_from_drawing.
%
% GEOMETRY OF PASS 12, all measured with \wd in the real preamble before a
% coordinate was written (figs/_preview.py licensing, 200dpi, read at each step)
% -------------------------------------------------------------------------
%   panel (a) origin x = 0.00 ; panels (b) and (c) origin x = 7.75 ;
%   panel (d) origin x = 12.70 ; divider x = 7.42 ; measure 0 to 17.30cm.
%   axis rule x = 6.90, y 2.70 to 6.90 ; ticks 6.62 to 6.90 ; numerals
%   anchor=east on 6.52 ; rotated "layer" at x = 5.44, wrapped scope note at
%   x = 5.76 and 6.04, all centred on the axis midspan y = 4.80.
%   Measured string widths used (cm, \footnotesize unless marked s = \scriptsize):
%     (a) title 4.969 | (b) title 9.495 | (c) title 4.081 |
%     (d) title, wrapped 4.157 over 1.921
%     prose 4.851 / 5.241
%     mode 1 5.842, gloss s 4.292 | mode 2 4.472, gloss s 2.093
%     outcomes s 2.727 / 2.147 and 2.921 / 1.712
%     (c) entries s 3.210 / 2.516 / 3.352 / 2.293
%     (d) entries s 3.804 / 3.303 / 3.689 / 3.399
%     count line s 6.557 | foot line s 10.860
%     cell labels s: 0.954 / 0.616 bold / 0.579 / 1.424 / 2.511
%   Widest object on the page is panel (b)'s title, 9.495cm, which ends at
%   x = 17.245, 0.055cm inside the measure. No text crosses the measure.
%
% C10. THE EMPTY BAND IN PANEL (a) IS STRUCTURAL AND IS LEFT EMPTY ON PURPOSE.
%   Measured on the 200dpi render it is about 5.2 x 4.0cm, at x 0.0-5.2 and
%   y 2.4-6.4. It is forced, not careless, and four ways of removing it were
%   built and rejected:
%     (i)  shortening the axis. The 8/9 pair is 0.30cm apart at the current
%          scale, which is the document's minimum label separation, so the
%          16-layer span cannot be under 4.80cm and the drawn 2-to-16 span
%          cannot be under 4.20cm without dropping a numeral or falsifying the
%          grid. Design decision 2 and 2b, both still standing.
%     (ii) narrowing the strip. Its five field names measure 6.084cm of ink and
%          hook.tex sets the same five, so the strip cannot go below about
%          6.7cm without either abbreviating a field or breaking the agreement
%          with fig:hook that both captions assert.
%     (iii)inverting the axis so the strip sits at the top and depth runs
%          downward. This moves the band, it does not remove it, and it would
%          draw the residual stream running the wrong way.
%     (iv) moving panels (c) and (d) into the band. They need 2 x 4.60cm plus a
%          gutter, which is 9.55cm; the band is 5.2cm wide.
%   What is left is the space between the prompt and the read, which is the
%   residual stream, and this figure deliberately draws nothing there: the
%   stream's contents are fig:mechanism's subject and the slab's geometry is
%   fig:slab's. Filling it with narration is exactly what pass 12 removed.
% ===========================================================================
% PASS 13 -- 2026-08-18. THE REFUTATION PASS.
% ===========================================================================
% Three adversarial refuters read the pass-12 render. Everything above is left
% intact; this block states what pass 13 changed and SUPERSEDES BY NAME every
% pass-12 claim it makes false. Each change is recorded at the code it changes
% as well, so a reader of the drawing does not have to come back here.
%
% R1. THE COORDINATE RECORD IN "GEOMETRY OF PASS 12" AND IN C5 WAS WRONG IN
%     SIX PLACES, AND THE HEADERS ARE WHAT THE THESIS IS AUDITED AGAINST.
%     Superseded, by name, with the values the code actually drew in pass 12:
%       divider x         header 7.42   -> DRAWN 7.40   (and now deleted)
%       caret right base  header 7.09   -> DRAWN 7.05   (and now 7.02)
%       divider gutters   header 0.33   -> DRAWN 0.35 both sides
%       tick inner end    header 6.62   -> DRAWN 6.74   (and now 6.78)
%       numerals anchor   header 6.52   -> DRAWN 6.64   (and now 6.70)
%       rotated "layer"   header 5.44   -> DRAWN 5.36   (and now 6.18)
%       scope line two    header 6.04   -> DRAWN 6.12   (and now deleted)
%     The same block recorded the two failure-mode gloss lines at their
%     \scriptsize widths ("gloss s 4.292", "gloss s 2.093") while pass 12's
%     code set both at \footnotesize (5.365 and 2.616cm) -- and it is the gloss
%     width that sets the strike's right end, so the header could not be used
%     to check the one mark the file calls its argument. Pass 13 sets the modes
%     at \scriptsize, so 4.292 and 2.093 are now the drawn widths and the
%     record and the code agree for the first time.
%     C4's "Two columns of equal width (4.60cm)" was also never true of the
%     drawing: pass 12's origins were 7.75 and 12.70 on a 17.30 measure, i.e.
%     4.95 and 4.60. See the (c)/(d) block for what pass 13 drew instead.
%
% R2. figs/_PENDING_CAPTIONS.tex DID NOT EXIST. D4 and the note above the
%     staged float both asserted that the cut epigram and the pass-12
%     replacement caption were "written to figs/_PENDING_CAPTIONS.tex"; the
%     file was never created, so every sentence cut under D1, D3, D4 and D7 was
%     unaccounted for outside a comment block in this file, and the caption
%     that actually prints in the thesis is still the pass-11 one, which is
%     false against the drawing in four places (it indexes the panels "Left,"
%     and "Right,", it prints four accuracies as "set here as text", and it
%     says the two licence blocks "are sized by their contents"). Pass 13
%     creates figs/_PENDING_CAPTIONS.tex and writes the pass-13 caption and
%     every deleted sentence into it, keyed by \label{fig:licensing}. This
%     builder still does not edit Sections/. The copy below is retained as the
%     audit trail and is marked a duplicate of the pending file.
%
% R3. THE STRIKES WERE BLACK SOLID 0.90pt, WHICH IS THE SLATE'S THRESHOLD
%     STROKE. Corrected to quietink 0.40pt. Full reasoning at the code.
%     SUPERSEDES C6's third weight; the figure now has two weights, 0.30 and
%     0.40pt, neither a singleton.
% R4. THE CARET'S APEX WAS THE AXIS BASE AND THE LAYER-2 TICK AT ONCE.
%     Broken apart; apex 2.50, base 6.78-7.02. SUPERSEDES C8's apex clause.
% R5. THE AXIS LABEL WAS THREE ROTATED LINES AND OVERRAN THE AXIS AT BOTH ENDS,
%     with the quantity name outermost. One \footnotesize "layer" beside the
%     numerals, one \scriptsize scope line outboard, 3.908cm inside a 4.20cm
%     axis. SUPERSEDES C1's wrap arithmetic and its "line 1 is the outermost".
% R6. PANELS (c) AND (d) COULD NOT BE READ ACROSS, and two of (c)'s four
%     entries were unsupported. Three paired rows; "no class of failure,"
%     deleted as contradicting this figure's own panel (b) and the claim box;
%     "an activation-side failure," deleted as a duplicate of panel (b)'s own
%     struck line; (d)'s first entry hedged to what a LINEAR probe's negative
%     licenses. SUPERSEDES C4 in full and the ledger's
%     entries_3_and_4_are_new_in_pass_12.
% R7. THE VERTICAL DIVIDER AND THE DEPTH AXIS WERE PARALLEL HAIRLINES 0.50cm
%     APART. The divider is deleted.
% R8. PANEL (b)'s MODES WERE \footnotesize, i.e. the same size as the panel
%     titles, and their outcome labels were \scriptsize on a different leading,
%     so nothing lined up across the gutter. Both are \scriptsize now, on one
%     leading, on one set of baselines. One leading per size across the figure:
%     0.40cm at \footnotesize, 0.35cm at \scriptsize.
% R9. SMALLER ITEMS: the mirrored brace's pointer aimed at empty white and is
%     now a plain span rule; "$480$" set three digits in URWPalladioL-Roma
%     against Pagella everywhere else and is now text; the prose block's
%     pronoun had no antecedent on the reading line and now names the residual
%     stream; the foot line conceded two counterexamples in its own second
%     clause and now states what the counts are; the word "seven", deleted with
%     the caret label in pass 12, is back on the drawing in that foot line,
%     which is what the still-standing pass-9 layout note asks for.
%
% WHAT PASS 13 OVERRULED, AND WHY
% -------------------------------
%  O1. "REBUILD PANEL (a): DROP THE DEPTH SWEEP, OR SET LAYER HORIZONTALLY, OR
%      SHORTEN THE AXIS BY TIGHTENING THE 8/9 PAIR." Overruled on all three
%      counts. The seven-tick depth scale at true spacing is the figure's only
%      scaled mark and its keep list names it; the 8/9 pair is already at true
%      spacing and 0.30cm apart, the document's minimum label separation, so
%      the span cannot shorten without dropping a numeral or falsifying the
%      grid (design decisions 2 and 2b, both standing); a horizontal layer axis
%      would draw the residual stream running the wrong way and would break the
%      caret's read-here relation to the strip. C10's four rejected removals
%      stand. What pass 13 did take from that finding is the head alignment:
%      panels (a) and (b) now put their first content line on one baseline.
%  O2. "RETITLE (a) 'where the drug name enters, and where the probe reads'."
%      Overruled on measurement. Panel (b)'s title is 9.495cm of the 9.55cm the
%      right half has and sits on the SAME baseline, so panel (a)'s title
%      cannot exceed about 7.0cm; the proposed title measures 8.5cm and the
%      shortest honest variant 7.7cm. The read is named at the axis instead,
%      which is where the style contract wants it.
%  O3. "DRAW THE CARET AS AN OPEN rulegrey CHEVRON." Overruled. rulegrey is
%      furniture and the caret is content -- it is the one mark of the read
%      position, and the keep list names the triangle. Breaking the coincidence
%      and shrinking it to its pass-9 size answers the defect; recolouring it
%      would demote a content mark to furniture.
%  O4. "PUT PANEL (b)'s OUTCOME COLUMN ON PANEL (d)'s ORIGIN (12.40)."
%      Overruled: the longer strike ends at 12.484, so the label column would
%      begin inside the strike. It sits at 13.44, 0.50cm clear of that
%      terminal, which is the separation the pass-9 note asks for.
%  O5. "KEEP THE STRIKES AT 0.90pt IN quietink." Overruled in favour of 0.40pt.
%      0.90pt is the slate's threshold weight whatever its colour, and at 0.90pt
%      against \scriptsize the struck line does not survive at print size --
%      which is the legibility failure the pass-12 ledger already records.
%  O6. "REACH S15's 6.0% INK AND 5.0 CHARACTERS PER 100mm^2." Not reached, and
%      the arithmetic is recorded rather than hidden: pass 13 measures 6.65%
%      and 5.42, against pass 12's 7.95% and 6.05 and pass 11's 11.21% and 9.9.
%      What is left is 697 characters of which 483 are mandated furniture --
%      four panel titles, the axis definition, the definitional prose, the five
%      field names, the count line and the one interval-style disclosure -- and
%      214 are the partition and the two licences, which are the figure. The
%      only cuts left would delete the axis definition that answers the
%      supervisor's first complaint, or one of the two failure modes.
% ===========================================================================
% PASS 14 -- 2026-08-19. THE SECOND REFUTATION PASS.
% ===========================================================================
% Three adversarial refuters read the pass-13 render and returned thirty
% findings. Everything above is left intact; this block states what pass 14
% changed, what it OVERRULED and why, and SUPERSEDES BY NAME every pass-12 and
% pass-13 claim it makes false. Each change is also recorded at the code it
% changes. NO NUMBER IN THIS FIGURE CHANGED VALUE IN PASS 14: the seven depth
% ticks are still y = 2.10 + (L/16)*4.80 -> 2.70/3.30/3.90/4.50/4.80/5.70/6.90,
% byte-identical to passes 11, 12 and 13, and the count line still prints the
% same 480, twelve and forty. The strip translated 0.14cm DOWN in y as one
% rigid body; not one cell edge in x moved, and x is the only axis on which
% that strip carries anything at all.
%
% S1. THE PENDING CAPTION FILE. All three refuters opened with the same
%     blocking finding, and it is correct: figs/_PENDING_CAPTIONS.tex contains
%     no \label{fig:licensing} block. R2 above asserts, in as many words, that
%     "Pass 13 creates figs/_PENDING_CAPTIONS.tex and writes the pass-13
%     caption and every deleted sentence into it"; the file exists and holds
%     fig:materiality, fig:comparators and fig:hook blocks, and this figure's
%     append was lost -- almost certainly clobbered by a concurrent
%     read-modify-write from another builder, which is the hazard CF12 names.
%     So pass 13 committed, one pass later, the identical failure it accused
%     pass 12 of. SUPERSEDES R2's last three sentences and the "duplicate of
%     the pending file" note above the staged float: they described an intent,
%     not the filesystem.
%     Pass 14 APPENDS the block (pure append, never a rewrite, so it cannot
%     clobber another builder's), then READS THE FILE BACK to confirm the block
%     survived, and records the byte offset it landed at in licensing.json. The
%     block carries the pass-14 caption, the four clauses of the live pass-11
%     caption at 04b:117-140 that are false against this drawing with their
%     line numbers, and every sentence passes 12, 13 and 14 cut from the
%     drawing, each against the caption sentence that now carries it.
%
% S2. "no class of fix" IS DELETED FROM PANEL (c), AND THIS IS THE ONE
%     CORRECTNESS DEFECT IN THE PASS-13 DRAWING. Full argument at the code.
%     In one line: 04b:31-36 is a two-member partition in which each member
%     names its own fix class, panel (b) strikes one member and says so on this
%     page, so the survivor and its fix class stand -- and pass 13 had already
%     deleted the sibling entry "no class of failure," for exactly this reason
%     without applying it here. (d)'s "target- and input-side fixes" goes with
%     it, because a row with one side is not a row. Both fix classes move to
%     the caption, which can state the condition a two-word slot cannot.
%     SUPERSEDES R6's list of surviving entries and C4 in full.
%
% S3. THE TWO OUTCOME LEADERS ARE DELETED. Measured: the quietink one left a
%     0.24cm gap at its left end and a 0.10cm gap at its right, touching
%     neither mark nor label, and lay midway between two quietink strikes so
%     the eye read the verdict "struck by the positive," as itself struck; the
%     accent one began 1.397cm PAST the mark it claimed to join. The pass-13
%     ledger recorded the opposite as a measured property, having measured only
%     the first of the two. Colour and exactly shared baselines carry the
%     pairing, which is 4.11's discipline. The outcome column moves 13.44 ->
%     13.00. SUPERSEDES C6's leader clause, O4's arithmetic, and pass 13's
%     lead/.style rationale; lead/.style itself is deleted.
%
% S4. THE CARET IS NAMED AT ITSELF and the three cues that made it an axis
%     arrowhead are broken: it gains a \scriptsize ink direct label "the last
%     prompt position" 0.08cm off its left base, the strip's 0.14cm drop puts
%     0.34cm between its apex and the axis base, and its base widens to 0.28cm
%     so it no longer shares an endpoint with the tick lane. Under the figure's
%     own map it now occupies layers 0.13 to 0.87, below the drawn scale.
%     SUPERSEDES R4 and licensing.json's read_position note.
%
% S5. THE ROTATED SCOPE LINE IS DELETED and "layer" moves off the layer-9 tick.
%     The scope line shared the five-word clause "at the last prompt position"
%     with the caption that actually prints, which S18 forbids, and it named at
%     90 degrees and 1.14cm away the one glyph that had no name of its own.
%     "layer" was centred on the axis midspan, which is the layer-9 position by
%     arithmetic (2.10 + (9/16)*4.80 = 4.80 = (2.70+6.90)/2), so it read
%     "layer 9"; 5.70-to-6.90 is the only tick gap on this axis long enough to
%     hold it clear. SUPERSEDES R5 and C1.
%
% S6. THE SPAN RULE UNDER THE STRIP IS DELETED. It duplicated the strip's own
%     bottom border at the same colour, weight and length 0.12cm away, and it
%     bracketed one schematic prompt under a line counting 480 of them.
%     SUPERSEDES R9's span-rule clause.
%
% S7. PANEL (b)'s TITLE IS A NAME, NOT A SENTENCE: 9.495cm -> 5.821cm, so the
%     widest object on the page is no longer a title that restated 04b:31-32.
%     Panel (c)'s row one loses its wrap and its own S18 collision ("the model
%     did not discard what it was handed" is seven words verbatim from
%     04b:107-108, six lines above the float) and becomes "the drug field is
%     not discarded", one line. The strikes start on the panel origin 7.75
%     rather than 0.06cm left of it. The (c)/(d) group drops 0.30cm and its row
%     pitch opens to 0.70cm, so deleting a row did not trade panel (a)'s hole
%     for a new one in the bottom-right corner. The foot line drops its second
%     clause. SUPERSEDES O2's measurement clause and R8's 0.52cm row gap.
%
% S8. THE PANEL-(a) BAND. The pass-13 drawing put one contiguous 56.2 x 39.2mm
%     void -- 17.2% of the figure -- under the prose block and left of the
%     axis, because pass 13 had moved that block from the band's middle up to
%     the panel head to buy a shared first-content baseline with panel (b).
%     Pass 14 moves it back to 5.00/4.60, centred on the axis midspan, and puts
%     the caret's new label in the band's lower part; the band is now two voids
%     of about 1.9 and 1.6cm and the block sits level with the axis it
%     describes. The two TITLE baselines still align at 7.45, which is the
%     alignment a reader sees. NOTHING FURTHER WAS PUT IN THE BAND, and the
%     reason is the contract, not indifference: S4 allows at most one prose
%     block per panel and panel (a) already has it, and C10's four ways of
%     REMOVING the band (shorten the axis, narrow the strip, invert the axis,
%     move (c)/(d) into it) were each rebuilt and each still fails. The band is
%     the space between the prompt and the read; this figure draws nothing
%     there because the stream's contents are 4.11's subject and the slab's
%     geometry is A.3's. SUPERSEDES pass 13's placement repair (ii).
%
% S9. MEASURED AFTER ALL OF THE ABOVE (200dpi, clipped to the content bbox):
%       content bbox   173.10 x 74.29mm   (S14 wants 170-174; cap 75mm)
%       ink            5.826%             (was 6.65 at pass 13, 11.21 at 11)
%       characters     596 non-space, 4.63 per 100mm^2  (was 5.44 at pass 13)
%                      694 with spaces,  5.40 per 100mm^2  (was 6.41)
%       stroke weights 0.30pt rulegrey furniture (4 ordinary cells, 7 ticks,
%                      the (c)/(d) row rule) and 0.40pt (axis rule ink, drug
%                      cell border accent, two strikes quietink). Two weights,
%                      neither a singleton, no 0.90pt stroke, no arrowhead,
%                      no leader, no dotted rule, no legend.
%       type           two sizes: \footnotesize 10.0pt (four panel titles, the
%                      prose block, the axis quantity name) and \scriptsize
%                      8.0pt (everything else). 0% italic. 0 all-caps banners.
%     S15's ink ceiling is met for the first time in this figure's history.
%     The with-spaces character count is still 5.40 against a 5.0 ceiling and
%     is recorded rather than hidden: the refuters' own measurement of that
%     rule was taken on non-space characters (5.44 against 5.0), which is met
%     at 4.63; closing the with-spaces gap needs 51 more characters and the
%     only strings left are the four titles, the axis name, the five field
%     names, the count line the spec's keep list mandates, the two failure
%     modes and their outcomes, the two licence rows and the one disclosure
%     line. Every one of them is mandated content.
%
% WHAT PASS 14 OVERRULED, AND WHY
% -------------------------------
%  P1. "SET THE SEVEN TICK NUMERALS IN INK AND LENGTHEN THE TICKS TO 0.20cm,
%      BECAUSE THE AXIS READS AS A PANEL DIVIDER." Overruled on both counts.
%      Quietink tick numerals are this corpus's convention, not an oversight:
%      comparators.tex:70-72, gate.tex and purify.tex all set them so, and
%      design decision 10 changed them from accent to quietink in pass 9 with
%      that reason recorded. S8 assigns ink to axis LABELS, which is where
%      "layer" is, and does not assign a colour to tick numerals; making this
%      one figure's numerals black would break it away from the four figures it
%      is being aligned to. The ticks cannot lengthen inward either: they run
%      6.78-6.90 and the numerals are right-anchored on 6.70, so 0.20cm ticks
%      would touch them. The divider reading the finding describes was created
%      by the divider itself and died with it in pass 13 (R7); what is left is
%      a rule with seven numbered ticks and a rotated quantity name, which is
%      an axis.
%  P2. "OPEN A GUTTER BETWEEN THE PROMPT CELLS, OR DROP rounded corners=1pt,
%      BECAUSE THE ROUNDINGS LEAVE WHITE NOTCHES AT EVERY INTERNAL JUNCTION."
%      The notches are real -- re-cropped at 1200dpi and seen -- and the fix is
%      still overruled here. The finding's own evidence is that figs/hook.tex
%      draws the identical construction at the identical radius, and both
%      captions assert the two strips are the same prompt seen twice; a strip
%      repaired in one file and not the other would break an agreement the text
%      states, to remove a 0.1mm artefact. This builder may not edit hook.tex.
%      Recorded in licensing.json as a coordinated repair the two files owe
%      together, so the next pass that touches both can make it once.
%  P3. "SET THE PROSE BLOCK AS ONE LINE AND DROP IT A RANK, BECAUSE IT IS THE
%      SAME SIZE AS THE PANEL TITLES." Overruled, and the finding concedes the
%      contract point itself: S4 mandates \footnotesize for definitional prose
%      and states in terms that prose OUTRANKS labels, which is A.3's ranking
%      and the one the contract chose to keep. A one-line rewrite either drops
%      "the residual stream", which is the antecedent R9 put there and the only
%      naming of the object the axis stands for, or drops "cross-validated",
%      which is why the accuracy is a presence test at all.
%  P4. "RESTORE THE PROBE'S TASK TO THE DRAWING: 'over the twelve drug
%      classes'." Overruled as to the drawing, accepted as to the record. The
%      figure is at its ink ceiling, the count line already prints "twelve
%      drugs" 2.4cm below, and naming the instrument does not require naming
%      its class count. The twelve-class task AND the fold caveat (04b:155-157,
%      cut in pass 12 as D7) are both in the pass-14 caption, and
%      licensing.json records under explicitly_not_drawn that the drawing names
%      the instrument without its target and why.
%  P5. "LEAVE PANEL (c)'s SLOTS BLANK, OR SET AN EN-DASH IN THEM, SO THE SMALL
%      LICENCE LOOKS SMALL." Overruled. An empty slot under a head reading
%      "what the positive buys" is ambiguous between "nothing" and "not drawn
%      yet", and design decision 7's own supersession at C4 gives the reason a
%      named absence beats a blank or the bare word "none". The finding's
%      underlying complaint -- that the two licences were drawn the same size,
%      with the smaller one in darker ink and more lines -- is answered by S2
%      and S7 above: (c) now sets 39 characters against (d)'s 58 and its second
%      entry is 0.991cm against (d)'s 3.233cm.
%  P6. "DRAW THE TIE BETWEEN THE DEPTH AXIS AND (d)'s 'the layer at which it is
%      lost', SO THE AXIS HAS A STATED ROLE." Overruled as a drawn mark. A
%      leader across 5.5cm of figure would be the only connector on the page
%      and would assert an argument, and a line saying "a negative would have
%      named one of these seven depths" is argument, not definition, which S4
%      and S18 put in the caption. The pass-14 caption says it in those words.
%      The axis is now fully DEFINED on the drawing -- quantity name, seven
%      numbered ticks, and the position it is read at, labelled at the caret --
%      which is what S5 asks of it.
%  P7. "MAKE THE VERTICAL RULE THE RESIDUAL STREAM, SO THE PROSE BLOCK HAS A
%      DRAWN REFERENT." Overruled. The rule is a LAYER axis; the residual
%      stream at one position is a 2048-dimensional vector at each of seven
%      depths, and renaming a depth coordinate after the object sampled along
%      it would put two names on one rule and cost S5 its clean reading. The
%      prose block was instead moved level with the axis and the caret named,
%      so the reading relation the block defines -- position, depths, probe --
%      is drawn beside it rather than four centimetres above it.
% ===========================================================================
% inner sep=0pt on every node, so that anchor=base west puts the START OF THE
% INK exactly on the panel's x origin and a node's width is exactly the \wd
% measured for its string. With TikZ's default 0.3333em the picture ran 5.43pt
% past the fullwidth measure (the harness reported exactly that Overfull hbox),
% because panel (b)'s title is 9.495cm of the 9.55cm the right half has.
% PASS 14: that title is now 5.821cm, so inner sep=0pt is no longer what keeps
% the picture inside the measure -- but it is kept, because it is also what
% makes anchor=base west put the START OF THE INK on the panel origin, which is
% what S13's one-x-origin check measures. The widest object is now panel (d)'s
% title, ending at x = 17.247, 0.053cm inside the 17.30cm measure.
%
% ===========================================================================
% v-BLOCK, 2026-08-19: WHOLE-SLATE PASS (fig:comparators, fig:licensing,
% fig:purify, fig:hook, fig:materiality reconciled against one another and
% against fig:mechanism, which is the document's reference standard and was
% NOT rebuilt). Nothing above this line is deleted; where an earlier claim is
% now false it is superseded EXPLICITLY below.
%
% WHY. The five figures were rebuilt independently and drifted. This pass owns
% only the differences BETWEEN them. No value changed in any of the five; the
% sibling .json ledgers carry the value-to-position maps, old and new.
%
% THE SLATE RULES THAT NOW HOLD ACROSS ALL FIVE (each stated once here, in
% every file, so any future single-figure edit can see what it would break):
%
%  R1 TICK NUMERALS are quietink, \scriptsize, inner sep 0, anchored north on
%     the tick and set 0.16cm below the axis rule, over a 0.10cm rulegrey
%     0.30pt tick. Before: comparators/licensing/purify quietink, hook and
%     materiality ink; ticks 0.10 / 0.12 / 0.14cm; three different gaps.
%  R2 A THRESHOLD'S OWN LABEL IS BLACK, matching the black solid 0.90pt rule
%     it names (S6). Before: comparators black, purify's "gate, 0.8" and
%     materiality's two margin labels ink.
%  R3 ONE NAME PER OBJECT (S19). The pre-registered 10% rule is "10%
%     materiality margin" in fig:comparators row (c), fig:materiality (a) and
%     (b), and in figs/family.py, which is where the name comes from.
%     fig:comparators' "pre-registered margin" is superseded; "pre-registered"
%     is now the caption's word, not the drawing's. Zero, where it is the null
%     a bar is read against, is "zero: no identity effect" in BOTH
%     fig:comparators row (c) and fig:materiality (b) -- one quantity, one
%     wording. V_pure is spelt with \textrm everywhere, as figs/slab.tex does.
%  R4 COLOUR KEY, one key for five figures (S8). accent = the drug-side object
%     under test. quietink = every comparator, null, control, replication or
%     discounted arm, AND its label, AND its leader. ink/black = thresholds,
%     axis rules, axis labels, panel heads, definitional prose. rulegrey =
%     furniture only. Within quietink the GLYPH separates two kinds: an OPEN
%     disc is a construction built to carry no signal (a null, a control, a
%     raw/before series); a FILLED disc is a real measurement that is not the
%     headline. fig:hook (d) drew both its null and its cell-line control in
%     ink and is corrected.
%  R5 PROJECTION AND ZERO RULES are quietink 0.50pt densely dotted (S7).
%     fig:hook's two panel-(c) projection lines were rulegrey 0.30pt densely
%     dashed and are corrected.
%  R6 DEFINITIONAL PROSE is \footnotesize ink and OUTRANKS every label (S4);
%     at most one block per figure, on the panel's own x origin. fig:hook's
%     one prose line was \scriptsize and is raised. fig:purify's panel-(b)
%     block carried an argument clause ("so causal claims use the purified
%     slab, not the raw") that its caption already carries, and is now
%     definitional only.
%  R7 ONE INTERVAL DECLARATION per figure, \scriptsize quietink, at the foot,
%     on the figure's x origin, opening with the word "Intervals:" and naming
%     EVERY panel including the ones with no measured mark (S17).
%     fig:comparators already did this and is the model; the other four now
%     match its form.
%  R8 THE FIVE-FIELD PROMPT STRIP is ONE object drawn in two figures, and the
%     two drawings are now congruent: cell boundaries 0 / 1.12 / 1.90 / 2.64 /
%     4.23 / 6.90 cm and cell height 0.42cm in both fig:licensing (a) and
%     fig:hook (a) and (b), labels centred in their cells, the drug cell
%     accent-bordered at 0.40pt on accent!12 and its label NOT bold. Before:
%     licensing 6.90cm wide with 0.62cm cells and a bold label, hook 7.44cm
%     wide with 0.42cm cells and a plain label -- the same five fields at two
%     sizes, four pages apart, with each caption pointing at the other.
%  R9 MEASURE. Every fullwidth figure draws 172.0-174.0mm; the one textwidth
%     figure draws 140.3mm. Ink at 200dpi inside the content bbox, counted as
%     greyscale luminance below 230 (the contract's own method), is now
%     5.63-5.99% on all five, under the 6.0% body ceiling, where before this
%     pass it ran 5.36-6.54%.
%
% WHAT THIS PASS DELIBERATELY DID NOT UNIFY, so that a later reader does not
% think it was missed:
%  * PANEL-HEAD CASE. fig:comparators sets "(a) Is it there?" and the other
%    four set "(a) where the drug name enters". comparators' heads are
%    interrogative SENTENCES and sentence case is correct for a sentence; the
%    v4 header argued this and the argument stands. fig:mechanism's lowercase
%    heads are noun phrases, not questions. One capital letter, four times.
%  * PANEL-HEAD POSITION. fig:comparators puts its heads in a right-aligned
%    left stub, the other four above the panel at its x origin. Moving them
%    above the rules costs about 20mm of height against a 118mm cap that the
%    figure already meets at 117.60mm. Not available.
%  * TWO-COLUMN GUTTER. purify and hook split at x = 8.40cm, licensing at
%    7.75cm. licensing cannot move right: "(d) what a negative would buy" is
%    4.85cm wide and its origin is already at 12.40 of a 17.30cm measure, with
%    1.5pt of slack. Moving purify and hook LEFT to 7.75 would be arbitrary.
% ===========================================================================
%
% CHANGED IN THIS FILE. (i) The prompt strip adopts fig:hook's cell height,
% 0.42cm, and fig:hook drops to this file's cell WIDTHS, so the two drawings
% of the one object are congruent (R8); the strip's top stays at y = 2.14, so
% the caret, the layer axis and everything above them are untouched, and the
% strip's bottom rises from 1.52 to 1.72. Labels move from anchor=base at
% y = 1.775 to centred at the cell's own mid-height, 1.92, with inner sep 0,
% which is \tokbox's placement. (ii) The drug cell's label loses \bfseries: it
% was the only bold glyph on the slate, and the accent colour plus the accent
% border already carry the emphasis. (iii) The layer ticks go from 0.12 to
% 0.10cm and the numerals from 6.70 to 6.74, a 0.06cm gap (R1). (iv) The
% "480 real tier-2 prompts" line goes ink -> quietink, matching the same class
% of line in fig:purify, and rises to y = 1.38 with the strip. (v) The prose
% block moves from y 5.00/4.60 to 5.55/5.15 so its first line sits on panel
% (b)'s second-mode baseline, 5.55, giving the two columns a shared
% horizontal. (vi) The foot line takes the slate's declaration form (R7).
% NOT FIXED, and it is this figure's real weakness: panel (a) is an L -- a
% 4.20cm vertical layer axis at the column's right edge and a 6.90cm strip
% along its bottom -- so about 6.0 x 3.4cm of the column is empty. The pass-12
% header records four rejected fills and this pass adds no fifth; the void is
% the (position x depth) plane in which this figure measures nothing, and the
% prose block sits in it to break it rather than to disguise it.
% RE-MEASURED: 173.10 x 74.30mm; ink 5.975%; 726 characters, 5.65 per 100mm^2,
% which is over S15's 5.0 and is the one contract clause this figure still
% misses -- see pass 12's arithmetic, unchanged by this pass. Exactly two type
% sizes, 9.96 and 7.97, in one face, no bold, no italic, no sans, no CM.
% One page, zero Overfull, zero Underfull.
% ===========================================================================
\begin{tikzpicture}[x=1cm, y=1cm, font=\small, inner sep=0pt]
  \tikzset{
    % 4.11's panel head: lowercase parenthesised letter, sentence-case
    % descriptive title, \footnotesize roman ink, left on the panel's origin.
    ttl/.style ={font=\footnotesize, text=ink, anchor=base west},
    % the figure's ONE prose register: definitional only, \footnotesize ink.
    def/.style ={font=\footnotesize, text=ink, anchor=base west},
    % the surviving branch -- the object the rest of the chapter measures.
    % PASS 13: \scriptsize, not \footnotesize. These are the two branches of a
    % partition, i.e. direct labels of drawn objects, not definitional prose;
    % under S1 the largest size in the figure must belong to a panel title, an
    % axis label or definitional prose, and 4.11 sets its in-plot labels one
    % step below its titles. They now share a size and a baseline grid with
    % their own outcome labels, which is what fixed the cross-gutter splay.
    liv/.style ={font=\scriptsize, text=accent, anchor=base west},
    % the struck branch, and every counterfactual.
    ded/.style ={font=\scriptsize, text=quietink, anchor=base west},
    % \scriptsize registers: labels, entries, scope and foot lines.
    lab/.style ={font=\scriptsize, text=ink, anchor=base west},
    qlab/.style={font=\scriptsize, text=quietink, anchor=base west},
    alab/.style={font=\scriptsize, text=accent, anchor=base west},
    tick/.style={font=\scriptsize, text=quietink, anchor=east},
    % The prompt cell is drawn as figs/hook.tex draws it (\tokbox): panelbg
    % fill, rounded corners=1pt. PASS 12: the border is 0.30pt, the document's
    % furniture weight, not the 0.35pt of pass 9, which is not one of the five
    % weights the slate uses.
    cell/.style={draw=rulegrey, line width=0.30pt, fill=panelbg,
                 rounded corners=1pt},
    % The two outcome connectors were IDENTICAL plain leaders. PASS 13 deleted
    % the Stealth heads: a single arrowhead is reserved for an operation acting
    % in a direction and its label must be a verb of action, and "struck by the
    % positive" and "survives" are outcomes, not operations.
    % PASS 14 DELETES lead/.style AND BOTH LEADERS. Reasoning at the two mode
    % blocks and in the pass-14 header at S3. The figure now contains no
    % arrowhead and no connector of any kind, and its only glyph that could be
    % mistaken for either -- the caret -- carries a direct label at itself.
    }

% ---------------------------------------------------------------------------
% ONE X ORIGIN PER PANEL. Every element below references its panel's origin,
% so the spread of text-line x0 inside a panel is exactly zero.
% ---------------------------------------------------------------------------
  \coordinate (OA) at (0.00,0);    % panel (a)
  \coordinate (OB) at (7.75,0);    % panels (b) and (c)
  \coordinate (OD) at (12.40,0);   % panel (d)
  \coordinate (OBL) at (13.44,0);  % panel (b)'s outcome column

% PASS 13: THE VERTICAL PANEL DIVIDER IS DELETED. It stood at x = 7.40 and the
% depth axis stands at x = 6.90, so the two were parallel hairlines 0.50cm
% apart with an empty channel between them; measured on the render, the band
% between them carried 0.16% ink over the full height. At print size they read
% as one doubled boundary, and panel (a)'s only scaled object was thereby read
% as frame furniture. The lettered panel heads and a 0.85cm gutter partition
% the figure without it, which is 4.11's own arrangement -- it draws no
% dividers at all. SUPERSEDES C5's divider clause and its gutter arithmetic.

% ###########################################################################
% # PANEL (a) -- where the drug name enters
% ###########################################################################
  \node[ttl] at (0.00,7.45) {(a) where the drug name enters};

  % --- the figure's only prose block: definitional, two lines -------------
  % PASS 13, TWO REPAIRS.
  %  (i) THE PRONOUN HAD NO ANTECEDENT ON THE READING LINE. Pass 12 read "what
  %      reads it: ...", and the only occurrence of "residual stream" was
  %      inside the rotated \scriptsize scope line, 4cm away and at 90 degrees.
  %      The block now names the object it defines. "twelve classes" leaves the
  %      block: the line "a cross-validated linear probe, twelve classes"
  %      measures 6.68cm and the panel's prose is capped at 5.48cm by the
  %      rotated label column at x = 5.62 (see the axis block below). "linear"
  %      is the word worth the space, because it is what makes panel (d)'s
  %      "the drug not linearly readable" the right counterfactual, and twelve
  %      still reaches the reader through the count line under the strip.
  % (ii) IT MOVED FROM THE PANEL'S VERTICAL MIDDLE TO THE PANEL'S HEAD, so that
  %      panels (a) and (b) put their first content line on ONE baseline, 6.70,
  %      0.75cm under their titles. Pass 12 set it at 5.90/5.50, which gave
  %      panel (a) a 1.55cm head-to-content gap against panel (b)'s 0.75cm and
  %      put panel (a)'s only substantial ink level with panels (c) and (d), so
  %      the lettering order and the spatial order disagreed.
  %      SUPERSEDES C10's placement paragraph ("three tiers with voids of 1.55
  %      and 3.22cm"); C10's four rejected ways of REMOVING the band all still
  %      stand, and the band is still deliberately empty, because it is the
  %      space between the prompt and the read, i.e. the residual stream.
  %      Clearance measured: line one is 4.880cm and line two 4.592cm, both
  %      clear of the rotated scope column's left edge at x = 5.62 by more than
  %      0.7cm at every y.
  % PASS 14: THE BLOCK MOVES BACK INTO THE BAND, at 5.00/4.60, centred on the
  % axis midspan 4.80. Pass 13's (ii) above bought a shared first-content
  % baseline with panel (b) and paid for it with ONE CONTIGUOUS VOID of
  % 56.2 x 39.2mm -- 17.2% of the figure's whole area, measured at 200dpi --
  % under the block and left of the axis, so panel (a) read as text at the top,
  % a stranded vertical stick in the middle and a hole between them. The two
  % TITLE baselines still align at 7.45, which is the alignment a reader
  % actually sees; the first-content alignment across a 7.75cm gutter is not.
  % Set here the block also sits level with the axis it describes, so "the
  % residual stream" has the rule and its seven ticks beside it rather than
  % 4cm above them. The band is now two voids of about 1.9 and 1.6cm rather
  % than one of 3.5cm, and its lower part carries the caret's direct label.
  % SUPERSEDES pass 13's repair (ii) at this block and C10's placement
  % paragraph as re-stated there; C10's four rejected ways of REMOVING the
  % band all still stand, and S4's "at most one block per panel" is why no
  % second prose block was written into it.
  \node[def] at (0.00,5.55) {what reads the residual stream:};
  \node[def] at (0.00,5.15) {a cross-validated linear probe};

  % --- the depth axis ------------------------------------------------------
  % PASS 12: a plain 0.40pt ink rule from the layer-2 tick to the layer-16
  % tick, terminating exactly on them. Before this pass it was a panelbg-filled
  % 0.36cm-wide column running from 2.10 to 7.05, i.e. past both extreme ticks,
  % which read as a bar whose length encoded nothing and implied depth beyond
  % the measured range.
  \draw[ink, line width=0.40pt] (6.90,2.70) -- (6.90,6.90);
  % y = 2.10 + (L/16)*4.80 for every one of the seven, unchanged from pass 11
  % and CHECKED against the map rather than against the eye:
  %   L   2    4    6    8    9    12   16
  %   y  2.70 3.30 3.90 4.50 4.80 5.70 6.90
  % Tick length 0.12cm = 3.4pt, PASS 13, down from 0.16cm: 4.11's own ticks are
  % about 3.5pt and the 0.04cm bought here is what lets the quantity name sit
  % beside the numerals without touching them.
  \foreach \L/\y in {2/2.70, 4/3.30, 6/3.90, 8/4.50, 9/4.80, 12/5.70, 16/6.90}{
    \draw[rulegrey, line width=0.30pt] (6.80,\y) -- (6.90,\y);
    \node[tick] at (6.74,\y) {\L};}
  % AXIS LABEL, in axis-label position: rotated 90 degrees reading upward,
  % centred on the axis midspan y = 4.80. PASS 13, TWO REPAIRS.
  %  (i) THE RANKING IS INVERTED BACK. Pass 12 put the quantity name OUTERMOST,
  %      at x = 5.36, with two wrapped \scriptsize grey lines between it and the
  %      numerals, so the thing immediately left of the tick numerals was a grey
  %      sentence and the quantity itself was 1.4cm further out -- the opposite
  %      of what S5 asks for. "layer" now sits immediately left of the numerals
  %      and the scope note outboard of it.
  % (ii) THE SCOPE NOTE IS ONE LINE, NOT A WRAPPED TWO. Pass 12's note measured
  %      9.09cm, wrapped onto two typeset lines of 4.817 and 4.202cm, and the
  %      longer of the two OVERRAN a 4.20cm axis at both ends -- 3.0mm past the
  %      layer-16 tick above and the same below -- which is the exact defect D6
  %      records itself as having removed from the column. It also broke the
  %      noun phrase across the wrap ("the last prompt / position"). The note is
  %      now "read at the last prompt position", 111.20pt = 3.908cm, which lies
  %      INSIDE the 4.20cm axis with 0.146cm of clearance at each end. "where
  %      generation begins" is caption material and is in the caption.
  %      SUPERSEDES C1's wrap arithmetic and its "line 1 is the outermost"
  %      clause in full.
  % x geometry, measured: at inner sep=0 a rotated \footnotesize line is 0.35cm
  % wide in x and a \scriptsize one 0.28cm; the widest numeral ("16") is 8.0pt =
  % 0.281cm and is right-anchored on 6.70, so it begins at x = 6.419.
  % scope 5.62-5.90 | layer 6.005-6.355 | numerals 6.419-6.70 | ticks 6.78-6.90.
  % PASS 14, TWO REPAIRS AT THIS BLOCK.
  %  (i) THE SCOPE LINE IS DELETED AND ITS CONTENT MOVES TO THE MARK IT NAMES.
  %      "read at the last prompt position" shared the five-word clause "at the
  %      last prompt position" with the caption that actually prints
  %      (04b:124-126, "the probe reads the residual stream at the last prompt
  %      position, where generation begins"), which S18 forbids; and it named,
  %      at 90 degrees and 1.14cm away, the one glyph on the page that had no
  %      name of its own -- the caret. 4.11's discipline is to direct-label at
  %      the mark, so the position is now labelled AT the caret, horizontally,
  %      in \scriptsize ink, and the axis keeps only its quantity name. S5
  %      makes the scope line optional; the mandatory \footnotesize roman ink
  %      quantity name in axis-label position is still there and is still the
  %      only thing between the tick numerals and panel (a)'s body.
  %      SUPERSEDES pass 13's repair (ii) at this block in full.
  % (ii) "layer" MOVES OFF THE LAYER-9 TICK. The axis midspan and the layer-9
  %      position are the SAME y by arithmetic, not by accident:
  %      2.10 + (9/16)*4.80 = 4.80 = (2.70 + 6.90)/2. So a rotated label
  %      centred on the midspan was centred on one numeral, 1.96mm from it, and
  %      read "layer 9". Measured, "layer" is 21.46pt = 0.754cm long, and the
  %      seven tick gaps are 0.60/0.60/0.60/0.30/0.90/1.20cm, so 5.70-to-6.90
  %      is the ONLY stretch of this axis long enough to hold the label with no
  %      numeral inside its span. Centred at 6.30 it spans 5.923 to 6.677 and
  %      the nearest numerals are 12 at 5.70 and 16 at 6.90. It stays in
  %      axis-label position: x 6.005-6.355, immediately left of the numeral
  %      column at 6.419-6.70, inside S5's 12pt band either side of that column.
  \node[font=\footnotesize, text=ink, rotate=90, anchor=center]
    at (6.18,6.30) {layer};

  % --- the prompt strip ----------------------------------------------------
  % Edges from the measured field-name widths: 6.084cm of ink, 0.816cm of
  % padding shared as 0.0816cm a side. See C8.
  % PASS 14: THE WHOLE STRIP DROPS 0.14cm, from y 1.66-2.28 to y 1.52-2.14.
  % NO CELL EDGE IN x MOVED, no field changed, no order changed and the strip
  % is still 6.90cm wide with the same five measured widths -- the change is
  % one rigid translation in y, and it buys the 0.14cm the caret's direct label
  % needs between the strip's top rule and the layer-2 numeral (see the read
  % position below). The room exists because pass 14 deletes the span rule that
  % occupied 1.46-1.54; the count line's baseline at 1.18 is unmoved and its
  % ascenders now clear the strip's bottom rule by 0.14cm.
  \draw[cell] (0.00,1.72) rectangle (1.12,2.14);
  \draw[cell] (1.90,1.72) rectangle (2.64,2.14);
  \draw[cell] (2.64,1.72) rectangle (4.23,2.14);
  \draw[cell] (4.23,1.72) rectangle (6.90,2.14);
  % the drug cell: the only accent mark in panel (a). Same construction as
  % hook.tex's \tokdrug -- accent border, accent!12 fill, rounded corners=1pt.
  \filldraw[fill=accent!12, draw=accent, line width=0.40pt, rounded corners=1pt]
    (1.12,1.72) rectangle (1.90,2.14);

  % baseline 1.775 = the strip's vertical centre 1.83 less 0.055, which centres
  % the x-height band rather than the box. (Pass 13: 1.915 on a centre of 1.97.)
  \node[font=\scriptsize, text=ink,    inner sep=0pt] at (0.560,1.92) {cell line};
  \node[font=\scriptsize, text=accent, inner sep=0pt] at (1.510,1.92) {drug};
  \node[font=\scriptsize, text=ink,    inner sep=0pt] at (2.270,1.92) {dose};
  \node[font=\scriptsize, text=ink,    inner sep=0pt] at (3.435,1.92) {mechanism};
  \node[font=\scriptsize, text=ink,    inner sep=0pt]
    at (5.565,1.92) {control cell sentence};

  % --- the read position ---------------------------------------------------
  % The caret straddles the strip's right edge, which is where the prompt ends.
  % It is ink, not accent: it is part of the axis apparatus, and accent in this
  % figure means the drug field and the surviving branch and nothing else.
  % PASS 13: IT NO LONGER TOUCHES THE AXIS AND NO LONGER SITS ON A NUMBERED
  % TICK. Pass 12 put its apex at (6.90, 2.70), which is simultaneously the
  % axis rule's base, the layer-2 tick and the largest filled black glyph in
  % panel (a); measured on the render the apex and the layer-2 tick coincided
  % to within a pixel. Three objects at one coordinate read two ways, both
  % false: as a filled arrowhead terminating the depth axis downward, and as a
  % mark AT LAYER 2 on the figure's only scaled axis -- which also falsified
  % the foot line, since the caret would then occupy a position on the depth
  % scale. It marks a token position, not a depth. The apex is now 2.50, a
  % clear 0.20cm below the axis base at 2.70, and the triangle is back to the
  % 0.24cm width of pass 9 (base 6.78 to 7.02, so it still straddles the
  % strip's right edge at 6.90 symmetrically and overhangs it by 0.12cm rather
  % than 0.15cm). It is named by the axis scope line 0.9cm to its left, which
  % now reads "read at the last prompt position".
  % SUPERSEDES C8's clause "the apex is the axis rule's own base at the layer-2
  % tick, y = 2.70", and restores the pass-9 note's own reading that the apex
  % belongs at a position the depth scale does not number.
  % PASS 14: THE CARET IS NAMED AT ITSELF, AND THE THREE CUES THAT MADE IT AN
  % ARROWHEAD ARE BROKEN ONE BY ONE. Pass 13 moved it 2mm and did not resolve
  % the reading; all three refuters read it, unaided, as the axis's terminal
  % arrowhead, and two of them read the layer-2 numeral as its label. Measured
  % on the pass-13 render, three cues did that: (1) it was the only
  % arrowhead-shaped glyph left in the figure and the largest solid black one
  % in the panel; (2) its apex stood on the axis rule's own x, 2.00mm below the
  % rule's bare square terminus; (3) its base ran 6.78 to 7.02 while the seven
  % ticks run 6.78 to 6.90, so its left half lay exactly in the tick lane, and
  % the "2" numeral's box bottom sat 0.25mm above its apex. Pass 14:
  %   * it acquires a \scriptsize ink DIRECT LABEL at itself, "the last prompt
  %     position", right-anchored at 6.68 so its ink stops 0.08cm from the
  %     caret's left base -- inside the document's ~0.3cm leader threshold, so
  %     no leader is owed. This is where that string was always meant to be:
  %     it was set rotated in axis-label position until pass 14, where a reader
  %     attached it to the axis rather than to the triangle;
  %   * the strip drops 0.14cm, so the apex is now 0.34cm below the axis base
  %     (was 0.20cm) and the label sits in the gap that opens;
  %   * the base widens from 0.24 to 0.28cm, 6.76 to 7.04, so it no longer
  %     shares an endpoint with the tick lane. It still straddles the strip's
  %     right edge at 6.90 symmetrically, which is the pass-9 relation the keep
  %     list names, and it is still ink, not rulegrey (O3 stands: rulegrey is
  %     furniture and this mark is content).
  % Under the figure's own map y = 2.10 + (L/16)*4.80 the triangle now occupies
  % y 2.14 to 2.36, i.e. layers 0.13 to 0.87 -- BELOW layer 1, off the drawn
  % scale, which is what "it marks a token position, not a depth" requires and
  % what the pass-13 ledger claimed without arithmetic. SUPERSEDES pass 13's
  % apex/base clause at this block and licensing.json's read_position note.
  \fill[ink] (6.90,2.36) -- (6.76,2.14) -- (7.04,2.14) -- cycle;
  \node[font=\scriptsize, text=ink, anchor=base east] at (6.68,2.28)
    {the last prompt position};

  % --- the span: the five fields are one prompt ----------------------------
  % PASS 13: A PLAIN SPAN RULE WITH DOWNTURNED ENDS REPLACES THE MIRRORED
  % BRACE. A brace has a central pointer, and pass 12 moved the count line from
  % centred-on-the-brace to left-aligned on the panel origin (S13) without
  % moving the pointer, so the notch at x = 3.45 aimed at 3.45cm of empty white
  % while its label began at x = 0.00. A span rule carries the same "these five
  % fields are one prompt" grouping, has no pointer to aim wrongly, and drops
  % the decorations.pathreplacing dependency.
  % PASS 14: THE SPAN RULE IS DELETED. All three refuters read it and all three
  % called it. Measured on the pass-13 render it was a rulegrey 0.30pt rule of
  % the strip's exact x-extent, 0.12cm below and parallel to the strip's own
  % rulegrey 0.30pt bottom border -- same colour, same weight, same length --
  % so it reproduced as a doubled edge or an offset re-impression rather than
  % as a bracket, its 0.08cm downturns vanishing at print size. Its left
  % downturn also ended 0.78mm directly above the "4" of "480" and read as a
  % stray tick on the numeral. And it grouped ONE schematic prompt's five
  % fields while the line beneath it counts 480 prompts, so on a strip that
  % carries no horizontal scale a full-width bracket labelled 480 invited
  % reading the width as the count. Nothing is lost: the five cells abut and
  % already read as one object, and the count line sits under the strip on the
  % strip's own left edge and names the population it is drawn from.
  % SUPERSEDES pass 13's R9 span-rule clause and the "the five fields are one
  % prompt" heading above it; the 0.14cm it frees is what lets the strip drop.
  % The count line names the population the spanned object is drawn from. It is
  % left-aligned to panel (a)'s origin, so panel (a) has exactly one x origin.
  % \scriptsize: at \footnotesize it measures 8.196cm and would run past the
  % panel's own measure. PASS 13: 480 is set as TEXT, not as $480$. In math
  % mode the three digits rendered in URWPalladioL-Roma while every other
  % numeral in the figure -- the seven tick labels, and the 2 of "tier-2" on
  % this very line -- rendered in TeXGyrePagella-Regular. One figure, two digit
  % faces, three characters. The measured width is 186.58pt = 6.5575cm either
  % way, so no coordinate moved and no value changed.
  \node[qlab] at (0.00,1.38)
    {480 real tier-2 prompts: twelve drugs, forty cells each};

% ###########################################################################
% # PANEL (b) -- the two ways a model can ignore a drug it was told about
% ###########################################################################
  % PASS 14: THE TITLE IS A NAME, NOT A SENTENCE. "the two ways a model can
  % ignore a drug it was told about" measured 270.16pt = 9.495cm, which made a
  % TITLE the widest object on the page, set the figure's whole right measure,
  % left 0.055cm of optical margin, and outweighed panel (a)'s title three to
  % one on the baseline they share, so the two heads did not read as peers. It
  % also restated 04b:31-32's own opening sentence ("A model that ignores a
  % drug it was told about can fail in two ways"), which is the paragraph the
  % float is attached to. 4.11's two titles run 29 and 38 characters; this one
  % ran 60. "two ways to ignore a named drug" is 165.62pt = 5.821cm, keeps the
  % "was told about" sense in the word NAMED -- which is also what panel (a)'s
  % title says enters -- and ends at x = 13.571. The full sentence is in the
  % caption. SUPERSEDES pass 13's O2 in its measurement clause: panel (a)'s
  % title is no longer constrained by this one. It is left at 4.969cm anyway,
  % because every longer variant either overruns panel (a)'s own 6.90cm measure
  % or says the probe reads the drug name, which it does not -- it reads the
  % stream, and the read is named at the caret and at the axis.
  \node[ttl] at (7.75,7.45)
    {(b) two ways to ignore a named drug};

  % --- failure mode one: struck --------------------------------------------
  % quietink, because it is eliminated. TWO strikes, one through each line: the
  % entry is two lines and a rule through its first line only left the second
  % standing. Each overhangs its line's ink by 0.06cm at each end, which is what
  % a strikethrough does. Endpoints from the measured \scriptsize widths,
  % 4.6741 and 4.2918cm: 7.69 to 12.484 and 7.69 to 12.102.
  % PASS 13, THE COLOUR AND THE WEIGHT. The strikes were ink at 0.90pt, and
  % that is the register S6 and CF2 reserve, across the whole slate, for a
  % threshold: chance in 4.11(a), the removal gate in 4.9, the pre-registered
  % materiality margin in 4.12 and 4.13, and style_v2.chance_line()'s own
  % hard-coded color='black', lw=0.9. C6's defence was that this figure
  % contains no threshold so nothing collides -- which answers S10's singleton
  % clause but not S6's, whose check is that NOTHING WHICH IS NOT A THRESHOLD
  % may be black solid 0.90pt. This figure prints six pages before the reader
  % meets the first real threshold in the arc, and it is the first figure of
  % the slate, so it was teaching the threshold stroke on a negation. Measured
  % against it: the struck text is quietink (grey 107) and the strikes were ink
  % (grey 25), so the branch the chapter DISCARDS carried the darkest ink on
  % the page while the surviving branch receded -- the S8 inversion C3 fixed in
  % colour, reintroduced in weight. Both strikes are now quietink 0.40pt, the
  % colour and weight of the text they cross. 0.40pt is already live here (the
  % axis rule and the drug cell's border), so no new register appears and no
  % singleton is created; 0.50pt was rejected because S10 reserves it for a
  % dotted coordinate. SUPERSEDES C6's third weight in full: the figure's
  % weights are now 0.30pt furniture and 0.40pt, and no 0.90pt stroke exists.
  % Strike y is the line's baseline + 0.065, half the \scriptsize x-height.
  % PASS 14, TWO REPAIRS.
  %  (i) THE STRIKES START ON THE PANEL'S OWN ORIGIN. They began at 7.69, which
  %      is 0.06cm (1.7pt) LEFT of the 7.75 every other text line in panels (b)
  %      and (c) starts on, so the panel's left margin was ragged and the item
  %      protruding furthest left was the one the figure discards. A
  %      strikethrough needs overhang at its TERMINAL, where the eye leaves the
  %      word, not at its initial; the right ends are unchanged at 12.484 and
  %      12.102, still 0.06cm past their lines' ink at 12.424 and 12.042.
  % (ii) THE OUTCOME LEADER IS DELETED (both of them; see the surviving mode).
  %      Measured on the pass-13 render this one ran 12.72 to 13.34 with a
  %      0.24cm gap at its left end and a 0.10cm gap at its right, so it
  %      touched neither the strike nor the label and read as a free-floating
  %      dash. Worse, it was a THIRD quietink horizontal lying midway between
  %      two quietink strikes 0.35cm apart, at the same colour and near the
  %      same weight, so the eye completed strike + gap + leader + label into
  %      one continuous rule and read the verdict "struck by the positive," as
  %      itself struck -- the exact "strike and leader read as one continuous
  %      pointer" fault the pass-9 note records itself as having killed on the
  %      diagonal. Widening the gap did not kill it because the stroke was
  %      still collinear in band, colour and register. Nothing is owed under
  %      S22: there is no MARK here, only a text block, and the pairing is
  %      carried twice over -- by colour (grey block to grey label, accent
  %      block to accent label) and by baselines shared to 0.00pt, which is
  %      4.11's own discipline. SUPERSEDES pass 13's lead/.style rationale and
  %      C6's "the two outcome leaders" clause: the figure now has one stroke
  %      weight for furniture (0.30pt) and one for the strikes and the axis
  %      (0.40pt), and no connector of any kind.
  \node[ded] at (7.75,6.70) {the drug never enters its computation};
  \node[ded] at (7.75,6.35) {an activation-side encoding failure};
  \draw[quietink, line width=0.40pt] (7.75,6.765) -- (12.484,6.765);
  \draw[quietink, line width=0.40pt] (7.75,6.415) -- (12.102,6.415);
  \node[qlab] at (13.00,6.70) {struck by the positive,};
  \node[qlab] at (13.00,6.35) {and only this one};

  % --- failure mode two: survives ------------------------------------------
  % accent, because it is the object the rest of the chapter measures. The two
  % leaders are identical in weight; only the label differs.
  % PASS 13: EACH OUTCOME LABEL NOW SITS ON ITS MODE'S OWN BASELINES, 6.70/6.35
  % and 5.55/5.20. Pass 12 set the modes \footnotesize on 0.40cm leading and
  % their labels \scriptsize on 0.30cm leading, each block centred on the other,
  % so no line matched its partner and the pair splayed in opposite directions
  % across a 55mm gutter -- measured, 0.41mm apart at the first line and 0.60mm
  % the other way at the second. One leading per size now holds across the whole
  % figure: 0.40cm at \footnotesize, 0.35cm at \scriptsize.
  % The outcome column moved from 14.28 to 13.44, which is what the \scriptsize
  % modes freed. Its leader starts 0.50cm clear of the longer strike's terminal
  % (12.484), against 0.17cm in pass 12, so strike and leader cannot read as one
  % continuous pointer -- the fault that killed the pass-1 diagonal.
  % PASS 14: THIS LEADER IS DELETED TOO, AND IT WAS THE WORSE OF THE PAIR.
  % Measured on the pass-13 render, the accent mode's longest line ends at
  % x = 11.328 and the leader began at 12.72, so it started 1.397cm -- 39.6pt,
  % 14mm -- past the mark it claimed to join and covered 0.61 of a 2.13cm gap:
  % at print size a 6.1mm dash floating in 14mm of white with nothing on its
  % left. The pass-13 ledger recorded the OPPOSITE as a measured property
  % ("the leader now starts 0.236cm past the longer strike's terminal"), which
  % is true of the struck mode's leader only; the accent one's 1.397cm gap was
  % never measured. Two leaders of visibly unequal reach were the alternative
  % and were rejected: the struck one cannot be extended leftward without
  % butting the strike and restoring the continuous-pointer reading.
  % The outcome column moves 13.44 -> 13.00, which halves the unlinked gap on
  % this mode from 2.11 to 1.67cm and still clears the longer strike's terminal
  % (12.484) by 0.516cm, the 0.5cm separation the pass-9 note asks for.
  % SUPERSEDES pass 13's O4 arithmetic (13.44) and its "the two leaders are
  % identical in weight" clause -- there are no leaders.
  \node[liv] at (7.75,5.55) {it enters and nothing reads it};
  \node[liv] at (7.75,5.20) {a decoder failure};
  \node[alab] at (13.00,5.55) {survives, and is the rest};
  \node[alab] at (13.00,5.20) {of the chapter};

% ###########################################################################
% # PANELS (c) AND (d) -- the two licences, as two columns read across
% ###########################################################################
  % PASS 14: THE WHOLE (c)/(d) GROUP DROPS 0.30cm AND ITS ROW PITCH OPENS FROM
  % 0.52 to 0.70cm. Deleting the third row (see the entries below) left the
  % right half ending at y = 2.55 against panel (a)'s count line at 1.18, i.e.
  % a fresh 9.5 x 1.8cm void in the bottom-right corner -- trading one hole for
  % another is not a repair. Moving the rule to 4.05 also makes the gap from
  % panel (b)'s last line (5.20) to the rule 1.15cm, which is exactly the gap
  % between panel (b)'s own two mode blocks, so the right half now has one
  % vertical rhythm. The title still sits 0.45cm under the rule, as before.
  \draw[rulegrey, line width=0.30pt] (7.75,4.05) -- (17.30,4.05);

  % PASS 13: BOTH TITLES ARE ONE LINE ON ONE BASELINE. (d) read "what a
  % negative would have bought", 6.167cm, which does not fit a column and
  % wrapped, so (c)'s head-to-body gap was 0.90cm against (d)'s 0.54cm on two
  % columns the ledger called equal. "would buy" is 4.847cm and the conditional
  % still carries the counterfactual against (c)'s indicative "buys"; the
  % perfect tense survives in the caption. Origins are 7.75 and 12.40: (d)'s
  % title then ends at 17.247, which is 0.053cm inside the measure, exactly the
  % clearance panel (b)'s title already runs at. Columns are 4.65 and 4.90cm,
  % and the content footprints are 4.081/4.847 (titles) and 3.210/3.697
  % (entries), so neither column is the larger and more detailed one -- which is
  % what C4's equalisation was for. SUPERSEDES C4's "equal width (4.60cm)",
  % which was never true of the drawing: pass 12's origins were 7.75 and 12.70,
  % i.e. 4.95 and 4.60.
  \node[ttl] at (7.75,3.60) {(c) what the positive buys};
  \node[ttl] at (12.40,3.60) {(d) what a negative would buy};

  % THREE ENTRIES EACH, PAIRED ROW BY ROW. PASS 13, AND THIS IS THE CORRECTNESS
  % REPAIR IN THIS PANEL, NOT A TYPOGRAPHIC ONE.
  %  (i) THE ACROSS-READING WAS OFF BY ONE. Pass 12 set four baselines in each
  %      column and asserted, in the ledger and in the caption, that they were
  %      to be read across entry by entry. They could not be: (c) was two
  %      sentences word-wrapped over four lines and (d) a single four-item
  %      comma list, so row 1 paired "the model did not discard" with "the drug
  %      lost inside the model," and the negation of the layer sat one row
  %      above the layer it negated. Three of the four pairings were nonsense.
  %      The rows are now grouped, not merely stacked: 0.32cm of leading WITHIN
  %      a row and 0.52cm BETWEEN rows, so a wrapped entry cannot be mistaken
  %      for the next entry. Each row's first line is on one baseline in both
  %      columns, and the three pairs are exact.
  % (ii) "no class of failure," IS DELETED, because it contradicts this
  %      figure's own panel (b) and the claim box it is drawn from. 04b:142-144
  %      reads "It rules out that the drug field is lost before the decoder
  %      runs, an encoding failure in the activations", and panel (b) says so
  %      on the page, 3.5cm above, in the words "struck by the positive, and
  %      only this one". A positive that strikes one of two modes has bought a
  %      class of failure ruled out; what it has not bought is a depth and a
  %      fix. The two entries that survive, "no layer" and "no class of fix",
  %      are the two the sources license: 04b:110-112, "what the probe buys is
  %      the depth at which one would appear", and the caption's own lede, that
  %      the negative "would have named the layer at which the drug was lost
  %      and the class of fix that could reach it".
  % (iii)"an activation-side failure, and" IS DELETED FROM (d), because panel
  %      (b)'s struck line already reads "an activation-side encoding failure"
  %      2.5cm away. One object, one name, drawn once; the fix classes stay,
  %      because they are named nowhere else in the drawing.
  % (iv) (d)'s first entry carries the hedge the instrument requires. The probe
  %      is linear -- panel (a) now says so -- and a linear decode failure
  %      licenses only that the drug is not linearly readable, not that it is
  %      absent. The claim box is careful about exactly this ("Its identity is
  %      linearly readable from the residual stream at every layer measured");
  %      pass 12's "the drug lost inside the model," was not, and since the
  %      figure's whole argument is a comparison of licence SIZES, an
  %      overstated counterfactual inflates the asymmetry it exists to draw.
  %      SUPERSEDES the pass-9 note "THE COUNTERFACTUAL BLOCK'S FIRST ENTRY
  %      LOST ITS FINITE VERB" as to its wording only; the reason that note
  %      gives -- that the entry must be a noun phrase, not an assertion that
  %      the drug IS lost -- is preserved and strengthened by the hedge.
  % (c) is ink: it happened. (d) is quietink: it did not.
  % Leading is 0.35cm WITHIN a row -- the one \scriptsize leading in the figure,
  % shared with panel (b) -- and 0.52cm BETWEEN rows.
  % PASS 14, THREE ENTRIES BECOME TWO PAIRED ROWS OF ONE LINE EACH, AND ONE OF
  % THE DELETIONS IS A CORRECTNESS REPAIR, NOT A TYPOGRAPHIC ONE.
  %  (i) "no class of fix" IS DELETED, BY EXACTLY THE ARGUMENT PASS 13 USED TO
  %      DELETE "no class of failure," ONE ROW ABOVE IT -- and that argument
  %      was never applied to this entry, in the header or in the ledger.
  %      04b:31-36 is a TWO-MEMBER partition and each member names its own fix
  %      class: mode one is "an activation-side encoding failure that target-
  %      and input-side fixes might reach", mode two "a decoder failure only an
  %      objective-side fix could reach". Panel (b) strikes mode one and says
  %      so on this page, 3.5cm above, in the words "struck by the positive,
  %      and only this one". If one member of an exhaustive pair is struck the
  %      survivor stands, and the survivor names its fix class -- so within
  %      this figure's own frame the positive DOES buy a class of fix, and the
  %      entry is false. It survived pass 13 only because panel (b) draws mode
  %      two's gloss truncated to "a decoder failure", dropping the very clause
  %      of its source sentence that refutes the entry. A drawing may not
  %      contradict, at 8pt, the box it cites as its source.
  %      The alternative repair -- rewriting the entry to the class the
  %      positive does buy, "objective-side fixes only", against (d)'s
  %      "target- and input-side fixes" -- was built and rejected: that licence
  %      holds only CONDITIONALLY, on the model in fact ignoring the drug,
  %      which is 4.4's and 4.5's finding and not this section's, and a
  %      two-word slot cannot carry the condition without overstating it. Both
  %      fix classes go to the caption, which can state the condition.
  %      (d)'s "target- and input-side fixes" goes with it, because a row with
  %      one side is not a row and the entry is named nowhere else it is owed.
  % (ii) ROW ONE LOSES ITS WRAP AND ITS S18 COLLISION. "the model did not
  %      discard what it was handed" is seven words verbatim from 04b:107-108,
  %      the paragraph the float is attached to ("recovering it from the
  %      activations establishes only that the model did not discard what it
  %      was handed"), six lines of source above the float -- which S18 forbids
  %      outright. It was also the only entry in either column that wrapped, so
  %      (c) read as three items against (d)'s two and the across-reading broke
  %      at row one: measured, the white gap between rows 1 and 2 was 6.76pt in
  %      (c) and 16.68pt in (d), a 2.47x difference in the same two rows.
  %      "the drug field is not discarded" is 106.48pt = 3.742cm, one line,
  %      inside (c)'s 4.65cm column, keeps the sense the positive licenses, and
  %      shares no five-word clause with the prose or with any caption.
  % (iii)WITH EVERY ENTRY ON ONE LINE THERE IS NO WITHIN-ROW LEADING LEFT, so
  %      the two rows sit at one pitch of 0.60cm in both columns and the
  %      pairing needs no rule and no cue. The asymmetry the figure exists to
  %      draw is now also visible as ink for the first time: (c) sets 39
  %      characters against (d)'s 58, and (c)'s second entry is 0.991cm against
  %      (d)'s 3.233cm, so the small licence is the small column.
  %      SUPERSEDES pass 13's (ii)/(iii) rationale at this block as to the
  %      entries that survive, and C4's "FOUR ENTRIES EACH" in full.
  \node[lab]  at (7.75,2.85) {the drug field is not discarded};
  \node[qlab] at (12.40,2.85) {the drug not linearly readable};

  \node[lab]  at (7.75,2.15) {no layer};
  \node[qlab] at (12.40,2.15) {the layer at which it is lost};

% --- the figure's one disclosure line ---------------------------------------
% PASS 13: REWORDED. "Nothing in this figure is a measured quantity" sat 0.9cm
% under a line printing 480, twelve and forty, and then exempted the depth
% ticks in its own second clause, so it read as a blanket claim followed
% immediately by two counterexamples. The counts are the probe's design, not
% its results, and saying so is shorter than conceding an exception. It also
% now carries the count "seven", which pass 12 deleted with the caret label and
% which the pass-9 layout note (still standing) argues is the one fact the
% depth column exists to carry: without it a reader cannot tell whether the
% ticks are the depths probed or a sample of a sweep over all sixteen.
% One line, 389.66pt = 13.695cm, left-aligned to the figure's x origin.
% PASS 14: THE SECOND CLAUSE GOES, AND THE FIRST IS WHAT WAS DOING THE WORK.
% Pass 13 added "the counts are the probe's design" because the pass-12 wording
% ("Nothing in this figure is a measured quantity") sat under a line printing
% 480, twelve and forty and read as a blanket claim with two counterexamples
% beneath it. But pass 13 made TWO changes at once, and it was the other one --
% "a measured quantity" becoming "an estimate" -- that removed the
% contradiction, since a count is not an estimate. The clause was therefore
% belt and braces on a line already sound, at 35 characters, in the longest
% string in the figure, on a drawing over S15's ink ceiling. At 76 characters
% the line still answers the acceptance question ("can the reader say, without
% leaving the figure, that nothing here is a measured quantity?") on its own
% and still carries the count "seven" the pass-9 layout note asks for. The line
% is now 9.60cm rather than 13.695cm and is no longer the widest \scriptsize
% object on the page. SUPERSEDES pass 13's R9 foot-line clause as to the second
% sentence only; its rewording of "measured quantity" to "estimate" stands.
  \node[qlab] at (0.00,0.38)
    {Intervals: none in any panel; nothing here is an estimate, and the seven
     depth ticks are the only scaled marks.};

  % the bounding box IS the fullwidth measure, less 1.8pt
  \path (0,0.16) rectangle (17.30,7.72);
\end{tikzpicture}

% ---------------------------------------------------------------------------
% READY TO PASTE into Sections/04b-representation.tex, at \beatsection
% sec:methods-probe, immediately after the paragraph "The probe, and why its
% positive is nearly free" (04b-representation.tex:108-113) and before
% "Forward passes only." Not written into the Section file by this build.
%
% THE PASS-11 CAPTION IS RETAINED BELOW, COMMENTED OUT AND MARKED SUPERSEDED,
% because the thesis is audited against these headers. It is FALSE against the
% pass-12 and pass-13 drawings in four places: it indexes the panels "Left,"
% and "Right," where the drawing sets four lowercase panel letters, it says the
% three accuracies and the floor are "set here as text" (they are not set here
% at all any more), it spends two sentences denying a reading the drawing no
% longer offers, and it says the two licence blocks "are sized by their
% contents", which was the footprint asymmetry pass 12 replaced.
%
% The PASS-13 REPLACEMENT IS WRITTEN TO figs/_PENDING_CAPTIONS.tex, keyed by
% \label{fig:licensing}, together with every sentence pass 12 and pass 13 cut
% from the drawing. PASS 13 CREATED THAT FILE: pass 12 asserted twice that it
% had, and it did not exist. This builder does not edit Sections/. The block
% below is a DUPLICATE of the pending file, kept so this file is self-contained
% and the audit trail is here; the pending file is the one an integrator reads.
%
% ###########################################################################
% PASS 14 SUPERSEDES THE WHOLE OF THE PARAGRAPH ABOVE, AND THE PASS-13 CAPTION
% BELOW IS NOW HISTORY, NOT THE PROPOSAL.
%
% (i) The paragraph above is FALSE AS WRITTEN. figs/_PENDING_CAPTIONS.tex did
%     exist when pass 13 ran (it was created by the pass-10 repair of
%     fig:materiality), and pass 13's APPEND TO IT DID NOT SURVIVE: read on
%     2026-08-19 the file held fig:materiality, fig:comparators and fig:hook
%     blocks and no fig:licensing block at all. All three pass-14 refuters
%     opened with that finding and all three were right. The likeliest cause is
%     a concurrent read-modify-write from another figure's builder, which is
%     the hazard CF12 names about this shared file.
% (ii) PASS 14 APPENDED the block, in place, to figs/_PENDING_CAPTIONS.tex, and
%     READ THE FILE BACK TWICE -- once immediately after the append and once at
%     the end of the pass. Both times the block was there and the
%     fig:materiality, fig:comparators and fig:hook blocks were all still
%     there. It was appended at the END rather than inserted in figure order
%     deliberately: an append cannot destroy another builder's block and an
%     ordered insert can, which is what happened to pass 13.
%     NO LINE NUMBER IS RECORDED FOR IT, and that is the point of the two
%     reads: between them another builder appended to the fig:comparators
%     block and every heading below it moved -- this block went from line 484
%     to 597 and fig:hook from 306 to 369 -- so a line number written into
%     this header would be false within the hour. The \label heading is the
%     stable address.
% (iii)THE PASS-14 CAPTION IS THE ONE IN THE PENDING FILE, and it is NOT
%     duplicated here. Keeping a second copy in this header is what let pass 13
%     believe it had staged a caption it had not; one copy, in the file an
%     integrator is told to read, is the safer arrangement. What is recorded
%     here instead is what the pass-14 caption differs from the pass-13 draft
%     below in, so the audit trail is complete without a duplicate:
%       * it prints NONE of the four accuracies, where the pass-13 draft set
%         all four. It names them ("the raw peak and the deepest, the
%         context-removed deepest and the $1/12$ floor") and points at
%         \cref{tab:workspace} and \cref{fig:mechanism}(a). This is the
%         cross-figure redundancy resolution applied to the caption as well as
%         to the drawing: 0.821 and 0.883 are owned by 4.11(a), the only place
%         the two series are drawn AS series with their crossing at layer 12
%         visible, and printing them a fourth time here is the redraw the
%         resolution exists to stop. The values stay in the prose at
%         04b:159-166, in the claim box, and in \cref{tab:workspace};
%       * it carries the two FIX CLASSES in full, with the premise they depend
%         on, because pass 14 deleted "no class of fix" and "target- and
%         input-side fixes" from panels (c) and (d);
%       * it restates each failure mode's gloss in full, because the drawing
%         abbreviates both;
%       * it drops "The two licence blocks are sized by their contents", which
%         the pass-13 draft had already dropped, and also drops "the figure
%         quantifies nothing about the size of a licence": with the boxes gone
%         and (c)/(d) set as plain text columns, nothing on the page reads as
%         an area, and a caption denying a reading the drawing no longer offers
%         is the defect that sentence was meant to fix;
%       * it says "the single position read, the last of the prompt" rather
%         than "at the last prompt position", because the drawing now sets
%         "the last prompt position" as the caret's direct label and S18 bars
%         a shared five-word clause.
%     MEASURED: 2554 characters, 22 typeset lines under the real preamble in a
%     fullwidth float with the pass-14 drawing above it, against the live
%     pass-11 caption's 1692 characters and 15 lines, and against the 19-22 the
%     rest of the 4.7-4.12 arc runs at. It compiles with zero errors and zero
%     Overfull boxes on one page.
% ###########################################################################
%
% ---- PASS 13 REPLACEMENT CAPTION, SUPERSEDED BY PASS 14, kept for the audit
% ---- trail. DO NOT PASTE THIS ONE. -----------------------------------------
% \begin{figure}[htbp]
% \begin{fullwidth}
% \centering
% \input{figs/licensing}
% \caption[The probe's cheap positive]{\textbf{The drug is named in plain text
% in the prompt, so a positive was near-certain before the probe was run: it
% eliminates one of the two ways the model could be ignoring the drug and
% leaves the other standing, whereas the negative that did not arrive would
% have named the depth at which the drug was lost and the class of fix that
% could reach it.} One run, two possible outcomes, and the one that could not
% fail is the one that happened; that asymmetry, and not any accuracy, is what
% this figure draws. \emph{(a)} The drug is one of five prompt fields
% (\cref{sec:methods-data}), and the probe reads the residual stream at one
% position, the last of the prompt, where generation begins; the caret marks
% that position and the axis beside it names the seven depths. The strip is a
% schematic of \emph{fields}: cells are sized to their typeset labels and not
% to their token counts, so it carries no horizontal scale, and the control
% cell's sentence, which runs to hundreds of tokens, is not drawn hundreds of
% times wider. Tick spacing on that axis is the one thing here drawn to scale,
% and the scale is layer index rather than accuracy --- the grid is even apart
% from the layer added at $9$. The classifier is cross-validated and
% multinomial over the twelve drug classes, with folds stratified over prompts
% and not grouped by well. \emph{(b)} The partition this chapter opens with.
% Accent marks the branch that survives, because it is what \cref{sec:res-used}
% and \cref{sec:res-mechanism} spend the rest of the chapter measuring; the
% eliminated branch and the counterfactual column are quiet. \emph{(c)} and
% \emph{(d)} are read across, row by row: each row sets what the positive
% bought beside what a negative would have bought in the same slot. Because the
% instrument is linear, a negative would have licensed only that the drug is
% not linearly recoverable, which is the form \emph{(d)} states it in.
% \textbf{No accuracy is drawn or printed here.} The probe reads $0.821$ at
% layer $9$ and $0.758$ at layer $16$ against a floor of $1/12 \approx 0.083$,
% and $0.883$ at layer $16$ once cell line, plate and dose are projected out;
% those four are \cref{tab:workspace}'s and \cref{fig:mechanism}(a)'s, which
% plots them against depth on a real accuracy axis and is the only place a
% reader can see that the context-removed series \emph{exceeds} the raw one at
% layers $12$ and $16$ rather than sitting under it. No mark on this page
% carries an interval and none could: the probe returns a cross-validated
% accuracy and \texttt{workspace\_probe.json} stores no dispersion for it. The
% two licence columns quantify nothing about the size of a licence.}
% \label{fig:licensing}
% \end{fullwidth}
% \end{figure}
%
% ---- WHAT THE PASS-12 CAPTION SAID, AND WHY PASS 13 CHANGED IT -------------
% The pass-12 replacement was never integrated (it was staged in this file and
% nowhere else). It differed from the block above in four clauses, all four
% superseded here by name rather than silently rewritten:
%   (i)   "\emph{(c)} and \emph{(d)} are two equal columns and are meant to be
%         read across, entry by entry" -- it described a correspondence the
%         pass-12 geometry did not draw: four baselines carrying a two-sentence
%         block against a four-item list, so three of the four pairings were
%         nonsense. Pass 13 draws three paired rows and the caption now says
%         what each row does.
%   (ii)  it did not say the instrument is linear, so it did not carry the
%         hedge a linear decode failure actually licenses. The new clause
%         "Because the instrument is linear ..." does.
%   (iii) it did not name the classifier's twelve classes or the fold design;
%         both were drawn up to pass 11, "twelve classes" left the drawing in
%         pass 13 for measure (see the prose block) and the folds line left it
%         in pass 12 under D7, so the caption now carries both.
%   (iv)  "nothing on the page is an estimate, none of it carries an interval"
%         shared a clause with the figure's own foot line, which S18 forbids.
%         Reworded to "No mark on this page carries an interval and none
%         could", against a foot line that now reads "Nothing here is an
%         estimate: ...".
%
% ---- SUPERSEDED, PASS 11 CAPTION, kept for the audit trail -----------------
% \caption[The probe's cheap positive]{\textbf{The positive rules out
% one of the two ways the model could be ignoring the drug and localises
% nothing, while the negative that did not arrive would have named the layer at
% which the drug was lost and the class of fix that could reach it.}
% \emph{Left,} the drug is one field of the prompt
% (\cref{sec:methods-data}), and the probe reads the residual stream at the
% last prompt position, where generation begins, at each of seven layers. The
% three accuracies and the floor are read from
% \texttt{workspace\_probe.json} and are also printed in
% \cref{tab:workspace}: $0.821$ at layer 9 and $0.758$ at the deepest against a
% $0.083$ shuffled-label floor, and $0.883$ at the deepest once cell line,
% plate and dose are projected out. They are set here as text, not as marks on
% a rule; the accuracy rule is \cref{fig:comparators}, row one. The prompt
% strip is a schematic of \emph{fields}: cells are sized to their labels and
% not to their token counts, so it carries no horizontal scale, and the control
% cell's sentence, which is hundreds of tokens, is not drawn hundreds of times
% wider. The tick spacing on the stack is the only scaled thing in
% the figure, and it is depth, not accuracy: the grid is even apart from the
% layer added at $9$. \emph{Right,} the partition this section opens with. A
% positive strikes the first mode and only the first; the second is what
% \cref{sec:res-used} and \cref{sec:res-mechanism} spend the rest of the
% chapter on. The two licence blocks are sized by their contents and by nothing
% else --- the figure quantifies nothing about the size of a licence, carries
% no interval, and no mark in it may be read as a measured quantity.}
% ---------------------------------------------------------------------------

% \caption[The probe's cheap positive]{\textbf{The drug is named in plain text
% in the prompt, so a positive was near-certain before the probe was run: it
% eliminates one of the two ways the model could be ignoring the drug and
% leaves the other standing, whereas the negative that did not arrive would
% have named the depth at which the drug was lost and the class of fix that
% could reach it.} One run, two possible outcomes, and the one that could not
% fail is the one that happened; that asymmetry, and not any accuracy, is what
% this figure draws. \emph{(a)} The drug is one of five prompt fields
% (\cref{sec:methods-data}), and the probe reads the residual stream at one
% position, the last of the prompt, where generation begins; the caret marks
% that position and the axis beside it names the seven depths. The strip is a
% schematic of \emph{fields}: cells are sized to their typeset labels and not
% to their token counts, so it carries no horizontal scale, and the control
% cell's sentence, which runs to hundreds of tokens, is not drawn hundreds of
% times wider. Tick spacing on that axis is the one thing here drawn to scale,
% and the scale is layer index rather than accuracy --- the grid is even apart
% from the layer added at $9$. The classifier is cross-validated and
% multinomial over the twelve drug classes, with folds stratified over prompts
% and not grouped by well. \emph{(b)} The partition this chapter opens with.
% Accent marks the branch that survives, because it is what \cref{sec:res-used}
% and \cref{sec:res-mechanism} spend the rest of the chapter measuring; the
% eliminated branch and the counterfactual column are quiet. \emph{(c)} and
% \emph{(d)} are read across, row by row: each row sets what the positive
% bought beside what a negative would have bought in the same slot. Because the
% instrument is linear, a negative would have licensed only that the drug is
% not linearly recoverable, which is the form \emph{(d)} states it in.
% \textbf{No accuracy is drawn or printed here.} The probe reads $0.821$ at
% layer $9$ and $0.758$ at layer $16$ against a floor of $1/12 \approx 0.083$,
% and $0.883$ at layer $16$ once cell line, plate and dose are projected out;
% those four are \cref{tab:workspace}'s and \cref{fig:mechanism}(a)'s, which
% plots them against depth on a real accuracy axis and is the only place a
% reader can see that the context-removed series \emph{exceeds} the raw one at
% layers $12$ and $16$ rather than sitting under it. No mark on this page
% carries an interval and none could: the probe returns a cross-validated
% accuracy and \texttt{workspace\_probe.json} stores no dispersion for it. The
% two licence columns quantify nothing about the size of a licence.}
% \label{fig:licensing}
% \end{fullwidth}
% \end{figure}
%
% ---- WHAT THE PASS-12 CAPTION SAID, AND WHY PASS 13 CHANGED IT -------------
% The pass-12 replacement was never integrated (it was staged in this file and
% nowhere else). It differed from the block above in four clauses, all four
% superseded here by name rather than silently rewritten:
%   (i)   "\emph{(c)} and \emph{(d)} are two equal columns and are meant to be
%         read across, entry by entry" -- it described a correspondence the
%         pass-12 geometry did not draw: four baselines carrying a two-sentence
%         block against a four-item list, so three of the four pairings were
%         nonsense. Pass 13 draws three paired rows and the caption now says
%         what each row does.
%   (ii)  it did not say the instrument is linear, so it did not carry the
%         hedge a linear decode failure actually licenses. The new clause
%         "Because the instrument is linear ..." does.
%   (iii) it did not name the classifier's twelve classes or the fold design;
%         both were drawn up to pass 11, "twelve classes" left the drawing in
%         pass 13 for measure (see the prose block) and the folds line left it
%         in pass 12 under D7, so the caption now carries both.
%   (iv)  "nothing on the page is an estimate, none of it carries an interval"
%         shared a clause with the figure's own foot line, which S18 forbids.
%         Reworded to "No mark on this page carries an interval and none
%         could", against a foot line that now reads "Nothing here is an
%         estimate: ...".
%
% ---- SUPERSEDED, PASS 11 CAPTION, kept for the audit trail -----------------
% \caption[The probe's cheap positive]{\textbf{The positive rules out
% one of the two ways the model could be ignoring the drug and localises
% nothing, while the negative that did not arrive would have named the layer at
% which the drug was lost and the class of fix that could reach it.}
% \emph{Left,} the drug is one field of the prompt
% (\cref{sec:methods-data}), and the probe reads the residual stream at the
% last prompt position, where generation begins, at each of seven layers. The
% three accuracies and the floor are read from
% \texttt{workspace\_probe.json} and are also printed in
% \cref{tab:workspace}: $0.821$ at layer 9 and $0.758$ at the deepest against a
% $0.083$ shuffled-label floor, and $0.883$ at the deepest once cell line,
% plate and dose are projected out. They are set here as text, not as marks on
% a rule; the accuracy rule is \cref{fig:comparators}, row one. The prompt
% strip is a schematic of \emph{fields}: cells are sized to their labels and
% not to their token counts, so it carries no horizontal scale, and the control
% cell's sentence, which is hundreds of tokens, is not drawn hundreds of times
% wider. The tick spacing on the stack is the only scaled thing in
% the figure, and it is depth, not accuracy: the grid is even apart from the
% layer added at $9$. \emph{Right,} the partition this section opens with. A
% positive strikes the first mode and only the first; the second is what
% \cref{sec:res-used} and \cref{sec:res-mechanism} spend the rest of the
% chapter on. The two licence blocks are sized by their contents and by nothing
% else --- the figure quantifies nothing about the size of a licence, carries
% no interval, and no mark in it may be read as a measured quantity.}
% ---------------------------------------------------------------------------

% \caption[The probe's cheap positive]{\textbf{The drug is named in plain text
% in the prompt, so a positive was near-certain before the probe was run: it
% eliminates one of the two ways the model could be ignoring the drug and
% leaves the other standing, whereas the negative that did not arrive would
% have named the depth at which the drug was lost and the class of fix that
% could reach it.} One run, two possible outcomes, and the one that could not
% fail is the one that happened; that asymmetry, and not any accuracy, is what
% this figure draws. \emph{(a)} The drug is one of five prompt fields
% (\cref{sec:methods-data}), and the probe reads the residual stream at one
% position, the last of the prompt, where generation begins; the caret marks
% that position and the axis beside it names the seven depths. The strip is a
% schematic of \emph{fields}: cells are sized to their typeset labels and not
% to their token counts, so it carries no horizontal scale, and the control
% cell's sentence, which runs to hundreds of tokens, is not drawn hundreds of
% times wider. Tick spacing on that axis is the one thing here drawn to scale,
% and the scale is layer index rather than accuracy --- the grid is even apart
% from the layer added at $9$. The classifier is cross-validated and
% multinomial over the twelve drug classes, with folds stratified over prompts
% and not grouped by well. \emph{(b)} The partition this chapter opens with.
% Accent marks the branch that survives, because it is what \cref{sec:res-used}
% and \cref{sec:res-mechanism} spend the rest of the chapter measuring; the
% eliminated branch and the counterfactual column are quiet. \emph{(c)} and
% \emph{(d)} are read across, row by row: each row sets what the positive
% bought beside what a negative would have bought in the same slot. Because the
% instrument is linear, a negative would have licensed only that the drug is
% not linearly recoverable, which is the form \emph{(d)} states it in.
% \textbf{No accuracy is drawn or printed here.} The probe reads $0.821$ at
% layer $9$ and $0.758$ at layer $16$ against a floor of $1/12 \approx 0.083$,
% and $0.883$ at layer $16$ once cell line, plate and dose are projected out;
% those four are \cref{tab:workspace}'s and \cref{fig:mechanism}(a)'s, which
% plots them against depth on a real accuracy axis and is the only place a
% reader can see that the context-removed series \emph{exceeds} the raw one at
% layers $12$ and $16$ rather than sitting under it. No mark on this page
% carries an interval and none could: the probe returns a cross-validated
% accuracy and \texttt{workspace\_probe.json} stores no dispersion for it. The
% two licence columns quantify nothing about the size of a licence.}
% \label{fig:licensing}
% \end{fullwidth}
% \end{figure}
%
% ---- WHAT THE PASS-12 CAPTION SAID, AND WHY PASS 13 CHANGED IT -------------
% The pass-12 replacement was never integrated (it was staged in this file and
% nowhere else). It differed from the block above in four clauses, all four
% superseded here by name rather than silently rewritten:
%   (i)   "\emph{(c)} and \emph{(d)} are two equal columns and are meant to be
%         read across, entry by entry" -- it described a correspondence the
%         pass-12 geometry did not draw: four baselines carrying a two-sentence
%         block against a four-item list, so three of the four pairings were
%         nonsense. Pass 13 draws three paired rows and the caption now says
%         what each row does.
%   (ii)  it did not say the instrument is linear, so it did not carry the
%         hedge a linear decode failure actually licenses. The new clause
%         "Because the instrument is linear ..." does.
%   (iii) it did not name the classifier's twelve classes or the fold design;
%         both were drawn up to pass 11, "twelve classes" left the drawing in
%         pass 13 for measure (see the prose block) and the folds line left it
%         in pass 12 under D7, so the caption now carries both.
%   (iv)  "nothing on the page is an estimate, none of it carries an interval"
%         shared a clause with the figure's own foot line, which S18 forbids.
%         Reworded to "No mark on this page carries an interval and none
%         could", against a foot line that now reads "Nothing here is an
%         estimate: ...".
%
% ---- SUPERSEDED, PASS 11 CAPTION, kept for the audit trail -----------------
% \caption[The probe's cheap positive]{\textbf{The positive rules out
% one of the two ways the model could be ignoring the drug and localises
% nothing, while the negative that did not arrive would have named the layer at
% which the drug was lost and the class of fix that could reach it.}
% \emph{Left,} the drug is one field of the prompt
% (\cref{sec:methods-data}), and the probe reads the residual stream at the
% last prompt position, where generation begins, at each of seven layers. The
% three accuracies and the floor are read from
% \texttt{workspace\_probe.json} and are also printed in
% \cref{tab:workspace}: $0.821$ at layer 9 and $0.758$ at the deepest against a
% $0.083$ shuffled-label floor, and $0.883$ at the deepest once cell line,
% plate and dose are projected out. They are set here as text, not as marks on
% a rule; the accuracy rule is \cref{fig:comparators}, row one. The prompt
% strip is a schematic of \emph{fields}: cells are sized to their labels and
% not to their token counts, so it carries no horizontal scale, and the control
% cell's sentence, which is hundreds of tokens, is not drawn hundreds of times
% wider. The tick spacing on the stack is the only scaled thing in
% the figure, and it is depth, not accuracy: the grid is even apart from the
% layer added at $9$. \emph{Right,} the partition this section opens with. A
% positive strikes the first mode and only the first; the second is what
% \cref{sec:res-used} and \cref{sec:res-mechanism} spend the rest of the
% chapter on. The two licence blocks are sized by their contents and by nothing
% else --- the figure quantifies nothing about the size of a licence, carries
% no interval, and no mark in it may be read as a measured quantity.}
% ---------------------------------------------------------------------------

\caption[The probe's cheap positive]{\textbf{The drug is named in plain text in
the prompt, so a positive was near-certain before the probe was run: it
eliminates one of the two ways the model could be ignoring the drug and leaves
the other standing, whereas the negative that did not arrive would have named
the depth at which the drug was lost and the class of fix that could reach
it.} One run, two possible outcomes, and the one that could not fail is the one
that happened; that asymmetry, and not any accuracy, is what this figure draws.
\emph{(a)}~The drug is one of five prompt fields (\cref{sec:methods-data}); the
caret marks the last prompt position, which is what the probe reads, and the
axis beside it names the seven depths at which it reads. The strip is a
schematic of \emph{fields}: cells are sized to their typeset labels and not to
their token counts, so it carries no horizontal scale, and the control cell's
sentence, which runs to hundreds of tokens, is not drawn hundreds of times
wider. It is the same object, at the same cell widths and the same cell height,
as the strip in \cref{fig:hook}. Tick spacing on the layer axis is the one
thing here drawn to scale, and the scale is layer index rather than accuracy
--- the grid is even apart from the layer added at $9$. \emph{(b)}~The
partition this chapter opens with. Accent marks the branch that survives,
because it is what \cref{sec:res-used} and \cref{sec:res-mechanism} spend the
rest of the chapter measuring; the struck branch and the counterfactual column
are quiet. \emph{(c)}~and~\emph{(d)} are two equal columns and are meant to be
read across, entry by entry. \textbf{No accuracy is drawn or printed here.} The
probe reads $0.821$ at layer $9$ and $0.758$ at layer $16$ against a floor of
$1/12 \approx 0.083$, and $0.883$ at layer $16$ once cell line, plate and dose
are projected out; those four are \cref{tab:workspace}'s and
\cref{fig:mechanism}(a)'s, and that panel is the only place a reader can see
that the context-removed series \emph{exceeds} the raw one at layers $12$ and
$16$ rather than sitting under it. What may not be read across: nothing on the
page is an estimate, none of it carries an interval --- the probe returns a
cross-validated accuracy and \texttt{workspace\_probe.json} stores no
dispersion for it --- and the two licence columns quantify nothing about the
size of a licence.}
\label{fig:licensing}
\end{fullwidth}
\end{figure}

Forward passes only. Real tier-2 evaluation prompts are grouped by drug; for
each, the residual-stream activation is captured at the last prompt position,
where generation begins, at every layer. A cross-validated multinomial
logistic-regression classifier learns which of the twelve drugs the prompt named,
against an analytic floor of $1/12 \approx 8\%$, the file storing no
permutation null. The twelve
are a random draw from tier-2 drugs with at least forty scored cells each; that
count sets the $8\%$ floor and, in \cref{sec:res-used}, bounds the drug subspace
at rank $11$. Folds are stratified over prompts and not grouped by well, a
limitation stated and not repaired, and the reason the accuracy is read as a
presence test, not a generalisation estimate.

\paragraph{The negative never arrives} The probe reads $82\%$ at layer 9 and
$76\%$ at layer 16, the deepest measured, against the $\approx 8\%$ floor. The
between-drug to within-drug variance ratio collapses about three-fold across the
final layers, so the drug is present throughout but geometrically de-emphasised
in the layers nearest generation. That de-emphasis is in the raw geometry; once
cell line, plate and dose are projected out, the series rises with depth instead, to
$0.883$ ($88\%$) at the deepest layer. The survival carries the weight, not the
peak, and that is the end of what a probe can say.

\begin{claim}
The drug is not missing from the computation. Its identity is linearly readable
from the residual stream at every layer measured, at the position the first
target token is emitted from. The raw probe reads $82\%$ at its layer-9 peak
and $76\%$ at the deepest against an $8\%$ chance floor; with cell line, plate
and dose projected out it rises with depth instead, to $88\%$
(\cref{sec:res-used}). The prompt names the drug, so this is a floor test and its
positive is close to free. It rules out that the drug field is lost before the
decoder runs, an encoding failure in the activations. It does not rule out that
the decoder ignores it. (\cref{sec:methods-residual} later uses \emph{target
encoding} for a different thing, how the target string is written, and nothing
here bears on that.)
\end{claim}

% ===========================================================================
\beatsection{E}{3}{The drug region is used, at a bit player's weight}{sec:res-used}

\begin{question}[3]
Knowing the drug is written somewhere says nothing about whether the network
reads it. Intervening on the region requires localising it, purifying it of the
batch structure it is confounded with, and proving it really holds the drug
before any null about it can be believed. Generation is sensitive to removing the
region, at a bit player's weight beside cell line.
\end{question}

\noindent
The second instrument is built in four steps, each because the unguarded version
of the same measurement misled us once already.

\paragraph{Step one, where the drug lives} On \num{480} tier-2 prompts (twelve
drugs, forty real cells each) at seven layers, the twelve per-drug mean
activation vectors are centred on the global mean and stacked, and their singular
value decomposition contributes its top ten orthonormal directions,
$V_{\textrm{drug}}$. Twelve centred class means span at most eleven directions,
so the slab keeps essentially the whole span, noisiest included, which for a
removal claim is conservative. The layer grid is even (2, 4, 6, 8, 12, 16) plus
layer 9, added after the raw probe peaked there. A cell-line subspace is built
identically as the positive control.\seealso{the grouping ladder in
\cref{sec:res-blind} is what makes cell line usable as a positive control}

\paragraph{Step two, purification} The raw drug slab is partly a batch slab.
Drugs are confounded with plate and dose by the atlas's design, and probing dose
from inside the drug subspace recovers it at $0.90$--$0.95$ at every layer. The
drug axis is therefore orthogonalised against a context subspace of cell line,
plate and dose directions, giving $V_{\textrm{pure}}$. Every causal claim below
is made on $V_{\textrm{pure}}$, never on the raw slab.

\paragraph{Step three, the removal gate} Before any ``ablating the drug changes
nothing'' result can be believed, the instrument must prove the slab actually
contains the drug. Accuracy is measured before and after projecting the drug
subspace out, and the fraction of above-chance signal removed must exceed $0.8$;
below that the null is unfalsifiable, a measurement of nothing. The gate passes
at every layer. Accuracy falls from $0.41$--$0.82$ to a constant $0.121$ floor,
removing $88$--$95\%$ of the above-chance signal, so the nulls below are real
properties of the model and not broken hooks.

% ===========================================================================
\begin{figure}[htbp]
\begin{fullwidth}
\centering
% ===========================================================================
% figs/purify.tex -- fig:purify, how the drug slab is BUILT, and the gate it
% has to pass before anything may be ablated with it.
%
% FIGURE BODY ONLY. No preamble, no document environment, no float, no caption.
% \input it from inside a figure/fullwidth float in
% Sections/04b-representation.tex; the ready-to-paste float sits commented at
% the foot of this file and is also staged in figs/_PENDING_FLOATS.tex.
% Compiles against thesis_v2/preamble.tex and nothing else: no \includegraphics,
% no external data, no TikZ library beyond the ones preamble.tex already loads
% (this file uses arrows.meta and calc only).
%
% WHAT THE FIGURE IS
% ------------------
% Steps one to three of sec:res-used, drawn as steps. Step one localises the
% drug: twelve centred class means, ten kept directions. Step two purifies it:
% the raw slab is orthogonalised off the context it is confounded with. Step
% three is the falsifiability check: projecting the slab out must destroy most
% of the decodable drug signal, or every null downstream is a measurement of
% nothing. fig:hook (a separate float, six paragraphs later) draws step four;
% the two are deliberately NOT merged -- each is about 100mm tall, the corpus
% ceiling is slab.tex at 117mm, and the house caption shape carries one bolded
% finding, not two.
%
% EVERY NUMBER DRAWN, AND WHERE IT COMES FROM
% -------------------------------------------
% Source for every measured quantity: RESULTS_cluster/workspace_probe.json.
% Config there: tier tier2_unseen_drugs, n_drugs 12, n_per_drug 40, n_dims 10,
% layers 2,4,6,8,9,12,16, chance 0.08333333333333333, seed 42.
% Sibling ledger with the full machine-precision values: figs/purify.json.
%
% (a) THE SEVEN PLOTTED DOTS -- the only marks in this figure with a scale.
%     layers.<L>.frac_signal_removed, layers 2,4,6,8,9,12,16:
%       0.8846153846153847  0.9253112033195021  0.9383561643835617
%       0.9461077844311377  0.9491525423728814  0.9440993788819876
%       0.9444444444444445
%     layers.<L>.gate_pass is true at all seven; the dots' position above the
%     0.8 rule is that fact, so no second mark asserts it.
%
% (b) THE GATE, 0.8. PROSE ONLY. Sections/04b-representation.tex:185 "the
%     fraction of above-chance signal removed must exceed $0.8$". There is no
%     config field for it in workspace_probe.json, so the ledger cites the .tex
%     line and not a key path. Drawn dashed accent, the house register for a
%     pre-set threshold a value is tested against (gate.tex:170-174).
%     SUPERSEDED IN v10 BELOW: it is now black, solid, 0.9pt.
%
% (c) THE THREE PRINTED ACCURACIES in the right-hand block of band two.
%       before  0.41--0.82   layers.2.drug_decode  = 0.4083333333333334
%                            layers.9.drug_decode  = 0.8208333333333334
%       after   0.121        layers.<L>.drug_decode_ablated
%                            = 0.12083333333333332, byte-identical at all seven
%       chance  0.083        top-level chance = 0.08333333333333333
%     Thesis 04b:186-188 "Accuracy falls from $0.41$--$0.82$ to a constant
%     $0.121$ floor, removing $88$--$95\%$ of the above-chance signal".
%     SUPERSEDED IN v10 BELOW as to PLACEMENT only: the three values are no
%     longer a right-hand gutter block, they are the figure's foot note. No
%     value changed.
%
% (d) THE IDENTITY, verified numerically before drawing, in python, over all
%     seven layers, exact float equality (not a tolerance):
%       (drug_decode - drug_decode_ablated) / (drug_decode - chance)
%         == frac_signal_removed
%     Layer 2: (0.4083333333333334 - 0.12083333333333332)
%              / (0.4083333333333334 - 0.08333333333333333)
%              = 0.8846153846153847, bit for bit.
%     WHICH SLAB THE GATE IS MEASURED ON. This identity fixes it, and the
%     answer is NOT V_pure. workspace_probe.json stores exactly one ablated
%     series, drug_decode_ablated, and the identity reproduces from the RAW
%     decode. The section-4.5 arm files prove the pipeline can distinguish the
%     two cases -- they store drug_decode_ablated AND drug_decode_pure_ablated,
%     frac_signal_removed AND frac_signal_removed_pure -- and this file stores
%     only the first of each pair. So the figure, the caption and the ledger all
%     say "the slab that is projected out" and nowhere claim the gate was run on
%     V_pure. Stage two's arrow lands on V_pure because that is where the CAUSAL
%     claims are made (04b:179-180); band two is a separate measurement and the
%     divider is what keeps the two from being read as one.
%     SUPERSEDED IN v10 BELOW as to the LAST CLAUSE ONLY: the divider is gone
%     and the panel letters do that work now. The caveat itself survives, in the
%     foot note.
%
% (e) THE FOUR PRINTED NUMBERS IN STAGE TWO, the only measured quantities above
%     the divider.
%       0.90--0.95   layers.<L>.confound_crossdecode.dose, seven layers:
%                    0.9020833333333333 0.9229166666666666 0.9416666666666667
%                    0.9395833333333332 0.95 0.8979166666666666
%                    0.8979166666666668. Thesis 04b:176-178.
%                    PRINTED AS TEXT AND NOT DRAWN, and there is deliberately NO
%                    0.90 reference rule anywhere in this figure: the stored
%                    minimum is 0.8979166666666666 at layers 12 and 16 and
%                    reaches the thesis's "0.90" only by rounding to two places.
%                    A rule at 0.90 would put two of seven points below it and
%                    would read as a failure the thesis does not report.
%       raw 0.620    layers.9.variance_share.drug      = 0.6201207831136324
%       purified     layers.9.variance_share.drug_pure = 0.030620247335072275
%         0.031
%     The RAW share is quoted in no sentence of the thesis; figs/fig-mechanism.json
%     already records it as drug_raw_not_drawn. Printing it here is a new reading
%     of a stored field, flagged as such in purify.json, on the precedent of
%     figs/slab.json's handling of variance_share.context. Layer 9 is named in
%     the figure text so the pair cannot be read as a seven-layer summary; it is
%     the layer at which the raw probe peaks (04b:170-171), which is why the
%     section uses it as its own reference depth.
%
% (f) THE COUNTS in stage one: 480 prompts, twelve drugs, forty cells each,
%     seven layers, ten kept directions, span at most eleven.
%     config.n_drugs 12, config.n_per_drug 40, config.n_dims 10,
%     layers.<L>.pure_drug_dims 10, config.layers "2,4,6,8,9,12,16".
%     480 = 12 x 40 and the rank-11 bound (twelve centred means span at most
%     eleven directions) are ARITHMETIC, not stored fields; the ledger says so.
%     Thesis 04b:164-171.
%
% AGREEMENT WITH THE TEXT (checked line by line, not assumed):
%   04b:164-165  "\num{480} tier-2 prompts (twelve drugs, forty real cells
%                each) at seven layers"            -> stage one, line 4.     OK
%   04b:168-170  "Twelve centred class means span at most eleven directions,
%                so the slab keeps essentially the whole span, noisiest
%                included, which for a removal claim is conservative"
%                -> stage one, lines 1-3, verbatim.                          OK
%   04b:175-178  "The raw drug slab is partly a batch slab. Drugs are
%                confounded with plate and dose by the atlas's design, and
%                probing dose from inside the drug subspace recovers it at
%                $0.90$--$0.95$ at every layer"    -> stage two, lines 1-3.  OK
%   04b:179-180  "Every causal claim below is made on $V_{\textrm{pure}}$,
%                never on the raw slab."           -> the V_pure callout.    OK
%   04b:185-186  "the fraction of above-chance signal removed must exceed
%                $0.8$; below that the null is unfalsifiable, a measurement
%                of nothing"                       -> the dashed rule, and
%                                                     G4's two italic lines,
%                                                     which are the section's
%                                                     own words.             OK
%   04b:186-188  "Accuracy falls from $0.41$--$0.82$ to a constant $0.121$
%                floor, removing $88$--$95\%$"     -> right-hand block and
%                                                     the set label.         OK
%   04b:212-213  tab:workspace caption "All seven layers pass the removal gate
%                (88--95\% ...); chance is $0.083$"                          OK
%
% AVAILABLE, DELIBERATELY NOT DRAWN, recorded here and in purify.json so a
% reviewer who asks for direct evidence of the batch overlap is answered
% without a redraw:
%   * layers.<L>.overlaps.context   0.5215396547181934 -- 0.6323768458353303
%   * layers.<L>.overlaps.cell_line 0.3043896441494730 -- 0.3534060289850234
%     These are the per-layer measure of "the raw drug slab is partly a batch
%     slab" and are quoted nowhere in the thesis. Drawing them would put a
%     second measured scale into a panel whose whole discipline line says the
%     geometry has no scale, and would be a new claim rather than a redraw.
%   * layers.<L>.confound_crossdecode.plate, 0.4541666666666667 --
%     0.5770833333333333, a materially weaker and entirely undiscussed
%     confound. Not printed, because the section names dose and not plate.
%   * layers.<L>.swap.{drug_b_kl, random_inject_kl} in this same JSON is the
%     RETIRED steering design (07-Appendix.tex, sec:app-identity). Nothing from
%     it is on this figure. figs/mechanism.py's header records that the v1
%     figure plotted exactly that pair while its caption described the ablation.
%
% DESIGN DECISIONS, each with its reason
% --------------------------------------
%  1. TWO BANDS AND A DIVIDER THAT IS A STATEMENT. Above the rule at y=4.72
%     nothing has a scale; below it everything does. No mark crosses the rule,
%     no leader connects a schematic object to a plotted one, and the italic
%     sentence under the rule sets band two up as the answer to the objection
%     stage two raises -- if purification strips the raw slab from 0.620 to
%     0.031 of the stream's variance, does anything survive in it? The
%     alternative, one continuous diagram from cloud to axis, was rejected: it
%     would invite the reader to measure the ellipse.
%     SUPERSEDED IN v10 BELOW. The divider and its sentence are deleted; the
%     schematic/measured distinction is now carried by the MARK (open rulegrey
%     ring against filled accent disc), which is where it belongs.
%  2. THE FAIL REGION IS DRAWN, THE PASS REGION IS NOT. The one non-obvious
%     mark in band two is the rulegrey!18 rectangle from 0 to 0.8. It is the
%     region of outcomes that would have made every null in this section a
%     measurement of nothing, and it is empty. A figure that drew only the seven
%     dots and the rule would be a bar of good news; drawing the region the
%     dots had to avoid is what makes it a falsifiability picture. Its emptiness
%     is also where the two italic lines fit, which is not a coincidence and is
%     the reason they are placed there rather than in the caption.
%     PARTLY SUPERSEDED IN v10: the region survives, at rulegrey!12, and the two
%     italic lines inside it do not. The emptiness is the message and text in it
%     was competing with the marks.
%  3. NO CONNECTING LINE THROUGH THE SEVEN DOTS. The series is flat to 0.1355cm
%     on this scale (0.8846 to 0.9492 at 2.10cm per unit of fraction). A line
%     would assert a trend across a 6-point spread the axis cannot resolve; the
%     seven marks are named as one ordered set by "all seven layers: 88--95%"
%     instead, which is comparators.tex's rule for marks too close to leader
%     individually.
%  4. NO INTERVALS ANYWHERE. frac_signal_removed is a single ratio of two
%     accuracies on 480 prompts and carries no stored dispersion of any kind;
%     the KL series in the same file carry normal-approximation SEs, which are
%     not this document's interval convention (figs/fig-mechanism.json records
%     intervals_drawn false with that reason). A bare solid dot with no bar is
%     how this document says a quantity has no interval (comparators.tex:149-151
%     changed a hollow ring to a solid dot for exactly this).
%  5. THE ORTHOGONALISATION IS DRAWN AS AN OPERATION, NOT AS A RESULT, AND THE
%     OPERATION IS A PROJECTION AND NOT A ROTATION. This is the one mark
%     figs/slab.tex does not have. There the drug wedge and the
%     context wedge are simply shown already orthogonal, which states the
%     outcome and hides the step; here the raw band sits at 38 degrees, visibly
%     crossing the context band, and the raw slab's tip is carried onto V_pure.
%     REPAIRED, PASS 9. It used to be carried by an accent ARC from 38 to 76
%     degrees, i.e. the raw band swung upright. A reviewer read that panel as
%     "the slab is turned until it is perpendicular to context" and had to go to
%     the prose to learn that a component is REMOVED, and they were right to:
%     two bands of equal width (0.32cm) and near-equal length (1.90 against
%     2.00) differing only in angle say "same slab, rotated, nothing lost",
%     against measured numbers 1.8cm to the right that say raw 0.620, purified
%     0.031, a twentyfold shrink. A curved arrow is the universal sign for
%     rotation. What replaces it is the subtraction itself: a dashed rulegrey
%     perpendicular from the raw band's tip down to the context band, the
%     component it lands on drawn as its own named mark, and a STRAIGHT accent
%     arrow at the tip's own height carrying the tip left by exactly that
%     component, landing on V_pure. Two further faults the arc had, both from
%     the render: its head stopped 1.5mm short of the V_pure bar and pointed at
%     white space, and the grey leader from the word "orthogonalise" stopped 4mm
%     short of the arc's tail -- the corpus puts the 0.06cm gap at the LABEL
%     end, never at the mark end, so that stub read as a stray hairline. The
%     label now sits 0.15cm from the arrow's tail and carries no leader, and the
%     arrow is 0.6pt: a schematic operation, not a measurement, where 1.0pt is
%     the corpus's "this IS the claim" register. V_pure was also narrowed from
%     0.32 to 0.18cm, so the purified band is visibly thinner than the raw one.
%     slab.tex is left alone and is not
%     superseded: three of its four panels are the failed identity designs of
%     sec:app-identity, which are the appendix's argument, and its magnified
%     interior is sec:res-mechanism's swap geometry. What it draws that belongs
%     to sec:res-used is one thing, the variance-share disc, and this figure
%     does not redraw it. The two figures share vocabulary (accent-on-white
%     bands, rulegrey leaders, one discipline line saying what is to scale) and
%     no coordinate.
%     The whole of this decision survives v10 except the arrow's WEIGHT, which
%     goes 0.6 -> 0.5pt to bring the file inside a four-weight budget; the
%     perpendicular goes dashed-rulegrey -> densely-dotted-quietink, which is
%     this document's reserved stroke for a projection line.
%  6. THE CONFOUND IS A NAMED MARK, NOT THE PRODUCT OF TWO TINTS.
%     REPAIRED, PASS 9, and this is the repair both reviewers called blocking.
%     It used to be the lozenge where the raw band's tint multiplied over the
%     context band's, with the whole claim carried by the caption ("that
%     darkened region is the confound step two exists to remove"). Three faults:
%     (a) it was the only element in the panel with no direct label, while every
%     other mark had one; (b) V_pure ran vertically through its centre and split
%     it in two -- sampled on the render, a row across it read grey (174,174,174),
%     then (119,143,164) for the tint product, then a hard cut to opaque accent
%     (11,82,147) for 24px, then (119,143,164) again -- so it was not a nameable
%     region at all; (c) it made a CROSSING mean "confounded" on a panel where
%     another crossing, V_pure through the same band, has to mean "orthogonal".
%     One glyph cannot carry both. It is now the orthogonal projection of the
%     raw direction onto context: a segment along the context band from the
%     shared origin to the foot of the perpendicular dropped from the raw band's
%     tip, drawn in the raw slab's own tint at a higher opacity with a 0.35pt
%     accent!70 outline, and directly labelled "the confound". It begins at
%     V_pure's right edge, so nothing bisects it, and it is exactly the length
%     of the removal arrow above it, which is the rhyme that makes the arrow
%     legible. With the tint-multiply claim gone, a crossing means nothing
%     anywhere in the panel and the right-angle mark is the only assertion of
%     orthogonality, which is what decision 8 already wanted.
%     The context band is now panelbg with a rulegrey outline and not ink!35:
%     see decision 11.
%  7. NO 0.90 REFERENCE RULE. See (e) above: the stored dose minimum is
%     0.8979166666666666 and rounds to the thesis's 0.90. Printing the range as
%     text states exactly what the section states; a rule would show two of
%     seven points below it and would manufacture a failure.
%  8. RIGHT ANGLE AS A MARK, NOT AS A WORD. The small square in the
%     upper-left quadrant at the shared origin is the only assertion that
%     V_pure and the context subspace share no direction. The word
%     "orthogonalise" labels the operation; the square labels its result.
%     PASS 9: ink 0.35pt, up from rulegrey 0.3pt. The caption calls it "the one
%     geometric fact true by construction" and it was the faintest mark in the
%     panel.
%     AMENDED IN v10: the square moves to the LOWER-left quadrant, doubles to
%     0.40cm arms, and acquires the direct label it never had.
%     SUPERSEDED IN v11: THIS LINE IS FALSE ABOUT ITS OWN DRAWING AND ALWAYS
%     WAS. v10 drew the square in the UPPER-left quadrant of the crossing --
%     (10.01,8.65)-(10.01,9.05)-(10.41,9.05) -- and no build has ever put it
%     lower-left. The direct label and the enlargement are true. v11 keeps the
%     upper-left quadrant, which is the only one of the four that is free (the
%     raw band at 38 degrees crosses lower-left and upper-right, and "the
%     confound" and its leader own lower-right), and takes the arms 0.40 ->
%     0.28cm to clear the removal arrow -- see v11 item 2.
%  9. THE GLOBAL MEAN IS A CROSS, NOT A DOT. The twelve class means are dots;
%     the thing they are centred on must not be a thirteenth dot. It is also
%     leadered out of the cloud rather than labelled in place, because at
%     8pt the label box collided with two of the twelve dots.
% 10. V_drug IS DRAWN WITH STAGE TWO'S OWN GLYPH, AND ALL TWELVE DISPLACEMENTS
%     ARE DRAWN. REPAIRED, PASS 9, two findings at once.
%     (a) V_drug used to be a dashed, accent!10 FILLED ELLIPSE enclosing the
%     twelve dots, while stage two drew the same name as an unbounded band. A
%     bounded blob around the class means says "the region where the drugs
%     live"; a subspace is not a region and has no edge. Read cold, the ellipse
%     was taken for the data cloud and stage two's band for a new object, so
%     stage one's output was never connected to stage two's input -- the one
%     join the figure exists to make, and the join the brief says readers lose.
%     V_drug is now an accent band at the SAME fill opacity stage two uses,
%     through the global mean, running past the data on both sides. The ellipse
%     is deleted rather than kept unlabelled: with the band tall enough to
%     contain all twelve dots (half-height 0.52 against a maximum |dy| of 0.44)
%     it had nothing left to say, and a class mean drawn OUTSIDE V_drug would
%     have been false, since the twelve centred means span it by construction.
%     (b) ONE displacement was drawn and eleven were not, so "twelve are
%     stacked" -- the sentence directly under the dots -- was invisible, and the
%     dots were unlabelled while the prose two lines below names a different
%     population ("480 prompts: twelve drugs, forty real cells each"), so the
%     scatter read as prompts and the cross as their grand mean. All twelve
%     stems are now drawn at 0.3pt quietink, with the one accent arrow on top
%     and still the only labelled one, and the dots carry the tag "the twelve
%     per-drug means". The stems make the count a picture and they are what the
%     "at most eleven directions" line is about.
%     The twelve offsets are written out in a \foreach list so no build can
%     randomise them and no reader can be told a different story by a different
%     compile. Opacity: the band is 0.20, not the 0.34 of stage two's first
%     draft, and stage two's raw band was brought to 0.20 with it. At 0.34 over
%     5.10 x 1.04cm the band became the single heaviest area on the page, which
%     is the thing this figure was already being marked down for; at 0.20 over
%     4.10 x 1.04 it lays down about the ink of stage two's 1.90 x 0.32 band and
%     the two are drawn at the SAME opacity, so the glyph rhymes exactly.
%     AMENDED IN v10: the twelve class means become OPEN RINGS and the one
%     accent displacement arrow loses its head -- see v10 items 3 and 6.
% 11. ONE GREY FAMILY, AND THE CONTEXT SUBSPACE IS THE COLOUR slab.tex GIVES IT.
%     PASS 9. Three greys were in play, derived two different ways: rules at
%     rulegrey, the context band at ink!35, the gate region at rulegrey!18.
%     ink!35 is not derived from rulegrey and is LIGHTER than a hairline rule,
%     so the band read as a heavy solid slab while the rules read as thin, and
%     it was the darkest region on either page. slab.tex draws this same named
%     subspace in panelbg (its \slabdisc, slab.tex:83-90), so the figure also
%     disagreed with the published figure about what colour the object is. The
%     band is now panelbg with a 0.35pt rulegrey outline; ink!35 no longer
%     appears in the document, and purify's context band and slab.tex's context
%     wedge are the same object drawn twice.
% 12. THE GATE PANEL KEEPS ITS FULL 0-TO-1 AXIS. A reviewer asked for the range
%     to be cut to 0.75-1.00 so the seven fractions (0.8846 to 0.9492, 0.1355cm
%     apart on this scale) would resolve, against a fail region 12.4x their
%     height. OVERRULED, and the reason is the same one that bans bars on a log
%     axis: the fail region's height is not a drawing choice, it is the share of
%     the outcome space that would have made every null below unfalsifiable, and
%     0 is the natural floor of a fraction, not an arbitrary cut. Cutting the
%     axis at 0.75 would shrink the one region whose SIZE is the argument, to
%     resolve an ordering the thesis makes no claim about -- tab:workspace is
%     the authority for the per-layer values and the section quotes only the
%     range. What was applied instead: the direct label "all seven layers:
%     88--95%" names the set (decision 3), and the 0.8 numeral was dropped from
%     the tick set because the accent label already owns that coordinate.
%     STILL IN FORCE IN v10, and it is the first thing v10 was told to protect.
%
% WHAT SIX RENDER-AND-READ PASSES CHANGED, with the measurement that forced it
% ---------------------------------------------------------------------------
%  i.   PASS 1, Overfull 88.29pt. The discipline line under stage two ran 146
%       characters at \footnotesize and measured 20.4cm against a 17.30cm
%       measure. Cut to 103 characters; the sentence now states only what it
%       has to, that the geometry has no scale and four numbers do.
%  ii.  PASS 1. "every causal claim is made here", set anchor=south at y=9.22,
%       put its box top at 9.63 against the stage-two subtitle's ink bottom of
%       9.37 and the two overprinted. There is no room above V_pure for a
%       second line, so the sentence became the fourth line of stage two's
%       prose block, in the section's own words ("never on the raw slab").
%  iii. PASS 1. "the global mean" and "$\mu_d-\bar\mu$, one drug's" collided
%       character to character: both wanted the space up and right of the
%       cloud, which is where the displacement arrow points. The global mean
%       was moved BELOW the cloud on a leader. Drawn straight down that leader
%       read as a dropped axis, so PASS 4 gave it a 0.25cm dogleg; the dogleg
%       clears the two nearest class means by 0.23cm and 0.25cm.
%  iv.  PASS 2, Overfull 6.20pt, unchanged from PASS 1 by three edits, which is
%       what identified it: the culprit was not the widest LINE but the
%       fraction axis's own tick label. "1.0" is anchor=east on a spine at
%       x=0.30, and its box ran 0.28cm LEFT of x=0. The bare \path pins a
%       minimum box, not a maximum, so nothing clipped it. The whole plot was
%       shifted right (spine 0.30 -> 0.75, layer 2 from 0.60 -> 1.05) rather
%       than the label being clipped or shrunk. boxes=0 from PASS 4 on.
%  v.   PASS 2. Band two carried two annotation columns 0.30cm apart at the
%       same baseline, so "88--95%" and "accuracy before" read as one line. The
%       plot was widened from 8.60 to 11.00cm of layer axis and the two columns
%       became one, which also removed 3.4 x 1.9cm of dead white under them.
%  vi.  PASS 4, at 520dpi. The leader to "orthogonalise" left the arc 0.20cm
%       from its own arrowhead and read as a grey continuation of the accent
%       arrow. It now leaves at 50 degrees, 0.42cm clear of the head, and stays
%       above the raw band's upper edge for its whole length. The right-angle
%       mark went from 0.16 to 0.20cm, and the global-mean cross from 0.30 to
%       0.45pt with 2mm arms, because at 0.3pt it was fainter than the leaders.
%  vii. PASS 5. Stage one left 4.2 x 1.9cm of white to the right of the cloud
%       while stage two's matching region was full. The ellipse went from
%       1.75 to 2.05 x-radius and the cloud from 1.45 to 1.72, which both fills
%       the panel and raises the aspect from 2.34:1 to 2.77:1. The ellipse's
%       y-radius was NOT raised with it: at 0.98 the ellipse reached y=6.92 and
%       left the global-mean label 0.06cm from the prose block below.
%  viii.PASS 5. The panel-three subtitle said "projecting the slab out", which
%       after stage two's "every causal claim is made on V_pure" invited the
%       reading that the gate was run on V_pure. It now says "projecting the
%       drug subspace out", the section's own noun at 04b:184-185, and the
%       string V_pure does not appear in panel three at all.
%
% LAYOUT NOTES, all forced by a build
% -----------------------------------
%   * fullwidth = textwidth 142mm + marginparwidth 28mm + marginparsep 4mm
%     = 174mm. (The "171mm" beside \newenvironment{fullwidth} in preamble.tex
%     is stale arithmetic.) The bounding box is pinned with a bare \path at
%     (0,-0.12) rectangle (17.30,10.20) = 173.0 x 103.2mm, inside the 17.35cm
%     hard edge circularity.tex records and under the corpus height ceiling,
%     slab.tex at 117mm. The harness reports pages=1, boxes=0. It grew 1.7mm in
%     pass 9: band two moved down 0.20cm to open a line under stage three's
%     subtitle for the "which slab" caveat, and the "four ... see fig:hook"
%     pointer sits on the bottom line. The float got much SHORTER overall in the
%     same pass, because the caption went from 484 words to about 280.
%   * x=1cm, y=1cm. Nothing is scaled, nothing is \resizebox'd; 1cm here is
%     1cm on the page.
%   * pgfmath: \lx and \fy expand to BRACED groups, so \lx{2}+0.05 is not a
%     coordinate. Every stand-off is either inside the braces or an xshift on
%     the node (gate.tex records the same trap).
%   * the tick rows use the \v/\t value/label pair form, so the fraction axis
%     prints 0.5 and 1.0 rather than \foreach's raw 0.5 and 1 tokens. The layer
%     row does not need it: its values are integers and print as themselves.
%   * no node uses align= or text width=: such a node builds a minipage whose
%     first-line height follows the AMBIENT \baselineskip, so its ink moves
%     with whatever prose precedes the float (frames.tex measured 3.7pt of
%     drift). Every multi-line block here is a stack of single
%     anchor=base west lines on explicit baselines: 0.28cm apart in band one's
%     two prose blocks, 0.33cm in band two's right-hand block, where the lines
%     are read one at a time rather than as a paragraph.
%   * the fraction axis carries no rotated title. The panel subtitle sits
%     directly above it and names the quantity, the direction and the frame in
%     one line, which is what gate.tex's fix 1 asks an axis to do; a second
%     title would repeat it 0.2cm away.
%     THIS NOTE IS NOW FALSE AND IS SUPERSEDED BY v10 ITEM 4. It was the
%     supervisor's first complaint about this arc of figures and it was wrong:
%     a quantity named in a panel subtitle 30mm from the ticks is not named at
%     the axis. The fraction axis carries a rotated title from v10 on.
%   * PASS 9, THE "WHICH SLAB" CAVEAT, three builds to place. The wording went
%     155 -> 121 -> 92 characters. At 155 the picture was 55.72pt over the
%     measure and the harness said so; at 121 it ran to x = 15.90 and
%     overprinted the accent label "all seven layers: 88--95%" at 12.60; at 92
%     it ends at 12.30. Vertically it was first set at baseline 3.24 and its
%     ascenders touched the subtitle's descenders with no white between -- the
%     subtitle is \footnotesize in a node whose TOP is pinned at 3.76, so its
%     descender line is 3.351 and a \scriptsize line at 3.24 reaches 3.444. It
%     is now at baseline 3.05 with 0.10cm of clearance above, and INDENTED to
%     x = 0.90: at x = 0 its descenders at 2.97 met the "1.0" tick numeral,
%     whose box top is 2.95, and the indent both clears the numeral and breaks
%     the note away from the subtitle so the two do not read as one paragraph.
%   * PASS 9, THE DIVIDER BETWEEN STAGES ONE AND TWO now ends at y = 5.28, the
%     descender line of the last prose row in both columns (baseline 5.38),
%     instead of at 5.40 where it stopped between two lines of type and read as
%     a broken rule. It clears the centred italic beneath it by 0.26cm.
%   * PASS 9, THE DISCIPLINE LINE IS CENTRED on the figure's own centre and is
%     no longer flush left at the prose column's x. Flush left, at the same
%     leading as the four-line block directly above it and with no blank
%     scan-line between, it parsed as that paragraph's fifth line even though it
%     ran full width past the vertical divider. Centred, it pairs with the
%     sentence centred under the rule, and the two bracket the divider.
%   * PASS 9, THE AXIS TITLE moved from x = 6.60 to 7.73. \lx{9} = 6.55, so
%     centred on 6.60 the word "layer" sat within 0.05cm of the layer-9 numeral,
%     one line below it, and read as a label of layer 9 rather than of the axis.
%     7.73 is the midpoint of the widest tick gap (\lx{9} = 6.55 to
%     \lx{12} = 8.91) and is 0.83cm clear of both.
%   * PASS 9, THE 0.8 TICK IS GONE. The coordinate was labelled twice at the
%     same height, once as a quietink numeral on the left spine and once as the
%     accent label "the gate: 0.8" on the right, so the plot appeared to carry
%     a left scale and a right scale. 0, 0.5, 1.0 is also a regular tick set.
%     A 0.9 tick was considered and rejected by measurement: at \fy it would
%     sit 0.21cm from both neighbours against a 0.203cm numeral height.
%
% TYPE. Ambient \small on the picture; \scriptsize (8pt) on every node except
% the \footnotesize (10pt) italic prose -- the three panel subtitles, the
% discipline line, the divider sentence and the two lines inside the fail
% region. Nothing is set at body size.
% SUPERSEDED IN v10 BELOW: two sizes only, \footnotesize and \scriptsize, and
% ZERO italic.
%
% ===========================================================================
% v10 -- 2026-08-18. REBUILD AGAINST THE STYLE CONTRACT AND THE SUPERVISOR'S
% VERDICT. Nine of the twelve design decisions above survive whole; the three
% that do not are superseded in place, above, rather than deleted, because the
% thesis is audited against this header.
%
% THE VERDICT THIS PASS ANSWERS, quoted: "For the other figures 4.8-12 go back
% over them and iron them out even more. Make it look more professional, right
% now it looks a bit sloppy and childish ... and moreover make sure that the
% figures are CORRECT and display the idea we want to convey in the best way
% possible." Plus, on 4.7 but binding on the whole arc: "there is no definition
% on the x axis for what we are measuring".
%
% NOT ONE PLOTTED OR PRINTED VALUE CHANGED IN THIS PASS. The seven gate
% fractions are the same seven floats, unrounded; 0.8, 0.41--0.82, 0.121,
% 0.083, 0.90--0.95, 0.620, 0.031, 480, twelve, forty, ten, eleven and seven
% are all unmoved. What changed is the y map (0.75 + f*2.10 -> 1.45 + f*3.95)
% and the x map (1.05 + (L-2)/14*11.00 -> 1.45 + (L-2)/14*11.75), both recorded
% in figs/purify.json beside the old ones. A map change is a position change,
% not a value change, and the ledger carries both maps so the audit can redo it.
% SUPERSEDED IN v11, AND THESE TWO SENTENCES CARRIED FOUR WRONG COEFFICIENTS.
% What v10 actually shipped, read off its own code and confirmed by measuring
% its render's dot centroids, was \fy = 1.45 + f*3.85 (not f*3.95) and
% \lx = 1.85 + (L-2)/14*12.00 (not 1.45 + (L-2)/14*11.75). figs/purify.json's
% value_to_position_maps.current_v10 held the correct pair all along, so the
% header and its own sibling ledger disagreed; under the header's numbers an
% auditor redoing the drawing would have put layer 16 at 13.20cm against the
% 13.85cm drawn and the gate at 4.61cm against 4.53cm. The v11 maps are
% recorded beside the code at the top of the picture, which is the only place
% they can be checked against what is drawn. No value changed in either pass.
%
%  1. THE THREE ALL-CAPS SANS ACCENT BANNERS ARE GONE, replaced by fig:mechanism's
%     own panel-head form: "(a) localise the slab", "(b) purify it against
%     context", "(c) the removal gate", \footnotesize roman ink, sentence case,
%     each left-aligned to its panel's x origin. fig:mechanism (4.11) carries
%     zero banners and sets "(a) encoded more, used no more" at 10.0pt roman
%     ink; figs/slab.tex (A.3) carries nine banners, but A.3 is an appendix
%     reference plate a reader stops at and this is read inside a running
%     argument, where the caps sans accent face belongs to the margin
%     apparatus (\marginlabel) and not to the figure. The ordinal now lives in
%     the panel letter, so the enumeration "one, two, three" that made a reader
%     hunt for a missing stage four is gone with it, and so is the
%     "four \quad what the ablation costs: \cref{fig:hook}" pointer that closed
%     it -- a cross-reference in a drawing renders as ?? in the harness and on
%     the page competed with the x-tick row it shared a baseline with, where
%     the word "four" read as a leftmost tick.
%  2. EVERY ITALIC IS GONE, AND WITH IT EVERY SENTENCE THAT WAS ARGUING RATHER
%     THAN NAMING. Deleted: the three \footnotesize italic panel subtitles;
%     the centred divider sentence "an ablation is only a test of the drug if
%     the slab it removes really holds the drug", which restates the body
%     paragraph directly above the float; the two lines inside the fail region,
%     "had any layer landed in here, the null below it / would have been
%     unfalsifiable --- a measurement of nothing", which restate the caption's
%     own "the shaded region ... is empty"; and the centred discipline line
%     "the dots and the bands are schematic and have no scale; only the four
%     numbers in stage two are measured". Italic went from 31% of the figure's
%     characters -- the highest in the five-figure set, against 0% in 4.11 --
%     to nil. All five sentences are staged in figs/_PENDING_CAPTIONS.tex under
%     \label{fig:purify} for a human to fold into the caption, which is where
%     this document keeps its argument; none is silently lost.
%  3. THE SCHEMATIC DISCLAIMER IS NOW CARRIED BY THE GLYPH AND NOT BY A
%     FOOTNOTE. The panel drew thirteen SCHEMATIC dots at the same radius and
%     the same fill as panel (c)'s seven MEASURED ones and then apologised for
%     the collision in an 8pt line. The twelve class means are now OPEN RINGS,
%     radius 1.4pt, white fill, 0.5pt rulegrey edge -- this document's mark for
%     a schematic, unscaled object -- and the seven gate fractions stay filled
%     accent discs at radius 1.4pt, which is its mark for a point estimate.
%     One look now separates them and the disclaimer line is deleted.
%  4. THE FRACTION AXIS IS DEFINED AT THE AXIS. This is the supervisor's first
%     complaint, answered mechanically. It carried bare numerals 1.0 / 0.5 / 0
%     with the quantity named only in an italic subtitle about 30mm above and
%     20mm left of the topmost tick, so a reader at the axis could not tell
%     whether 1.0 meant "all the signal removed" or "accuracy 1.0". It now
%     carries, rotated 90 degrees reading upward and immediately left of the
%     tick numerals, exactly fig:mechanism's form: line one \footnotesize roman
%     ink naming the measured quantity, line two \scriptsize quietink giving its
%     construction. The layer axis keeps the label "layer" it already had --
%     one of only two axes in 4.7-4.10 that already met 4.11's standard -- and
%     is now left-aligned to the axis's own left end rather than centred on the
%     widest tick gap.
%     MEASURED, and this is what fixed the wording. A rotated line cannot be
%     longer than the lane left of the tick column, and that lane is pinned at
%     both ends by things that may not move: below by the foot note's ascender
%     line at y = 0.52, above by the descender line of the prose at x = 0, which
%     after the 1.5mm lift of item 15 sits at y = 6.44. The axis's own midpoint
%     is 3.375, so the usable half-length is 2.855cm and a line may run 5.71cm.
%     Measured off this figure's own render, \footnotesize sets at 0.166
%     cm/character in this class and \scriptsize at 0.119, so the lane holds 34
%     \footnotesize or 47 \scriptsize characters.
%     "fraction of above-chance drug signal removed" is 43 characters, needs
%     7.14cm, and fits no layout under the 105mm height cap. What is set instead
%     is "fraction of drug signal removed" (31 characters, 5.15cm) over
%     "(before minus after)/(before minus chance)" (42 characters, 5.00cm). The
%     second line carries "above chance" in the only form that is auditable
%     against the JSON anyway, the construction itself: a reader standing at the
%     axis can now see that 1.0 means all of it, and can recompute the quantity
%     from the three accuracies in the foot note. The minus signs are spelled as
%     words and not set as $-$: in math mode CMSY10 came out at 8.3pt, a THIRD
%     type size above 7pt, which the style contract's first clause forbids.
%     SUPERSEDED IN v11, AND THIS PARAGRAPH IS THE ERROR v11 EXISTS TO CORRECT.
%     Three things in it are wrong. (i) THE LANE ARITHMETIC DISAGREES WITH ITS
%     OWN LEDGER AND WITH THE CODE: this block says lane 5.71cm at 0.166
%     cm/character with the axis midpoint at 3.375, while the inline comment at
%     the axis said 5.44cm at 0.173 with midpoint 3.24, and figs/purify.json
%     said 5.71/0.166/3.48. Remeasured on the v11 render from the PDF's own
%     text spans, \footnotesize sets at 0.161 cm/character in this class, and
%     the one figure that matters is now recorded once, beside the axis.
%     (ii) THE CONCLUSION DOES NOT FOLLOW. That a 43-character STRING does not
%     fit rules out the string, not the QUANTITY; dropping "above-chance" left
%     the \footnotesize line naming a quantity the marks are not, and that is a
%     correctness defect, not a compromise. The name is now set over two rotated
%     \footnotesize lines and the \scriptsize construction line, which S5 makes
%     optional, is deleted -- it had become the tallest object in the panel,
%     rising 6.0pt above the panel head's own top, so the smallest type in the
%     block outranked the head. The construction moved to the foot note.
%     (iii) THE $-$ NOTE IS RIGHT AND WAS APPLIED ONE CHARACTER SHORT: the v10
%     foot note still set "chance $0.083 = 1/12$", whose "=" came out of CMR10
%     at 8.3pt, a third size above 7pt, exactly the fault the note describes.
%     v11 sets "chance $0.083$, exactly $1/12$" and the render carries two
%     sizes above 7pt and no CMR10 at all.
%  5. THE DISPLAY FRACTION IS OUT OF THE PLOT. A \displaystyle fraction stacked
%     in a plot's label gutter was the most slide-like object in the five-figure
%     set. Its content is not lost: it IS the axis's second line now, which is
%     where a definition of the measured quantity belongs. The three accuracies
%     it explained (before 0.41--0.82, after 0.121 at every layer, chance 0.083)
%     move to a two-line \scriptsize foot note, and the six-line right gutter --
%     which mixed two direct labels with four notes at three different left
%     edges, 320.7, 402.9 and 428.4pt on the printed page -- is dissolved. What
%     stays beside the plot is exactly the two DIRECT labels: "all seven
%     layers: 88--95%" at the height of the dot cloud it names, and "gate, 0.8"
%     at the right end of the rule it names.
%  6. THE 0.8 GATE IS BLACK, SOLID, 0.9pt. It was accent 0.6pt dashed. This is
%     a correctness repair and not a preference: the same pre-registered class
%     of object -- a threshold a measurement is tested against -- was drawn
%     three different ways inside nine printed pages (accent dashed here,
%     accent dashed in fig:materiality, solid black in figs/family.pdf), while
%     fig:mechanism draws the chance floor solid black 0.9pt and style_v2.py's
%     chance_line() is hard-coded color='black', lw=0.9, solid, with the
%     docstring "Thin, solid, black. Never dashed, never omitted". Accent in
%     this document means the object under test; a threshold is not the object
%     under test, so a threshold may not be accent. Its label goes with it:
%     \scriptsize roman black, reading "gate, 0.8".
%  7. THE FAIL REGION IS LIGHTENED TO rulegrey!12 AND EMPTIED OF TEXT. It was
%     rulegrey!18 carrying two italic lines, it filled about 80% of the plotting
%     area, and a cold reader's eye went to it first and read it as the data
%     while the seven dots sat in a 12% band at the top. rulegrey!12 over white
%     renders at grey 243, above the 230 threshold the ink budget counts, so
%     the region now costs nothing against the 6.0% ceiling and still shows. Its
%     emptiness is the message; nothing is written inside it.
%  8. THE CONFOUND AND THE RIGHT ANGLE ARE LABELLED AT THE MARK. A cold reader
%     attached the label "the confound" to the diagonal band rather than to the
%     intersection patch, and did not see the right-angle glyph at all until
%     zooming: it was 2mm. The glyph is now 4mm on each arm and it carries the
%     direct label it never had.
%     THE LABEL IS ONE WORD, AND THAT WAS FORCED BY A RENDER. "orthogonal by
%     construction" is 3.6cm at \scriptsize, and panel (b) has no 3.6cm lane at
%     the glyph's own height that does not cross V_pure, the context band or the
%     removal arrow. Set anchor=base east at x = 9.71 it ran back to x = 6.20,
%     i.e. it sat in the gutter BETWEEN the two panels and a cold read could not
%     tell which panel owned it. Set base west on panel (b)'s x origin and
%     broken over two lines it ran to x = 10.20 and butted the V_pure label at
%     10.22, so the two read as one string, "orthogonal byV_pure" -- measured at
%     500dpi on the pass-4 render. What is drawn is the one word that has to be
%     AT the mark, set adjacent to the glyph's vertical arm with no leader.
%     "by construction" is the claim and not the name, so it goes to the
%     caption, and figs/purify.json records the move.
%     The confound's label acquires a 0.30pt leader to the patch. The dose
%     reading joins the two subspace fractions in one measurement strip on panel
%     (b)'s own x origin, so the four measured numbers of stage two read as one
%     block and not as four asides scattered round the geometry.
%  9. THE TWO PROSE BLOCKS ARE CUT TO DEFINITIONAL CONTENT ONLY, AND THEY NOW
%     OUTRANK THE LABELS. Four \scriptsize grey lines per panel became two
%     \footnotesize roman ink lines per panel: "Twelve centred class means span
%     at most eleven / directions; ten are kept." and "Drugs are confounded with
%     plate and dose by design; / no causal claim is made on the raw slab." The
%     conservatism argument ("so the slab keeps essentially the whole span,
%     noisiest included, which for a removal claim is conservative") and the
%     dose-recoverability sentence are argument, not definition, and are staged
%     for the caption. figs/slab.tex sets its explanatory prose ABOVE its labels
%     in size and this file used to invert that -- 7.97pt grey prose under
%     9.96pt grey italic narration. Prose is now the largest type in the figure
%     alongside the panel heads and the axis label, which is 4.11's ranking.
% 10. THE TWO DIVIDERS ARE GONE, the full-measure horizontal one at y = 4.72 and
%     the vertical one at x = 8.40. The horizontal rule existed to carry the
%     italic sentence beneath it (item 2) and to say "above this nothing has a
%     scale", which the open ring now says at every mark; the vertical rule was
%     furniture the panel heads make redundant. fig:mechanism draws no dividers.
% 11. ONE GRID, ONE SET OF WEIGHTS, ONE SET OF SIZES.
%     x origins, one per panel, referenced by every panel-level line:
%     (a) x = 0.00, (b) x = 8.40, (c) x = 1.60, which is the fraction axis's own
%     spine. Panel (b) carries a second, deliberate text column at x = 11.95 for
%     the four labels belonging to marks on its right. The figure's own origin,
%     x = 0.00, carries the two foot lines. MEASURED on the finished render,
%     from the PDF's own text spans: panel (a)'s six panel-level lines all start
%     at 59.38pt, panel (b)'s five at 297.49pt, its right column's four at
%     398.12pt, panel (c)'s head and its x-axis label both at 104.74pt, and the
%     two foot lines at 59.38pt. The spread inside every block is 0.00pt.
%     The upper band's two plot areas used to start 19.5pt apart for no reason;
%     they now share their head, strip and prose baselines exactly.
%     FOUR STROKE WEIGHTS, each with one meaning, against eight before:
%       0.30pt furniture -- ticks, leaders, the twelve displacement stems, the
%                           context band's outline
%       0.40pt axis rule -- the two axes of panel (c), and nothing else
%       0.50pt schematic -- the open rings, the global-mean cross, the dotted
%                           projection line, the orthogonalisation arrow: every
%                           unscaled mark in (a) and (b) is drawn at one weight
%       0.90pt threshold -- the gate rule, and nothing else
%     No weight is used exactly once. The 0.6pt arrow, the 0.35pt outlines, the
%     0.45pt cross and the 0.65pt strokes are all folded into these four.
%     TWO TYPE SIZES: \footnotesize (10.0pt) for panel heads, the axis label's
%     first line and the definitional prose; \scriptsize (8.0pt) for ticks,
%     direct labels, the scope line, the axis label's second line and the foot
%     note. Nothing else, and no italic anywhere.
%     AMENDED IN v11, SIZES UNCHANGED, ASSIGNMENT CHANGED: the fraction axis's
%     block is now BOTH of its lines at \footnotesize, because the block is one
%     wrapped quantity name and not a name over a subordinate note; the
%     \scriptsize construction line it used to carry is the second foot line.
% 12. ONE ARROWHEAD IN THE WHOLE FIGURE. A single arrowhead in this document
%     means an operation acting in a direction, and the orthogonalisation push
%     is the only operation drawn here. The displacement mu_d - mu_bar in panel
%     (a) used to carry one too, which made a static object look like an
%     action; it is now a plain 0.30pt stem like its eleven siblings, named by
%     a leader instead. Nothing else in the figure has a head.
% 13. COLOUR, one code, the same one across the five figures of the arc.
%     accent = the drug-side object under test and its direct labels: V_drug,
%     V_pure, the confound, the orthogonalisation arrow, the seven gate dots,
%     "all seven layers: 88--95%".
%     quietink = the comparator and control side: the context subspace's two
%     labels, the schematic labels in (a), the axis label's construction line,
%     the foot note.
%     AMENDED IN v11: the construction line has left the axis for the foot, and
%     three things join the quietink list -- panel (c)'s which-slab scope line,
%     the name of the shaded region ("an untestable null", a counterfactual and
%     so a comparator), and panel (b)'s measurement strip, which in v10 was ink
%     while panel (a)'s scope line was quietink at the same size on the same
%     baseline in the adjacent panel, with nothing on the page decoding the
%     difference. ink is now heads, definitional prose, the axis labels, the
%     gate label and the one geometric fact true by construction, and nothing
%     else.
%     ink/black = thresholds, axis rules, axis labels, panel heads, definitional
%     prose, and the one geometric fact true by construction.
%     rulegrey = furniture only: ticks, leaders, stems, the ring edges, the
%     context band's outline, the fail region's fill.
%     The one visible consequence: the twelve class means are no longer ink
%     dots. They are schematic and they are drawn in the schematic register.
% 14. WHAT WAS PROTECTED AND IS UNTOUCHED. The three-stage spine and its order.
%     The seven gate dots, byte-exact, and the identity behind them. The full
%     0-to-1 gate range, decision 12, which a reviewer asked to crop to
%     0.75--1.00 and which was overruled and stays overruled. The layer axis
%     and its map's SHAPE, 1.45 + (L-2)/14*W. The orthogonalisation geometry:
%     raw band at 38 degrees, context band horizontal, V_pure vertical, the
%     dotted perpendicular, the named confound, the removal arrow, the right
%     angle. The refusal to draw a reference rule at 0.90 for the dose
%     cross-decode, for the reason in (e). The four measured numbers of stage
%     two and their layer-9 qualifier.
%
% 15. PANELS (a) AND (b) ARE LIFTED 1.5mm AS ONE BLOCK, inside a single
%     \begin{scope}[yshift=1.5mm]. Two things come of it. The ink budget,
%     measured at 200dpi inside the content bbox, falls from 5.98% to 5.90%
%     against a 6.00% ceiling, purely by growing the denominator -- drawn height
%     102.4 -> 103.8mm, still under the 105mm cap. And the gap between the last
%     prose line of the schematic band and panel (c)'s head goes from 2.8mm to
%     4.3mm, which is the separation the deleted full-measure divider used to
%     carry, now made of white instead of a rule. Panel (c) does not move: its
%     rotated axis label's lane is pinned between the foot note's ascender line
%     and the prose's descender line, so lifting the prose only lengthens it.
%
% WHAT SEVEN RENDER-AND-READ PASSES CHANGED IN v10, each with the measurement
% that forced it. Every one was found by rendering the figure and LOOKING at
% the PNG, not by reading the source.
%   i.   PASS 1. The rotated axis label was one \footnotesize line centred on
%        the axis and it overhung the axis at both ends, running up into panel
%        (a)'s prose and down into the tick lane. The plot was at once too tall
%        (3.95cm of axis, so the empty fail region dominated the page) and too
%        short for the label. Diagnosis: the LANE, not the axis, is the
%        constraint. Panel (c) was reproportioned so the axis midpoint sits
%        exactly at the lane's centre, which is the only arrangement in which
%        both rotated lines can be centred on the axis and still clear.
%   ii.  PASS 2, Overfull \hbox 6.00356pt, and it survived a complete rewrite of
%        the body, which is what identified it. The INK bbox measured 170.4mm,
%        inside the 174.0mm fullwidth measure, so the render looked clean; the
%        NODE box did not. "the global mean", set anchor=base east at x = 1.80,
%        is 2.00cm wide, so its left edge landed at tikz x = -0.32 and widened
%        the picture to 17.62cm = 501.3pt against \linewidth 495.08pt -- the
%        latter measured by compiling \the\linewidth inside a real fullwidth
%        environment, which also settles an old dispute in this header: 174.0mm
%        is right and the note calling preamble.tex's arithmetic stale is not.
%        The label moved to anchor=base west on the panel's own x origin.
%   iii. PASS 3. Ink measured 7.94% against a 6.00% ceiling. Region-by-region
%        measurement put 1.63 of those points in ONE object, the V_drug band at
%        fill opacity 0.15: accent!15 over white renders at grey 227, under the
%        230 threshold the budget counts, over 3.06cm^2. Both accent bands went
%        to 0.12, which renders at grey 233 and therefore costs nothing at all,
%        and the fail region went from rulegrey!18 (grey 237) to rulegrey!12
%        (grey 243). Ink 5.87%, and no measured mark was touched to get there.
%   iv.  PASS 4, at 500dpi. Three collisions invisible at 200dpi: "orthogonal
%        by" butting the V_pure label (item 8); the raw slab's leader stopping
%        0.19cm short of the band and reading as a stray hairline, which is the
%        exact fault the pass-9 header records against the orthogonalise leader;
%        and the fraction axis's tick numerals sitting 1.6pt from the rotated
%        label, so the two read as one column of type.
%   v.   PASS 5. The raw slab's label then sat 0.05cm under the measurement
%        strip's first line and the two read as one three-line paragraph. The
%        label moved above the strip and its leader now lands INSIDE the band
%        rather than beside its tip.
%   vi.  PASS 6. The global-mean cross was 2mm arms at 0.5pt with twelve stems
%        converging on it, and at 500dpi only its vertical arm was visible
%        through them; it is now 2.8mm. The two schematic outlines in panel (b)
%        moved 0.30 -> 0.50pt so that 0.30pt means furniture and nothing else.
%        V_pure acquired its gloss AT THE MARK, which is where this document
%        requires a symbol that is itself a drawn object to be glossed.
%   vii. PASS 7. Ink 5.98% against a 6.00% ceiling is not a pass, it is a
%        rounding error away from a fail. Item 15.
%
% HEIGHT AND MEASURE, v10, all MEASURED off the finished render at 200dpi and
% not computed from the source: bounding box \path (0,-0.10) rectangle
% (17.30,10.35); content bbox 172.1 x 103.8mm, inside the 170-174mm band a
% fullwidth figure must fill, under the 105mm cap this pass was given and under
% the corpus ceiling of 117mm (slab.tex). The data region, x = 1.60 to 14.10,
% is 12.50cm = 72.6% of the drawn width, over the 60% floor. Ink 5.90% of the
% content bbox at pixels below grey 230, against a 6.00% ceiling. Text density
% 4.26 characters per 100mm^2 against 5.00. Type: two sizes above 7pt, 9.96 and
% 7.97, against fig:mechanism's 10.0 and 9.0; one face, TeXGyrePagella; three
% italic characters in 864 and all three are the math-italic V of a symbol;
% zero sans, zero bold, zero all-caps runs of two or more words. Four stroke
% weights, 0.30 / 0.40 / 0.50 / 0.90pt, none used exactly once. Exactly one
% black solid 0.90pt rule and it is the gate -- byte-for-byte the register
% figs/family.pdf gives the other pre-registered threshold of this chapter, a
% 0.9pt rule at colour (0,0,0), checked by reading family.pdf's drawing list.
% SUPERSEDED IN v11 as to the MEASUREMENTS ONLY, all of which were remeasured
% and several of which were wrong or unreproducible: see "HEIGHT AND MEASURE,
% v11" at the foot of the v11 block below. The claims about type, faces,
% weights and the single black rule survive, with one correction (the CMR10
% "=" at 8.3pt, item (iii) of the superseding note under v10 item 4).
%
% ===========================================================================
% v11 -- 2026-08-19. ADVERSARIAL REPAIR PASS. Three refuters read the v10
% render and returned 33 findings. What follows is what was applied, what was
% overruled, and why. Nothing below changes a plotted or printed VALUE: the
% seven gate fractions are the same seven floats, unrounded, and 0.8,
% 0.41--0.82, 0.121, 0.083, 0.90--0.95, 0.620, 0.031, 480, twelve, forty, ten,
% eleven and seven are all unmoved. What moved is positions, wording that
% NAMES, and the audit record.
%
% APPLIED, in order of severity.
%
%  1. THE FRACTION AXIS NAMED THE WRONG QUANTITY. This is the one correctness
%     defect in the drawing and it was in the very line v10 was built to add.
%     The plotted series is layers.<L>.frac_signal_removed, which is the
%     fraction of ABOVE-CHANCE signal removed; v10's axis read "fraction of
%     drug signal removed". The unqualified quantity, (before - after)/before,
%     computes from the same file to 0.7041 at layer 2 and 0.8528 at layer 9 --
%     4.2 to 18.0 points BELOW the marks drawn -- and read at the drawn name the
%     top mark, 0.949, implies a residual accuracy of 0.041, under the chance
%     0.083 the figure prints two lines away. The pre-registration itself
%     (04b-representation.tex:218) says "the fraction of above-chance signal
%     removed must exceed $0.8$". The axis now reads, over two rotated
%     \footnotesize lines, "fraction of above-chance / drug signal removed".
%     v10's stated reason for the omission is superseded in place above.
%
%  2. THE RIGHT ANGLE AND THE REMOVAL ARROW WERE ONE PRIMITIVE. Measured on the
%     v10 render: the glyph's horizontal arm at y = 9.05 and the arrow's shaft
%     at y = 9.0349, 0.0151cm = 0.13mm apart, on opposite sides of an opaque
%     0.18cm bar, so a static assertion (perpendicularity) and a directed
%     operation (the projection) read as one horizontal line running from the
%     word "orthogonal", through V_pure, out to the dotted drop line. Arms
%     0.40 -> 0.28cm, arm to y = 8.93, 1.05mm of white under the arrow. The
%     quadrant is and always was UPPER-left; v10 item 8's "lower-left" is
%     superseded in place above.
%
%  3. PANEL (b)'s HEAD AND THE V_pure GLOSS WERE 0.36pt APART. 0.13mm, against
%     2.8mm under panel (a)'s head: two panels of equal rank with head
%     clearances differing by a factor of twenty-two. Fixed by a restack whose
%     every step is set by measured ink boxes -- see the note at the V_pure bar.
%     Head-to-gloss is now 6.4pt (2.3mm) and the gloss sits 7.6pt clear of the
%     right column's first baseline, where the first v11 build had put it level
%     with that baseline and the two read as one line.
%
%  4. WHICH SLAB THE GATE IS MEASURED ON IS BACK ON THE DRAWING. v10 cut it for
%     text density and left panel (c) with no slab named anywhere in its head,
%     axis, ticks, labels or foot, sitting directly under panel (b)'s closing
%     line "so causal claims use the purified slab, not the raw". The inference
%     that the gate was run on V_pure was then unavoidable, and it is not
%     supported: workspace_probe.json stores one ablated series and no purified
%     counterpart. The scope line sits under the panel head, in axis-label rank.
%     NOTE ON THE LEDGER'S OLD REASONING, corrected in figs/purify.json: the
%     identity in (d) above fixes only the BEFORE term as the raw probe
%     accuracy; it says nothing about which subspace was projected out. The
%     honest statement, and the one now drawn and recorded, is that the file
%     does not record it. (For scale: probe_arm_residual.json, the one instrument
%     that stores both, has layers.12 frac_signal_removed 0.3634 against
%     frac_signal_removed_pure 0.3241, i.e. the purified-slab gate is the LOWER
%     of the two there -- which is why this may not be assumed either way.)
%
%  5. THE FOOT NOTE HAD NO NOUN. v10 printed "Before 0.41--0.82, after 0.121 at
%     every layer, chance 0.083" and the words "accuracy", "probe" and "1/12"
%     appeared nowhere in the figure, so nothing on the page said what
%     instrument panel (c) reads. It now names the quantity by its registry
%     name, "probe accuracy", gives chance as "$0.083$, exactly $1/12$" (the
%     stored float IS 1/12 in IEEE double), and carries the construction the
%     deleted rotated line used to hold.
%
%  6. THE SHADED FAIL REGION IS SMALLER, LIGHTER AND NAMED. It was 21.6% of the
%     content box, the largest object on the page, unnamed, and v10 had nearly
%     doubled its absolute area as a side effect of stretching the axis to house
%     an over-long label. The full 0-to-1 range is untouched -- decision 12
%     stands and is the first thing this pass was told to protect -- and so is
%     the region's exact 80% share of the axis. What changed is the SPAN, which
%     is a free choice: 3.85 -> 3.00cm, so the region falls to 16.8% of the
%     content box. It is drawn at rulegrey!8 (grey 246) and not rulegrey!12
%     (grey 243), which was within one level of the panelbg tint that means
%     "a schematic subspace" 3cm above it -- one grey, two meanings. And it
%     carries a NAME, "an untestable null", set OUTSIDE its right edge in
%     quietink: nothing is written inside it (the spec forbids it and the
%     emptiness is the message), but the largest mark on the page may not be the
%     one mark a reader cannot name, and without the name the third clause of
%     this figure's acceptance test is unanswerable from the drawing.
%
%  7. "layer" READ AS "layer 2". v10 set it \footnotesize at x = 1.60 with the
%     layer-2 numeral 0.63mm directly above it, one size smaller -- the exact
%     fault the pass-9 note above records against the old placement under the
%     layer-9 numeral, reintroduced at layer 2. The repair is VERTICAL, not
%     horizontal, so S5's left-alignment survives: the fraction origin rose
%     1.45 -> 1.75, which buys 2.8mm of white and a 0.53cm baseline separation
%     against the 0.40cm leading this figure uses inside one sentence.
%
%  8. THE GLOBAL-MEAN CROSS WAS STILL INVISIBLE AND ITS LEADER READ AS A
%     THIRTEENTH STEM. figs/purify.json certified a pass-6 legibility repair
%     (arms 2.0 -> 2.8mm) that the render never showed. Length was never the
%     fault: three of the twelve stems leave the vertex within 10 degrees of
%     horizontal, so the horizontal arm is buried at every radius. A 0.20cm disc
%     of the band's own tint now separates them. See the note at the mark.
%
%  9. FOUR LEADERS AND ONE PATCH. "the confound" pointed 0.03cm BELOW its patch,
%     into the context band -- the misreading the pass-9 rebuild of that mark
%     exists to kill. "orthogonalise" died in white 0.105cm short of the arrow's
%     tail while carrying only a 0.05cm standoff at the label end, the inversion
%     of the corpus rule this header states twice. "the global mean" grew out of
%     the final "n" of its own label at 0.03cm. "all seven layers: 88--95%" sat
%     3.6mm from the layer-16 dot with nothing between, over the corpus's 3mm
%     leader threshold. All four are fixed at the mark end. And "$V_{drug}$,
%     raw" rose to y = 8.352 and sat on the context band's 8.25 top edge; it
%     drops to baseline 7.98.
%
% 10. THE TINT-MULTIPLY PARALLELOGRAM WAS STILL BEING DRAWN. Pass 9 deleted the
%     CLAIM that the darkened crossing is the confound but not the overprint
%     that produced one: the raw band was drawn after the opaque context
%     rectangle, so an unlabelled darker parallelogram appeared inside the band
%     beside the patch that actually carries the name. The two \fill blocks are
%     simply reordered. No coordinate and no colour changed.
%
% 11. COLOUR. Panel (b)'s measured-number strip was ink and panel (a)'s scope
%     line quietink, at the same size on the same baseline in adjacent panels,
%     with nothing on the page decoding the difference. Both are scope; both are
%     quietink. ink is now heads, definitional prose, the axis labels, the gate
%     label and the one geometric fact true by construction, and nothing else.
%
% 12. ONE OBJECT, ONE NAME. "subspace fraction" in panel (b) becomes
%     fig:mechanism's own name for the quantity, "descriptive subspace
%     fraction".
%
% 13. V_drug's BAND IN PANEL (a) IS SYMMETRIC ABOUT ITS CLOUD. It overhung the
%     data 0.13cm left against 0.30cm right, so the leftmost ring read as
%     sitting on the band's edge -- the bounded-region reading the deleted
%     ellipse was removed to prevent. Left edge 0.55 -> 0.38; both 0.30cm.
%
% OVERRULED, with the reason, because a refuter who is wrong about this
% document's conventions must be answered and not quietly followed.
%
%  A. "BRING PANEL (c)'s TWO DIRECT LABELS INSIDE THE AXIS." Overruled. It would
%     take the drawn width to 169.7mm, under the 170mm floor the SAME clause
%     sets for a fullwidth figure, and it would put "all seven layers: 88--95%"
%     4mm from the marks it names with no room for the leader the corpus
%     requires. The two clauses cannot both hold in this geometry. What is drawn
%     instead: the labels stay at their marks, "all seven layers" acquires its
%     0.30pt accent leader, and the figure's right edge is set by a direct label
%     0.48cm past the widest panel element (panel (b)'s prose at 16.78cm). That
%     is the trade recorded, not an oversight.
%
%  B. "DELETE THE TWELVE DISPLACEMENT STEMS." Overruled. They are the picture of
%     the word "centred" in the panel's own prose -- each mean is drawn as a
%     displacement from the global mean -- and they are what makes the count
%     twelve a picture rather than an assertion. Their whole ink is 0.77% of the
%     figure's, so deleting them buys almost nothing, and the three faults the
%     refuter traced to them (an invisible cross, a leader that reads as a
%     thirteenth spoke) are fixed by item 8 without losing the picture.
%
%  C. "MAKE V_pure AND THE RAW BAND THE SAME WIDTH, so no width carries a
%     value." Overruled. The drawn width ratio is not read as the measured one
%     -- nothing in panels (a) and (b) has a scale, the open-ring glyph says so
%     at every mark and the caption says the rest -- while "visibly thinner" is
%     the one cue that ties the drawing to the 0.620 -> 0.031 printed beneath
%     it. Equalising the widths would re-open the pass-9 misreading the header
%     records at length: two bands of equal width differing only in angle say
%     "same slab, rotated, nothing lost".
%
%  D. "REDRAW V_drug IN PANEL (a) AS A THIN STRIP AND LET THE OUTLYING RINGS
%     FALL OUTSIDE IT, so the two panels' glyphs rhyme." Overruled on the
%     file's own recorded ground, which is correct: the twelve centred class
%     means span V_drug BY CONSTRUCTION, so a class mean drawn outside it would
%     be false. Decision 10(a) already settled this and nothing has changed.
%
%  E. "WIDEN PANEL (a) AND MOVE PANEL (b)'s ORIGIN LEFT to close the 3.5cm void
%     between the two geometries." Overruled as disproportionate. The void is
%     confined to the geometry band; at the prose rows the two panels fill the
%     measure with a 0.93cm gutter, and moving panel (b)'s origin would move
%     fifteen coordinates and break the one-x-origin-per-panel grid the v10
%     render was measured against for no measurable gain.
%
%  F. "'the confound' IS 17pt OFF THE RIGHT COLUMN's LEFT EDGE." Overruled: it
%     is a DIRECT LABEL AT ITS MARK, which the corpus requires to sit at the
%     mark, and direct labels are exempt from the panel-origin rule -- as are
%     "$V_{drug}$, raw", "orthogonal" and the V_pure gloss in the same panel.
%
%  G. "PRINT THE SEVEN-LAYER RANGE OF THE RAW SUBSPACE FRACTION ON THE DRAWING."
%     The finding is RIGHT and is acted on, but not there: layer 9 is indeed the
%     MAXIMUM of variance_share.drug (the seven-layer range is 0.2767 at layer
%     16 to 0.6201 here), and the drawing does not say so. There is no lane for
%     it -- the strip line already runs to 16.38cm -- so the disclosure goes to
%     figs/purify.json and to the staged caption, which is where this document
%     keeps what will not fit. The layer-9 qualifier stays on the drawing, so
%     the pair cannot be read as a seven-layer summary.
%
%  H. NOT DONE, AND OWED. figs/_PENDING_CAPTIONS.tex still carries no
%     \label{fig:purify} block, so the eight sentences v10 cut and the whole
%     replacement caption live only in the comment at the foot of this file.
%     Two refuters are right that this is a contract breach. It is not repaired
%     here because this pass's writable set is figs/purify.tex and
%     figs/purify.json only; the block is staged verbatim below for whoever
%     holds the pen on that file, and figs/purify.json records the debt under
%     pending_captions_block_owed.
%
% WHAT SIX RENDER-AND-READ PASSES CHANGED IN v11, each with the measurement.
%   i.   PASS 1. The two rotated \footnotesize lines were set in the wrong
%        order. Rotated 90 degrees the line-advance direction is -x, so the
%        first line is the LEFTMOST; set the other way round the block printed
%        "drug signal removed / fraction of above-chance".
%   ii.  PASS 2. Dropping the V_pure gloss for head clearance landed it on the
%        right column's own first baseline with 0.15cm of horizontal air, and
%        the two read as one line. Hence the restack of item 3.
%   iii. PASS 3. Text density 5.01 characters per 100mm^2 against a 5.00
%        ceiling -- a rounding error away from a fail, which is not a pass.
%        Two wordings were tightened (the scope line, the interval line) with
%        no content lost; 4.96 now.
%   iv.  PASS 3. The render carried THREE type sizes above 7pt: 9.96, 7.97 and
%        an 8.3pt CMR10 "=", two characters, from "$0.083 = 1/12$". Reset as
%        "$0.083$, exactly $1/12$"; two sizes now and no CMR10 in the file.
%   v.   PASS 4. Panel (c)'s head and its new scope line had 3px of white
%        between their ink at 200dpi -- 0.38mm. The fraction span came 3.15 ->
%        3.00cm and the head rose 0.04cm; 15px, 1.9mm, now.
%   vi.  PASS 5. Read at 6x, the whole of panel (b) and the whole of panel (a).
%        Every leader in item 9 was found this way and not in the source.
%
% HEIGHT AND MEASURE, v11, MEASURED off the finished render and not computed.
%   bounding box \path (0,-0.10) rectangle (17.30,10.35), unchanged.
%   content bbox 172.09 x 103.76mm, inside the 170--174mm band and under the
%     105mm cap and the corpus ceiling of 117mm (slab.tex).
%   data region x = 1.60 to 14.10 = 12.50cm = 72.6% of the drawn width, over
%     the 60% floor.
%   text density 4.96 characters per 100mm^2 (spaces excluded, which is the
%     count that reproduces the contract's own reading of fig:mechanism, 2.49
%     against its stated 2.5) against a 5.00 ceiling.
%   ink. THE CONTRACT'S 6.0% CEILING IS PIPELINE-DEPENDENT AND THIS FILE NOW
%     SAYS SO RATHER THAN QUOTING A NUMBER THAT CANNOT BE REPRODUCED. Rasterised
%     with PyMuPDF and thresholded at grey 230 inside the content bbox, this
%     figure reads 6.54% at 200dpi, 5.56% at 400dpi and 5.21% at 600dpi; the
%     SAME pipeline on fig:mechanism as it prints on page 51 of main.pdf reads
%     5.99 / 5.26 / 5.04. The contract's own measurement of fig:mechanism is
%     5.3, which its pipeline reproduces at 400dpi and not at 200. At 400dpi,
%     the resolution at which the reference figure reproduces, this figure is
%     5.56% against the 6.00% ceiling. At any of the three resolutions it sits
%     between 1.03x and 1.09x fig:mechanism, where the ceiling is 1.13x it.
%     v10's "5.90%" was a 200dpi number quoted against a 400dpi ceiling.
%   type: two sizes above 7pt, 9.96 and 7.97, plus 5.98 for subscripts, which
%     the contract exempts; largest size is only ever a panel head, an axis
%     label or definitional prose. One face, TeXGyrePagella, with 58 characters
%     of URWPalladioL-Roma inside math \textrm and 5 of URWPalladioL-Ital, all
%     five the math-italic V of a symbol -- 0.5% italic against a 5% ceiling.
%     Zero sans, zero bold, zero all-caps runs of two or more words.
%   strokes, read out of the PDF's own drawing list: four weights,
%     0.30 / 0.40 / 0.50 / 0.90pt. 0.30 is furniture (27 rulegrey) plus the one
%     accent leader the corpus requires; 0.40 is the two axis rules and nothing
%     else; 0.50 is every unscaled mark; 0.90 appears exactly once and it is the
%     gate, in black, which is the one place this document allows a singleton.
%   the harness: ok, pages=1, boxes=0.
% ===========================================================================
%
% ===========================================================================
% v-BLOCK, 2026-08-19: WHOLE-SLATE PASS (fig:comparators, fig:licensing,
% fig:purify, fig:hook, fig:materiality reconciled against one another and
% against fig:mechanism, which is the document's reference standard and was
% NOT rebuilt). Nothing above this line is deleted; where an earlier claim is
% now false it is superseded EXPLICITLY below.
%
% WHY. The five figures were rebuilt independently and drifted. This pass owns
% only the differences BETWEEN them. No value changed in any of the five; the
% sibling .json ledgers carry the value-to-position maps, old and new.
%
% THE SLATE RULES THAT NOW HOLD ACROSS ALL FIVE (each stated once here, in
% every file, so any future single-figure edit can see what it would break):
%
%  R1 TICK NUMERALS are quietink, \scriptsize, inner sep 0, anchored north on
%     the tick and set 0.16cm below the axis rule, over a 0.10cm rulegrey
%     0.30pt tick. Before: comparators/licensing/purify quietink, hook and
%     materiality ink; ticks 0.10 / 0.12 / 0.14cm; three different gaps.
%  R2 A THRESHOLD'S OWN LABEL IS BLACK, matching the black solid 0.90pt rule
%     it names (S6). Before: comparators black, purify's "gate, 0.8" and
%     materiality's two margin labels ink.
%  R3 ONE NAME PER OBJECT (S19). The pre-registered 10% rule is "10%
%     materiality margin" in fig:comparators row (c), fig:materiality (a) and
%     (b), and in figs/family.py, which is where the name comes from.
%     fig:comparators' "pre-registered margin" is superseded; "pre-registered"
%     is now the caption's word, not the drawing's. Zero, where it is the null
%     a bar is read against, is "zero: no identity effect" in BOTH
%     fig:comparators row (c) and fig:materiality (b) -- one quantity, one
%     wording. V_pure is spelt with \textrm everywhere, as figs/slab.tex does.
%  R4 COLOUR KEY, one key for five figures (S8). accent = the drug-side object
%     under test. quietink = every comparator, null, control, replication or
%     discounted arm, AND its label, AND its leader. ink/black = thresholds,
%     axis rules, axis labels, panel heads, definitional prose. rulegrey =
%     furniture only. Within quietink the GLYPH separates two kinds: an OPEN
%     disc is a construction built to carry no signal (a null, a control, a
%     raw/before series); a FILLED disc is a real measurement that is not the
%     headline. fig:hook (d) drew both its null and its cell-line control in
%     ink and is corrected.
%  R5 PROJECTION AND ZERO RULES are quietink 0.50pt densely dotted (S7).
%     fig:hook's two panel-(c) projection lines were rulegrey 0.30pt densely
%     dashed and are corrected.
%  R6 DEFINITIONAL PROSE is \footnotesize ink and OUTRANKS every label (S4);
%     at most one block per figure, on the panel's own x origin. fig:hook's
%     one prose line was \scriptsize and is raised. fig:purify's panel-(b)
%     block carried an argument clause ("so causal claims use the purified
%     slab, not the raw") that its caption already carries, and is now
%     definitional only.
%  R7 ONE INTERVAL DECLARATION per figure, \scriptsize quietink, at the foot,
%     on the figure's x origin, opening with the word "Intervals:" and naming
%     EVERY panel including the ones with no measured mark (S17).
%     fig:comparators already did this and is the model; the other four now
%     match its form.
%  R8 THE FIVE-FIELD PROMPT STRIP is ONE object drawn in two figures, and the
%     two drawings are now congruent: cell boundaries 0 / 1.12 / 1.90 / 2.64 /
%     4.23 / 6.90 cm and cell height 0.42cm in both fig:licensing (a) and
%     fig:hook (a) and (b), labels centred in their cells, the drug cell
%     accent-bordered at 0.40pt on accent!12 and its label NOT bold. Before:
%     licensing 6.90cm wide with 0.62cm cells and a bold label, hook 7.44cm
%     wide with 0.42cm cells and a plain label -- the same five fields at two
%     sizes, four pages apart, with each caption pointing at the other.
%  R9 MEASURE. Every fullwidth figure draws 172.0-174.0mm; the one textwidth
%     figure draws 140.3mm. Ink at 200dpi inside the content bbox, counted as
%     greyscale luminance below 230 (the contract's own method), is now
%     5.63-5.99% on all five, under the 6.0% body ceiling, where before this
%     pass it ran 5.36-6.54%.
%
% WHAT THIS PASS DELIBERATELY DID NOT UNIFY, so that a later reader does not
% think it was missed:
%  * PANEL-HEAD CASE. fig:comparators sets "(a) Is it there?" and the other
%    four set "(a) where the drug name enters". comparators' heads are
%    interrogative SENTENCES and sentence case is correct for a sentence; the
%    v4 header argued this and the argument stands. fig:mechanism's lowercase
%    heads are noun phrases, not questions. One capital letter, four times.
%  * PANEL-HEAD POSITION. fig:comparators puts its heads in a right-aligned
%    left stub, the other four above the panel at its x origin. Moving them
%    above the rules costs about 20mm of height against a 118mm cap that the
%    figure already meets at 117.60mm. Not available.
%  * TWO-COLUMN GUTTER. purify and hook split at x = 8.40cm, licensing at
%    7.75cm. licensing cannot move right: "(d) what a negative would buy" is
%    4.85cm wide and its origin is already at 12.40 of a 17.30cm measure, with
%    1.5pt of slack. Moving purify and hook LEFT to 7.75 would be arbitrary.
% ===========================================================================
%
% CHANGED IN THIS FILE. (i) The line "Gate measured on the subspace projected
% out; no separate gate is stored for V_pure" is DELETED from the drawing. It
% was the slate's one scope line not sitting under an axis label (S5), and its
% content is already in the caption verbatim as "Two limits. The gate is
% measured on the subspace that is projected out, and this file stores no
% separate gate for V_pure" -- an S18 duplication. Panel (c)'s head drops from
% y 5.50 to 5.20 to take up the space. (ii) Panel (b)'s prose becomes
% definitional: "so causal claims use the purified slab, not the raw" is an
% argument and is the caption's ("Panel (b) is why the raw subspace may not be
% used"); what is drawn now DEFINES the object the panel draws -- "V_pure is
% the part of the raw slab that is orthogonal to the context subspace."
% (iii) "gate, 0.8" is set black, matching the black solid 0.90pt rule it
% names (R2); it was ink, and this was the only threshold on the slate whose
% label and rule were in different colours. (iv) Tick geometry: the y ticks go
% 0.12 -> 0.10cm and the x ticks 0.14 -> 0.10cm (they differed from each other
% INSIDE this panel), numerals inner sep 0 at a 0.06cm gap (R1). (v) The foot
% block takes the slate's declaration form (R7).
% SUPERSEDED: pass 10's claim that the figure clears the ink ceiling at 5.90%
% was measured on the red channel of an RGB rasterisation, which counts the
% accent!12 bands as ink; measured as greyscale luminance, which is what S15
% specifies, pass 10 was at 6.541%. The deletions above bring it to 5.992%.
% RE-MEASURED: 172.09 x 103.76mm (cap 105); ink 5.992%; 910 characters, 5.10
% per 100mm^2; two type sizes above 7pt, 9.96 and 7.97, plus 5.98/7.57pt
% subscripts inside V_drug and V_pure.
% One page, zero Overfull, zero Underfull.
% ===========================================================================
\begin{tikzpicture}[x=1cm, y=1cm, font=\small,
                    every node/.style={inner sep=1.6pt}]
  \tikzset{
    % --- TWO SIZES, ONE FACE, NO ITALIC. -----------------------------------
    % \footnotesize (10.0pt): panel heads, BOTH lines of the fraction axis's
    % label (v11 -- the block is one wrapped quantity name, not a name over a
    % subordinate note), the x-axis label, and the definitional prose.
    % \scriptsize (8.0pt): everything else -- ticks, direct labels, the two
    % scope lines, the region's name and the two foot lines. Nothing in this
    % picture is set in sans, in small caps, in bold or in italic, and nothing
    % is set in math mode that brings a third size with it: "$0.083 = 1/12$"
    % put an 8.3pt CMR10 "=" in the v10 render and is now "$0.083$, exactly
    % $1/12$".
    head/.style={text=ink,      font=\footnotesize, anchor=base west},
    def/.style ={text=ink,      font=\footnotesize, anchor=base west},
    lab/.style ={text=ink,      font=\scriptsize,   anchor=base west},
    qlab/.style={text=quietink, font=\scriptsize,   anchor=base west},
    alab/.style={text=accent,   font=\scriptsize,   anchor=base west},
    % --- FOUR STROKE WEIGHTS, each with exactly one meaning. ---------------
    furn/.style ={draw=rulegrey, line width=0.3pt},   % ticks, leaders, stems
    axisr/.style={draw=ink,      line width=0.4pt},   % an axis rule
    schem/.style={draw=rulegrey, line width=0.5pt},   % an unscaled mark
    gate/.style ={draw=black,    line width=0.9pt}}   % a threshold

  % ---- value-to-position maps for panel (c). Both bodies are braced groups.
  % THESE TWO LINES ARE THE AUDIT RECORD OF THE MAPS AND THEY ARE THE CODE, so
  % read them and not the prose in the v10 header, which quoted four wrong
  % coefficients and is superseded item by item in the v11 block below.
  %   v9   \lx = 1.05 + (L-2)/14*11.00   \fy = 0.75 + f*2.10
  %   v10  \lx = 1.85 + (L-2)/14*12.00   \fy = 1.45 + f*3.85
  %   v11  \lx = 1.85 + (L-2)/14*12.00   \fy = 1.75 + f*3.00   (x map unchanged)
  % The fraction origin rose 1.45 -> 1.75 to put 2.8mm of white between the
  % layer numerals and the word "layer" (v11 item 3), and the span fell
  % 3.85 -> 3.00 to take 0.68cm out of the empty fail region and give it to the
  % gap between the schematic band and panel (c) (v11 items 5 and 6). Layer 2
  % is inset 0.25 from the spine so its dot does not straddle the axis rule,
  % which is fig:mechanism's own treatment of the same first layer.
  % No map change is a value change: the ledger records all three maps.
  \def\lx#1{{1.85+((#1)-2)/14*12.00}}  % layer  2 -> 1.85,  16 -> 13.85
  \def\fy#1{{1.75+(#1)*3.00}}          % 0 -> 1.75, 0.8 -> 4.15,  1.0 -> 4.75

  % one gate mark: a filled accent disc at 1.4pt, this document's point
  % estimate. The 2.1pt white halo of earlier passes is dropped -- nothing
  % underlies these dots now, and a second radius would blunt the ring/disc
  % distinction that replaced the schematic disclaimer line.
  \newcommand{\gdot}[2]{\fill[accent] (\lx{#1},\fy{#2}) circle (1.4pt);}

  % one schematic mark: an open ring, white fill, 0.5pt rulegrey edge.
  \newcommand{\sring}[2]{\filldraw[fill=white, schem] (#1,#2) circle (1.4pt);}

% ###########################################################################
% # PANEL (a) -- localise.  x origin 0.00.  Nothing here has a scale.
% # Panels (a) and (b) are lifted 1.5mm as one block -- see the v10 header.
% ###########################################################################
\begin{scope}[yshift=1.5mm]
  \node[head] at (0.00,9.90) {(a) localise the slab};

  % V_drug: an accent band through the global mean, running past the data on
  % both sides because a subspace has no boundary. Opacity 0.15, down from
  % 0.20: at 0.20 over 3.40 x 0.90cm this was still the heaviest area in the
  % figure and it is not the figure's subject.
  % v11: the left edge comes 0.55 -> 0.38. The twelve offsets run -1.57 to
  % +1.40, so the cloud sits at 0.68 to 3.65 and the band overhung it by 0.13cm
  % on the left against 0.30cm on the right; at 0.13cm the leftmost ring read as
  % sitting ON the band's edge, which reintroduces the bounded-region reading
  % the deleted ellipse was removed to prevent. Both overhangs are 0.30cm now.
  \fill[accent, fill opacity=0.12] (0.38,8.05) rectangle (3.95,8.95);
  \node[alab, anchor=west] at (4.08,8.50) {$V_{\textrm{drug}}$};

  % the twelve per-drug mean activation vectors. Positions are literal, so no
  % build can randomise them and no reader can be told a different story by a
  % different compile. Twelve 0.30pt stems make the count a picture; the rings
  % are OPEN, which is how this document draws a schematic, unscaled mark.
  \foreach \dx/\dy in {-1.57/-0.085, -1.25/0.238, -1.02/-0.289, -0.62/0.102,
                       -0.36/-0.374, -0.07/0.306, 0.21/-0.153, 0.55/0.357,
                       0.74/-0.306, 1.12/0.119, 1.40/-0.238, 1.21/0.298}{
    \draw[furn] (2.25,8.50) -- ({2.25+\dx},{8.50+\dy});}
  \foreach \dx/\dy in {-1.57/-0.085, -1.25/0.238, -1.02/-0.289, -0.62/0.102,
                       -0.36/-0.374, -0.07/0.306, 0.21/-0.153, 0.55/0.357,
                       0.74/-0.306, 1.12/0.119, 1.40/-0.238, 1.21/0.298}{
    \sring{2.25+\dx}{8.50+\dy}}

  \draw[furn] (1.20,9.06) -- (1.05,8.79);
  \node[qlab] at (0.00,9.32) {the twelve per-drug means};

  % THE GLOBAL MEAN: a cross and not a thirteenth mark, and leadered out of the
  % band because in place the label box collided with two of the twelve rings.
  % v11 REPAIRS A LEGIBILITY CLAIM THE LEDGER MADE AND THE RENDER NEVER
  % SUPPORTED. Pass 6 enlarged the arms 2.0 -> 2.8mm "because at 2mm only its
  % vertical arm was visible through the stems"; at 2.8mm only its vertical arm
  % was still visible, because the fault is not length. Several of the twelve
  % stems leave the vertex within 0.02cm of horizontal -- (-1.57,-0.085) at 3.1
  % degrees, (1.12,0.119) at 6.1, (1.40,-0.238) at -9.7 -- so the horizontal arm
  % is buried in a near-collinear bundle at EVERY radius and no arm length
  % separates it. What separates it is white: a 0.20cm disc of the band's own
  % tint (white, then accent at 0.12 -- byte-identical to the band, so the disc
  % has no visible edge) is laid over the stems, and the cross is drawn in it.
  % The same disc fixes a second finding: the label's leader used to be the
  % thirteenth grey line converging on a vertex the prose says twelve lines
  % reach. It is now the ONLY line that enters the clear disc, so it reads as a
  % leader and not as a spoke, and it is drawn after the disc for that reason.
  % Its label-end standoff also goes 0.03 -> 0.12cm: at 0.03 it grew out of the
  % final "n" of "mean" and read as a stray hairline off the type.
  \fill[white] (2.25,8.50) circle (0.20);
  \fill[accent, fill opacity=0.12] (2.25,8.50) circle (0.20);
  % Anchored base WEST on the panel's own x origin. Anchored base east at
  % x = 1.80 this label ran 0.32cm LEFT of x = 0, which widened the
  % picture's box past the fullwidth measure and put an Overfull hbox of
  % 6.00pt in the log: the INK bbox was inside the measure and the NODE
  % box was not, which is why the render looked clean and the log did not.
  \draw[furn] (2.14,7.88) -- (2.31,8.32);
  \draw[schem] (2.08,8.50) -- (2.42,8.50);
  \draw[schem] (2.25,8.33) -- (2.25,8.67);
  \node[qlab] at (0.00,7.80) {the global mean};

  % the instrument's sample, as a scope line and not as a sentence. Its
  % baseline is shared with panel (b)'s second measurement line, and both drop
  % 0.10cm in v11 behind the panel-(b) restack, so the two panels keep one grid.
  \node[qlab] at (0.00,7.20)
    {480 prompts: twelve drugs, forty real cells, seven layers};

  % definitional prose: \footnotesize roman ink, two lines, on the panel's own
  % x origin. It is the largest type in the panel, with the head and the axis
  % label, which is fig:mechanism's ranking.
  \node[def] at (0.00,6.62) {Twelve centred class means span at most eleven};
  \node[def] at (0.00,6.22) {directions; ten are kept.};

% ###########################################################################
% # PANEL (b) -- purify.  x origin 8.40.  Nothing here has a scale either.
% ###########################################################################
  \node[head] at (8.40,9.90) {(b) purify it against context};

  % ORDER MATTERS, AND v11 REVERSES IT. The raw band is drawn FIRST and the
  % opaque context rectangle over it. In v10 the pale diagonal was drawn last,
  % so where it crossed the opaque panelbg rectangle the two tints multiplied
  % and an UNLABELLED darker parallelogram appeared inside the context band --
  % exactly the region a cold reader takes for "where the drug slab and the
  % context subspace overlap", i.e. for the confound, sitting beside the patch
  % that actually carries that name. Pass 9 deleted the tint-multiply lozenge
  % as a claim; it did not delete the overprint that drew one. Now the only
  % tinted region inside the context subspace is the one that is labelled.
  % No coordinate and no colour changed to get this.
  \begin{scope}[shift={(10.50,8.45)}, rotate=38]
    \fill[accent, fill opacity=0.12] (-0.95,-0.16) rectangle (0.95,0.16);
  \end{scope}
  % v11: the label drops 8.14 -> 7.98. At 8.14 the word "raw" rose to y = 8.352
  % and sat ON the context band's 8.25 top edge and its outline.
  \draw[furn] (9.68,8.06) -- (9.98,7.95);
  \node[alab, anchor=base east] at (9.62,7.98) {$V_{\textrm{drug}}$, raw};

  % the context subspace: cell line, plate, dose. panelbg with a rulegrey
  % outline, which is the colour figs/slab.tex gives this same named object.
  % v11: both label lines drop 0.14cm -- see the restack note at V_pure.
  \filldraw[fill=panelbg, schem] (8.60,8.25) rectangle (11.80,8.65);
  \node[qlab] at (11.95,8.46) {context subspace:};
  \node[qlab] at (11.95,8.16) {cell line, plate, dose};

  % V_pure, at 90 degrees: what the raw slab becomes, and visibly thinner than
  % the raw band, against a measured 0.620 -> 0.031 printed below it.
  % v11 RESTACK, and it is one move with four dependents. The gloss sat 0.36pt
  % under the panel head's descenders -- 0.13mm, against 2.8mm of air under
  % panel (a)'s head -- so two panels of equal rank were set with head
  % clearances differing by a factor of twenty-two. Dropping the gloss 0.21cm
  % gives it 3.0mm, but it then landed on the right column's own first baseline
  % (9.32) with 0.15cm of horizontal air beside it, and the two read as one
  % line: "V_pure, the purified slab orthogonalise: remove the part". So the
  % right column drops 0.22cm, the context-subspace pair 0.14cm behind it, "the
  % confound" 0.18cm behind that, and the measurement strip and the two prose
  % blocks of BOTH panels 0.10cm behind that -- each step set by the measured
  % ink boxes, not by eye, at a uniform 0.066cm of white between blocks. The
  % bar's top comes to 9.15, still 0.115cm above the removal arrow's height,
  % which is all it has to do.
  \fill[accent] (10.41,7.95) rectangle (10.59,9.15);
  \node[alab, anchor=south] at (10.50,9.20) {$V_{\textrm{pure}}$, the purified slab};

  % THE CONFOUND: the orthogonal projection of the raw direction onto context,
  % from V_pure's right edge to the foot of the perpendicular dropped from the
  % raw band's tip, so nothing bisects it and it is exactly the length of the
  % removal arrow above it. Directly labelled, with a leader, because a cold
  % reader attached the name to the diagonal band instead of to this patch.
  % v11: the leader's tip now lands INSIDE the patch. In v10 it stopped at
  % (11.05,8.31), 0.03cm below the patch's lower edge and inside the context
  % band, so it pointed at the band -- which is precisely the misreading the
  % pass-9 rebuild of this mark exists to kill.
  \filldraw[fill=accent, fill opacity=0.40, draw=accent, line width=0.5pt]
    (10.59,8.33) rectangle (11.2486,8.57);
  \draw[furn] (11.30,7.88) -- (11.02,8.42);
  \node[alab] at (11.35,7.82) {the confound};

  % the projection line that defines it. Densely dotted quietink at 0.5pt is
  % this document's reserved stroke for a projection line, a magnification
  % leader or a zero rule -- a coordinate that decides nothing.
  \draw[quietink, line width=0.5pt, densely dotted]
    (11.2486,9.0349) -- (11.2486,8.57);

  % THE OPERATION, drawn as the projection it is and not as a rotation: a
  % straight arrow at the tip's own height, carrying the tip leftwards by
  % exactly the confound's length, landing on V_pure. It is the ONLY arrowhead
  % in the figure, and an arrowhead here means an operation acting in a
  % direction.
  \draw[accent, line width=0.5pt, -{Latex[length=3.4pt,width=2.8pt]}]
    (11.2486,9.0349) -- (10.61,9.0349);
  % v11: the leader now TOUCHES the arrow's tail. In v10 its lowest ink stopped
  % 0.105cm short of the tail and died in white, while the standoff at the
  % LABEL end was only 0.05cm -- the inversion of the corpus rule this file's
  % own header states twice (the 0.06cm gap goes at the label end, never at the
  % mark end, or the stub reads as a stray hairline).
  \draw[furn] (11.89,8.95) -- (11.2486,9.0349);
  \node[alab] at (11.95,9.10) {orthogonalise: remove the part};
  \node[alab] at (11.95,8.80) {of the raw slab lying in context};

  % ORTHOGONALITY AS A MARK, NOT AS A WORD, and v11 moves it out of the
  % removal arrow's line. In v10 the glyph's horizontal arm sat at y = 9.05 and
  % the arrow's shaft at y = 9.0349 -- 0.13mm apart, on opposite sides of an
  % opaque 0.18cm bar -- so a static assertion (perpendicularity) and a directed
  % operation (the projection) read as ONE horizontal primitive running from the
  % word "orthogonal", through V_pure, out to the dotted drop line. That is the
  % glyph collision S11 exists to forbid, and it was the only one left in the
  % figure. The arms come 0.40 -> 0.28cm and the arm drops to y = 8.93, which
  % puts 1.05mm of white under the arrow; 2.8mm is still 40% larger than the
  % 2mm a cold reader did not see until zooming, and the quadrant is unchanged.
  % THE QUADRANT IS UPPER-LEFT and always has been: v10 item 8 says "lower-left"
  % and is false about its own drawing. The other three quadrants are ruled out
  % by measurement -- lower-left and upper-right are crossed by the raw band at
  % 38 degrees, and lower-right is where "the confound" and its leader sit.
  \draw[ink, line width=0.5pt] (10.13,8.65) -- (10.13,8.93) -- (10.41,8.93);
  \node[lab, anchor=base east] at (9.99,8.73) {orthogonal};

  % the four measured numbers of stage two, as one strip on the panel's own x
  % origin. These are the only quantities above panel (c) with a source.
  % v11, two repairs. (i) COLOUR: the strip was ink and panel (a)'s scope line
  % is quietink, on the SAME baseline in adjacent panels at the same size, with
  % nothing on the page decoding the difference; both are scope, so both are
  % quietink now and ink is left to heads, definitional prose, the axis label
  % and the gate. (ii) NAME: "subspace fraction" becomes fig:mechanism's own
  % name for the quantity, "descriptive subspace fraction" (one object, one
  % name). The layer-9 qualifier stays and is load-bearing: layer 9 is also the
  % MAXIMUM of variance_share.drug over the seven depths (the range is 0.277 at
  % layer 16 to 0.620 here), which the drawing has no lane to state and which
  % the caption and figs/purify.json now carry.
  \node[qlab] at (8.40,7.48)
    {Dose decodes from inside the raw slab at $0.90$--$0.95$;};
  \node[qlab] at (8.40,7.20)
    {descriptive subspace fraction at layer 9: raw $0.620$, purified $0.031$.};

  \node[def] at (8.40,6.62) {$V_{\textrm{pure}}$ is the part of the raw slab that is};
  \node[def] at (8.40,6.22) {orthogonal to the context subspace.};

\end{scope}

% ###########################################################################
% # PANEL (c) -- the removal gate.  x origin 1.60, the fraction axis's spine.
% # The only measured marks in the figure.
% ###########################################################################
  \node[head] at (1.60,5.20) {(c) the removal gate};

  % WHICH SLAB THE GATE IS MEASURED ON, restored to the drawing in v11 and set
  % in axis-label rank rather than at the foot. Without it the last line a
  % reader crosses before entering panel (c) is panel (b)'s "so causal claims
  % use the purified slab, not the raw", and the inference that the gate was
  % run on V_pure is unavoidable. workspace_probe.json does not record which
  % subspace was projected out; it stores one ablated series and no purified
  % counterpart, so the line states only what the file supports.

  % the region of outcomes that would have made every null in this section a
  % measurement of nothing. Drawn FIRST so every rule and dot sits on top of
  % it, at rulegrey!8 (grey 247) and not rulegrey!12 (grey 243): at 12 it was
  % within one grey level of panelbg, the tint that means "a schematic
  % subspace" 3cm above it in panel (b), so one tint carried two meanings.
  % Nothing is written inside it: its emptiness is the whole point, and the
  % two italic lines that used to sit in it competed with the marks. It is
  % NAMED instead, outside its own right edge and under the rule that bounds
  % it, because the largest mark on the page may not be the one mark a reader
  % cannot name (and because the name makes the third clause of this figure's
  % acceptance test answerable from the drawing).
  \fill[rulegrey!8] (1.60,1.75) rectangle (14.10,\fy{0.8});

  % ---- fraction axis, DEFINED AT THE AXIS ---------------------------------
  % Rotated 90 degrees reading upward, immediately left of the tick numerals,
  % which is fig:mechanism's form and the answer to the supervisor's first
  % complaint. Nothing else is in this block: no result, no interval note, no
  % cross-reference.
  % v11, AND THIS IS A CORRECTNESS REPAIR AND NOT A TYPOGRAPHIC ONE. v10 set
  % one \footnotesize line reading "fraction of drug signal removed" over one
  % \scriptsize construction line. Two faults, both fatal.
  %   (a) THE NAME WAS FALSE. The plotted series is frac_signal_removed, which
  %       is ABOVE-CHANCE signal: (before - after)/(before - chance). The
  %       unqualified quantity (before - after)/before is 0.7041 at layer 2 and
  %       0.8528 at layer 9, four to eighteen points BELOW the marks drawn. Read
  %       at the drawn name the top mark, 0.949, implies a residual accuracy of
  %       0.041, under the chance 0.083 this figure prints in its own foot note.
  %       The thesis's pre-registration (04b:218) says "the fraction of
  %       above-chance signal removed must exceed $0.8$". v10's stated reason for
  %       the omission -- that the 43-character form needs 7.14cm against a
  %       5.71cm lane -- ruled out one STRING, not the correct QUANTITY, and is
  %       superseded: the name is broken over two rotated \footnotesize lines.
  %   (b) THE SUBORDINATE LINE WAS THE TALLEST OBJECT IN THE PANEL. Measured on
  %       the v10 render, the \scriptsize construction line ran 5.52cm against a
  %       3.85cm axis (143%) and its top sat 6.0pt ABOVE the panel head's, so the
  %       smallest type in the block outranked the head. S5 makes that line
  %       optional; it is deleted, and the construction it carried moves to the
  %       foot note, where the three accuracies it is built from already sit and
  %       where the missing noun ("probe accuracy") costs nothing.
  % LANE, remeasured on the v11 render and reconciled with figs/purify.json,
  % which the v10 header disagreed with: \footnotesize sets at 0.158 cm per
  % character in this class, so "fraction of above-chance" (24 characters) runs
  % 3.86cm against a 3.00cm axis and "drug signal removed" (19) runs 3.03cm.
  % The block is centred on the axis midpoint, y = 3.25.
  \draw[axisr] (1.60,1.75) -- (1.60,4.75);
  \foreach \v/\t in {0/0, 0.5/0.5, 1/1.0}{
    \draw[furn] (1.60,\fy{\v}) -- (1.50,\fy{\v});
    \node[qlab, anchor=east, inner sep=0pt] at (1.44,\fy{\v}) {\t};}
  % LINE ORDER. Rotated 90 degrees the line-advance direction is -x, so the
  % FIRST line is the leftmost and the LAST line is the one against the tick
  % numerals, which is also matplotlib's own stacking for a two-line ylabel.
  % Set the other way round the block reads "drug signal removed / fraction of
  % above-chance", which is what the first v11 build printed.
  \node[text=ink, font=\footnotesize, rotate=90, anchor=center]
    at (0.36,3.25) {fraction of above-chance};
  \node[text=ink, font=\footnotesize, rotate=90, anchor=center]
    at (0.80,3.25) {drug signal removed};

  % ---- layer axis ---------------------------------------------------------
  % "layer" is left-aligned to the axis's own left end, per the contract, and
  % v11 puts 2.8mm of white under the tick numerals instead of v10's 0.63mm.
  % At 0.63mm the word sat directly under the numeral 2, one size LARGER than
  % it, and the pair read as "layer 2" -- the identical fault the pass-9 note
  % above records against the old placement under the layer-9 numeral. The
  % repair is vertical, not horizontal: the fraction origin rose 1.45 -> 1.75,
  % which is where the 2.8mm came from. Baseline to baseline the numerals and
  % the word are now 0.53cm apart, against the 0.40cm leading this figure uses
  % between two lines of one sentence, so they can no longer read as one block.
  \draw[axisr] (1.60,1.75) -- (14.10,1.75);
  \foreach \l in {2,4,6,8,9,12,16}{
    \draw[furn] (\lx{\l},1.75) -- (\lx{\l},1.65);
    \node[qlab, anchor=north, inner sep=0pt] at (\lx{\l},1.59) {\l};}
  \node[def] at (1.60,0.85) {layer};

  % ---- the gate: a pre-set threshold, so BLACK, SOLID, 0.9pt --------------
  % It was accent 0.6pt dashed. Accent in this document means the object under
  % test, and a threshold is not the object under test.
  \draw[gate] (1.60,\fy{0.8}) -- (14.10,\fy{0.8});
  \node[lab, anchor=west, text=black] at (14.25,\fy{0.8}) {gate, $0.8$};
  % the name of the shaded region, set outside it and directly under the rule
  % that bounds it. Quietink, because the region is a counterfactual and not
  % the object under test.
  \node[qlab, anchor=west] at (14.25,2.95) {an untestable null};

  % ---- the seven measured fractions. No connecting line: the series is flat.
  \gdot{2}{0.8846153846153847}
  \gdot{4}{0.9253112033195021}
  \gdot{6}{0.9383561643835617}
  \gdot{8}{0.9461077844311377}
  \gdot{9}{0.9491525423728814}
  \gdot{12}{0.9440993788819876}
  \gdot{16}{0.9444444444444445}
  % direct label for the set, on the last mark's own baseline, with the 0.30pt
  % leader the corpus requires beyond 3mm: the gap from the layer-16 dot to the
  % label's ink is 3.6mm, and in v10 that gap was open with nothing in it.
  \draw[accent, line width=0.3pt] (14.19,\fy{0.9444444444444445})
                              -- (13.95,\fy{0.9444444444444445});
  \node[alab, anchor=west] at (14.25,\fy{0.9444444444444445})
    {all seven layers: $88$--$95\%$};

% ###########################################################################
% # the figure's foot: one two-line \scriptsize block on the figure's own x
% # origin. Line one names the instrument and the three accuracies the axis is
% # built from; line two gives the construction and this figure's single
% # interval declaration.
% # v11: the noun was missing. v10 printed "Before 0.41--0.82, after 0.121 ...,
% # chance 0.083" with no quantity attached to any of the three, and the words
% # "accuracy", "probe" and "1/12" appeared nowhere in the figure, so panel (c)
% # never said what instrument it was reading. The registry name is "probe
% # accuracy" and chance is "0.083 = 1/12", which is the stored float exactly
% # (workspace_probe.json :: chance == 1/12 in IEEE double).
% ###########################################################################
  \node[qlab] at (0.00,0.32)
    {Probe accuracy before $0.41$--$0.82$, after $0.121$ at every layer,
     chance $0.083$, exactly $1/12$.};
  \node[qlab] at (0.00,0.02)
    {The fraction is (before minus after)$/$(before minus chance). Intervals:
     none returned in (c); (a) and (b) carry no measured mark.};

  % the bounding box IS the fullwidth measure; nothing may stick out of it
  \path (0,-0.10) rectangle (17.30,10.35);
\end{tikzpicture}
%
% ---------------------------------------------------------------------------
% READY TO PASTE into Sections/04b-representation.tex, at sec:res-used,
% immediately after the paragraph "Step three, the removal gate"
% (04b-representation.tex:182-190) and BEFORE tab:workspace. The same block is
% staged in figs/_PENDING_FLOATS.tex for a human to integrate.
%
% Do NOT merge this float with fig:hook. The two are read six paragraphs apart;
% each is about 100mm tall against a corpus ceiling of 117mm, so stacked they
% spill the page; and the house caption shape carries one bolded finding, not
% two.
%
% CAPTION, v10 AND v11. The drawing lost eight sentences in v10 and one more in
% v11, and every one is folded into the caption below, because in this document
% the argument lives in the caption and the drawing names objects.
% What was cut from the drawing, in the order it used to read:
%   1. the three \footnotesize italic panel subtitles -- "where the twelve drugs
%      differ, at one layer", "the raw slab is partly a batch slab", and "the
%      fraction of above-chance drug signal that projecting the drug subspace
%      out destroys". The third is now the axis label, which is where it always
%      belonged; the other two are caption clauses.
%   2. the centred divider sentence, "an ablation is only a test of the drug if
%      the slab it removes really holds the drug", which restated the body
%      paragraph directly above the float.
%   3. the two lines inside the fail region, "had any layer landed in here, the
%      null below it / would have been unfalsifiable --- a measurement of
%      nothing", which restated the caption's own "the shaded region ... is
%      empty" and sat inside the data area competing with the marks. v11 gives
%      the region the NAME "an untestable null", set outside its right edge;
%      the argument stays in the caption.
%   4. the centred discipline line, "the dots and the bands are schematic and
%      have no scale; only the four numbers in stage two are measured". Its job
%      is done by the GLYPH now -- open ring against filled disc -- and the
%      caption says the rest.
%   5. the conservatism argument, "so the slab keeps essentially the whole span,
%      noisiest included, which for a removal claim is conservative", cut from
%      panel (a)'s prose, which is now definitional only.
%   6. "by construction", cut from the right angle's label for want of a lane
%      3.6cm wide at the glyph's own height (v10 item 8).
%   7. "measured on the slab that is projected out; the stored file has no
%      separate purified-slab gate", cut from the foot note in v10 for text
%      density. RESTORED TO THE DRAWING IN v11, reworded, under panel (c)'s
%      head, because without it panel (c) named no slab at all and sat directly
%      under panel (b)'s "so causal claims use the purified slab, not the raw".
%      The caption's own statement of the same limit is reworded below so that
%      no clause of five or more words appears in both places.
%   8. "four \quad what the ablation costs: \cref{fig:hook}", the forward
%      pointer that used to sit on the x-tick baseline where the word "four"
%      read as a leftmost tick.
%   9. NEW IN v11: the axis's \scriptsize construction line, "(before minus
%      after)/(before minus chance)". It did not vanish -- it is the second foot
%      line now -- but it left the axis block, which S5 allows and which the
%      render forced: at 5.52cm against a 3.85cm axis it was the tallest object
%      in the panel and rose above the panel head's own top.
%
% CLAUSES OF THE PRINTED CAPTION THAT ARE NOW FALSE AGAINST THE DRAWING, with
% their lines in Sections/04b-representation.tex. Nobody may edit that file from
% this pass, so they are listed here and in figs/purify.json, and the
% replacement for each is in the caption below.
%   04b:234  "\emph{Above the rule, nothing has a scale.}" -- there is no rule.
%            The full-measure divider at y = 4.72 was deleted in v10 and the
%            panel letters carry the scoping now.
%   04b:234  "The dots, the bands, the arrow and the right angle are schematic"
%            -- v10 turned panel (a)'s twelve schematic marks into OPEN RINGS,
%            so the only dots in the figure are panel (c)'s seven MEASURED
%            discs, and the sentence as printed now calls the measurements
%            schematic. Replacement: "open rings and untitled bands".
%   04b:236  "the four numbers printed in stage two" -- there is no stage two;
%            the panel is lettered (b) and the word "stage" went with the
%            banners.
%   04b:243  "\emph{Below the rule, everything has a scale.}" -- same rule.
%   04b:250  "Two limits, both stated on the drawing" -- only one of the two is
%            on the drawing (the which-slab limit, restored in v11); the
%            twelve-class against twenty-four-class limit is caption-only and
%            always was.
%   04b:250-252  the which-slab limit as printed shares a clause of more than
%            five words with the drawing's new scope line, and it asserts more
%            than the file supports; it is reworded below.
% AND ONE THING THE CAPTION OWES THE READER AND HAS NEVER SAID: the printed
% pair "raw $0.620$, purified $0.031$" is at layer 9, and layer 9 is the
% MAXIMUM of layers.<L>.variance_share.drug over the seven depths. The stored
% range is 0.2766847391691891 at layer 16 to 0.6201207831136324 here, so the
% drawn pair is the largest raw-to-purified contrast the file holds. The raw
% series is drawn in no other figure and quoted in no sentence of the thesis
% (figs/fig-mechanism.json records it as drug_raw_not_drawn), so the caption is
% the reader's only chance to learn the spread. It is added below.
% The dose qualifier the v10 proposal added is kept: "$0.90$--$0.95$" is true
% only to two places, the stored minimum being 0.8979166666666666 at layers 12
% and 16, which is also why no reference rule is drawn at $0.90$.
%
% NOTHING IN THE CAPTION RESTATES A CLAUSE OF FIVE OR MORE WORDS FROM THE
% DRAWING. Checked string by string against the v11 render's own text spans: the
% drawing names objects and prints numbers, the caption argues. The drawing says
% "the subspace projected out"; the caption says "whichever subspace the
% ablation removed", and the longest run shared with any drawn string is four
% words.
%
% Two caption details, checked against a build of the whole float (returncode 0,
% one page, zero Overfull/Underfull boxes; the \cref targets resolve in the
% document and show as ?? only in an isolated harness):
%   * the two numerals in the bold lead are set in TEXT and not as $88$/$95$.
%     In math mode they come out upright and unbolded inside \textbf, so a
%     bolded finding would have had two unbolded numbers in the middle of it.
%   * the FIGURE BODY carries no \cref. The "four ... fig:hook" pointer was
%     deleted in v10, so the caption is the only place the sibling float is
%     named.
%
% \begin{figure}[htbp]
% \begin{fullwidth}
% \centering
% % ===========================================================================
% figs/purify.tex -- fig:purify, how the drug slab is BUILT, and the gate it
% has to pass before anything may be ablated with it.
%
% FIGURE BODY ONLY. No preamble, no document environment, no float, no caption.
% \input it from inside a figure/fullwidth float in
% Sections/04b-representation.tex; the ready-to-paste float sits commented at
% the foot of this file and is also staged in figs/_PENDING_FLOATS.tex.
% Compiles against thesis_v2/preamble.tex and nothing else: no \includegraphics,
% no external data, no TikZ library beyond the ones preamble.tex already loads
% (this file uses arrows.meta and calc only).
%
% WHAT THE FIGURE IS
% ------------------
% Steps one to three of sec:res-used, drawn as steps. Step one localises the
% drug: twelve centred class means, ten kept directions. Step two purifies it:
% the raw slab is orthogonalised off the context it is confounded with. Step
% three is the falsifiability check: projecting the slab out must destroy most
% of the decodable drug signal, or every null downstream is a measurement of
% nothing. fig:hook (a separate float, six paragraphs later) draws step four;
% the two are deliberately NOT merged -- each is about 100mm tall, the corpus
% ceiling is slab.tex at 117mm, and the house caption shape carries one bolded
% finding, not two.
%
% EVERY NUMBER DRAWN, AND WHERE IT COMES FROM
% -------------------------------------------
% Source for every measured quantity: RESULTS_cluster/workspace_probe.json.
% Config there: tier tier2_unseen_drugs, n_drugs 12, n_per_drug 40, n_dims 10,
% layers 2,4,6,8,9,12,16, chance 0.08333333333333333, seed 42.
% Sibling ledger with the full machine-precision values: figs/purify.json.
%
% (a) THE SEVEN PLOTTED DOTS -- the only marks in this figure with a scale.
%     layers.<L>.frac_signal_removed, layers 2,4,6,8,9,12,16:
%       0.8846153846153847  0.9253112033195021  0.9383561643835617
%       0.9461077844311377  0.9491525423728814  0.9440993788819876
%       0.9444444444444445
%     layers.<L>.gate_pass is true at all seven; the dots' position above the
%     0.8 rule is that fact, so no second mark asserts it.
%
% (b) THE GATE, 0.8. PROSE ONLY. Sections/04b-representation.tex:185 "the
%     fraction of above-chance signal removed must exceed $0.8$". There is no
%     config field for it in workspace_probe.json, so the ledger cites the .tex
%     line and not a key path. Drawn dashed accent, the house register for a
%     pre-set threshold a value is tested against (gate.tex:170-174).
%     SUPERSEDED IN v10 BELOW: it is now black, solid, 0.9pt.
%
% (c) THE THREE PRINTED ACCURACIES in the right-hand block of band two.
%       before  0.41--0.82   layers.2.drug_decode  = 0.4083333333333334
%                            layers.9.drug_decode  = 0.8208333333333334
%       after   0.121        layers.<L>.drug_decode_ablated
%                            = 0.12083333333333332, byte-identical at all seven
%       chance  0.083        top-level chance = 0.08333333333333333
%     Thesis 04b:186-188 "Accuracy falls from $0.41$--$0.82$ to a constant
%     $0.121$ floor, removing $88$--$95\%$ of the above-chance signal".
%     SUPERSEDED IN v10 BELOW as to PLACEMENT only: the three values are no
%     longer a right-hand gutter block, they are the figure's foot note. No
%     value changed.
%
% (d) THE IDENTITY, verified numerically before drawing, in python, over all
%     seven layers, exact float equality (not a tolerance):
%       (drug_decode - drug_decode_ablated) / (drug_decode - chance)
%         == frac_signal_removed
%     Layer 2: (0.4083333333333334 - 0.12083333333333332)
%              / (0.4083333333333334 - 0.08333333333333333)
%              = 0.8846153846153847, bit for bit.
%     WHICH SLAB THE GATE IS MEASURED ON. This identity fixes it, and the
%     answer is NOT V_pure. workspace_probe.json stores exactly one ablated
%     series, drug_decode_ablated, and the identity reproduces from the RAW
%     decode. The section-4.5 arm files prove the pipeline can distinguish the
%     two cases -- they store drug_decode_ablated AND drug_decode_pure_ablated,
%     frac_signal_removed AND frac_signal_removed_pure -- and this file stores
%     only the first of each pair. So the figure, the caption and the ledger all
%     say "the slab that is projected out" and nowhere claim the gate was run on
%     V_pure. Stage two's arrow lands on V_pure because that is where the CAUSAL
%     claims are made (04b:179-180); band two is a separate measurement and the
%     divider is what keeps the two from being read as one.
%     SUPERSEDED IN v10 BELOW as to the LAST CLAUSE ONLY: the divider is gone
%     and the panel letters do that work now. The caveat itself survives, in the
%     foot note.
%
% (e) THE FOUR PRINTED NUMBERS IN STAGE TWO, the only measured quantities above
%     the divider.
%       0.90--0.95   layers.<L>.confound_crossdecode.dose, seven layers:
%                    0.9020833333333333 0.9229166666666666 0.9416666666666667
%                    0.9395833333333332 0.95 0.8979166666666666
%                    0.8979166666666668. Thesis 04b:176-178.
%                    PRINTED AS TEXT AND NOT DRAWN, and there is deliberately NO
%                    0.90 reference rule anywhere in this figure: the stored
%                    minimum is 0.8979166666666666 at layers 12 and 16 and
%                    reaches the thesis's "0.90" only by rounding to two places.
%                    A rule at 0.90 would put two of seven points below it and
%                    would read as a failure the thesis does not report.
%       raw 0.620    layers.9.variance_share.drug      = 0.6201207831136324
%       purified     layers.9.variance_share.drug_pure = 0.030620247335072275
%         0.031
%     The RAW share is quoted in no sentence of the thesis; figs/fig-mechanism.json
%     already records it as drug_raw_not_drawn. Printing it here is a new reading
%     of a stored field, flagged as such in purify.json, on the precedent of
%     figs/slab.json's handling of variance_share.context. Layer 9 is named in
%     the figure text so the pair cannot be read as a seven-layer summary; it is
%     the layer at which the raw probe peaks (04b:170-171), which is why the
%     section uses it as its own reference depth.
%
% (f) THE COUNTS in stage one: 480 prompts, twelve drugs, forty cells each,
%     seven layers, ten kept directions, span at most eleven.
%     config.n_drugs 12, config.n_per_drug 40, config.n_dims 10,
%     layers.<L>.pure_drug_dims 10, config.layers "2,4,6,8,9,12,16".
%     480 = 12 x 40 and the rank-11 bound (twelve centred means span at most
%     eleven directions) are ARITHMETIC, not stored fields; the ledger says so.
%     Thesis 04b:164-171.
%
% AGREEMENT WITH THE TEXT (checked line by line, not assumed):
%   04b:164-165  "\num{480} tier-2 prompts (twelve drugs, forty real cells
%                each) at seven layers"            -> stage one, line 4.     OK
%   04b:168-170  "Twelve centred class means span at most eleven directions,
%                so the slab keeps essentially the whole span, noisiest
%                included, which for a removal claim is conservative"
%                -> stage one, lines 1-3, verbatim.                          OK
%   04b:175-178  "The raw drug slab is partly a batch slab. Drugs are
%                confounded with plate and dose by the atlas's design, and
%                probing dose from inside the drug subspace recovers it at
%                $0.90$--$0.95$ at every layer"    -> stage two, lines 1-3.  OK
%   04b:179-180  "Every causal claim below is made on $V_{\textrm{pure}}$,
%                never on the raw slab."           -> the V_pure callout.    OK
%   04b:185-186  "the fraction of above-chance signal removed must exceed
%                $0.8$; below that the null is unfalsifiable, a measurement
%                of nothing"                       -> the dashed rule, and
%                                                     G4's two italic lines,
%                                                     which are the section's
%                                                     own words.             OK
%   04b:186-188  "Accuracy falls from $0.41$--$0.82$ to a constant $0.121$
%                floor, removing $88$--$95\%$"     -> right-hand block and
%                                                     the set label.         OK
%   04b:212-213  tab:workspace caption "All seven layers pass the removal gate
%                (88--95\% ...); chance is $0.083$"                          OK
%
% AVAILABLE, DELIBERATELY NOT DRAWN, recorded here and in purify.json so a
% reviewer who asks for direct evidence of the batch overlap is answered
% without a redraw:
%   * layers.<L>.overlaps.context   0.5215396547181934 -- 0.6323768458353303
%   * layers.<L>.overlaps.cell_line 0.3043896441494730 -- 0.3534060289850234
%     These are the per-layer measure of "the raw drug slab is partly a batch
%     slab" and are quoted nowhere in the thesis. Drawing them would put a
%     second measured scale into a panel whose whole discipline line says the
%     geometry has no scale, and would be a new claim rather than a redraw.
%   * layers.<L>.confound_crossdecode.plate, 0.4541666666666667 --
%     0.5770833333333333, a materially weaker and entirely undiscussed
%     confound. Not printed, because the section names dose and not plate.
%   * layers.<L>.swap.{drug_b_kl, random_inject_kl} in this same JSON is the
%     RETIRED steering design (07-Appendix.tex, sec:app-identity). Nothing from
%     it is on this figure. figs/mechanism.py's header records that the v1
%     figure plotted exactly that pair while its caption described the ablation.
%
% DESIGN DECISIONS, each with its reason
% --------------------------------------
%  1. TWO BANDS AND A DIVIDER THAT IS A STATEMENT. Above the rule at y=4.72
%     nothing has a scale; below it everything does. No mark crosses the rule,
%     no leader connects a schematic object to a plotted one, and the italic
%     sentence under the rule sets band two up as the answer to the objection
%     stage two raises -- if purification strips the raw slab from 0.620 to
%     0.031 of the stream's variance, does anything survive in it? The
%     alternative, one continuous diagram from cloud to axis, was rejected: it
%     would invite the reader to measure the ellipse.
%     SUPERSEDED IN v10 BELOW. The divider and its sentence are deleted; the
%     schematic/measured distinction is now carried by the MARK (open rulegrey
%     ring against filled accent disc), which is where it belongs.
%  2. THE FAIL REGION IS DRAWN, THE PASS REGION IS NOT. The one non-obvious
%     mark in band two is the rulegrey!18 rectangle from 0 to 0.8. It is the
%     region of outcomes that would have made every null in this section a
%     measurement of nothing, and it is empty. A figure that drew only the seven
%     dots and the rule would be a bar of good news; drawing the region the
%     dots had to avoid is what makes it a falsifiability picture. Its emptiness
%     is also where the two italic lines fit, which is not a coincidence and is
%     the reason they are placed there rather than in the caption.
%     PARTLY SUPERSEDED IN v10: the region survives, at rulegrey!12, and the two
%     italic lines inside it do not. The emptiness is the message and text in it
%     was competing with the marks.
%  3. NO CONNECTING LINE THROUGH THE SEVEN DOTS. The series is flat to 0.1355cm
%     on this scale (0.8846 to 0.9492 at 2.10cm per unit of fraction). A line
%     would assert a trend across a 6-point spread the axis cannot resolve; the
%     seven marks are named as one ordered set by "all seven layers: 88--95%"
%     instead, which is comparators.tex's rule for marks too close to leader
%     individually.
%  4. NO INTERVALS ANYWHERE. frac_signal_removed is a single ratio of two
%     accuracies on 480 prompts and carries no stored dispersion of any kind;
%     the KL series in the same file carry normal-approximation SEs, which are
%     not this document's interval convention (figs/fig-mechanism.json records
%     intervals_drawn false with that reason). A bare solid dot with no bar is
%     how this document says a quantity has no interval (comparators.tex:149-151
%     changed a hollow ring to a solid dot for exactly this).
%  5. THE ORTHOGONALISATION IS DRAWN AS AN OPERATION, NOT AS A RESULT, AND THE
%     OPERATION IS A PROJECTION AND NOT A ROTATION. This is the one mark
%     figs/slab.tex does not have. There the drug wedge and the
%     context wedge are simply shown already orthogonal, which states the
%     outcome and hides the step; here the raw band sits at 38 degrees, visibly
%     crossing the context band, and the raw slab's tip is carried onto V_pure.
%     REPAIRED, PASS 9. It used to be carried by an accent ARC from 38 to 76
%     degrees, i.e. the raw band swung upright. A reviewer read that panel as
%     "the slab is turned until it is perpendicular to context" and had to go to
%     the prose to learn that a component is REMOVED, and they were right to:
%     two bands of equal width (0.32cm) and near-equal length (1.90 against
%     2.00) differing only in angle say "same slab, rotated, nothing lost",
%     against measured numbers 1.8cm to the right that say raw 0.620, purified
%     0.031, a twentyfold shrink. A curved arrow is the universal sign for
%     rotation. What replaces it is the subtraction itself: a dashed rulegrey
%     perpendicular from the raw band's tip down to the context band, the
%     component it lands on drawn as its own named mark, and a STRAIGHT accent
%     arrow at the tip's own height carrying the tip left by exactly that
%     component, landing on V_pure. Two further faults the arc had, both from
%     the render: its head stopped 1.5mm short of the V_pure bar and pointed at
%     white space, and the grey leader from the word "orthogonalise" stopped 4mm
%     short of the arc's tail -- the corpus puts the 0.06cm gap at the LABEL
%     end, never at the mark end, so that stub read as a stray hairline. The
%     label now sits 0.15cm from the arrow's tail and carries no leader, and the
%     arrow is 0.6pt: a schematic operation, not a measurement, where 1.0pt is
%     the corpus's "this IS the claim" register. V_pure was also narrowed from
%     0.32 to 0.18cm, so the purified band is visibly thinner than the raw one.
%     slab.tex is left alone and is not
%     superseded: three of its four panels are the failed identity designs of
%     sec:app-identity, which are the appendix's argument, and its magnified
%     interior is sec:res-mechanism's swap geometry. What it draws that belongs
%     to sec:res-used is one thing, the variance-share disc, and this figure
%     does not redraw it. The two figures share vocabulary (accent-on-white
%     bands, rulegrey leaders, one discipline line saying what is to scale) and
%     no coordinate.
%     The whole of this decision survives v10 except the arrow's WEIGHT, which
%     goes 0.6 -> 0.5pt to bring the file inside a four-weight budget; the
%     perpendicular goes dashed-rulegrey -> densely-dotted-quietink, which is
%     this document's reserved stroke for a projection line.
%  6. THE CONFOUND IS A NAMED MARK, NOT THE PRODUCT OF TWO TINTS.
%     REPAIRED, PASS 9, and this is the repair both reviewers called blocking.
%     It used to be the lozenge where the raw band's tint multiplied over the
%     context band's, with the whole claim carried by the caption ("that
%     darkened region is the confound step two exists to remove"). Three faults:
%     (a) it was the only element in the panel with no direct label, while every
%     other mark had one; (b) V_pure ran vertically through its centre and split
%     it in two -- sampled on the render, a row across it read grey (174,174,174),
%     then (119,143,164) for the tint product, then a hard cut to opaque accent
%     (11,82,147) for 24px, then (119,143,164) again -- so it was not a nameable
%     region at all; (c) it made a CROSSING mean "confounded" on a panel where
%     another crossing, V_pure through the same band, has to mean "orthogonal".
%     One glyph cannot carry both. It is now the orthogonal projection of the
%     raw direction onto context: a segment along the context band from the
%     shared origin to the foot of the perpendicular dropped from the raw band's
%     tip, drawn in the raw slab's own tint at a higher opacity with a 0.35pt
%     accent!70 outline, and directly labelled "the confound". It begins at
%     V_pure's right edge, so nothing bisects it, and it is exactly the length
%     of the removal arrow above it, which is the rhyme that makes the arrow
%     legible. With the tint-multiply claim gone, a crossing means nothing
%     anywhere in the panel and the right-angle mark is the only assertion of
%     orthogonality, which is what decision 8 already wanted.
%     The context band is now panelbg with a rulegrey outline and not ink!35:
%     see decision 11.
%  7. NO 0.90 REFERENCE RULE. See (e) above: the stored dose minimum is
%     0.8979166666666666 and rounds to the thesis's 0.90. Printing the range as
%     text states exactly what the section states; a rule would show two of
%     seven points below it and would manufacture a failure.
%  8. RIGHT ANGLE AS A MARK, NOT AS A WORD. The small square in the
%     upper-left quadrant at the shared origin is the only assertion that
%     V_pure and the context subspace share no direction. The word
%     "orthogonalise" labels the operation; the square labels its result.
%     PASS 9: ink 0.35pt, up from rulegrey 0.3pt. The caption calls it "the one
%     geometric fact true by construction" and it was the faintest mark in the
%     panel.
%     AMENDED IN v10: the square moves to the LOWER-left quadrant, doubles to
%     0.40cm arms, and acquires the direct label it never had.
%     SUPERSEDED IN v11: THIS LINE IS FALSE ABOUT ITS OWN DRAWING AND ALWAYS
%     WAS. v10 drew the square in the UPPER-left quadrant of the crossing --
%     (10.01,8.65)-(10.01,9.05)-(10.41,9.05) -- and no build has ever put it
%     lower-left. The direct label and the enlargement are true. v11 keeps the
%     upper-left quadrant, which is the only one of the four that is free (the
%     raw band at 38 degrees crosses lower-left and upper-right, and "the
%     confound" and its leader own lower-right), and takes the arms 0.40 ->
%     0.28cm to clear the removal arrow -- see v11 item 2.
%  9. THE GLOBAL MEAN IS A CROSS, NOT A DOT. The twelve class means are dots;
%     the thing they are centred on must not be a thirteenth dot. It is also
%     leadered out of the cloud rather than labelled in place, because at
%     8pt the label box collided with two of the twelve dots.
% 10. V_drug IS DRAWN WITH STAGE TWO'S OWN GLYPH, AND ALL TWELVE DISPLACEMENTS
%     ARE DRAWN. REPAIRED, PASS 9, two findings at once.
%     (a) V_drug used to be a dashed, accent!10 FILLED ELLIPSE enclosing the
%     twelve dots, while stage two drew the same name as an unbounded band. A
%     bounded blob around the class means says "the region where the drugs
%     live"; a subspace is not a region and has no edge. Read cold, the ellipse
%     was taken for the data cloud and stage two's band for a new object, so
%     stage one's output was never connected to stage two's input -- the one
%     join the figure exists to make, and the join the brief says readers lose.
%     V_drug is now an accent band at the SAME fill opacity stage two uses,
%     through the global mean, running past the data on both sides. The ellipse
%     is deleted rather than kept unlabelled: with the band tall enough to
%     contain all twelve dots (half-height 0.52 against a maximum |dy| of 0.44)
%     it had nothing left to say, and a class mean drawn OUTSIDE V_drug would
%     have been false, since the twelve centred means span it by construction.
%     (b) ONE displacement was drawn and eleven were not, so "twelve are
%     stacked" -- the sentence directly under the dots -- was invisible, and the
%     dots were unlabelled while the prose two lines below names a different
%     population ("480 prompts: twelve drugs, forty real cells each"), so the
%     scatter read as prompts and the cross as their grand mean. All twelve
%     stems are now drawn at 0.3pt quietink, with the one accent arrow on top
%     and still the only labelled one, and the dots carry the tag "the twelve
%     per-drug means". The stems make the count a picture and they are what the
%     "at most eleven directions" line is about.
%     The twelve offsets are written out in a \foreach list so no build can
%     randomise them and no reader can be told a different story by a different
%     compile. Opacity: the band is 0.20, not the 0.34 of stage two's first
%     draft, and stage two's raw band was brought to 0.20 with it. At 0.34 over
%     5.10 x 1.04cm the band became the single heaviest area on the page, which
%     is the thing this figure was already being marked down for; at 0.20 over
%     4.10 x 1.04 it lays down about the ink of stage two's 1.90 x 0.32 band and
%     the two are drawn at the SAME opacity, so the glyph rhymes exactly.
%     AMENDED IN v10: the twelve class means become OPEN RINGS and the one
%     accent displacement arrow loses its head -- see v10 items 3 and 6.
% 11. ONE GREY FAMILY, AND THE CONTEXT SUBSPACE IS THE COLOUR slab.tex GIVES IT.
%     PASS 9. Three greys were in play, derived two different ways: rules at
%     rulegrey, the context band at ink!35, the gate region at rulegrey!18.
%     ink!35 is not derived from rulegrey and is LIGHTER than a hairline rule,
%     so the band read as a heavy solid slab while the rules read as thin, and
%     it was the darkest region on either page. slab.tex draws this same named
%     subspace in panelbg (its \slabdisc, slab.tex:83-90), so the figure also
%     disagreed with the published figure about what colour the object is. The
%     band is now panelbg with a 0.35pt rulegrey outline; ink!35 no longer
%     appears in the document, and purify's context band and slab.tex's context
%     wedge are the same object drawn twice.
% 12. THE GATE PANEL KEEPS ITS FULL 0-TO-1 AXIS. A reviewer asked for the range
%     to be cut to 0.75-1.00 so the seven fractions (0.8846 to 0.9492, 0.1355cm
%     apart on this scale) would resolve, against a fail region 12.4x their
%     height. OVERRULED, and the reason is the same one that bans bars on a log
%     axis: the fail region's height is not a drawing choice, it is the share of
%     the outcome space that would have made every null below unfalsifiable, and
%     0 is the natural floor of a fraction, not an arbitrary cut. Cutting the
%     axis at 0.75 would shrink the one region whose SIZE is the argument, to
%     resolve an ordering the thesis makes no claim about -- tab:workspace is
%     the authority for the per-layer values and the section quotes only the
%     range. What was applied instead: the direct label "all seven layers:
%     88--95%" names the set (decision 3), and the 0.8 numeral was dropped from
%     the tick set because the accent label already owns that coordinate.
%     STILL IN FORCE IN v10, and it is the first thing v10 was told to protect.
%
% WHAT SIX RENDER-AND-READ PASSES CHANGED, with the measurement that forced it
% ---------------------------------------------------------------------------
%  i.   PASS 1, Overfull 88.29pt. The discipline line under stage two ran 146
%       characters at \footnotesize and measured 20.4cm against a 17.30cm
%       measure. Cut to 103 characters; the sentence now states only what it
%       has to, that the geometry has no scale and four numbers do.
%  ii.  PASS 1. "every causal claim is made here", set anchor=south at y=9.22,
%       put its box top at 9.63 against the stage-two subtitle's ink bottom of
%       9.37 and the two overprinted. There is no room above V_pure for a
%       second line, so the sentence became the fourth line of stage two's
%       prose block, in the section's own words ("never on the raw slab").
%  iii. PASS 1. "the global mean" and "$\mu_d-\bar\mu$, one drug's" collided
%       character to character: both wanted the space up and right of the
%       cloud, which is where the displacement arrow points. The global mean
%       was moved BELOW the cloud on a leader. Drawn straight down that leader
%       read as a dropped axis, so PASS 4 gave it a 0.25cm dogleg; the dogleg
%       clears the two nearest class means by 0.23cm and 0.25cm.
%  iv.  PASS 2, Overfull 6.20pt, unchanged from PASS 1 by three edits, which is
%       what identified it: the culprit was not the widest LINE but the
%       fraction axis's own tick label. "1.0" is anchor=east on a spine at
%       x=0.30, and its box ran 0.28cm LEFT of x=0. The bare \path pins a
%       minimum box, not a maximum, so nothing clipped it. The whole plot was
%       shifted right (spine 0.30 -> 0.75, layer 2 from 0.60 -> 1.05) rather
%       than the label being clipped or shrunk. boxes=0 from PASS 4 on.
%  v.   PASS 2. Band two carried two annotation columns 0.30cm apart at the
%       same baseline, so "88--95%" and "accuracy before" read as one line. The
%       plot was widened from 8.60 to 11.00cm of layer axis and the two columns
%       became one, which also removed 3.4 x 1.9cm of dead white under them.
%  vi.  PASS 4, at 520dpi. The leader to "orthogonalise" left the arc 0.20cm
%       from its own arrowhead and read as a grey continuation of the accent
%       arrow. It now leaves at 50 degrees, 0.42cm clear of the head, and stays
%       above the raw band's upper edge for its whole length. The right-angle
%       mark went from 0.16 to 0.20cm, and the global-mean cross from 0.30 to
%       0.45pt with 2mm arms, because at 0.3pt it was fainter than the leaders.
%  vii. PASS 5. Stage one left 4.2 x 1.9cm of white to the right of the cloud
%       while stage two's matching region was full. The ellipse went from
%       1.75 to 2.05 x-radius and the cloud from 1.45 to 1.72, which both fills
%       the panel and raises the aspect from 2.34:1 to 2.77:1. The ellipse's
%       y-radius was NOT raised with it: at 0.98 the ellipse reached y=6.92 and
%       left the global-mean label 0.06cm from the prose block below.
%  viii.PASS 5. The panel-three subtitle said "projecting the slab out", which
%       after stage two's "every causal claim is made on V_pure" invited the
%       reading that the gate was run on V_pure. It now says "projecting the
%       drug subspace out", the section's own noun at 04b:184-185, and the
%       string V_pure does not appear in panel three at all.
%
% LAYOUT NOTES, all forced by a build
% -----------------------------------
%   * fullwidth = textwidth 142mm + marginparwidth 28mm + marginparsep 4mm
%     = 174mm. (The "171mm" beside \newenvironment{fullwidth} in preamble.tex
%     is stale arithmetic.) The bounding box is pinned with a bare \path at
%     (0,-0.12) rectangle (17.30,10.20) = 173.0 x 103.2mm, inside the 17.35cm
%     hard edge circularity.tex records and under the corpus height ceiling,
%     slab.tex at 117mm. The harness reports pages=1, boxes=0. It grew 1.7mm in
%     pass 9: band two moved down 0.20cm to open a line under stage three's
%     subtitle for the "which slab" caveat, and the "four ... see fig:hook"
%     pointer sits on the bottom line. The float got much SHORTER overall in the
%     same pass, because the caption went from 484 words to about 280.
%   * x=1cm, y=1cm. Nothing is scaled, nothing is \resizebox'd; 1cm here is
%     1cm on the page.
%   * pgfmath: \lx and \fy expand to BRACED groups, so \lx{2}+0.05 is not a
%     coordinate. Every stand-off is either inside the braces or an xshift on
%     the node (gate.tex records the same trap).
%   * the tick rows use the \v/\t value/label pair form, so the fraction axis
%     prints 0.5 and 1.0 rather than \foreach's raw 0.5 and 1 tokens. The layer
%     row does not need it: its values are integers and print as themselves.
%   * no node uses align= or text width=: such a node builds a minipage whose
%     first-line height follows the AMBIENT \baselineskip, so its ink moves
%     with whatever prose precedes the float (frames.tex measured 3.7pt of
%     drift). Every multi-line block here is a stack of single
%     anchor=base west lines on explicit baselines: 0.28cm apart in band one's
%     two prose blocks, 0.33cm in band two's right-hand block, where the lines
%     are read one at a time rather than as a paragraph.
%   * the fraction axis carries no rotated title. The panel subtitle sits
%     directly above it and names the quantity, the direction and the frame in
%     one line, which is what gate.tex's fix 1 asks an axis to do; a second
%     title would repeat it 0.2cm away.
%     THIS NOTE IS NOW FALSE AND IS SUPERSEDED BY v10 ITEM 4. It was the
%     supervisor's first complaint about this arc of figures and it was wrong:
%     a quantity named in a panel subtitle 30mm from the ticks is not named at
%     the axis. The fraction axis carries a rotated title from v10 on.
%   * PASS 9, THE "WHICH SLAB" CAVEAT, three builds to place. The wording went
%     155 -> 121 -> 92 characters. At 155 the picture was 55.72pt over the
%     measure and the harness said so; at 121 it ran to x = 15.90 and
%     overprinted the accent label "all seven layers: 88--95%" at 12.60; at 92
%     it ends at 12.30. Vertically it was first set at baseline 3.24 and its
%     ascenders touched the subtitle's descenders with no white between -- the
%     subtitle is \footnotesize in a node whose TOP is pinned at 3.76, so its
%     descender line is 3.351 and a \scriptsize line at 3.24 reaches 3.444. It
%     is now at baseline 3.05 with 0.10cm of clearance above, and INDENTED to
%     x = 0.90: at x = 0 its descenders at 2.97 met the "1.0" tick numeral,
%     whose box top is 2.95, and the indent both clears the numeral and breaks
%     the note away from the subtitle so the two do not read as one paragraph.
%   * PASS 9, THE DIVIDER BETWEEN STAGES ONE AND TWO now ends at y = 5.28, the
%     descender line of the last prose row in both columns (baseline 5.38),
%     instead of at 5.40 where it stopped between two lines of type and read as
%     a broken rule. It clears the centred italic beneath it by 0.26cm.
%   * PASS 9, THE DISCIPLINE LINE IS CENTRED on the figure's own centre and is
%     no longer flush left at the prose column's x. Flush left, at the same
%     leading as the four-line block directly above it and with no blank
%     scan-line between, it parsed as that paragraph's fifth line even though it
%     ran full width past the vertical divider. Centred, it pairs with the
%     sentence centred under the rule, and the two bracket the divider.
%   * PASS 9, THE AXIS TITLE moved from x = 6.60 to 7.73. \lx{9} = 6.55, so
%     centred on 6.60 the word "layer" sat within 0.05cm of the layer-9 numeral,
%     one line below it, and read as a label of layer 9 rather than of the axis.
%     7.73 is the midpoint of the widest tick gap (\lx{9} = 6.55 to
%     \lx{12} = 8.91) and is 0.83cm clear of both.
%   * PASS 9, THE 0.8 TICK IS GONE. The coordinate was labelled twice at the
%     same height, once as a quietink numeral on the left spine and once as the
%     accent label "the gate: 0.8" on the right, so the plot appeared to carry
%     a left scale and a right scale. 0, 0.5, 1.0 is also a regular tick set.
%     A 0.9 tick was considered and rejected by measurement: at \fy it would
%     sit 0.21cm from both neighbours against a 0.203cm numeral height.
%
% TYPE. Ambient \small on the picture; \scriptsize (8pt) on every node except
% the \footnotesize (10pt) italic prose -- the three panel subtitles, the
% discipline line, the divider sentence and the two lines inside the fail
% region. Nothing is set at body size.
% SUPERSEDED IN v10 BELOW: two sizes only, \footnotesize and \scriptsize, and
% ZERO italic.
%
% ===========================================================================
% v10 -- 2026-08-18. REBUILD AGAINST THE STYLE CONTRACT AND THE SUPERVISOR'S
% VERDICT. Nine of the twelve design decisions above survive whole; the three
% that do not are superseded in place, above, rather than deleted, because the
% thesis is audited against this header.
%
% THE VERDICT THIS PASS ANSWERS, quoted: "For the other figures 4.8-12 go back
% over them and iron them out even more. Make it look more professional, right
% now it looks a bit sloppy and childish ... and moreover make sure that the
% figures are CORRECT and display the idea we want to convey in the best way
% possible." Plus, on 4.7 but binding on the whole arc: "there is no definition
% on the x axis for what we are measuring".
%
% NOT ONE PLOTTED OR PRINTED VALUE CHANGED IN THIS PASS. The seven gate
% fractions are the same seven floats, unrounded; 0.8, 0.41--0.82, 0.121,
% 0.083, 0.90--0.95, 0.620, 0.031, 480, twelve, forty, ten, eleven and seven
% are all unmoved. What changed is the y map (0.75 + f*2.10 -> 1.45 + f*3.95)
% and the x map (1.05 + (L-2)/14*11.00 -> 1.45 + (L-2)/14*11.75), both recorded
% in figs/purify.json beside the old ones. A map change is a position change,
% not a value change, and the ledger carries both maps so the audit can redo it.
% SUPERSEDED IN v11, AND THESE TWO SENTENCES CARRIED FOUR WRONG COEFFICIENTS.
% What v10 actually shipped, read off its own code and confirmed by measuring
% its render's dot centroids, was \fy = 1.45 + f*3.85 (not f*3.95) and
% \lx = 1.85 + (L-2)/14*12.00 (not 1.45 + (L-2)/14*11.75). figs/purify.json's
% value_to_position_maps.current_v10 held the correct pair all along, so the
% header and its own sibling ledger disagreed; under the header's numbers an
% auditor redoing the drawing would have put layer 16 at 13.20cm against the
% 13.85cm drawn and the gate at 4.61cm against 4.53cm. The v11 maps are
% recorded beside the code at the top of the picture, which is the only place
% they can be checked against what is drawn. No value changed in either pass.
%
%  1. THE THREE ALL-CAPS SANS ACCENT BANNERS ARE GONE, replaced by fig:mechanism's
%     own panel-head form: "(a) localise the slab", "(b) purify it against
%     context", "(c) the removal gate", \footnotesize roman ink, sentence case,
%     each left-aligned to its panel's x origin. fig:mechanism (4.11) carries
%     zero banners and sets "(a) encoded more, used no more" at 10.0pt roman
%     ink; figs/slab.tex (A.3) carries nine banners, but A.3 is an appendix
%     reference plate a reader stops at and this is read inside a running
%     argument, where the caps sans accent face belongs to the margin
%     apparatus (\marginlabel) and not to the figure. The ordinal now lives in
%     the panel letter, so the enumeration "one, two, three" that made a reader
%     hunt for a missing stage four is gone with it, and so is the
%     "four \quad what the ablation costs: \cref{fig:hook}" pointer that closed
%     it -- a cross-reference in a drawing renders as ?? in the harness and on
%     the page competed with the x-tick row it shared a baseline with, where
%     the word "four" read as a leftmost tick.
%  2. EVERY ITALIC IS GONE, AND WITH IT EVERY SENTENCE THAT WAS ARGUING RATHER
%     THAN NAMING. Deleted: the three \footnotesize italic panel subtitles;
%     the centred divider sentence "an ablation is only a test of the drug if
%     the slab it removes really holds the drug", which restates the body
%     paragraph directly above the float; the two lines inside the fail region,
%     "had any layer landed in here, the null below it / would have been
%     unfalsifiable --- a measurement of nothing", which restate the caption's
%     own "the shaded region ... is empty"; and the centred discipline line
%     "the dots and the bands are schematic and have no scale; only the four
%     numbers in stage two are measured". Italic went from 31% of the figure's
%     characters -- the highest in the five-figure set, against 0% in 4.11 --
%     to nil. All five sentences are staged in figs/_PENDING_CAPTIONS.tex under
%     \label{fig:purify} for a human to fold into the caption, which is where
%     this document keeps its argument; none is silently lost.
%  3. THE SCHEMATIC DISCLAIMER IS NOW CARRIED BY THE GLYPH AND NOT BY A
%     FOOTNOTE. The panel drew thirteen SCHEMATIC dots at the same radius and
%     the same fill as panel (c)'s seven MEASURED ones and then apologised for
%     the collision in an 8pt line. The twelve class means are now OPEN RINGS,
%     radius 1.4pt, white fill, 0.5pt rulegrey edge -- this document's mark for
%     a schematic, unscaled object -- and the seven gate fractions stay filled
%     accent discs at radius 1.4pt, which is its mark for a point estimate.
%     One look now separates them and the disclaimer line is deleted.
%  4. THE FRACTION AXIS IS DEFINED AT THE AXIS. This is the supervisor's first
%     complaint, answered mechanically. It carried bare numerals 1.0 / 0.5 / 0
%     with the quantity named only in an italic subtitle about 30mm above and
%     20mm left of the topmost tick, so a reader at the axis could not tell
%     whether 1.0 meant "all the signal removed" or "accuracy 1.0". It now
%     carries, rotated 90 degrees reading upward and immediately left of the
%     tick numerals, exactly fig:mechanism's form: line one \footnotesize roman
%     ink naming the measured quantity, line two \scriptsize quietink giving its
%     construction. The layer axis keeps the label "layer" it already had --
%     one of only two axes in 4.7-4.10 that already met 4.11's standard -- and
%     is now left-aligned to the axis's own left end rather than centred on the
%     widest tick gap.
%     MEASURED, and this is what fixed the wording. A rotated line cannot be
%     longer than the lane left of the tick column, and that lane is pinned at
%     both ends by things that may not move: below by the foot note's ascender
%     line at y = 0.52, above by the descender line of the prose at x = 0, which
%     after the 1.5mm lift of item 15 sits at y = 6.44. The axis's own midpoint
%     is 3.375, so the usable half-length is 2.855cm and a line may run 5.71cm.
%     Measured off this figure's own render, \footnotesize sets at 0.166
%     cm/character in this class and \scriptsize at 0.119, so the lane holds 34
%     \footnotesize or 47 \scriptsize characters.
%     "fraction of above-chance drug signal removed" is 43 characters, needs
%     7.14cm, and fits no layout under the 105mm height cap. What is set instead
%     is "fraction of drug signal removed" (31 characters, 5.15cm) over
%     "(before minus after)/(before minus chance)" (42 characters, 5.00cm). The
%     second line carries "above chance" in the only form that is auditable
%     against the JSON anyway, the construction itself: a reader standing at the
%     axis can now see that 1.0 means all of it, and can recompute the quantity
%     from the three accuracies in the foot note. The minus signs are spelled as
%     words and not set as $-$: in math mode CMSY10 came out at 8.3pt, a THIRD
%     type size above 7pt, which the style contract's first clause forbids.
%     SUPERSEDED IN v11, AND THIS PARAGRAPH IS THE ERROR v11 EXISTS TO CORRECT.
%     Three things in it are wrong. (i) THE LANE ARITHMETIC DISAGREES WITH ITS
%     OWN LEDGER AND WITH THE CODE: this block says lane 5.71cm at 0.166
%     cm/character with the axis midpoint at 3.375, while the inline comment at
%     the axis said 5.44cm at 0.173 with midpoint 3.24, and figs/purify.json
%     said 5.71/0.166/3.48. Remeasured on the v11 render from the PDF's own
%     text spans, \footnotesize sets at 0.161 cm/character in this class, and
%     the one figure that matters is now recorded once, beside the axis.
%     (ii) THE CONCLUSION DOES NOT FOLLOW. That a 43-character STRING does not
%     fit rules out the string, not the QUANTITY; dropping "above-chance" left
%     the \footnotesize line naming a quantity the marks are not, and that is a
%     correctness defect, not a compromise. The name is now set over two rotated
%     \footnotesize lines and the \scriptsize construction line, which S5 makes
%     optional, is deleted -- it had become the tallest object in the panel,
%     rising 6.0pt above the panel head's own top, so the smallest type in the
%     block outranked the head. The construction moved to the foot note.
%     (iii) THE $-$ NOTE IS RIGHT AND WAS APPLIED ONE CHARACTER SHORT: the v10
%     foot note still set "chance $0.083 = 1/12$", whose "=" came out of CMR10
%     at 8.3pt, a third size above 7pt, exactly the fault the note describes.
%     v11 sets "chance $0.083$, exactly $1/12$" and the render carries two
%     sizes above 7pt and no CMR10 at all.
%  5. THE DISPLAY FRACTION IS OUT OF THE PLOT. A \displaystyle fraction stacked
%     in a plot's label gutter was the most slide-like object in the five-figure
%     set. Its content is not lost: it IS the axis's second line now, which is
%     where a definition of the measured quantity belongs. The three accuracies
%     it explained (before 0.41--0.82, after 0.121 at every layer, chance 0.083)
%     move to a two-line \scriptsize foot note, and the six-line right gutter --
%     which mixed two direct labels with four notes at three different left
%     edges, 320.7, 402.9 and 428.4pt on the printed page -- is dissolved. What
%     stays beside the plot is exactly the two DIRECT labels: "all seven
%     layers: 88--95%" at the height of the dot cloud it names, and "gate, 0.8"
%     at the right end of the rule it names.
%  6. THE 0.8 GATE IS BLACK, SOLID, 0.9pt. It was accent 0.6pt dashed. This is
%     a correctness repair and not a preference: the same pre-registered class
%     of object -- a threshold a measurement is tested against -- was drawn
%     three different ways inside nine printed pages (accent dashed here,
%     accent dashed in fig:materiality, solid black in figs/family.pdf), while
%     fig:mechanism draws the chance floor solid black 0.9pt and style_v2.py's
%     chance_line() is hard-coded color='black', lw=0.9, solid, with the
%     docstring "Thin, solid, black. Never dashed, never omitted". Accent in
%     this document means the object under test; a threshold is not the object
%     under test, so a threshold may not be accent. Its label goes with it:
%     \scriptsize roman black, reading "gate, 0.8".
%  7. THE FAIL REGION IS LIGHTENED TO rulegrey!12 AND EMPTIED OF TEXT. It was
%     rulegrey!18 carrying two italic lines, it filled about 80% of the plotting
%     area, and a cold reader's eye went to it first and read it as the data
%     while the seven dots sat in a 12% band at the top. rulegrey!12 over white
%     renders at grey 243, above the 230 threshold the ink budget counts, so
%     the region now costs nothing against the 6.0% ceiling and still shows. Its
%     emptiness is the message; nothing is written inside it.
%  8. THE CONFOUND AND THE RIGHT ANGLE ARE LABELLED AT THE MARK. A cold reader
%     attached the label "the confound" to the diagonal band rather than to the
%     intersection patch, and did not see the right-angle glyph at all until
%     zooming: it was 2mm. The glyph is now 4mm on each arm and it carries the
%     direct label it never had.
%     THE LABEL IS ONE WORD, AND THAT WAS FORCED BY A RENDER. "orthogonal by
%     construction" is 3.6cm at \scriptsize, and panel (b) has no 3.6cm lane at
%     the glyph's own height that does not cross V_pure, the context band or the
%     removal arrow. Set anchor=base east at x = 9.71 it ran back to x = 6.20,
%     i.e. it sat in the gutter BETWEEN the two panels and a cold read could not
%     tell which panel owned it. Set base west on panel (b)'s x origin and
%     broken over two lines it ran to x = 10.20 and butted the V_pure label at
%     10.22, so the two read as one string, "orthogonal byV_pure" -- measured at
%     500dpi on the pass-4 render. What is drawn is the one word that has to be
%     AT the mark, set adjacent to the glyph's vertical arm with no leader.
%     "by construction" is the claim and not the name, so it goes to the
%     caption, and figs/purify.json records the move.
%     The confound's label acquires a 0.30pt leader to the patch. The dose
%     reading joins the two subspace fractions in one measurement strip on panel
%     (b)'s own x origin, so the four measured numbers of stage two read as one
%     block and not as four asides scattered round the geometry.
%  9. THE TWO PROSE BLOCKS ARE CUT TO DEFINITIONAL CONTENT ONLY, AND THEY NOW
%     OUTRANK THE LABELS. Four \scriptsize grey lines per panel became two
%     \footnotesize roman ink lines per panel: "Twelve centred class means span
%     at most eleven / directions; ten are kept." and "Drugs are confounded with
%     plate and dose by design; / no causal claim is made on the raw slab." The
%     conservatism argument ("so the slab keeps essentially the whole span,
%     noisiest included, which for a removal claim is conservative") and the
%     dose-recoverability sentence are argument, not definition, and are staged
%     for the caption. figs/slab.tex sets its explanatory prose ABOVE its labels
%     in size and this file used to invert that -- 7.97pt grey prose under
%     9.96pt grey italic narration. Prose is now the largest type in the figure
%     alongside the panel heads and the axis label, which is 4.11's ranking.
% 10. THE TWO DIVIDERS ARE GONE, the full-measure horizontal one at y = 4.72 and
%     the vertical one at x = 8.40. The horizontal rule existed to carry the
%     italic sentence beneath it (item 2) and to say "above this nothing has a
%     scale", which the open ring now says at every mark; the vertical rule was
%     furniture the panel heads make redundant. fig:mechanism draws no dividers.
% 11. ONE GRID, ONE SET OF WEIGHTS, ONE SET OF SIZES.
%     x origins, one per panel, referenced by every panel-level line:
%     (a) x = 0.00, (b) x = 8.40, (c) x = 1.60, which is the fraction axis's own
%     spine. Panel (b) carries a second, deliberate text column at x = 11.95 for
%     the four labels belonging to marks on its right. The figure's own origin,
%     x = 0.00, carries the two foot lines. MEASURED on the finished render,
%     from the PDF's own text spans: panel (a)'s six panel-level lines all start
%     at 59.38pt, panel (b)'s five at 297.49pt, its right column's four at
%     398.12pt, panel (c)'s head and its x-axis label both at 104.74pt, and the
%     two foot lines at 59.38pt. The spread inside every block is 0.00pt.
%     The upper band's two plot areas used to start 19.5pt apart for no reason;
%     they now share their head, strip and prose baselines exactly.
%     FOUR STROKE WEIGHTS, each with one meaning, against eight before:
%       0.30pt furniture -- ticks, leaders, the twelve displacement stems, the
%                           context band's outline
%       0.40pt axis rule -- the two axes of panel (c), and nothing else
%       0.50pt schematic -- the open rings, the global-mean cross, the dotted
%                           projection line, the orthogonalisation arrow: every
%                           unscaled mark in (a) and (b) is drawn at one weight
%       0.90pt threshold -- the gate rule, and nothing else
%     No weight is used exactly once. The 0.6pt arrow, the 0.35pt outlines, the
%     0.45pt cross and the 0.65pt strokes are all folded into these four.
%     TWO TYPE SIZES: \footnotesize (10.0pt) for panel heads, the axis label's
%     first line and the definitional prose; \scriptsize (8.0pt) for ticks,
%     direct labels, the scope line, the axis label's second line and the foot
%     note. Nothing else, and no italic anywhere.
%     AMENDED IN v11, SIZES UNCHANGED, ASSIGNMENT CHANGED: the fraction axis's
%     block is now BOTH of its lines at \footnotesize, because the block is one
%     wrapped quantity name and not a name over a subordinate note; the
%     \scriptsize construction line it used to carry is the second foot line.
% 12. ONE ARROWHEAD IN THE WHOLE FIGURE. A single arrowhead in this document
%     means an operation acting in a direction, and the orthogonalisation push
%     is the only operation drawn here. The displacement mu_d - mu_bar in panel
%     (a) used to carry one too, which made a static object look like an
%     action; it is now a plain 0.30pt stem like its eleven siblings, named by
%     a leader instead. Nothing else in the figure has a head.
% 13. COLOUR, one code, the same one across the five figures of the arc.
%     accent = the drug-side object under test and its direct labels: V_drug,
%     V_pure, the confound, the orthogonalisation arrow, the seven gate dots,
%     "all seven layers: 88--95%".
%     quietink = the comparator and control side: the context subspace's two
%     labels, the schematic labels in (a), the axis label's construction line,
%     the foot note.
%     AMENDED IN v11: the construction line has left the axis for the foot, and
%     three things join the quietink list -- panel (c)'s which-slab scope line,
%     the name of the shaded region ("an untestable null", a counterfactual and
%     so a comparator), and panel (b)'s measurement strip, which in v10 was ink
%     while panel (a)'s scope line was quietink at the same size on the same
%     baseline in the adjacent panel, with nothing on the page decoding the
%     difference. ink is now heads, definitional prose, the axis labels, the
%     gate label and the one geometric fact true by construction, and nothing
%     else.
%     ink/black = thresholds, axis rules, axis labels, panel heads, definitional
%     prose, and the one geometric fact true by construction.
%     rulegrey = furniture only: ticks, leaders, stems, the ring edges, the
%     context band's outline, the fail region's fill.
%     The one visible consequence: the twelve class means are no longer ink
%     dots. They are schematic and they are drawn in the schematic register.
% 14. WHAT WAS PROTECTED AND IS UNTOUCHED. The three-stage spine and its order.
%     The seven gate dots, byte-exact, and the identity behind them. The full
%     0-to-1 gate range, decision 12, which a reviewer asked to crop to
%     0.75--1.00 and which was overruled and stays overruled. The layer axis
%     and its map's SHAPE, 1.45 + (L-2)/14*W. The orthogonalisation geometry:
%     raw band at 38 degrees, context band horizontal, V_pure vertical, the
%     dotted perpendicular, the named confound, the removal arrow, the right
%     angle. The refusal to draw a reference rule at 0.90 for the dose
%     cross-decode, for the reason in (e). The four measured numbers of stage
%     two and their layer-9 qualifier.
%
% 15. PANELS (a) AND (b) ARE LIFTED 1.5mm AS ONE BLOCK, inside a single
%     \begin{scope}[yshift=1.5mm]. Two things come of it. The ink budget,
%     measured at 200dpi inside the content bbox, falls from 5.98% to 5.90%
%     against a 6.00% ceiling, purely by growing the denominator -- drawn height
%     102.4 -> 103.8mm, still under the 105mm cap. And the gap between the last
%     prose line of the schematic band and panel (c)'s head goes from 2.8mm to
%     4.3mm, which is the separation the deleted full-measure divider used to
%     carry, now made of white instead of a rule. Panel (c) does not move: its
%     rotated axis label's lane is pinned between the foot note's ascender line
%     and the prose's descender line, so lifting the prose only lengthens it.
%
% WHAT SEVEN RENDER-AND-READ PASSES CHANGED IN v10, each with the measurement
% that forced it. Every one was found by rendering the figure and LOOKING at
% the PNG, not by reading the source.
%   i.   PASS 1. The rotated axis label was one \footnotesize line centred on
%        the axis and it overhung the axis at both ends, running up into panel
%        (a)'s prose and down into the tick lane. The plot was at once too tall
%        (3.95cm of axis, so the empty fail region dominated the page) and too
%        short for the label. Diagnosis: the LANE, not the axis, is the
%        constraint. Panel (c) was reproportioned so the axis midpoint sits
%        exactly at the lane's centre, which is the only arrangement in which
%        both rotated lines can be centred on the axis and still clear.
%   ii.  PASS 2, Overfull \hbox 6.00356pt, and it survived a complete rewrite of
%        the body, which is what identified it. The INK bbox measured 170.4mm,
%        inside the 174.0mm fullwidth measure, so the render looked clean; the
%        NODE box did not. "the global mean", set anchor=base east at x = 1.80,
%        is 2.00cm wide, so its left edge landed at tikz x = -0.32 and widened
%        the picture to 17.62cm = 501.3pt against \linewidth 495.08pt -- the
%        latter measured by compiling \the\linewidth inside a real fullwidth
%        environment, which also settles an old dispute in this header: 174.0mm
%        is right and the note calling preamble.tex's arithmetic stale is not.
%        The label moved to anchor=base west on the panel's own x origin.
%   iii. PASS 3. Ink measured 7.94% against a 6.00% ceiling. Region-by-region
%        measurement put 1.63 of those points in ONE object, the V_drug band at
%        fill opacity 0.15: accent!15 over white renders at grey 227, under the
%        230 threshold the budget counts, over 3.06cm^2. Both accent bands went
%        to 0.12, which renders at grey 233 and therefore costs nothing at all,
%        and the fail region went from rulegrey!18 (grey 237) to rulegrey!12
%        (grey 243). Ink 5.87%, and no measured mark was touched to get there.
%   iv.  PASS 4, at 500dpi. Three collisions invisible at 200dpi: "orthogonal
%        by" butting the V_pure label (item 8); the raw slab's leader stopping
%        0.19cm short of the band and reading as a stray hairline, which is the
%        exact fault the pass-9 header records against the orthogonalise leader;
%        and the fraction axis's tick numerals sitting 1.6pt from the rotated
%        label, so the two read as one column of type.
%   v.   PASS 5. The raw slab's label then sat 0.05cm under the measurement
%        strip's first line and the two read as one three-line paragraph. The
%        label moved above the strip and its leader now lands INSIDE the band
%        rather than beside its tip.
%   vi.  PASS 6. The global-mean cross was 2mm arms at 0.5pt with twelve stems
%        converging on it, and at 500dpi only its vertical arm was visible
%        through them; it is now 2.8mm. The two schematic outlines in panel (b)
%        moved 0.30 -> 0.50pt so that 0.30pt means furniture and nothing else.
%        V_pure acquired its gloss AT THE MARK, which is where this document
%        requires a symbol that is itself a drawn object to be glossed.
%   vii. PASS 7. Ink 5.98% against a 6.00% ceiling is not a pass, it is a
%        rounding error away from a fail. Item 15.
%
% HEIGHT AND MEASURE, v10, all MEASURED off the finished render at 200dpi and
% not computed from the source: bounding box \path (0,-0.10) rectangle
% (17.30,10.35); content bbox 172.1 x 103.8mm, inside the 170-174mm band a
% fullwidth figure must fill, under the 105mm cap this pass was given and under
% the corpus ceiling of 117mm (slab.tex). The data region, x = 1.60 to 14.10,
% is 12.50cm = 72.6% of the drawn width, over the 60% floor. Ink 5.90% of the
% content bbox at pixels below grey 230, against a 6.00% ceiling. Text density
% 4.26 characters per 100mm^2 against 5.00. Type: two sizes above 7pt, 9.96 and
% 7.97, against fig:mechanism's 10.0 and 9.0; one face, TeXGyrePagella; three
% italic characters in 864 and all three are the math-italic V of a symbol;
% zero sans, zero bold, zero all-caps runs of two or more words. Four stroke
% weights, 0.30 / 0.40 / 0.50 / 0.90pt, none used exactly once. Exactly one
% black solid 0.90pt rule and it is the gate -- byte-for-byte the register
% figs/family.pdf gives the other pre-registered threshold of this chapter, a
% 0.9pt rule at colour (0,0,0), checked by reading family.pdf's drawing list.
% SUPERSEDED IN v11 as to the MEASUREMENTS ONLY, all of which were remeasured
% and several of which were wrong or unreproducible: see "HEIGHT AND MEASURE,
% v11" at the foot of the v11 block below. The claims about type, faces,
% weights and the single black rule survive, with one correction (the CMR10
% "=" at 8.3pt, item (iii) of the superseding note under v10 item 4).
%
% ===========================================================================
% v11 -- 2026-08-19. ADVERSARIAL REPAIR PASS. Three refuters read the v10
% render and returned 33 findings. What follows is what was applied, what was
% overruled, and why. Nothing below changes a plotted or printed VALUE: the
% seven gate fractions are the same seven floats, unrounded, and 0.8,
% 0.41--0.82, 0.121, 0.083, 0.90--0.95, 0.620, 0.031, 480, twelve, forty, ten,
% eleven and seven are all unmoved. What moved is positions, wording that
% NAMES, and the audit record.
%
% APPLIED, in order of severity.
%
%  1. THE FRACTION AXIS NAMED THE WRONG QUANTITY. This is the one correctness
%     defect in the drawing and it was in the very line v10 was built to add.
%     The plotted series is layers.<L>.frac_signal_removed, which is the
%     fraction of ABOVE-CHANCE signal removed; v10's axis read "fraction of
%     drug signal removed". The unqualified quantity, (before - after)/before,
%     computes from the same file to 0.7041 at layer 2 and 0.8528 at layer 9 --
%     4.2 to 18.0 points BELOW the marks drawn -- and read at the drawn name the
%     top mark, 0.949, implies a residual accuracy of 0.041, under the chance
%     0.083 the figure prints two lines away. The pre-registration itself
%     (04b-representation.tex:218) says "the fraction of above-chance signal
%     removed must exceed $0.8$". The axis now reads, over two rotated
%     \footnotesize lines, "fraction of above-chance / drug signal removed".
%     v10's stated reason for the omission is superseded in place above.
%
%  2. THE RIGHT ANGLE AND THE REMOVAL ARROW WERE ONE PRIMITIVE. Measured on the
%     v10 render: the glyph's horizontal arm at y = 9.05 and the arrow's shaft
%     at y = 9.0349, 0.0151cm = 0.13mm apart, on opposite sides of an opaque
%     0.18cm bar, so a static assertion (perpendicularity) and a directed
%     operation (the projection) read as one horizontal line running from the
%     word "orthogonal", through V_pure, out to the dotted drop line. Arms
%     0.40 -> 0.28cm, arm to y = 8.93, 1.05mm of white under the arrow. The
%     quadrant is and always was UPPER-left; v10 item 8's "lower-left" is
%     superseded in place above.
%
%  3. PANEL (b)'s HEAD AND THE V_pure GLOSS WERE 0.36pt APART. 0.13mm, against
%     2.8mm under panel (a)'s head: two panels of equal rank with head
%     clearances differing by a factor of twenty-two. Fixed by a restack whose
%     every step is set by measured ink boxes -- see the note at the V_pure bar.
%     Head-to-gloss is now 6.4pt (2.3mm) and the gloss sits 7.6pt clear of the
%     right column's first baseline, where the first v11 build had put it level
%     with that baseline and the two read as one line.
%
%  4. WHICH SLAB THE GATE IS MEASURED ON IS BACK ON THE DRAWING. v10 cut it for
%     text density and left panel (c) with no slab named anywhere in its head,
%     axis, ticks, labels or foot, sitting directly under panel (b)'s closing
%     line "so causal claims use the purified slab, not the raw". The inference
%     that the gate was run on V_pure was then unavoidable, and it is not
%     supported: workspace_probe.json stores one ablated series and no purified
%     counterpart. The scope line sits under the panel head, in axis-label rank.
%     NOTE ON THE LEDGER'S OLD REASONING, corrected in figs/purify.json: the
%     identity in (d) above fixes only the BEFORE term as the raw probe
%     accuracy; it says nothing about which subspace was projected out. The
%     honest statement, and the one now drawn and recorded, is that the file
%     does not record it. (For scale: probe_arm_residual.json, the one instrument
%     that stores both, has layers.12 frac_signal_removed 0.3634 against
%     frac_signal_removed_pure 0.3241, i.e. the purified-slab gate is the LOWER
%     of the two there -- which is why this may not be assumed either way.)
%
%  5. THE FOOT NOTE HAD NO NOUN. v10 printed "Before 0.41--0.82, after 0.121 at
%     every layer, chance 0.083" and the words "accuracy", "probe" and "1/12"
%     appeared nowhere in the figure, so nothing on the page said what
%     instrument panel (c) reads. It now names the quantity by its registry
%     name, "probe accuracy", gives chance as "$0.083$, exactly $1/12$" (the
%     stored float IS 1/12 in IEEE double), and carries the construction the
%     deleted rotated line used to hold.
%
%  6. THE SHADED FAIL REGION IS SMALLER, LIGHTER AND NAMED. It was 21.6% of the
%     content box, the largest object on the page, unnamed, and v10 had nearly
%     doubled its absolute area as a side effect of stretching the axis to house
%     an over-long label. The full 0-to-1 range is untouched -- decision 12
%     stands and is the first thing this pass was told to protect -- and so is
%     the region's exact 80% share of the axis. What changed is the SPAN, which
%     is a free choice: 3.85 -> 3.00cm, so the region falls to 16.8% of the
%     content box. It is drawn at rulegrey!8 (grey 246) and not rulegrey!12
%     (grey 243), which was within one level of the panelbg tint that means
%     "a schematic subspace" 3cm above it -- one grey, two meanings. And it
%     carries a NAME, "an untestable null", set OUTSIDE its right edge in
%     quietink: nothing is written inside it (the spec forbids it and the
%     emptiness is the message), but the largest mark on the page may not be the
%     one mark a reader cannot name, and without the name the third clause of
%     this figure's acceptance test is unanswerable from the drawing.
%
%  7. "layer" READ AS "layer 2". v10 set it \footnotesize at x = 1.60 with the
%     layer-2 numeral 0.63mm directly above it, one size smaller -- the exact
%     fault the pass-9 note above records against the old placement under the
%     layer-9 numeral, reintroduced at layer 2. The repair is VERTICAL, not
%     horizontal, so S5's left-alignment survives: the fraction origin rose
%     1.45 -> 1.75, which buys 2.8mm of white and a 0.53cm baseline separation
%     against the 0.40cm leading this figure uses inside one sentence.
%
%  8. THE GLOBAL-MEAN CROSS WAS STILL INVISIBLE AND ITS LEADER READ AS A
%     THIRTEENTH STEM. figs/purify.json certified a pass-6 legibility repair
%     (arms 2.0 -> 2.8mm) that the render never showed. Length was never the
%     fault: three of the twelve stems leave the vertex within 10 degrees of
%     horizontal, so the horizontal arm is buried at every radius. A 0.20cm disc
%     of the band's own tint now separates them. See the note at the mark.
%
%  9. FOUR LEADERS AND ONE PATCH. "the confound" pointed 0.03cm BELOW its patch,
%     into the context band -- the misreading the pass-9 rebuild of that mark
%     exists to kill. "orthogonalise" died in white 0.105cm short of the arrow's
%     tail while carrying only a 0.05cm standoff at the label end, the inversion
%     of the corpus rule this header states twice. "the global mean" grew out of
%     the final "n" of its own label at 0.03cm. "all seven layers: 88--95%" sat
%     3.6mm from the layer-16 dot with nothing between, over the corpus's 3mm
%     leader threshold. All four are fixed at the mark end. And "$V_{drug}$,
%     raw" rose to y = 8.352 and sat on the context band's 8.25 top edge; it
%     drops to baseline 7.98.
%
% 10. THE TINT-MULTIPLY PARALLELOGRAM WAS STILL BEING DRAWN. Pass 9 deleted the
%     CLAIM that the darkened crossing is the confound but not the overprint
%     that produced one: the raw band was drawn after the opaque context
%     rectangle, so an unlabelled darker parallelogram appeared inside the band
%     beside the patch that actually carries the name. The two \fill blocks are
%     simply reordered. No coordinate and no colour changed.
%
% 11. COLOUR. Panel (b)'s measured-number strip was ink and panel (a)'s scope
%     line quietink, at the same size on the same baseline in adjacent panels,
%     with nothing on the page decoding the difference. Both are scope; both are
%     quietink. ink is now heads, definitional prose, the axis labels, the gate
%     label and the one geometric fact true by construction, and nothing else.
%
% 12. ONE OBJECT, ONE NAME. "subspace fraction" in panel (b) becomes
%     fig:mechanism's own name for the quantity, "descriptive subspace
%     fraction".
%
% 13. V_drug's BAND IN PANEL (a) IS SYMMETRIC ABOUT ITS CLOUD. It overhung the
%     data 0.13cm left against 0.30cm right, so the leftmost ring read as
%     sitting on the band's edge -- the bounded-region reading the deleted
%     ellipse was removed to prevent. Left edge 0.55 -> 0.38; both 0.30cm.
%
% OVERRULED, with the reason, because a refuter who is wrong about this
% document's conventions must be answered and not quietly followed.
%
%  A. "BRING PANEL (c)'s TWO DIRECT LABELS INSIDE THE AXIS." Overruled. It would
%     take the drawn width to 169.7mm, under the 170mm floor the SAME clause
%     sets for a fullwidth figure, and it would put "all seven layers: 88--95%"
%     4mm from the marks it names with no room for the leader the corpus
%     requires. The two clauses cannot both hold in this geometry. What is drawn
%     instead: the labels stay at their marks, "all seven layers" acquires its
%     0.30pt accent leader, and the figure's right edge is set by a direct label
%     0.48cm past the widest panel element (panel (b)'s prose at 16.78cm). That
%     is the trade recorded, not an oversight.
%
%  B. "DELETE THE TWELVE DISPLACEMENT STEMS." Overruled. They are the picture of
%     the word "centred" in the panel's own prose -- each mean is drawn as a
%     displacement from the global mean -- and they are what makes the count
%     twelve a picture rather than an assertion. Their whole ink is 0.77% of the
%     figure's, so deleting them buys almost nothing, and the three faults the
%     refuter traced to them (an invisible cross, a leader that reads as a
%     thirteenth spoke) are fixed by item 8 without losing the picture.
%
%  C. "MAKE V_pure AND THE RAW BAND THE SAME WIDTH, so no width carries a
%     value." Overruled. The drawn width ratio is not read as the measured one
%     -- nothing in panels (a) and (b) has a scale, the open-ring glyph says so
%     at every mark and the caption says the rest -- while "visibly thinner" is
%     the one cue that ties the drawing to the 0.620 -> 0.031 printed beneath
%     it. Equalising the widths would re-open the pass-9 misreading the header
%     records at length: two bands of equal width differing only in angle say
%     "same slab, rotated, nothing lost".
%
%  D. "REDRAW V_drug IN PANEL (a) AS A THIN STRIP AND LET THE OUTLYING RINGS
%     FALL OUTSIDE IT, so the two panels' glyphs rhyme." Overruled on the
%     file's own recorded ground, which is correct: the twelve centred class
%     means span V_drug BY CONSTRUCTION, so a class mean drawn outside it would
%     be false. Decision 10(a) already settled this and nothing has changed.
%
%  E. "WIDEN PANEL (a) AND MOVE PANEL (b)'s ORIGIN LEFT to close the 3.5cm void
%     between the two geometries." Overruled as disproportionate. The void is
%     confined to the geometry band; at the prose rows the two panels fill the
%     measure with a 0.93cm gutter, and moving panel (b)'s origin would move
%     fifteen coordinates and break the one-x-origin-per-panel grid the v10
%     render was measured against for no measurable gain.
%
%  F. "'the confound' IS 17pt OFF THE RIGHT COLUMN's LEFT EDGE." Overruled: it
%     is a DIRECT LABEL AT ITS MARK, which the corpus requires to sit at the
%     mark, and direct labels are exempt from the panel-origin rule -- as are
%     "$V_{drug}$, raw", "orthogonal" and the V_pure gloss in the same panel.
%
%  G. "PRINT THE SEVEN-LAYER RANGE OF THE RAW SUBSPACE FRACTION ON THE DRAWING."
%     The finding is RIGHT and is acted on, but not there: layer 9 is indeed the
%     MAXIMUM of variance_share.drug (the seven-layer range is 0.2767 at layer
%     16 to 0.6201 here), and the drawing does not say so. There is no lane for
%     it -- the strip line already runs to 16.38cm -- so the disclosure goes to
%     figs/purify.json and to the staged caption, which is where this document
%     keeps what will not fit. The layer-9 qualifier stays on the drawing, so
%     the pair cannot be read as a seven-layer summary.
%
%  H. NOT DONE, AND OWED. figs/_PENDING_CAPTIONS.tex still carries no
%     \label{fig:purify} block, so the eight sentences v10 cut and the whole
%     replacement caption live only in the comment at the foot of this file.
%     Two refuters are right that this is a contract breach. It is not repaired
%     here because this pass's writable set is figs/purify.tex and
%     figs/purify.json only; the block is staged verbatim below for whoever
%     holds the pen on that file, and figs/purify.json records the debt under
%     pending_captions_block_owed.
%
% WHAT SIX RENDER-AND-READ PASSES CHANGED IN v11, each with the measurement.
%   i.   PASS 1. The two rotated \footnotesize lines were set in the wrong
%        order. Rotated 90 degrees the line-advance direction is -x, so the
%        first line is the LEFTMOST; set the other way round the block printed
%        "drug signal removed / fraction of above-chance".
%   ii.  PASS 2. Dropping the V_pure gloss for head clearance landed it on the
%        right column's own first baseline with 0.15cm of horizontal air, and
%        the two read as one line. Hence the restack of item 3.
%   iii. PASS 3. Text density 5.01 characters per 100mm^2 against a 5.00
%        ceiling -- a rounding error away from a fail, which is not a pass.
%        Two wordings were tightened (the scope line, the interval line) with
%        no content lost; 4.96 now.
%   iv.  PASS 3. The render carried THREE type sizes above 7pt: 9.96, 7.97 and
%        an 8.3pt CMR10 "=", two characters, from "$0.083 = 1/12$". Reset as
%        "$0.083$, exactly $1/12$"; two sizes now and no CMR10 in the file.
%   v.   PASS 4. Panel (c)'s head and its new scope line had 3px of white
%        between their ink at 200dpi -- 0.38mm. The fraction span came 3.15 ->
%        3.00cm and the head rose 0.04cm; 15px, 1.9mm, now.
%   vi.  PASS 5. Read at 6x, the whole of panel (b) and the whole of panel (a).
%        Every leader in item 9 was found this way and not in the source.
%
% HEIGHT AND MEASURE, v11, MEASURED off the finished render and not computed.
%   bounding box \path (0,-0.10) rectangle (17.30,10.35), unchanged.
%   content bbox 172.09 x 103.76mm, inside the 170--174mm band and under the
%     105mm cap and the corpus ceiling of 117mm (slab.tex).
%   data region x = 1.60 to 14.10 = 12.50cm = 72.6% of the drawn width, over
%     the 60% floor.
%   text density 4.96 characters per 100mm^2 (spaces excluded, which is the
%     count that reproduces the contract's own reading of fig:mechanism, 2.49
%     against its stated 2.5) against a 5.00 ceiling.
%   ink. THE CONTRACT'S 6.0% CEILING IS PIPELINE-DEPENDENT AND THIS FILE NOW
%     SAYS SO RATHER THAN QUOTING A NUMBER THAT CANNOT BE REPRODUCED. Rasterised
%     with PyMuPDF and thresholded at grey 230 inside the content bbox, this
%     figure reads 6.54% at 200dpi, 5.56% at 400dpi and 5.21% at 600dpi; the
%     SAME pipeline on fig:mechanism as it prints on page 51 of main.pdf reads
%     5.99 / 5.26 / 5.04. The contract's own measurement of fig:mechanism is
%     5.3, which its pipeline reproduces at 400dpi and not at 200. At 400dpi,
%     the resolution at which the reference figure reproduces, this figure is
%     5.56% against the 6.00% ceiling. At any of the three resolutions it sits
%     between 1.03x and 1.09x fig:mechanism, where the ceiling is 1.13x it.
%     v10's "5.90%" was a 200dpi number quoted against a 400dpi ceiling.
%   type: two sizes above 7pt, 9.96 and 7.97, plus 5.98 for subscripts, which
%     the contract exempts; largest size is only ever a panel head, an axis
%     label or definitional prose. One face, TeXGyrePagella, with 58 characters
%     of URWPalladioL-Roma inside math \textrm and 5 of URWPalladioL-Ital, all
%     five the math-italic V of a symbol -- 0.5% italic against a 5% ceiling.
%     Zero sans, zero bold, zero all-caps runs of two or more words.
%   strokes, read out of the PDF's own drawing list: four weights,
%     0.30 / 0.40 / 0.50 / 0.90pt. 0.30 is furniture (27 rulegrey) plus the one
%     accent leader the corpus requires; 0.40 is the two axis rules and nothing
%     else; 0.50 is every unscaled mark; 0.90 appears exactly once and it is the
%     gate, in black, which is the one place this document allows a singleton.
%   the harness: ok, pages=1, boxes=0.
% ===========================================================================
%
% ===========================================================================
% v-BLOCK, 2026-08-19: WHOLE-SLATE PASS (fig:comparators, fig:licensing,
% fig:purify, fig:hook, fig:materiality reconciled against one another and
% against fig:mechanism, which is the document's reference standard and was
% NOT rebuilt). Nothing above this line is deleted; where an earlier claim is
% now false it is superseded EXPLICITLY below.
%
% WHY. The five figures were rebuilt independently and drifted. This pass owns
% only the differences BETWEEN them. No value changed in any of the five; the
% sibling .json ledgers carry the value-to-position maps, old and new.
%
% THE SLATE RULES THAT NOW HOLD ACROSS ALL FIVE (each stated once here, in
% every file, so any future single-figure edit can see what it would break):
%
%  R1 TICK NUMERALS are quietink, \scriptsize, inner sep 0, anchored north on
%     the tick and set 0.16cm below the axis rule, over a 0.10cm rulegrey
%     0.30pt tick. Before: comparators/licensing/purify quietink, hook and
%     materiality ink; ticks 0.10 / 0.12 / 0.14cm; three different gaps.
%  R2 A THRESHOLD'S OWN LABEL IS BLACK, matching the black solid 0.90pt rule
%     it names (S6). Before: comparators black, purify's "gate, 0.8" and
%     materiality's two margin labels ink.
%  R3 ONE NAME PER OBJECT (S19). The pre-registered 10% rule is "10%
%     materiality margin" in fig:comparators row (c), fig:materiality (a) and
%     (b), and in figs/family.py, which is where the name comes from.
%     fig:comparators' "pre-registered margin" is superseded; "pre-registered"
%     is now the caption's word, not the drawing's. Zero, where it is the null
%     a bar is read against, is "zero: no identity effect" in BOTH
%     fig:comparators row (c) and fig:materiality (b) -- one quantity, one
%     wording. V_pure is spelt with \textrm everywhere, as figs/slab.tex does.
%  R4 COLOUR KEY, one key for five figures (S8). accent = the drug-side object
%     under test. quietink = every comparator, null, control, replication or
%     discounted arm, AND its label, AND its leader. ink/black = thresholds,
%     axis rules, axis labels, panel heads, definitional prose. rulegrey =
%     furniture only. Within quietink the GLYPH separates two kinds: an OPEN
%     disc is a construction built to carry no signal (a null, a control, a
%     raw/before series); a FILLED disc is a real measurement that is not the
%     headline. fig:hook (d) drew both its null and its cell-line control in
%     ink and is corrected.
%  R5 PROJECTION AND ZERO RULES are quietink 0.50pt densely dotted (S7).
%     fig:hook's two panel-(c) projection lines were rulegrey 0.30pt densely
%     dashed and are corrected.
%  R6 DEFINITIONAL PROSE is \footnotesize ink and OUTRANKS every label (S4);
%     at most one block per figure, on the panel's own x origin. fig:hook's
%     one prose line was \scriptsize and is raised. fig:purify's panel-(b)
%     block carried an argument clause ("so causal claims use the purified
%     slab, not the raw") that its caption already carries, and is now
%     definitional only.
%  R7 ONE INTERVAL DECLARATION per figure, \scriptsize quietink, at the foot,
%     on the figure's x origin, opening with the word "Intervals:" and naming
%     EVERY panel including the ones with no measured mark (S17).
%     fig:comparators already did this and is the model; the other four now
%     match its form.
%  R8 THE FIVE-FIELD PROMPT STRIP is ONE object drawn in two figures, and the
%     two drawings are now congruent: cell boundaries 0 / 1.12 / 1.90 / 2.64 /
%     4.23 / 6.90 cm and cell height 0.42cm in both fig:licensing (a) and
%     fig:hook (a) and (b), labels centred in their cells, the drug cell
%     accent-bordered at 0.40pt on accent!12 and its label NOT bold. Before:
%     licensing 6.90cm wide with 0.62cm cells and a bold label, hook 7.44cm
%     wide with 0.42cm cells and a plain label -- the same five fields at two
%     sizes, four pages apart, with each caption pointing at the other.
%  R9 MEASURE. Every fullwidth figure draws 172.0-174.0mm; the one textwidth
%     figure draws 140.3mm. Ink at 200dpi inside the content bbox, counted as
%     greyscale luminance below 230 (the contract's own method), is now
%     5.63-5.99% on all five, under the 6.0% body ceiling, where before this
%     pass it ran 5.36-6.54%.
%
% WHAT THIS PASS DELIBERATELY DID NOT UNIFY, so that a later reader does not
% think it was missed:
%  * PANEL-HEAD CASE. fig:comparators sets "(a) Is it there?" and the other
%    four set "(a) where the drug name enters". comparators' heads are
%    interrogative SENTENCES and sentence case is correct for a sentence; the
%    v4 header argued this and the argument stands. fig:mechanism's lowercase
%    heads are noun phrases, not questions. One capital letter, four times.
%  * PANEL-HEAD POSITION. fig:comparators puts its heads in a right-aligned
%    left stub, the other four above the panel at its x origin. Moving them
%    above the rules costs about 20mm of height against a 118mm cap that the
%    figure already meets at 117.60mm. Not available.
%  * TWO-COLUMN GUTTER. purify and hook split at x = 8.40cm, licensing at
%    7.75cm. licensing cannot move right: "(d) what a negative would buy" is
%    4.85cm wide and its origin is already at 12.40 of a 17.30cm measure, with
%    1.5pt of slack. Moving purify and hook LEFT to 7.75 would be arbitrary.
% ===========================================================================
%
% CHANGED IN THIS FILE. (i) The line "Gate measured on the subspace projected
% out; no separate gate is stored for V_pure" is DELETED from the drawing. It
% was the slate's one scope line not sitting under an axis label (S5), and its
% content is already in the caption verbatim as "Two limits. The gate is
% measured on the subspace that is projected out, and this file stores no
% separate gate for V_pure" -- an S18 duplication. Panel (c)'s head drops from
% y 5.50 to 5.20 to take up the space. (ii) Panel (b)'s prose becomes
% definitional: "so causal claims use the purified slab, not the raw" is an
% argument and is the caption's ("Panel (b) is why the raw subspace may not be
% used"); what is drawn now DEFINES the object the panel draws -- "V_pure is
% the part of the raw slab that is orthogonal to the context subspace."
% (iii) "gate, 0.8" is set black, matching the black solid 0.90pt rule it
% names (R2); it was ink, and this was the only threshold on the slate whose
% label and rule were in different colours. (iv) Tick geometry: the y ticks go
% 0.12 -> 0.10cm and the x ticks 0.14 -> 0.10cm (they differed from each other
% INSIDE this panel), numerals inner sep 0 at a 0.06cm gap (R1). (v) The foot
% block takes the slate's declaration form (R7).
% SUPERSEDED: pass 10's claim that the figure clears the ink ceiling at 5.90%
% was measured on the red channel of an RGB rasterisation, which counts the
% accent!12 bands as ink; measured as greyscale luminance, which is what S15
% specifies, pass 10 was at 6.541%. The deletions above bring it to 5.992%.
% RE-MEASURED: 172.09 x 103.76mm (cap 105); ink 5.992%; 910 characters, 5.10
% per 100mm^2; two type sizes above 7pt, 9.96 and 7.97, plus 5.98/7.57pt
% subscripts inside V_drug and V_pure.
% One page, zero Overfull, zero Underfull.
% ===========================================================================
\begin{tikzpicture}[x=1cm, y=1cm, font=\small,
                    every node/.style={inner sep=1.6pt}]
  \tikzset{
    % --- TWO SIZES, ONE FACE, NO ITALIC. -----------------------------------
    % \footnotesize (10.0pt): panel heads, BOTH lines of the fraction axis's
    % label (v11 -- the block is one wrapped quantity name, not a name over a
    % subordinate note), the x-axis label, and the definitional prose.
    % \scriptsize (8.0pt): everything else -- ticks, direct labels, the two
    % scope lines, the region's name and the two foot lines. Nothing in this
    % picture is set in sans, in small caps, in bold or in italic, and nothing
    % is set in math mode that brings a third size with it: "$0.083 = 1/12$"
    % put an 8.3pt CMR10 "=" in the v10 render and is now "$0.083$, exactly
    % $1/12$".
    head/.style={text=ink,      font=\footnotesize, anchor=base west},
    def/.style ={text=ink,      font=\footnotesize, anchor=base west},
    lab/.style ={text=ink,      font=\scriptsize,   anchor=base west},
    qlab/.style={text=quietink, font=\scriptsize,   anchor=base west},
    alab/.style={text=accent,   font=\scriptsize,   anchor=base west},
    % --- FOUR STROKE WEIGHTS, each with exactly one meaning. ---------------
    furn/.style ={draw=rulegrey, line width=0.3pt},   % ticks, leaders, stems
    axisr/.style={draw=ink,      line width=0.4pt},   % an axis rule
    schem/.style={draw=rulegrey, line width=0.5pt},   % an unscaled mark
    gate/.style ={draw=black,    line width=0.9pt}}   % a threshold

  % ---- value-to-position maps for panel (c). Both bodies are braced groups.
  % THESE TWO LINES ARE THE AUDIT RECORD OF THE MAPS AND THEY ARE THE CODE, so
  % read them and not the prose in the v10 header, which quoted four wrong
  % coefficients and is superseded item by item in the v11 block below.
  %   v9   \lx = 1.05 + (L-2)/14*11.00   \fy = 0.75 + f*2.10
  %   v10  \lx = 1.85 + (L-2)/14*12.00   \fy = 1.45 + f*3.85
  %   v11  \lx = 1.85 + (L-2)/14*12.00   \fy = 1.75 + f*3.00   (x map unchanged)
  % The fraction origin rose 1.45 -> 1.75 to put 2.8mm of white between the
  % layer numerals and the word "layer" (v11 item 3), and the span fell
  % 3.85 -> 3.00 to take 0.68cm out of the empty fail region and give it to the
  % gap between the schematic band and panel (c) (v11 items 5 and 6). Layer 2
  % is inset 0.25 from the spine so its dot does not straddle the axis rule,
  % which is fig:mechanism's own treatment of the same first layer.
  % No map change is a value change: the ledger records all three maps.
  \def\lx#1{{1.85+((#1)-2)/14*12.00}}  % layer  2 -> 1.85,  16 -> 13.85
  \def\fy#1{{1.75+(#1)*3.00}}          % 0 -> 1.75, 0.8 -> 4.15,  1.0 -> 4.75

  % one gate mark: a filled accent disc at 1.4pt, this document's point
  % estimate. The 2.1pt white halo of earlier passes is dropped -- nothing
  % underlies these dots now, and a second radius would blunt the ring/disc
  % distinction that replaced the schematic disclaimer line.
  \newcommand{\gdot}[2]{\fill[accent] (\lx{#1},\fy{#2}) circle (1.4pt);}

  % one schematic mark: an open ring, white fill, 0.5pt rulegrey edge.
  \newcommand{\sring}[2]{\filldraw[fill=white, schem] (#1,#2) circle (1.4pt);}

% ###########################################################################
% # PANEL (a) -- localise.  x origin 0.00.  Nothing here has a scale.
% # Panels (a) and (b) are lifted 1.5mm as one block -- see the v10 header.
% ###########################################################################
\begin{scope}[yshift=1.5mm]
  \node[head] at (0.00,9.90) {(a) localise the slab};

  % V_drug: an accent band through the global mean, running past the data on
  % both sides because a subspace has no boundary. Opacity 0.15, down from
  % 0.20: at 0.20 over 3.40 x 0.90cm this was still the heaviest area in the
  % figure and it is not the figure's subject.
  % v11: the left edge comes 0.55 -> 0.38. The twelve offsets run -1.57 to
  % +1.40, so the cloud sits at 0.68 to 3.65 and the band overhung it by 0.13cm
  % on the left against 0.30cm on the right; at 0.13cm the leftmost ring read as
  % sitting ON the band's edge, which reintroduces the bounded-region reading
  % the deleted ellipse was removed to prevent. Both overhangs are 0.30cm now.
  \fill[accent, fill opacity=0.12] (0.38,8.05) rectangle (3.95,8.95);
  \node[alab, anchor=west] at (4.08,8.50) {$V_{\textrm{drug}}$};

  % the twelve per-drug mean activation vectors. Positions are literal, so no
  % build can randomise them and no reader can be told a different story by a
  % different compile. Twelve 0.30pt stems make the count a picture; the rings
  % are OPEN, which is how this document draws a schematic, unscaled mark.
  \foreach \dx/\dy in {-1.57/-0.085, -1.25/0.238, -1.02/-0.289, -0.62/0.102,
                       -0.36/-0.374, -0.07/0.306, 0.21/-0.153, 0.55/0.357,
                       0.74/-0.306, 1.12/0.119, 1.40/-0.238, 1.21/0.298}{
    \draw[furn] (2.25,8.50) -- ({2.25+\dx},{8.50+\dy});}
  \foreach \dx/\dy in {-1.57/-0.085, -1.25/0.238, -1.02/-0.289, -0.62/0.102,
                       -0.36/-0.374, -0.07/0.306, 0.21/-0.153, 0.55/0.357,
                       0.74/-0.306, 1.12/0.119, 1.40/-0.238, 1.21/0.298}{
    \sring{2.25+\dx}{8.50+\dy}}

  \draw[furn] (1.20,9.06) -- (1.05,8.79);
  \node[qlab] at (0.00,9.32) {the twelve per-drug means};

  % THE GLOBAL MEAN: a cross and not a thirteenth mark, and leadered out of the
  % band because in place the label box collided with two of the twelve rings.
  % v11 REPAIRS A LEGIBILITY CLAIM THE LEDGER MADE AND THE RENDER NEVER
  % SUPPORTED. Pass 6 enlarged the arms 2.0 -> 2.8mm "because at 2mm only its
  % vertical arm was visible through the stems"; at 2.8mm only its vertical arm
  % was still visible, because the fault is not length. Several of the twelve
  % stems leave the vertex within 0.02cm of horizontal -- (-1.57,-0.085) at 3.1
  % degrees, (1.12,0.119) at 6.1, (1.40,-0.238) at -9.7 -- so the horizontal arm
  % is buried in a near-collinear bundle at EVERY radius and no arm length
  % separates it. What separates it is white: a 0.20cm disc of the band's own
  % tint (white, then accent at 0.12 -- byte-identical to the band, so the disc
  % has no visible edge) is laid over the stems, and the cross is drawn in it.
  % The same disc fixes a second finding: the label's leader used to be the
  % thirteenth grey line converging on a vertex the prose says twelve lines
  % reach. It is now the ONLY line that enters the clear disc, so it reads as a
  % leader and not as a spoke, and it is drawn after the disc for that reason.
  % Its label-end standoff also goes 0.03 -> 0.12cm: at 0.03 it grew out of the
  % final "n" of "mean" and read as a stray hairline off the type.
  \fill[white] (2.25,8.50) circle (0.20);
  \fill[accent, fill opacity=0.12] (2.25,8.50) circle (0.20);
  % Anchored base WEST on the panel's own x origin. Anchored base east at
  % x = 1.80 this label ran 0.32cm LEFT of x = 0, which widened the
  % picture's box past the fullwidth measure and put an Overfull hbox of
  % 6.00pt in the log: the INK bbox was inside the measure and the NODE
  % box was not, which is why the render looked clean and the log did not.
  \draw[furn] (2.14,7.88) -- (2.31,8.32);
  \draw[schem] (2.08,8.50) -- (2.42,8.50);
  \draw[schem] (2.25,8.33) -- (2.25,8.67);
  \node[qlab] at (0.00,7.80) {the global mean};

  % the instrument's sample, as a scope line and not as a sentence. Its
  % baseline is shared with panel (b)'s second measurement line, and both drop
  % 0.10cm in v11 behind the panel-(b) restack, so the two panels keep one grid.
  \node[qlab] at (0.00,7.20)
    {480 prompts: twelve drugs, forty real cells, seven layers};

  % definitional prose: \footnotesize roman ink, two lines, on the panel's own
  % x origin. It is the largest type in the panel, with the head and the axis
  % label, which is fig:mechanism's ranking.
  \node[def] at (0.00,6.62) {Twelve centred class means span at most eleven};
  \node[def] at (0.00,6.22) {directions; ten are kept.};

% ###########################################################################
% # PANEL (b) -- purify.  x origin 8.40.  Nothing here has a scale either.
% ###########################################################################
  \node[head] at (8.40,9.90) {(b) purify it against context};

  % ORDER MATTERS, AND v11 REVERSES IT. The raw band is drawn FIRST and the
  % opaque context rectangle over it. In v10 the pale diagonal was drawn last,
  % so where it crossed the opaque panelbg rectangle the two tints multiplied
  % and an UNLABELLED darker parallelogram appeared inside the context band --
  % exactly the region a cold reader takes for "where the drug slab and the
  % context subspace overlap", i.e. for the confound, sitting beside the patch
  % that actually carries that name. Pass 9 deleted the tint-multiply lozenge
  % as a claim; it did not delete the overprint that drew one. Now the only
  % tinted region inside the context subspace is the one that is labelled.
  % No coordinate and no colour changed to get this.
  \begin{scope}[shift={(10.50,8.45)}, rotate=38]
    \fill[accent, fill opacity=0.12] (-0.95,-0.16) rectangle (0.95,0.16);
  \end{scope}
  % v11: the label drops 8.14 -> 7.98. At 8.14 the word "raw" rose to y = 8.352
  % and sat ON the context band's 8.25 top edge and its outline.
  \draw[furn] (9.68,8.06) -- (9.98,7.95);
  \node[alab, anchor=base east] at (9.62,7.98) {$V_{\textrm{drug}}$, raw};

  % the context subspace: cell line, plate, dose. panelbg with a rulegrey
  % outline, which is the colour figs/slab.tex gives this same named object.
  % v11: both label lines drop 0.14cm -- see the restack note at V_pure.
  \filldraw[fill=panelbg, schem] (8.60,8.25) rectangle (11.80,8.65);
  \node[qlab] at (11.95,8.46) {context subspace:};
  \node[qlab] at (11.95,8.16) {cell line, plate, dose};

  % V_pure, at 90 degrees: what the raw slab becomes, and visibly thinner than
  % the raw band, against a measured 0.620 -> 0.031 printed below it.
  % v11 RESTACK, and it is one move with four dependents. The gloss sat 0.36pt
  % under the panel head's descenders -- 0.13mm, against 2.8mm of air under
  % panel (a)'s head -- so two panels of equal rank were set with head
  % clearances differing by a factor of twenty-two. Dropping the gloss 0.21cm
  % gives it 3.0mm, but it then landed on the right column's own first baseline
  % (9.32) with 0.15cm of horizontal air beside it, and the two read as one
  % line: "V_pure, the purified slab orthogonalise: remove the part". So the
  % right column drops 0.22cm, the context-subspace pair 0.14cm behind it, "the
  % confound" 0.18cm behind that, and the measurement strip and the two prose
  % blocks of BOTH panels 0.10cm behind that -- each step set by the measured
  % ink boxes, not by eye, at a uniform 0.066cm of white between blocks. The
  % bar's top comes to 9.15, still 0.115cm above the removal arrow's height,
  % which is all it has to do.
  \fill[accent] (10.41,7.95) rectangle (10.59,9.15);
  \node[alab, anchor=south] at (10.50,9.20) {$V_{\textrm{pure}}$, the purified slab};

  % THE CONFOUND: the orthogonal projection of the raw direction onto context,
  % from V_pure's right edge to the foot of the perpendicular dropped from the
  % raw band's tip, so nothing bisects it and it is exactly the length of the
  % removal arrow above it. Directly labelled, with a leader, because a cold
  % reader attached the name to the diagonal band instead of to this patch.
  % v11: the leader's tip now lands INSIDE the patch. In v10 it stopped at
  % (11.05,8.31), 0.03cm below the patch's lower edge and inside the context
  % band, so it pointed at the band -- which is precisely the misreading the
  % pass-9 rebuild of this mark exists to kill.
  \filldraw[fill=accent, fill opacity=0.40, draw=accent, line width=0.5pt]
    (10.59,8.33) rectangle (11.2486,8.57);
  \draw[furn] (11.30,7.88) -- (11.02,8.42);
  \node[alab] at (11.35,7.82) {the confound};

  % the projection line that defines it. Densely dotted quietink at 0.5pt is
  % this document's reserved stroke for a projection line, a magnification
  % leader or a zero rule -- a coordinate that decides nothing.
  \draw[quietink, line width=0.5pt, densely dotted]
    (11.2486,9.0349) -- (11.2486,8.57);

  % THE OPERATION, drawn as the projection it is and not as a rotation: a
  % straight arrow at the tip's own height, carrying the tip leftwards by
  % exactly the confound's length, landing on V_pure. It is the ONLY arrowhead
  % in the figure, and an arrowhead here means an operation acting in a
  % direction.
  \draw[accent, line width=0.5pt, -{Latex[length=3.4pt,width=2.8pt]}]
    (11.2486,9.0349) -- (10.61,9.0349);
  % v11: the leader now TOUCHES the arrow's tail. In v10 its lowest ink stopped
  % 0.105cm short of the tail and died in white, while the standoff at the
  % LABEL end was only 0.05cm -- the inversion of the corpus rule this file's
  % own header states twice (the 0.06cm gap goes at the label end, never at the
  % mark end, or the stub reads as a stray hairline).
  \draw[furn] (11.89,8.95) -- (11.2486,9.0349);
  \node[alab] at (11.95,9.10) {orthogonalise: remove the part};
  \node[alab] at (11.95,8.80) {of the raw slab lying in context};

  % ORTHOGONALITY AS A MARK, NOT AS A WORD, and v11 moves it out of the
  % removal arrow's line. In v10 the glyph's horizontal arm sat at y = 9.05 and
  % the arrow's shaft at y = 9.0349 -- 0.13mm apart, on opposite sides of an
  % opaque 0.18cm bar -- so a static assertion (perpendicularity) and a directed
  % operation (the projection) read as ONE horizontal primitive running from the
  % word "orthogonal", through V_pure, out to the dotted drop line. That is the
  % glyph collision S11 exists to forbid, and it was the only one left in the
  % figure. The arms come 0.40 -> 0.28cm and the arm drops to y = 8.93, which
  % puts 1.05mm of white under the arrow; 2.8mm is still 40% larger than the
  % 2mm a cold reader did not see until zooming, and the quadrant is unchanged.
  % THE QUADRANT IS UPPER-LEFT and always has been: v10 item 8 says "lower-left"
  % and is false about its own drawing. The other three quadrants are ruled out
  % by measurement -- lower-left and upper-right are crossed by the raw band at
  % 38 degrees, and lower-right is where "the confound" and its leader sit.
  \draw[ink, line width=0.5pt] (10.13,8.65) -- (10.13,8.93) -- (10.41,8.93);
  \node[lab, anchor=base east] at (9.99,8.73) {orthogonal};

  % the four measured numbers of stage two, as one strip on the panel's own x
  % origin. These are the only quantities above panel (c) with a source.
  % v11, two repairs. (i) COLOUR: the strip was ink and panel (a)'s scope line
  % is quietink, on the SAME baseline in adjacent panels at the same size, with
  % nothing on the page decoding the difference; both are scope, so both are
  % quietink now and ink is left to heads, definitional prose, the axis label
  % and the gate. (ii) NAME: "subspace fraction" becomes fig:mechanism's own
  % name for the quantity, "descriptive subspace fraction" (one object, one
  % name). The layer-9 qualifier stays and is load-bearing: layer 9 is also the
  % MAXIMUM of variance_share.drug over the seven depths (the range is 0.277 at
  % layer 16 to 0.620 here), which the drawing has no lane to state and which
  % the caption and figs/purify.json now carry.
  \node[qlab] at (8.40,7.48)
    {Dose decodes from inside the raw slab at $0.90$--$0.95$;};
  \node[qlab] at (8.40,7.20)
    {descriptive subspace fraction at layer 9: raw $0.620$, purified $0.031$.};

  \node[def] at (8.40,6.62) {$V_{\textrm{pure}}$ is the part of the raw slab that is};
  \node[def] at (8.40,6.22) {orthogonal to the context subspace.};

\end{scope}

% ###########################################################################
% # PANEL (c) -- the removal gate.  x origin 1.60, the fraction axis's spine.
% # The only measured marks in the figure.
% ###########################################################################
  \node[head] at (1.60,5.20) {(c) the removal gate};

  % WHICH SLAB THE GATE IS MEASURED ON, restored to the drawing in v11 and set
  % in axis-label rank rather than at the foot. Without it the last line a
  % reader crosses before entering panel (c) is panel (b)'s "so causal claims
  % use the purified slab, not the raw", and the inference that the gate was
  % run on V_pure is unavoidable. workspace_probe.json does not record which
  % subspace was projected out; it stores one ablated series and no purified
  % counterpart, so the line states only what the file supports.

  % the region of outcomes that would have made every null in this section a
  % measurement of nothing. Drawn FIRST so every rule and dot sits on top of
  % it, at rulegrey!8 (grey 247) and not rulegrey!12 (grey 243): at 12 it was
  % within one grey level of panelbg, the tint that means "a schematic
  % subspace" 3cm above it in panel (b), so one tint carried two meanings.
  % Nothing is written inside it: its emptiness is the whole point, and the
  % two italic lines that used to sit in it competed with the marks. It is
  % NAMED instead, outside its own right edge and under the rule that bounds
  % it, because the largest mark on the page may not be the one mark a reader
  % cannot name (and because the name makes the third clause of this figure's
  % acceptance test answerable from the drawing).
  \fill[rulegrey!8] (1.60,1.75) rectangle (14.10,\fy{0.8});

  % ---- fraction axis, DEFINED AT THE AXIS ---------------------------------
  % Rotated 90 degrees reading upward, immediately left of the tick numerals,
  % which is fig:mechanism's form and the answer to the supervisor's first
  % complaint. Nothing else is in this block: no result, no interval note, no
  % cross-reference.
  % v11, AND THIS IS A CORRECTNESS REPAIR AND NOT A TYPOGRAPHIC ONE. v10 set
  % one \footnotesize line reading "fraction of drug signal removed" over one
  % \scriptsize construction line. Two faults, both fatal.
  %   (a) THE NAME WAS FALSE. The plotted series is frac_signal_removed, which
  %       is ABOVE-CHANCE signal: (before - after)/(before - chance). The
  %       unqualified quantity (before - after)/before is 0.7041 at layer 2 and
  %       0.8528 at layer 9, four to eighteen points BELOW the marks drawn. Read
  %       at the drawn name the top mark, 0.949, implies a residual accuracy of
  %       0.041, under the chance 0.083 this figure prints in its own foot note.
  %       The thesis's pre-registration (04b:218) says "the fraction of
  %       above-chance signal removed must exceed $0.8$". v10's stated reason for
  %       the omission -- that the 43-character form needs 7.14cm against a
  %       5.71cm lane -- ruled out one STRING, not the correct QUANTITY, and is
  %       superseded: the name is broken over two rotated \footnotesize lines.
  %   (b) THE SUBORDINATE LINE WAS THE TALLEST OBJECT IN THE PANEL. Measured on
  %       the v10 render, the \scriptsize construction line ran 5.52cm against a
  %       3.85cm axis (143%) and its top sat 6.0pt ABOVE the panel head's, so the
  %       smallest type in the block outranked the head. S5 makes that line
  %       optional; it is deleted, and the construction it carried moves to the
  %       foot note, where the three accuracies it is built from already sit and
  %       where the missing noun ("probe accuracy") costs nothing.
  % LANE, remeasured on the v11 render and reconciled with figs/purify.json,
  % which the v10 header disagreed with: \footnotesize sets at 0.158 cm per
  % character in this class, so "fraction of above-chance" (24 characters) runs
  % 3.86cm against a 3.00cm axis and "drug signal removed" (19) runs 3.03cm.
  % The block is centred on the axis midpoint, y = 3.25.
  \draw[axisr] (1.60,1.75) -- (1.60,4.75);
  \foreach \v/\t in {0/0, 0.5/0.5, 1/1.0}{
    \draw[furn] (1.60,\fy{\v}) -- (1.50,\fy{\v});
    \node[qlab, anchor=east, inner sep=0pt] at (1.44,\fy{\v}) {\t};}
  % LINE ORDER. Rotated 90 degrees the line-advance direction is -x, so the
  % FIRST line is the leftmost and the LAST line is the one against the tick
  % numerals, which is also matplotlib's own stacking for a two-line ylabel.
  % Set the other way round the block reads "drug signal removed / fraction of
  % above-chance", which is what the first v11 build printed.
  \node[text=ink, font=\footnotesize, rotate=90, anchor=center]
    at (0.36,3.25) {fraction of above-chance};
  \node[text=ink, font=\footnotesize, rotate=90, anchor=center]
    at (0.80,3.25) {drug signal removed};

  % ---- layer axis ---------------------------------------------------------
  % "layer" is left-aligned to the axis's own left end, per the contract, and
  % v11 puts 2.8mm of white under the tick numerals instead of v10's 0.63mm.
  % At 0.63mm the word sat directly under the numeral 2, one size LARGER than
  % it, and the pair read as "layer 2" -- the identical fault the pass-9 note
  % above records against the old placement under the layer-9 numeral. The
  % repair is vertical, not horizontal: the fraction origin rose 1.45 -> 1.75,
  % which is where the 2.8mm came from. Baseline to baseline the numerals and
  % the word are now 0.53cm apart, against the 0.40cm leading this figure uses
  % between two lines of one sentence, so they can no longer read as one block.
  \draw[axisr] (1.60,1.75) -- (14.10,1.75);
  \foreach \l in {2,4,6,8,9,12,16}{
    \draw[furn] (\lx{\l},1.75) -- (\lx{\l},1.65);
    \node[qlab, anchor=north, inner sep=0pt] at (\lx{\l},1.59) {\l};}
  \node[def] at (1.60,0.85) {layer};

  % ---- the gate: a pre-set threshold, so BLACK, SOLID, 0.9pt --------------
  % It was accent 0.6pt dashed. Accent in this document means the object under
  % test, and a threshold is not the object under test.
  \draw[gate] (1.60,\fy{0.8}) -- (14.10,\fy{0.8});
  \node[lab, anchor=west, text=black] at (14.25,\fy{0.8}) {gate, $0.8$};
  % the name of the shaded region, set outside it and directly under the rule
  % that bounds it. Quietink, because the region is a counterfactual and not
  % the object under test.
  \node[qlab, anchor=west] at (14.25,2.95) {an untestable null};

  % ---- the seven measured fractions. No connecting line: the series is flat.
  \gdot{2}{0.8846153846153847}
  \gdot{4}{0.9253112033195021}
  \gdot{6}{0.9383561643835617}
  \gdot{8}{0.9461077844311377}
  \gdot{9}{0.9491525423728814}
  \gdot{12}{0.9440993788819876}
  \gdot{16}{0.9444444444444445}
  % direct label for the set, on the last mark's own baseline, with the 0.30pt
  % leader the corpus requires beyond 3mm: the gap from the layer-16 dot to the
  % label's ink is 3.6mm, and in v10 that gap was open with nothing in it.
  \draw[accent, line width=0.3pt] (14.19,\fy{0.9444444444444445})
                              -- (13.95,\fy{0.9444444444444445});
  \node[alab, anchor=west] at (14.25,\fy{0.9444444444444445})
    {all seven layers: $88$--$95\%$};

% ###########################################################################
% # the figure's foot: one two-line \scriptsize block on the figure's own x
% # origin. Line one names the instrument and the three accuracies the axis is
% # built from; line two gives the construction and this figure's single
% # interval declaration.
% # v11: the noun was missing. v10 printed "Before 0.41--0.82, after 0.121 ...,
% # chance 0.083" with no quantity attached to any of the three, and the words
% # "accuracy", "probe" and "1/12" appeared nowhere in the figure, so panel (c)
% # never said what instrument it was reading. The registry name is "probe
% # accuracy" and chance is "0.083 = 1/12", which is the stored float exactly
% # (workspace_probe.json :: chance == 1/12 in IEEE double).
% ###########################################################################
  \node[qlab] at (0.00,0.32)
    {Probe accuracy before $0.41$--$0.82$, after $0.121$ at every layer,
     chance $0.083$, exactly $1/12$.};
  \node[qlab] at (0.00,0.02)
    {The fraction is (before minus after)$/$(before minus chance). Intervals:
     none returned in (c); (a) and (b) carry no measured mark.};

  % the bounding box IS the fullwidth measure; nothing may stick out of it
  \path (0,-0.10) rectangle (17.30,10.35);
\end{tikzpicture}
%
% ---------------------------------------------------------------------------
% READY TO PASTE into Sections/04b-representation.tex, at sec:res-used,
% immediately after the paragraph "Step three, the removal gate"
% (04b-representation.tex:182-190) and BEFORE tab:workspace. The same block is
% staged in figs/_PENDING_FLOATS.tex for a human to integrate.
%
% Do NOT merge this float with fig:hook. The two are read six paragraphs apart;
% each is about 100mm tall against a corpus ceiling of 117mm, so stacked they
% spill the page; and the house caption shape carries one bolded finding, not
% two.
%
% CAPTION, v10 AND v11. The drawing lost eight sentences in v10 and one more in
% v11, and every one is folded into the caption below, because in this document
% the argument lives in the caption and the drawing names objects.
% What was cut from the drawing, in the order it used to read:
%   1. the three \footnotesize italic panel subtitles -- "where the twelve drugs
%      differ, at one layer", "the raw slab is partly a batch slab", and "the
%      fraction of above-chance drug signal that projecting the drug subspace
%      out destroys". The third is now the axis label, which is where it always
%      belonged; the other two are caption clauses.
%   2. the centred divider sentence, "an ablation is only a test of the drug if
%      the slab it removes really holds the drug", which restated the body
%      paragraph directly above the float.
%   3. the two lines inside the fail region, "had any layer landed in here, the
%      null below it / would have been unfalsifiable --- a measurement of
%      nothing", which restated the caption's own "the shaded region ... is
%      empty" and sat inside the data area competing with the marks. v11 gives
%      the region the NAME "an untestable null", set outside its right edge;
%      the argument stays in the caption.
%   4. the centred discipline line, "the dots and the bands are schematic and
%      have no scale; only the four numbers in stage two are measured". Its job
%      is done by the GLYPH now -- open ring against filled disc -- and the
%      caption says the rest.
%   5. the conservatism argument, "so the slab keeps essentially the whole span,
%      noisiest included, which for a removal claim is conservative", cut from
%      panel (a)'s prose, which is now definitional only.
%   6. "by construction", cut from the right angle's label for want of a lane
%      3.6cm wide at the glyph's own height (v10 item 8).
%   7. "measured on the slab that is projected out; the stored file has no
%      separate purified-slab gate", cut from the foot note in v10 for text
%      density. RESTORED TO THE DRAWING IN v11, reworded, under panel (c)'s
%      head, because without it panel (c) named no slab at all and sat directly
%      under panel (b)'s "so causal claims use the purified slab, not the raw".
%      The caption's own statement of the same limit is reworded below so that
%      no clause of five or more words appears in both places.
%   8. "four \quad what the ablation costs: \cref{fig:hook}", the forward
%      pointer that used to sit on the x-tick baseline where the word "four"
%      read as a leftmost tick.
%   9. NEW IN v11: the axis's \scriptsize construction line, "(before minus
%      after)/(before minus chance)". It did not vanish -- it is the second foot
%      line now -- but it left the axis block, which S5 allows and which the
%      render forced: at 5.52cm against a 3.85cm axis it was the tallest object
%      in the panel and rose above the panel head's own top.
%
% CLAUSES OF THE PRINTED CAPTION THAT ARE NOW FALSE AGAINST THE DRAWING, with
% their lines in Sections/04b-representation.tex. Nobody may edit that file from
% this pass, so they are listed here and in figs/purify.json, and the
% replacement for each is in the caption below.
%   04b:234  "\emph{Above the rule, nothing has a scale.}" -- there is no rule.
%            The full-measure divider at y = 4.72 was deleted in v10 and the
%            panel letters carry the scoping now.
%   04b:234  "The dots, the bands, the arrow and the right angle are schematic"
%            -- v10 turned panel (a)'s twelve schematic marks into OPEN RINGS,
%            so the only dots in the figure are panel (c)'s seven MEASURED
%            discs, and the sentence as printed now calls the measurements
%            schematic. Replacement: "open rings and untitled bands".
%   04b:236  "the four numbers printed in stage two" -- there is no stage two;
%            the panel is lettered (b) and the word "stage" went with the
%            banners.
%   04b:243  "\emph{Below the rule, everything has a scale.}" -- same rule.
%   04b:250  "Two limits, both stated on the drawing" -- only one of the two is
%            on the drawing (the which-slab limit, restored in v11); the
%            twelve-class against twenty-four-class limit is caption-only and
%            always was.
%   04b:250-252  the which-slab limit as printed shares a clause of more than
%            five words with the drawing's new scope line, and it asserts more
%            than the file supports; it is reworded below.
% AND ONE THING THE CAPTION OWES THE READER AND HAS NEVER SAID: the printed
% pair "raw $0.620$, purified $0.031$" is at layer 9, and layer 9 is the
% MAXIMUM of layers.<L>.variance_share.drug over the seven depths. The stored
% range is 0.2766847391691891 at layer 16 to 0.6201207831136324 here, so the
% drawn pair is the largest raw-to-purified contrast the file holds. The raw
% series is drawn in no other figure and quoted in no sentence of the thesis
% (figs/fig-mechanism.json records it as drug_raw_not_drawn), so the caption is
% the reader's only chance to learn the spread. It is added below.
% The dose qualifier the v10 proposal added is kept: "$0.90$--$0.95$" is true
% only to two places, the stored minimum being 0.8979166666666666 at layers 12
% and 16, which is also why no reference rule is drawn at $0.90$.
%
% NOTHING IN THE CAPTION RESTATES A CLAUSE OF FIVE OR MORE WORDS FROM THE
% DRAWING. Checked string by string against the v11 render's own text spans: the
% drawing names objects and prints numbers, the caption argues. The drawing says
% "the subspace projected out"; the caption says "whichever subspace the
% ablation removed", and the longest run shared with any drawn string is four
% words.
%
% Two caption details, checked against a build of the whole float (returncode 0,
% one page, zero Overfull/Underfull boxes; the \cref targets resolve in the
% document and show as ?? only in an isolated harness):
%   * the two numerals in the bold lead are set in TEXT and not as $88$/$95$.
%     In math mode they come out upright and unbolded inside \textbf, so a
%     bolded finding would have had two unbolded numbers in the middle of it.
%   * the FIGURE BODY carries no \cref. The "four ... fig:hook" pointer was
%     deleted in v10, so the caption is the only place the sibling float is
%     named.
%
% \begin{figure}[htbp]
% \begin{fullwidth}
% \centering
% % ===========================================================================
% figs/purify.tex -- fig:purify, how the drug slab is BUILT, and the gate it
% has to pass before anything may be ablated with it.
%
% FIGURE BODY ONLY. No preamble, no document environment, no float, no caption.
% \input it from inside a figure/fullwidth float in
% Sections/04b-representation.tex; the ready-to-paste float sits commented at
% the foot of this file and is also staged in figs/_PENDING_FLOATS.tex.
% Compiles against thesis_v2/preamble.tex and nothing else: no \includegraphics,
% no external data, no TikZ library beyond the ones preamble.tex already loads
% (this file uses arrows.meta and calc only).
%
% WHAT THE FIGURE IS
% ------------------
% Steps one to three of sec:res-used, drawn as steps. Step one localises the
% drug: twelve centred class means, ten kept directions. Step two purifies it:
% the raw slab is orthogonalised off the context it is confounded with. Step
% three is the falsifiability check: projecting the slab out must destroy most
% of the decodable drug signal, or every null downstream is a measurement of
% nothing. fig:hook (a separate float, six paragraphs later) draws step four;
% the two are deliberately NOT merged -- each is about 100mm tall, the corpus
% ceiling is slab.tex at 117mm, and the house caption shape carries one bolded
% finding, not two.
%
% EVERY NUMBER DRAWN, AND WHERE IT COMES FROM
% -------------------------------------------
% Source for every measured quantity: RESULTS_cluster/workspace_probe.json.
% Config there: tier tier2_unseen_drugs, n_drugs 12, n_per_drug 40, n_dims 10,
% layers 2,4,6,8,9,12,16, chance 0.08333333333333333, seed 42.
% Sibling ledger with the full machine-precision values: figs/purify.json.
%
% (a) THE SEVEN PLOTTED DOTS -- the only marks in this figure with a scale.
%     layers.<L>.frac_signal_removed, layers 2,4,6,8,9,12,16:
%       0.8846153846153847  0.9253112033195021  0.9383561643835617
%       0.9461077844311377  0.9491525423728814  0.9440993788819876
%       0.9444444444444445
%     layers.<L>.gate_pass is true at all seven; the dots' position above the
%     0.8 rule is that fact, so no second mark asserts it.
%
% (b) THE GATE, 0.8. PROSE ONLY. Sections/04b-representation.tex:185 "the
%     fraction of above-chance signal removed must exceed $0.8$". There is no
%     config field for it in workspace_probe.json, so the ledger cites the .tex
%     line and not a key path. Drawn dashed accent, the house register for a
%     pre-set threshold a value is tested against (gate.tex:170-174).
%     SUPERSEDED IN v10 BELOW: it is now black, solid, 0.9pt.
%
% (c) THE THREE PRINTED ACCURACIES in the right-hand block of band two.
%       before  0.41--0.82   layers.2.drug_decode  = 0.4083333333333334
%                            layers.9.drug_decode  = 0.8208333333333334
%       after   0.121        layers.<L>.drug_decode_ablated
%                            = 0.12083333333333332, byte-identical at all seven
%       chance  0.083        top-level chance = 0.08333333333333333
%     Thesis 04b:186-188 "Accuracy falls from $0.41$--$0.82$ to a constant
%     $0.121$ floor, removing $88$--$95\%$ of the above-chance signal".
%     SUPERSEDED IN v10 BELOW as to PLACEMENT only: the three values are no
%     longer a right-hand gutter block, they are the figure's foot note. No
%     value changed.
%
% (d) THE IDENTITY, verified numerically before drawing, in python, over all
%     seven layers, exact float equality (not a tolerance):
%       (drug_decode - drug_decode_ablated) / (drug_decode - chance)
%         == frac_signal_removed
%     Layer 2: (0.4083333333333334 - 0.12083333333333332)
%              / (0.4083333333333334 - 0.08333333333333333)
%              = 0.8846153846153847, bit for bit.
%     WHICH SLAB THE GATE IS MEASURED ON. This identity fixes it, and the
%     answer is NOT V_pure. workspace_probe.json stores exactly one ablated
%     series, drug_decode_ablated, and the identity reproduces from the RAW
%     decode. The section-4.5 arm files prove the pipeline can distinguish the
%     two cases -- they store drug_decode_ablated AND drug_decode_pure_ablated,
%     frac_signal_removed AND frac_signal_removed_pure -- and this file stores
%     only the first of each pair. So the figure, the caption and the ledger all
%     say "the slab that is projected out" and nowhere claim the gate was run on
%     V_pure. Stage two's arrow lands on V_pure because that is where the CAUSAL
%     claims are made (04b:179-180); band two is a separate measurement and the
%     divider is what keeps the two from being read as one.
%     SUPERSEDED IN v10 BELOW as to the LAST CLAUSE ONLY: the divider is gone
%     and the panel letters do that work now. The caveat itself survives, in the
%     foot note.
%
% (e) THE FOUR PRINTED NUMBERS IN STAGE TWO, the only measured quantities above
%     the divider.
%       0.90--0.95   layers.<L>.confound_crossdecode.dose, seven layers:
%                    0.9020833333333333 0.9229166666666666 0.9416666666666667
%                    0.9395833333333332 0.95 0.8979166666666666
%                    0.8979166666666668. Thesis 04b:176-178.
%                    PRINTED AS TEXT AND NOT DRAWN, and there is deliberately NO
%                    0.90 reference rule anywhere in this figure: the stored
%                    minimum is 0.8979166666666666 at layers 12 and 16 and
%                    reaches the thesis's "0.90" only by rounding to two places.
%                    A rule at 0.90 would put two of seven points below it and
%                    would read as a failure the thesis does not report.
%       raw 0.620    layers.9.variance_share.drug      = 0.6201207831136324
%       purified     layers.9.variance_share.drug_pure = 0.030620247335072275
%         0.031
%     The RAW share is quoted in no sentence of the thesis; figs/fig-mechanism.json
%     already records it as drug_raw_not_drawn. Printing it here is a new reading
%     of a stored field, flagged as such in purify.json, on the precedent of
%     figs/slab.json's handling of variance_share.context. Layer 9 is named in
%     the figure text so the pair cannot be read as a seven-layer summary; it is
%     the layer at which the raw probe peaks (04b:170-171), which is why the
%     section uses it as its own reference depth.
%
% (f) THE COUNTS in stage one: 480 prompts, twelve drugs, forty cells each,
%     seven layers, ten kept directions, span at most eleven.
%     config.n_drugs 12, config.n_per_drug 40, config.n_dims 10,
%     layers.<L>.pure_drug_dims 10, config.layers "2,4,6,8,9,12,16".
%     480 = 12 x 40 and the rank-11 bound (twelve centred means span at most
%     eleven directions) are ARITHMETIC, not stored fields; the ledger says so.
%     Thesis 04b:164-171.
%
% AGREEMENT WITH THE TEXT (checked line by line, not assumed):
%   04b:164-165  "\num{480} tier-2 prompts (twelve drugs, forty real cells
%                each) at seven layers"            -> stage one, line 4.     OK
%   04b:168-170  "Twelve centred class means span at most eleven directions,
%                so the slab keeps essentially the whole span, noisiest
%                included, which for a removal claim is conservative"
%                -> stage one, lines 1-3, verbatim.                          OK
%   04b:175-178  "The raw drug slab is partly a batch slab. Drugs are
%                confounded with plate and dose by the atlas's design, and
%                probing dose from inside the drug subspace recovers it at
%                $0.90$--$0.95$ at every layer"    -> stage two, lines 1-3.  OK
%   04b:179-180  "Every causal claim below is made on $V_{\textrm{pure}}$,
%                never on the raw slab."           -> the V_pure callout.    OK
%   04b:185-186  "the fraction of above-chance signal removed must exceed
%                $0.8$; below that the null is unfalsifiable, a measurement
%                of nothing"                       -> the dashed rule, and
%                                                     G4's two italic lines,
%                                                     which are the section's
%                                                     own words.             OK
%   04b:186-188  "Accuracy falls from $0.41$--$0.82$ to a constant $0.121$
%                floor, removing $88$--$95\%$"     -> right-hand block and
%                                                     the set label.         OK
%   04b:212-213  tab:workspace caption "All seven layers pass the removal gate
%                (88--95\% ...); chance is $0.083$"                          OK
%
% AVAILABLE, DELIBERATELY NOT DRAWN, recorded here and in purify.json so a
% reviewer who asks for direct evidence of the batch overlap is answered
% without a redraw:
%   * layers.<L>.overlaps.context   0.5215396547181934 -- 0.6323768458353303
%   * layers.<L>.overlaps.cell_line 0.3043896441494730 -- 0.3534060289850234
%     These are the per-layer measure of "the raw drug slab is partly a batch
%     slab" and are quoted nowhere in the thesis. Drawing them would put a
%     second measured scale into a panel whose whole discipline line says the
%     geometry has no scale, and would be a new claim rather than a redraw.
%   * layers.<L>.confound_crossdecode.plate, 0.4541666666666667 --
%     0.5770833333333333, a materially weaker and entirely undiscussed
%     confound. Not printed, because the section names dose and not plate.
%   * layers.<L>.swap.{drug_b_kl, random_inject_kl} in this same JSON is the
%     RETIRED steering design (07-Appendix.tex, sec:app-identity). Nothing from
%     it is on this figure. figs/mechanism.py's header records that the v1
%     figure plotted exactly that pair while its caption described the ablation.
%
% DESIGN DECISIONS, each with its reason
% --------------------------------------
%  1. TWO BANDS AND A DIVIDER THAT IS A STATEMENT. Above the rule at y=4.72
%     nothing has a scale; below it everything does. No mark crosses the rule,
%     no leader connects a schematic object to a plotted one, and the italic
%     sentence under the rule sets band two up as the answer to the objection
%     stage two raises -- if purification strips the raw slab from 0.620 to
%     0.031 of the stream's variance, does anything survive in it? The
%     alternative, one continuous diagram from cloud to axis, was rejected: it
%     would invite the reader to measure the ellipse.
%     SUPERSEDED IN v10 BELOW. The divider and its sentence are deleted; the
%     schematic/measured distinction is now carried by the MARK (open rulegrey
%     ring against filled accent disc), which is where it belongs.
%  2. THE FAIL REGION IS DRAWN, THE PASS REGION IS NOT. The one non-obvious
%     mark in band two is the rulegrey!18 rectangle from 0 to 0.8. It is the
%     region of outcomes that would have made every null in this section a
%     measurement of nothing, and it is empty. A figure that drew only the seven
%     dots and the rule would be a bar of good news; drawing the region the
%     dots had to avoid is what makes it a falsifiability picture. Its emptiness
%     is also where the two italic lines fit, which is not a coincidence and is
%     the reason they are placed there rather than in the caption.
%     PARTLY SUPERSEDED IN v10: the region survives, at rulegrey!12, and the two
%     italic lines inside it do not. The emptiness is the message and text in it
%     was competing with the marks.
%  3. NO CONNECTING LINE THROUGH THE SEVEN DOTS. The series is flat to 0.1355cm
%     on this scale (0.8846 to 0.9492 at 2.10cm per unit of fraction). A line
%     would assert a trend across a 6-point spread the axis cannot resolve; the
%     seven marks are named as one ordered set by "all seven layers: 88--95%"
%     instead, which is comparators.tex's rule for marks too close to leader
%     individually.
%  4. NO INTERVALS ANYWHERE. frac_signal_removed is a single ratio of two
%     accuracies on 480 prompts and carries no stored dispersion of any kind;
%     the KL series in the same file carry normal-approximation SEs, which are
%     not this document's interval convention (figs/fig-mechanism.json records
%     intervals_drawn false with that reason). A bare solid dot with no bar is
%     how this document says a quantity has no interval (comparators.tex:149-151
%     changed a hollow ring to a solid dot for exactly this).
%  5. THE ORTHOGONALISATION IS DRAWN AS AN OPERATION, NOT AS A RESULT, AND THE
%     OPERATION IS A PROJECTION AND NOT A ROTATION. This is the one mark
%     figs/slab.tex does not have. There the drug wedge and the
%     context wedge are simply shown already orthogonal, which states the
%     outcome and hides the step; here the raw band sits at 38 degrees, visibly
%     crossing the context band, and the raw slab's tip is carried onto V_pure.
%     REPAIRED, PASS 9. It used to be carried by an accent ARC from 38 to 76
%     degrees, i.e. the raw band swung upright. A reviewer read that panel as
%     "the slab is turned until it is perpendicular to context" and had to go to
%     the prose to learn that a component is REMOVED, and they were right to:
%     two bands of equal width (0.32cm) and near-equal length (1.90 against
%     2.00) differing only in angle say "same slab, rotated, nothing lost",
%     against measured numbers 1.8cm to the right that say raw 0.620, purified
%     0.031, a twentyfold shrink. A curved arrow is the universal sign for
%     rotation. What replaces it is the subtraction itself: a dashed rulegrey
%     perpendicular from the raw band's tip down to the context band, the
%     component it lands on drawn as its own named mark, and a STRAIGHT accent
%     arrow at the tip's own height carrying the tip left by exactly that
%     component, landing on V_pure. Two further faults the arc had, both from
%     the render: its head stopped 1.5mm short of the V_pure bar and pointed at
%     white space, and the grey leader from the word "orthogonalise" stopped 4mm
%     short of the arc's tail -- the corpus puts the 0.06cm gap at the LABEL
%     end, never at the mark end, so that stub read as a stray hairline. The
%     label now sits 0.15cm from the arrow's tail and carries no leader, and the
%     arrow is 0.6pt: a schematic operation, not a measurement, where 1.0pt is
%     the corpus's "this IS the claim" register. V_pure was also narrowed from
%     0.32 to 0.18cm, so the purified band is visibly thinner than the raw one.
%     slab.tex is left alone and is not
%     superseded: three of its four panels are the failed identity designs of
%     sec:app-identity, which are the appendix's argument, and its magnified
%     interior is sec:res-mechanism's swap geometry. What it draws that belongs
%     to sec:res-used is one thing, the variance-share disc, and this figure
%     does not redraw it. The two figures share vocabulary (accent-on-white
%     bands, rulegrey leaders, one discipline line saying what is to scale) and
%     no coordinate.
%     The whole of this decision survives v10 except the arrow's WEIGHT, which
%     goes 0.6 -> 0.5pt to bring the file inside a four-weight budget; the
%     perpendicular goes dashed-rulegrey -> densely-dotted-quietink, which is
%     this document's reserved stroke for a projection line.
%  6. THE CONFOUND IS A NAMED MARK, NOT THE PRODUCT OF TWO TINTS.
%     REPAIRED, PASS 9, and this is the repair both reviewers called blocking.
%     It used to be the lozenge where the raw band's tint multiplied over the
%     context band's, with the whole claim carried by the caption ("that
%     darkened region is the confound step two exists to remove"). Three faults:
%     (a) it was the only element in the panel with no direct label, while every
%     other mark had one; (b) V_pure ran vertically through its centre and split
%     it in two -- sampled on the render, a row across it read grey (174,174,174),
%     then (119,143,164) for the tint product, then a hard cut to opaque accent
%     (11,82,147) for 24px, then (119,143,164) again -- so it was not a nameable
%     region at all; (c) it made a CROSSING mean "confounded" on a panel where
%     another crossing, V_pure through the same band, has to mean "orthogonal".
%     One glyph cannot carry both. It is now the orthogonal projection of the
%     raw direction onto context: a segment along the context band from the
%     shared origin to the foot of the perpendicular dropped from the raw band's
%     tip, drawn in the raw slab's own tint at a higher opacity with a 0.35pt
%     accent!70 outline, and directly labelled "the confound". It begins at
%     V_pure's right edge, so nothing bisects it, and it is exactly the length
%     of the removal arrow above it, which is the rhyme that makes the arrow
%     legible. With the tint-multiply claim gone, a crossing means nothing
%     anywhere in the panel and the right-angle mark is the only assertion of
%     orthogonality, which is what decision 8 already wanted.
%     The context band is now panelbg with a rulegrey outline and not ink!35:
%     see decision 11.
%  7. NO 0.90 REFERENCE RULE. See (e) above: the stored dose minimum is
%     0.8979166666666666 and rounds to the thesis's 0.90. Printing the range as
%     text states exactly what the section states; a rule would show two of
%     seven points below it and would manufacture a failure.
%  8. RIGHT ANGLE AS A MARK, NOT AS A WORD. The small square in the
%     upper-left quadrant at the shared origin is the only assertion that
%     V_pure and the context subspace share no direction. The word
%     "orthogonalise" labels the operation; the square labels its result.
%     PASS 9: ink 0.35pt, up from rulegrey 0.3pt. The caption calls it "the one
%     geometric fact true by construction" and it was the faintest mark in the
%     panel.
%     AMENDED IN v10: the square moves to the LOWER-left quadrant, doubles to
%     0.40cm arms, and acquires the direct label it never had.
%     SUPERSEDED IN v11: THIS LINE IS FALSE ABOUT ITS OWN DRAWING AND ALWAYS
%     WAS. v10 drew the square in the UPPER-left quadrant of the crossing --
%     (10.01,8.65)-(10.01,9.05)-(10.41,9.05) -- and no build has ever put it
%     lower-left. The direct label and the enlargement are true. v11 keeps the
%     upper-left quadrant, which is the only one of the four that is free (the
%     raw band at 38 degrees crosses lower-left and upper-right, and "the
%     confound" and its leader own lower-right), and takes the arms 0.40 ->
%     0.28cm to clear the removal arrow -- see v11 item 2.
%  9. THE GLOBAL MEAN IS A CROSS, NOT A DOT. The twelve class means are dots;
%     the thing they are centred on must not be a thirteenth dot. It is also
%     leadered out of the cloud rather than labelled in place, because at
%     8pt the label box collided with two of the twelve dots.
% 10. V_drug IS DRAWN WITH STAGE TWO'S OWN GLYPH, AND ALL TWELVE DISPLACEMENTS
%     ARE DRAWN. REPAIRED, PASS 9, two findings at once.
%     (a) V_drug used to be a dashed, accent!10 FILLED ELLIPSE enclosing the
%     twelve dots, while stage two drew the same name as an unbounded band. A
%     bounded blob around the class means says "the region where the drugs
%     live"; a subspace is not a region and has no edge. Read cold, the ellipse
%     was taken for the data cloud and stage two's band for a new object, so
%     stage one's output was never connected to stage two's input -- the one
%     join the figure exists to make, and the join the brief says readers lose.
%     V_drug is now an accent band at the SAME fill opacity stage two uses,
%     through the global mean, running past the data on both sides. The ellipse
%     is deleted rather than kept unlabelled: with the band tall enough to
%     contain all twelve dots (half-height 0.52 against a maximum |dy| of 0.44)
%     it had nothing left to say, and a class mean drawn OUTSIDE V_drug would
%     have been false, since the twelve centred means span it by construction.
%     (b) ONE displacement was drawn and eleven were not, so "twelve are
%     stacked" -- the sentence directly under the dots -- was invisible, and the
%     dots were unlabelled while the prose two lines below names a different
%     population ("480 prompts: twelve drugs, forty real cells each"), so the
%     scatter read as prompts and the cross as their grand mean. All twelve
%     stems are now drawn at 0.3pt quietink, with the one accent arrow on top
%     and still the only labelled one, and the dots carry the tag "the twelve
%     per-drug means". The stems make the count a picture and they are what the
%     "at most eleven directions" line is about.
%     The twelve offsets are written out in a \foreach list so no build can
%     randomise them and no reader can be told a different story by a different
%     compile. Opacity: the band is 0.20, not the 0.34 of stage two's first
%     draft, and stage two's raw band was brought to 0.20 with it. At 0.34 over
%     5.10 x 1.04cm the band became the single heaviest area on the page, which
%     is the thing this figure was already being marked down for; at 0.20 over
%     4.10 x 1.04 it lays down about the ink of stage two's 1.90 x 0.32 band and
%     the two are drawn at the SAME opacity, so the glyph rhymes exactly.
%     AMENDED IN v10: the twelve class means become OPEN RINGS and the one
%     accent displacement arrow loses its head -- see v10 items 3 and 6.
% 11. ONE GREY FAMILY, AND THE CONTEXT SUBSPACE IS THE COLOUR slab.tex GIVES IT.
%     PASS 9. Three greys were in play, derived two different ways: rules at
%     rulegrey, the context band at ink!35, the gate region at rulegrey!18.
%     ink!35 is not derived from rulegrey and is LIGHTER than a hairline rule,
%     so the band read as a heavy solid slab while the rules read as thin, and
%     it was the darkest region on either page. slab.tex draws this same named
%     subspace in panelbg (its \slabdisc, slab.tex:83-90), so the figure also
%     disagreed with the published figure about what colour the object is. The
%     band is now panelbg with a 0.35pt rulegrey outline; ink!35 no longer
%     appears in the document, and purify's context band and slab.tex's context
%     wedge are the same object drawn twice.
% 12. THE GATE PANEL KEEPS ITS FULL 0-TO-1 AXIS. A reviewer asked for the range
%     to be cut to 0.75-1.00 so the seven fractions (0.8846 to 0.9492, 0.1355cm
%     apart on this scale) would resolve, against a fail region 12.4x their
%     height. OVERRULED, and the reason is the same one that bans bars on a log
%     axis: the fail region's height is not a drawing choice, it is the share of
%     the outcome space that would have made every null below unfalsifiable, and
%     0 is the natural floor of a fraction, not an arbitrary cut. Cutting the
%     axis at 0.75 would shrink the one region whose SIZE is the argument, to
%     resolve an ordering the thesis makes no claim about -- tab:workspace is
%     the authority for the per-layer values and the section quotes only the
%     range. What was applied instead: the direct label "all seven layers:
%     88--95%" names the set (decision 3), and the 0.8 numeral was dropped from
%     the tick set because the accent label already owns that coordinate.
%     STILL IN FORCE IN v10, and it is the first thing v10 was told to protect.
%
% WHAT SIX RENDER-AND-READ PASSES CHANGED, with the measurement that forced it
% ---------------------------------------------------------------------------
%  i.   PASS 1, Overfull 88.29pt. The discipline line under stage two ran 146
%       characters at \footnotesize and measured 20.4cm against a 17.30cm
%       measure. Cut to 103 characters; the sentence now states only what it
%       has to, that the geometry has no scale and four numbers do.
%  ii.  PASS 1. "every causal claim is made here", set anchor=south at y=9.22,
%       put its box top at 9.63 against the stage-two subtitle's ink bottom of
%       9.37 and the two overprinted. There is no room above V_pure for a
%       second line, so the sentence became the fourth line of stage two's
%       prose block, in the section's own words ("never on the raw slab").
%  iii. PASS 1. "the global mean" and "$\mu_d-\bar\mu$, one drug's" collided
%       character to character: both wanted the space up and right of the
%       cloud, which is where the displacement arrow points. The global mean
%       was moved BELOW the cloud on a leader. Drawn straight down that leader
%       read as a dropped axis, so PASS 4 gave it a 0.25cm dogleg; the dogleg
%       clears the two nearest class means by 0.23cm and 0.25cm.
%  iv.  PASS 2, Overfull 6.20pt, unchanged from PASS 1 by three edits, which is
%       what identified it: the culprit was not the widest LINE but the
%       fraction axis's own tick label. "1.0" is anchor=east on a spine at
%       x=0.30, and its box ran 0.28cm LEFT of x=0. The bare \path pins a
%       minimum box, not a maximum, so nothing clipped it. The whole plot was
%       shifted right (spine 0.30 -> 0.75, layer 2 from 0.60 -> 1.05) rather
%       than the label being clipped or shrunk. boxes=0 from PASS 4 on.
%  v.   PASS 2. Band two carried two annotation columns 0.30cm apart at the
%       same baseline, so "88--95%" and "accuracy before" read as one line. The
%       plot was widened from 8.60 to 11.00cm of layer axis and the two columns
%       became one, which also removed 3.4 x 1.9cm of dead white under them.
%  vi.  PASS 4, at 520dpi. The leader to "orthogonalise" left the arc 0.20cm
%       from its own arrowhead and read as a grey continuation of the accent
%       arrow. It now leaves at 50 degrees, 0.42cm clear of the head, and stays
%       above the raw band's upper edge for its whole length. The right-angle
%       mark went from 0.16 to 0.20cm, and the global-mean cross from 0.30 to
%       0.45pt with 2mm arms, because at 0.3pt it was fainter than the leaders.
%  vii. PASS 5. Stage one left 4.2 x 1.9cm of white to the right of the cloud
%       while stage two's matching region was full. The ellipse went from
%       1.75 to 2.05 x-radius and the cloud from 1.45 to 1.72, which both fills
%       the panel and raises the aspect from 2.34:1 to 2.77:1. The ellipse's
%       y-radius was NOT raised with it: at 0.98 the ellipse reached y=6.92 and
%       left the global-mean label 0.06cm from the prose block below.
%  viii.PASS 5. The panel-three subtitle said "projecting the slab out", which
%       after stage two's "every causal claim is made on V_pure" invited the
%       reading that the gate was run on V_pure. It now says "projecting the
%       drug subspace out", the section's own noun at 04b:184-185, and the
%       string V_pure does not appear in panel three at all.
%
% LAYOUT NOTES, all forced by a build
% -----------------------------------
%   * fullwidth = textwidth 142mm + marginparwidth 28mm + marginparsep 4mm
%     = 174mm. (The "171mm" beside \newenvironment{fullwidth} in preamble.tex
%     is stale arithmetic.) The bounding box is pinned with a bare \path at
%     (0,-0.12) rectangle (17.30,10.20) = 173.0 x 103.2mm, inside the 17.35cm
%     hard edge circularity.tex records and under the corpus height ceiling,
%     slab.tex at 117mm. The harness reports pages=1, boxes=0. It grew 1.7mm in
%     pass 9: band two moved down 0.20cm to open a line under stage three's
%     subtitle for the "which slab" caveat, and the "four ... see fig:hook"
%     pointer sits on the bottom line. The float got much SHORTER overall in the
%     same pass, because the caption went from 484 words to about 280.
%   * x=1cm, y=1cm. Nothing is scaled, nothing is \resizebox'd; 1cm here is
%     1cm on the page.
%   * pgfmath: \lx and \fy expand to BRACED groups, so \lx{2}+0.05 is not a
%     coordinate. Every stand-off is either inside the braces or an xshift on
%     the node (gate.tex records the same trap).
%   * the tick rows use the \v/\t value/label pair form, so the fraction axis
%     prints 0.5 and 1.0 rather than \foreach's raw 0.5 and 1 tokens. The layer
%     row does not need it: its values are integers and print as themselves.
%   * no node uses align= or text width=: such a node builds a minipage whose
%     first-line height follows the AMBIENT \baselineskip, so its ink moves
%     with whatever prose precedes the float (frames.tex measured 3.7pt of
%     drift). Every multi-line block here is a stack of single
%     anchor=base west lines on explicit baselines: 0.28cm apart in band one's
%     two prose blocks, 0.33cm in band two's right-hand block, where the lines
%     are read one at a time rather than as a paragraph.
%   * the fraction axis carries no rotated title. The panel subtitle sits
%     directly above it and names the quantity, the direction and the frame in
%     one line, which is what gate.tex's fix 1 asks an axis to do; a second
%     title would repeat it 0.2cm away.
%     THIS NOTE IS NOW FALSE AND IS SUPERSEDED BY v10 ITEM 4. It was the
%     supervisor's first complaint about this arc of figures and it was wrong:
%     a quantity named in a panel subtitle 30mm from the ticks is not named at
%     the axis. The fraction axis carries a rotated title from v10 on.
%   * PASS 9, THE "WHICH SLAB" CAVEAT, three builds to place. The wording went
%     155 -> 121 -> 92 characters. At 155 the picture was 55.72pt over the
%     measure and the harness said so; at 121 it ran to x = 15.90 and
%     overprinted the accent label "all seven layers: 88--95%" at 12.60; at 92
%     it ends at 12.30. Vertically it was first set at baseline 3.24 and its
%     ascenders touched the subtitle's descenders with no white between -- the
%     subtitle is \footnotesize in a node whose TOP is pinned at 3.76, so its
%     descender line is 3.351 and a \scriptsize line at 3.24 reaches 3.444. It
%     is now at baseline 3.05 with 0.10cm of clearance above, and INDENTED to
%     x = 0.90: at x = 0 its descenders at 2.97 met the "1.0" tick numeral,
%     whose box top is 2.95, and the indent both clears the numeral and breaks
%     the note away from the subtitle so the two do not read as one paragraph.
%   * PASS 9, THE DIVIDER BETWEEN STAGES ONE AND TWO now ends at y = 5.28, the
%     descender line of the last prose row in both columns (baseline 5.38),
%     instead of at 5.40 where it stopped between two lines of type and read as
%     a broken rule. It clears the centred italic beneath it by 0.26cm.
%   * PASS 9, THE DISCIPLINE LINE IS CENTRED on the figure's own centre and is
%     no longer flush left at the prose column's x. Flush left, at the same
%     leading as the four-line block directly above it and with no blank
%     scan-line between, it parsed as that paragraph's fifth line even though it
%     ran full width past the vertical divider. Centred, it pairs with the
%     sentence centred under the rule, and the two bracket the divider.
%   * PASS 9, THE AXIS TITLE moved from x = 6.60 to 7.73. \lx{9} = 6.55, so
%     centred on 6.60 the word "layer" sat within 0.05cm of the layer-9 numeral,
%     one line below it, and read as a label of layer 9 rather than of the axis.
%     7.73 is the midpoint of the widest tick gap (\lx{9} = 6.55 to
%     \lx{12} = 8.91) and is 0.83cm clear of both.
%   * PASS 9, THE 0.8 TICK IS GONE. The coordinate was labelled twice at the
%     same height, once as a quietink numeral on the left spine and once as the
%     accent label "the gate: 0.8" on the right, so the plot appeared to carry
%     a left scale and a right scale. 0, 0.5, 1.0 is also a regular tick set.
%     A 0.9 tick was considered and rejected by measurement: at \fy it would
%     sit 0.21cm from both neighbours against a 0.203cm numeral height.
%
% TYPE. Ambient \small on the picture; \scriptsize (8pt) on every node except
% the \footnotesize (10pt) italic prose -- the three panel subtitles, the
% discipline line, the divider sentence and the two lines inside the fail
% region. Nothing is set at body size.
% SUPERSEDED IN v10 BELOW: two sizes only, \footnotesize and \scriptsize, and
% ZERO italic.
%
% ===========================================================================
% v10 -- 2026-08-18. REBUILD AGAINST THE STYLE CONTRACT AND THE SUPERVISOR'S
% VERDICT. Nine of the twelve design decisions above survive whole; the three
% that do not are superseded in place, above, rather than deleted, because the
% thesis is audited against this header.
%
% THE VERDICT THIS PASS ANSWERS, quoted: "For the other figures 4.8-12 go back
% over them and iron them out even more. Make it look more professional, right
% now it looks a bit sloppy and childish ... and moreover make sure that the
% figures are CORRECT and display the idea we want to convey in the best way
% possible." Plus, on 4.7 but binding on the whole arc: "there is no definition
% on the x axis for what we are measuring".
%
% NOT ONE PLOTTED OR PRINTED VALUE CHANGED IN THIS PASS. The seven gate
% fractions are the same seven floats, unrounded; 0.8, 0.41--0.82, 0.121,
% 0.083, 0.90--0.95, 0.620, 0.031, 480, twelve, forty, ten, eleven and seven
% are all unmoved. What changed is the y map (0.75 + f*2.10 -> 1.45 + f*3.95)
% and the x map (1.05 + (L-2)/14*11.00 -> 1.45 + (L-2)/14*11.75), both recorded
% in figs/purify.json beside the old ones. A map change is a position change,
% not a value change, and the ledger carries both maps so the audit can redo it.
% SUPERSEDED IN v11, AND THESE TWO SENTENCES CARRIED FOUR WRONG COEFFICIENTS.
% What v10 actually shipped, read off its own code and confirmed by measuring
% its render's dot centroids, was \fy = 1.45 + f*3.85 (not f*3.95) and
% \lx = 1.85 + (L-2)/14*12.00 (not 1.45 + (L-2)/14*11.75). figs/purify.json's
% value_to_position_maps.current_v10 held the correct pair all along, so the
% header and its own sibling ledger disagreed; under the header's numbers an
% auditor redoing the drawing would have put layer 16 at 13.20cm against the
% 13.85cm drawn and the gate at 4.61cm against 4.53cm. The v11 maps are
% recorded beside the code at the top of the picture, which is the only place
% they can be checked against what is drawn. No value changed in either pass.
%
%  1. THE THREE ALL-CAPS SANS ACCENT BANNERS ARE GONE, replaced by fig:mechanism's
%     own panel-head form: "(a) localise the slab", "(b) purify it against
%     context", "(c) the removal gate", \footnotesize roman ink, sentence case,
%     each left-aligned to its panel's x origin. fig:mechanism (4.11) carries
%     zero banners and sets "(a) encoded more, used no more" at 10.0pt roman
%     ink; figs/slab.tex (A.3) carries nine banners, but A.3 is an appendix
%     reference plate a reader stops at and this is read inside a running
%     argument, where the caps sans accent face belongs to the margin
%     apparatus (\marginlabel) and not to the figure. The ordinal now lives in
%     the panel letter, so the enumeration "one, two, three" that made a reader
%     hunt for a missing stage four is gone with it, and so is the
%     "four \quad what the ablation costs: \cref{fig:hook}" pointer that closed
%     it -- a cross-reference in a drawing renders as ?? in the harness and on
%     the page competed with the x-tick row it shared a baseline with, where
%     the word "four" read as a leftmost tick.
%  2. EVERY ITALIC IS GONE, AND WITH IT EVERY SENTENCE THAT WAS ARGUING RATHER
%     THAN NAMING. Deleted: the three \footnotesize italic panel subtitles;
%     the centred divider sentence "an ablation is only a test of the drug if
%     the slab it removes really holds the drug", which restates the body
%     paragraph directly above the float; the two lines inside the fail region,
%     "had any layer landed in here, the null below it / would have been
%     unfalsifiable --- a measurement of nothing", which restate the caption's
%     own "the shaded region ... is empty"; and the centred discipline line
%     "the dots and the bands are schematic and have no scale; only the four
%     numbers in stage two are measured". Italic went from 31% of the figure's
%     characters -- the highest in the five-figure set, against 0% in 4.11 --
%     to nil. All five sentences are staged in figs/_PENDING_CAPTIONS.tex under
%     \label{fig:purify} for a human to fold into the caption, which is where
%     this document keeps its argument; none is silently lost.
%  3. THE SCHEMATIC DISCLAIMER IS NOW CARRIED BY THE GLYPH AND NOT BY A
%     FOOTNOTE. The panel drew thirteen SCHEMATIC dots at the same radius and
%     the same fill as panel (c)'s seven MEASURED ones and then apologised for
%     the collision in an 8pt line. The twelve class means are now OPEN RINGS,
%     radius 1.4pt, white fill, 0.5pt rulegrey edge -- this document's mark for
%     a schematic, unscaled object -- and the seven gate fractions stay filled
%     accent discs at radius 1.4pt, which is its mark for a point estimate.
%     One look now separates them and the disclaimer line is deleted.
%  4. THE FRACTION AXIS IS DEFINED AT THE AXIS. This is the supervisor's first
%     complaint, answered mechanically. It carried bare numerals 1.0 / 0.5 / 0
%     with the quantity named only in an italic subtitle about 30mm above and
%     20mm left of the topmost tick, so a reader at the axis could not tell
%     whether 1.0 meant "all the signal removed" or "accuracy 1.0". It now
%     carries, rotated 90 degrees reading upward and immediately left of the
%     tick numerals, exactly fig:mechanism's form: line one \footnotesize roman
%     ink naming the measured quantity, line two \scriptsize quietink giving its
%     construction. The layer axis keeps the label "layer" it already had --
%     one of only two axes in 4.7-4.10 that already met 4.11's standard -- and
%     is now left-aligned to the axis's own left end rather than centred on the
%     widest tick gap.
%     MEASURED, and this is what fixed the wording. A rotated line cannot be
%     longer than the lane left of the tick column, and that lane is pinned at
%     both ends by things that may not move: below by the foot note's ascender
%     line at y = 0.52, above by the descender line of the prose at x = 0, which
%     after the 1.5mm lift of item 15 sits at y = 6.44. The axis's own midpoint
%     is 3.375, so the usable half-length is 2.855cm and a line may run 5.71cm.
%     Measured off this figure's own render, \footnotesize sets at 0.166
%     cm/character in this class and \scriptsize at 0.119, so the lane holds 34
%     \footnotesize or 47 \scriptsize characters.
%     "fraction of above-chance drug signal removed" is 43 characters, needs
%     7.14cm, and fits no layout under the 105mm height cap. What is set instead
%     is "fraction of drug signal removed" (31 characters, 5.15cm) over
%     "(before minus after)/(before minus chance)" (42 characters, 5.00cm). The
%     second line carries "above chance" in the only form that is auditable
%     against the JSON anyway, the construction itself: a reader standing at the
%     axis can now see that 1.0 means all of it, and can recompute the quantity
%     from the three accuracies in the foot note. The minus signs are spelled as
%     words and not set as $-$: in math mode CMSY10 came out at 8.3pt, a THIRD
%     type size above 7pt, which the style contract's first clause forbids.
%     SUPERSEDED IN v11, AND THIS PARAGRAPH IS THE ERROR v11 EXISTS TO CORRECT.
%     Three things in it are wrong. (i) THE LANE ARITHMETIC DISAGREES WITH ITS
%     OWN LEDGER AND WITH THE CODE: this block says lane 5.71cm at 0.166
%     cm/character with the axis midpoint at 3.375, while the inline comment at
%     the axis said 5.44cm at 0.173 with midpoint 3.24, and figs/purify.json
%     said 5.71/0.166/3.48. Remeasured on the v11 render from the PDF's own
%     text spans, \footnotesize sets at 0.161 cm/character in this class, and
%     the one figure that matters is now recorded once, beside the axis.
%     (ii) THE CONCLUSION DOES NOT FOLLOW. That a 43-character STRING does not
%     fit rules out the string, not the QUANTITY; dropping "above-chance" left
%     the \footnotesize line naming a quantity the marks are not, and that is a
%     correctness defect, not a compromise. The name is now set over two rotated
%     \footnotesize lines and the \scriptsize construction line, which S5 makes
%     optional, is deleted -- it had become the tallest object in the panel,
%     rising 6.0pt above the panel head's own top, so the smallest type in the
%     block outranked the head. The construction moved to the foot note.
%     (iii) THE $-$ NOTE IS RIGHT AND WAS APPLIED ONE CHARACTER SHORT: the v10
%     foot note still set "chance $0.083 = 1/12$", whose "=" came out of CMR10
%     at 8.3pt, a third size above 7pt, exactly the fault the note describes.
%     v11 sets "chance $0.083$, exactly $1/12$" and the render carries two
%     sizes above 7pt and no CMR10 at all.
%  5. THE DISPLAY FRACTION IS OUT OF THE PLOT. A \displaystyle fraction stacked
%     in a plot's label gutter was the most slide-like object in the five-figure
%     set. Its content is not lost: it IS the axis's second line now, which is
%     where a definition of the measured quantity belongs. The three accuracies
%     it explained (before 0.41--0.82, after 0.121 at every layer, chance 0.083)
%     move to a two-line \scriptsize foot note, and the six-line right gutter --
%     which mixed two direct labels with four notes at three different left
%     edges, 320.7, 402.9 and 428.4pt on the printed page -- is dissolved. What
%     stays beside the plot is exactly the two DIRECT labels: "all seven
%     layers: 88--95%" at the height of the dot cloud it names, and "gate, 0.8"
%     at the right end of the rule it names.
%  6. THE 0.8 GATE IS BLACK, SOLID, 0.9pt. It was accent 0.6pt dashed. This is
%     a correctness repair and not a preference: the same pre-registered class
%     of object -- a threshold a measurement is tested against -- was drawn
%     three different ways inside nine printed pages (accent dashed here,
%     accent dashed in fig:materiality, solid black in figs/family.pdf), while
%     fig:mechanism draws the chance floor solid black 0.9pt and style_v2.py's
%     chance_line() is hard-coded color='black', lw=0.9, solid, with the
%     docstring "Thin, solid, black. Never dashed, never omitted". Accent in
%     this document means the object under test; a threshold is not the object
%     under test, so a threshold may not be accent. Its label goes with it:
%     \scriptsize roman black, reading "gate, 0.8".
%  7. THE FAIL REGION IS LIGHTENED TO rulegrey!12 AND EMPTIED OF TEXT. It was
%     rulegrey!18 carrying two italic lines, it filled about 80% of the plotting
%     area, and a cold reader's eye went to it first and read it as the data
%     while the seven dots sat in a 12% band at the top. rulegrey!12 over white
%     renders at grey 243, above the 230 threshold the ink budget counts, so
%     the region now costs nothing against the 6.0% ceiling and still shows. Its
%     emptiness is the message; nothing is written inside it.
%  8. THE CONFOUND AND THE RIGHT ANGLE ARE LABELLED AT THE MARK. A cold reader
%     attached the label "the confound" to the diagonal band rather than to the
%     intersection patch, and did not see the right-angle glyph at all until
%     zooming: it was 2mm. The glyph is now 4mm on each arm and it carries the
%     direct label it never had.
%     THE LABEL IS ONE WORD, AND THAT WAS FORCED BY A RENDER. "orthogonal by
%     construction" is 3.6cm at \scriptsize, and panel (b) has no 3.6cm lane at
%     the glyph's own height that does not cross V_pure, the context band or the
%     removal arrow. Set anchor=base east at x = 9.71 it ran back to x = 6.20,
%     i.e. it sat in the gutter BETWEEN the two panels and a cold read could not
%     tell which panel owned it. Set base west on panel (b)'s x origin and
%     broken over two lines it ran to x = 10.20 and butted the V_pure label at
%     10.22, so the two read as one string, "orthogonal byV_pure" -- measured at
%     500dpi on the pass-4 render. What is drawn is the one word that has to be
%     AT the mark, set adjacent to the glyph's vertical arm with no leader.
%     "by construction" is the claim and not the name, so it goes to the
%     caption, and figs/purify.json records the move.
%     The confound's label acquires a 0.30pt leader to the patch. The dose
%     reading joins the two subspace fractions in one measurement strip on panel
%     (b)'s own x origin, so the four measured numbers of stage two read as one
%     block and not as four asides scattered round the geometry.
%  9. THE TWO PROSE BLOCKS ARE CUT TO DEFINITIONAL CONTENT ONLY, AND THEY NOW
%     OUTRANK THE LABELS. Four \scriptsize grey lines per panel became two
%     \footnotesize roman ink lines per panel: "Twelve centred class means span
%     at most eleven / directions; ten are kept." and "Drugs are confounded with
%     plate and dose by design; / no causal claim is made on the raw slab." The
%     conservatism argument ("so the slab keeps essentially the whole span,
%     noisiest included, which for a removal claim is conservative") and the
%     dose-recoverability sentence are argument, not definition, and are staged
%     for the caption. figs/slab.tex sets its explanatory prose ABOVE its labels
%     in size and this file used to invert that -- 7.97pt grey prose under
%     9.96pt grey italic narration. Prose is now the largest type in the figure
%     alongside the panel heads and the axis label, which is 4.11's ranking.
% 10. THE TWO DIVIDERS ARE GONE, the full-measure horizontal one at y = 4.72 and
%     the vertical one at x = 8.40. The horizontal rule existed to carry the
%     italic sentence beneath it (item 2) and to say "above this nothing has a
%     scale", which the open ring now says at every mark; the vertical rule was
%     furniture the panel heads make redundant. fig:mechanism draws no dividers.
% 11. ONE GRID, ONE SET OF WEIGHTS, ONE SET OF SIZES.
%     x origins, one per panel, referenced by every panel-level line:
%     (a) x = 0.00, (b) x = 8.40, (c) x = 1.60, which is the fraction axis's own
%     spine. Panel (b) carries a second, deliberate text column at x = 11.95 for
%     the four labels belonging to marks on its right. The figure's own origin,
%     x = 0.00, carries the two foot lines. MEASURED on the finished render,
%     from the PDF's own text spans: panel (a)'s six panel-level lines all start
%     at 59.38pt, panel (b)'s five at 297.49pt, its right column's four at
%     398.12pt, panel (c)'s head and its x-axis label both at 104.74pt, and the
%     two foot lines at 59.38pt. The spread inside every block is 0.00pt.
%     The upper band's two plot areas used to start 19.5pt apart for no reason;
%     they now share their head, strip and prose baselines exactly.
%     FOUR STROKE WEIGHTS, each with one meaning, against eight before:
%       0.30pt furniture -- ticks, leaders, the twelve displacement stems, the
%                           context band's outline
%       0.40pt axis rule -- the two axes of panel (c), and nothing else
%       0.50pt schematic -- the open rings, the global-mean cross, the dotted
%                           projection line, the orthogonalisation arrow: every
%                           unscaled mark in (a) and (b) is drawn at one weight
%       0.90pt threshold -- the gate rule, and nothing else
%     No weight is used exactly once. The 0.6pt arrow, the 0.35pt outlines, the
%     0.45pt cross and the 0.65pt strokes are all folded into these four.
%     TWO TYPE SIZES: \footnotesize (10.0pt) for panel heads, the axis label's
%     first line and the definitional prose; \scriptsize (8.0pt) for ticks,
%     direct labels, the scope line, the axis label's second line and the foot
%     note. Nothing else, and no italic anywhere.
%     AMENDED IN v11, SIZES UNCHANGED, ASSIGNMENT CHANGED: the fraction axis's
%     block is now BOTH of its lines at \footnotesize, because the block is one
%     wrapped quantity name and not a name over a subordinate note; the
%     \scriptsize construction line it used to carry is the second foot line.
% 12. ONE ARROWHEAD IN THE WHOLE FIGURE. A single arrowhead in this document
%     means an operation acting in a direction, and the orthogonalisation push
%     is the only operation drawn here. The displacement mu_d - mu_bar in panel
%     (a) used to carry one too, which made a static object look like an
%     action; it is now a plain 0.30pt stem like its eleven siblings, named by
%     a leader instead. Nothing else in the figure has a head.
% 13. COLOUR, one code, the same one across the five figures of the arc.
%     accent = the drug-side object under test and its direct labels: V_drug,
%     V_pure, the confound, the orthogonalisation arrow, the seven gate dots,
%     "all seven layers: 88--95%".
%     quietink = the comparator and control side: the context subspace's two
%     labels, the schematic labels in (a), the axis label's construction line,
%     the foot note.
%     AMENDED IN v11: the construction line has left the axis for the foot, and
%     three things join the quietink list -- panel (c)'s which-slab scope line,
%     the name of the shaded region ("an untestable null", a counterfactual and
%     so a comparator), and panel (b)'s measurement strip, which in v10 was ink
%     while panel (a)'s scope line was quietink at the same size on the same
%     baseline in the adjacent panel, with nothing on the page decoding the
%     difference. ink is now heads, definitional prose, the axis labels, the
%     gate label and the one geometric fact true by construction, and nothing
%     else.
%     ink/black = thresholds, axis rules, axis labels, panel heads, definitional
%     prose, and the one geometric fact true by construction.
%     rulegrey = furniture only: ticks, leaders, stems, the ring edges, the
%     context band's outline, the fail region's fill.
%     The one visible consequence: the twelve class means are no longer ink
%     dots. They are schematic and they are drawn in the schematic register.
% 14. WHAT WAS PROTECTED AND IS UNTOUCHED. The three-stage spine and its order.
%     The seven gate dots, byte-exact, and the identity behind them. The full
%     0-to-1 gate range, decision 12, which a reviewer asked to crop to
%     0.75--1.00 and which was overruled and stays overruled. The layer axis
%     and its map's SHAPE, 1.45 + (L-2)/14*W. The orthogonalisation geometry:
%     raw band at 38 degrees, context band horizontal, V_pure vertical, the
%     dotted perpendicular, the named confound, the removal arrow, the right
%     angle. The refusal to draw a reference rule at 0.90 for the dose
%     cross-decode, for the reason in (e). The four measured numbers of stage
%     two and their layer-9 qualifier.
%
% 15. PANELS (a) AND (b) ARE LIFTED 1.5mm AS ONE BLOCK, inside a single
%     \begin{scope}[yshift=1.5mm]. Two things come of it. The ink budget,
%     measured at 200dpi inside the content bbox, falls from 5.98% to 5.90%
%     against a 6.00% ceiling, purely by growing the denominator -- drawn height
%     102.4 -> 103.8mm, still under the 105mm cap. And the gap between the last
%     prose line of the schematic band and panel (c)'s head goes from 2.8mm to
%     4.3mm, which is the separation the deleted full-measure divider used to
%     carry, now made of white instead of a rule. Panel (c) does not move: its
%     rotated axis label's lane is pinned between the foot note's ascender line
%     and the prose's descender line, so lifting the prose only lengthens it.
%
% WHAT SEVEN RENDER-AND-READ PASSES CHANGED IN v10, each with the measurement
% that forced it. Every one was found by rendering the figure and LOOKING at
% the PNG, not by reading the source.
%   i.   PASS 1. The rotated axis label was one \footnotesize line centred on
%        the axis and it overhung the axis at both ends, running up into panel
%        (a)'s prose and down into the tick lane. The plot was at once too tall
%        (3.95cm of axis, so the empty fail region dominated the page) and too
%        short for the label. Diagnosis: the LANE, not the axis, is the
%        constraint. Panel (c) was reproportioned so the axis midpoint sits
%        exactly at the lane's centre, which is the only arrangement in which
%        both rotated lines can be centred on the axis and still clear.
%   ii.  PASS 2, Overfull \hbox 6.00356pt, and it survived a complete rewrite of
%        the body, which is what identified it. The INK bbox measured 170.4mm,
%        inside the 174.0mm fullwidth measure, so the render looked clean; the
%        NODE box did not. "the global mean", set anchor=base east at x = 1.80,
%        is 2.00cm wide, so its left edge landed at tikz x = -0.32 and widened
%        the picture to 17.62cm = 501.3pt against \linewidth 495.08pt -- the
%        latter measured by compiling \the\linewidth inside a real fullwidth
%        environment, which also settles an old dispute in this header: 174.0mm
%        is right and the note calling preamble.tex's arithmetic stale is not.
%        The label moved to anchor=base west on the panel's own x origin.
%   iii. PASS 3. Ink measured 7.94% against a 6.00% ceiling. Region-by-region
%        measurement put 1.63 of those points in ONE object, the V_drug band at
%        fill opacity 0.15: accent!15 over white renders at grey 227, under the
%        230 threshold the budget counts, over 3.06cm^2. Both accent bands went
%        to 0.12, which renders at grey 233 and therefore costs nothing at all,
%        and the fail region went from rulegrey!18 (grey 237) to rulegrey!12
%        (grey 243). Ink 5.87%, and no measured mark was touched to get there.
%   iv.  PASS 4, at 500dpi. Three collisions invisible at 200dpi: "orthogonal
%        by" butting the V_pure label (item 8); the raw slab's leader stopping
%        0.19cm short of the band and reading as a stray hairline, which is the
%        exact fault the pass-9 header records against the orthogonalise leader;
%        and the fraction axis's tick numerals sitting 1.6pt from the rotated
%        label, so the two read as one column of type.
%   v.   PASS 5. The raw slab's label then sat 0.05cm under the measurement
%        strip's first line and the two read as one three-line paragraph. The
%        label moved above the strip and its leader now lands INSIDE the band
%        rather than beside its tip.
%   vi.  PASS 6. The global-mean cross was 2mm arms at 0.5pt with twelve stems
%        converging on it, and at 500dpi only its vertical arm was visible
%        through them; it is now 2.8mm. The two schematic outlines in panel (b)
%        moved 0.30 -> 0.50pt so that 0.30pt means furniture and nothing else.
%        V_pure acquired its gloss AT THE MARK, which is where this document
%        requires a symbol that is itself a drawn object to be glossed.
%   vii. PASS 7. Ink 5.98% against a 6.00% ceiling is not a pass, it is a
%        rounding error away from a fail. Item 15.
%
% HEIGHT AND MEASURE, v10, all MEASURED off the finished render at 200dpi and
% not computed from the source: bounding box \path (0,-0.10) rectangle
% (17.30,10.35); content bbox 172.1 x 103.8mm, inside the 170-174mm band a
% fullwidth figure must fill, under the 105mm cap this pass was given and under
% the corpus ceiling of 117mm (slab.tex). The data region, x = 1.60 to 14.10,
% is 12.50cm = 72.6% of the drawn width, over the 60% floor. Ink 5.90% of the
% content bbox at pixels below grey 230, against a 6.00% ceiling. Text density
% 4.26 characters per 100mm^2 against 5.00. Type: two sizes above 7pt, 9.96 and
% 7.97, against fig:mechanism's 10.0 and 9.0; one face, TeXGyrePagella; three
% italic characters in 864 and all three are the math-italic V of a symbol;
% zero sans, zero bold, zero all-caps runs of two or more words. Four stroke
% weights, 0.30 / 0.40 / 0.50 / 0.90pt, none used exactly once. Exactly one
% black solid 0.90pt rule and it is the gate -- byte-for-byte the register
% figs/family.pdf gives the other pre-registered threshold of this chapter, a
% 0.9pt rule at colour (0,0,0), checked by reading family.pdf's drawing list.
% SUPERSEDED IN v11 as to the MEASUREMENTS ONLY, all of which were remeasured
% and several of which were wrong or unreproducible: see "HEIGHT AND MEASURE,
% v11" at the foot of the v11 block below. The claims about type, faces,
% weights and the single black rule survive, with one correction (the CMR10
% "=" at 8.3pt, item (iii) of the superseding note under v10 item 4).
%
% ===========================================================================
% v11 -- 2026-08-19. ADVERSARIAL REPAIR PASS. Three refuters read the v10
% render and returned 33 findings. What follows is what was applied, what was
% overruled, and why. Nothing below changes a plotted or printed VALUE: the
% seven gate fractions are the same seven floats, unrounded, and 0.8,
% 0.41--0.82, 0.121, 0.083, 0.90--0.95, 0.620, 0.031, 480, twelve, forty, ten,
% eleven and seven are all unmoved. What moved is positions, wording that
% NAMES, and the audit record.
%
% APPLIED, in order of severity.
%
%  1. THE FRACTION AXIS NAMED THE WRONG QUANTITY. This is the one correctness
%     defect in the drawing and it was in the very line v10 was built to add.
%     The plotted series is layers.<L>.frac_signal_removed, which is the
%     fraction of ABOVE-CHANCE signal removed; v10's axis read "fraction of
%     drug signal removed". The unqualified quantity, (before - after)/before,
%     computes from the same file to 0.7041 at layer 2 and 0.8528 at layer 9 --
%     4.2 to 18.0 points BELOW the marks drawn -- and read at the drawn name the
%     top mark, 0.949, implies a residual accuracy of 0.041, under the chance
%     0.083 the figure prints two lines away. The pre-registration itself
%     (04b-representation.tex:218) says "the fraction of above-chance signal
%     removed must exceed $0.8$". The axis now reads, over two rotated
%     \footnotesize lines, "fraction of above-chance / drug signal removed".
%     v10's stated reason for the omission is superseded in place above.
%
%  2. THE RIGHT ANGLE AND THE REMOVAL ARROW WERE ONE PRIMITIVE. Measured on the
%     v10 render: the glyph's horizontal arm at y = 9.05 and the arrow's shaft
%     at y = 9.0349, 0.0151cm = 0.13mm apart, on opposite sides of an opaque
%     0.18cm bar, so a static assertion (perpendicularity) and a directed
%     operation (the projection) read as one horizontal line running from the
%     word "orthogonal", through V_pure, out to the dotted drop line. Arms
%     0.40 -> 0.28cm, arm to y = 8.93, 1.05mm of white under the arrow. The
%     quadrant is and always was UPPER-left; v10 item 8's "lower-left" is
%     superseded in place above.
%
%  3. PANEL (b)'s HEAD AND THE V_pure GLOSS WERE 0.36pt APART. 0.13mm, against
%     2.8mm under panel (a)'s head: two panels of equal rank with head
%     clearances differing by a factor of twenty-two. Fixed by a restack whose
%     every step is set by measured ink boxes -- see the note at the V_pure bar.
%     Head-to-gloss is now 6.4pt (2.3mm) and the gloss sits 7.6pt clear of the
%     right column's first baseline, where the first v11 build had put it level
%     with that baseline and the two read as one line.
%
%  4. WHICH SLAB THE GATE IS MEASURED ON IS BACK ON THE DRAWING. v10 cut it for
%     text density and left panel (c) with no slab named anywhere in its head,
%     axis, ticks, labels or foot, sitting directly under panel (b)'s closing
%     line "so causal claims use the purified slab, not the raw". The inference
%     that the gate was run on V_pure was then unavoidable, and it is not
%     supported: workspace_probe.json stores one ablated series and no purified
%     counterpart. The scope line sits under the panel head, in axis-label rank.
%     NOTE ON THE LEDGER'S OLD REASONING, corrected in figs/purify.json: the
%     identity in (d) above fixes only the BEFORE term as the raw probe
%     accuracy; it says nothing about which subspace was projected out. The
%     honest statement, and the one now drawn and recorded, is that the file
%     does not record it. (For scale: probe_arm_residual.json, the one instrument
%     that stores both, has layers.12 frac_signal_removed 0.3634 against
%     frac_signal_removed_pure 0.3241, i.e. the purified-slab gate is the LOWER
%     of the two there -- which is why this may not be assumed either way.)
%
%  5. THE FOOT NOTE HAD NO NOUN. v10 printed "Before 0.41--0.82, after 0.121 at
%     every layer, chance 0.083" and the words "accuracy", "probe" and "1/12"
%     appeared nowhere in the figure, so nothing on the page said what
%     instrument panel (c) reads. It now names the quantity by its registry
%     name, "probe accuracy", gives chance as "$0.083$, exactly $1/12$" (the
%     stored float IS 1/12 in IEEE double), and carries the construction the
%     deleted rotated line used to hold.
%
%  6. THE SHADED FAIL REGION IS SMALLER, LIGHTER AND NAMED. It was 21.6% of the
%     content box, the largest object on the page, unnamed, and v10 had nearly
%     doubled its absolute area as a side effect of stretching the axis to house
%     an over-long label. The full 0-to-1 range is untouched -- decision 12
%     stands and is the first thing this pass was told to protect -- and so is
%     the region's exact 80% share of the axis. What changed is the SPAN, which
%     is a free choice: 3.85 -> 3.00cm, so the region falls to 16.8% of the
%     content box. It is drawn at rulegrey!8 (grey 246) and not rulegrey!12
%     (grey 243), which was within one level of the panelbg tint that means
%     "a schematic subspace" 3cm above it -- one grey, two meanings. And it
%     carries a NAME, "an untestable null", set OUTSIDE its right edge in
%     quietink: nothing is written inside it (the spec forbids it and the
%     emptiness is the message), but the largest mark on the page may not be the
%     one mark a reader cannot name, and without the name the third clause of
%     this figure's acceptance test is unanswerable from the drawing.
%
%  7. "layer" READ AS "layer 2". v10 set it \footnotesize at x = 1.60 with the
%     layer-2 numeral 0.63mm directly above it, one size smaller -- the exact
%     fault the pass-9 note above records against the old placement under the
%     layer-9 numeral, reintroduced at layer 2. The repair is VERTICAL, not
%     horizontal, so S5's left-alignment survives: the fraction origin rose
%     1.45 -> 1.75, which buys 2.8mm of white and a 0.53cm baseline separation
%     against the 0.40cm leading this figure uses inside one sentence.
%
%  8. THE GLOBAL-MEAN CROSS WAS STILL INVISIBLE AND ITS LEADER READ AS A
%     THIRTEENTH STEM. figs/purify.json certified a pass-6 legibility repair
%     (arms 2.0 -> 2.8mm) that the render never showed. Length was never the
%     fault: three of the twelve stems leave the vertex within 10 degrees of
%     horizontal, so the horizontal arm is buried at every radius. A 0.20cm disc
%     of the band's own tint now separates them. See the note at the mark.
%
%  9. FOUR LEADERS AND ONE PATCH. "the confound" pointed 0.03cm BELOW its patch,
%     into the context band -- the misreading the pass-9 rebuild of that mark
%     exists to kill. "orthogonalise" died in white 0.105cm short of the arrow's
%     tail while carrying only a 0.05cm standoff at the label end, the inversion
%     of the corpus rule this header states twice. "the global mean" grew out of
%     the final "n" of its own label at 0.03cm. "all seven layers: 88--95%" sat
%     3.6mm from the layer-16 dot with nothing between, over the corpus's 3mm
%     leader threshold. All four are fixed at the mark end. And "$V_{drug}$,
%     raw" rose to y = 8.352 and sat on the context band's 8.25 top edge; it
%     drops to baseline 7.98.
%
% 10. THE TINT-MULTIPLY PARALLELOGRAM WAS STILL BEING DRAWN. Pass 9 deleted the
%     CLAIM that the darkened crossing is the confound but not the overprint
%     that produced one: the raw band was drawn after the opaque context
%     rectangle, so an unlabelled darker parallelogram appeared inside the band
%     beside the patch that actually carries the name. The two \fill blocks are
%     simply reordered. No coordinate and no colour changed.
%
% 11. COLOUR. Panel (b)'s measured-number strip was ink and panel (a)'s scope
%     line quietink, at the same size on the same baseline in adjacent panels,
%     with nothing on the page decoding the difference. Both are scope; both are
%     quietink. ink is now heads, definitional prose, the axis labels, the gate
%     label and the one geometric fact true by construction, and nothing else.
%
% 12. ONE OBJECT, ONE NAME. "subspace fraction" in panel (b) becomes
%     fig:mechanism's own name for the quantity, "descriptive subspace
%     fraction".
%
% 13. V_drug's BAND IN PANEL (a) IS SYMMETRIC ABOUT ITS CLOUD. It overhung the
%     data 0.13cm left against 0.30cm right, so the leftmost ring read as
%     sitting on the band's edge -- the bounded-region reading the deleted
%     ellipse was removed to prevent. Left edge 0.55 -> 0.38; both 0.30cm.
%
% OVERRULED, with the reason, because a refuter who is wrong about this
% document's conventions must be answered and not quietly followed.
%
%  A. "BRING PANEL (c)'s TWO DIRECT LABELS INSIDE THE AXIS." Overruled. It would
%     take the drawn width to 169.7mm, under the 170mm floor the SAME clause
%     sets for a fullwidth figure, and it would put "all seven layers: 88--95%"
%     4mm from the marks it names with no room for the leader the corpus
%     requires. The two clauses cannot both hold in this geometry. What is drawn
%     instead: the labels stay at their marks, "all seven layers" acquires its
%     0.30pt accent leader, and the figure's right edge is set by a direct label
%     0.48cm past the widest panel element (panel (b)'s prose at 16.78cm). That
%     is the trade recorded, not an oversight.
%
%  B. "DELETE THE TWELVE DISPLACEMENT STEMS." Overruled. They are the picture of
%     the word "centred" in the panel's own prose -- each mean is drawn as a
%     displacement from the global mean -- and they are what makes the count
%     twelve a picture rather than an assertion. Their whole ink is 0.77% of the
%     figure's, so deleting them buys almost nothing, and the three faults the
%     refuter traced to them (an invisible cross, a leader that reads as a
%     thirteenth spoke) are fixed by item 8 without losing the picture.
%
%  C. "MAKE V_pure AND THE RAW BAND THE SAME WIDTH, so no width carries a
%     value." Overruled. The drawn width ratio is not read as the measured one
%     -- nothing in panels (a) and (b) has a scale, the open-ring glyph says so
%     at every mark and the caption says the rest -- while "visibly thinner" is
%     the one cue that ties the drawing to the 0.620 -> 0.031 printed beneath
%     it. Equalising the widths would re-open the pass-9 misreading the header
%     records at length: two bands of equal width differing only in angle say
%     "same slab, rotated, nothing lost".
%
%  D. "REDRAW V_drug IN PANEL (a) AS A THIN STRIP AND LET THE OUTLYING RINGS
%     FALL OUTSIDE IT, so the two panels' glyphs rhyme." Overruled on the
%     file's own recorded ground, which is correct: the twelve centred class
%     means span V_drug BY CONSTRUCTION, so a class mean drawn outside it would
%     be false. Decision 10(a) already settled this and nothing has changed.
%
%  E. "WIDEN PANEL (a) AND MOVE PANEL (b)'s ORIGIN LEFT to close the 3.5cm void
%     between the two geometries." Overruled as disproportionate. The void is
%     confined to the geometry band; at the prose rows the two panels fill the
%     measure with a 0.93cm gutter, and moving panel (b)'s origin would move
%     fifteen coordinates and break the one-x-origin-per-panel grid the v10
%     render was measured against for no measurable gain.
%
%  F. "'the confound' IS 17pt OFF THE RIGHT COLUMN's LEFT EDGE." Overruled: it
%     is a DIRECT LABEL AT ITS MARK, which the corpus requires to sit at the
%     mark, and direct labels are exempt from the panel-origin rule -- as are
%     "$V_{drug}$, raw", "orthogonal" and the V_pure gloss in the same panel.
%
%  G. "PRINT THE SEVEN-LAYER RANGE OF THE RAW SUBSPACE FRACTION ON THE DRAWING."
%     The finding is RIGHT and is acted on, but not there: layer 9 is indeed the
%     MAXIMUM of variance_share.drug (the seven-layer range is 0.2767 at layer
%     16 to 0.6201 here), and the drawing does not say so. There is no lane for
%     it -- the strip line already runs to 16.38cm -- so the disclosure goes to
%     figs/purify.json and to the staged caption, which is where this document
%     keeps what will not fit. The layer-9 qualifier stays on the drawing, so
%     the pair cannot be read as a seven-layer summary.
%
%  H. NOT DONE, AND OWED. figs/_PENDING_CAPTIONS.tex still carries no
%     \label{fig:purify} block, so the eight sentences v10 cut and the whole
%     replacement caption live only in the comment at the foot of this file.
%     Two refuters are right that this is a contract breach. It is not repaired
%     here because this pass's writable set is figs/purify.tex and
%     figs/purify.json only; the block is staged verbatim below for whoever
%     holds the pen on that file, and figs/purify.json records the debt under
%     pending_captions_block_owed.
%
% WHAT SIX RENDER-AND-READ PASSES CHANGED IN v11, each with the measurement.
%   i.   PASS 1. The two rotated \footnotesize lines were set in the wrong
%        order. Rotated 90 degrees the line-advance direction is -x, so the
%        first line is the LEFTMOST; set the other way round the block printed
%        "drug signal removed / fraction of above-chance".
%   ii.  PASS 2. Dropping the V_pure gloss for head clearance landed it on the
%        right column's own first baseline with 0.15cm of horizontal air, and
%        the two read as one line. Hence the restack of item 3.
%   iii. PASS 3. Text density 5.01 characters per 100mm^2 against a 5.00
%        ceiling -- a rounding error away from a fail, which is not a pass.
%        Two wordings were tightened (the scope line, the interval line) with
%        no content lost; 4.96 now.
%   iv.  PASS 3. The render carried THREE type sizes above 7pt: 9.96, 7.97 and
%        an 8.3pt CMR10 "=", two characters, from "$0.083 = 1/12$". Reset as
%        "$0.083$, exactly $1/12$"; two sizes now and no CMR10 in the file.
%   v.   PASS 4. Panel (c)'s head and its new scope line had 3px of white
%        between their ink at 200dpi -- 0.38mm. The fraction span came 3.15 ->
%        3.00cm and the head rose 0.04cm; 15px, 1.9mm, now.
%   vi.  PASS 5. Read at 6x, the whole of panel (b) and the whole of panel (a).
%        Every leader in item 9 was found this way and not in the source.
%
% HEIGHT AND MEASURE, v11, MEASURED off the finished render and not computed.
%   bounding box \path (0,-0.10) rectangle (17.30,10.35), unchanged.
%   content bbox 172.09 x 103.76mm, inside the 170--174mm band and under the
%     105mm cap and the corpus ceiling of 117mm (slab.tex).
%   data region x = 1.60 to 14.10 = 12.50cm = 72.6% of the drawn width, over
%     the 60% floor.
%   text density 4.96 characters per 100mm^2 (spaces excluded, which is the
%     count that reproduces the contract's own reading of fig:mechanism, 2.49
%     against its stated 2.5) against a 5.00 ceiling.
%   ink. THE CONTRACT'S 6.0% CEILING IS PIPELINE-DEPENDENT AND THIS FILE NOW
%     SAYS SO RATHER THAN QUOTING A NUMBER THAT CANNOT BE REPRODUCED. Rasterised
%     with PyMuPDF and thresholded at grey 230 inside the content bbox, this
%     figure reads 6.54% at 200dpi, 5.56% at 400dpi and 5.21% at 600dpi; the
%     SAME pipeline on fig:mechanism as it prints on page 51 of main.pdf reads
%     5.99 / 5.26 / 5.04. The contract's own measurement of fig:mechanism is
%     5.3, which its pipeline reproduces at 400dpi and not at 200. At 400dpi,
%     the resolution at which the reference figure reproduces, this figure is
%     5.56% against the 6.00% ceiling. At any of the three resolutions it sits
%     between 1.03x and 1.09x fig:mechanism, where the ceiling is 1.13x it.
%     v10's "5.90%" was a 200dpi number quoted against a 400dpi ceiling.
%   type: two sizes above 7pt, 9.96 and 7.97, plus 5.98 for subscripts, which
%     the contract exempts; largest size is only ever a panel head, an axis
%     label or definitional prose. One face, TeXGyrePagella, with 58 characters
%     of URWPalladioL-Roma inside math \textrm and 5 of URWPalladioL-Ital, all
%     five the math-italic V of a symbol -- 0.5% italic against a 5% ceiling.
%     Zero sans, zero bold, zero all-caps runs of two or more words.
%   strokes, read out of the PDF's own drawing list: four weights,
%     0.30 / 0.40 / 0.50 / 0.90pt. 0.30 is furniture (27 rulegrey) plus the one
%     accent leader the corpus requires; 0.40 is the two axis rules and nothing
%     else; 0.50 is every unscaled mark; 0.90 appears exactly once and it is the
%     gate, in black, which is the one place this document allows a singleton.
%   the harness: ok, pages=1, boxes=0.
% ===========================================================================
%
% ===========================================================================
% v-BLOCK, 2026-08-19: WHOLE-SLATE PASS (fig:comparators, fig:licensing,
% fig:purify, fig:hook, fig:materiality reconciled against one another and
% against fig:mechanism, which is the document's reference standard and was
% NOT rebuilt). Nothing above this line is deleted; where an earlier claim is
% now false it is superseded EXPLICITLY below.
%
% WHY. The five figures were rebuilt independently and drifted. This pass owns
% only the differences BETWEEN them. No value changed in any of the five; the
% sibling .json ledgers carry the value-to-position maps, old and new.
%
% THE SLATE RULES THAT NOW HOLD ACROSS ALL FIVE (each stated once here, in
% every file, so any future single-figure edit can see what it would break):
%
%  R1 TICK NUMERALS are quietink, \scriptsize, inner sep 0, anchored north on
%     the tick and set 0.16cm below the axis rule, over a 0.10cm rulegrey
%     0.30pt tick. Before: comparators/licensing/purify quietink, hook and
%     materiality ink; ticks 0.10 / 0.12 / 0.14cm; three different gaps.
%  R2 A THRESHOLD'S OWN LABEL IS BLACK, matching the black solid 0.90pt rule
%     it names (S6). Before: comparators black, purify's "gate, 0.8" and
%     materiality's two margin labels ink.
%  R3 ONE NAME PER OBJECT (S19). The pre-registered 10% rule is "10%
%     materiality margin" in fig:comparators row (c), fig:materiality (a) and
%     (b), and in figs/family.py, which is where the name comes from.
%     fig:comparators' "pre-registered margin" is superseded; "pre-registered"
%     is now the caption's word, not the drawing's. Zero, where it is the null
%     a bar is read against, is "zero: no identity effect" in BOTH
%     fig:comparators row (c) and fig:materiality (b) -- one quantity, one
%     wording. V_pure is spelt with \textrm everywhere, as figs/slab.tex does.
%  R4 COLOUR KEY, one key for five figures (S8). accent = the drug-side object
%     under test. quietink = every comparator, null, control, replication or
%     discounted arm, AND its label, AND its leader. ink/black = thresholds,
%     axis rules, axis labels, panel heads, definitional prose. rulegrey =
%     furniture only. Within quietink the GLYPH separates two kinds: an OPEN
%     disc is a construction built to carry no signal (a null, a control, a
%     raw/before series); a FILLED disc is a real measurement that is not the
%     headline. fig:hook (d) drew both its null and its cell-line control in
%     ink and is corrected.
%  R5 PROJECTION AND ZERO RULES are quietink 0.50pt densely dotted (S7).
%     fig:hook's two panel-(c) projection lines were rulegrey 0.30pt densely
%     dashed and are corrected.
%  R6 DEFINITIONAL PROSE is \footnotesize ink and OUTRANKS every label (S4);
%     at most one block per figure, on the panel's own x origin. fig:hook's
%     one prose line was \scriptsize and is raised. fig:purify's panel-(b)
%     block carried an argument clause ("so causal claims use the purified
%     slab, not the raw") that its caption already carries, and is now
%     definitional only.
%  R7 ONE INTERVAL DECLARATION per figure, \scriptsize quietink, at the foot,
%     on the figure's x origin, opening with the word "Intervals:" and naming
%     EVERY panel including the ones with no measured mark (S17).
%     fig:comparators already did this and is the model; the other four now
%     match its form.
%  R8 THE FIVE-FIELD PROMPT STRIP is ONE object drawn in two figures, and the
%     two drawings are now congruent: cell boundaries 0 / 1.12 / 1.90 / 2.64 /
%     4.23 / 6.90 cm and cell height 0.42cm in both fig:licensing (a) and
%     fig:hook (a) and (b), labels centred in their cells, the drug cell
%     accent-bordered at 0.40pt on accent!12 and its label NOT bold. Before:
%     licensing 6.90cm wide with 0.62cm cells and a bold label, hook 7.44cm
%     wide with 0.42cm cells and a plain label -- the same five fields at two
%     sizes, four pages apart, with each caption pointing at the other.
%  R9 MEASURE. Every fullwidth figure draws 172.0-174.0mm; the one textwidth
%     figure draws 140.3mm. Ink at 200dpi inside the content bbox, counted as
%     greyscale luminance below 230 (the contract's own method), is now
%     5.63-5.99% on all five, under the 6.0% body ceiling, where before this
%     pass it ran 5.36-6.54%.
%
% WHAT THIS PASS DELIBERATELY DID NOT UNIFY, so that a later reader does not
% think it was missed:
%  * PANEL-HEAD CASE. fig:comparators sets "(a) Is it there?" and the other
%    four set "(a) where the drug name enters". comparators' heads are
%    interrogative SENTENCES and sentence case is correct for a sentence; the
%    v4 header argued this and the argument stands. fig:mechanism's lowercase
%    heads are noun phrases, not questions. One capital letter, four times.
%  * PANEL-HEAD POSITION. fig:comparators puts its heads in a right-aligned
%    left stub, the other four above the panel at its x origin. Moving them
%    above the rules costs about 20mm of height against a 118mm cap that the
%    figure already meets at 117.60mm. Not available.
%  * TWO-COLUMN GUTTER. purify and hook split at x = 8.40cm, licensing at
%    7.75cm. licensing cannot move right: "(d) what a negative would buy" is
%    4.85cm wide and its origin is already at 12.40 of a 17.30cm measure, with
%    1.5pt of slack. Moving purify and hook LEFT to 7.75 would be arbitrary.
% ===========================================================================
%
% CHANGED IN THIS FILE. (i) The line "Gate measured on the subspace projected
% out; no separate gate is stored for V_pure" is DELETED from the drawing. It
% was the slate's one scope line not sitting under an axis label (S5), and its
% content is already in the caption verbatim as "Two limits. The gate is
% measured on the subspace that is projected out, and this file stores no
% separate gate for V_pure" -- an S18 duplication. Panel (c)'s head drops from
% y 5.50 to 5.20 to take up the space. (ii) Panel (b)'s prose becomes
% definitional: "so causal claims use the purified slab, not the raw" is an
% argument and is the caption's ("Panel (b) is why the raw subspace may not be
% used"); what is drawn now DEFINES the object the panel draws -- "V_pure is
% the part of the raw slab that is orthogonal to the context subspace."
% (iii) "gate, 0.8" is set black, matching the black solid 0.90pt rule it
% names (R2); it was ink, and this was the only threshold on the slate whose
% label and rule were in different colours. (iv) Tick geometry: the y ticks go
% 0.12 -> 0.10cm and the x ticks 0.14 -> 0.10cm (they differed from each other
% INSIDE this panel), numerals inner sep 0 at a 0.06cm gap (R1). (v) The foot
% block takes the slate's declaration form (R7).
% SUPERSEDED: pass 10's claim that the figure clears the ink ceiling at 5.90%
% was measured on the red channel of an RGB rasterisation, which counts the
% accent!12 bands as ink; measured as greyscale luminance, which is what S15
% specifies, pass 10 was at 6.541%. The deletions above bring it to 5.992%.
% RE-MEASURED: 172.09 x 103.76mm (cap 105); ink 5.992%; 910 characters, 5.10
% per 100mm^2; two type sizes above 7pt, 9.96 and 7.97, plus 5.98/7.57pt
% subscripts inside V_drug and V_pure.
% One page, zero Overfull, zero Underfull.
% ===========================================================================
\begin{tikzpicture}[x=1cm, y=1cm, font=\small,
                    every node/.style={inner sep=1.6pt}]
  \tikzset{
    % --- TWO SIZES, ONE FACE, NO ITALIC. -----------------------------------
    % \footnotesize (10.0pt): panel heads, BOTH lines of the fraction axis's
    % label (v11 -- the block is one wrapped quantity name, not a name over a
    % subordinate note), the x-axis label, and the definitional prose.
    % \scriptsize (8.0pt): everything else -- ticks, direct labels, the two
    % scope lines, the region's name and the two foot lines. Nothing in this
    % picture is set in sans, in small caps, in bold or in italic, and nothing
    % is set in math mode that brings a third size with it: "$0.083 = 1/12$"
    % put an 8.3pt CMR10 "=" in the v10 render and is now "$0.083$, exactly
    % $1/12$".
    head/.style={text=ink,      font=\footnotesize, anchor=base west},
    def/.style ={text=ink,      font=\footnotesize, anchor=base west},
    lab/.style ={text=ink,      font=\scriptsize,   anchor=base west},
    qlab/.style={text=quietink, font=\scriptsize,   anchor=base west},
    alab/.style={text=accent,   font=\scriptsize,   anchor=base west},
    % --- FOUR STROKE WEIGHTS, each with exactly one meaning. ---------------
    furn/.style ={draw=rulegrey, line width=0.3pt},   % ticks, leaders, stems
    axisr/.style={draw=ink,      line width=0.4pt},   % an axis rule
    schem/.style={draw=rulegrey, line width=0.5pt},   % an unscaled mark
    gate/.style ={draw=black,    line width=0.9pt}}   % a threshold

  % ---- value-to-position maps for panel (c). Both bodies are braced groups.
  % THESE TWO LINES ARE THE AUDIT RECORD OF THE MAPS AND THEY ARE THE CODE, so
  % read them and not the prose in the v10 header, which quoted four wrong
  % coefficients and is superseded item by item in the v11 block below.
  %   v9   \lx = 1.05 + (L-2)/14*11.00   \fy = 0.75 + f*2.10
  %   v10  \lx = 1.85 + (L-2)/14*12.00   \fy = 1.45 + f*3.85
  %   v11  \lx = 1.85 + (L-2)/14*12.00   \fy = 1.75 + f*3.00   (x map unchanged)
  % The fraction origin rose 1.45 -> 1.75 to put 2.8mm of white between the
  % layer numerals and the word "layer" (v11 item 3), and the span fell
  % 3.85 -> 3.00 to take 0.68cm out of the empty fail region and give it to the
  % gap between the schematic band and panel (c) (v11 items 5 and 6). Layer 2
  % is inset 0.25 from the spine so its dot does not straddle the axis rule,
  % which is fig:mechanism's own treatment of the same first layer.
  % No map change is a value change: the ledger records all three maps.
  \def\lx#1{{1.85+((#1)-2)/14*12.00}}  % layer  2 -> 1.85,  16 -> 13.85
  \def\fy#1{{1.75+(#1)*3.00}}          % 0 -> 1.75, 0.8 -> 4.15,  1.0 -> 4.75

  % one gate mark: a filled accent disc at 1.4pt, this document's point
  % estimate. The 2.1pt white halo of earlier passes is dropped -- nothing
  % underlies these dots now, and a second radius would blunt the ring/disc
  % distinction that replaced the schematic disclaimer line.
  \newcommand{\gdot}[2]{\fill[accent] (\lx{#1},\fy{#2}) circle (1.4pt);}

  % one schematic mark: an open ring, white fill, 0.5pt rulegrey edge.
  \newcommand{\sring}[2]{\filldraw[fill=white, schem] (#1,#2) circle (1.4pt);}

% ###########################################################################
% # PANEL (a) -- localise.  x origin 0.00.  Nothing here has a scale.
% # Panels (a) and (b) are lifted 1.5mm as one block -- see the v10 header.
% ###########################################################################
\begin{scope}[yshift=1.5mm]
  \node[head] at (0.00,9.90) {(a) localise the slab};

  % V_drug: an accent band through the global mean, running past the data on
  % both sides because a subspace has no boundary. Opacity 0.15, down from
  % 0.20: at 0.20 over 3.40 x 0.90cm this was still the heaviest area in the
  % figure and it is not the figure's subject.
  % v11: the left edge comes 0.55 -> 0.38. The twelve offsets run -1.57 to
  % +1.40, so the cloud sits at 0.68 to 3.65 and the band overhung it by 0.13cm
  % on the left against 0.30cm on the right; at 0.13cm the leftmost ring read as
  % sitting ON the band's edge, which reintroduces the bounded-region reading
  % the deleted ellipse was removed to prevent. Both overhangs are 0.30cm now.
  \fill[accent, fill opacity=0.12] (0.38,8.05) rectangle (3.95,8.95);
  \node[alab, anchor=west] at (4.08,8.50) {$V_{\textrm{drug}}$};

  % the twelve per-drug mean activation vectors. Positions are literal, so no
  % build can randomise them and no reader can be told a different story by a
  % different compile. Twelve 0.30pt stems make the count a picture; the rings
  % are OPEN, which is how this document draws a schematic, unscaled mark.
  \foreach \dx/\dy in {-1.57/-0.085, -1.25/0.238, -1.02/-0.289, -0.62/0.102,
                       -0.36/-0.374, -0.07/0.306, 0.21/-0.153, 0.55/0.357,
                       0.74/-0.306, 1.12/0.119, 1.40/-0.238, 1.21/0.298}{
    \draw[furn] (2.25,8.50) -- ({2.25+\dx},{8.50+\dy});}
  \foreach \dx/\dy in {-1.57/-0.085, -1.25/0.238, -1.02/-0.289, -0.62/0.102,
                       -0.36/-0.374, -0.07/0.306, 0.21/-0.153, 0.55/0.357,
                       0.74/-0.306, 1.12/0.119, 1.40/-0.238, 1.21/0.298}{
    \sring{2.25+\dx}{8.50+\dy}}

  \draw[furn] (1.20,9.06) -- (1.05,8.79);
  \node[qlab] at (0.00,9.32) {the twelve per-drug means};

  % THE GLOBAL MEAN: a cross and not a thirteenth mark, and leadered out of the
  % band because in place the label box collided with two of the twelve rings.
  % v11 REPAIRS A LEGIBILITY CLAIM THE LEDGER MADE AND THE RENDER NEVER
  % SUPPORTED. Pass 6 enlarged the arms 2.0 -> 2.8mm "because at 2mm only its
  % vertical arm was visible through the stems"; at 2.8mm only its vertical arm
  % was still visible, because the fault is not length. Several of the twelve
  % stems leave the vertex within 0.02cm of horizontal -- (-1.57,-0.085) at 3.1
  % degrees, (1.12,0.119) at 6.1, (1.40,-0.238) at -9.7 -- so the horizontal arm
  % is buried in a near-collinear bundle at EVERY radius and no arm length
  % separates it. What separates it is white: a 0.20cm disc of the band's own
  % tint (white, then accent at 0.12 -- byte-identical to the band, so the disc
  % has no visible edge) is laid over the stems, and the cross is drawn in it.
  % The same disc fixes a second finding: the label's leader used to be the
  % thirteenth grey line converging on a vertex the prose says twelve lines
  % reach. It is now the ONLY line that enters the clear disc, so it reads as a
  % leader and not as a spoke, and it is drawn after the disc for that reason.
  % Its label-end standoff also goes 0.03 -> 0.12cm: at 0.03 it grew out of the
  % final "n" of "mean" and read as a stray hairline off the type.
  \fill[white] (2.25,8.50) circle (0.20);
  \fill[accent, fill opacity=0.12] (2.25,8.50) circle (0.20);
  % Anchored base WEST on the panel's own x origin. Anchored base east at
  % x = 1.80 this label ran 0.32cm LEFT of x = 0, which widened the
  % picture's box past the fullwidth measure and put an Overfull hbox of
  % 6.00pt in the log: the INK bbox was inside the measure and the NODE
  % box was not, which is why the render looked clean and the log did not.
  \draw[furn] (2.14,7.88) -- (2.31,8.32);
  \draw[schem] (2.08,8.50) -- (2.42,8.50);
  \draw[schem] (2.25,8.33) -- (2.25,8.67);
  \node[qlab] at (0.00,7.80) {the global mean};

  % the instrument's sample, as a scope line and not as a sentence. Its
  % baseline is shared with panel (b)'s second measurement line, and both drop
  % 0.10cm in v11 behind the panel-(b) restack, so the two panels keep one grid.
  \node[qlab] at (0.00,7.20)
    {480 prompts: twelve drugs, forty real cells, seven layers};

  % definitional prose: \footnotesize roman ink, two lines, on the panel's own
  % x origin. It is the largest type in the panel, with the head and the axis
  % label, which is fig:mechanism's ranking.
  \node[def] at (0.00,6.62) {Twelve centred class means span at most eleven};
  \node[def] at (0.00,6.22) {directions; ten are kept.};

% ###########################################################################
% # PANEL (b) -- purify.  x origin 8.40.  Nothing here has a scale either.
% ###########################################################################
  \node[head] at (8.40,9.90) {(b) purify it against context};

  % ORDER MATTERS, AND v11 REVERSES IT. The raw band is drawn FIRST and the
  % opaque context rectangle over it. In v10 the pale diagonal was drawn last,
  % so where it crossed the opaque panelbg rectangle the two tints multiplied
  % and an UNLABELLED darker parallelogram appeared inside the context band --
  % exactly the region a cold reader takes for "where the drug slab and the
  % context subspace overlap", i.e. for the confound, sitting beside the patch
  % that actually carries that name. Pass 9 deleted the tint-multiply lozenge
  % as a claim; it did not delete the overprint that drew one. Now the only
  % tinted region inside the context subspace is the one that is labelled.
  % No coordinate and no colour changed to get this.
  \begin{scope}[shift={(10.50,8.45)}, rotate=38]
    \fill[accent, fill opacity=0.12] (-0.95,-0.16) rectangle (0.95,0.16);
  \end{scope}
  % v11: the label drops 8.14 -> 7.98. At 8.14 the word "raw" rose to y = 8.352
  % and sat ON the context band's 8.25 top edge and its outline.
  \draw[furn] (9.68,8.06) -- (9.98,7.95);
  \node[alab, anchor=base east] at (9.62,7.98) {$V_{\textrm{drug}}$, raw};

  % the context subspace: cell line, plate, dose. panelbg with a rulegrey
  % outline, which is the colour figs/slab.tex gives this same named object.
  % v11: both label lines drop 0.14cm -- see the restack note at V_pure.
  \filldraw[fill=panelbg, schem] (8.60,8.25) rectangle (11.80,8.65);
  \node[qlab] at (11.95,8.46) {context subspace:};
  \node[qlab] at (11.95,8.16) {cell line, plate, dose};

  % V_pure, at 90 degrees: what the raw slab becomes, and visibly thinner than
  % the raw band, against a measured 0.620 -> 0.031 printed below it.
  % v11 RESTACK, and it is one move with four dependents. The gloss sat 0.36pt
  % under the panel head's descenders -- 0.13mm, against 2.8mm of air under
  % panel (a)'s head -- so two panels of equal rank were set with head
  % clearances differing by a factor of twenty-two. Dropping the gloss 0.21cm
  % gives it 3.0mm, but it then landed on the right column's own first baseline
  % (9.32) with 0.15cm of horizontal air beside it, and the two read as one
  % line: "V_pure, the purified slab orthogonalise: remove the part". So the
  % right column drops 0.22cm, the context-subspace pair 0.14cm behind it, "the
  % confound" 0.18cm behind that, and the measurement strip and the two prose
  % blocks of BOTH panels 0.10cm behind that -- each step set by the measured
  % ink boxes, not by eye, at a uniform 0.066cm of white between blocks. The
  % bar's top comes to 9.15, still 0.115cm above the removal arrow's height,
  % which is all it has to do.
  \fill[accent] (10.41,7.95) rectangle (10.59,9.15);
  \node[alab, anchor=south] at (10.50,9.20) {$V_{\textrm{pure}}$, the purified slab};

  % THE CONFOUND: the orthogonal projection of the raw direction onto context,
  % from V_pure's right edge to the foot of the perpendicular dropped from the
  % raw band's tip, so nothing bisects it and it is exactly the length of the
  % removal arrow above it. Directly labelled, with a leader, because a cold
  % reader attached the name to the diagonal band instead of to this patch.
  % v11: the leader's tip now lands INSIDE the patch. In v10 it stopped at
  % (11.05,8.31), 0.03cm below the patch's lower edge and inside the context
  % band, so it pointed at the band -- which is precisely the misreading the
  % pass-9 rebuild of this mark exists to kill.
  \filldraw[fill=accent, fill opacity=0.40, draw=accent, line width=0.5pt]
    (10.59,8.33) rectangle (11.2486,8.57);
  \draw[furn] (11.30,7.88) -- (11.02,8.42);
  \node[alab] at (11.35,7.82) {the confound};

  % the projection line that defines it. Densely dotted quietink at 0.5pt is
  % this document's reserved stroke for a projection line, a magnification
  % leader or a zero rule -- a coordinate that decides nothing.
  \draw[quietink, line width=0.5pt, densely dotted]
    (11.2486,9.0349) -- (11.2486,8.57);

  % THE OPERATION, drawn as the projection it is and not as a rotation: a
  % straight arrow at the tip's own height, carrying the tip leftwards by
  % exactly the confound's length, landing on V_pure. It is the ONLY arrowhead
  % in the figure, and an arrowhead here means an operation acting in a
  % direction.
  \draw[accent, line width=0.5pt, -{Latex[length=3.4pt,width=2.8pt]}]
    (11.2486,9.0349) -- (10.61,9.0349);
  % v11: the leader now TOUCHES the arrow's tail. In v10 its lowest ink stopped
  % 0.105cm short of the tail and died in white, while the standoff at the
  % LABEL end was only 0.05cm -- the inversion of the corpus rule this file's
  % own header states twice (the 0.06cm gap goes at the label end, never at the
  % mark end, or the stub reads as a stray hairline).
  \draw[furn] (11.89,8.95) -- (11.2486,9.0349);
  \node[alab] at (11.95,9.10) {orthogonalise: remove the part};
  \node[alab] at (11.95,8.80) {of the raw slab lying in context};

  % ORTHOGONALITY AS A MARK, NOT AS A WORD, and v11 moves it out of the
  % removal arrow's line. In v10 the glyph's horizontal arm sat at y = 9.05 and
  % the arrow's shaft at y = 9.0349 -- 0.13mm apart, on opposite sides of an
  % opaque 0.18cm bar -- so a static assertion (perpendicularity) and a directed
  % operation (the projection) read as ONE horizontal primitive running from the
  % word "orthogonal", through V_pure, out to the dotted drop line. That is the
  % glyph collision S11 exists to forbid, and it was the only one left in the
  % figure. The arms come 0.40 -> 0.28cm and the arm drops to y = 8.93, which
  % puts 1.05mm of white under the arrow; 2.8mm is still 40% larger than the
  % 2mm a cold reader did not see until zooming, and the quadrant is unchanged.
  % THE QUADRANT IS UPPER-LEFT and always has been: v10 item 8 says "lower-left"
  % and is false about its own drawing. The other three quadrants are ruled out
  % by measurement -- lower-left and upper-right are crossed by the raw band at
  % 38 degrees, and lower-right is where "the confound" and its leader sit.
  \draw[ink, line width=0.5pt] (10.13,8.65) -- (10.13,8.93) -- (10.41,8.93);
  \node[lab, anchor=base east] at (9.99,8.73) {orthogonal};

  % the four measured numbers of stage two, as one strip on the panel's own x
  % origin. These are the only quantities above panel (c) with a source.
  % v11, two repairs. (i) COLOUR: the strip was ink and panel (a)'s scope line
  % is quietink, on the SAME baseline in adjacent panels at the same size, with
  % nothing on the page decoding the difference; both are scope, so both are
  % quietink now and ink is left to heads, definitional prose, the axis label
  % and the gate. (ii) NAME: "subspace fraction" becomes fig:mechanism's own
  % name for the quantity, "descriptive subspace fraction" (one object, one
  % name). The layer-9 qualifier stays and is load-bearing: layer 9 is also the
  % MAXIMUM of variance_share.drug over the seven depths (the range is 0.277 at
  % layer 16 to 0.620 here), which the drawing has no lane to state and which
  % the caption and figs/purify.json now carry.
  \node[qlab] at (8.40,7.48)
    {Dose decodes from inside the raw slab at $0.90$--$0.95$;};
  \node[qlab] at (8.40,7.20)
    {descriptive subspace fraction at layer 9: raw $0.620$, purified $0.031$.};

  \node[def] at (8.40,6.62) {$V_{\textrm{pure}}$ is the part of the raw slab that is};
  \node[def] at (8.40,6.22) {orthogonal to the context subspace.};

\end{scope}

% ###########################################################################
% # PANEL (c) -- the removal gate.  x origin 1.60, the fraction axis's spine.
% # The only measured marks in the figure.
% ###########################################################################
  \node[head] at (1.60,5.20) {(c) the removal gate};

  % WHICH SLAB THE GATE IS MEASURED ON, restored to the drawing in v11 and set
  % in axis-label rank rather than at the foot. Without it the last line a
  % reader crosses before entering panel (c) is panel (b)'s "so causal claims
  % use the purified slab, not the raw", and the inference that the gate was
  % run on V_pure is unavoidable. workspace_probe.json does not record which
  % subspace was projected out; it stores one ablated series and no purified
  % counterpart, so the line states only what the file supports.

  % the region of outcomes that would have made every null in this section a
  % measurement of nothing. Drawn FIRST so every rule and dot sits on top of
  % it, at rulegrey!8 (grey 247) and not rulegrey!12 (grey 243): at 12 it was
  % within one grey level of panelbg, the tint that means "a schematic
  % subspace" 3cm above it in panel (b), so one tint carried two meanings.
  % Nothing is written inside it: its emptiness is the whole point, and the
  % two italic lines that used to sit in it competed with the marks. It is
  % NAMED instead, outside its own right edge and under the rule that bounds
  % it, because the largest mark on the page may not be the one mark a reader
  % cannot name (and because the name makes the third clause of this figure's
  % acceptance test answerable from the drawing).
  \fill[rulegrey!8] (1.60,1.75) rectangle (14.10,\fy{0.8});

  % ---- fraction axis, DEFINED AT THE AXIS ---------------------------------
  % Rotated 90 degrees reading upward, immediately left of the tick numerals,
  % which is fig:mechanism's form and the answer to the supervisor's first
  % complaint. Nothing else is in this block: no result, no interval note, no
  % cross-reference.
  % v11, AND THIS IS A CORRECTNESS REPAIR AND NOT A TYPOGRAPHIC ONE. v10 set
  % one \footnotesize line reading "fraction of drug signal removed" over one
  % \scriptsize construction line. Two faults, both fatal.
  %   (a) THE NAME WAS FALSE. The plotted series is frac_signal_removed, which
  %       is ABOVE-CHANCE signal: (before - after)/(before - chance). The
  %       unqualified quantity (before - after)/before is 0.7041 at layer 2 and
  %       0.8528 at layer 9, four to eighteen points BELOW the marks drawn. Read
  %       at the drawn name the top mark, 0.949, implies a residual accuracy of
  %       0.041, under the chance 0.083 this figure prints in its own foot note.
  %       The thesis's pre-registration (04b:218) says "the fraction of
  %       above-chance signal removed must exceed $0.8$". v10's stated reason for
  %       the omission -- that the 43-character form needs 7.14cm against a
  %       5.71cm lane -- ruled out one STRING, not the correct QUANTITY, and is
  %       superseded: the name is broken over two rotated \footnotesize lines.
  %   (b) THE SUBORDINATE LINE WAS THE TALLEST OBJECT IN THE PANEL. Measured on
  %       the v10 render, the \scriptsize construction line ran 5.52cm against a
  %       3.85cm axis (143%) and its top sat 6.0pt ABOVE the panel head's, so the
  %       smallest type in the block outranked the head. S5 makes that line
  %       optional; it is deleted, and the construction it carried moves to the
  %       foot note, where the three accuracies it is built from already sit and
  %       where the missing noun ("probe accuracy") costs nothing.
  % LANE, remeasured on the v11 render and reconciled with figs/purify.json,
  % which the v10 header disagreed with: \footnotesize sets at 0.158 cm per
  % character in this class, so "fraction of above-chance" (24 characters) runs
  % 3.86cm against a 3.00cm axis and "drug signal removed" (19) runs 3.03cm.
  % The block is centred on the axis midpoint, y = 3.25.
  \draw[axisr] (1.60,1.75) -- (1.60,4.75);
  \foreach \v/\t in {0/0, 0.5/0.5, 1/1.0}{
    \draw[furn] (1.60,\fy{\v}) -- (1.50,\fy{\v});
    \node[qlab, anchor=east, inner sep=0pt] at (1.44,\fy{\v}) {\t};}
  % LINE ORDER. Rotated 90 degrees the line-advance direction is -x, so the
  % FIRST line is the leftmost and the LAST line is the one against the tick
  % numerals, which is also matplotlib's own stacking for a two-line ylabel.
  % Set the other way round the block reads "drug signal removed / fraction of
  % above-chance", which is what the first v11 build printed.
  \node[text=ink, font=\footnotesize, rotate=90, anchor=center]
    at (0.36,3.25) {fraction of above-chance};
  \node[text=ink, font=\footnotesize, rotate=90, anchor=center]
    at (0.80,3.25) {drug signal removed};

  % ---- layer axis ---------------------------------------------------------
  % "layer" is left-aligned to the axis's own left end, per the contract, and
  % v11 puts 2.8mm of white under the tick numerals instead of v10's 0.63mm.
  % At 0.63mm the word sat directly under the numeral 2, one size LARGER than
  % it, and the pair read as "layer 2" -- the identical fault the pass-9 note
  % above records against the old placement under the layer-9 numeral. The
  % repair is vertical, not horizontal: the fraction origin rose 1.45 -> 1.75,
  % which is where the 2.8mm came from. Baseline to baseline the numerals and
  % the word are now 0.53cm apart, against the 0.40cm leading this figure uses
  % between two lines of one sentence, so they can no longer read as one block.
  \draw[axisr] (1.60,1.75) -- (14.10,1.75);
  \foreach \l in {2,4,6,8,9,12,16}{
    \draw[furn] (\lx{\l},1.75) -- (\lx{\l},1.65);
    \node[qlab, anchor=north, inner sep=0pt] at (\lx{\l},1.59) {\l};}
  \node[def] at (1.60,0.85) {layer};

  % ---- the gate: a pre-set threshold, so BLACK, SOLID, 0.9pt --------------
  % It was accent 0.6pt dashed. Accent in this document means the object under
  % test, and a threshold is not the object under test.
  \draw[gate] (1.60,\fy{0.8}) -- (14.10,\fy{0.8});
  \node[lab, anchor=west, text=black] at (14.25,\fy{0.8}) {gate, $0.8$};
  % the name of the shaded region, set outside it and directly under the rule
  % that bounds it. Quietink, because the region is a counterfactual and not
  % the object under test.
  \node[qlab, anchor=west] at (14.25,2.95) {an untestable null};

  % ---- the seven measured fractions. No connecting line: the series is flat.
  \gdot{2}{0.8846153846153847}
  \gdot{4}{0.9253112033195021}
  \gdot{6}{0.9383561643835617}
  \gdot{8}{0.9461077844311377}
  \gdot{9}{0.9491525423728814}
  \gdot{12}{0.9440993788819876}
  \gdot{16}{0.9444444444444445}
  % direct label for the set, on the last mark's own baseline, with the 0.30pt
  % leader the corpus requires beyond 3mm: the gap from the layer-16 dot to the
  % label's ink is 3.6mm, and in v10 that gap was open with nothing in it.
  \draw[accent, line width=0.3pt] (14.19,\fy{0.9444444444444445})
                              -- (13.95,\fy{0.9444444444444445});
  \node[alab, anchor=west] at (14.25,\fy{0.9444444444444445})
    {all seven layers: $88$--$95\%$};

% ###########################################################################
% # the figure's foot: one two-line \scriptsize block on the figure's own x
% # origin. Line one names the instrument and the three accuracies the axis is
% # built from; line two gives the construction and this figure's single
% # interval declaration.
% # v11: the noun was missing. v10 printed "Before 0.41--0.82, after 0.121 ...,
% # chance 0.083" with no quantity attached to any of the three, and the words
% # "accuracy", "probe" and "1/12" appeared nowhere in the figure, so panel (c)
% # never said what instrument it was reading. The registry name is "probe
% # accuracy" and chance is "0.083 = 1/12", which is the stored float exactly
% # (workspace_probe.json :: chance == 1/12 in IEEE double).
% ###########################################################################
  \node[qlab] at (0.00,0.32)
    {Probe accuracy before $0.41$--$0.82$, after $0.121$ at every layer,
     chance $0.083$, exactly $1/12$.};
  \node[qlab] at (0.00,0.02)
    {The fraction is (before minus after)$/$(before minus chance). Intervals:
     none returned in (c); (a) and (b) carry no measured mark.};

  % the bounding box IS the fullwidth measure; nothing may stick out of it
  \path (0,-0.10) rectangle (17.30,10.35);
\end{tikzpicture}
%
% ---------------------------------------------------------------------------
% READY TO PASTE into Sections/04b-representation.tex, at sec:res-used,
% immediately after the paragraph "Step three, the removal gate"
% (04b-representation.tex:182-190) and BEFORE tab:workspace. The same block is
% staged in figs/_PENDING_FLOATS.tex for a human to integrate.
%
% Do NOT merge this float with fig:hook. The two are read six paragraphs apart;
% each is about 100mm tall against a corpus ceiling of 117mm, so stacked they
% spill the page; and the house caption shape carries one bolded finding, not
% two.
%
% CAPTION, v10 AND v11. The drawing lost eight sentences in v10 and one more in
% v11, and every one is folded into the caption below, because in this document
% the argument lives in the caption and the drawing names objects.
% What was cut from the drawing, in the order it used to read:
%   1. the three \footnotesize italic panel subtitles -- "where the twelve drugs
%      differ, at one layer", "the raw slab is partly a batch slab", and "the
%      fraction of above-chance drug signal that projecting the drug subspace
%      out destroys". The third is now the axis label, which is where it always
%      belonged; the other two are caption clauses.
%   2. the centred divider sentence, "an ablation is only a test of the drug if
%      the slab it removes really holds the drug", which restated the body
%      paragraph directly above the float.
%   3. the two lines inside the fail region, "had any layer landed in here, the
%      null below it / would have been unfalsifiable --- a measurement of
%      nothing", which restated the caption's own "the shaded region ... is
%      empty" and sat inside the data area competing with the marks. v11 gives
%      the region the NAME "an untestable null", set outside its right edge;
%      the argument stays in the caption.
%   4. the centred discipline line, "the dots and the bands are schematic and
%      have no scale; only the four numbers in stage two are measured". Its job
%      is done by the GLYPH now -- open ring against filled disc -- and the
%      caption says the rest.
%   5. the conservatism argument, "so the slab keeps essentially the whole span,
%      noisiest included, which for a removal claim is conservative", cut from
%      panel (a)'s prose, which is now definitional only.
%   6. "by construction", cut from the right angle's label for want of a lane
%      3.6cm wide at the glyph's own height (v10 item 8).
%   7. "measured on the slab that is projected out; the stored file has no
%      separate purified-slab gate", cut from the foot note in v10 for text
%      density. RESTORED TO THE DRAWING IN v11, reworded, under panel (c)'s
%      head, because without it panel (c) named no slab at all and sat directly
%      under panel (b)'s "so causal claims use the purified slab, not the raw".
%      The caption's own statement of the same limit is reworded below so that
%      no clause of five or more words appears in both places.
%   8. "four \quad what the ablation costs: \cref{fig:hook}", the forward
%      pointer that used to sit on the x-tick baseline where the word "four"
%      read as a leftmost tick.
%   9. NEW IN v11: the axis's \scriptsize construction line, "(before minus
%      after)/(before minus chance)". It did not vanish -- it is the second foot
%      line now -- but it left the axis block, which S5 allows and which the
%      render forced: at 5.52cm against a 3.85cm axis it was the tallest object
%      in the panel and rose above the panel head's own top.
%
% CLAUSES OF THE PRINTED CAPTION THAT ARE NOW FALSE AGAINST THE DRAWING, with
% their lines in Sections/04b-representation.tex. Nobody may edit that file from
% this pass, so they are listed here and in figs/purify.json, and the
% replacement for each is in the caption below.
%   04b:234  "\emph{Above the rule, nothing has a scale.}" -- there is no rule.
%            The full-measure divider at y = 4.72 was deleted in v10 and the
%            panel letters carry the scoping now.
%   04b:234  "The dots, the bands, the arrow and the right angle are schematic"
%            -- v10 turned panel (a)'s twelve schematic marks into OPEN RINGS,
%            so the only dots in the figure are panel (c)'s seven MEASURED
%            discs, and the sentence as printed now calls the measurements
%            schematic. Replacement: "open rings and untitled bands".
%   04b:236  "the four numbers printed in stage two" -- there is no stage two;
%            the panel is lettered (b) and the word "stage" went with the
%            banners.
%   04b:243  "\emph{Below the rule, everything has a scale.}" -- same rule.
%   04b:250  "Two limits, both stated on the drawing" -- only one of the two is
%            on the drawing (the which-slab limit, restored in v11); the
%            twelve-class against twenty-four-class limit is caption-only and
%            always was.
%   04b:250-252  the which-slab limit as printed shares a clause of more than
%            five words with the drawing's new scope line, and it asserts more
%            than the file supports; it is reworded below.
% AND ONE THING THE CAPTION OWES THE READER AND HAS NEVER SAID: the printed
% pair "raw $0.620$, purified $0.031$" is at layer 9, and layer 9 is the
% MAXIMUM of layers.<L>.variance_share.drug over the seven depths. The stored
% range is 0.2766847391691891 at layer 16 to 0.6201207831136324 here, so the
% drawn pair is the largest raw-to-purified contrast the file holds. The raw
% series is drawn in no other figure and quoted in no sentence of the thesis
% (figs/fig-mechanism.json records it as drug_raw_not_drawn), so the caption is
% the reader's only chance to learn the spread. It is added below.
% The dose qualifier the v10 proposal added is kept: "$0.90$--$0.95$" is true
% only to two places, the stored minimum being 0.8979166666666666 at layers 12
% and 16, which is also why no reference rule is drawn at $0.90$.
%
% NOTHING IN THE CAPTION RESTATES A CLAUSE OF FIVE OR MORE WORDS FROM THE
% DRAWING. Checked string by string against the v11 render's own text spans: the
% drawing names objects and prints numbers, the caption argues. The drawing says
% "the subspace projected out"; the caption says "whichever subspace the
% ablation removed", and the longest run shared with any drawn string is four
% words.
%
% Two caption details, checked against a build of the whole float (returncode 0,
% one page, zero Overfull/Underfull boxes; the \cref targets resolve in the
% document and show as ?? only in an isolated harness):
%   * the two numerals in the bold lead are set in TEXT and not as $88$/$95$.
%     In math mode they come out upright and unbolded inside \textbf, so a
%     bolded finding would have had two unbolded numbers in the middle of it.
%   * the FIGURE BODY carries no \cref. The "four ... fig:hook" pointer was
%     deleted in v10, so the caption is the only place the sibling float is
%     named.
%
% \begin{figure}[htbp]
% \begin{fullwidth}
% \centering
% \input{figs/purify}
% \caption[The drug slab and its removal gate]{\textbf{The slab is built, not
% found --- ten directions fitted to twelve centred class means, then
% orthogonalised against the batch structure they are confounded with --- and it
% earns the right to be ablated only because projecting it out destroys 88 to 95
% per cent of the decodable drug signal at every depth.} An ablation tests the
% drug only if the subspace it removes really carries the drug, which is what
% panel (c) settles before \cref{fig:hook} spends it.
% \emph{Panels (a) and (b) have no scale.} Their marks are open rings and
% untitled bands for that reason: nothing in them is a measurement, and the only
% quantities there with a source are the four numbers printed beneath panel (b),
% all from \texttt{workspace\_probe.json}. Panel (a) is where the twelve drug
% classes differ at one depth: ten directions are kept out of the eleven such a
% set can span, so the noisiest are kept too, which for a removal claim errs the
% safe way. Panel (b) is why the raw subspace may not be used --- it is partly a
% batch subspace --- and the darker patch on the horizontal band is the
% component projected away, the arrow carrying the raw tip left by exactly that
% much. The right angle is not measured; it is true by construction. Both
% printed readings there are quoted at their least flattering: the dose range
% $0.90$--$0.95$ is rounded to two places, the stored minimum being $0.8979$ at
% layers $12$ and $16$, which is why no reference rule is drawn at $0.90$; and
% layer $9$ is the depth at which the raw subspace fraction is largest, the
% seven-layer range running $0.277$ to $0.620$, so the shrink to $0.031$ printed
% there is the widest the file holds and not a typical one.
% \emph{Panel (c) has a scale, and it is the only one that does.} Each mark is
% one depth's share of above-chance drug-label accuracy destroyed when the
% subspace is projected out, over twelve drug classes at chance $1/12$. The
% shaded region under the rule is the set of outcomes that would have left this
% section's nulls with nothing to measure, and it is empty at all seven depths.
% \textbf{Why no interval is drawn:} the quantity is a single ratio of two
% accuracies over \num{480} prompts and the instrument stores no dispersion of
% any kind for it. Two limits. \texttt{workspace\_probe.json} does not record
% whichever subspace the ablation removed --- it holds one ablated series and no
% purified counterpart --- so neither the drawing nor this caption attributes
% the gate to $V_{\textrm{pure}}$; where both are stored, in
% \cref{sec:res-mechanism}'s residual arm, the purified figure is the smaller of
% the two. And these twelve-class accuracies may not be read against the
% twenty-four-class sweep over the five arms of \cref{sec:res-mechanism}, whose
% chance is $1/24$ on a different label set. Step four is \cref{fig:hook}.}
% \label{fig:purify}
% \end{fullwidth}
% \end{figure}
% ---------------------------------------------------------------------------

% \caption[The drug slab and its removal gate]{\textbf{The slab is built, not
% found --- ten directions fitted to twelve centred class means, then
% orthogonalised against the batch structure they are confounded with --- and it
% earns the right to be ablated only because projecting it out destroys 88 to 95
% per cent of the decodable drug signal at every depth.} An ablation tests the
% drug only if the subspace it removes really carries the drug, which is what
% panel (c) settles before \cref{fig:hook} spends it.
% \emph{Panels (a) and (b) have no scale.} Their marks are open rings and
% untitled bands for that reason: nothing in them is a measurement, and the only
% quantities there with a source are the four numbers printed beneath panel (b),
% all from \texttt{workspace\_probe.json}. Panel (a) is where the twelve drug
% classes differ at one depth: ten directions are kept out of the eleven such a
% set can span, so the noisiest are kept too, which for a removal claim errs the
% safe way. Panel (b) is why the raw subspace may not be used --- it is partly a
% batch subspace --- and the darker patch on the horizontal band is the
% component projected away, the arrow carrying the raw tip left by exactly that
% much. The right angle is not measured; it is true by construction. Both
% printed readings there are quoted at their least flattering: the dose range
% $0.90$--$0.95$ is rounded to two places, the stored minimum being $0.8979$ at
% layers $12$ and $16$, which is why no reference rule is drawn at $0.90$; and
% layer $9$ is the depth at which the raw subspace fraction is largest, the
% seven-layer range running $0.277$ to $0.620$, so the shrink to $0.031$ printed
% there is the widest the file holds and not a typical one.
% \emph{Panel (c) has a scale, and it is the only one that does.} Each mark is
% one depth's share of above-chance drug-label accuracy destroyed when the
% subspace is projected out, over twelve drug classes at chance $1/12$. The
% shaded region under the rule is the set of outcomes that would have left this
% section's nulls with nothing to measure, and it is empty at all seven depths.
% \textbf{Why no interval is drawn:} the quantity is a single ratio of two
% accuracies over \num{480} prompts and the instrument stores no dispersion of
% any kind for it. Two limits. \texttt{workspace\_probe.json} does not record
% whichever subspace the ablation removed --- it holds one ablated series and no
% purified counterpart --- so neither the drawing nor this caption attributes
% the gate to $V_{\textrm{pure}}$; where both are stored, in
% \cref{sec:res-mechanism}'s residual arm, the purified figure is the smaller of
% the two. And these twelve-class accuracies may not be read against the
% twenty-four-class sweep over the five arms of \cref{sec:res-mechanism}, whose
% chance is $1/24$ on a different label set. Step four is \cref{fig:hook}.}
% \label{fig:purify}
% \end{fullwidth}
% \end{figure}
% ---------------------------------------------------------------------------

% \caption[The drug slab and its removal gate]{\textbf{The slab is built, not
% found --- ten directions fitted to twelve centred class means, then
% orthogonalised against the batch structure they are confounded with --- and it
% earns the right to be ablated only because projecting it out destroys 88 to 95
% per cent of the decodable drug signal at every depth.} An ablation tests the
% drug only if the subspace it removes really carries the drug, which is what
% panel (c) settles before \cref{fig:hook} spends it.
% \emph{Panels (a) and (b) have no scale.} Their marks are open rings and
% untitled bands for that reason: nothing in them is a measurement, and the only
% quantities there with a source are the four numbers printed beneath panel (b),
% all from \texttt{workspace\_probe.json}. Panel (a) is where the twelve drug
% classes differ at one depth: ten directions are kept out of the eleven such a
% set can span, so the noisiest are kept too, which for a removal claim errs the
% safe way. Panel (b) is why the raw subspace may not be used --- it is partly a
% batch subspace --- and the darker patch on the horizontal band is the
% component projected away, the arrow carrying the raw tip left by exactly that
% much. The right angle is not measured; it is true by construction. Both
% printed readings there are quoted at their least flattering: the dose range
% $0.90$--$0.95$ is rounded to two places, the stored minimum being $0.8979$ at
% layers $12$ and $16$, which is why no reference rule is drawn at $0.90$; and
% layer $9$ is the depth at which the raw subspace fraction is largest, the
% seven-layer range running $0.277$ to $0.620$, so the shrink to $0.031$ printed
% there is the widest the file holds and not a typical one.
% \emph{Panel (c) has a scale, and it is the only one that does.} Each mark is
% one depth's share of above-chance drug-label accuracy destroyed when the
% subspace is projected out, over twelve drug classes at chance $1/12$. The
% shaded region under the rule is the set of outcomes that would have left this
% section's nulls with nothing to measure, and it is empty at all seven depths.
% \textbf{Why no interval is drawn:} the quantity is a single ratio of two
% accuracies over \num{480} prompts and the instrument stores no dispersion of
% any kind for it. Two limits. \texttt{workspace\_probe.json} does not record
% whichever subspace the ablation removed --- it holds one ablated series and no
% purified counterpart --- so neither the drawing nor this caption attributes
% the gate to $V_{\textrm{pure}}$; where both are stored, in
% \cref{sec:res-mechanism}'s residual arm, the purified figure is the smaller of
% the two. And these twelve-class accuracies may not be read against the
% twenty-four-class sweep over the five arms of \cref{sec:res-mechanism}, whose
% chance is $1/24$ on a different label set. Step four is \cref{fig:hook}.}
% \label{fig:purify}
% \end{fullwidth}
% \end{figure}
% ---------------------------------------------------------------------------

\caption[The drug slab and its removal gate]{\textbf{The slab is built, not
found --- ten directions fitted to twelve centred class means, then
orthogonalised against the batch structure they are confounded with --- and it
earns the right to be ablated only because projecting it out destroys $88$ to
$95$ per cent of the decodable drug signal at every depth.} An ablation is a
test of the drug only if the subspace it removes really carries the drug, which
is what panel~(c) settles before \cref{fig:hook} spends it.
\emph{Panels (a) and (b) have no scale.} Their marks are open rings and
untitled bands for that reason: nothing in them is a measurement, and the only
quantities there with a source are the four numbers printed beneath panel~(b),
all from \texttt{workspace\_probe.json}. Panel~(a) is where the twelve drug
classes differ at one depth: ten directions are kept out of the eleven such a
set can span, so the noisiest are kept too, which for a removal claim errs the
safe way. Panel~(b) is why the raw subspace may not be used --- it is partly a
batch subspace, and every causal claim in this chapter is made on
$V_{\textrm{pure}}$ and never on the raw slab. The darker patch on the
horizontal band is the component projected away, the arrow carrying the raw tip
left by exactly that much; the right angle is not measured, it is true by
construction. The printed dose range is rounded to two places: the stored
minimum is $0.8979$, at layers $12$ and $16$, which is also why no reference
rule is drawn at $0.90$.
\emph{Panel (c) has a scale, and it is the only one that does.} Each mark is
one depth's share of above-chance drug-label accuracy destroyed when the
subspace is projected out, over twelve drug classes at chance $1/12$. The
shaded region under the rule is the set of outcomes that would have left every
null in this section untestable, a measurement of nothing, and it is empty at
all seven depths. \textbf{Why no interval is drawn:} the quantity is a single
ratio of two accuracies over \num{480} prompts and the instrument stores no
dispersion of any kind for it. Two limits, both recorded in
\texttt{figs/purify.json} and neither drawn. The gate is measured on the
subspace that is projected out, and this file stores no separate gate for
$V_{\textrm{pure}}$. And these twelve-class accuracies may not be read against
the twenty-four-class sweep over the five arms of \cref{sec:res-mechanism},
whose chance is $1/24$ on a different label set. Step four is \cref{fig:hook}.}
\label{fig:purify}
\end{fullwidth}
\end{figure}

\paragraph{Step four, the causal measurement in logit space} The true sentence is
teacher-forced and its per-token log-probabilities read off clean; a forward hook
then projects the slab out at every position ($h \leftarrow h - (hV)V^{\top}$)
and they are read again. The measure is the mean KL divergence from clean to
ablated over the target tokens, on $60$ prompts.\gloss{teacher forcing}{the
true sentence is scored token by token, never generated} Logit space repairs the
project's first causal attempt, which injected a steering vector and generated;
that vector was small against the residual-stream norm, its effect sat under the
sampling-noise floor, and its ``no effect'' was uninterpretable.\corrmark{the
steering-vector attempt this replaces, \cref{sec:app-identity}} The ablation runs
for $V_{\textrm{pure}}$ (headline), a matched-dimension random subspace (the
null), and the cell-line slab (the positive control). The prompt count is small
and no interval is quoted; what intervals the probe arc has come from the
identity test of \cref{sec:res-mechanism}, which is discounted there in turn, so
no conclusion in this chapter rests on the probe arc alone.

% ===========================================================================
\begin{figure}[htbp]
\begin{fullwidth}
\centering
% ===========================================================================
% figs/hook.tex -- fig:hook, step four of sec:res-used drawn as what it
% physically is: one true sentence, laid out as a token strip, scored twice by
% the same weights, with a forward hook that zeroes the drug slab at every
% position in between, and the KL read only over the target tokens.
%
% FIGURE BODY ONLY. No preamble, no document environment, no float, no caption.
% \input it from inside a figure/fullwidth float in
% Sections/04b-representation.tex, at sec:res-used, immediately after the
% paragraph "Step four, the causal measurement in logit space" (04b:191-203).
% The complete float block sits COMMENTED at the foot of this file.
%
% Compiles against thesis_v2/preamble.tex and nothing else: no \includegraphics,
% no external data, no TikZ library beyond the ones preamble.tex already loads
% (this file uses arrows.meta and decorations.pathreplacing).
%
% THE MISREADING THIS FIGURE EXISTS TO KILL
% -----------------------------------------
% A reader who has met "ablation" in a generative setting arrives with a default
% picture: the model writes a sentence, the hook is switched on, the model
% writes a DIFFERENT sentence, and the KL measures how far the two sentences
% drifted apart. Every clause of that is false here. Nothing is sampled. The
% sentence is fixed, real, and identical in both passes. And the hook acts on
% the prompt as well as on the target, so it is not an intervention "on
% generation" at all. The figure is built so that all three corrections are
% legible from the marks alone: the two strips are drawn by ONE macro so they
% are provably identical box for box; the nine hook ticks include the four that
% land on the prompt; and the two braces span the target boxes only and stop
% dead at the prompt/target divider.
%
% EVERY NUMBER DRAWN, AND WHERE IT COMES FROM
% -------------------------------------------
% Source: RESULTS_cluster/workspace_probe.json, key layers.4.kl. Config there:
% tier tier2_unseen_drugs, n_drugs 12, n_per_drug 40, n_dims 10,
% n_kl_prompts 60, seed 42. Sibling ledger of the exact values drawn:
% thesis_v2/figs/hook.json.
%
%   mark                    key path                       source value      drawn
%   purified drug slab      layers.4.kl.pure_drug[0]       1.5631485521793365  1.563
%   matched random null     layers.4.kl.random_matched[0]  0.08388670384883881 0.084
%   cell-line control       layers.4.kl.cell_line[0]       4.179661695162455   4.180
%   n per entry             layers.4.kl.*[2] = config.n_kl_prompts             60
%   slab rank               layers.4.pure_drug_dims = config.n_dims            10
%
%   The three dots are the ONLY data-bearing coordinates in the figure. Their
%   positions are \kx{v} = 8.20 + v/4.5*8.70, a linear map of nats to cm:
%     0.084 -> 8.3624 ,  1.563 -> 11.2218 ,  4.180 -> 16.2813
%
%   COMPUTED, not stored as floats:
%   ratio to the null       pure_drug[0]/random_matched[0] 18.634044258027835  18.6x
%   share of the control    pure_drug[0]/cell_line[0]      0.3739892522853049  37%
%   Both rounded forms are also embedded verbatim in the prose string
%   layers.4.ablation_note, which reads "...18.6x the functional variance of a
%   matched-random slab (37% of the cell-line control)...".
%
%   NOT drawn: layers.4.kl.pure_drug[1] = 0.021199743463084047, and its two
%   siblings. These are normal-approximation standard errors over the 60
%   prompts, not this document's interval convention (95% bootstrap clustered
%   over cell lines). figs/fig-mechanism.json already records intervals_drawn
%   false for exactly this reason, and mechanism.py's header states it: drawing
%   a +/-0.003 whisker beside the thesis's genuine +/-0.03 intervals would
%   invite a reader to distrust all of them. The three marks are therefore bare
%   solid dots, which is how this document says "no interval", and the caption
%   says why.
%
%   2048 (the residual width) is prose only: 04b-representation.tex:383 and
%   07-Appendix.tex, sec:app-identity. It appears here as a LABEL ("ten of the
%   2048 directions"), never as a coordinate. Nothing in the vector panel is to
%   scale and the caption says so.
%
% AGREEMENT WITH THE TEXT (checked, not assumed):
%   04b-representation.tex:229  tab:workspace layer-4 row "1.563 & 0.084 &
%                               $18.6\times$"                            -> OK
%   04b-representation.tex:193  "projects the slab out at every position
%                               ($h \leftarrow h - (hV)V^{\top}$)"
%                               -> the formula is set verbatim and the nine
%                                  ticks are what "every position" means.  OK
%   04b-representation.tex:195  "mean KL divergence from clean to ablated over
%                               the target tokens, on $60$ prompts"
%                               -> the braces span the target only; n = 60 is
%                                  printed under the axis.                OK
%   04b-representation.tex:195  \gloss{teacher forcing}{the true sentence is
%                               scored token by token, never generated}
%                               -> the band-one sentence and both subtitles. OK
%   04b-representation.tex:200-202 "The ablation runs for $V_{\textrm{pure}}$
%                               (headline), a matched-dimension random subspace
%                               (the null), and the cell-line slab (the positive
%                               control)" -> the three dots are named in exactly
%                               those three roles, in that order of size.  OK
%   04b-representation.tex:254  "3.6--18.6x" -> layer 4 IS the 18.6.        OK
%   04b-representation.tex:260  "9--37% of what removing cell line costs at
%                               layers 2--9" -> layer 4 IS the 37%.        OK
%   04b-representation.tex:286-288 "No intervals are drawn ... normal-
%                               approximation standard errors over the 60
%                               prompts" -> inherited here.                OK
%   03-Apparatus.tex:249-252    "its expressed panel genes in descending order,
%                               terminated by \ENDCELL. The prompt supplies cell
%                               line, drug, dose, mechanism and the control
%                               cell's sentence" -> the five prompt fields, in
%                               that order, and the terminator.            OK
%   The layer-4 cell-line KL, 4.180, is printed in NO sentence of the thesis;
%   the section quotes only the ratios built from it (04b:258-262). It is a new
%   reading of a stored field and hook.json flags it as such. Precedent:
%   figs/slab.json does the same for variance_share.context.
%
% DESIGN DECISIONS, each with its reason
% --------------------------------------
%  1. BOTH STRIPS ARE DRAWN BY ONE MACRO, \tokstrip, taking only the y of the
%     box bottoms. Two hand-written strips would be two chances to differ, and
%     the single claim the figure has to carry is that they do NOT differ. The
%     macro makes "identical box for box" a property of the source, not of my
%     care.
%  2. THE FIVE PROMPT FIELDS AND THEIR ORDER ARE NOT MINE. They are
%     03-Apparatus.tex:250-252 verbatim, and the same five in the same order as
%     figs/licensing.tex's strip, so the two figures read as one object seen
%     twice. A sixth field would be unsourced and would break that.
%     PASS 9: the two strips now also share a GLYPH. A reviewer found that the
%     words agreed and the drawing did not -- hook drew rounded panelbg boxes
%     with rounded gutters and licensing drew square-cornered WHITE cells
%     sharing edges, so a reader met the same object twice and was shown two
%     objects. hook's treatment was judged the better mark and licensing.tex
%     adopted it (its cell/.style now carries fill=panelbg, rounded
%     corners=1pt, and its drug cell \tokdrug's accent 0.4pt border). Nothing
%     in THIS file changed for it; the cell edges still differ, because the two
%     strips are 8.20cm and 7.20cm and each distributes its own slack.
%  3. THE ELLIPSIS BOX is the mitigation for this figure's largest honesty
%     exposure. The real target is a cell sentence of hundreds of gene tokens;
%     four boxes is a schematic. "no token count is implied" appears in the
%     caption and in the ledger. Rejected alternative: drawing twenty boxes,
%     which implies a count just as firmly and only moves the lie.
%  4. THE DRUG BOX IS THE ONE TINTED PROMPT BOX (accent!12, accent rule). It is
%     the one place the prompt names the drug in plain text -- the fact that
%     makes sec:methods-probe's positive nearly free -- carried forward here so
%     the reader sees the hook deleting downstream of it. It is also the only
%     accent inside either strip, so the strips themselves stay furniture.
%  5. THE HOOK RULE SITS ABOVE THE PASS-TWO STRIP, not below it, so the eye
%     meets the intervention before it meets the boxes it acts on. The RULE is
%     what says "everywhere along here": it spans the strip's full width,
%     prompt half included, which is the correction to the reader's default
%     assumption that an intervention on generation touches only what is
%     generated.
%     REPAIRED, PASS 9: the nine descending TICKS are gone, replaced by one
%     Latex arrow into each half at 0.5pt, stopping 0.06cm above the box tops.
%     Two faults, found independently by both reviewers. (a) A tick per BOX is
%     not a tick per POSITION. "control cell sentence" is one box and hundreds
%     of positions; the ellipsis box is an unknown number. Counting a comb and
%     mapping it to positions is the first thing a reader does with a comb, and
%     it gave the wrong answer -- while the caption simultaneously had to say
%     "the token strip is schematic and implies no token count". (b) The ticks
%     were accent at nearly the rule's own weight and each terminated exactly
%     on its box's top border with no gap; for the DRUG box, whose border is
%     also accent, tick and border fused into one continuous stroke, so the one
%     cell that must be found first sprouted a flagpole and stopped reading as
%     a highlighted field. Marks running into rules uninterrupted is precisely
%     what the house rule forbids. The comb also read as a categorical axis,
%     which is exactly what the WHAT COMES OUT panel two bands below is. Two
%     arrows say "down into both halves" and assert no count; the words "at
%     every position, prompt and target alike" under the strip carry the claim.
%  5b.NO SEPARATE TAG ON THE RULE. One reading "the hook" was drawn at
%     (8.70,7.40) and taken out on the next build: PASS TWO's subtitle now ends
%     "...the hook is on one layer's residual stream" on baseline 7.58 and runs
%     to x = 13.5, so a tag 0.18cm beneath it collided with it and named the
%     same object a third time. The rule is named above by that subtitle and
%     below by the accent formula.
%  6. THE RULE RUNS UNBROKEN ACROSS THE 1.20cm GAP between the prompt strip and
%     the target strip. The gap is a reading device (it is where the divider
%     is); the hook does not know about it. Breaking the rule there would draw
%     the exact exemption the figure denies.
%  7. THE TWO BRACES SPAN THE TARGET SIDE ONLY. That the brace covers the
%     target and no more is the whole answer to "why is the KL taken over the
%     target tokens", and it is drawn rather than said.
%     PASS 9: they now start at the first target box's own left edge, 9.20, and
%     not at the prompt/target divider, 8.70. The caption says "the braces span
%     the target tokens only" and a reader is invited to test that against
%     "prompt boxes included" for the hook rule immediately above, so a brace
%     that began 0.50cm before the first token was contradicting its own label.
%     They are rulegrey, not accent:
%     they are furniture that delimits, and the colour budget is spent on the
%     hook, the KL arrow and the in-slab component.
%  8. THE KL ARROW IS THE ONLY STROKE ABOVE 1.0pt IN THE FIGURE (1.1pt).
%     Nothing else exceeds 1.0pt, so weight alone says which mark is the claim.
%     It runs between the two brace LABELS' baselines, so it visibly joins two
%     readings of the same tokens rather than two sentences.
%     REPAIRED, PASS 9: it was DOUBLE-headed and is now single-headed, tail at
%     the pass-one readout and head at the pass-two readout, matching the
%     direction its own label carries. KL(clean || ablated) is asymmetric; a
%     two-way head says "the difference between", i.e. that the two passes are
%     interchangeable, which is the one thing the measure is not. It was also
%     the figure's heaviest stroke saying the opposite of its label. No
%     double-headed arrow exists anywhere in the corpus: arrows.meta is used
%     only as -{Latex} for a displacement and -{Stealth} for a pipeline edge.
%  9. THE VECTOR PANEL DRAWS THE SLAB AS A BAND WITH THICKNESS, not as a line.
%     A line would say "one direction"; the slab is ten. Nothing in that panel
%     is to scale (a true 10-of-2048 slab is 0.5% of the squared norm and
%     cannot hold a labelled decomposition -- figs/slab.tex draws that true
%     angle, once, and this panel does not pretend to repeat it).
%     REPAIRED, PASS 9, AND THIS WAS THE PANEL'S WORST FAULT: THE RATIO WAS
%     THE WRONG WAY ROUND. h was drawn at 30 degrees above the band, so the
%     REMOVED component (hV)V^T measured 2.90cm and the SURVIVING remainder
%     h-(hV)V^T measured 1.70cm -- on the render, 228px against 130px, the
%     ten-dimensional slab drawn 1.75x LONGER than the 2038-dimensional
%     remainder, under a subtitle reading "ten of the 2048 directions". Cold,
%     the slab read as the dominant subspace and the residual as a leftover,
%     the exact inversion of the section's finding. h is now steep: component
%     0.48cm, complement 1.90cm, 4.0:1 the other way. It is deliberately NOT
%     the true 14.3:1 (86 degrees, cos^2(86) = 0.0049 = 10/2048): at that angle
%     the component is 0.13cm, shorter than its own arrowhead. slab.tex inset
%     one draws that geometry true and this panel defers to it in the caption.
%     Three smaller repairs in the same pass. (a) (hV)V^T was the only one of
%     the three vectors drawn HEADLESS, so in a decomposition the component
%     that is removed was the one not drawn as a vector and read as a baseline;
%     it now carries the same Latex head as its siblings, and its label carries
%     its fate, "-> 0", so the zeroing is a mark and not only the sentence
%     underneath. (b) Its label's 0.1cm rulegrey stub leader is deleted: it was
%     shorter than the 0.3cm threshold that triggers a leader and in the wrong
%     colour. (c) The construction is centred on the band's midpoint at 3.40,
%     because at 2.30 it crowded into the band's left third with 3.1cm of empty
%     band to its right.
%  9b.THE BAND IS panelbg WITH A rulegrey OUTLINE, not accent!15.
%     PASS 9, a cross-figure fix. Across this slate a PALE accent tint meant
%     "raw, still confounded" in fig:purify and "purified" here, while SOLID
%     accent meant V_pure in fig:purify and in the published slab.tex but "the
%     component being removed" here: the tint semantics were inverted between
%     two figures a few pages apart. And inside this figure alone, accent!12
%     (the drug field, both strips) and accent!15 (this band) sampled as
%     (225,234,242) and (218,229,238) -- a 3% tint difference that will be
%     indistinguishable on press, two meanings in one apparent colour. panelbg
%     with a rulegrey outline is slab.tex's own treatment for "a region of the
%     stream" (its \slabdisc, slab.tex:83-90), it leaves accent!12 uniquely
%     "the drug field" across licensing and hook, and it keeps solid accent for
%     the slab itself across the slate.
%     ONE RESIDUAL DIFFERENCE, AND IT IS DELIBERATE: fig:purify draws
%     V_pure as a SOLID accent bar and this figure draws it as a cream region.
%     The two are doing different jobs. In fig:purify V_pure is one subspace
%     among three and is the object every causal claim is made on, so it takes
%     the accent, exactly as slab.tex's published wedge does; here it is the
%     backdrop a vector is decomposed against, and an accent component has to
%     be legible ON it, which a solid accent fill would make impossible. The
%     tie between the two figures is the direct label, which is accent
%     $V_{\textrm{pure}}$ in both. Giving this band an accent OUTLINE instead
%     was considered and rejected: at panelbg fill with an accent 0.35pt border
%     it is the drug box's own treatment (accent!12 fill, accent 0.4pt border)
%     to within a tint, and one glyph would then mean two things again.
%  9c.ONE NAME FOR ONE OBJECT. The panel said V, the arrows said (hV)V^T and
%     the band said V_pure -- three names for the same ten directions, in a
%     figure that sits a page from fig:purify where V_pure is built. The
%     subtitle now states the identity once: "V = V_pure: ten of the 2048
%     directions".
% 10. THE RIGHT-ANGLE MARK AT THE ORIGIN is load-bearing, not decoration. It is
%     what distinguishes "the slab component is set to ZERO" from "the slab
%     component is reduced". Without it the vertical arrow is just a shorter
%     arrow.
% 11. THE READOUT AXIS IS LINEAR AND CARRIES NO BARS. It is one layer, so the
%     reason mechanism.py forbids a bar on a log axis (a bar's length is set by
%     wherever the axis is cut) does not arise -- and there are no bars in any
%     case, only three dots. Rejected alternative: seven layers on this axis,
%     which would rebuild fig:mechanism panel (b) inside a mechanism figure and
%     make the readout compete with the mechanism for the reader.
%     REPAIRED, PASS 9: EVERY DOT NOW CARRIES ITS VALUE. None of 1.563, 0.084
%     and 4.180 was printed on the drawing, and \kx{0.084} = 8.362 against a
%     zero tick at 8.20, so the null sat 1.6mm from zero and, under a 2.1pt
%     halo and a 1.4pt fill, was visually AT zero -- while the strip below
%     asserted "18.6x the null". A reader who takes the null for zero takes
%     18.6x for a claim they cannot check. Both reviewers proposed the same two
%     remedies, printing the values or moving to a log axis with a joining
%     stem; the first is taken, because it is the document's own no-legend
%     rule, it makes the ratio checkable from the marks, and it keeps the axis
%     comparable in kind with tab:workspace's nats column. The three
%     descriptions were also three near-full-width stacked lines in an order
%     that was neither the axis order nor a reading order -- the top line named
%     the rightmost dot, and the middle line, spanning crop x 275-1630, sat
%     directly over the accent dot at 770, so "the null" attached to the drug
%     slab on first pass. They are now short and value-first, each near its own
%     dot; the long descriptions live in the caption, which already carried
%     them. The three dots already sat over 2.1pt white halos in ladder.tex's
%     order (spine, halos, fills, labels) and still do; a reviewer reported
%     them as halo-less and that finding was overruled on the source.
% 12. NO MARK JOINS THE KL AXIS TO ANY ACCURACY OR SUBSPACE-FRACTION SCALE, and
%     no double-headed arrow spans two rules anywhere below the divider. Those
%     are different instruments with different comparators; fig:comparators is
%     the figure that owns that argument.
% 13. LAYER 4 IS THE ONE LAYER DRAWN because it is the slab's strongest showing
%     of the seven against the null (18.6x, the top of the thesis's 3.6-18.6x)
%     and the top of the 9-37% band against cell line. Drawing the best case
%     and still showing the accent dot sitting at a third of the ink dot is the
%     honest form of the claim. Layer 12 was the alternative and is not drawn:
%     there the drug and cell-line costs are nearly equal (0.133 against 0.135,
%     98%), which is the section's own named exception and would read as the
%     rule if it were the only layer shown. The caption says both, and points
%     at tab:workspace for all seven. hook.json records layer 12's three values.
% 14. THE TOP-RIGHT NOTE SAYS "AT ONE LAYER". The word "every" in "at every
%     position" ranges over the SEQUENCE, not over depth, and a reader who
%     meets "layer 4" only in the readout panel has no way to tell which. The
%     misreading -- that the slab is projected out at every layer as well as at
%     every position -- is exactly the kind this figure exists to prevent, and
%     no mark can carry it, so it is stated in words. It is a design fact from
%     the method (04b:164-171, one measurement per layer, one row per layer in
%     tab:workspace; 04b:191-195, the hook itself), not a stored quantity, and
%     hook.json records it as such. Set quietink at \scriptsize, because it is
%     furniture; that it also fills the figure's one empty quadrant is a
%     consequence and not the reason.
%
% WHAT SIX BUILDS CHANGED, with the measurement that forced each
% -------------------------------------------------------------
% (a) EVERY TEXT NODE IS anchor=base, so each y below IS a baseline. Build 1
%     used anchor=north west with an assumed 0.34cm node height. The real
%     height of a \footnotesize node at TikZ's default inner sep (0.3333em,
%     2.7pt a side) is 12 + 5.4 = 17.4pt = 0.61cm, nearly twice that, and the
%     consequences were visible: "the true sentence is teacher-forced, never
%     generated" put its descenders through the pass-one box tops, and the hook
%     rule at y = 7.52 ran clean through the middle of "the same prompt, the
%     same sentence, the same weights". Baseline anchoring removes the guess:
%     ink sits about 0.20cm above and 0.06cm below a \scriptsize baseline, and
%     0.25/0.08 at \footnotesize. frames.tex records the same fix for the same
%     reason, plus a second one -- a node with align= or text width builds a
%     minipage whose first line follows the AMBIENT \baselineskip, so its ink
%     moves with the prose before the float. No node here uses either.
% (b) THE PROMPT/TARGET DIVIDER MOVED OFF THE BOX EDGE. Build 1 put it at
%     x = 7.80, which is exactly the right-hand border of the "control cell
%     sentence" box: a 0.3pt rulegrey line drawn on top of a 0.35pt rulegrey
%     box border is invisible as a separate mark, and the requirement is that
%     the braces stop VISIBLY at it. It now sits at x = 8.40, the midpoint of
%     the 1.20cm gap between the strips, where nothing else is drawn, and it is
%     set as `ink, 0.5pt, densely dotted' -- this document's mark for a mere
%     coordinate worth naming, which is exactly what the boundary is. Both
%     braces now begin at 8.40 and their left tips land on it.
% (c) THE BRACE AMPLITUDE WENT FROM 2.5pt TO 3.5pt AT 0.35pt. At 2.5pt and
%     0.3pt, under a rounded-corner box border 0.10cm above it, the brace read
%     as a second box outline rather than as a brace; slab.tex and ladder.tex
%     both use 3.0-3.5pt for braces that have to be seen.
% (d) THE THREE READOUT LABELS WERE STAGGERED BY 0.50cm, not 0.33cm. At 0.33cm
%     "the purified drug slab" and "a matched-dimension random subspace --- the
%     null" overlapped by about 0.08cm of ink, because both labels are wide
%     enough to share x and the node boxes are taller than a line. 0.50cm gives
%     0.24cm of clear space between successive label blocks, which is also what
%     lets the two leaders leave their dots at distinguishable angles.
% (e) "ten of the 2048 directions" AND THE VECTOR PANEL'S CLOSING SENTENCE
%     overlapped outright in build 1 (the first sat at y 1.05-1.35, the second
%     at 0.95-1.15). The label was DELETED rather than moved: the panel's own
%     subtitle already reads "$V$ holds ten of the 2048 directions", so the
%     node under the band was saying a second time what the subtitle says
%     first, and deleting it also freed the space the closing pair needed.
% (f) THE STRIP WAS WIDENED FROM 13.20 TO 13.55 and the KL arrow moved from
%     13.75 to 14.05. At the narrower width the top-right quadrant, roughly
%     4.1 x 1.7cm, held nothing at all, and the KL label was the only mark
%     right of the strip. The right column now carries the "one layer" note as
%     well (decision 14) and band one reaches x = 16.5, which is where band
%     two's readout already ended.
% (g) THE IN-SLAB COMPONENT'S LABEL GOT A LEADER. $(hV)V^{\top}$ names a
%     segment that lies INSIDE the slab band, so the nearest place its label
%     can sit is 0.28cm above the band's top rule and 0.44cm above the mark
%     itself -- well past the ~0.3cm at which this document's rule requires a
%     leader. Built without one it read as floating in the interior of the
%     triangle. It now has a 0.14cm accent!45 hairline from just above the
%     segment to just under the label, which is the same device comparators.tex
%     uses for its 0.121 mark.
% (h) THE 18.6x / 37% LINE MOVED FROM quietink TO ink. Under the axis it was
%     the third of four grey centred lines and read as a fourth piece of axis
%     apparatus. It is a reading of the three marks, so it takes the register
%     this document gives an estimate rather than the one it gives furniture.
% (i) THE AXIS SPINE WAS CUT FROM 16.90 TO 16.60. The map's range top is 4.5
%     nats, but the largest mark is 4.180 and the last tick is 4, so a full-
%     range spine left 0.97cm of empty rule past the last tick and read as an
%     axis that had run out rather than as headroom. No coordinate moved.
% (j) THE PASS-ONE BOUNDARY RULE GREW A TAIL, from 8.72 down to 8.55. Its
%     label sits to its left on the brace label's baseline, and with the rule
%     stopping 0.14cm under the boxes the label read as a caption of the prompt
%     half rather than as a pointer at the boundary. The pass-two copy keeps
%     the symmetric 0.14/0.14, because it carries no label.
%
% LAYOUT NOTES, all forced by a build
% -----------------------------------
%   * fullwidth = textwidth 142mm + marginparwidth 28mm + marginparsep 4mm
%     = 174mm. (preamble.tex's own comment beside \newenvironment{fullwidth}
%     says 171mm and is stale arithmetic; the geometry keys give 174mm.) The
%     box is PINNED at 17.30 x 10.00cm = 173.0 x 100.0mm, 1.0mm inside the
%     measure. Rightmost ink is the axis spine at 16.90.
%   * the spec's token x-edges (0.00/1.20/2.20/3.00/4.30/7.30) do not hold
%     their labels at \scriptsize: "cell line" measures about 1.27cm of ink and
%     the box was 1.20cm, and "mechanism" the same against 1.30cm. The five
%     boxes were rewidened to 1.55/0.85/0.85/1.70/3.25, keeping the order and
%     keeping "control cell sentence" the widest field, and the prompt strip
%     now ends at 8.20 rather than 7.30. Everything keyed to the boundary (both
%     dividers, both braces, the target strip, the KL arrow) moved with it.
%     The five widths are set against figs/licensing.tex's measured ink for the
%     SAME five labels at the same size (cell line 0.954, drug 0.616 bold, dose
%     0.579, mechanism 1.424, control cell sentence 2.511; total 6.084cm), so
%     the clear space each side runs 0.117 to 0.370cm and no cell is near the
%     0.06cm floor. A first pass gave mechanism only 1.55cm, i.e. 0.063 a side,
%     while control cell sentence had 0.445; 0.15cm moved from the last cell to
%     the fourth evens that out. \ENDCELL's cell keeps 0.22 a side.
%   * the target strip is four boxes, so the strip is nine boxes and the hook
%     carries nine ticks. \ENDCELL is \texttt{[END\_CELL]}: ten inconsolata
%     characters, about 1.41cm at \scriptsize, so its box is 1.85cm.
%   * the prompt/target GAP is 1.00cm (8.20 to 9.20) and the boundary rule sits
%     at its midpoint, 8.70. Both braces begin there and their left tips land
%     on the rule, which is what "the brace stops dead at the divider" means
%     when it is drawn rather than described.
%   * nothing is scaled. 1cm here is 1cm on the page; there is no scale= key
%     and no \resizebox. style_v2.py records that v1's figures were silently
%     scaled to 0.887 by LaTeX and that this, not the drawing, read as "off".
%   * the pgfmath trap: \kx expands to a BRACED group, so \kx{1.563}+0.05 is
%     not a coordinate. Where a label needs a shim it is applied as an xshift
%     on the node, or the shim lives inside the braces. (gate.tex's header
%     records the same trap costing a compile.)
%   * \foreach prints its raw token, so the tick labels are given in the
%     \v/\t value/label pair form. Here every tick is a non-negative integer so
%     the pair form is not strictly needed, but it is the corpus idiom and it
%     keeps the axis editable without a surprise.
%
% TYPE
% ----
%   Ambient font=\small on the picture; every node sets its own size. The
%   dominant size is \scriptsize (8pt), with \footnotesize (10pt) only for the
%   four panel subtitles, the band-one sentence and the vector panel's two
%   closing lines. Nothing is set at body size. Registers: \sffamily bold
%   accent = a panel tag; roman accent = the intervention; roman ink = the
%   object being scored or the comparator; roman quietink = furniture and
%   annotation; \ttfamily = only \ENDCELL, which names a real token.
%
% ===========================================================================
% PASS 10 -- 2026-08-18 -- THE STYLE-CONTRACT REBUILD
% ===========================================================================
% EVERYTHING ABOVE THIS LINE IS LEFT INTACT and is the design record of passes
% 1-9. This section supersedes, BY NAME, only the claims above that this pass
% made false; a decision above that is not named here still stands, and the
% reasoning that produced it is still the reasoning this figure is audited
% against. NO NUMBER CHANGED IN THIS PASS: 1.563, 0.084, 4.180, n = 60, rank
% 10, residual width 2048 and layer 4 are byte-identical to pass 9, and the two
% ratios 18.6x and 37% were not altered but MOVED -- off the drawing and into
% the caption, which is where this document's contract puts a result.
%
% WHY THIS PASS EXISTS
% --------------------
% The supervisor's verdict on the 4.7-4.12 slate: the figures look "sloppy and
% childish", they must be CORRECT, and every axis must define what it measures.
% A binding style contract (S1-S25) and a cross-figure code (CF1-CF14) were then
% written against measurements of this document's own two reference figures --
% 4.11, figs/mechanism.py, the BODY standard, and A.3, figs/slab.tex, the
% appendix plate -- and the ruling was: in the body, take 4.11. Measured against
% it, this figure was the worst offender of the five.
%
%                          pass 9            4.11        pass 10 (measured)
%   type sizes >7pt        9                 2           2 chosen + 1 artefact
%   embedded faces         9                 1           6, four of them math/tt
%   italic characters      34%               0%          1.8%
%   stroke weights         9 (0.17, 0.68     --          5
%                          used once each)
%   all-caps sans banners  6                 0           0
%   ink at 200dpi          8.55%             5.3%        5.98%   (cap 6.0)
%   drawn box              173.1 x 97.0mm    --          173.99 x 99.95mm
%
% THE CORRECTNESS DEFECT, WHICH IS WHY THE FIGURE WAS REOPENED AT ALL
% -------------------------------------------------------------------
% THE READOUT AXIS IS NOW LOGARITHMIC. Design decision 11 above chose a linear
% nats axis, and pass 9 mitigated its consequence by printing every dot's value.
% BOTH ARE SUPERSEDED. The mitigation was not enough, and the header itself said
% so in as many words: on the linear map \kx = 8.20 + v/4.5*8.70 the null at
% 0.084 landed 0.162cm -- 1.6mm -- from the zero tick, under a 2.1pt halo and a
% 1.4pt fill, so the null dot WAS the origin at print size while the strip below
% asserted "18.6x the null". A ratio the drawing renders as infinite is not a
% mark a reader can check, whatever is printed beside it. Both quantities this
% figure quotes are RATIOS; on a log axis a horizontal gap IS a ratio, so the
% mark now means what the sentence means. 4.11 panel (b) published exactly this
% solution for exactly these numbers -- log axis, open and filled dots, the
% ratio read as a gap -- and this figure now uses it.
%   NEW MAP  \kx#1 = 8.40 + (ln(#1) - ln(0.05))/(ln(6.0) - ln(0.05)) * 8.98
%            domain [0.05, 6.0] nats, spine 8.40 to 17.38, labelled ticks at
%            0.1, 0.3, 1, 3
%   0.084 -> 9.3731 , 1.563 -> 14.8569 , 4.180 -> 16.7020
%   The null-to-slab gap is 5.4838cm and IS 18.6x. The slab-to-cell-line gap is
%   1.8452cm and IS 37% read the other way. Neither is annotated: the caption
%   carries both, which is where the contract puts a result.
%   The domain is deliberately wider than the data at BOTH ends, so no mark sits
%   on the cut and no bar is possible; 4.11(b) does the same, its ticks running
%   0.01-1.00 inside a taller axis. There are still no bars and there must not
%   be: on a log axis a bar's length is set by wherever the axis is cut.
%   SUPERSEDED: the linear map (0.084 -> 8.3624, 1.563 -> 11.2218,
%   4.180 -> 16.2813); decision 11's argument that "one layer, so a log axis is
%   not needed"; and note (i)'s spine crop, which no longer applies.
%
% THE FOUR MARKS AND LINES THAT WERE DELETED, EACH WITH ITS REASON
% ----------------------------------------------------------------
% (i)   THE TWO SHORT HOOK ARROWS (pass 9's replacement for the nine-tick comb)
%       are gone. They implied a count of two and the caption had to deny it in
%       a sentence of its own ("The two arrows off the rule mark that it acts in
%       both halves and imply no count of positions"). A drawing that needs its
%       caption to unsay it is drawing the wrong mark. The RULE's span is the
%       "everywhere" claim -- decision 5 above is right and survives -- and the
%       accent formula now sits ON the rule, 0.20cm above it, instead of 1.1cm
%       below the strip. This also empties the figure of every arrowhead except
%       the comparison arrow's two and the vector panel's three, which is what
%       S11 asks for: one arrowhead, one meaning.
% (ii)  THE ITALIC EPIGRAM "nothing is generated; the two passes differ only in
%       the hook" and the 17.30cm rulegrey divider that existed only to carry it
%       are deleted. It was 9.96pt Pagella-Italic quietink centred on the
%       figure's own x-centre -- the LARGEST type on the drawing, 25% larger
%       than the 7.97pt at which every measured number was set -- and its first
%       clause restates the caption's opening words, "Nothing is generated:".
%       It is the caption's sentence and the caption already had it.
% (iii) THE PRINTED RESULT "18.6x the null; 37% of what removing cell line
%       costs" is deleted from under the axis. A result may not sit in the
%       axis-label lane, and it was the third of four stacked centred lines
%       there. NOTE (h) ABOVE IS SUPERSEDED: moving that line from quietink to
%       ink so it would stop reading as axis apparatus was the wrong repair --
%       the line does not belong under the axis in any colour. Both ratios are
%       in the caption, which now also carries the qualifier the old caption
%       dropped: "9-37%" is the range over layers 2-9 ONLY; the true seven-layer
%       range of kl.pure_drug[0]/kl.cell_line[0] is 8.90% (layer 9) to 98.45%
%       (layer 12).
% (iv)  "the marks carry no interval; the caption says why" is deleted. It
%       restated the caption's own bold "no interval is drawn" and it was
%       italic. It is replaced by the single figure-foot declaration S17
%       requires, set \scriptsize quietink on the figure's x origin, which NAMES
%       the construction instead of pointing at the caption.
% (v)   "generation would begin here", the prompt/target divider's label, is
%       deleted, and this was the one deletion made for the ink budget rather
%       than for a defect. Recorded here because it was load-bearing and a
%       future pass may want it back. Its claim is drawn three other ways --
%       panel (a)'s head is "pass one, clean", both braces read "scored tokens"
%       so the prompt half is visibly not scored, and the caption's lede opens
%       "Nothing is generated" -- and the divider itself survives as what it
%       always was, a piece of the strip's structure. Cost measured: 27mm2 of
%       ink, against a 6.0% ceiling the figure clears at 5.98%.
%
% WHAT REPLACED THE SIX BANNERS AND THE FOUR ITALIC SUBTITLES
% -----------------------------------------------------------
% Every \sffamily\scriptsize\bfseries accent banner is gone, and with it the
% only sans face in the figure. The apparatus face belongs to the margin column,
% not to a figure read inside a running argument; A.3 may use it because A.3 is
% a reference plate a reader stops at. Panel heads are now 4.11's form exactly
% -- a lowercase parenthesised letter and a sentence-case descriptive title,
% \footnotesize roman ink, on the panel's own x origin, no colon, no bold, no
% caps:
%     (a) pass one, clean
%     (b) pass two, the slab projected out
%     (c) what the hook does to one vector
%     (d) what comes out
% The four \footnotesize\itshape quietink subtitles under them are deleted.
% Their content was distributed, not lost:
%   "the true sentence is teacher-forced, never generated"  -> already the
%       caption's opening clause; it was duplicating it (S18).
%   "the same prompt, the same sentence, the same weights"  -> now the drawn
%       tie between the strips (below).
%   "the hook is on ONE layer's residual stream"            -> now the axis
%       block's scope line, where the layer is named anyway.
%   "V = V_pure: ten of the 2048 directions"                -> the identity is
%       now the band's own direct label; the rank moved to the caption.
%   "layer 4: mean KL from the clean pass, over the target tokens, 60 prompts"
%       -> now the axis block itself, split into a quantity line and a scope
%       line, which is the position S5 requires.
% DECISION 14 ABOVE IS PARTLY SUPERSEDED: the "at one layer" fact is still
% stated in words, because no mark can carry it, but it is no longer a clause of
% PASS TWO's subtitle -- there is no subtitle -- and lives in the axis block's
% second line, the place S5 reserves for scope.
%
% THE TIE BETWEEN THE STRIPS, WHICH IS THE ONLY NEW TEXT IN THE FIGURE
% ---------------------------------------------------------------------
% The two strips are emitted by one macro and are therefore provably identical
% box for box -- the figure's central claim, and its best-executed feature.
% Nothing on the drawing said so. Worse, the two braces beneath them carried
% DIFFERENT labels ("per-token log p, read here only" against "the same tokens,
% read again"), which sends a reader looking for a difference between two
% objects whose sameness is the point. Both braces now carry the SAME label,
% "scored tokens", and one \scriptsize roman ink line reads "the same sentence,
% box for box". It sits UNDER both strips, on band one's own x origin and on the
% same baseline as pass two's brace label, so it is a statement about a pair the
% reader has just seen and it costs the drawing no extra row of ink.
%
% THE COMPARISON IS NO LONGER AN ARROW WITH ONE HEAD
% ---------------------------------------------------
% DECISION 8 ABOVE AND ITS PASS-9 REPAIR ARE SUPERSEDED. Pass 9 made the KL mark
% single-headed, reasoning that KL(clean || ablated) is directed and asymmetric
% and that a two-way head would say the passes are interchangeable. That
% reasoning is about the MEASURE; the MARK is about the two passes, and a naive
% reader read the single head as pass one FEEDING pass two -- generation, the
% precise misreading this figure exists to kill, drawn by the figure itself in
% its heaviest stroke. The two passes are independent scorings of one fixed
% sentence by one set of weights: that is a comparison, and S11 gives a
% comparison a double head or a brace. It is now a DOUBLE-headed ink arrow at
% 1.0pt, and the single accent arrowhead is reserved for an operation acting in
% a direction. It is no longer the figure's heaviest stroke either -- the 1.1pt
% was a tenth-of-a-point singleton in a nine-weight vocabulary -- and its label
% is the registry name, "KL from the clean forward pass", the same words the
% axis in panel (d) carries, because it is the same quantity and S19 gives an
% object exactly one name.
%
% THE AXIS DEFINITION, WHICH IS THE SUPERVISOR'S FIRST COMPLAINT
% --------------------------------------------------------------
% This figure was the only one of the five that already named quantity and unit
% in axis-label position, and note (h) plus decision 11 built that. The WORDS
% survive; the position, the rank and the stack around them change.
%   LINE 1  \footnotesize roman ink, on the axis's own left end, x = 8.40:
%           "KL from the clean forward pass (nats), logarithmic scale"
%   LINE 2  \scriptsize quietink, same x:
%           "mean over the target tokens of 60 teacher-forced prompts, layer 4"
% "Mean KL divergence, clean -> ablated (nats)" is RETIRED. It was centred and
% bold under the rule; more importantly it was a SECOND name for 4.11(b)'s "KL
% from the clean forward pass", and CF3 gives the quantity one name across the
% slate. The old panel subtitle's "layer 4 ... over the target tokens, 60
% prompts" is folded into line 2, so the scope is stated once instead of twice.
%
% THE READOUT IS THREE ROWS, NOT ONE, AND WHY
% -------------------------------------------
% S22 requires that within a row every direct label share one baseline and that
% no label cross the drawn measure. Three labels naming three objects at one
% layer -- "0.084 - matched-dimension random slab" (5.2cm), "1.563 - the
% purified drug slab" (4.2cm) and "4.180 - cell line, the positive control"
% (5.4cm) -- cannot share one baseline: the last two marks sit 1.85cm apart on
% the log map. Staggering labels off one rule is what S22 forbids; giving each
% object its OWN row is not. Each row therefore carries exactly one mark and one
% label on one baseline; the three rows share one log x axis; and they are
% ordered by value with the largest on top, so vertical position carries the
% same order as horizontal position and can assert nothing the axis does not.
% That ordering is 4.11(b)'s own, the filled slab dot above the open null dot.
% The rows sit 0.38, 0.82 and 1.26cm above the spine so the three read as marks
% over one scale. Each label is 0.26cm from its own mark, just inside the 3mm at
% which this document requires a leader, so no leader is drawn. THE LABEL SIDE
% DIFFERS BY ROW AND IT HAS TO: the two large marks are 0.7 and 2.5cm from the
% right end of the measure and their labels can only go left, while the null at
% 9.37 has 0.97cm to its left and its label can only go right. NO DROP LINES
% JOIN THE MARKS TO THE SPINE: a hairline from a dot down to a horizontal axis
% reads as a bar, and a bar on a log axis is forbidden in this document.
%
% MARKS, WEIGHTS, TYPE AND PROSE
% -------------------------------
% MARKS, S9/CF5, three glyphs and no more:
%   filled disc, r = 1.4pt, accent   the purified drug slab -- the object under
%                                    test, and the only accent mark in panel (d)
%   filled disc, r = 1.4pt, ink      cell line, the positive control
%   open disc,  r = 1.4pt, white     matched-dimension random slab, the null.
%               fill, 0.8pt ink edge This is 4.11(b)'s own encoding for this
%                                    same object (mechanism.py: mfc="white",
%                                    mew=1.1, color=style.INK), and adopting it
%                                    is CF5's stated requirement. The 0.8pt edge
%                                    is the mark vocabulary's own weight, not a
%                                    path weight, and it is the only 0.8pt here.
%   The three 2.1pt white halos are DELETED together with the two leaders they
%   existed to cover. Pass 9's overruling of the "no halo" finding stands as a
%   fact about pass 9; the halos are simply no longer needed, because no leader
%   now runs into a dot.
% WEIGHTS, S10. Measured in the built PDF: 0.30 x26, 0.40 x5, 0.80 x1, 0.90 x1,
% 1.00 x4. Five weights, against pass 9's nine.
%   0.30pt  furniture: every token-box border, both braces, the slab band's
%           outline, the right-angle mark, the two dashed construction lines,
%           the four axis ticks
%   0.40pt  the drug cell's accent border, both prompt/target dividers, the
%           axis rule
%   0.80pt  the open disc's edge (S9's own weight for that glyph; see above)
%   0.90pt  the hook rule -- USED ONCE, and justified here as S10 expressly
%           permits: it is a distinct kind of object, the intervention, the one
%           thing that differs between the two passes. It is accent and not
%           black, so it cannot be confused with the chance/gate/margin register
%           S6 reserves for a black solid 0.9pt rule, and this figure draws no
%           threshold at all.
%   1.00pt  panel (c)'s three vectors and the comparison arrow
%   RETIRED: 0.35pt (box borders and braces, folded into 0.30), 0.50pt (the
%   dotted dividers, see below), 1.10pt (the KL arrow), and the 0.17pt and
%   0.68pt singletons, which were never chosen at all -- they were the Latex
%   arrow tips' own stroke, and the tips are now declared line width=0pt so they
%   are pure fills and contribute no stroke weight.
% THE PROMPT/TARGET DIVIDERS are no longer "ink, 0.5pt, densely dotted".
%   NOTE (b) ABOVE IS SUPERSEDED. S7 reserves the densely-dotted 0.5pt stroke
%   for a zero rule, a magnification leader or a projection line, and requires a
%   dotted rule's label to be "0", to be a leader, or to be absent; this one
%   carried "generation would begin here", which is none of those. The line is
%   what it always was -- a piece of the strip's structure that delimits two
%   halves -- so it is now rulegrey 0.40pt solid, a strip weight, still at 7.94
%   in the middle of the empty 1.00cm gap where nothing else is drawn. Note
%   (b)'s actual finding, that a rulegrey line drawn ON a rulegrey box border is
%   invisible as a separate mark, is why it stays at mid-gap.
% TYPE, S1. Measured in the built PDF: 7.97pt x587 characters and 9.96pt x165,
% plus three exempt classes -- 6.85pt x20 (\ENDCELL, Inconsolata, the \texttt
% token name S1 exempts), 5.98pt x8 and 6.23pt x3 (subscripts and the \top
% superscript, which S1 exempts), and 8.30pt x10.
%   THE 8.30pt SPANS ARE NOT A THIRD CHOSEN SIZE and are recorded here so an
%   auditor need not rediscover them: they are ten glyphs -- \leftarrow, two
%   minus signs, six parentheses and \to -- set from CMSY10 and CMR10 inside the
%   two formulas $h \leftarrow h-(hV)V^{\top}$ and $(hV)V^{\top} \to 0$.
%   mathpazo substitutes Palatino for math letters but leaves the CM symbol
%   fonts in place and scales them by 1.041 to match Palatino's x-height, so a
%   \scriptsize (7.97pt) math operator prints at 7.97 x 1.041 = 8.30pt. It is
%   the same nominal size in a differently scaled font, it is a property of this
%   document's preamble and not of this figure, and every figure in the thesis
%   that sets inline math shows it. The one instance that COULD be removed was:
%   the ellipsis token box was $\cdots$ (CMSY10 at 8.30) and is now \dots, set
%   from Pagella at 7.97 like every other token label.
%   Registers: \footnotesize roman ink = the four panel heads and the axis's
%   quantity line, and nothing else, so the largest type on the drawing always
%   names a panel or a measured quantity and never narrates. \scriptsize carries
%   every label, tick, formula, tie and foot line. Zero \sffamily, zero
%   \bfseries, zero \itshape. Italic characters: 14 of 793, 1.8%, and all
%   fourteen are the math variables h and V, which are a symbol's own face and
%   not a voice. Pass 9 measured 34%.
% PROSE, S4: THE FIGURE NOW CARRIES NONE, which is 4.11's own austerity -- 4.11
%   has axis labels and direct labels and nothing else. Pass 9's two closing
%   lines, "the slab component is set to zero; every other direction is
%   untouched", are deleted because the drawing already carries both claims as
%   MARKS: the accent component's label ends "-> 0", the right angle says the
%   component is zeroed and not merely shortened, and the surviving vector is
%   labelled $h-(hV)V^{\top}$. A sentence whose deletion leaves no drawn object
%   unnamed is not definitional and S4 does not permit it. Decision 9c's "one
%   name for one object" survives, carried now by the band's own direct label,
%   "$V_{\textrm{pure}}$, the purified drug slab", which is also the gloss S19
%   requires at the symbol's first use.
%
% THE RIGHT ANGLE, AND THE ONE GEOMETRIC RATIO THAT MOVED
% --------------------------------------------------------
% The right-angle mark at the foot of the surviving vector was 0.16cm (1.6mm)
% and a naive reader did not see it until zooming. It is now 0.40cm (4.0mm) on
% both legs. Decision 10 above is why it may not simply be dropped: it is what
% distinguishes "the slab component is set to ZERO" from "the slab component is
% reduced". A 4mm glyph cannot sit on a 0.48cm arm without swallowing it, so the
% in-slab component grew from 0.48cm to 0.90cm and the surviving complement from
% 1.90cm to 3.14cm. THE RATIO IS UNCHANGED IN DIRECTION AND BARELY CHANGED IN
% SIZE: 3.49:1 in favour of the survivor against pass 9's 4.0:1, still the right
% way round, and still deliberately NOT the true 14.3:1 (86 degrees,
% cos^2 86 = 0.0049 = 10/2048), at which the component would be 0.13cm, shorter
% than its own arrowhead. figs/slab.tex inset one draws that geometry true, and
% the caption still sends a reader there. Nothing in this panel is a coordinate
% and nothing in it is measured.
%
% THE STRIP: SAME MACRO, SAME FIELDS, EVEN CLEAR SPACE
% -----------------------------------------------------
% The nine box edges were reset so that EVERY cell carries the same 0.13cm of
% clear space each side of its label's measured ink, against pass 9's 0.117 to
% 0.370cm. The strip therefore runs 0.00 to 12.65 instead of 0.00 to 13.55, the
% prompt half ends at 7.44 and the target half begins at 8.44, and everything
% keyed to the boundary (both dividers at \bnd = 7.94, both braces, the brace
% labels at 10.545, the hook rule, the comparison arrow at 13.40) moved with it.
% Two reasons, and both matter: the clear space is now a constant rather than a
% number that varied by a factor of three between cells, and the strip's own
% rules are 3.6cm shorter, which is 5mm2 of the 6.0% ink ceiling. The five
% prompt fields, their order, the [END_CELL] terminator and the ONE-MACRO
% construction are untouched. Ink widths the edges are set from, measured in
% figs/licensing.tex at the same size: cell line 0.954, drug 0.616, dose 0.579,
% mechanism 1.424, control cell sentence 2.511, [END_CELL] 1.41.
%
% GEOMETRY, MEASURED IN THE 200dpi RENDER
% ----------------------------------------
% Content box 173.99 x 99.95mm, against a 170-174mm fullwidth measure and a
% 100mm height cap. Ink 5.98% of the content box, against a 6.0% ceiling and
% pass 9's 8.55%; 793 non-space characters over 17393mm2 is 4.6 per 100mm2,
% against a 5.0 ceiling. Picture bbox 17.38 x 10.04cm (y 0.58 to 10.62).
% Leftmost ink is x = 0.00 -- the four panel heads, both strips, the hook rule
% and its formula, the tie, panel (c)'s band and label, and the foot line, all
% on it. Rightmost ink is the axis spine at 17.38. PANEL X ORIGINS, S13: band
% one and panel (c) at 0.00, panel (d) at 8.40. Measured in the built PDF, the
% left edges of the panel-level text lines are 58.34, 58.34, 58.48, 58.34,
% 58.34, 58.36, 58.34pt for the 0.00 panels -- a spread of 0.14pt -- and
% 296.45pt three times over for panel (d), a spread of 0.00pt.
%
% ===========================================================================
% PASS 11 -- 2026-08-18 -- THE ADVERSARIAL-REFUTATION REPAIR
% ===========================================================================
% EVERYTHING ABOVE THIS LINE IS LEFT INTACT. Three independent refuters read the
% pass-10 build and returned 33 findings. This section supersedes BY NAME only
% the claims above that this pass made false, records every finding APPLIED and
% every finding OVERRULED with its reason, and re-records the measurements.
% NO NUMBER CHANGED IN THIS PASS. 1.563, 0.084, 4.180, n = 60, rank 10, residual
% width 2048 and layer 4 are byte-identical to passes 9 and 10, the log map
% \kx is byte-identical to pass 10, and the three marks stand at the same three
% x to within the 0.02pt pgfmath residual recorded below. Nothing measured moved.
%
% TWO CLAIMS ABOVE ARE FALSE AND ARE SUPERSEDED HERE
% ---------------------------------------------------
% (S-1) The pass-10 note at the foot of this file said the replacement caption
%       "is also staged, keyed by \label, in figs/_PENDING_CAPTIONS.tex for a
%       human to integrate". THAT FILE DID NOT EXIST. It does now, this pass
%       created it, and it carries the caption plus an explicit list of the
%       clauses the integrator must DELETE from the live caption -- above all
%       Sections/04b-representation.tex:289-290, "The two arrows off the rule
%       mark that it acts in both halves and imply no count of positions",
%       which names two marks pass 10 deleted. Until that deletion is made the
%       figure ships contradicting its own caption. hook.json's .shape carried
%       the same false claim and is corrected.
% (S-2) The pass 1-9 block's AGREEMENT WITH THE TEXT cites 04b-representation
%       lines 229, 254, 260 and 286-288, and note 5/decision 5 cites 193, 195
%       and 200-202. Four of those are stale: tab:workspace's layer-4 row is now
%       line 338 (229 is fig:purify's caption), "3.6--18.6x" is 363, the
%       "9--37% ... at layers 2--9" sentence is 368-370, and "No intervals are
%       drawn in either panel" is 395. The 193/195/200-202 set is now 263, 264,
%       265 and 270-272 and pass 10 had already refreshed it in hook.json. The
%       "2048 is prose only at 04b:383" citation is also stale: 383 is
%       fig:mechanism's caption and the residual width is stated at 04b:648.
%       hook.json's citations are repaired; the stale text above is left in
%       place, superseded here, because S25 forbids rewriting a header.
%
% THE FIFTEEN REPAIRS, EACH WITH THE MEASUREMENT THAT FORCED IT
% --------------------------------------------------------------
% (1) PANEL (d), THE HEADLINE GAP WAS COVERED BY ITS OWN LABEL. This is the
%     serious one, because the log axis exists for one reason -- on a log scale
%     a horizontal gap IS a ratio -- and the gap was drawn to the pixel and then
%     papered over. The null's label ran RIGHT from its mark and occupied 100%
%     of the 5.48cm between the open null disc and the filled accent disc. It is
%     now two right-aligned lines running LEFT into the gutter, so the null's own
%     row is clear from its mark rightward and the ratio is a visible span.
% (2) PANEL (d), EVERY VALUE NOW ABUTS ITS OWN MARK. "4.180 --- cell line, the
%     positive control" put its numeral at the far end of the label from the dot
%     it names: measured in the pass-10 PDF, "4.180" began at x = 393.28, the
%     axis interval 0.32-0.43 nats, with its leading digit overlapping the "0.3"
%     tick numeral's column 12.5mm below it, while its own mark stood 48mm away.
%     On a value axis a numeral is read at the position it occupies (S21). Each
%     label is reversed to "name --- value"; no character is added.
% (3) PANEL (d), ALL THREE LABELS ARE NOW ON THE SAME SIDE OF THEIR MARKS. The
%     side flipped between rows two and three, so the disc trailed the text
%     twice and led it once, and at print a trailing 1.4pt disc reads as a
%     bullet ending a line -- three key entries, i.e. a legend, which is the one
%     construction this document forbids outright. All three labels now sit
%     left, all three discs right, and (1) is what made that possible.
% (4) PANEL (d), THE RULE NOW BEGINS AND ENDS ON A DRAWN TICK. At four labelled
%     ticks (0.1, 0.3, 1, 3) the rule overran its tick lane by 36.87pt (13.0mm)
%     at EACH end -- 29% of the axis carrying no numeral, which is the exact
%     defect S13 measures and condemns in fig:comparators row (d) -- and BOTH
%     extreme marks fell in those bare zones, so neither the null at 0.084 nor
%     the positive control at 4.180 could be placed against a printed number.
%     The domain's own ends are now ticks: 0.05, 0.1, 0.3, 1, 3, 6. The two end
%     numerals are set flush INSIDE the rule's ends rather than centred on their
%     ticks, because centred they overhung x = 17.38 by 0.09cm and made the
%     float overfull by 1.88pt; every other numeral is centred on its tick.
% (5) PANEL (d), THE TICK NUMERALS ARE INK, NOT quietink. quietink is this
%     document's comparator-and-null colour (S8/CF4). Spending it on the scale
%     while the figure's actual null was labelled in ink inverted the code.
%     fig:mechanism sets every tick numeral in ink. THE NULL'S OWN INK LABEL AND
%     ITS ink-EDGED OPEN DISC ARE DELIBERATE AND OVERRIDE CF4's quietink slot:
%     CF5 instructs this figure by name to adopt 4.11(b)'s encoding for this
%     object, and 4.11(b) draws it mfc="white", mew=1.1, color=style.INK. The
%     distinction between the null and the positive control is carried by the
%     GLYPH -- open against filled -- which is S9's own device, not by hue.
% (6) THE COMPARISON ARROW'S HEADS POINTED AT NOTHING. It ran between the two
%     brace LABELS' baselines, 9.04 to 6.88, so its upper head sat 0.50cm BELOW
%     the strip it was meant to span and both heads terminated in white space.
%     It now runs 9.75 to 7.59, the two strips' own vertical centres, at x =
%     13.10 rather than 13.40.
% (7) THE COMPARISON ARROW'S LABEL IS NOW ONE WORD, "compared". It had set
%     panel (d)'s axis quantity string a second time 130pt away, so a reader
%     could take this vertical 1.0pt rule for an axis OF that quantity; the axis
%     names the quantity once and this mark's job is to say that the two passes
%     are compared and not chained. S20 is satisfied -- the role is stated in a
%     verb of comparison beside the double head S11 gives a comparison. The 17
%     characters this frees are what pay for (8), (9) and (4) inside the 6.0%
%     ink ceiling, which the pass-10 build cleared by 2.8mm2.
% (8) THE TIE MOVED AND GAINED THE WEIGHTS. It sat at 6.88 -- the EXACT baseline
%     of pass two's brace label, measured to 0.00pt -- and directly under pass
%     two's prompt half, the one half of the strip that carries no brace, so it
%     read as that half's missing brace label, i.e. as a statement about ONE
%     strip when its whole job is to be a statement about TWO. It is now band
%     one's foot line on its own baseline at 6.44, 0.44cm below the brace label
%     and 0.76cm above band two's heads. And it read "the same sentence, box for
%     box", which carries only half of the question this figure must answer:
%     nothing anywhere on the drawing named the WEIGHTS, and the pass-10 note
%     claiming the deleted subtitle "the same prompt, the same sentence, the
%     same weights" had been inherited by the tie was wrong. It now reads "one
%     sentence, one set of weights, scored twice".
% (9) THE LAYER IS NAMED AT THE INTERVENTION. "at every position" ranges over
%     the SEQUENCE, not over depth, and the only "layer 4" on the drawing was in
%     the axis block 40mm away, attached to the READOUT. Panel (b)'s head now
%     reads "(b) pass two, the slab projected out at layer 4". DECISION 14 above
%     and pass 10's partial supersession of it are superseded again: the fact is
%     still stated in words because no mark can carry it, but it is stated where
%     the hook is drawn, and the axis block keeps its own "layer 4" because CF10
%     requires the layer named where the drawn layer's marks are.
% (10) THE HOOK RULE NO LONGER READS AS AN UNDERLINE. At y = 7.92 it sat 3.42pt
%     under an accent line of type and 3.40pt over the box tops -- equidistant,
%     same colour above and below, which is the heading-and-underline idiom. It
%     is at 7.88: 0.18cm of clear space above, 0.08cm below, so it binds to the
%     strip it acts on. It is NOT dropped onto the boxes, as one refuter
%     proposed: a mark running into a rule uninterrupted is what this file's own
%     note 5(b) forbids, and the drug cell's border is accent too.
% (11) PANEL (c), A PRINT COLLISION. The label $(hV)V^{\top} \to 0$ cleared the
%     top edge of the V_pure band by 0.72pt -- two pixels at 200dpi -- so its
%     parenthesis descenders would have closed on press. The band thins from
%     0.32 to 0.24cm and the label rises 0.02cm; the label is now 0.286cm from
%     its own mark, still inside the 3mm at which this document requires a
%     leader, and clears the band by 0.08cm.
% (12) PANEL (c), THE RIGHT ANGLE IS MIRRORED TO THE LEFT. On the right it was
%     unreadable: the 1.0pt black h vector crossed its 0.30pt horizontal leg
%     0.12cm from the corner, and its vertical leg terminated exactly on the
%     accent component's shaft. h rises to the RIGHT of the surviving vector, so
%     the corner between the survivor and the slab on the LEFT is empty. Same 90
%     degrees, same two arms, nothing crossing it. DECISION 10 above stands: the
%     mark is what distinguishes "set to ZERO" from "reduced".
% (13) PANEL (c), THE VERTICAL CONSTRUCTION GUIDE STOPS 0.20cm SHORT of h's tip.
%     It used to end inside the arrowhead. The HORIZONTAL guide still touches
%     both tips and must: it is the closing side of the parallelogram.
% (14) PANEL (c), THE BAND RUNS 0.00 TO 5.80 AND HAS LOST ITS OUTLINE. Two
%     findings, one repair. The lower band held a rectangle 38 x 46mm between
%     this panel and panel (d) with literally no ink in it -- 10% of the content
%     area and the largest void on the drawing -- and the band was byte-
%     identical in stroke and fill to the eighteen token cells above it
%     (rulegrey 0.30pt on panelbg), so ONE rectangle treatment carried two
%     meanings: "a field of the prompt", eighteen times, and "the purified drug
%     slab", once. Lengthening the band closes half the void at no ink cost
%     (panelbg renders above 230 grey and S15 does not count it); dropping the
%     outline makes the band the only fill-without-stroke region in the figure
%     and frees 8.7mm2. AN ACCENT OUTLINE WAS CONSIDERED AND REJECTED AGAIN, for
%     note 9b's reason: accent fill plus accent border is the drug cell's own
%     treatment to within a tint, and one glyph would mean two things again.
%     NOTE 9b IS THEREFORE AMENDED, NOT SUPERSEDED: the band is still panelbg
%     and still not accent; it simply no longer carries slab.tex's rulegrey
%     outline, because in THIS figure that outline is the token cell's own.
% (15) EVERY PANEL HAD TWO LEFT EDGES 1.20pt APART. TikZ's default 1pt inner
%     xsep meant a base-west node's ink began 1.20pt right of the coordinate it
%     was placed on, so the drawn rules started at one x and the panel's text at
%     another, visible under a flat-left glyph like the "K" of "KL". inner
%     xsep=0pt now. MEASURED IN THE BUILT PDF: the panel-level text lines of the
%     0.00 panels start at 57.3, 57.3, 57.5, 57.3, 57.3, 57.4, 57.3pt -- spread
%     0.2pt, and the residual is glyph side bearing, not node placement -- and
%     panel (d)'s head, its first tick numeral and both axis lines start at
%     295.5pt against an axis rule whose left end is 295.25pt.
%
% WHAT WAS PROPOSED AND OVERRULED, AND WHY
% -----------------------------------------
% (O-1) DELETE PANEL (c) AND RE-LETTER (d) AS (c). Refused. The spec's keep-list
%     protects the panel by name -- "h, h - (hV)V^T, the removed component going
%     to zero, and the deferral of the true 10/2048 angle to fig:slab" -- and
%     CF14 forbids losing load-bearing work. The finding's real content, the
%     void, is answered by (14) without deleting anything.
% (O-2) PUT ALL THREE MARKS ON ONE BASELINE AND STAGGER THE LABELS ABOVE WITH
%     LEADERS. Refused: staggering labels off one rule is exactly what S22
%     forbids, and the three names measure about 11cm against a 9cm axis, so one
%     baseline is not available. The three-row construction, one mark and one
%     label per row, is S22's own reading and it survives; (1), (2) and (3) fix
%     what was actually wrong with it.
% (O-3) DRAW DROP LINES OR PROJECTION LINES FROM THE MARKS TO THE SPINE.
%     Refused, and not for the reason pass 10 gave. S7 does sanction a densely
%     dotted 0.5pt quietink projection line, so pass 10's blanket objection was
%     too strong. The decisive objection is arithmetic: the three rows sit 0.38,
%     0.82 and 1.26cm above the spine in VALUE order, so three stems dropped to
%     the axis would have lengths that rise with the value they sit over, and
%     the panel would read as a bar chart of the three numbers against a
%     vertical scale that does not exist. A mark vocabulary must not encode a
%     quantity twice, once truly and once by accident.
% (O-4) ANNOTATE THE TWO GAPS WITH "18.6x" AND "37%" AS 4.11(b) DOES. Refused.
%     The spec deleted both printed ratios from this figure to the caption, and
%     a result is caption material here. 4.11(b) owns the annotated ratios --
%     CF1 says so -- and this figure now shows the gap and lets the caption
%     name it.
% (O-5) SET THE HOOK'S FORMULA IN ink AS DEFINITIONAL PROSE. Refused. It is the
%     direct label of an accent object, and S8 gives accent to "the drug-side
%     object under test AND ITS DIRECT LABELS". It is also \scriptsize, and S4's
%     prose is \footnotesize; calling it prose would break S4's size rule to fix
%     a colour it does not break. (10) fixes the underline reading instead.
% (O-6) DROP THE ACCENT FROM THE "drug" PROMPT CELL. Refused. It is a cross-
%     figure tie recorded in two files: decision 2 above, and figs/licensing.tex,
%     which ADOPTED this figure's cell treatment in pass 9 so that the two
%     strips read as one object seen twice. Removing it here would desynchronise
%     the pair. S8's own check lists what may not be accent -- a comparator, a
%     null, a discounted arm, a struck branch -- and the drug field is none of
%     them.
% (O-7) DRAW THE SECOND STORED RANDOM DRAW AS A FOURTH MARK. Refused as a
%     DRAWING change and ACCEPTED as a disclosure. workspace_probe.json really
%     does store layers.4.kl.random = 0.19796918605764707 beside the drawn
%     layers.4.kl.random_matched = 0.08388670384883881, and with
%     layers.4.raw_drug_dims = layers.4.pure_drug_dims = 10 the two are draws of
%     the same rank-10 construction, so the headline 18.6x would be 7.9x on the
%     other draw. That is a real and undisclosed exposure and the caption now
%     carries it. But the mark itself would be a fourth dot with a role no
%     section of the thesis states, drawn from a numeral that appears in no
%     table and in no sentence -- S20's own failure condition -- and tab:workspace
%     inherits the same single draw, so the figure would then disagree with the
%     table it is read against. hook.json records the second draw beside marks[1].
% (O-8) RENAME THE MARKS TO "ablating the purified drug slab" AND THE AXIS TO
%     "KL from the clean forward pass when a slab is ablated". Refused: S19's
%     registry and the spec both fix these strings, and line 1 is 4.11(b)'s
%     exact quantity name, which CF3 requires the slate to share. The verb the
%     finding wants is in the caption's lede, where a result belongs.
% (O-9) REVERSE OR BEHEAD THE (hV)V^T ARROW BECAUSE IT "POINTS AWAY FROM THE
%     ZERO IT GOES TO". Refused. It is a component of a vector decomposition and
%     points along the slab, which is its direction; "-> 0" is its fate, carried
%     by the label and by the right angle. Pass 9 changed it FROM headless TO
%     headed for a recorded reason -- the removed component was the only one of
%     three not drawn as a vector and read as a baseline -- and reversing the
%     head would draw a vector backwards.
%
% GEOMETRY AND STYLE, RE-MEASURED AT PASS 11
% -------------------------------------------
% Content box 173.99 x 99.95mm (fullwidth measure 170-174mm, height cap 100mm),
% unchanged from pass 10 to the hundredth of a millimetre. Ink 5.970% of the
% content box against a 6.0% ceiling, down from pass 10's 5.984% and pass 9's
% 8.55%. 709 non-space characters over 17390mm2 is 4.08 per 100mm2 against a 5.0
% ceiling. (The 793 recorded at pass 10 was counted by a different method; 709
% is the count of non-space characters in the built PDF's own text spans, and
% that method is now the one hook.json records.)
% TYPE, measured in the built PDF: 7.97pt x522 and 9.96pt x146 -- two chosen
% sizes -- plus the exempt classes 6.85 x20 (\ENDCELL, Inconsolata), 5.98 x8 and
% 6.23 x3 (sub- and superscripts) and 8.30 x10 (the CMSY10/CMR10 math symbols
% pass 10 explains and mathpazo causes). Faces: Pagella-Regular 653,
% PalladioL-Roma 9, Inconsolata 20, PalladioL-Ital 14, CMSY10 7, CMR10 6. Zero
% sans, zero bold, zero all-caps banners. Italic 14 of 709 = 2.0%, and all
% fourteen are the math variables h and V.
% STROKES, measured: 0.30 x27, 0.40 x5, 0.80 x1, 0.90 x1, 1.00 x4 -- the same
% five as pass 10. RECTANGLES: nineteen, and pass 10's count of them was wrong
% in its own arithmetic. Sixteen panelbg token cells, TWO accent!12 drug cells,
% eighteen cells in all, plus the V_pure band -- which is now the only one drawn
% without a stroke. ARROWHEADS: five, three on panel (c)'s vectors and two on
% the comparison arrow.
% THE AXIS, measured: spine x = 295.25 to 549.81pt, ticks at 295.25 (0.05),
% 332.10, 390.48, 454.51, 512.94 and 549.81 (6), i.e. the first and last drawn
% tick sit exactly on the rule's endpoints and there is no unnumbered overrun.
% Marks at 323.06 (null, between the 0.05 and 0.1 ticks), 478.26 (purified slab,
% between 1 and 3) and 530.58 (cell line, between 3 and 6). Each label ends
% 7.56-7.58pt from its own mark. The null-to-slab gap is 155.20pt = 5.474cm and
% carries no text on the null's own row.
%
% ===========================================================================
% PASS 12 -- 2026-08-19 -- THE SECOND ADVERSARIAL-REFUTATION REPAIR
% ===========================================================================
% EVERYTHING ABOVE THIS LINE IS LEFT INTACT. Three further refuters read the
% pass-11 build and returned 30 findings, two of them marked blocking and both
% blocking findings naming the same object. This section supersedes BY NAME the
% claims above that this pass made false, records every finding APPLIED and
% every finding OVERRULED with its reason, and re-measures the drawing.
%
% NO NUMBER CHANGED IN THIS PASS EITHER. 1.563, 0.084, 4.180, n = 60, rank 10,
% residual width 2048 and layer 4 are byte-identical to passes 9, 10 and 11.
% What changed is the value-to-POSITION map (\kx's domain, recorded below and
% in hook.json beside the two maps it supersedes) and where labels and marks
% sit. No mark's VALUE moved; every mark's x moved because the scale did.
%
% CLAIMS ABOVE THAT ARE NOW FALSE, SUPERSEDED HERE BY NAME
% ---------------------------------------------------------
% (T-1) PASS 11's O-2 ("put all three marks on one baseline" -- refused) and
%       O-3 ("drop lines from the marks" -- refused). Both are superseded; see
%       repair (1). O-2's arithmetic was about the LABELS, which still stagger;
%       it wrongly carried that conclusion over to the MARKS, which are 4.52cm
%       and 1.52cm apart and cannot collide. O-3's bar-chart objection is
%       answered by de-correlating leader length from value.
% (T-2) PASS 11's repair (2), "puts every numeral within a few per cent of its
%       own axis position". Measured, it did not: it moved each numeral one
%       decade along and left "4.180" sharing a column with the "3" tick, "1.563"
%       with the "1" tick and "0.084" with the "0.05" tick. Repair (2) below
%       replaces the approximation with an exact construction.
% (T-3) PASS 11's repair (3), "all three labels now sit left of their marks ...
%       a trailing 1.4pt disc reads as a bullet ending a line". The note
%       diagnosed the legend hazard correctly and then made ALL THREE rows
%       trail into a disc, which is the hazard. Superseded by repair (1).
% (T-4) PASS 11's repair (4), which set the domain-end tick numerals at
%       0.05 and 6. Superseded by repair (3): the ladder is now 0.03 / 0.1 /
%       0.3 / 1 / 3 / 10. The REASON recorded in repair (4) -- that the rule
%       must begin and end on a drawn tick and that every mark must have a
%       printed numeral either side of it -- stands and is preserved.
% (T-5) PASS 11's repair (13), which stopped the vertical construction guide
%       0.20cm short of h's tip "at the top end only". Its foot is repaired in
%       this pass; see (7).
% (T-6) PASS 11's repair (12), which mirrored panel (c)'s right angle to the
%       LEFT of the surviving vector. Superseded by (6): the origin cannot
%       carry that mark at all, whichever way it is mirrored.
% (T-7) LAYOUT NOTE (j) in the pass 1-9 block, "THE PASS-ONE BOUNDARY RULE GREW
%       A TAIL, from 8.72 down to 8.55", is superseded by (9): the label that
%       tail was grown to carry was deleted in pass 10, and the two dividers
%       are now identical.
% (T-8) PASS 11's claim at the head of its S-2 that "hook.json's citations are
%       repaired". They were not. Twelve of eighteen line citations in the
%       ledger did not resolve against the current Sections/04b-representation
%       .tex, and five ledger fields still described the pass-10 drawing. All
%       are repaired in this pass and hook.json now carries a pass-12
%       measurement block that supersedes the pass-10 style block by name.
%
% THE ELEVEN REPAIRS, EACH WITH THE MEASUREMENT THAT FORCED IT
% -------------------------------------------------------------
% (1) PANEL (d) IS REBUILT: THE THREE MARKS NOW SHARE ONE BASELINE AND EACH
%     NAME IS LEADERED TO ITS OWN MARK. This is the pass's one structural
%     change and it answers eight findings at once, including both blocking
%     ones. Measured in the pass-11 PDF: the three dot centres stood at
%     (323.06, 268.30), (478.26, 255.83) and (530.58, 243.36)pt -- three
%     baselines 0.44cm apart -- so the 155.20pt "gap" the log axis exists to
%     show was not a distance any reader could lay a ruler along, and the two
%     labels of the rows above ran x 363.68-470.69 and 394.28-523.00, i.e.
%     wholly inside it. Each row was also a text line ENDING in its own disc,
%     so at print three 1.4pt marks read as three bullets and the block read as
%     a key. And the null's two label lines began at x = 245.04 and 241.61
%     against panel (d)'s origin at 295.45 -- 19.0mm outside their own panel,
%     in panel (c)'s column, with "random slab" set 71pt above a band labelled
%     "the purified drug slab" in the same size and face.
%     Now: marks at y = 4.02 at their own three x, nothing else on that line;
%     names on rows 5.26 / 4.86 / 4.46; a 0.30pt leader in each mark's own
%     colour from just under its name to just above its disc. The row order is
%     null / cell line / purified slab, chosen for two reasons at once -- it is
%     the only order in which no leader crosses another row's label, and it
%     makes leader length (1.07 / 0.67 / 0.27cm) non-monotone in value (0.084 /
%     4.180 / 1.563), which is what kills O-3's bar-chart reading. The null's
%     name runs RIGHT from its mark because its mark is 1.59cm from the axis's
%     left end and the name measures 5.17cm; every text line in the panel now
%     lies inside 8.40-17.38.
% (2) EVERY PRINTED VALUE NOW STANDS EXACTLY AT ITS OWN AXIS POSITION. Each
%     label is aligned so its numeral's INNER edge sits on its mark's x:
%     "0.084" begins at \kx{0.084} = 9.9916, "4.180" ends at \kx{4.180} =
%     16.0316, "1.563" ends at \kx{1.563} = 14.5110. That is why one row is
%     value-first and two are value-last -- the value always abuts the mark and
%     the side is set by the measure, not by taste. It also lifts every printed
%     value 1.36cm or more clear of the spine, so the rule no longer carries
%     numerals within a 15pt band on both sides (S21).
% (3) THE TICK LADDER IS UNIFORM. 0.05 / 0.1 / 0.3 / 1 / 3 / 6 stepped x2, x3,
%     x3.33, x3, x2 and printed at 1.300 / 2.061 / 2.258 / 2.061 / 1.300cm -- a
%     74% spread, crowding a pair of numerals at each end of the rule, with the
%     null inside the crowded left pair. 0.03 / 0.1 / 0.3 / 1 / 3 / 10 is the
%     textbook 1-3 ladder and prints at 1.861 / 1.698 / 1.861 / 1.698 / 1.861,
%     a 10% spread. The rule still begins and ends on a drawn tick, no mark
%     sits on a tick, and every mark keeps a printed numeral on each side.
% (4) THE HOOK RULE NO LONGER HEADS A FLUSH-LEFT STACK. Title (9.96pt ink,
%     x0 57.34pt), formula (7.97pt accent, x0 57.48), accent 0.90pt rule (x0
%     57.14), boxes -- four decks on one left edge with 4.5pt of clear space
%     above the rule and 2.8pt below is the heading-and-underline idiom, and
%     the one mark whose whole job is "this operation runs the length of the
%     stream" was reading as "panel (b) starts here". The formula is now
%     right-aligned to the rule's own right end. It is a DIRECT LABEL of a mark
%     (S22), not a panel-level element (S13), and a rule's direct label belongs
%     at one of its ends.
% (5) "compared" IS ROTATED 90 DEGREES, READING UPWARD, BESIDE THE ARROW. Set
%     horizontally its baseline was 8.62 at \scriptsize while panel (b)'s head
%     was 8.54 at \footnotesize, 13.3cm away: two lines 2.27pt apart in
%     baseline, at two sizes, neither aligned nor separated. Rotated there is
%     no second baseline to near-miss and the label sits on the mark it names.
% (6) PANEL (c)'S RIGHT ANGLE MOVES TO THE FOOT OF THE PERPENDICULAR. See the
%     note at the mark itself. The short form: h = survivor + component with
%     both terms positive, so h always lies inside the quadrant the ORIGIN's
%     right angle would enclose, and no mirroring escapes that; the
%     decomposition is a rectangle, so the foot of the perpendicular carries
%     the same 90 degrees and is the one corner that is empty and whose two
%     arms both end on drawn strokes.
% (7) THE VERTICAL CONSTRUCTION GUIDE'S FOOT STOPS 0.14cm SHORT of the accent
%     component's arrow tip, the same treatment pass 11 gave its top end.
% (8) THE SURVIVING VECTOR IS LABELLED AT MID-SHAFT. At the tip, its baseline
%     (5.16) and h's (5.20) sat either side of the horizontal closing guide
%     (5.24), so the row read as "h - (hV)V^T - - - - - h", a dashed RELATION
%     between two expressions -- an equality or comparison the panel does not
%     assert. The guide cannot move; it is the closing side of the
%     parallelogram. One of the two labels had to leave its line, and the
%     survivor's is the one that has a shaft to sit beside.
% (9) THE TWO PROMPT/TARGET DIVIDERS ARE NOW IDENTICAL. Measured: pass one
%     0.686cm, symmetric; pass two 0.572cm, asymmetric. In a figure whose
%     single claim is that the two strips are the same object, the one
%     structural mark that visibly differed between them was a vestige.
% (10) \ENDCELL'S CELL IS 1.47cm, NOT 1.70. Its Inconsolata ink measures
%     1.204cm and every other cell in the strip carries 0.131-0.150cm of clear
%     space a side; this one carried 0.248. The strip end, the braces, the hook
%     rule and the comparison arrow are keyed to \stripend and moved with it.
% (11) THE INTERVAL DECLARATION IS SET ON TWO LINES. At 15.89cm it was the
%     widest object in the figure -- wider than the axis and wider than band
%     one -- which S14 forbids, and an 8pt grey line that wide, stranded a
%     centimetre below everything else, reads as a second and smaller caption
%     above the real one. Broken at the colon it measures 8.68cm and 7.10cm.
%
% WHAT WAS PROPOSED AND OVERRULED, AND WHY
% -----------------------------------------
% (P-1) DELETE PANEL (c) AND RE-LETTER (d) AS (c), giving the readout the full
%     17.38cm. Refused again, for O-1's reason: the spec's keep-list protects
%     the panel by name and CF14 forbids losing load-bearing work. The
%     finding's substantive point -- that the panel overstates the deletion --
%     is conceded in the caption, which also defers the true 10/2048 geometry
%     to fig:slab, where it is drawn at 86 degrees.
% (P-2) WIDEN THE IN-SLAB COMPONENT FROM 0.90 TO 1.35cm so the two black
%     vectors separate at the origin; and, from a different refuter, NARROW IT
%     TO 0.30cm so the panel stops overstating the deletion. Both refused,
%     because they are contradictory demands on one parameter: the angle
%     between h and the survivor IS the drawn ratio, so every millimetre of
%     separation at the origin is bought by overstating what the hook removes.
%     Pass 9 chose that trade-off deliberately at 4.0:1 and recorded it; the
%     built ratio is 3.49:1 and the fusion is confined to the bottom 1.2mm of a
%     31.4mm vector. The relocated right angle (6) also moves the corner a
%     reader has to read away from the fused base.
% (P-3) GIVE THE V_pure BAND AN ACCENT OUTLINE, so the object under test is not
%     drawn in the furniture colour while its own label is accent. The
%     COMPLAINT IS CONCEDED and the repair is refused on arithmetic. The
%     drawing measures 5.972% non-white at 200dpi inside a 173.99 x 99.95mm
%     content box against S15's 6.0% ceiling: 4.9mm2 of headroom. A 0.30pt
%     outline round the 5.80 x 0.24cm band costs 12.8mm2, and even a band
%     shortened to the extent its construction occupies costs more than the
%     ceiling allows -- the longest band an outline could be bought for is
%     2.07cm, which is shorter than the construction standing on it. Pass 9's
%     note 9b and pass 11's repair (14) also both refused an accent border, on
%     the ground that accent-on-panelbg is the drug cell's treatment to within
%     a tint. The band keeps its accent DIRECT LABEL, the accent component is
%     drawn lying inside it, and the same object is an accent filled disc in
%     panel (d), which is where the colour code is carried. RECORDED AS AN
%     OWED REPAIR for any future pass that finds ink to spend.
% (P-4) REBUILD \tokstrip AS TWO OUTER RECTANGLES PLUS INTERNAL VERTICALS, to
%     close the ~0.4mm notches the abutting cells' rounded corners leave in the
%     strip's top and bottom edges. Refused, and it is the closest call in this
%     pass. The notches are real and measurable (13 missing pixels across four
%     boundaries at 200dpi), and the repair would also SAVE 6.2mm2. But the
%     nine-rounded-cells glyph is a recorded cross-figure tie: decision 2 above
%     records that figs/licensing.tex ADOPTED this figure's cell treatment in
%     pass 9 precisely so a reader who meets the prompt strip twice meets one
%     object, and licensing.tex is out of this task's scope, so the repair
%     could not be mirrored and the two figures would diverge again. RECORDED
%     AS AN OWED REPAIR to be made in hook.tex and licensing.tex together.
% (P-5) CENTRE ALL SIX TICK NUMERALS AND CUT THE RULE TO 17.10 TO PAY FOR IT.
%     Refused; the reason is at the ladder itself. In short: at the LEFT end
%     centring puts "0.03" 0.27cm outside panel (d)'s x origin, and S13's
%     one-grid rule, measured at a 0.2pt spread, outranks the optical centring
%     of a domain-end numeral that labels no mark.
% (P-6) EXTEND THE HOOK RULE 0.25cm PAST BOTH STRIP ENDS. Refused: the
%     picture's bounding path is pinned at x = 0.00 and a rule at -0.25 would
%     put the content box at 176.3mm, past S14's 174mm ceiling. (4) breaks the
%     heading reading without moving the rule.
% (P-7) MOVE THE COMPARISON ARROW OUT TO x ~ 15.2 WITH LEADERS RUNNING BACK TO
%     THE STRIPS, to close the empty top-right quadrant. Refused. It would
%     detach the figure's comparison mark from the two objects it compares,
%     buy the composition with two 2.1cm horizontal leaders (4.5mm2, more than
%     the whole ink headroom) and re-open exactly the flow reading -- one strip
%     feeding a mark feeding the other -- that the double head exists to kill.
%     The empty quadrant is accepted and recorded.
% (P-8) RESTORE "generation would begin here" AT THE PROMPT/TARGET DIVIDER, or
%     delete the divider as an unlabelled duplicate of the brace's boundary.
%     Refused both ways. The spec's keep-list protects the divider by name; a
%     label there would restate what the 1.00cm gap already shows and cost ink
%     the ceiling does not have; and the divider is named where an inherited
%     phrase belongs, in the caption, which now reads "the divider marks only
%     where a generated continuation would have begun".
% (P-9) MIRROR PANEL (c)'S CONSTRUCTION so the removed component points LEFT
%     and h tips right. Refused: that is not a mirror, it is an error. With h
%     tipping right and the component pointing left, h is no longer the sum of
%     the two drawn vectors. The finding's DIAGNOSIS -- that the pass-11 mark
%     did not enclose the angle it claimed -- is correct and is applied as (6).
% (P-10) SET THE RIGHT ANGLE IN ink 0.40pt. Refused: the two construction
%     guides beside it are rulegrey 0.30pt, the mark is the third element of
%     that one apparatus, and splitting it off would give 0.40pt a sixth
%     meaning inside S10's five-weight table.
%
% GEOMETRY AND STYLE, RE-MEASURED AT PASS 12
% -------------------------------------------
% Content box 173.99 x 99.95mm -- fullwidth measure 170-174mm, height cap
% 100mm -- unchanged to the hundredth of a millimetre from passes 10 and 11.
% Ink 5.972% at 200dpi against a 6.0% ceiling. 710 non-space characters in the
% built PDF's own text spans over 17390mm2 is 4.08 per 100mm2 against 5.0.
% TYPE, measured in the built PDF: 7.97pt x523 and 9.96pt x146 -- two chosen
% sizes -- plus the exempt classes 6.85 x20 (\ENDCELL, Inconsolata), 5.98 x8 and
% 6.23 x3 (sub- and superscripts) and 8.30 x10 (CMSY10/CMR10 math symbols).
% Faces: Pagella-Regular 654, PalladioL-Roma 9, Inconsolata 20, PalladioL-Ital
% 14, CMSY10 7, CMR10 6. Zero sans, zero bold, zero all-caps banners. Italic
% 14 of 710 = 2.0%, and all fourteen are the math variables h and V.
% STROKES: 0.30 (furniture -- 18 cell borders' worth of rulegrey, both braces,
% six axis ticks, two construction guides, the right angle, three label
% leaders), 0.40 (two prompt/target dividers, the axis rule), 0.80 (ONE, the
% null's open-disc edge: CF5 instructs this figure by name to take 4.11(b)'s
% encoding for this object and mechanism.py draws it mew=1.1, so the singleton
% is a mandated distinct kind of object and not a stray weight), 0.90 (ONE, the
% hook rule: it is the figure's intervention and the single heaviest non-data
% stroke; at 1.00pt it would merge with the vectors and the comparison arrow,
% which are data paths, and at 0.40pt with the axis), 1.00 (three vectors and
% the comparison arrow). Five weights, both singletons justified.
% RECTANGLES: nineteen -- sixteen panelbg token cells, two accent!12 drug
% cells, and the V_pure band, which is still the only one drawn without a
% stroke. ARROWHEADS: five, three on panel (c)'s vectors and two on the
% comparison arrow; no other mark in the figure carries one.
% THE AXIS, measured in the built PDF: spine 295.25 to 549.81pt; ticks at
% 295.25 (0.03), 348.07 (0.1), 396.15 (0.3), 448.91 (1), 497.05 (3) and 549.81
% (10), i.e. tick spacings 52.82 / 48.08 / 52.76 / 48.14 / 52.76pt, a 10%
% spread against pass 11's 74%, and the first and last drawn tick still sit
% exactly on the rule's endpoints. The four interior numerals are centred on
% their ticks; the two domain-end numerals are set flush inside the rule's ends
% for the reason recorded at the ladder. THE THREE MARKS ARE AT 340.61 (null,
% white-filled open disc, 0.797pt edge), 468.48 (purified slab, accent filled)
% and 511.58 (cell line, ink filled) -- ALL THREE AT cy = 262.63, one baseline
% to the hundredth of a point. The null-to-slab gap is 127.87pt = 4.510cm and
% carries no ink at all; the gap the map predicts for a ratio of
% 18.634044258027835 is 4.521cm, so the drawn span IS the ratio to within
% 0.011cm, which is a third of the width of the stroke that draws the discs.
% The slab-to-cell-line gap is 43.10pt = 1.520cm against a predicted 1.5202cm.
% THE THREE LEADERS: 0.299pt, at x = 340.61 (ink, 30.33pt), 511.58 (ink,
% 19.00pt) and 468.48 (accent, 7.65pt) -- lengths 1.07 / 0.67 / 0.27cm against
% values 0.084 / 4.180 / 1.563, deliberately not monotone.
% ===========================================================================
%
% ===========================================================================
% v-BLOCK, 2026-08-19: WHOLE-SLATE PASS (fig:comparators, fig:licensing,
% fig:purify, fig:hook, fig:materiality reconciled against one another and
% against fig:mechanism, which is the document's reference standard and was
% NOT rebuilt). Nothing above this line is deleted; where an earlier claim is
% now false it is superseded EXPLICITLY below.
%
% WHY. The five figures were rebuilt independently and drifted. This pass owns
% only the differences BETWEEN them. No value changed in any of the five; the
% sibling .json ledgers carry the value-to-position maps, old and new.
%
% THE SLATE RULES THAT NOW HOLD ACROSS ALL FIVE (each stated once here, in
% every file, so any future single-figure edit can see what it would break):
%
%  R1 TICK NUMERALS are quietink, \scriptsize, inner sep 0, anchored north on
%     the tick and set 0.16cm below the axis rule, over a 0.10cm rulegrey
%     0.30pt tick. Before: comparators/licensing/purify quietink, hook and
%     materiality ink; ticks 0.10 / 0.12 / 0.14cm; three different gaps.
%  R2 A THRESHOLD'S OWN LABEL IS BLACK, matching the black solid 0.90pt rule
%     it names (S6). Before: comparators black, purify's "gate, 0.8" and
%     materiality's two margin labels ink.
%  R3 ONE NAME PER OBJECT (S19). The pre-registered 10% rule is "10%
%     materiality margin" in fig:comparators row (c), fig:materiality (a) and
%     (b), and in figs/family.py, which is where the name comes from.
%     fig:comparators' "pre-registered margin" is superseded; "pre-registered"
%     is now the caption's word, not the drawing's. Zero, where it is the null
%     a bar is read against, is "zero: no identity effect" in BOTH
%     fig:comparators row (c) and fig:materiality (b) -- one quantity, one
%     wording. V_pure is spelt with \textrm everywhere, as figs/slab.tex does.
%  R4 COLOUR KEY, one key for five figures (S8). accent = the drug-side object
%     under test. quietink = every comparator, null, control, replication or
%     discounted arm, AND its label, AND its leader. ink/black = thresholds,
%     axis rules, axis labels, panel heads, definitional prose. rulegrey =
%     furniture only. Within quietink the GLYPH separates two kinds: an OPEN
%     disc is a construction built to carry no signal (a null, a control, a
%     raw/before series); a FILLED disc is a real measurement that is not the
%     headline. fig:hook (d) drew both its null and its cell-line control in
%     ink and is corrected.
%  R5 PROJECTION AND ZERO RULES are quietink 0.50pt densely dotted (S7).
%     fig:hook's two panel-(c) projection lines were rulegrey 0.30pt densely
%     dashed and are corrected.
%  R6 DEFINITIONAL PROSE is \footnotesize ink and OUTRANKS every label (S4);
%     at most one block per figure, on the panel's own x origin. fig:hook's
%     one prose line was \scriptsize and is raised. fig:purify's panel-(b)
%     block carried an argument clause ("so causal claims use the purified
%     slab, not the raw") that its caption already carries, and is now
%     definitional only.
%  R7 ONE INTERVAL DECLARATION per figure, \scriptsize quietink, at the foot,
%     on the figure's x origin, opening with the word "Intervals:" and naming
%     EVERY panel including the ones with no measured mark (S17).
%     fig:comparators already did this and is the model; the other four now
%     match its form.
%  R8 THE FIVE-FIELD PROMPT STRIP is ONE object drawn in two figures, and the
%     two drawings are now congruent: cell boundaries 0 / 1.12 / 1.90 / 2.64 /
%     4.23 / 6.90 cm and cell height 0.42cm in both fig:licensing (a) and
%     fig:hook (a) and (b), labels centred in their cells, the drug cell
%     accent-bordered at 0.40pt on accent!12 and its label NOT bold. Before:
%     licensing 6.90cm wide with 0.62cm cells and a bold label, hook 7.44cm
%     wide with 0.42cm cells and a plain label -- the same five fields at two
%     sizes, four pages apart, with each caption pointing at the other.
%  R9 MEASURE. Every fullwidth figure draws 172.0-174.0mm; the one textwidth
%     figure draws 140.3mm. Ink at 200dpi inside the content bbox, counted as
%     greyscale luminance below 230 (the contract's own method), is now
%     5.63-5.99% on all five, under the 6.0% body ceiling, where before this
%     pass it ran 5.36-6.54%.
%
% WHAT THIS PASS DELIBERATELY DID NOT UNIFY, so that a later reader does not
% think it was missed:
%  * PANEL-HEAD CASE. fig:comparators sets "(a) Is it there?" and the other
%    four set "(a) where the drug name enters". comparators' heads are
%    interrogative SENTENCES and sentence case is correct for a sentence; the
%    v4 header argued this and the argument stands. fig:mechanism's lowercase
%    heads are noun phrases, not questions. One capital letter, four times.
%  * PANEL-HEAD POSITION. fig:comparators puts its heads in a right-aligned
%    left stub, the other four above the panel at its x origin. Moving them
%    above the rules costs about 20mm of height against a 118mm cap that the
%    figure already meets at 117.60mm. Not available.
%  * TWO-COLUMN GUTTER. purify and hook split at x = 8.40cm, licensing at
%    7.75cm. licensing cannot move right: "(d) what a negative would buy" is
%    4.85cm wide and its origin is already at 12.40 of a 17.30cm measure, with
%    1.5pt of slack. Moving purify and hook LEFT to 7.75 would be arbitrary.
% ===========================================================================
%
% CHANGED IN THIS FILE. (i) The prompt strip adopts fig:licensing's cell
% boundaries, 0 / 1.12 / 1.90 / 2.64 / 4.23 / 6.90cm, so the one object is one
% width in both figures (R8). The divider \bnd follows it from 7.94 to 7.40,
% the target strip from 8.44-12.42 to \tgt = 7.90 - \stripend = 11.88 with its
% four cell widths unchanged, the two braces and their "scored tokens" labels
% with it, and the comparison arrow from x 12.87 to 12.33 with its rotated
% label from 13.27 to 12.73. Panel (d), the axis and the foot do not move.
% (ii) Panel (d): the matched-dimension random slab and the cell-line positive
% control are comparators, so their marks, labels and leaders go ink ->
% quietink (R4). The null keeps the open disc and the control the filled one,
% which is the slate's glyph distinction between a construction built to carry
% no signal and a real measurement that is not the headline. accent now
% appears exactly once in panel (d), on the object under test.
% (iii) The six x-tick numerals go ink -> quietink, anchor=base -> anchor=north
% at a 0.06cm gap (R1); the end pair keep their north-west / north-east pull.
% (iv) The two panel-(c) projection lines go from rulegrey 0.30pt densely
% dashed to quietink 0.50pt densely dotted, which is the register
% fig:comparators, fig:purify and fig:materiality all use for a projection and
% which S7 reserves for it (R5).
% (v) "one sentence, one set of weights, scored twice" goes \scriptsize ->
% \footnotesize ink, so this figure's one prose block outranks its labels as
% fig:purify's and fig:licensing's do (R6).
% (vi) The hook formula's tail, "prompt and target alike", is deleted; the
% accent rule visibly spans the prompt half and runs across the divider, and
% the caption states it in full. This is the one deletion made for the ink
% ceiling rather than for a defect, and it is the first thing to restore.
% (vii) V_pure is respelt with \textrm, as figs/slab.tex and figs/purify.tex
% spell it. (viii) The foot block takes the slate's declaration form (R7).
% STILL TRUE AND STILL DECLARED: the 0.90pt accent hook rule is a singleton
% weight, kept under S10's own permission; it is accent, not black, so it
% cannot be read into the threshold register, and this figure draws no
% threshold. The ten 8.30pt glyphs are CMR10/CMSY10 math signs inside the two
% formulas, scaled 1.041 by mathpazo; they are not a third chosen type size
% and cannot be removed without deleting the definition of the hook.
% RE-MEASURED: 173.99 x 99.95mm; ink 5.955%; 806 characters, 4.63 per 100mm^2.
% One page, zero Overfull, zero Underfull.
% ===========================================================================
\begin{tikzpicture}[x=1cm, y=1cm, font=\small,
                    every node/.style={inner xsep=0pt, inner ysep=1pt}]
  % TWO SIZES, ONE FACE. \footnotesize (10.0pt) is reserved for the four panel
  % heads and the axis's quantity line; \scriptsize (8.0pt) carries everything
  % else. No \sffamily, no \bfseries, no \itshape anywhere in the picture.
  % inner xsep=0pt, PASS 11: with TikZ's 1pt x padding a base-west node's ink
  % began 1.20pt right of the coordinate it was placed on, so in every panel the
  % drawn rules started at one x and the text at another (S5, S13). The padding
  % is now zero in x and the text sits ON the panel's origin.
  \tikzset{
    ttl/.style  ={font=\footnotesize, text=ink,      anchor=base west},
    lab/.style  ={font=\scriptsize,   text=ink,      anchor=base west},
    labe/.style ={font=\scriptsize,   text=ink,      anchor=base east},
    qlab/.style ={font=\scriptsize,   text=quietink, anchor=base west},
    qlabe/.style={font=\scriptsize,   text=quietink, anchor=base east},
    tick/.style ={font=\scriptsize,   text=quietink, anchor=north,
                  inner sep=0pt},
    def/.style  ={font=\footnotesize, text=ink,      anchor=base west},
    alab/.style ={font=\scriptsize,   text=accent,   anchor=base west},
    brc/.style  ={draw=rulegrey, line width=0.30pt, decorate,
                  decoration={brace, amplitude=3.5pt, mirror}},
    arw/.style  ={-{Latex[length=3.4pt,width=2.8pt,line width=0pt]}},
    cmp/.style  ={{Latex[length=3.4pt,width=2.8pt,line width=0pt]}-%
                  {Latex[length=3.4pt,width=2.8pt,line width=0pt]}}}

  % --- one token box, and the drug box, so the two strips cannot differ ------
  % Border 0.30pt: a cell is a container, and containers are furniture here.
  % The drug cell keeps its 0.40pt accent border, because it is the one field
  % of the five that is an object in the argument -- the place the prompt names
  % the drug in plain text -- and it is the only accent inside either strip.
  \newcommand{\tokbox}[4]{%
    \draw[rulegrey, line width=0.30pt, fill=panelbg, rounded corners=1pt]
      (#1,#3) rectangle (#2,{#3+0.42});
    \node[font=\scriptsize, text=ink, inner sep=0pt]
      at ({(#1+#2)/2},{#3+0.20}) {#4};}
  \newcommand{\tokdrug}[3]{%
    \draw[accent, line width=0.40pt, fill=accent!12, rounded corners=1pt]
      (#1,#3) rectangle (#2,{#3+0.42});
    \node[font=\scriptsize, text=accent, inner sep=0pt]
      at ({(#1+#2)/2},{#3+0.20}) {drug};}
  % THE strip. One argument: the y of the box bottoms. Called twice, and that
  % is the whole guarantee that pass one and pass two are the same sentence.
  \newcommand{\tokstrip}[1]{%
    \tokbox{0.00}{1.12}{#1}{cell line}%
    \tokdrug{1.12}{1.90}{#1}%
    \tokbox{1.90}{2.64}{#1}{dose}%
    \tokbox{2.64}{4.23}{#1}{mechanism}%
    \tokbox{4.23}{6.90}{#1}{control cell sentence}%
    \tokbox{\tgt}{8.88}{#1}{gene$_1$}%
    \tokbox{8.88}{9.86}{#1}{gene$_2$}%
    \tokbox{9.86}{10.41}{#1}{\dots}%
    \tokbox{10.41}{\stripend}{#1}{\ENDCELL}}
  % the prompt/target boundary: mid-gap, where nothing else is drawn.
  \def\bnd{7.40}
  \def\tgt{7.90}
  % PASS 12: 12.42, not 12.65. \ENDCELL's cell was the one field of nine whose
  % clear space did not match the other eight: measured ink 1.204cm (Inconsolata
  % at \scriptsize, built PDF x 374.49-408.76pt) inside a 1.70cm cell is 0.248cm
  % a side against 0.131-0.150cm everywhere else, so at print it was the one
  % roomy cell in a strip whose whole claim is regularity. 1.204 + 2 x 0.133 =
  % 1.470. The strip end is a \def and the brace, the hook rule and the
  % comparison arrow are all keyed to it, so they move with it.
  \def\stripend{11.88}

% ###########################################################################
% # BAND ONE -- the two passes.  x origin 0.00
% ###########################################################################

  % ---------------- (a) pass one, clean -------------------------------------
  \node[ttl] at (0.00,10.34) {(a) pass one, clean};

  \tokstrip{9.54}

  % The divider is rulegrey 0.40pt SOLID, not ink 0.5pt densely dotted: S7
  % reserves the dotted stroke for a zero rule, a leader or a projection line
  % and forbids it a label of its own, and this one is labelled. It is a piece
  % of the strip's own structure, so it takes a strip weight, and it sits at
  % mid-gap because a rulegrey line on a rulegrey box border is invisible.
  % PASS 12: 9.40 to 9.98, not 9.40 to 10.10. The two dividers were drawn at two
  % different lengths and two different symmetries -- pass one 0.686cm with
  % 0.133cm clear above and below its boxes, pass two 0.572cm with 0.134 below
  % and 0.018 above (clipped to clear the hook rule at 7.88). The one structural
  % mark that visibly DIFFERED between two strips whose whole claim is that they
  % are identical was a leftover: layout note (j) below grew that tail to carry
  % the label "generation would begin here", which pass 10 deleted. NOTE (j) IS
  % THEREFORE SUPERSEDED -- its reason no longer exists. Both dividers now run
  % 0.58cm with the same 0.14/0.02 offsets from their own boxes.
  \draw[rulegrey, line width=0.40pt] (\bnd,9.40) -- (\bnd,9.98);

  % The brace spans the TARGET side only and starts at the first target box's
  % own left edge, 9.20. Both braces carry the SAME label: the two spans are
  % the same tokens read twice, and two different labels sent a reader looking
  % for a difference between them.
  \draw[brc] (\tgt,9.44) -- (\stripend,9.44);
  \node[font=\scriptsize, text=ink, anchor=base] at (9.89,9.04)
    {scored tokens};

  % ---------------- (b) pass two, the slab projected out --------------------
  % PASS 11: "at layer 4" is added to the head. The word "every" in "at every
  % position" ranges over the SEQUENCE, not over depth, and the only place the
  % depth was named was the axis block 40mm away, attached to the READOUT rather
  % than to the intervention -- so a reader could take the hook to run at every
  % layer as well as at every position, which is exactly the misreading decision
  % 14 above says the figure exists to prevent. The axis block keeps its own
  % "layer 4" because CF10 requires the layer named where the one drawn layer's
  % marks are.
  \node[ttl] at (0.00,8.54) {(b) pass two, the slab projected out at layer 4};

  % THE HOOK. The RULE is what says "everywhere along here": it spans the whole
  % strip, prompt half and target half, and runs straight across the
  % prompt/target boundary, which is the point. Its formula sits ON the rule
  % rather than 1.1cm under the strip, so the label names the mark it is beside.
  % The two short arrows of pass 9 are gone: they implied a count of two and
  % the caption had to deny it.
  % PASS 11: the rule drops from 7.92 to 7.88, so its leading is ASYMMETRIC --
  % 0.18cm of clear space above it to the formula's descenders, 0.08cm below it
  % to the box tops. At 7.92 it was equidistant between an accent line of type
  % above and the boxes below, and accent type with an accent rule tight beneath
  % it is the heading-and-underline idiom; the rule now binds to the strip it
  % acts on. It is NOT dropped onto the boxes: a mark running into a rule
  % uninterrupted is what the house rule forbids, and the drug cell's border is
  % accent too.
  % PASS 12: THE FORMULA IS RIGHT-ALIGNED TO THE RULE'S OWN RIGHT END. Flush
  % left it made a three-deck flush-left stack -- title (9.96pt ink, x0 57.34pt),
  % formula (7.97pt accent, x0 57.48pt), accent 0.90pt rule (x0 57.14pt), boxes
  % -- with 4.5pt of clear space above the rule and 2.8pt below, and an accent
  % rule tight under an accent line of type sharing its left edge is the
  % heading-and-underline idiom. The mark that has to say "this operation runs
  % along the stream" was reading as "panel (b) starts here". The formula is a
  % DIRECT LABEL of the rule, not a panel-level element, so S22 governs it and
  % not S13: a rule's direct label sits at one of its ends, and putting it at
  % the far end from the panel's title breaks the stack without moving the rule.
  % The rule is NOT extended past the strip ends, as one refuter proposed: the
  % picture's bounding path is pinned at x = 0.00 and a rule at -0.25 would put
  % the content box at 176.3mm, past S14's 174mm ceiling.
  \node[font=\scriptsize, text=accent, anchor=base east] at (\stripend,8.12)
    {$h \leftarrow h-(hV)V^{\top}$ at every position};
  \draw[accent, line width=0.90pt] (0.00,7.88) -- (\stripend,7.88);

  \tokstrip{7.38}

  \draw[rulegrey, line width=0.40pt] (\bnd,7.24) -- (\bnd,7.82);
  \draw[brc] (\tgt,7.28) -- (\stripend,7.28);
  \node[font=\scriptsize, text=ink, anchor=base] at (9.89,6.88)
    {scored tokens};

  % ---------------- the tie: the figure's central claim, drawn --------------
  % It sits UNDER both strips, on band one's own x origin, and it is now band
  % one's FOOT LINE: on its own baseline, 0.44cm below pass two's brace label
  % and 0.76cm above the two panel heads of band two, so it belongs to the band
  % it closes.
  % PASS 11, TWO REPAIRS. (i) POSITION. It was at 6.88, the EXACT baseline of
  % pass two's brace label, and directly under pass two's prompt half -- the one
  % half of the strip that carries no brace. A reader therefore read it as the
  % missing brace label of that half ("prompt half: the same sentence, box for
  % box | target half: scored tokens"), i.e. as a statement about ONE strip,
  % which is the opposite of what it says. Nothing else is now set on 6.44.
  % (ii) WORDING. "the same sentence, box for box" carried only half of the
  % question this figure has to answer -- are the two passes scoring the same
  % sentence WITH THE SAME WEIGHTS -- and nothing anywhere on the drawing named
  % the weights at all; a reader could infer them only from the formula being an
  % activation edit. The deleted pass-9 subtitle "the same prompt, the same
  % sentence, the same weights" did carry it and pass 10's note claimed the tie
  % had inherited it, which was not true. It now reads "one sentence, one set of
  % weights, scored twice". This also clears the S18 duplication with the
  % caption's "are the same sentence box for box", which is reworded in the
  % staged caption to "emitted by one macro, so no difference between them is
  % possible except the hook".
  \node[def] at (0.00,6.36) {one sentence, one set of weights, scored twice};

  % ---------------- the comparison -----------------------------------------
  % DOUBLE-headed, and ink. The two passes are independent scorings of one
  % fixed sentence by one set of weights; that is a comparison, not a flow, and
  % the single head read as pass one feeding pass two. The single accent
  % arrowhead is reserved for an operation acting in a direction.
  % PASS 11, TWO REPAIRS. (i) GEOMETRY. Its two heads pointed at nothing: it ran
  % between the two brace LABELS' baselines (9.04 to 6.88), so the upper head sat
  % 0.50cm BELOW the strip it was meant to span and both heads terminated in
  % white space. It now runs between the two strips' own vertical centres, 9.75
  % to 7.59, and sits 0.45cm off the strip end instead of 0.75cm, so each head is
  % level with the band it names.
  % (ii) LABEL. "KL from the clean / forward pass" is deleted and replaced by one
  % word. It set panel (d)'s axis quantity string a second time, 130pt away, so a
  % reader could take this vertical rule for an axis OF that quantity; the axis
  % names the quantity, and this mark's job is to say that the two passes are
  % compared and not chained. S20 is satisfied -- the role is stated in one verb
  % of comparison, beside a double head that S11 gives a comparison. The 17
  % characters this frees are what pay for "at layer 4", the reworded tie and the
  % two domain-end tick numerals inside the 6.0% ink ceiling.
  % PASS 12, TWO MOVES. (i) x = 12.87, not 13.10: the strip end moved from 12.65
  % to 12.42 and the arrow keeps pass 11's 0.45cm standoff from it.
  % (ii) THE LABEL IS ROTATED 90 DEGREES, READING UPWARD, BESIDE THE ARROW'S
  % SHAFT AND CENTRED ON ITS MIDPOINT. Set horizontally it sat on baseline 8.62
  % at \scriptsize while panel (b)'s head sat on 8.54 at \footnotesize, 13.3cm
  % away: two lines of type 2.27pt apart in baseline, at two sizes, neither
  % sharing a baseline nor clearly separated, which is the near-alignment a
  % reader reads as carelessness. Rotated, there is no second baseline to
  % near-miss, the label sits ON the mark it names instead of beside it, and it
  % adopts fig:mechanism's own rotated-label idiom. It reads upward, as 4.11's
  % y-labels do, so its ascenders point left and its ink clears the 1.00pt shaft
  % by 0.18cm.
  \draw[ink, line width=1.00pt, cmp] (12.33,9.75) -- (12.33,7.59);
  \node[font=\scriptsize, text=ink, anchor=base, rotate=90]
    at (12.73,8.67) {compared};

% ###########################################################################
% # BAND TWO, LEFT -- (c) what the hook does to one vector.  x origin 0.00
% ###########################################################################
  \node[ttl] at (0.00,5.68) {(c) what the hook does to one vector};

  % THE SLAB. PASS 11, TWO CHANGES, AND THE SECOND PAYS FOR MOST OF THIS PASS.
  % (i) IT RUNS 0.00 TO 5.80, not 0.00 to 3.80. The lower band held a rectangle
  % 38 x 46mm between this panel and panel (d) containing literally no ink --
  % 10% of the content area, and the largest single void on the drawing. The
  % band is the panel's own object and lengthening it is what a subspace does;
  % the construction moves with it to stay on the band's midpoint.
  % (ii) THE rulegrey OUTLINE IS GONE; the band is now a fill alone. It was
  % byte-identical in stroke and fill to the eighteen token cells above it
  % (rulegrey 0.30pt on panelbg), so one rectangle treatment carried two
  % meanings -- "a field of the prompt", eighteen times, and "the purified drug
  % slab, a ten-dimensional subspace", once -- and the only difference was
  % square against rounded corners on a 9pt-tall shape. An ACCENT outline was
  % considered and rejected again for pass 9's reason (accent fill plus accent
  % border is the drug cell's own treatment to within a tint), so the border
  % goes rather than changes colour. The band keeps its accent direct label, is
  % 24 times as wide as it is tall, and is the only fill-without-stroke region
  % in the figure. It also frees 8.7mm2 of the 6.0% ink ceiling, which is what
  % buys "at layer 4", the reworded tie and the two new domain-end ticks.
  \fill[panelbg] (0.00,1.98) rectangle (5.80,2.22);

  % construction lines first, so the arrows and the right angle cover them.
  % PASS 11: the vertical guide stops 0.20cm short of h's tip. It used to end
  % ON the tip, inside the arrowhead, so a leader ran into the mark it points
  % from. The horizontal guide still touches both tips: it IS the closing side
  % of the parallelogram and a gap there would break the construction.
  % PASS 12: the vertical guide now stops at 2.24, 0.14cm above the component's
  % line and just clear of the V_pure band's top edge. It used to run to 2.10,
  % which is exactly the accent component's arrow TIP, so a dashed leader ended
  % inside a mark and at 6x zoom the guide read as a continuation of the
  % component rather than as the side of a rectangle. The same 0.20cm-standoff
  % treatment pass 11 gave the guide's top end, applied to its foot.
  \draw[quietink, line width=0.50pt, densely dotted] (3.40,5.04) -- (3.40,2.24);
  \draw[quietink, line width=0.50pt, densely dotted] (2.50,5.24) -- (3.40,5.24);

  \draw[ink, line width=1.00pt, arw] (2.50,2.10) -- (3.40,5.24);
  \node[lab] at (3.50,5.20) {$h$};

  % PASS 12: THE SURVIVOR IS LABELLED AT MID-SHAFT, NOT AT ITS TIP. At the tip
  % its baseline was 5.16 and $h$'s was 5.20, with the horizontal closing guide
  % between them on 5.24 -- so the row read left to right as
  % "h - (hV)V^T - - - - - h", a dashed RELATION joining two expressions, i.e.
  % an assertion of equality or comparison the panel does not make (the two are
  % a vector and its component, not two sides of anything). Moving one of the
  % two labels off the guide's line is what kills that reading; the guide itself
  % cannot move, because it IS the closing side of the parallelogram and must
  % touch both tips. The label now sits beside the shaft it names, 0.12cm from
  % it and 0.56cm from $h$'s shaft at the same height, so the attachment is
  % unambiguous and no leader is required.
  \draw[ink, line width=1.00pt, arw] (2.50,2.10) -- (2.50,5.24);
  \node[labe] at (2.38,3.62) {$h-(hV)V^{\top}$};

  % the in-slab component, and its fate carried as a mark, not only in words.
  % PASS 11: its label rises from +0.26cm to +0.28cm above the segment and the
  % band thins from 0.32 to 0.24cm, because at the old geometry the label's
  % parenthesis descenders cleared the band's top rule by 0.72pt -- two pixels
  % at 200dpi -- and would have closed on press. The label is still 0.286cm
  % from its own mark, inside the 3mm at which this document requires a leader.
  % PASS 12: the label's x origin moves 3.46 -> 3.58. At 3.46 its opening
  % parenthesis began 0.06cm (2.0pt) right of the vertical guide at 3.40, so at
  % print the dashed guide and the "(" read as one glyph cluster. 0.18cm now.
  \draw[accent, line width=1.00pt, arw] (2.50,2.10) -- (3.40,2.10);
  \node[alab] at (3.58,2.38) {$(hV)V^{\top} \to 0$};

  % the right angle: the component is set to ZERO, not merely reduced. 0.40cm
  % on both legs, against pass 9's 0.16cm, which a reader did not see.
  % PASS 11: IT IS MIRRORED TO THE LEFT OF THE SURVIVING VECTOR. On the right it
  % was unreadable: the 1.0pt black h vector crossed its horizontal leg 0.12cm
  % from the corner, and its vertical leg terminated exactly on the accent
  % component's shaft, so a 0.30pt rulegrey glyph sat under two heavier strokes
  % and read as an artefact. h rises to the RIGHT of the survivor, so the corner
  % between the survivor and the slab on the LEFT is empty; the same 90 degrees,
  % the same two arms, nothing crossing it.
  % PASS 12: THE MARK MOVES TO THE FOOT OF THE PERPENDICULAR, (3.40,2.10).
  % Pass 11 mirrored it to the LEFT of the survivor to escape a collision, and
  % in doing so put it in a quadrant where NEITHER of its arms lies on a drawn
  % ray: the two rays that leave the shared origin (2.50,2.10) are the survivor
  % (up) and the accent component (right), so they bound the up-RIGHT quadrant,
  % while the mark enclosed the up-LEFT one -- its upper arm dead-ended on the
  % survivor's 1.0pt shaft and its lower arm ended at (2.10,2.10) inside the
  % V_pure band's fill, on the backward extension of a ray that is not drawn.
  % At 6x zoom it read as a detached grey bracket, and three refuters read it
  % that way independently.
  % THE ORIGIN CANNOT CARRY THE MARK AT ALL, and that is why this is a move and
  % not a re-mirror: h = survivor + component with both terms positive, so h
  % ALWAYS lies inside the very quadrant the origin's right angle encloses, and
  % at 0.40cm legs h crosses the far arm 0.115cm from its corner whichever way
  % the construction is mirrored. The decomposition is a RECTANGLE -- origin,
  % (3.40,2.10), (3.40,5.24), (2.50,5.24) -- so all four corners carry the same
  % 90 degrees, and the foot of the perpendicular is the one corner that is
  % empty. Its two arms end on drawn strokes: (3.40,2.50) on the vertical
  % construction guide, (3.00,2.10) on the component's own shaft. It asserts
  % what decision 10 says it must -- the surviving vector is orthogonal to the
  % slab, so the slab component is set to ZERO and not merely reduced.
  % IT STAYS rulegrey 0.30pt. One refuter asked for ink 0.40pt on the ground
  % that a load-bearing assertion should not share a stroke with a box border.
  % Overruled: the two construction guides in this same panel are rulegrey
  % 0.30pt, this mark is the third element of that one apparatus, and splitting
  % it off would give 0.40pt a sixth meaning inside S10's five-weight table.
  \draw[rulegrey, line width=0.30pt]
    (3.40,2.50) -- (3.00,2.50) -- (3.00,2.10);

  % ONE label for the band, 0.32cm under its own left edge: the symbol and
  % its gloss in one place, which is the only licence S19 gives the symbol.
  \node[alab] at (0.00,1.66) {$V_{\textrm{pure}}$, the purified drug slab};

% ###########################################################################
% # BAND TWO, RIGHT -- (d) what comes out.  x origin 8.40
% ###########################################################################
  \node[ttl] at (8.40,5.68) {(d) what comes out};

  % VALUE-TO-POSITION MAP, LOGARITHMIC. Braced: it is a pgfmath group, and a
  % shim added outside the braces is not a coordinate, so a label that needs a
  % shim takes it as an xshift on the node.
  %   \kx{v} = 8.40 + (ln v - ln 0.03)/(ln 10 - ln 0.03) * 8.98
  %   domain [0.03, 10] nats; spine 8.40 to 17.38; ticks 0.03 0.1 0.3 1 3 10
  %   0.084 -> 9.9916   1.563 -> 14.5110   4.180 -> 16.0316
  % A horizontal gap is a ratio: null to slab is 4.5194cm and IS 18.6x; slab to
  % cell line is 1.5206cm and IS 37% read the other way.
  %
  % PASS 12, THE DOMAIN CHANGES FROM [0.05, 6.0] TO [0.03, 10] AND NO VALUE
  % MOVES. The pass-11 ladder was 0.05 / 0.1 / 0.3 / 1 / 3 / 6: steps of x2, x3,
  % x3.33, x3, x2, which printed as spacings 1.300 / 2.061 / 2.258 / 2.061 /
  % 1.300cm -- a 74% spread, with a crowded pair of numerals jammed at each end
  % of the rule and the null sitting inside the crowded left pair. A reader
  % cannot infer the scale's rule from a ladder whose steps are not constant,
  % and fig:mechanism(b), the figure this axis was built to imitate, uses a
  % clean decade ladder at uniform spacing. 0.03 / 0.1 / 0.3 / 1 / 3 / 10 is the
  % textbook 1-3 ladder: steps x3.33, x3, x3.33, x3, x3.33, printing as
  % 1.861 / 1.698 / 1.861 / 1.698 / 1.861cm, a 10% spread. The rule still begins
  % and ends on a drawn tick (S13), every mark still has a printed numeral on
  % each side of it, and no mark sits on a tick. Only the value-to-position map
  % changes; 0.084, 1.563 and 4.180 are byte-identical to passes 9, 10 and 11.
  \def\kx#1{{8.40+(ln(#1)-ln(0.03))/(ln(10)-ln(0.03))*8.98}}

  % PASS 12 REBUILDS THIS PANEL, AND IT IS THE PASS'S ONE STRUCTURAL CHANGE.
  % ------------------------------------------------------------------------
  % WHAT WAS WRONG, measured in the pass-11 PDF and reported independently by
  % all three refuters, two of them as blocking:
  %  (a) THE GAP WAS NOT A GAP. The log axis exists for exactly one reason --
  %      on a log scale a horizontal distance IS a ratio -- and the three marks
  %      stood on three different baselines 0.44cm apart (268.30, 255.83,
  %      243.36pt), so there was no single line to lay a ruler along. Worse, the
  %      155.20pt between the null and the purified slab was filled, one row up,
  %      by "the purified drug slab --- 1.563" (x 363.68-470.69) and, two rows
  %      up, by "cell line, the positive control --- 4.180" (394.28-523.00).
  %      The caption asserted a span the drawing did not draw.
  %  (b) IT READ AS A LEGEND. Each row was a text line TERMINATING in its own
  %      disc, so at print each 1.4pt mark read as a bullet ending a line: three
  %      key entries, stacked in descending value order, i.e. the one
  %      construction S22 forbids outright. Pass 11's own note (3) names this
  %      hazard verbatim and then made all three rows trail instead of one.
  %  (c) THE NULL'S LABEL WAS OUTSIDE ITS OWN PANEL. Its two lines began at
  %      x = 245.04 and 241.61pt against a panel origin of 295.45 -- 19.0mm to
  %      the LEFT of the axis's own left end, in panel (c)'s column, 20pt from
  %      the end of panel (c)'s V_pure band, so "random slab" sat above a band
  %      labelled "the purified drug slab" in the same size and face.
  %  (d) THE ROWS WERE MIS-SET VERTICALLY. The null's first line (baseline 4.05)
  %      and the whole slab row (4.19) fell in one connected ink row; and panel
  %      (d)'s head stood 0.98cm above its first row where panel (c)'s stood
  %      0.48cm above its first content, so two side-by-side panels began at
  %      two depths.
  % ------------------------------------------------------------------------
  % WHAT IS DRAWN NOW. THE THREE MARKS SHARE ONE BASELINE, y = 4.02, at their
  % own three x. Nothing else is set on that line, so the null-to-slab distance
  % is a single uninterrupted 4.52cm horizontal span between two discs and the
  % caption's "a horizontal gap IS their ratio" is true of the drawing. Each
  % NAME sits on its own row above, ending (or, for the null, beginning) exactly
  % at its mark's x, and is joined to its mark by a 0.30pt leader in the mark's
  % own colour -- which is S22's own prescription for a label more than 3mm from
  % its mark, and which is what stops any row reading as a key entry, because a
  % mark tied to its name by a drawn leader is not a bullet.
  % PASS 11'S O-2 IS SUPERSEDED, NOT REVERSED. O-2 refused "one baseline with
  % the labels staggered above" on the ground that the three names measure about
  % 11cm against a 9cm axis. That is still true and the labels are still
  % staggered on three rows; what O-2 got wrong is that the MARKS have to
  % stagger with them. They do not: they are 4.52cm and 1.52cm apart and cannot
  % collide.
  % PASS 11'S O-3 IS SUPERSEDED. O-3 refused drop lines because three stems to
  % the spine would have lengths rising with value and would read as a bar
  % chart. These leaders do not touch the spine and their lengths are 1.07,
  % 0.67 and 0.27cm against values 0.084, 4.180 and 1.563 -- deliberately NOT
  % monotone in value, which is why the row order is null / cell line /
  % purified slab and not the pass-11 descending-value order. The row order is
  % also the only one in which no leader crosses another row's label: the
  % null's leader at 9.99 passes left of both rows below it, and the cell line's
  % at 16.03 passes right of the slab row's right end at 14.51.
  % S21, ANSWERED EXACTLY RATHER THAN APPROXIMATELY. A numeral printed near a
  % value axis is read at the position it occupies. Because each label is
  % aligned so that its numeral's INNER edge sits on its own mark's x, each
  % printed value now stands at its own axis position to the character:
  % "0.084" begins at \kx{0.084}, "4.180" ends at \kx{4.180}, "1.563" ends at
  % \kx{1.563}. That is why the null's row is value-first and the other two are
  % value-last -- the value always abuts the mark, and which side the name falls
  % on is set by the measure, the null's mark being only 1.59cm from the axis's
  % left end. No printed value is now within 1.36cm of the spine, so the rule
  % does not carry numerals on both sides.
  \def\my{4.02}
  % row one: the null.  Label runs RIGHT from its mark; the value leads.
  \node[qlab] at (\kx{0.084},5.26) {0.084 --- matched-dimension random slab};
  \draw[quietink, line width=0.30pt] (\kx{0.084},5.16) -- (\kx{0.084},4.09);
  % the null is an OPEN disc, white fill, 0.8pt edge -- 4.11(b)'s own encoding
  % for this same object (mechanism.py: mfc="white", mew=1.1, color=INK).
  \draw[quietink, line width=0.8pt, fill=white] (\kx{0.084},\my) circle (1.4pt);

  % row two: the positive control.  Label runs LEFT; the value trails.
  \node[qlabe] at (\kx{4.180},4.86) {cell line, the positive control --- 4.180};
  \draw[quietink, line width=0.30pt] (\kx{4.180},4.76) -- (\kx{4.180},4.09);
  \fill[quietink] (\kx{4.180},\my) circle (1.4pt);

  % row three: the object under test.
  \node[font=\scriptsize, text=accent, anchor=base east]
    at (\kx{1.563},4.46) {the purified drug slab --- 1.563};
  \draw[accent, line width=0.30pt] (\kx{1.563},4.36) -- (\kx{1.563},4.09);
  \fill[accent] (\kx{1.563},\my) circle (1.4pt);

  % THE TICK LADDER. PASS 11: six numerals, not four, and the outer two are the
  % domain's own ends, so the rule now BEGINS and ENDS on a drawn tick. At four
  % ticks the rule overran its tick lane by 36.87pt (13.0mm) at each end -- 29%
  % of the axis carrying no numeral, the exact defect S13 measures and condemns
  % in fig:comparators row (d) -- and BOTH extreme marks fell in those bare
  % zones, so neither the null nor the positive control could be placed against
  % a printed number. 0.05, 0.1, 0.3, 1, 3, 6 still reads as a log ladder and
  % puts a numeral either side of every mark.
  % The numerals are INK, not quietink. quietink is this document's comparator
  % and null colour (S8/CF4); spending it on the scale while the figure's actual
  % null is labelled in ink inverted the code. fig:mechanism sets every tick
  % numeral in ink and reserves quietink for its annotations.
  \draw[ink, line width=0.40pt] (8.40,3.44) -- (17.38,3.44);
  \foreach \v in {0.03, 0.1, 0.3, 1, 3, 10}{
    \draw[rulegrey, line width=0.30pt] (\kx{\v},3.44) -- (\kx{\v},3.34);}
  \foreach \v/\t in {0.1/0.1, 0.3/0.3, 1/1, 3/3}{
    \node[tick] at (\kx{\v},3.28) {\t};}
  % the two domain-end numerals are pulled flush INSIDE the rule's own ends, so
  % no ink crosses the drawn measure (S14). Centred on their ticks they would
  % have overhung x = 17.38 by 0.09cm and the float overfull by 1.88pt.
  % PASS 12 CONSIDERED AND OVERRULED centring all six and cutting the rule to
  % 17.10 to pay for it. At the LEFT end that is not available: "0.03" centred
  % on 8.40 begins at 8.13, which is 0.27cm outside panel (d)'s x origin, and
  % S13's one-grid rule -- measured at a 0.2pt spread by pass 11's repair (15)
  % -- outranks the optical centring of a domain-end numeral that labels no
  % mark. The residual is 0.275cm on "0.03" and 0.16cm on "10", which on this
  % ladder is 7% and 4% of a decade.
  \node[tick, anchor=north west] at (8.40,3.28) {0.03};
  \node[tick, anchor=north east] at (17.38,3.28) {10};

  % THE AXIS BLOCK. Quantity and unit on line one, scope on line two, both
  % left-aligned to the axis's own left end, directly under the tick numerals.
  % No result, no interval disclaimer and no cross-reference may sit here.
  \node[ttl]  at (8.40,2.68)
    {KL from the clean forward pass (nats), logarithmic scale};
  \node[qlab] at (8.40,2.34)
    {mean over the target tokens of 60 teacher-forced prompts, layer 4};

% ###########################################################################
% # THE FIGURE'S ONE INTERVAL DECLARATION.  x origin 0.00
% ###########################################################################
  % PASS 12: SET ON TWO LINES, NOT ONE. As one line it measured 15.89cm and was
  % the WIDEST OBJECT IN THE FIGURE -- wider than the axis (8.98), wider than
  % band one (14.54) -- which S14's check forbids outright ("the widest object
  % must be an axis or a panel, never a label"), and a 15.89cm line of 8pt grey
  % stranded a centimetre below everything else reads as a second, smaller
  % caption set above the real one. S17 expressly allows two lines. Broken at
  % the colon, the block measures about 7.1cm, and its two baselines are set so
  % the picture's lowest ink is unchanged and the height cap is not touched.
  \node[qlab] at (0.00,1.04)
    {Intervals: none in (d) --- the stored dispersion is a normal-approximation};
  \node[qlab] at (0.00,0.70)
    {standard error, not this document's convention. (a) to (c) carry no measured mark.};

  % bounding box, so the picture is exactly the measure it claims
  \path (0,0.58) rectangle (17.38,10.62);
\end{tikzpicture}
%
% ---------------------------------------------------------------------------
% READY TO PASTE into Sections/04b-representation.tex, at sec:res-used,
% immediately after the paragraph "Step four, the causal measurement in logit
% space" (04b:261-275). The caption below is the PASS 11 replacement AS REVISED
% AT PASS 12. It is also staged, keyed by \label{fig:hook}, in
% figs/_PENDING_CAPTIONS.tex -- a file PASS 11 CREATED, because the pass-10 note
% in this position asserted it existed and it did not. NO BUILDER EDITS
% Sections/. The staged copy carries the pass-12 revision note; the two copies
% are kept byte-equal in the \caption itself and figs/_PENDING_CAPTIONS.tex is
% the authoritative one.
% PASS 12 REVISED THREE THINGS IN IT, none of them a numeral already there:
%   * panel (d)'s marks are described as standing on ONE LINE, tied to their
%     names by leaders, because pass 11's three baselines meant the "horizontal
%     gap IS the ratio" sentence described a span that was not on the page;
%   * the POSITIVE CONTROL's size confound is disclosed for the first time
%     (cell line's subspace fraction is 20.9x the purified slab's at layer 4,
%     15-22x over the seven -- 04b:370), beside the null's 5.6x, because the
%     37% was being quoted against a comparator matched on nothing at all;
%   * the random-draw disclosure is strengthened from "one random draw" to the
%     fact it understated: with two stored draws at seven layers the file holds
%     fourteen ratios and the drawn 18.634 is the largest of them, and the two
%     draws swap order with depth (at layer 6, 14.7x undrawn against 5.7x
%     drawn).
% Two figures inside the DELETIONS block of the staged copy are superseded by
% the new log domain: the null-to-slab gap is 4.510cm, not 5.474, and the
% slab-to-cell-line gap is 1.520cm, not 1.845. No value moved; the map did.
%
% \begin{figure}[htbp]
% \begin{fullwidth}
% \centering
% % ===========================================================================
% figs/hook.tex -- fig:hook, step four of sec:res-used drawn as what it
% physically is: one true sentence, laid out as a token strip, scored twice by
% the same weights, with a forward hook that zeroes the drug slab at every
% position in between, and the KL read only over the target tokens.
%
% FIGURE BODY ONLY. No preamble, no document environment, no float, no caption.
% \input it from inside a figure/fullwidth float in
% Sections/04b-representation.tex, at sec:res-used, immediately after the
% paragraph "Step four, the causal measurement in logit space" (04b:191-203).
% The complete float block sits COMMENTED at the foot of this file.
%
% Compiles against thesis_v2/preamble.tex and nothing else: no \includegraphics,
% no external data, no TikZ library beyond the ones preamble.tex already loads
% (this file uses arrows.meta and decorations.pathreplacing).
%
% THE MISREADING THIS FIGURE EXISTS TO KILL
% -----------------------------------------
% A reader who has met "ablation" in a generative setting arrives with a default
% picture: the model writes a sentence, the hook is switched on, the model
% writes a DIFFERENT sentence, and the KL measures how far the two sentences
% drifted apart. Every clause of that is false here. Nothing is sampled. The
% sentence is fixed, real, and identical in both passes. And the hook acts on
% the prompt as well as on the target, so it is not an intervention "on
% generation" at all. The figure is built so that all three corrections are
% legible from the marks alone: the two strips are drawn by ONE macro so they
% are provably identical box for box; the nine hook ticks include the four that
% land on the prompt; and the two braces span the target boxes only and stop
% dead at the prompt/target divider.
%
% EVERY NUMBER DRAWN, AND WHERE IT COMES FROM
% -------------------------------------------
% Source: RESULTS_cluster/workspace_probe.json, key layers.4.kl. Config there:
% tier tier2_unseen_drugs, n_drugs 12, n_per_drug 40, n_dims 10,
% n_kl_prompts 60, seed 42. Sibling ledger of the exact values drawn:
% thesis_v2/figs/hook.json.
%
%   mark                    key path                       source value      drawn
%   purified drug slab      layers.4.kl.pure_drug[0]       1.5631485521793365  1.563
%   matched random null     layers.4.kl.random_matched[0]  0.08388670384883881 0.084
%   cell-line control       layers.4.kl.cell_line[0]       4.179661695162455   4.180
%   n per entry             layers.4.kl.*[2] = config.n_kl_prompts             60
%   slab rank               layers.4.pure_drug_dims = config.n_dims            10
%
%   The three dots are the ONLY data-bearing coordinates in the figure. Their
%   positions are \kx{v} = 8.20 + v/4.5*8.70, a linear map of nats to cm:
%     0.084 -> 8.3624 ,  1.563 -> 11.2218 ,  4.180 -> 16.2813
%
%   COMPUTED, not stored as floats:
%   ratio to the null       pure_drug[0]/random_matched[0] 18.634044258027835  18.6x
%   share of the control    pure_drug[0]/cell_line[0]      0.3739892522853049  37%
%   Both rounded forms are also embedded verbatim in the prose string
%   layers.4.ablation_note, which reads "...18.6x the functional variance of a
%   matched-random slab (37% of the cell-line control)...".
%
%   NOT drawn: layers.4.kl.pure_drug[1] = 0.021199743463084047, and its two
%   siblings. These are normal-approximation standard errors over the 60
%   prompts, not this document's interval convention (95% bootstrap clustered
%   over cell lines). figs/fig-mechanism.json already records intervals_drawn
%   false for exactly this reason, and mechanism.py's header states it: drawing
%   a +/-0.003 whisker beside the thesis's genuine +/-0.03 intervals would
%   invite a reader to distrust all of them. The three marks are therefore bare
%   solid dots, which is how this document says "no interval", and the caption
%   says why.
%
%   2048 (the residual width) is prose only: 04b-representation.tex:383 and
%   07-Appendix.tex, sec:app-identity. It appears here as a LABEL ("ten of the
%   2048 directions"), never as a coordinate. Nothing in the vector panel is to
%   scale and the caption says so.
%
% AGREEMENT WITH THE TEXT (checked, not assumed):
%   04b-representation.tex:229  tab:workspace layer-4 row "1.563 & 0.084 &
%                               $18.6\times$"                            -> OK
%   04b-representation.tex:193  "projects the slab out at every position
%                               ($h \leftarrow h - (hV)V^{\top}$)"
%                               -> the formula is set verbatim and the nine
%                                  ticks are what "every position" means.  OK
%   04b-representation.tex:195  "mean KL divergence from clean to ablated over
%                               the target tokens, on $60$ prompts"
%                               -> the braces span the target only; n = 60 is
%                                  printed under the axis.                OK
%   04b-representation.tex:195  \gloss{teacher forcing}{the true sentence is
%                               scored token by token, never generated}
%                               -> the band-one sentence and both subtitles. OK
%   04b-representation.tex:200-202 "The ablation runs for $V_{\textrm{pure}}$
%                               (headline), a matched-dimension random subspace
%                               (the null), and the cell-line slab (the positive
%                               control)" -> the three dots are named in exactly
%                               those three roles, in that order of size.  OK
%   04b-representation.tex:254  "3.6--18.6x" -> layer 4 IS the 18.6.        OK
%   04b-representation.tex:260  "9--37% of what removing cell line costs at
%                               layers 2--9" -> layer 4 IS the 37%.        OK
%   04b-representation.tex:286-288 "No intervals are drawn ... normal-
%                               approximation standard errors over the 60
%                               prompts" -> inherited here.                OK
%   03-Apparatus.tex:249-252    "its expressed panel genes in descending order,
%                               terminated by \ENDCELL. The prompt supplies cell
%                               line, drug, dose, mechanism and the control
%                               cell's sentence" -> the five prompt fields, in
%                               that order, and the terminator.            OK
%   The layer-4 cell-line KL, 4.180, is printed in NO sentence of the thesis;
%   the section quotes only the ratios built from it (04b:258-262). It is a new
%   reading of a stored field and hook.json flags it as such. Precedent:
%   figs/slab.json does the same for variance_share.context.
%
% DESIGN DECISIONS, each with its reason
% --------------------------------------
%  1. BOTH STRIPS ARE DRAWN BY ONE MACRO, \tokstrip, taking only the y of the
%     box bottoms. Two hand-written strips would be two chances to differ, and
%     the single claim the figure has to carry is that they do NOT differ. The
%     macro makes "identical box for box" a property of the source, not of my
%     care.
%  2. THE FIVE PROMPT FIELDS AND THEIR ORDER ARE NOT MINE. They are
%     03-Apparatus.tex:250-252 verbatim, and the same five in the same order as
%     figs/licensing.tex's strip, so the two figures read as one object seen
%     twice. A sixth field would be unsourced and would break that.
%     PASS 9: the two strips now also share a GLYPH. A reviewer found that the
%     words agreed and the drawing did not -- hook drew rounded panelbg boxes
%     with rounded gutters and licensing drew square-cornered WHITE cells
%     sharing edges, so a reader met the same object twice and was shown two
%     objects. hook's treatment was judged the better mark and licensing.tex
%     adopted it (its cell/.style now carries fill=panelbg, rounded
%     corners=1pt, and its drug cell \tokdrug's accent 0.4pt border). Nothing
%     in THIS file changed for it; the cell edges still differ, because the two
%     strips are 8.20cm and 7.20cm and each distributes its own slack.
%  3. THE ELLIPSIS BOX is the mitigation for this figure's largest honesty
%     exposure. The real target is a cell sentence of hundreds of gene tokens;
%     four boxes is a schematic. "no token count is implied" appears in the
%     caption and in the ledger. Rejected alternative: drawing twenty boxes,
%     which implies a count just as firmly and only moves the lie.
%  4. THE DRUG BOX IS THE ONE TINTED PROMPT BOX (accent!12, accent rule). It is
%     the one place the prompt names the drug in plain text -- the fact that
%     makes sec:methods-probe's positive nearly free -- carried forward here so
%     the reader sees the hook deleting downstream of it. It is also the only
%     accent inside either strip, so the strips themselves stay furniture.
%  5. THE HOOK RULE SITS ABOVE THE PASS-TWO STRIP, not below it, so the eye
%     meets the intervention before it meets the boxes it acts on. The RULE is
%     what says "everywhere along here": it spans the strip's full width,
%     prompt half included, which is the correction to the reader's default
%     assumption that an intervention on generation touches only what is
%     generated.
%     REPAIRED, PASS 9: the nine descending TICKS are gone, replaced by one
%     Latex arrow into each half at 0.5pt, stopping 0.06cm above the box tops.
%     Two faults, found independently by both reviewers. (a) A tick per BOX is
%     not a tick per POSITION. "control cell sentence" is one box and hundreds
%     of positions; the ellipsis box is an unknown number. Counting a comb and
%     mapping it to positions is the first thing a reader does with a comb, and
%     it gave the wrong answer -- while the caption simultaneously had to say
%     "the token strip is schematic and implies no token count". (b) The ticks
%     were accent at nearly the rule's own weight and each terminated exactly
%     on its box's top border with no gap; for the DRUG box, whose border is
%     also accent, tick and border fused into one continuous stroke, so the one
%     cell that must be found first sprouted a flagpole and stopped reading as
%     a highlighted field. Marks running into rules uninterrupted is precisely
%     what the house rule forbids. The comb also read as a categorical axis,
%     which is exactly what the WHAT COMES OUT panel two bands below is. Two
%     arrows say "down into both halves" and assert no count; the words "at
%     every position, prompt and target alike" under the strip carry the claim.
%  5b.NO SEPARATE TAG ON THE RULE. One reading "the hook" was drawn at
%     (8.70,7.40) and taken out on the next build: PASS TWO's subtitle now ends
%     "...the hook is on one layer's residual stream" on baseline 7.58 and runs
%     to x = 13.5, so a tag 0.18cm beneath it collided with it and named the
%     same object a third time. The rule is named above by that subtitle and
%     below by the accent formula.
%  6. THE RULE RUNS UNBROKEN ACROSS THE 1.20cm GAP between the prompt strip and
%     the target strip. The gap is a reading device (it is where the divider
%     is); the hook does not know about it. Breaking the rule there would draw
%     the exact exemption the figure denies.
%  7. THE TWO BRACES SPAN THE TARGET SIDE ONLY. That the brace covers the
%     target and no more is the whole answer to "why is the KL taken over the
%     target tokens", and it is drawn rather than said.
%     PASS 9: they now start at the first target box's own left edge, 9.20, and
%     not at the prompt/target divider, 8.70. The caption says "the braces span
%     the target tokens only" and a reader is invited to test that against
%     "prompt boxes included" for the hook rule immediately above, so a brace
%     that began 0.50cm before the first token was contradicting its own label.
%     They are rulegrey, not accent:
%     they are furniture that delimits, and the colour budget is spent on the
%     hook, the KL arrow and the in-slab component.
%  8. THE KL ARROW IS THE ONLY STROKE ABOVE 1.0pt IN THE FIGURE (1.1pt).
%     Nothing else exceeds 1.0pt, so weight alone says which mark is the claim.
%     It runs between the two brace LABELS' baselines, so it visibly joins two
%     readings of the same tokens rather than two sentences.
%     REPAIRED, PASS 9: it was DOUBLE-headed and is now single-headed, tail at
%     the pass-one readout and head at the pass-two readout, matching the
%     direction its own label carries. KL(clean || ablated) is asymmetric; a
%     two-way head says "the difference between", i.e. that the two passes are
%     interchangeable, which is the one thing the measure is not. It was also
%     the figure's heaviest stroke saying the opposite of its label. No
%     double-headed arrow exists anywhere in the corpus: arrows.meta is used
%     only as -{Latex} for a displacement and -{Stealth} for a pipeline edge.
%  9. THE VECTOR PANEL DRAWS THE SLAB AS A BAND WITH THICKNESS, not as a line.
%     A line would say "one direction"; the slab is ten. Nothing in that panel
%     is to scale (a true 10-of-2048 slab is 0.5% of the squared norm and
%     cannot hold a labelled decomposition -- figs/slab.tex draws that true
%     angle, once, and this panel does not pretend to repeat it).
%     REPAIRED, PASS 9, AND THIS WAS THE PANEL'S WORST FAULT: THE RATIO WAS
%     THE WRONG WAY ROUND. h was drawn at 30 degrees above the band, so the
%     REMOVED component (hV)V^T measured 2.90cm and the SURVIVING remainder
%     h-(hV)V^T measured 1.70cm -- on the render, 228px against 130px, the
%     ten-dimensional slab drawn 1.75x LONGER than the 2038-dimensional
%     remainder, under a subtitle reading "ten of the 2048 directions". Cold,
%     the slab read as the dominant subspace and the residual as a leftover,
%     the exact inversion of the section's finding. h is now steep: component
%     0.48cm, complement 1.90cm, 4.0:1 the other way. It is deliberately NOT
%     the true 14.3:1 (86 degrees, cos^2(86) = 0.0049 = 10/2048): at that angle
%     the component is 0.13cm, shorter than its own arrowhead. slab.tex inset
%     one draws that geometry true and this panel defers to it in the caption.
%     Three smaller repairs in the same pass. (a) (hV)V^T was the only one of
%     the three vectors drawn HEADLESS, so in a decomposition the component
%     that is removed was the one not drawn as a vector and read as a baseline;
%     it now carries the same Latex head as its siblings, and its label carries
%     its fate, "-> 0", so the zeroing is a mark and not only the sentence
%     underneath. (b) Its label's 0.1cm rulegrey stub leader is deleted: it was
%     shorter than the 0.3cm threshold that triggers a leader and in the wrong
%     colour. (c) The construction is centred on the band's midpoint at 3.40,
%     because at 2.30 it crowded into the band's left third with 3.1cm of empty
%     band to its right.
%  9b.THE BAND IS panelbg WITH A rulegrey OUTLINE, not accent!15.
%     PASS 9, a cross-figure fix. Across this slate a PALE accent tint meant
%     "raw, still confounded" in fig:purify and "purified" here, while SOLID
%     accent meant V_pure in fig:purify and in the published slab.tex but "the
%     component being removed" here: the tint semantics were inverted between
%     two figures a few pages apart. And inside this figure alone, accent!12
%     (the drug field, both strips) and accent!15 (this band) sampled as
%     (225,234,242) and (218,229,238) -- a 3% tint difference that will be
%     indistinguishable on press, two meanings in one apparent colour. panelbg
%     with a rulegrey outline is slab.tex's own treatment for "a region of the
%     stream" (its \slabdisc, slab.tex:83-90), it leaves accent!12 uniquely
%     "the drug field" across licensing and hook, and it keeps solid accent for
%     the slab itself across the slate.
%     ONE RESIDUAL DIFFERENCE, AND IT IS DELIBERATE: fig:purify draws
%     V_pure as a SOLID accent bar and this figure draws it as a cream region.
%     The two are doing different jobs. In fig:purify V_pure is one subspace
%     among three and is the object every causal claim is made on, so it takes
%     the accent, exactly as slab.tex's published wedge does; here it is the
%     backdrop a vector is decomposed against, and an accent component has to
%     be legible ON it, which a solid accent fill would make impossible. The
%     tie between the two figures is the direct label, which is accent
%     $V_{\textrm{pure}}$ in both. Giving this band an accent OUTLINE instead
%     was considered and rejected: at panelbg fill with an accent 0.35pt border
%     it is the drug box's own treatment (accent!12 fill, accent 0.4pt border)
%     to within a tint, and one glyph would then mean two things again.
%  9c.ONE NAME FOR ONE OBJECT. The panel said V, the arrows said (hV)V^T and
%     the band said V_pure -- three names for the same ten directions, in a
%     figure that sits a page from fig:purify where V_pure is built. The
%     subtitle now states the identity once: "V = V_pure: ten of the 2048
%     directions".
% 10. THE RIGHT-ANGLE MARK AT THE ORIGIN is load-bearing, not decoration. It is
%     what distinguishes "the slab component is set to ZERO" from "the slab
%     component is reduced". Without it the vertical arrow is just a shorter
%     arrow.
% 11. THE READOUT AXIS IS LINEAR AND CARRIES NO BARS. It is one layer, so the
%     reason mechanism.py forbids a bar on a log axis (a bar's length is set by
%     wherever the axis is cut) does not arise -- and there are no bars in any
%     case, only three dots. Rejected alternative: seven layers on this axis,
%     which would rebuild fig:mechanism panel (b) inside a mechanism figure and
%     make the readout compete with the mechanism for the reader.
%     REPAIRED, PASS 9: EVERY DOT NOW CARRIES ITS VALUE. None of 1.563, 0.084
%     and 4.180 was printed on the drawing, and \kx{0.084} = 8.362 against a
%     zero tick at 8.20, so the null sat 1.6mm from zero and, under a 2.1pt
%     halo and a 1.4pt fill, was visually AT zero -- while the strip below
%     asserted "18.6x the null". A reader who takes the null for zero takes
%     18.6x for a claim they cannot check. Both reviewers proposed the same two
%     remedies, printing the values or moving to a log axis with a joining
%     stem; the first is taken, because it is the document's own no-legend
%     rule, it makes the ratio checkable from the marks, and it keeps the axis
%     comparable in kind with tab:workspace's nats column. The three
%     descriptions were also three near-full-width stacked lines in an order
%     that was neither the axis order nor a reading order -- the top line named
%     the rightmost dot, and the middle line, spanning crop x 275-1630, sat
%     directly over the accent dot at 770, so "the null" attached to the drug
%     slab on first pass. They are now short and value-first, each near its own
%     dot; the long descriptions live in the caption, which already carried
%     them. The three dots already sat over 2.1pt white halos in ladder.tex's
%     order (spine, halos, fills, labels) and still do; a reviewer reported
%     them as halo-less and that finding was overruled on the source.
% 12. NO MARK JOINS THE KL AXIS TO ANY ACCURACY OR SUBSPACE-FRACTION SCALE, and
%     no double-headed arrow spans two rules anywhere below the divider. Those
%     are different instruments with different comparators; fig:comparators is
%     the figure that owns that argument.
% 13. LAYER 4 IS THE ONE LAYER DRAWN because it is the slab's strongest showing
%     of the seven against the null (18.6x, the top of the thesis's 3.6-18.6x)
%     and the top of the 9-37% band against cell line. Drawing the best case
%     and still showing the accent dot sitting at a third of the ink dot is the
%     honest form of the claim. Layer 12 was the alternative and is not drawn:
%     there the drug and cell-line costs are nearly equal (0.133 against 0.135,
%     98%), which is the section's own named exception and would read as the
%     rule if it were the only layer shown. The caption says both, and points
%     at tab:workspace for all seven. hook.json records layer 12's three values.
% 14. THE TOP-RIGHT NOTE SAYS "AT ONE LAYER". The word "every" in "at every
%     position" ranges over the SEQUENCE, not over depth, and a reader who
%     meets "layer 4" only in the readout panel has no way to tell which. The
%     misreading -- that the slab is projected out at every layer as well as at
%     every position -- is exactly the kind this figure exists to prevent, and
%     no mark can carry it, so it is stated in words. It is a design fact from
%     the method (04b:164-171, one measurement per layer, one row per layer in
%     tab:workspace; 04b:191-195, the hook itself), not a stored quantity, and
%     hook.json records it as such. Set quietink at \scriptsize, because it is
%     furniture; that it also fills the figure's one empty quadrant is a
%     consequence and not the reason.
%
% WHAT SIX BUILDS CHANGED, with the measurement that forced each
% -------------------------------------------------------------
% (a) EVERY TEXT NODE IS anchor=base, so each y below IS a baseline. Build 1
%     used anchor=north west with an assumed 0.34cm node height. The real
%     height of a \footnotesize node at TikZ's default inner sep (0.3333em,
%     2.7pt a side) is 12 + 5.4 = 17.4pt = 0.61cm, nearly twice that, and the
%     consequences were visible: "the true sentence is teacher-forced, never
%     generated" put its descenders through the pass-one box tops, and the hook
%     rule at y = 7.52 ran clean through the middle of "the same prompt, the
%     same sentence, the same weights". Baseline anchoring removes the guess:
%     ink sits about 0.20cm above and 0.06cm below a \scriptsize baseline, and
%     0.25/0.08 at \footnotesize. frames.tex records the same fix for the same
%     reason, plus a second one -- a node with align= or text width builds a
%     minipage whose first line follows the AMBIENT \baselineskip, so its ink
%     moves with the prose before the float. No node here uses either.
% (b) THE PROMPT/TARGET DIVIDER MOVED OFF THE BOX EDGE. Build 1 put it at
%     x = 7.80, which is exactly the right-hand border of the "control cell
%     sentence" box: a 0.3pt rulegrey line drawn on top of a 0.35pt rulegrey
%     box border is invisible as a separate mark, and the requirement is that
%     the braces stop VISIBLY at it. It now sits at x = 8.40, the midpoint of
%     the 1.20cm gap between the strips, where nothing else is drawn, and it is
%     set as `ink, 0.5pt, densely dotted' -- this document's mark for a mere
%     coordinate worth naming, which is exactly what the boundary is. Both
%     braces now begin at 8.40 and their left tips land on it.
% (c) THE BRACE AMPLITUDE WENT FROM 2.5pt TO 3.5pt AT 0.35pt. At 2.5pt and
%     0.3pt, under a rounded-corner box border 0.10cm above it, the brace read
%     as a second box outline rather than as a brace; slab.tex and ladder.tex
%     both use 3.0-3.5pt for braces that have to be seen.
% (d) THE THREE READOUT LABELS WERE STAGGERED BY 0.50cm, not 0.33cm. At 0.33cm
%     "the purified drug slab" and "a matched-dimension random subspace --- the
%     null" overlapped by about 0.08cm of ink, because both labels are wide
%     enough to share x and the node boxes are taller than a line. 0.50cm gives
%     0.24cm of clear space between successive label blocks, which is also what
%     lets the two leaders leave their dots at distinguishable angles.
% (e) "ten of the 2048 directions" AND THE VECTOR PANEL'S CLOSING SENTENCE
%     overlapped outright in build 1 (the first sat at y 1.05-1.35, the second
%     at 0.95-1.15). The label was DELETED rather than moved: the panel's own
%     subtitle already reads "$V$ holds ten of the 2048 directions", so the
%     node under the band was saying a second time what the subtitle says
%     first, and deleting it also freed the space the closing pair needed.
% (f) THE STRIP WAS WIDENED FROM 13.20 TO 13.55 and the KL arrow moved from
%     13.75 to 14.05. At the narrower width the top-right quadrant, roughly
%     4.1 x 1.7cm, held nothing at all, and the KL label was the only mark
%     right of the strip. The right column now carries the "one layer" note as
%     well (decision 14) and band one reaches x = 16.5, which is where band
%     two's readout already ended.
% (g) THE IN-SLAB COMPONENT'S LABEL GOT A LEADER. $(hV)V^{\top}$ names a
%     segment that lies INSIDE the slab band, so the nearest place its label
%     can sit is 0.28cm above the band's top rule and 0.44cm above the mark
%     itself -- well past the ~0.3cm at which this document's rule requires a
%     leader. Built without one it read as floating in the interior of the
%     triangle. It now has a 0.14cm accent!45 hairline from just above the
%     segment to just under the label, which is the same device comparators.tex
%     uses for its 0.121 mark.
% (h) THE 18.6x / 37% LINE MOVED FROM quietink TO ink. Under the axis it was
%     the third of four grey centred lines and read as a fourth piece of axis
%     apparatus. It is a reading of the three marks, so it takes the register
%     this document gives an estimate rather than the one it gives furniture.
% (i) THE AXIS SPINE WAS CUT FROM 16.90 TO 16.60. The map's range top is 4.5
%     nats, but the largest mark is 4.180 and the last tick is 4, so a full-
%     range spine left 0.97cm of empty rule past the last tick and read as an
%     axis that had run out rather than as headroom. No coordinate moved.
% (j) THE PASS-ONE BOUNDARY RULE GREW A TAIL, from 8.72 down to 8.55. Its
%     label sits to its left on the brace label's baseline, and with the rule
%     stopping 0.14cm under the boxes the label read as a caption of the prompt
%     half rather than as a pointer at the boundary. The pass-two copy keeps
%     the symmetric 0.14/0.14, because it carries no label.
%
% LAYOUT NOTES, all forced by a build
% -----------------------------------
%   * fullwidth = textwidth 142mm + marginparwidth 28mm + marginparsep 4mm
%     = 174mm. (preamble.tex's own comment beside \newenvironment{fullwidth}
%     says 171mm and is stale arithmetic; the geometry keys give 174mm.) The
%     box is PINNED at 17.30 x 10.00cm = 173.0 x 100.0mm, 1.0mm inside the
%     measure. Rightmost ink is the axis spine at 16.90.
%   * the spec's token x-edges (0.00/1.20/2.20/3.00/4.30/7.30) do not hold
%     their labels at \scriptsize: "cell line" measures about 1.27cm of ink and
%     the box was 1.20cm, and "mechanism" the same against 1.30cm. The five
%     boxes were rewidened to 1.55/0.85/0.85/1.70/3.25, keeping the order and
%     keeping "control cell sentence" the widest field, and the prompt strip
%     now ends at 8.20 rather than 7.30. Everything keyed to the boundary (both
%     dividers, both braces, the target strip, the KL arrow) moved with it.
%     The five widths are set against figs/licensing.tex's measured ink for the
%     SAME five labels at the same size (cell line 0.954, drug 0.616 bold, dose
%     0.579, mechanism 1.424, control cell sentence 2.511; total 6.084cm), so
%     the clear space each side runs 0.117 to 0.370cm and no cell is near the
%     0.06cm floor. A first pass gave mechanism only 1.55cm, i.e. 0.063 a side,
%     while control cell sentence had 0.445; 0.15cm moved from the last cell to
%     the fourth evens that out. \ENDCELL's cell keeps 0.22 a side.
%   * the target strip is four boxes, so the strip is nine boxes and the hook
%     carries nine ticks. \ENDCELL is \texttt{[END\_CELL]}: ten inconsolata
%     characters, about 1.41cm at \scriptsize, so its box is 1.85cm.
%   * the prompt/target GAP is 1.00cm (8.20 to 9.20) and the boundary rule sits
%     at its midpoint, 8.70. Both braces begin there and their left tips land
%     on the rule, which is what "the brace stops dead at the divider" means
%     when it is drawn rather than described.
%   * nothing is scaled. 1cm here is 1cm on the page; there is no scale= key
%     and no \resizebox. style_v2.py records that v1's figures were silently
%     scaled to 0.887 by LaTeX and that this, not the drawing, read as "off".
%   * the pgfmath trap: \kx expands to a BRACED group, so \kx{1.563}+0.05 is
%     not a coordinate. Where a label needs a shim it is applied as an xshift
%     on the node, or the shim lives inside the braces. (gate.tex's header
%     records the same trap costing a compile.)
%   * \foreach prints its raw token, so the tick labels are given in the
%     \v/\t value/label pair form. Here every tick is a non-negative integer so
%     the pair form is not strictly needed, but it is the corpus idiom and it
%     keeps the axis editable without a surprise.
%
% TYPE
% ----
%   Ambient font=\small on the picture; every node sets its own size. The
%   dominant size is \scriptsize (8pt), with \footnotesize (10pt) only for the
%   four panel subtitles, the band-one sentence and the vector panel's two
%   closing lines. Nothing is set at body size. Registers: \sffamily bold
%   accent = a panel tag; roman accent = the intervention; roman ink = the
%   object being scored or the comparator; roman quietink = furniture and
%   annotation; \ttfamily = only \ENDCELL, which names a real token.
%
% ===========================================================================
% PASS 10 -- 2026-08-18 -- THE STYLE-CONTRACT REBUILD
% ===========================================================================
% EVERYTHING ABOVE THIS LINE IS LEFT INTACT and is the design record of passes
% 1-9. This section supersedes, BY NAME, only the claims above that this pass
% made false; a decision above that is not named here still stands, and the
% reasoning that produced it is still the reasoning this figure is audited
% against. NO NUMBER CHANGED IN THIS PASS: 1.563, 0.084, 4.180, n = 60, rank
% 10, residual width 2048 and layer 4 are byte-identical to pass 9, and the two
% ratios 18.6x and 37% were not altered but MOVED -- off the drawing and into
% the caption, which is where this document's contract puts a result.
%
% WHY THIS PASS EXISTS
% --------------------
% The supervisor's verdict on the 4.7-4.12 slate: the figures look "sloppy and
% childish", they must be CORRECT, and every axis must define what it measures.
% A binding style contract (S1-S25) and a cross-figure code (CF1-CF14) were then
% written against measurements of this document's own two reference figures --
% 4.11, figs/mechanism.py, the BODY standard, and A.3, figs/slab.tex, the
% appendix plate -- and the ruling was: in the body, take 4.11. Measured against
% it, this figure was the worst offender of the five.
%
%                          pass 9            4.11        pass 10 (measured)
%   type sizes >7pt        9                 2           2 chosen + 1 artefact
%   embedded faces         9                 1           6, four of them math/tt
%   italic characters      34%               0%          1.8%
%   stroke weights         9 (0.17, 0.68     --          5
%                          used once each)
%   all-caps sans banners  6                 0           0
%   ink at 200dpi          8.55%             5.3%        5.98%   (cap 6.0)
%   drawn box              173.1 x 97.0mm    --          173.99 x 99.95mm
%
% THE CORRECTNESS DEFECT, WHICH IS WHY THE FIGURE WAS REOPENED AT ALL
% -------------------------------------------------------------------
% THE READOUT AXIS IS NOW LOGARITHMIC. Design decision 11 above chose a linear
% nats axis, and pass 9 mitigated its consequence by printing every dot's value.
% BOTH ARE SUPERSEDED. The mitigation was not enough, and the header itself said
% so in as many words: on the linear map \kx = 8.20 + v/4.5*8.70 the null at
% 0.084 landed 0.162cm -- 1.6mm -- from the zero tick, under a 2.1pt halo and a
% 1.4pt fill, so the null dot WAS the origin at print size while the strip below
% asserted "18.6x the null". A ratio the drawing renders as infinite is not a
% mark a reader can check, whatever is printed beside it. Both quantities this
% figure quotes are RATIOS; on a log axis a horizontal gap IS a ratio, so the
% mark now means what the sentence means. 4.11 panel (b) published exactly this
% solution for exactly these numbers -- log axis, open and filled dots, the
% ratio read as a gap -- and this figure now uses it.
%   NEW MAP  \kx#1 = 8.40 + (ln(#1) - ln(0.05))/(ln(6.0) - ln(0.05)) * 8.98
%            domain [0.05, 6.0] nats, spine 8.40 to 17.38, labelled ticks at
%            0.1, 0.3, 1, 3
%   0.084 -> 9.3731 , 1.563 -> 14.8569 , 4.180 -> 16.7020
%   The null-to-slab gap is 5.4838cm and IS 18.6x. The slab-to-cell-line gap is
%   1.8452cm and IS 37% read the other way. Neither is annotated: the caption
%   carries both, which is where the contract puts a result.
%   The domain is deliberately wider than the data at BOTH ends, so no mark sits
%   on the cut and no bar is possible; 4.11(b) does the same, its ticks running
%   0.01-1.00 inside a taller axis. There are still no bars and there must not
%   be: on a log axis a bar's length is set by wherever the axis is cut.
%   SUPERSEDED: the linear map (0.084 -> 8.3624, 1.563 -> 11.2218,
%   4.180 -> 16.2813); decision 11's argument that "one layer, so a log axis is
%   not needed"; and note (i)'s spine crop, which no longer applies.
%
% THE FOUR MARKS AND LINES THAT WERE DELETED, EACH WITH ITS REASON
% ----------------------------------------------------------------
% (i)   THE TWO SHORT HOOK ARROWS (pass 9's replacement for the nine-tick comb)
%       are gone. They implied a count of two and the caption had to deny it in
%       a sentence of its own ("The two arrows off the rule mark that it acts in
%       both halves and imply no count of positions"). A drawing that needs its
%       caption to unsay it is drawing the wrong mark. The RULE's span is the
%       "everywhere" claim -- decision 5 above is right and survives -- and the
%       accent formula now sits ON the rule, 0.20cm above it, instead of 1.1cm
%       below the strip. This also empties the figure of every arrowhead except
%       the comparison arrow's two and the vector panel's three, which is what
%       S11 asks for: one arrowhead, one meaning.
% (ii)  THE ITALIC EPIGRAM "nothing is generated; the two passes differ only in
%       the hook" and the 17.30cm rulegrey divider that existed only to carry it
%       are deleted. It was 9.96pt Pagella-Italic quietink centred on the
%       figure's own x-centre -- the LARGEST type on the drawing, 25% larger
%       than the 7.97pt at which every measured number was set -- and its first
%       clause restates the caption's opening words, "Nothing is generated:".
%       It is the caption's sentence and the caption already had it.
% (iii) THE PRINTED RESULT "18.6x the null; 37% of what removing cell line
%       costs" is deleted from under the axis. A result may not sit in the
%       axis-label lane, and it was the third of four stacked centred lines
%       there. NOTE (h) ABOVE IS SUPERSEDED: moving that line from quietink to
%       ink so it would stop reading as axis apparatus was the wrong repair --
%       the line does not belong under the axis in any colour. Both ratios are
%       in the caption, which now also carries the qualifier the old caption
%       dropped: "9-37%" is the range over layers 2-9 ONLY; the true seven-layer
%       range of kl.pure_drug[0]/kl.cell_line[0] is 8.90% (layer 9) to 98.45%
%       (layer 12).
% (iv)  "the marks carry no interval; the caption says why" is deleted. It
%       restated the caption's own bold "no interval is drawn" and it was
%       italic. It is replaced by the single figure-foot declaration S17
%       requires, set \scriptsize quietink on the figure's x origin, which NAMES
%       the construction instead of pointing at the caption.
% (v)   "generation would begin here", the prompt/target divider's label, is
%       deleted, and this was the one deletion made for the ink budget rather
%       than for a defect. Recorded here because it was load-bearing and a
%       future pass may want it back. Its claim is drawn three other ways --
%       panel (a)'s head is "pass one, clean", both braces read "scored tokens"
%       so the prompt half is visibly not scored, and the caption's lede opens
%       "Nothing is generated" -- and the divider itself survives as what it
%       always was, a piece of the strip's structure. Cost measured: 27mm2 of
%       ink, against a 6.0% ceiling the figure clears at 5.98%.
%
% WHAT REPLACED THE SIX BANNERS AND THE FOUR ITALIC SUBTITLES
% -----------------------------------------------------------
% Every \sffamily\scriptsize\bfseries accent banner is gone, and with it the
% only sans face in the figure. The apparatus face belongs to the margin column,
% not to a figure read inside a running argument; A.3 may use it because A.3 is
% a reference plate a reader stops at. Panel heads are now 4.11's form exactly
% -- a lowercase parenthesised letter and a sentence-case descriptive title,
% \footnotesize roman ink, on the panel's own x origin, no colon, no bold, no
% caps:
%     (a) pass one, clean
%     (b) pass two, the slab projected out
%     (c) what the hook does to one vector
%     (d) what comes out
% The four \footnotesize\itshape quietink subtitles under them are deleted.
% Their content was distributed, not lost:
%   "the true sentence is teacher-forced, never generated"  -> already the
%       caption's opening clause; it was duplicating it (S18).
%   "the same prompt, the same sentence, the same weights"  -> now the drawn
%       tie between the strips (below).
%   "the hook is on ONE layer's residual stream"            -> now the axis
%       block's scope line, where the layer is named anyway.
%   "V = V_pure: ten of the 2048 directions"                -> the identity is
%       now the band's own direct label; the rank moved to the caption.
%   "layer 4: mean KL from the clean pass, over the target tokens, 60 prompts"
%       -> now the axis block itself, split into a quantity line and a scope
%       line, which is the position S5 requires.
% DECISION 14 ABOVE IS PARTLY SUPERSEDED: the "at one layer" fact is still
% stated in words, because no mark can carry it, but it is no longer a clause of
% PASS TWO's subtitle -- there is no subtitle -- and lives in the axis block's
% second line, the place S5 reserves for scope.
%
% THE TIE BETWEEN THE STRIPS, WHICH IS THE ONLY NEW TEXT IN THE FIGURE
% ---------------------------------------------------------------------
% The two strips are emitted by one macro and are therefore provably identical
% box for box -- the figure's central claim, and its best-executed feature.
% Nothing on the drawing said so. Worse, the two braces beneath them carried
% DIFFERENT labels ("per-token log p, read here only" against "the same tokens,
% read again"), which sends a reader looking for a difference between two
% objects whose sameness is the point. Both braces now carry the SAME label,
% "scored tokens", and one \scriptsize roman ink line reads "the same sentence,
% box for box". It sits UNDER both strips, on band one's own x origin and on the
% same baseline as pass two's brace label, so it is a statement about a pair the
% reader has just seen and it costs the drawing no extra row of ink.
%
% THE COMPARISON IS NO LONGER AN ARROW WITH ONE HEAD
% ---------------------------------------------------
% DECISION 8 ABOVE AND ITS PASS-9 REPAIR ARE SUPERSEDED. Pass 9 made the KL mark
% single-headed, reasoning that KL(clean || ablated) is directed and asymmetric
% and that a two-way head would say the passes are interchangeable. That
% reasoning is about the MEASURE; the MARK is about the two passes, and a naive
% reader read the single head as pass one FEEDING pass two -- generation, the
% precise misreading this figure exists to kill, drawn by the figure itself in
% its heaviest stroke. The two passes are independent scorings of one fixed
% sentence by one set of weights: that is a comparison, and S11 gives a
% comparison a double head or a brace. It is now a DOUBLE-headed ink arrow at
% 1.0pt, and the single accent arrowhead is reserved for an operation acting in
% a direction. It is no longer the figure's heaviest stroke either -- the 1.1pt
% was a tenth-of-a-point singleton in a nine-weight vocabulary -- and its label
% is the registry name, "KL from the clean forward pass", the same words the
% axis in panel (d) carries, because it is the same quantity and S19 gives an
% object exactly one name.
%
% THE AXIS DEFINITION, WHICH IS THE SUPERVISOR'S FIRST COMPLAINT
% --------------------------------------------------------------
% This figure was the only one of the five that already named quantity and unit
% in axis-label position, and note (h) plus decision 11 built that. The WORDS
% survive; the position, the rank and the stack around them change.
%   LINE 1  \footnotesize roman ink, on the axis's own left end, x = 8.40:
%           "KL from the clean forward pass (nats), logarithmic scale"
%   LINE 2  \scriptsize quietink, same x:
%           "mean over the target tokens of 60 teacher-forced prompts, layer 4"
% "Mean KL divergence, clean -> ablated (nats)" is RETIRED. It was centred and
% bold under the rule; more importantly it was a SECOND name for 4.11(b)'s "KL
% from the clean forward pass", and CF3 gives the quantity one name across the
% slate. The old panel subtitle's "layer 4 ... over the target tokens, 60
% prompts" is folded into line 2, so the scope is stated once instead of twice.
%
% THE READOUT IS THREE ROWS, NOT ONE, AND WHY
% -------------------------------------------
% S22 requires that within a row every direct label share one baseline and that
% no label cross the drawn measure. Three labels naming three objects at one
% layer -- "0.084 - matched-dimension random slab" (5.2cm), "1.563 - the
% purified drug slab" (4.2cm) and "4.180 - cell line, the positive control"
% (5.4cm) -- cannot share one baseline: the last two marks sit 1.85cm apart on
% the log map. Staggering labels off one rule is what S22 forbids; giving each
% object its OWN row is not. Each row therefore carries exactly one mark and one
% label on one baseline; the three rows share one log x axis; and they are
% ordered by value with the largest on top, so vertical position carries the
% same order as horizontal position and can assert nothing the axis does not.
% That ordering is 4.11(b)'s own, the filled slab dot above the open null dot.
% The rows sit 0.38, 0.82 and 1.26cm above the spine so the three read as marks
% over one scale. Each label is 0.26cm from its own mark, just inside the 3mm at
% which this document requires a leader, so no leader is drawn. THE LABEL SIDE
% DIFFERS BY ROW AND IT HAS TO: the two large marks are 0.7 and 2.5cm from the
% right end of the measure and their labels can only go left, while the null at
% 9.37 has 0.97cm to its left and its label can only go right. NO DROP LINES
% JOIN THE MARKS TO THE SPINE: a hairline from a dot down to a horizontal axis
% reads as a bar, and a bar on a log axis is forbidden in this document.
%
% MARKS, WEIGHTS, TYPE AND PROSE
% -------------------------------
% MARKS, S9/CF5, three glyphs and no more:
%   filled disc, r = 1.4pt, accent   the purified drug slab -- the object under
%                                    test, and the only accent mark in panel (d)
%   filled disc, r = 1.4pt, ink      cell line, the positive control
%   open disc,  r = 1.4pt, white     matched-dimension random slab, the null.
%               fill, 0.8pt ink edge This is 4.11(b)'s own encoding for this
%                                    same object (mechanism.py: mfc="white",
%                                    mew=1.1, color=style.INK), and adopting it
%                                    is CF5's stated requirement. The 0.8pt edge
%                                    is the mark vocabulary's own weight, not a
%                                    path weight, and it is the only 0.8pt here.
%   The three 2.1pt white halos are DELETED together with the two leaders they
%   existed to cover. Pass 9's overruling of the "no halo" finding stands as a
%   fact about pass 9; the halos are simply no longer needed, because no leader
%   now runs into a dot.
% WEIGHTS, S10. Measured in the built PDF: 0.30 x26, 0.40 x5, 0.80 x1, 0.90 x1,
% 1.00 x4. Five weights, against pass 9's nine.
%   0.30pt  furniture: every token-box border, both braces, the slab band's
%           outline, the right-angle mark, the two dashed construction lines,
%           the four axis ticks
%   0.40pt  the drug cell's accent border, both prompt/target dividers, the
%           axis rule
%   0.80pt  the open disc's edge (S9's own weight for that glyph; see above)
%   0.90pt  the hook rule -- USED ONCE, and justified here as S10 expressly
%           permits: it is a distinct kind of object, the intervention, the one
%           thing that differs between the two passes. It is accent and not
%           black, so it cannot be confused with the chance/gate/margin register
%           S6 reserves for a black solid 0.9pt rule, and this figure draws no
%           threshold at all.
%   1.00pt  panel (c)'s three vectors and the comparison arrow
%   RETIRED: 0.35pt (box borders and braces, folded into 0.30), 0.50pt (the
%   dotted dividers, see below), 1.10pt (the KL arrow), and the 0.17pt and
%   0.68pt singletons, which were never chosen at all -- they were the Latex
%   arrow tips' own stroke, and the tips are now declared line width=0pt so they
%   are pure fills and contribute no stroke weight.
% THE PROMPT/TARGET DIVIDERS are no longer "ink, 0.5pt, densely dotted".
%   NOTE (b) ABOVE IS SUPERSEDED. S7 reserves the densely-dotted 0.5pt stroke
%   for a zero rule, a magnification leader or a projection line, and requires a
%   dotted rule's label to be "0", to be a leader, or to be absent; this one
%   carried "generation would begin here", which is none of those. The line is
%   what it always was -- a piece of the strip's structure that delimits two
%   halves -- so it is now rulegrey 0.40pt solid, a strip weight, still at 7.94
%   in the middle of the empty 1.00cm gap where nothing else is drawn. Note
%   (b)'s actual finding, that a rulegrey line drawn ON a rulegrey box border is
%   invisible as a separate mark, is why it stays at mid-gap.
% TYPE, S1. Measured in the built PDF: 7.97pt x587 characters and 9.96pt x165,
% plus three exempt classes -- 6.85pt x20 (\ENDCELL, Inconsolata, the \texttt
% token name S1 exempts), 5.98pt x8 and 6.23pt x3 (subscripts and the \top
% superscript, which S1 exempts), and 8.30pt x10.
%   THE 8.30pt SPANS ARE NOT A THIRD CHOSEN SIZE and are recorded here so an
%   auditor need not rediscover them: they are ten glyphs -- \leftarrow, two
%   minus signs, six parentheses and \to -- set from CMSY10 and CMR10 inside the
%   two formulas $h \leftarrow h-(hV)V^{\top}$ and $(hV)V^{\top} \to 0$.
%   mathpazo substitutes Palatino for math letters but leaves the CM symbol
%   fonts in place and scales them by 1.041 to match Palatino's x-height, so a
%   \scriptsize (7.97pt) math operator prints at 7.97 x 1.041 = 8.30pt. It is
%   the same nominal size in a differently scaled font, it is a property of this
%   document's preamble and not of this figure, and every figure in the thesis
%   that sets inline math shows it. The one instance that COULD be removed was:
%   the ellipsis token box was $\cdots$ (CMSY10 at 8.30) and is now \dots, set
%   from Pagella at 7.97 like every other token label.
%   Registers: \footnotesize roman ink = the four panel heads and the axis's
%   quantity line, and nothing else, so the largest type on the drawing always
%   names a panel or a measured quantity and never narrates. \scriptsize carries
%   every label, tick, formula, tie and foot line. Zero \sffamily, zero
%   \bfseries, zero \itshape. Italic characters: 14 of 793, 1.8%, and all
%   fourteen are the math variables h and V, which are a symbol's own face and
%   not a voice. Pass 9 measured 34%.
% PROSE, S4: THE FIGURE NOW CARRIES NONE, which is 4.11's own austerity -- 4.11
%   has axis labels and direct labels and nothing else. Pass 9's two closing
%   lines, "the slab component is set to zero; every other direction is
%   untouched", are deleted because the drawing already carries both claims as
%   MARKS: the accent component's label ends "-> 0", the right angle says the
%   component is zeroed and not merely shortened, and the surviving vector is
%   labelled $h-(hV)V^{\top}$. A sentence whose deletion leaves no drawn object
%   unnamed is not definitional and S4 does not permit it. Decision 9c's "one
%   name for one object" survives, carried now by the band's own direct label,
%   "$V_{\textrm{pure}}$, the purified drug slab", which is also the gloss S19
%   requires at the symbol's first use.
%
% THE RIGHT ANGLE, AND THE ONE GEOMETRIC RATIO THAT MOVED
% --------------------------------------------------------
% The right-angle mark at the foot of the surviving vector was 0.16cm (1.6mm)
% and a naive reader did not see it until zooming. It is now 0.40cm (4.0mm) on
% both legs. Decision 10 above is why it may not simply be dropped: it is what
% distinguishes "the slab component is set to ZERO" from "the slab component is
% reduced". A 4mm glyph cannot sit on a 0.48cm arm without swallowing it, so the
% in-slab component grew from 0.48cm to 0.90cm and the surviving complement from
% 1.90cm to 3.14cm. THE RATIO IS UNCHANGED IN DIRECTION AND BARELY CHANGED IN
% SIZE: 3.49:1 in favour of the survivor against pass 9's 4.0:1, still the right
% way round, and still deliberately NOT the true 14.3:1 (86 degrees,
% cos^2 86 = 0.0049 = 10/2048), at which the component would be 0.13cm, shorter
% than its own arrowhead. figs/slab.tex inset one draws that geometry true, and
% the caption still sends a reader there. Nothing in this panel is a coordinate
% and nothing in it is measured.
%
% THE STRIP: SAME MACRO, SAME FIELDS, EVEN CLEAR SPACE
% -----------------------------------------------------
% The nine box edges were reset so that EVERY cell carries the same 0.13cm of
% clear space each side of its label's measured ink, against pass 9's 0.117 to
% 0.370cm. The strip therefore runs 0.00 to 12.65 instead of 0.00 to 13.55, the
% prompt half ends at 7.44 and the target half begins at 8.44, and everything
% keyed to the boundary (both dividers at \bnd = 7.94, both braces, the brace
% labels at 10.545, the hook rule, the comparison arrow at 13.40) moved with it.
% Two reasons, and both matter: the clear space is now a constant rather than a
% number that varied by a factor of three between cells, and the strip's own
% rules are 3.6cm shorter, which is 5mm2 of the 6.0% ink ceiling. The five
% prompt fields, their order, the [END_CELL] terminator and the ONE-MACRO
% construction are untouched. Ink widths the edges are set from, measured in
% figs/licensing.tex at the same size: cell line 0.954, drug 0.616, dose 0.579,
% mechanism 1.424, control cell sentence 2.511, [END_CELL] 1.41.
%
% GEOMETRY, MEASURED IN THE 200dpi RENDER
% ----------------------------------------
% Content box 173.99 x 99.95mm, against a 170-174mm fullwidth measure and a
% 100mm height cap. Ink 5.98% of the content box, against a 6.0% ceiling and
% pass 9's 8.55%; 793 non-space characters over 17393mm2 is 4.6 per 100mm2,
% against a 5.0 ceiling. Picture bbox 17.38 x 10.04cm (y 0.58 to 10.62).
% Leftmost ink is x = 0.00 -- the four panel heads, both strips, the hook rule
% and its formula, the tie, panel (c)'s band and label, and the foot line, all
% on it. Rightmost ink is the axis spine at 17.38. PANEL X ORIGINS, S13: band
% one and panel (c) at 0.00, panel (d) at 8.40. Measured in the built PDF, the
% left edges of the panel-level text lines are 58.34, 58.34, 58.48, 58.34,
% 58.34, 58.36, 58.34pt for the 0.00 panels -- a spread of 0.14pt -- and
% 296.45pt three times over for panel (d), a spread of 0.00pt.
%
% ===========================================================================
% PASS 11 -- 2026-08-18 -- THE ADVERSARIAL-REFUTATION REPAIR
% ===========================================================================
% EVERYTHING ABOVE THIS LINE IS LEFT INTACT. Three independent refuters read the
% pass-10 build and returned 33 findings. This section supersedes BY NAME only
% the claims above that this pass made false, records every finding APPLIED and
% every finding OVERRULED with its reason, and re-records the measurements.
% NO NUMBER CHANGED IN THIS PASS. 1.563, 0.084, 4.180, n = 60, rank 10, residual
% width 2048 and layer 4 are byte-identical to passes 9 and 10, the log map
% \kx is byte-identical to pass 10, and the three marks stand at the same three
% x to within the 0.02pt pgfmath residual recorded below. Nothing measured moved.
%
% TWO CLAIMS ABOVE ARE FALSE AND ARE SUPERSEDED HERE
% ---------------------------------------------------
% (S-1) The pass-10 note at the foot of this file said the replacement caption
%       "is also staged, keyed by \label, in figs/_PENDING_CAPTIONS.tex for a
%       human to integrate". THAT FILE DID NOT EXIST. It does now, this pass
%       created it, and it carries the caption plus an explicit list of the
%       clauses the integrator must DELETE from the live caption -- above all
%       Sections/04b-representation.tex:289-290, "The two arrows off the rule
%       mark that it acts in both halves and imply no count of positions",
%       which names two marks pass 10 deleted. Until that deletion is made the
%       figure ships contradicting its own caption. hook.json's .shape carried
%       the same false claim and is corrected.
% (S-2) The pass 1-9 block's AGREEMENT WITH THE TEXT cites 04b-representation
%       lines 229, 254, 260 and 286-288, and note 5/decision 5 cites 193, 195
%       and 200-202. Four of those are stale: tab:workspace's layer-4 row is now
%       line 338 (229 is fig:purify's caption), "3.6--18.6x" is 363, the
%       "9--37% ... at layers 2--9" sentence is 368-370, and "No intervals are
%       drawn in either panel" is 395. The 193/195/200-202 set is now 263, 264,
%       265 and 270-272 and pass 10 had already refreshed it in hook.json. The
%       "2048 is prose only at 04b:383" citation is also stale: 383 is
%       fig:mechanism's caption and the residual width is stated at 04b:648.
%       hook.json's citations are repaired; the stale text above is left in
%       place, superseded here, because S25 forbids rewriting a header.
%
% THE FIFTEEN REPAIRS, EACH WITH THE MEASUREMENT THAT FORCED IT
% --------------------------------------------------------------
% (1) PANEL (d), THE HEADLINE GAP WAS COVERED BY ITS OWN LABEL. This is the
%     serious one, because the log axis exists for one reason -- on a log scale
%     a horizontal gap IS a ratio -- and the gap was drawn to the pixel and then
%     papered over. The null's label ran RIGHT from its mark and occupied 100%
%     of the 5.48cm between the open null disc and the filled accent disc. It is
%     now two right-aligned lines running LEFT into the gutter, so the null's own
%     row is clear from its mark rightward and the ratio is a visible span.
% (2) PANEL (d), EVERY VALUE NOW ABUTS ITS OWN MARK. "4.180 --- cell line, the
%     positive control" put its numeral at the far end of the label from the dot
%     it names: measured in the pass-10 PDF, "4.180" began at x = 393.28, the
%     axis interval 0.32-0.43 nats, with its leading digit overlapping the "0.3"
%     tick numeral's column 12.5mm below it, while its own mark stood 48mm away.
%     On a value axis a numeral is read at the position it occupies (S21). Each
%     label is reversed to "name --- value"; no character is added.
% (3) PANEL (d), ALL THREE LABELS ARE NOW ON THE SAME SIDE OF THEIR MARKS. The
%     side flipped between rows two and three, so the disc trailed the text
%     twice and led it once, and at print a trailing 1.4pt disc reads as a
%     bullet ending a line -- three key entries, i.e. a legend, which is the one
%     construction this document forbids outright. All three labels now sit
%     left, all three discs right, and (1) is what made that possible.
% (4) PANEL (d), THE RULE NOW BEGINS AND ENDS ON A DRAWN TICK. At four labelled
%     ticks (0.1, 0.3, 1, 3) the rule overran its tick lane by 36.87pt (13.0mm)
%     at EACH end -- 29% of the axis carrying no numeral, which is the exact
%     defect S13 measures and condemns in fig:comparators row (d) -- and BOTH
%     extreme marks fell in those bare zones, so neither the null at 0.084 nor
%     the positive control at 4.180 could be placed against a printed number.
%     The domain's own ends are now ticks: 0.05, 0.1, 0.3, 1, 3, 6. The two end
%     numerals are set flush INSIDE the rule's ends rather than centred on their
%     ticks, because centred they overhung x = 17.38 by 0.09cm and made the
%     float overfull by 1.88pt; every other numeral is centred on its tick.
% (5) PANEL (d), THE TICK NUMERALS ARE INK, NOT quietink. quietink is this
%     document's comparator-and-null colour (S8/CF4). Spending it on the scale
%     while the figure's actual null was labelled in ink inverted the code.
%     fig:mechanism sets every tick numeral in ink. THE NULL'S OWN INK LABEL AND
%     ITS ink-EDGED OPEN DISC ARE DELIBERATE AND OVERRIDE CF4's quietink slot:
%     CF5 instructs this figure by name to adopt 4.11(b)'s encoding for this
%     object, and 4.11(b) draws it mfc="white", mew=1.1, color=style.INK. The
%     distinction between the null and the positive control is carried by the
%     GLYPH -- open against filled -- which is S9's own device, not by hue.
% (6) THE COMPARISON ARROW'S HEADS POINTED AT NOTHING. It ran between the two
%     brace LABELS' baselines, 9.04 to 6.88, so its upper head sat 0.50cm BELOW
%     the strip it was meant to span and both heads terminated in white space.
%     It now runs 9.75 to 7.59, the two strips' own vertical centres, at x =
%     13.10 rather than 13.40.
% (7) THE COMPARISON ARROW'S LABEL IS NOW ONE WORD, "compared". It had set
%     panel (d)'s axis quantity string a second time 130pt away, so a reader
%     could take this vertical 1.0pt rule for an axis OF that quantity; the axis
%     names the quantity once and this mark's job is to say that the two passes
%     are compared and not chained. S20 is satisfied -- the role is stated in a
%     verb of comparison beside the double head S11 gives a comparison. The 17
%     characters this frees are what pay for (8), (9) and (4) inside the 6.0%
%     ink ceiling, which the pass-10 build cleared by 2.8mm2.
% (8) THE TIE MOVED AND GAINED THE WEIGHTS. It sat at 6.88 -- the EXACT baseline
%     of pass two's brace label, measured to 0.00pt -- and directly under pass
%     two's prompt half, the one half of the strip that carries no brace, so it
%     read as that half's missing brace label, i.e. as a statement about ONE
%     strip when its whole job is to be a statement about TWO. It is now band
%     one's foot line on its own baseline at 6.44, 0.44cm below the brace label
%     and 0.76cm above band two's heads. And it read "the same sentence, box for
%     box", which carries only half of the question this figure must answer:
%     nothing anywhere on the drawing named the WEIGHTS, and the pass-10 note
%     claiming the deleted subtitle "the same prompt, the same sentence, the
%     same weights" had been inherited by the tie was wrong. It now reads "one
%     sentence, one set of weights, scored twice".
% (9) THE LAYER IS NAMED AT THE INTERVENTION. "at every position" ranges over
%     the SEQUENCE, not over depth, and the only "layer 4" on the drawing was in
%     the axis block 40mm away, attached to the READOUT. Panel (b)'s head now
%     reads "(b) pass two, the slab projected out at layer 4". DECISION 14 above
%     and pass 10's partial supersession of it are superseded again: the fact is
%     still stated in words because no mark can carry it, but it is stated where
%     the hook is drawn, and the axis block keeps its own "layer 4" because CF10
%     requires the layer named where the drawn layer's marks are.
% (10) THE HOOK RULE NO LONGER READS AS AN UNDERLINE. At y = 7.92 it sat 3.42pt
%     under an accent line of type and 3.40pt over the box tops -- equidistant,
%     same colour above and below, which is the heading-and-underline idiom. It
%     is at 7.88: 0.18cm of clear space above, 0.08cm below, so it binds to the
%     strip it acts on. It is NOT dropped onto the boxes, as one refuter
%     proposed: a mark running into a rule uninterrupted is what this file's own
%     note 5(b) forbids, and the drug cell's border is accent too.
% (11) PANEL (c), A PRINT COLLISION. The label $(hV)V^{\top} \to 0$ cleared the
%     top edge of the V_pure band by 0.72pt -- two pixels at 200dpi -- so its
%     parenthesis descenders would have closed on press. The band thins from
%     0.32 to 0.24cm and the label rises 0.02cm; the label is now 0.286cm from
%     its own mark, still inside the 3mm at which this document requires a
%     leader, and clears the band by 0.08cm.
% (12) PANEL (c), THE RIGHT ANGLE IS MIRRORED TO THE LEFT. On the right it was
%     unreadable: the 1.0pt black h vector crossed its 0.30pt horizontal leg
%     0.12cm from the corner, and its vertical leg terminated exactly on the
%     accent component's shaft. h rises to the RIGHT of the surviving vector, so
%     the corner between the survivor and the slab on the LEFT is empty. Same 90
%     degrees, same two arms, nothing crossing it. DECISION 10 above stands: the
%     mark is what distinguishes "set to ZERO" from "reduced".
% (13) PANEL (c), THE VERTICAL CONSTRUCTION GUIDE STOPS 0.20cm SHORT of h's tip.
%     It used to end inside the arrowhead. The HORIZONTAL guide still touches
%     both tips and must: it is the closing side of the parallelogram.
% (14) PANEL (c), THE BAND RUNS 0.00 TO 5.80 AND HAS LOST ITS OUTLINE. Two
%     findings, one repair. The lower band held a rectangle 38 x 46mm between
%     this panel and panel (d) with literally no ink in it -- 10% of the content
%     area and the largest void on the drawing -- and the band was byte-
%     identical in stroke and fill to the eighteen token cells above it
%     (rulegrey 0.30pt on panelbg), so ONE rectangle treatment carried two
%     meanings: "a field of the prompt", eighteen times, and "the purified drug
%     slab", once. Lengthening the band closes half the void at no ink cost
%     (panelbg renders above 230 grey and S15 does not count it); dropping the
%     outline makes the band the only fill-without-stroke region in the figure
%     and frees 8.7mm2. AN ACCENT OUTLINE WAS CONSIDERED AND REJECTED AGAIN, for
%     note 9b's reason: accent fill plus accent border is the drug cell's own
%     treatment to within a tint, and one glyph would mean two things again.
%     NOTE 9b IS THEREFORE AMENDED, NOT SUPERSEDED: the band is still panelbg
%     and still not accent; it simply no longer carries slab.tex's rulegrey
%     outline, because in THIS figure that outline is the token cell's own.
% (15) EVERY PANEL HAD TWO LEFT EDGES 1.20pt APART. TikZ's default 1pt inner
%     xsep meant a base-west node's ink began 1.20pt right of the coordinate it
%     was placed on, so the drawn rules started at one x and the panel's text at
%     another, visible under a flat-left glyph like the "K" of "KL". inner
%     xsep=0pt now. MEASURED IN THE BUILT PDF: the panel-level text lines of the
%     0.00 panels start at 57.3, 57.3, 57.5, 57.3, 57.3, 57.4, 57.3pt -- spread
%     0.2pt, and the residual is glyph side bearing, not node placement -- and
%     panel (d)'s head, its first tick numeral and both axis lines start at
%     295.5pt against an axis rule whose left end is 295.25pt.
%
% WHAT WAS PROPOSED AND OVERRULED, AND WHY
% -----------------------------------------
% (O-1) DELETE PANEL (c) AND RE-LETTER (d) AS (c). Refused. The spec's keep-list
%     protects the panel by name -- "h, h - (hV)V^T, the removed component going
%     to zero, and the deferral of the true 10/2048 angle to fig:slab" -- and
%     CF14 forbids losing load-bearing work. The finding's real content, the
%     void, is answered by (14) without deleting anything.
% (O-2) PUT ALL THREE MARKS ON ONE BASELINE AND STAGGER THE LABELS ABOVE WITH
%     LEADERS. Refused: staggering labels off one rule is exactly what S22
%     forbids, and the three names measure about 11cm against a 9cm axis, so one
%     baseline is not available. The three-row construction, one mark and one
%     label per row, is S22's own reading and it survives; (1), (2) and (3) fix
%     what was actually wrong with it.
% (O-3) DRAW DROP LINES OR PROJECTION LINES FROM THE MARKS TO THE SPINE.
%     Refused, and not for the reason pass 10 gave. S7 does sanction a densely
%     dotted 0.5pt quietink projection line, so pass 10's blanket objection was
%     too strong. The decisive objection is arithmetic: the three rows sit 0.38,
%     0.82 and 1.26cm above the spine in VALUE order, so three stems dropped to
%     the axis would have lengths that rise with the value they sit over, and
%     the panel would read as a bar chart of the three numbers against a
%     vertical scale that does not exist. A mark vocabulary must not encode a
%     quantity twice, once truly and once by accident.
% (O-4) ANNOTATE THE TWO GAPS WITH "18.6x" AND "37%" AS 4.11(b) DOES. Refused.
%     The spec deleted both printed ratios from this figure to the caption, and
%     a result is caption material here. 4.11(b) owns the annotated ratios --
%     CF1 says so -- and this figure now shows the gap and lets the caption
%     name it.
% (O-5) SET THE HOOK'S FORMULA IN ink AS DEFINITIONAL PROSE. Refused. It is the
%     direct label of an accent object, and S8 gives accent to "the drug-side
%     object under test AND ITS DIRECT LABELS". It is also \scriptsize, and S4's
%     prose is \footnotesize; calling it prose would break S4's size rule to fix
%     a colour it does not break. (10) fixes the underline reading instead.
% (O-6) DROP THE ACCENT FROM THE "drug" PROMPT CELL. Refused. It is a cross-
%     figure tie recorded in two files: decision 2 above, and figs/licensing.tex,
%     which ADOPTED this figure's cell treatment in pass 9 so that the two
%     strips read as one object seen twice. Removing it here would desynchronise
%     the pair. S8's own check lists what may not be accent -- a comparator, a
%     null, a discounted arm, a struck branch -- and the drug field is none of
%     them.
% (O-7) DRAW THE SECOND STORED RANDOM DRAW AS A FOURTH MARK. Refused as a
%     DRAWING change and ACCEPTED as a disclosure. workspace_probe.json really
%     does store layers.4.kl.random = 0.19796918605764707 beside the drawn
%     layers.4.kl.random_matched = 0.08388670384883881, and with
%     layers.4.raw_drug_dims = layers.4.pure_drug_dims = 10 the two are draws of
%     the same rank-10 construction, so the headline 18.6x would be 7.9x on the
%     other draw. That is a real and undisclosed exposure and the caption now
%     carries it. But the mark itself would be a fourth dot with a role no
%     section of the thesis states, drawn from a numeral that appears in no
%     table and in no sentence -- S20's own failure condition -- and tab:workspace
%     inherits the same single draw, so the figure would then disagree with the
%     table it is read against. hook.json records the second draw beside marks[1].
% (O-8) RENAME THE MARKS TO "ablating the purified drug slab" AND THE AXIS TO
%     "KL from the clean forward pass when a slab is ablated". Refused: S19's
%     registry and the spec both fix these strings, and line 1 is 4.11(b)'s
%     exact quantity name, which CF3 requires the slate to share. The verb the
%     finding wants is in the caption's lede, where a result belongs.
% (O-9) REVERSE OR BEHEAD THE (hV)V^T ARROW BECAUSE IT "POINTS AWAY FROM THE
%     ZERO IT GOES TO". Refused. It is a component of a vector decomposition and
%     points along the slab, which is its direction; "-> 0" is its fate, carried
%     by the label and by the right angle. Pass 9 changed it FROM headless TO
%     headed for a recorded reason -- the removed component was the only one of
%     three not drawn as a vector and read as a baseline -- and reversing the
%     head would draw a vector backwards.
%
% GEOMETRY AND STYLE, RE-MEASURED AT PASS 11
% -------------------------------------------
% Content box 173.99 x 99.95mm (fullwidth measure 170-174mm, height cap 100mm),
% unchanged from pass 10 to the hundredth of a millimetre. Ink 5.970% of the
% content box against a 6.0% ceiling, down from pass 10's 5.984% and pass 9's
% 8.55%. 709 non-space characters over 17390mm2 is 4.08 per 100mm2 against a 5.0
% ceiling. (The 793 recorded at pass 10 was counted by a different method; 709
% is the count of non-space characters in the built PDF's own text spans, and
% that method is now the one hook.json records.)
% TYPE, measured in the built PDF: 7.97pt x522 and 9.96pt x146 -- two chosen
% sizes -- plus the exempt classes 6.85 x20 (\ENDCELL, Inconsolata), 5.98 x8 and
% 6.23 x3 (sub- and superscripts) and 8.30 x10 (the CMSY10/CMR10 math symbols
% pass 10 explains and mathpazo causes). Faces: Pagella-Regular 653,
% PalladioL-Roma 9, Inconsolata 20, PalladioL-Ital 14, CMSY10 7, CMR10 6. Zero
% sans, zero bold, zero all-caps banners. Italic 14 of 709 = 2.0%, and all
% fourteen are the math variables h and V.
% STROKES, measured: 0.30 x27, 0.40 x5, 0.80 x1, 0.90 x1, 1.00 x4 -- the same
% five as pass 10. RECTANGLES: nineteen, and pass 10's count of them was wrong
% in its own arithmetic. Sixteen panelbg token cells, TWO accent!12 drug cells,
% eighteen cells in all, plus the V_pure band -- which is now the only one drawn
% without a stroke. ARROWHEADS: five, three on panel (c)'s vectors and two on
% the comparison arrow.
% THE AXIS, measured: spine x = 295.25 to 549.81pt, ticks at 295.25 (0.05),
% 332.10, 390.48, 454.51, 512.94 and 549.81 (6), i.e. the first and last drawn
% tick sit exactly on the rule's endpoints and there is no unnumbered overrun.
% Marks at 323.06 (null, between the 0.05 and 0.1 ticks), 478.26 (purified slab,
% between 1 and 3) and 530.58 (cell line, between 3 and 6). Each label ends
% 7.56-7.58pt from its own mark. The null-to-slab gap is 155.20pt = 5.474cm and
% carries no text on the null's own row.
%
% ===========================================================================
% PASS 12 -- 2026-08-19 -- THE SECOND ADVERSARIAL-REFUTATION REPAIR
% ===========================================================================
% EVERYTHING ABOVE THIS LINE IS LEFT INTACT. Three further refuters read the
% pass-11 build and returned 30 findings, two of them marked blocking and both
% blocking findings naming the same object. This section supersedes BY NAME the
% claims above that this pass made false, records every finding APPLIED and
% every finding OVERRULED with its reason, and re-measures the drawing.
%
% NO NUMBER CHANGED IN THIS PASS EITHER. 1.563, 0.084, 4.180, n = 60, rank 10,
% residual width 2048 and layer 4 are byte-identical to passes 9, 10 and 11.
% What changed is the value-to-POSITION map (\kx's domain, recorded below and
% in hook.json beside the two maps it supersedes) and where labels and marks
% sit. No mark's VALUE moved; every mark's x moved because the scale did.
%
% CLAIMS ABOVE THAT ARE NOW FALSE, SUPERSEDED HERE BY NAME
% ---------------------------------------------------------
% (T-1) PASS 11's O-2 ("put all three marks on one baseline" -- refused) and
%       O-3 ("drop lines from the marks" -- refused). Both are superseded; see
%       repair (1). O-2's arithmetic was about the LABELS, which still stagger;
%       it wrongly carried that conclusion over to the MARKS, which are 4.52cm
%       and 1.52cm apart and cannot collide. O-3's bar-chart objection is
%       answered by de-correlating leader length from value.
% (T-2) PASS 11's repair (2), "puts every numeral within a few per cent of its
%       own axis position". Measured, it did not: it moved each numeral one
%       decade along and left "4.180" sharing a column with the "3" tick, "1.563"
%       with the "1" tick and "0.084" with the "0.05" tick. Repair (2) below
%       replaces the approximation with an exact construction.
% (T-3) PASS 11's repair (3), "all three labels now sit left of their marks ...
%       a trailing 1.4pt disc reads as a bullet ending a line". The note
%       diagnosed the legend hazard correctly and then made ALL THREE rows
%       trail into a disc, which is the hazard. Superseded by repair (1).
% (T-4) PASS 11's repair (4), which set the domain-end tick numerals at
%       0.05 and 6. Superseded by repair (3): the ladder is now 0.03 / 0.1 /
%       0.3 / 1 / 3 / 10. The REASON recorded in repair (4) -- that the rule
%       must begin and end on a drawn tick and that every mark must have a
%       printed numeral either side of it -- stands and is preserved.
% (T-5) PASS 11's repair (13), which stopped the vertical construction guide
%       0.20cm short of h's tip "at the top end only". Its foot is repaired in
%       this pass; see (7).
% (T-6) PASS 11's repair (12), which mirrored panel (c)'s right angle to the
%       LEFT of the surviving vector. Superseded by (6): the origin cannot
%       carry that mark at all, whichever way it is mirrored.
% (T-7) LAYOUT NOTE (j) in the pass 1-9 block, "THE PASS-ONE BOUNDARY RULE GREW
%       A TAIL, from 8.72 down to 8.55", is superseded by (9): the label that
%       tail was grown to carry was deleted in pass 10, and the two dividers
%       are now identical.
% (T-8) PASS 11's claim at the head of its S-2 that "hook.json's citations are
%       repaired". They were not. Twelve of eighteen line citations in the
%       ledger did not resolve against the current Sections/04b-representation
%       .tex, and five ledger fields still described the pass-10 drawing. All
%       are repaired in this pass and hook.json now carries a pass-12
%       measurement block that supersedes the pass-10 style block by name.
%
% THE ELEVEN REPAIRS, EACH WITH THE MEASUREMENT THAT FORCED IT
% -------------------------------------------------------------
% (1) PANEL (d) IS REBUILT: THE THREE MARKS NOW SHARE ONE BASELINE AND EACH
%     NAME IS LEADERED TO ITS OWN MARK. This is the pass's one structural
%     change and it answers eight findings at once, including both blocking
%     ones. Measured in the pass-11 PDF: the three dot centres stood at
%     (323.06, 268.30), (478.26, 255.83) and (530.58, 243.36)pt -- three
%     baselines 0.44cm apart -- so the 155.20pt "gap" the log axis exists to
%     show was not a distance any reader could lay a ruler along, and the two
%     labels of the rows above ran x 363.68-470.69 and 394.28-523.00, i.e.
%     wholly inside it. Each row was also a text line ENDING in its own disc,
%     so at print three 1.4pt marks read as three bullets and the block read as
%     a key. And the null's two label lines began at x = 245.04 and 241.61
%     against panel (d)'s origin at 295.45 -- 19.0mm outside their own panel,
%     in panel (c)'s column, with "random slab" set 71pt above a band labelled
%     "the purified drug slab" in the same size and face.
%     Now: marks at y = 4.02 at their own three x, nothing else on that line;
%     names on rows 5.26 / 4.86 / 4.46; a 0.30pt leader in each mark's own
%     colour from just under its name to just above its disc. The row order is
%     null / cell line / purified slab, chosen for two reasons at once -- it is
%     the only order in which no leader crosses another row's label, and it
%     makes leader length (1.07 / 0.67 / 0.27cm) non-monotone in value (0.084 /
%     4.180 / 1.563), which is what kills O-3's bar-chart reading. The null's
%     name runs RIGHT from its mark because its mark is 1.59cm from the axis's
%     left end and the name measures 5.17cm; every text line in the panel now
%     lies inside 8.40-17.38.
% (2) EVERY PRINTED VALUE NOW STANDS EXACTLY AT ITS OWN AXIS POSITION. Each
%     label is aligned so its numeral's INNER edge sits on its mark's x:
%     "0.084" begins at \kx{0.084} = 9.9916, "4.180" ends at \kx{4.180} =
%     16.0316, "1.563" ends at \kx{1.563} = 14.5110. That is why one row is
%     value-first and two are value-last -- the value always abuts the mark and
%     the side is set by the measure, not by taste. It also lifts every printed
%     value 1.36cm or more clear of the spine, so the rule no longer carries
%     numerals within a 15pt band on both sides (S21).
% (3) THE TICK LADDER IS UNIFORM. 0.05 / 0.1 / 0.3 / 1 / 3 / 6 stepped x2, x3,
%     x3.33, x3, x2 and printed at 1.300 / 2.061 / 2.258 / 2.061 / 1.300cm -- a
%     74% spread, crowding a pair of numerals at each end of the rule, with the
%     null inside the crowded left pair. 0.03 / 0.1 / 0.3 / 1 / 3 / 10 is the
%     textbook 1-3 ladder and prints at 1.861 / 1.698 / 1.861 / 1.698 / 1.861,
%     a 10% spread. The rule still begins and ends on a drawn tick, no mark
%     sits on a tick, and every mark keeps a printed numeral on each side.
% (4) THE HOOK RULE NO LONGER HEADS A FLUSH-LEFT STACK. Title (9.96pt ink,
%     x0 57.34pt), formula (7.97pt accent, x0 57.48), accent 0.90pt rule (x0
%     57.14), boxes -- four decks on one left edge with 4.5pt of clear space
%     above the rule and 2.8pt below is the heading-and-underline idiom, and
%     the one mark whose whole job is "this operation runs the length of the
%     stream" was reading as "panel (b) starts here". The formula is now
%     right-aligned to the rule's own right end. It is a DIRECT LABEL of a mark
%     (S22), not a panel-level element (S13), and a rule's direct label belongs
%     at one of its ends.
% (5) "compared" IS ROTATED 90 DEGREES, READING UPWARD, BESIDE THE ARROW. Set
%     horizontally its baseline was 8.62 at \scriptsize while panel (b)'s head
%     was 8.54 at \footnotesize, 13.3cm away: two lines 2.27pt apart in
%     baseline, at two sizes, neither aligned nor separated. Rotated there is
%     no second baseline to near-miss and the label sits on the mark it names.
% (6) PANEL (c)'S RIGHT ANGLE MOVES TO THE FOOT OF THE PERPENDICULAR. See the
%     note at the mark itself. The short form: h = survivor + component with
%     both terms positive, so h always lies inside the quadrant the ORIGIN's
%     right angle would enclose, and no mirroring escapes that; the
%     decomposition is a rectangle, so the foot of the perpendicular carries
%     the same 90 degrees and is the one corner that is empty and whose two
%     arms both end on drawn strokes.
% (7) THE VERTICAL CONSTRUCTION GUIDE'S FOOT STOPS 0.14cm SHORT of the accent
%     component's arrow tip, the same treatment pass 11 gave its top end.
% (8) THE SURVIVING VECTOR IS LABELLED AT MID-SHAFT. At the tip, its baseline
%     (5.16) and h's (5.20) sat either side of the horizontal closing guide
%     (5.24), so the row read as "h - (hV)V^T - - - - - h", a dashed RELATION
%     between two expressions -- an equality or comparison the panel does not
%     assert. The guide cannot move; it is the closing side of the
%     parallelogram. One of the two labels had to leave its line, and the
%     survivor's is the one that has a shaft to sit beside.
% (9) THE TWO PROMPT/TARGET DIVIDERS ARE NOW IDENTICAL. Measured: pass one
%     0.686cm, symmetric; pass two 0.572cm, asymmetric. In a figure whose
%     single claim is that the two strips are the same object, the one
%     structural mark that visibly differed between them was a vestige.
% (10) \ENDCELL'S CELL IS 1.47cm, NOT 1.70. Its Inconsolata ink measures
%     1.204cm and every other cell in the strip carries 0.131-0.150cm of clear
%     space a side; this one carried 0.248. The strip end, the braces, the hook
%     rule and the comparison arrow are keyed to \stripend and moved with it.
% (11) THE INTERVAL DECLARATION IS SET ON TWO LINES. At 15.89cm it was the
%     widest object in the figure -- wider than the axis and wider than band
%     one -- which S14 forbids, and an 8pt grey line that wide, stranded a
%     centimetre below everything else, reads as a second and smaller caption
%     above the real one. Broken at the colon it measures 8.68cm and 7.10cm.
%
% WHAT WAS PROPOSED AND OVERRULED, AND WHY
% -----------------------------------------
% (P-1) DELETE PANEL (c) AND RE-LETTER (d) AS (c), giving the readout the full
%     17.38cm. Refused again, for O-1's reason: the spec's keep-list protects
%     the panel by name and CF14 forbids losing load-bearing work. The
%     finding's substantive point -- that the panel overstates the deletion --
%     is conceded in the caption, which also defers the true 10/2048 geometry
%     to fig:slab, where it is drawn at 86 degrees.
% (P-2) WIDEN THE IN-SLAB COMPONENT FROM 0.90 TO 1.35cm so the two black
%     vectors separate at the origin; and, from a different refuter, NARROW IT
%     TO 0.30cm so the panel stops overstating the deletion. Both refused,
%     because they are contradictory demands on one parameter: the angle
%     between h and the survivor IS the drawn ratio, so every millimetre of
%     separation at the origin is bought by overstating what the hook removes.
%     Pass 9 chose that trade-off deliberately at 4.0:1 and recorded it; the
%     built ratio is 3.49:1 and the fusion is confined to the bottom 1.2mm of a
%     31.4mm vector. The relocated right angle (6) also moves the corner a
%     reader has to read away from the fused base.
% (P-3) GIVE THE V_pure BAND AN ACCENT OUTLINE, so the object under test is not
%     drawn in the furniture colour while its own label is accent. The
%     COMPLAINT IS CONCEDED and the repair is refused on arithmetic. The
%     drawing measures 5.972% non-white at 200dpi inside a 173.99 x 99.95mm
%     content box against S15's 6.0% ceiling: 4.9mm2 of headroom. A 0.30pt
%     outline round the 5.80 x 0.24cm band costs 12.8mm2, and even a band
%     shortened to the extent its construction occupies costs more than the
%     ceiling allows -- the longest band an outline could be bought for is
%     2.07cm, which is shorter than the construction standing on it. Pass 9's
%     note 9b and pass 11's repair (14) also both refused an accent border, on
%     the ground that accent-on-panelbg is the drug cell's treatment to within
%     a tint. The band keeps its accent DIRECT LABEL, the accent component is
%     drawn lying inside it, and the same object is an accent filled disc in
%     panel (d), which is where the colour code is carried. RECORDED AS AN
%     OWED REPAIR for any future pass that finds ink to spend.
% (P-4) REBUILD \tokstrip AS TWO OUTER RECTANGLES PLUS INTERNAL VERTICALS, to
%     close the ~0.4mm notches the abutting cells' rounded corners leave in the
%     strip's top and bottom edges. Refused, and it is the closest call in this
%     pass. The notches are real and measurable (13 missing pixels across four
%     boundaries at 200dpi), and the repair would also SAVE 6.2mm2. But the
%     nine-rounded-cells glyph is a recorded cross-figure tie: decision 2 above
%     records that figs/licensing.tex ADOPTED this figure's cell treatment in
%     pass 9 precisely so a reader who meets the prompt strip twice meets one
%     object, and licensing.tex is out of this task's scope, so the repair
%     could not be mirrored and the two figures would diverge again. RECORDED
%     AS AN OWED REPAIR to be made in hook.tex and licensing.tex together.
% (P-5) CENTRE ALL SIX TICK NUMERALS AND CUT THE RULE TO 17.10 TO PAY FOR IT.
%     Refused; the reason is at the ladder itself. In short: at the LEFT end
%     centring puts "0.03" 0.27cm outside panel (d)'s x origin, and S13's
%     one-grid rule, measured at a 0.2pt spread, outranks the optical centring
%     of a domain-end numeral that labels no mark.
% (P-6) EXTEND THE HOOK RULE 0.25cm PAST BOTH STRIP ENDS. Refused: the
%     picture's bounding path is pinned at x = 0.00 and a rule at -0.25 would
%     put the content box at 176.3mm, past S14's 174mm ceiling. (4) breaks the
%     heading reading without moving the rule.
% (P-7) MOVE THE COMPARISON ARROW OUT TO x ~ 15.2 WITH LEADERS RUNNING BACK TO
%     THE STRIPS, to close the empty top-right quadrant. Refused. It would
%     detach the figure's comparison mark from the two objects it compares,
%     buy the composition with two 2.1cm horizontal leaders (4.5mm2, more than
%     the whole ink headroom) and re-open exactly the flow reading -- one strip
%     feeding a mark feeding the other -- that the double head exists to kill.
%     The empty quadrant is accepted and recorded.
% (P-8) RESTORE "generation would begin here" AT THE PROMPT/TARGET DIVIDER, or
%     delete the divider as an unlabelled duplicate of the brace's boundary.
%     Refused both ways. The spec's keep-list protects the divider by name; a
%     label there would restate what the 1.00cm gap already shows and cost ink
%     the ceiling does not have; and the divider is named where an inherited
%     phrase belongs, in the caption, which now reads "the divider marks only
%     where a generated continuation would have begun".
% (P-9) MIRROR PANEL (c)'S CONSTRUCTION so the removed component points LEFT
%     and h tips right. Refused: that is not a mirror, it is an error. With h
%     tipping right and the component pointing left, h is no longer the sum of
%     the two drawn vectors. The finding's DIAGNOSIS -- that the pass-11 mark
%     did not enclose the angle it claimed -- is correct and is applied as (6).
% (P-10) SET THE RIGHT ANGLE IN ink 0.40pt. Refused: the two construction
%     guides beside it are rulegrey 0.30pt, the mark is the third element of
%     that one apparatus, and splitting it off would give 0.40pt a sixth
%     meaning inside S10's five-weight table.
%
% GEOMETRY AND STYLE, RE-MEASURED AT PASS 12
% -------------------------------------------
% Content box 173.99 x 99.95mm -- fullwidth measure 170-174mm, height cap
% 100mm -- unchanged to the hundredth of a millimetre from passes 10 and 11.
% Ink 5.972% at 200dpi against a 6.0% ceiling. 710 non-space characters in the
% built PDF's own text spans over 17390mm2 is 4.08 per 100mm2 against 5.0.
% TYPE, measured in the built PDF: 7.97pt x523 and 9.96pt x146 -- two chosen
% sizes -- plus the exempt classes 6.85 x20 (\ENDCELL, Inconsolata), 5.98 x8 and
% 6.23 x3 (sub- and superscripts) and 8.30 x10 (CMSY10/CMR10 math symbols).
% Faces: Pagella-Regular 654, PalladioL-Roma 9, Inconsolata 20, PalladioL-Ital
% 14, CMSY10 7, CMR10 6. Zero sans, zero bold, zero all-caps banners. Italic
% 14 of 710 = 2.0%, and all fourteen are the math variables h and V.
% STROKES: 0.30 (furniture -- 18 cell borders' worth of rulegrey, both braces,
% six axis ticks, two construction guides, the right angle, three label
% leaders), 0.40 (two prompt/target dividers, the axis rule), 0.80 (ONE, the
% null's open-disc edge: CF5 instructs this figure by name to take 4.11(b)'s
% encoding for this object and mechanism.py draws it mew=1.1, so the singleton
% is a mandated distinct kind of object and not a stray weight), 0.90 (ONE, the
% hook rule: it is the figure's intervention and the single heaviest non-data
% stroke; at 1.00pt it would merge with the vectors and the comparison arrow,
% which are data paths, and at 0.40pt with the axis), 1.00 (three vectors and
% the comparison arrow). Five weights, both singletons justified.
% RECTANGLES: nineteen -- sixteen panelbg token cells, two accent!12 drug
% cells, and the V_pure band, which is still the only one drawn without a
% stroke. ARROWHEADS: five, three on panel (c)'s vectors and two on the
% comparison arrow; no other mark in the figure carries one.
% THE AXIS, measured in the built PDF: spine 295.25 to 549.81pt; ticks at
% 295.25 (0.03), 348.07 (0.1), 396.15 (0.3), 448.91 (1), 497.05 (3) and 549.81
% (10), i.e. tick spacings 52.82 / 48.08 / 52.76 / 48.14 / 52.76pt, a 10%
% spread against pass 11's 74%, and the first and last drawn tick still sit
% exactly on the rule's endpoints. The four interior numerals are centred on
% their ticks; the two domain-end numerals are set flush inside the rule's ends
% for the reason recorded at the ladder. THE THREE MARKS ARE AT 340.61 (null,
% white-filled open disc, 0.797pt edge), 468.48 (purified slab, accent filled)
% and 511.58 (cell line, ink filled) -- ALL THREE AT cy = 262.63, one baseline
% to the hundredth of a point. The null-to-slab gap is 127.87pt = 4.510cm and
% carries no ink at all; the gap the map predicts for a ratio of
% 18.634044258027835 is 4.521cm, so the drawn span IS the ratio to within
% 0.011cm, which is a third of the width of the stroke that draws the discs.
% The slab-to-cell-line gap is 43.10pt = 1.520cm against a predicted 1.5202cm.
% THE THREE LEADERS: 0.299pt, at x = 340.61 (ink, 30.33pt), 511.58 (ink,
% 19.00pt) and 468.48 (accent, 7.65pt) -- lengths 1.07 / 0.67 / 0.27cm against
% values 0.084 / 4.180 / 1.563, deliberately not monotone.
% ===========================================================================
%
% ===========================================================================
% v-BLOCK, 2026-08-19: WHOLE-SLATE PASS (fig:comparators, fig:licensing,
% fig:purify, fig:hook, fig:materiality reconciled against one another and
% against fig:mechanism, which is the document's reference standard and was
% NOT rebuilt). Nothing above this line is deleted; where an earlier claim is
% now false it is superseded EXPLICITLY below.
%
% WHY. The five figures were rebuilt independently and drifted. This pass owns
% only the differences BETWEEN them. No value changed in any of the five; the
% sibling .json ledgers carry the value-to-position maps, old and new.
%
% THE SLATE RULES THAT NOW HOLD ACROSS ALL FIVE (each stated once here, in
% every file, so any future single-figure edit can see what it would break):
%
%  R1 TICK NUMERALS are quietink, \scriptsize, inner sep 0, anchored north on
%     the tick and set 0.16cm below the axis rule, over a 0.10cm rulegrey
%     0.30pt tick. Before: comparators/licensing/purify quietink, hook and
%     materiality ink; ticks 0.10 / 0.12 / 0.14cm; three different gaps.
%  R2 A THRESHOLD'S OWN LABEL IS BLACK, matching the black solid 0.90pt rule
%     it names (S6). Before: comparators black, purify's "gate, 0.8" and
%     materiality's two margin labels ink.
%  R3 ONE NAME PER OBJECT (S19). The pre-registered 10% rule is "10%
%     materiality margin" in fig:comparators row (c), fig:materiality (a) and
%     (b), and in figs/family.py, which is where the name comes from.
%     fig:comparators' "pre-registered margin" is superseded; "pre-registered"
%     is now the caption's word, not the drawing's. Zero, where it is the null
%     a bar is read against, is "zero: no identity effect" in BOTH
%     fig:comparators row (c) and fig:materiality (b) -- one quantity, one
%     wording. V_pure is spelt with \textrm everywhere, as figs/slab.tex does.
%  R4 COLOUR KEY, one key for five figures (S8). accent = the drug-side object
%     under test. quietink = every comparator, null, control, replication or
%     discounted arm, AND its label, AND its leader. ink/black = thresholds,
%     axis rules, axis labels, panel heads, definitional prose. rulegrey =
%     furniture only. Within quietink the GLYPH separates two kinds: an OPEN
%     disc is a construction built to carry no signal (a null, a control, a
%     raw/before series); a FILLED disc is a real measurement that is not the
%     headline. fig:hook (d) drew both its null and its cell-line control in
%     ink and is corrected.
%  R5 PROJECTION AND ZERO RULES are quietink 0.50pt densely dotted (S7).
%     fig:hook's two panel-(c) projection lines were rulegrey 0.30pt densely
%     dashed and are corrected.
%  R6 DEFINITIONAL PROSE is \footnotesize ink and OUTRANKS every label (S4);
%     at most one block per figure, on the panel's own x origin. fig:hook's
%     one prose line was \scriptsize and is raised. fig:purify's panel-(b)
%     block carried an argument clause ("so causal claims use the purified
%     slab, not the raw") that its caption already carries, and is now
%     definitional only.
%  R7 ONE INTERVAL DECLARATION per figure, \scriptsize quietink, at the foot,
%     on the figure's x origin, opening with the word "Intervals:" and naming
%     EVERY panel including the ones with no measured mark (S17).
%     fig:comparators already did this and is the model; the other four now
%     match its form.
%  R8 THE FIVE-FIELD PROMPT STRIP is ONE object drawn in two figures, and the
%     two drawings are now congruent: cell boundaries 0 / 1.12 / 1.90 / 2.64 /
%     4.23 / 6.90 cm and cell height 0.42cm in both fig:licensing (a) and
%     fig:hook (a) and (b), labels centred in their cells, the drug cell
%     accent-bordered at 0.40pt on accent!12 and its label NOT bold. Before:
%     licensing 6.90cm wide with 0.62cm cells and a bold label, hook 7.44cm
%     wide with 0.42cm cells and a plain label -- the same five fields at two
%     sizes, four pages apart, with each caption pointing at the other.
%  R9 MEASURE. Every fullwidth figure draws 172.0-174.0mm; the one textwidth
%     figure draws 140.3mm. Ink at 200dpi inside the content bbox, counted as
%     greyscale luminance below 230 (the contract's own method), is now
%     5.63-5.99% on all five, under the 6.0% body ceiling, where before this
%     pass it ran 5.36-6.54%.
%
% WHAT THIS PASS DELIBERATELY DID NOT UNIFY, so that a later reader does not
% think it was missed:
%  * PANEL-HEAD CASE. fig:comparators sets "(a) Is it there?" and the other
%    four set "(a) where the drug name enters". comparators' heads are
%    interrogative SENTENCES and sentence case is correct for a sentence; the
%    v4 header argued this and the argument stands. fig:mechanism's lowercase
%    heads are noun phrases, not questions. One capital letter, four times.
%  * PANEL-HEAD POSITION. fig:comparators puts its heads in a right-aligned
%    left stub, the other four above the panel at its x origin. Moving them
%    above the rules costs about 20mm of height against a 118mm cap that the
%    figure already meets at 117.60mm. Not available.
%  * TWO-COLUMN GUTTER. purify and hook split at x = 8.40cm, licensing at
%    7.75cm. licensing cannot move right: "(d) what a negative would buy" is
%    4.85cm wide and its origin is already at 12.40 of a 17.30cm measure, with
%    1.5pt of slack. Moving purify and hook LEFT to 7.75 would be arbitrary.
% ===========================================================================
%
% CHANGED IN THIS FILE. (i) The prompt strip adopts fig:licensing's cell
% boundaries, 0 / 1.12 / 1.90 / 2.64 / 4.23 / 6.90cm, so the one object is one
% width in both figures (R8). The divider \bnd follows it from 7.94 to 7.40,
% the target strip from 8.44-12.42 to \tgt = 7.90 - \stripend = 11.88 with its
% four cell widths unchanged, the two braces and their "scored tokens" labels
% with it, and the comparison arrow from x 12.87 to 12.33 with its rotated
% label from 13.27 to 12.73. Panel (d), the axis and the foot do not move.
% (ii) Panel (d): the matched-dimension random slab and the cell-line positive
% control are comparators, so their marks, labels and leaders go ink ->
% quietink (R4). The null keeps the open disc and the control the filled one,
% which is the slate's glyph distinction between a construction built to carry
% no signal and a real measurement that is not the headline. accent now
% appears exactly once in panel (d), on the object under test.
% (iii) The six x-tick numerals go ink -> quietink, anchor=base -> anchor=north
% at a 0.06cm gap (R1); the end pair keep their north-west / north-east pull.
% (iv) The two panel-(c) projection lines go from rulegrey 0.30pt densely
% dashed to quietink 0.50pt densely dotted, which is the register
% fig:comparators, fig:purify and fig:materiality all use for a projection and
% which S7 reserves for it (R5).
% (v) "one sentence, one set of weights, scored twice" goes \scriptsize ->
% \footnotesize ink, so this figure's one prose block outranks its labels as
% fig:purify's and fig:licensing's do (R6).
% (vi) The hook formula's tail, "prompt and target alike", is deleted; the
% accent rule visibly spans the prompt half and runs across the divider, and
% the caption states it in full. This is the one deletion made for the ink
% ceiling rather than for a defect, and it is the first thing to restore.
% (vii) V_pure is respelt with \textrm, as figs/slab.tex and figs/purify.tex
% spell it. (viii) The foot block takes the slate's declaration form (R7).
% STILL TRUE AND STILL DECLARED: the 0.90pt accent hook rule is a singleton
% weight, kept under S10's own permission; it is accent, not black, so it
% cannot be read into the threshold register, and this figure draws no
% threshold. The ten 8.30pt glyphs are CMR10/CMSY10 math signs inside the two
% formulas, scaled 1.041 by mathpazo; they are not a third chosen type size
% and cannot be removed without deleting the definition of the hook.
% RE-MEASURED: 173.99 x 99.95mm; ink 5.955%; 806 characters, 4.63 per 100mm^2.
% One page, zero Overfull, zero Underfull.
% ===========================================================================
\begin{tikzpicture}[x=1cm, y=1cm, font=\small,
                    every node/.style={inner xsep=0pt, inner ysep=1pt}]
  % TWO SIZES, ONE FACE. \footnotesize (10.0pt) is reserved for the four panel
  % heads and the axis's quantity line; \scriptsize (8.0pt) carries everything
  % else. No \sffamily, no \bfseries, no \itshape anywhere in the picture.
  % inner xsep=0pt, PASS 11: with TikZ's 1pt x padding a base-west node's ink
  % began 1.20pt right of the coordinate it was placed on, so in every panel the
  % drawn rules started at one x and the text at another (S5, S13). The padding
  % is now zero in x and the text sits ON the panel's origin.
  \tikzset{
    ttl/.style  ={font=\footnotesize, text=ink,      anchor=base west},
    lab/.style  ={font=\scriptsize,   text=ink,      anchor=base west},
    labe/.style ={font=\scriptsize,   text=ink,      anchor=base east},
    qlab/.style ={font=\scriptsize,   text=quietink, anchor=base west},
    qlabe/.style={font=\scriptsize,   text=quietink, anchor=base east},
    tick/.style ={font=\scriptsize,   text=quietink, anchor=north,
                  inner sep=0pt},
    def/.style  ={font=\footnotesize, text=ink,      anchor=base west},
    alab/.style ={font=\scriptsize,   text=accent,   anchor=base west},
    brc/.style  ={draw=rulegrey, line width=0.30pt, decorate,
                  decoration={brace, amplitude=3.5pt, mirror}},
    arw/.style  ={-{Latex[length=3.4pt,width=2.8pt,line width=0pt]}},
    cmp/.style  ={{Latex[length=3.4pt,width=2.8pt,line width=0pt]}-%
                  {Latex[length=3.4pt,width=2.8pt,line width=0pt]}}}

  % --- one token box, and the drug box, so the two strips cannot differ ------
  % Border 0.30pt: a cell is a container, and containers are furniture here.
  % The drug cell keeps its 0.40pt accent border, because it is the one field
  % of the five that is an object in the argument -- the place the prompt names
  % the drug in plain text -- and it is the only accent inside either strip.
  \newcommand{\tokbox}[4]{%
    \draw[rulegrey, line width=0.30pt, fill=panelbg, rounded corners=1pt]
      (#1,#3) rectangle (#2,{#3+0.42});
    \node[font=\scriptsize, text=ink, inner sep=0pt]
      at ({(#1+#2)/2},{#3+0.20}) {#4};}
  \newcommand{\tokdrug}[3]{%
    \draw[accent, line width=0.40pt, fill=accent!12, rounded corners=1pt]
      (#1,#3) rectangle (#2,{#3+0.42});
    \node[font=\scriptsize, text=accent, inner sep=0pt]
      at ({(#1+#2)/2},{#3+0.20}) {drug};}
  % THE strip. One argument: the y of the box bottoms. Called twice, and that
  % is the whole guarantee that pass one and pass two are the same sentence.
  \newcommand{\tokstrip}[1]{%
    \tokbox{0.00}{1.12}{#1}{cell line}%
    \tokdrug{1.12}{1.90}{#1}%
    \tokbox{1.90}{2.64}{#1}{dose}%
    \tokbox{2.64}{4.23}{#1}{mechanism}%
    \tokbox{4.23}{6.90}{#1}{control cell sentence}%
    \tokbox{\tgt}{8.88}{#1}{gene$_1$}%
    \tokbox{8.88}{9.86}{#1}{gene$_2$}%
    \tokbox{9.86}{10.41}{#1}{\dots}%
    \tokbox{10.41}{\stripend}{#1}{\ENDCELL}}
  % the prompt/target boundary: mid-gap, where nothing else is drawn.
  \def\bnd{7.40}
  \def\tgt{7.90}
  % PASS 12: 12.42, not 12.65. \ENDCELL's cell was the one field of nine whose
  % clear space did not match the other eight: measured ink 1.204cm (Inconsolata
  % at \scriptsize, built PDF x 374.49-408.76pt) inside a 1.70cm cell is 0.248cm
  % a side against 0.131-0.150cm everywhere else, so at print it was the one
  % roomy cell in a strip whose whole claim is regularity. 1.204 + 2 x 0.133 =
  % 1.470. The strip end is a \def and the brace, the hook rule and the
  % comparison arrow are all keyed to it, so they move with it.
  \def\stripend{11.88}

% ###########################################################################
% # BAND ONE -- the two passes.  x origin 0.00
% ###########################################################################

  % ---------------- (a) pass one, clean -------------------------------------
  \node[ttl] at (0.00,10.34) {(a) pass one, clean};

  \tokstrip{9.54}

  % The divider is rulegrey 0.40pt SOLID, not ink 0.5pt densely dotted: S7
  % reserves the dotted stroke for a zero rule, a leader or a projection line
  % and forbids it a label of its own, and this one is labelled. It is a piece
  % of the strip's own structure, so it takes a strip weight, and it sits at
  % mid-gap because a rulegrey line on a rulegrey box border is invisible.
  % PASS 12: 9.40 to 9.98, not 9.40 to 10.10. The two dividers were drawn at two
  % different lengths and two different symmetries -- pass one 0.686cm with
  % 0.133cm clear above and below its boxes, pass two 0.572cm with 0.134 below
  % and 0.018 above (clipped to clear the hook rule at 7.88). The one structural
  % mark that visibly DIFFERED between two strips whose whole claim is that they
  % are identical was a leftover: layout note (j) below grew that tail to carry
  % the label "generation would begin here", which pass 10 deleted. NOTE (j) IS
  % THEREFORE SUPERSEDED -- its reason no longer exists. Both dividers now run
  % 0.58cm with the same 0.14/0.02 offsets from their own boxes.
  \draw[rulegrey, line width=0.40pt] (\bnd,9.40) -- (\bnd,9.98);

  % The brace spans the TARGET side only and starts at the first target box's
  % own left edge, 9.20. Both braces carry the SAME label: the two spans are
  % the same tokens read twice, and two different labels sent a reader looking
  % for a difference between them.
  \draw[brc] (\tgt,9.44) -- (\stripend,9.44);
  \node[font=\scriptsize, text=ink, anchor=base] at (9.89,9.04)
    {scored tokens};

  % ---------------- (b) pass two, the slab projected out --------------------
  % PASS 11: "at layer 4" is added to the head. The word "every" in "at every
  % position" ranges over the SEQUENCE, not over depth, and the only place the
  % depth was named was the axis block 40mm away, attached to the READOUT rather
  % than to the intervention -- so a reader could take the hook to run at every
  % layer as well as at every position, which is exactly the misreading decision
  % 14 above says the figure exists to prevent. The axis block keeps its own
  % "layer 4" because CF10 requires the layer named where the one drawn layer's
  % marks are.
  \node[ttl] at (0.00,8.54) {(b) pass two, the slab projected out at layer 4};

  % THE HOOK. The RULE is what says "everywhere along here": it spans the whole
  % strip, prompt half and target half, and runs straight across the
  % prompt/target boundary, which is the point. Its formula sits ON the rule
  % rather than 1.1cm under the strip, so the label names the mark it is beside.
  % The two short arrows of pass 9 are gone: they implied a count of two and
  % the caption had to deny it.
  % PASS 11: the rule drops from 7.92 to 7.88, so its leading is ASYMMETRIC --
  % 0.18cm of clear space above it to the formula's descenders, 0.08cm below it
  % to the box tops. At 7.92 it was equidistant between an accent line of type
  % above and the boxes below, and accent type with an accent rule tight beneath
  % it is the heading-and-underline idiom; the rule now binds to the strip it
  % acts on. It is NOT dropped onto the boxes: a mark running into a rule
  % uninterrupted is what the house rule forbids, and the drug cell's border is
  % accent too.
  % PASS 12: THE FORMULA IS RIGHT-ALIGNED TO THE RULE'S OWN RIGHT END. Flush
  % left it made a three-deck flush-left stack -- title (9.96pt ink, x0 57.34pt),
  % formula (7.97pt accent, x0 57.48pt), accent 0.90pt rule (x0 57.14pt), boxes
  % -- with 4.5pt of clear space above the rule and 2.8pt below, and an accent
  % rule tight under an accent line of type sharing its left edge is the
  % heading-and-underline idiom. The mark that has to say "this operation runs
  % along the stream" was reading as "panel (b) starts here". The formula is a
  % DIRECT LABEL of the rule, not a panel-level element, so S22 governs it and
  % not S13: a rule's direct label sits at one of its ends, and putting it at
  % the far end from the panel's title breaks the stack without moving the rule.
  % The rule is NOT extended past the strip ends, as one refuter proposed: the
  % picture's bounding path is pinned at x = 0.00 and a rule at -0.25 would put
  % the content box at 176.3mm, past S14's 174mm ceiling.
  \node[font=\scriptsize, text=accent, anchor=base east] at (\stripend,8.12)
    {$h \leftarrow h-(hV)V^{\top}$ at every position};
  \draw[accent, line width=0.90pt] (0.00,7.88) -- (\stripend,7.88);

  \tokstrip{7.38}

  \draw[rulegrey, line width=0.40pt] (\bnd,7.24) -- (\bnd,7.82);
  \draw[brc] (\tgt,7.28) -- (\stripend,7.28);
  \node[font=\scriptsize, text=ink, anchor=base] at (9.89,6.88)
    {scored tokens};

  % ---------------- the tie: the figure's central claim, drawn --------------
  % It sits UNDER both strips, on band one's own x origin, and it is now band
  % one's FOOT LINE: on its own baseline, 0.44cm below pass two's brace label
  % and 0.76cm above the two panel heads of band two, so it belongs to the band
  % it closes.
  % PASS 11, TWO REPAIRS. (i) POSITION. It was at 6.88, the EXACT baseline of
  % pass two's brace label, and directly under pass two's prompt half -- the one
  % half of the strip that carries no brace. A reader therefore read it as the
  % missing brace label of that half ("prompt half: the same sentence, box for
  % box | target half: scored tokens"), i.e. as a statement about ONE strip,
  % which is the opposite of what it says. Nothing else is now set on 6.44.
  % (ii) WORDING. "the same sentence, box for box" carried only half of the
  % question this figure has to answer -- are the two passes scoring the same
  % sentence WITH THE SAME WEIGHTS -- and nothing anywhere on the drawing named
  % the weights at all; a reader could infer them only from the formula being an
  % activation edit. The deleted pass-9 subtitle "the same prompt, the same
  % sentence, the same weights" did carry it and pass 10's note claimed the tie
  % had inherited it, which was not true. It now reads "one sentence, one set of
  % weights, scored twice". This also clears the S18 duplication with the
  % caption's "are the same sentence box for box", which is reworded in the
  % staged caption to "emitted by one macro, so no difference between them is
  % possible except the hook".
  \node[def] at (0.00,6.36) {one sentence, one set of weights, scored twice};

  % ---------------- the comparison -----------------------------------------
  % DOUBLE-headed, and ink. The two passes are independent scorings of one
  % fixed sentence by one set of weights; that is a comparison, not a flow, and
  % the single head read as pass one feeding pass two. The single accent
  % arrowhead is reserved for an operation acting in a direction.
  % PASS 11, TWO REPAIRS. (i) GEOMETRY. Its two heads pointed at nothing: it ran
  % between the two brace LABELS' baselines (9.04 to 6.88), so the upper head sat
  % 0.50cm BELOW the strip it was meant to span and both heads terminated in
  % white space. It now runs between the two strips' own vertical centres, 9.75
  % to 7.59, and sits 0.45cm off the strip end instead of 0.75cm, so each head is
  % level with the band it names.
  % (ii) LABEL. "KL from the clean / forward pass" is deleted and replaced by one
  % word. It set panel (d)'s axis quantity string a second time, 130pt away, so a
  % reader could take this vertical rule for an axis OF that quantity; the axis
  % names the quantity, and this mark's job is to say that the two passes are
  % compared and not chained. S20 is satisfied -- the role is stated in one verb
  % of comparison, beside a double head that S11 gives a comparison. The 17
  % characters this frees are what pay for "at layer 4", the reworded tie and the
  % two domain-end tick numerals inside the 6.0% ink ceiling.
  % PASS 12, TWO MOVES. (i) x = 12.87, not 13.10: the strip end moved from 12.65
  % to 12.42 and the arrow keeps pass 11's 0.45cm standoff from it.
  % (ii) THE LABEL IS ROTATED 90 DEGREES, READING UPWARD, BESIDE THE ARROW'S
  % SHAFT AND CENTRED ON ITS MIDPOINT. Set horizontally it sat on baseline 8.62
  % at \scriptsize while panel (b)'s head sat on 8.54 at \footnotesize, 13.3cm
  % away: two lines of type 2.27pt apart in baseline, at two sizes, neither
  % sharing a baseline nor clearly separated, which is the near-alignment a
  % reader reads as carelessness. Rotated, there is no second baseline to
  % near-miss, the label sits ON the mark it names instead of beside it, and it
  % adopts fig:mechanism's own rotated-label idiom. It reads upward, as 4.11's
  % y-labels do, so its ascenders point left and its ink clears the 1.00pt shaft
  % by 0.18cm.
  \draw[ink, line width=1.00pt, cmp] (12.33,9.75) -- (12.33,7.59);
  \node[font=\scriptsize, text=ink, anchor=base, rotate=90]
    at (12.73,8.67) {compared};

% ###########################################################################
% # BAND TWO, LEFT -- (c) what the hook does to one vector.  x origin 0.00
% ###########################################################################
  \node[ttl] at (0.00,5.68) {(c) what the hook does to one vector};

  % THE SLAB. PASS 11, TWO CHANGES, AND THE SECOND PAYS FOR MOST OF THIS PASS.
  % (i) IT RUNS 0.00 TO 5.80, not 0.00 to 3.80. The lower band held a rectangle
  % 38 x 46mm between this panel and panel (d) containing literally no ink --
  % 10% of the content area, and the largest single void on the drawing. The
  % band is the panel's own object and lengthening it is what a subspace does;
  % the construction moves with it to stay on the band's midpoint.
  % (ii) THE rulegrey OUTLINE IS GONE; the band is now a fill alone. It was
  % byte-identical in stroke and fill to the eighteen token cells above it
  % (rulegrey 0.30pt on panelbg), so one rectangle treatment carried two
  % meanings -- "a field of the prompt", eighteen times, and "the purified drug
  % slab, a ten-dimensional subspace", once -- and the only difference was
  % square against rounded corners on a 9pt-tall shape. An ACCENT outline was
  % considered and rejected again for pass 9's reason (accent fill plus accent
  % border is the drug cell's own treatment to within a tint), so the border
  % goes rather than changes colour. The band keeps its accent direct label, is
  % 24 times as wide as it is tall, and is the only fill-without-stroke region
  % in the figure. It also frees 8.7mm2 of the 6.0% ink ceiling, which is what
  % buys "at layer 4", the reworded tie and the two new domain-end ticks.
  \fill[panelbg] (0.00,1.98) rectangle (5.80,2.22);

  % construction lines first, so the arrows and the right angle cover them.
  % PASS 11: the vertical guide stops 0.20cm short of h's tip. It used to end
  % ON the tip, inside the arrowhead, so a leader ran into the mark it points
  % from. The horizontal guide still touches both tips: it IS the closing side
  % of the parallelogram and a gap there would break the construction.
  % PASS 12: the vertical guide now stops at 2.24, 0.14cm above the component's
  % line and just clear of the V_pure band's top edge. It used to run to 2.10,
  % which is exactly the accent component's arrow TIP, so a dashed leader ended
  % inside a mark and at 6x zoom the guide read as a continuation of the
  % component rather than as the side of a rectangle. The same 0.20cm-standoff
  % treatment pass 11 gave the guide's top end, applied to its foot.
  \draw[quietink, line width=0.50pt, densely dotted] (3.40,5.04) -- (3.40,2.24);
  \draw[quietink, line width=0.50pt, densely dotted] (2.50,5.24) -- (3.40,5.24);

  \draw[ink, line width=1.00pt, arw] (2.50,2.10) -- (3.40,5.24);
  \node[lab] at (3.50,5.20) {$h$};

  % PASS 12: THE SURVIVOR IS LABELLED AT MID-SHAFT, NOT AT ITS TIP. At the tip
  % its baseline was 5.16 and $h$'s was 5.20, with the horizontal closing guide
  % between them on 5.24 -- so the row read left to right as
  % "h - (hV)V^T - - - - - h", a dashed RELATION joining two expressions, i.e.
  % an assertion of equality or comparison the panel does not make (the two are
  % a vector and its component, not two sides of anything). Moving one of the
  % two labels off the guide's line is what kills that reading; the guide itself
  % cannot move, because it IS the closing side of the parallelogram and must
  % touch both tips. The label now sits beside the shaft it names, 0.12cm from
  % it and 0.56cm from $h$'s shaft at the same height, so the attachment is
  % unambiguous and no leader is required.
  \draw[ink, line width=1.00pt, arw] (2.50,2.10) -- (2.50,5.24);
  \node[labe] at (2.38,3.62) {$h-(hV)V^{\top}$};

  % the in-slab component, and its fate carried as a mark, not only in words.
  % PASS 11: its label rises from +0.26cm to +0.28cm above the segment and the
  % band thins from 0.32 to 0.24cm, because at the old geometry the label's
  % parenthesis descenders cleared the band's top rule by 0.72pt -- two pixels
  % at 200dpi -- and would have closed on press. The label is still 0.286cm
  % from its own mark, inside the 3mm at which this document requires a leader.
  % PASS 12: the label's x origin moves 3.46 -> 3.58. At 3.46 its opening
  % parenthesis began 0.06cm (2.0pt) right of the vertical guide at 3.40, so at
  % print the dashed guide and the "(" read as one glyph cluster. 0.18cm now.
  \draw[accent, line width=1.00pt, arw] (2.50,2.10) -- (3.40,2.10);
  \node[alab] at (3.58,2.38) {$(hV)V^{\top} \to 0$};

  % the right angle: the component is set to ZERO, not merely reduced. 0.40cm
  % on both legs, against pass 9's 0.16cm, which a reader did not see.
  % PASS 11: IT IS MIRRORED TO THE LEFT OF THE SURVIVING VECTOR. On the right it
  % was unreadable: the 1.0pt black h vector crossed its horizontal leg 0.12cm
  % from the corner, and its vertical leg terminated exactly on the accent
  % component's shaft, so a 0.30pt rulegrey glyph sat under two heavier strokes
  % and read as an artefact. h rises to the RIGHT of the survivor, so the corner
  % between the survivor and the slab on the LEFT is empty; the same 90 degrees,
  % the same two arms, nothing crossing it.
  % PASS 12: THE MARK MOVES TO THE FOOT OF THE PERPENDICULAR, (3.40,2.10).
  % Pass 11 mirrored it to the LEFT of the survivor to escape a collision, and
  % in doing so put it in a quadrant where NEITHER of its arms lies on a drawn
  % ray: the two rays that leave the shared origin (2.50,2.10) are the survivor
  % (up) and the accent component (right), so they bound the up-RIGHT quadrant,
  % while the mark enclosed the up-LEFT one -- its upper arm dead-ended on the
  % survivor's 1.0pt shaft and its lower arm ended at (2.10,2.10) inside the
  % V_pure band's fill, on the backward extension of a ray that is not drawn.
  % At 6x zoom it read as a detached grey bracket, and three refuters read it
  % that way independently.
  % THE ORIGIN CANNOT CARRY THE MARK AT ALL, and that is why this is a move and
  % not a re-mirror: h = survivor + component with both terms positive, so h
  % ALWAYS lies inside the very quadrant the origin's right angle encloses, and
  % at 0.40cm legs h crosses the far arm 0.115cm from its corner whichever way
  % the construction is mirrored. The decomposition is a RECTANGLE -- origin,
  % (3.40,2.10), (3.40,5.24), (2.50,5.24) -- so all four corners carry the same
  % 90 degrees, and the foot of the perpendicular is the one corner that is
  % empty. Its two arms end on drawn strokes: (3.40,2.50) on the vertical
  % construction guide, (3.00,2.10) on the component's own shaft. It asserts
  % what decision 10 says it must -- the surviving vector is orthogonal to the
  % slab, so the slab component is set to ZERO and not merely reduced.
  % IT STAYS rulegrey 0.30pt. One refuter asked for ink 0.40pt on the ground
  % that a load-bearing assertion should not share a stroke with a box border.
  % Overruled: the two construction guides in this same panel are rulegrey
  % 0.30pt, this mark is the third element of that one apparatus, and splitting
  % it off would give 0.40pt a sixth meaning inside S10's five-weight table.
  \draw[rulegrey, line width=0.30pt]
    (3.40,2.50) -- (3.00,2.50) -- (3.00,2.10);

  % ONE label for the band, 0.32cm under its own left edge: the symbol and
  % its gloss in one place, which is the only licence S19 gives the symbol.
  \node[alab] at (0.00,1.66) {$V_{\textrm{pure}}$, the purified drug slab};

% ###########################################################################
% # BAND TWO, RIGHT -- (d) what comes out.  x origin 8.40
% ###########################################################################
  \node[ttl] at (8.40,5.68) {(d) what comes out};

  % VALUE-TO-POSITION MAP, LOGARITHMIC. Braced: it is a pgfmath group, and a
  % shim added outside the braces is not a coordinate, so a label that needs a
  % shim takes it as an xshift on the node.
  %   \kx{v} = 8.40 + (ln v - ln 0.03)/(ln 10 - ln 0.03) * 8.98
  %   domain [0.03, 10] nats; spine 8.40 to 17.38; ticks 0.03 0.1 0.3 1 3 10
  %   0.084 -> 9.9916   1.563 -> 14.5110   4.180 -> 16.0316
  % A horizontal gap is a ratio: null to slab is 4.5194cm and IS 18.6x; slab to
  % cell line is 1.5206cm and IS 37% read the other way.
  %
  % PASS 12, THE DOMAIN CHANGES FROM [0.05, 6.0] TO [0.03, 10] AND NO VALUE
  % MOVES. The pass-11 ladder was 0.05 / 0.1 / 0.3 / 1 / 3 / 6: steps of x2, x3,
  % x3.33, x3, x2, which printed as spacings 1.300 / 2.061 / 2.258 / 2.061 /
  % 1.300cm -- a 74% spread, with a crowded pair of numerals jammed at each end
  % of the rule and the null sitting inside the crowded left pair. A reader
  % cannot infer the scale's rule from a ladder whose steps are not constant,
  % and fig:mechanism(b), the figure this axis was built to imitate, uses a
  % clean decade ladder at uniform spacing. 0.03 / 0.1 / 0.3 / 1 / 3 / 10 is the
  % textbook 1-3 ladder: steps x3.33, x3, x3.33, x3, x3.33, printing as
  % 1.861 / 1.698 / 1.861 / 1.698 / 1.861cm, a 10% spread. The rule still begins
  % and ends on a drawn tick (S13), every mark still has a printed numeral on
  % each side of it, and no mark sits on a tick. Only the value-to-position map
  % changes; 0.084, 1.563 and 4.180 are byte-identical to passes 9, 10 and 11.
  \def\kx#1{{8.40+(ln(#1)-ln(0.03))/(ln(10)-ln(0.03))*8.98}}

  % PASS 12 REBUILDS THIS PANEL, AND IT IS THE PASS'S ONE STRUCTURAL CHANGE.
  % ------------------------------------------------------------------------
  % WHAT WAS WRONG, measured in the pass-11 PDF and reported independently by
  % all three refuters, two of them as blocking:
  %  (a) THE GAP WAS NOT A GAP. The log axis exists for exactly one reason --
  %      on a log scale a horizontal distance IS a ratio -- and the three marks
  %      stood on three different baselines 0.44cm apart (268.30, 255.83,
  %      243.36pt), so there was no single line to lay a ruler along. Worse, the
  %      155.20pt between the null and the purified slab was filled, one row up,
  %      by "the purified drug slab --- 1.563" (x 363.68-470.69) and, two rows
  %      up, by "cell line, the positive control --- 4.180" (394.28-523.00).
  %      The caption asserted a span the drawing did not draw.
  %  (b) IT READ AS A LEGEND. Each row was a text line TERMINATING in its own
  %      disc, so at print each 1.4pt mark read as a bullet ending a line: three
  %      key entries, stacked in descending value order, i.e. the one
  %      construction S22 forbids outright. Pass 11's own note (3) names this
  %      hazard verbatim and then made all three rows trail instead of one.
  %  (c) THE NULL'S LABEL WAS OUTSIDE ITS OWN PANEL. Its two lines began at
  %      x = 245.04 and 241.61pt against a panel origin of 295.45 -- 19.0mm to
  %      the LEFT of the axis's own left end, in panel (c)'s column, 20pt from
  %      the end of panel (c)'s V_pure band, so "random slab" sat above a band
  %      labelled "the purified drug slab" in the same size and face.
  %  (d) THE ROWS WERE MIS-SET VERTICALLY. The null's first line (baseline 4.05)
  %      and the whole slab row (4.19) fell in one connected ink row; and panel
  %      (d)'s head stood 0.98cm above its first row where panel (c)'s stood
  %      0.48cm above its first content, so two side-by-side panels began at
  %      two depths.
  % ------------------------------------------------------------------------
  % WHAT IS DRAWN NOW. THE THREE MARKS SHARE ONE BASELINE, y = 4.02, at their
  % own three x. Nothing else is set on that line, so the null-to-slab distance
  % is a single uninterrupted 4.52cm horizontal span between two discs and the
  % caption's "a horizontal gap IS their ratio" is true of the drawing. Each
  % NAME sits on its own row above, ending (or, for the null, beginning) exactly
  % at its mark's x, and is joined to its mark by a 0.30pt leader in the mark's
  % own colour -- which is S22's own prescription for a label more than 3mm from
  % its mark, and which is what stops any row reading as a key entry, because a
  % mark tied to its name by a drawn leader is not a bullet.
  % PASS 11'S O-2 IS SUPERSEDED, NOT REVERSED. O-2 refused "one baseline with
  % the labels staggered above" on the ground that the three names measure about
  % 11cm against a 9cm axis. That is still true and the labels are still
  % staggered on three rows; what O-2 got wrong is that the MARKS have to
  % stagger with them. They do not: they are 4.52cm and 1.52cm apart and cannot
  % collide.
  % PASS 11'S O-3 IS SUPERSEDED. O-3 refused drop lines because three stems to
  % the spine would have lengths rising with value and would read as a bar
  % chart. These leaders do not touch the spine and their lengths are 1.07,
  % 0.67 and 0.27cm against values 0.084, 4.180 and 1.563 -- deliberately NOT
  % monotone in value, which is why the row order is null / cell line /
  % purified slab and not the pass-11 descending-value order. The row order is
  % also the only one in which no leader crosses another row's label: the
  % null's leader at 9.99 passes left of both rows below it, and the cell line's
  % at 16.03 passes right of the slab row's right end at 14.51.
  % S21, ANSWERED EXACTLY RATHER THAN APPROXIMATELY. A numeral printed near a
  % value axis is read at the position it occupies. Because each label is
  % aligned so that its numeral's INNER edge sits on its own mark's x, each
  % printed value now stands at its own axis position to the character:
  % "0.084" begins at \kx{0.084}, "4.180" ends at \kx{4.180}, "1.563" ends at
  % \kx{1.563}. That is why the null's row is value-first and the other two are
  % value-last -- the value always abuts the mark, and which side the name falls
  % on is set by the measure, the null's mark being only 1.59cm from the axis's
  % left end. No printed value is now within 1.36cm of the spine, so the rule
  % does not carry numerals on both sides.
  \def\my{4.02}
  % row one: the null.  Label runs RIGHT from its mark; the value leads.
  \node[qlab] at (\kx{0.084},5.26) {0.084 --- matched-dimension random slab};
  \draw[quietink, line width=0.30pt] (\kx{0.084},5.16) -- (\kx{0.084},4.09);
  % the null is an OPEN disc, white fill, 0.8pt edge -- 4.11(b)'s own encoding
  % for this same object (mechanism.py: mfc="white", mew=1.1, color=INK).
  \draw[quietink, line width=0.8pt, fill=white] (\kx{0.084},\my) circle (1.4pt);

  % row two: the positive control.  Label runs LEFT; the value trails.
  \node[qlabe] at (\kx{4.180},4.86) {cell line, the positive control --- 4.180};
  \draw[quietink, line width=0.30pt] (\kx{4.180},4.76) -- (\kx{4.180},4.09);
  \fill[quietink] (\kx{4.180},\my) circle (1.4pt);

  % row three: the object under test.
  \node[font=\scriptsize, text=accent, anchor=base east]
    at (\kx{1.563},4.46) {the purified drug slab --- 1.563};
  \draw[accent, line width=0.30pt] (\kx{1.563},4.36) -- (\kx{1.563},4.09);
  \fill[accent] (\kx{1.563},\my) circle (1.4pt);

  % THE TICK LADDER. PASS 11: six numerals, not four, and the outer two are the
  % domain's own ends, so the rule now BEGINS and ENDS on a drawn tick. At four
  % ticks the rule overran its tick lane by 36.87pt (13.0mm) at each end -- 29%
  % of the axis carrying no numeral, the exact defect S13 measures and condemns
  % in fig:comparators row (d) -- and BOTH extreme marks fell in those bare
  % zones, so neither the null nor the positive control could be placed against
  % a printed number. 0.05, 0.1, 0.3, 1, 3, 6 still reads as a log ladder and
  % puts a numeral either side of every mark.
  % The numerals are INK, not quietink. quietink is this document's comparator
  % and null colour (S8/CF4); spending it on the scale while the figure's actual
  % null is labelled in ink inverted the code. fig:mechanism sets every tick
  % numeral in ink and reserves quietink for its annotations.
  \draw[ink, line width=0.40pt] (8.40,3.44) -- (17.38,3.44);
  \foreach \v in {0.03, 0.1, 0.3, 1, 3, 10}{
    \draw[rulegrey, line width=0.30pt] (\kx{\v},3.44) -- (\kx{\v},3.34);}
  \foreach \v/\t in {0.1/0.1, 0.3/0.3, 1/1, 3/3}{
    \node[tick] at (\kx{\v},3.28) {\t};}
  % the two domain-end numerals are pulled flush INSIDE the rule's own ends, so
  % no ink crosses the drawn measure (S14). Centred on their ticks they would
  % have overhung x = 17.38 by 0.09cm and the float overfull by 1.88pt.
  % PASS 12 CONSIDERED AND OVERRULED centring all six and cutting the rule to
  % 17.10 to pay for it. At the LEFT end that is not available: "0.03" centred
  % on 8.40 begins at 8.13, which is 0.27cm outside panel (d)'s x origin, and
  % S13's one-grid rule -- measured at a 0.2pt spread by pass 11's repair (15)
  % -- outranks the optical centring of a domain-end numeral that labels no
  % mark. The residual is 0.275cm on "0.03" and 0.16cm on "10", which on this
  % ladder is 7% and 4% of a decade.
  \node[tick, anchor=north west] at (8.40,3.28) {0.03};
  \node[tick, anchor=north east] at (17.38,3.28) {10};

  % THE AXIS BLOCK. Quantity and unit on line one, scope on line two, both
  % left-aligned to the axis's own left end, directly under the tick numerals.
  % No result, no interval disclaimer and no cross-reference may sit here.
  \node[ttl]  at (8.40,2.68)
    {KL from the clean forward pass (nats), logarithmic scale};
  \node[qlab] at (8.40,2.34)
    {mean over the target tokens of 60 teacher-forced prompts, layer 4};

% ###########################################################################
% # THE FIGURE'S ONE INTERVAL DECLARATION.  x origin 0.00
% ###########################################################################
  % PASS 12: SET ON TWO LINES, NOT ONE. As one line it measured 15.89cm and was
  % the WIDEST OBJECT IN THE FIGURE -- wider than the axis (8.98), wider than
  % band one (14.54) -- which S14's check forbids outright ("the widest object
  % must be an axis or a panel, never a label"), and a 15.89cm line of 8pt grey
  % stranded a centimetre below everything else reads as a second, smaller
  % caption set above the real one. S17 expressly allows two lines. Broken at
  % the colon, the block measures about 7.1cm, and its two baselines are set so
  % the picture's lowest ink is unchanged and the height cap is not touched.
  \node[qlab] at (0.00,1.04)
    {Intervals: none in (d) --- the stored dispersion is a normal-approximation};
  \node[qlab] at (0.00,0.70)
    {standard error, not this document's convention. (a) to (c) carry no measured mark.};

  % bounding box, so the picture is exactly the measure it claims
  \path (0,0.58) rectangle (17.38,10.62);
\end{tikzpicture}
%
% ---------------------------------------------------------------------------
% READY TO PASTE into Sections/04b-representation.tex, at sec:res-used,
% immediately after the paragraph "Step four, the causal measurement in logit
% space" (04b:261-275). The caption below is the PASS 11 replacement AS REVISED
% AT PASS 12. It is also staged, keyed by \label{fig:hook}, in
% figs/_PENDING_CAPTIONS.tex -- a file PASS 11 CREATED, because the pass-10 note
% in this position asserted it existed and it did not. NO BUILDER EDITS
% Sections/. The staged copy carries the pass-12 revision note; the two copies
% are kept byte-equal in the \caption itself and figs/_PENDING_CAPTIONS.tex is
% the authoritative one.
% PASS 12 REVISED THREE THINGS IN IT, none of them a numeral already there:
%   * panel (d)'s marks are described as standing on ONE LINE, tied to their
%     names by leaders, because pass 11's three baselines meant the "horizontal
%     gap IS the ratio" sentence described a span that was not on the page;
%   * the POSITIVE CONTROL's size confound is disclosed for the first time
%     (cell line's subspace fraction is 20.9x the purified slab's at layer 4,
%     15-22x over the seven -- 04b:370), beside the null's 5.6x, because the
%     37% was being quoted against a comparator matched on nothing at all;
%   * the random-draw disclosure is strengthened from "one random draw" to the
%     fact it understated: with two stored draws at seven layers the file holds
%     fourteen ratios and the drawn 18.634 is the largest of them, and the two
%     draws swap order with depth (at layer 6, 14.7x undrawn against 5.7x
%     drawn).
% Two figures inside the DELETIONS block of the staged copy are superseded by
% the new log domain: the null-to-slab gap is 4.510cm, not 5.474, and the
% slab-to-cell-line gap is 1.520cm, not 1.845. No value moved; the map did.
%
% \begin{figure}[htbp]
% \begin{fullwidth}
% \centering
% % ===========================================================================
% figs/hook.tex -- fig:hook, step four of sec:res-used drawn as what it
% physically is: one true sentence, laid out as a token strip, scored twice by
% the same weights, with a forward hook that zeroes the drug slab at every
% position in between, and the KL read only over the target tokens.
%
% FIGURE BODY ONLY. No preamble, no document environment, no float, no caption.
% \input it from inside a figure/fullwidth float in
% Sections/04b-representation.tex, at sec:res-used, immediately after the
% paragraph "Step four, the causal measurement in logit space" (04b:191-203).
% The complete float block sits COMMENTED at the foot of this file.
%
% Compiles against thesis_v2/preamble.tex and nothing else: no \includegraphics,
% no external data, no TikZ library beyond the ones preamble.tex already loads
% (this file uses arrows.meta and decorations.pathreplacing).
%
% THE MISREADING THIS FIGURE EXISTS TO KILL
% -----------------------------------------
% A reader who has met "ablation" in a generative setting arrives with a default
% picture: the model writes a sentence, the hook is switched on, the model
% writes a DIFFERENT sentence, and the KL measures how far the two sentences
% drifted apart. Every clause of that is false here. Nothing is sampled. The
% sentence is fixed, real, and identical in both passes. And the hook acts on
% the prompt as well as on the target, so it is not an intervention "on
% generation" at all. The figure is built so that all three corrections are
% legible from the marks alone: the two strips are drawn by ONE macro so they
% are provably identical box for box; the nine hook ticks include the four that
% land on the prompt; and the two braces span the target boxes only and stop
% dead at the prompt/target divider.
%
% EVERY NUMBER DRAWN, AND WHERE IT COMES FROM
% -------------------------------------------
% Source: RESULTS_cluster/workspace_probe.json, key layers.4.kl. Config there:
% tier tier2_unseen_drugs, n_drugs 12, n_per_drug 40, n_dims 10,
% n_kl_prompts 60, seed 42. Sibling ledger of the exact values drawn:
% thesis_v2/figs/hook.json.
%
%   mark                    key path                       source value      drawn
%   purified drug slab      layers.4.kl.pure_drug[0]       1.5631485521793365  1.563
%   matched random null     layers.4.kl.random_matched[0]  0.08388670384883881 0.084
%   cell-line control       layers.4.kl.cell_line[0]       4.179661695162455   4.180
%   n per entry             layers.4.kl.*[2] = config.n_kl_prompts             60
%   slab rank               layers.4.pure_drug_dims = config.n_dims            10
%
%   The three dots are the ONLY data-bearing coordinates in the figure. Their
%   positions are \kx{v} = 8.20 + v/4.5*8.70, a linear map of nats to cm:
%     0.084 -> 8.3624 ,  1.563 -> 11.2218 ,  4.180 -> 16.2813
%
%   COMPUTED, not stored as floats:
%   ratio to the null       pure_drug[0]/random_matched[0] 18.634044258027835  18.6x
%   share of the control    pure_drug[0]/cell_line[0]      0.3739892522853049  37%
%   Both rounded forms are also embedded verbatim in the prose string
%   layers.4.ablation_note, which reads "...18.6x the functional variance of a
%   matched-random slab (37% of the cell-line control)...".
%
%   NOT drawn: layers.4.kl.pure_drug[1] = 0.021199743463084047, and its two
%   siblings. These are normal-approximation standard errors over the 60
%   prompts, not this document's interval convention (95% bootstrap clustered
%   over cell lines). figs/fig-mechanism.json already records intervals_drawn
%   false for exactly this reason, and mechanism.py's header states it: drawing
%   a +/-0.003 whisker beside the thesis's genuine +/-0.03 intervals would
%   invite a reader to distrust all of them. The three marks are therefore bare
%   solid dots, which is how this document says "no interval", and the caption
%   says why.
%
%   2048 (the residual width) is prose only: 04b-representation.tex:383 and
%   07-Appendix.tex, sec:app-identity. It appears here as a LABEL ("ten of the
%   2048 directions"), never as a coordinate. Nothing in the vector panel is to
%   scale and the caption says so.
%
% AGREEMENT WITH THE TEXT (checked, not assumed):
%   04b-representation.tex:229  tab:workspace layer-4 row "1.563 & 0.084 &
%                               $18.6\times$"                            -> OK
%   04b-representation.tex:193  "projects the slab out at every position
%                               ($h \leftarrow h - (hV)V^{\top}$)"
%                               -> the formula is set verbatim and the nine
%                                  ticks are what "every position" means.  OK
%   04b-representation.tex:195  "mean KL divergence from clean to ablated over
%                               the target tokens, on $60$ prompts"
%                               -> the braces span the target only; n = 60 is
%                                  printed under the axis.                OK
%   04b-representation.tex:195  \gloss{teacher forcing}{the true sentence is
%                               scored token by token, never generated}
%                               -> the band-one sentence and both subtitles. OK
%   04b-representation.tex:200-202 "The ablation runs for $V_{\textrm{pure}}$
%                               (headline), a matched-dimension random subspace
%                               (the null), and the cell-line slab (the positive
%                               control)" -> the three dots are named in exactly
%                               those three roles, in that order of size.  OK
%   04b-representation.tex:254  "3.6--18.6x" -> layer 4 IS the 18.6.        OK
%   04b-representation.tex:260  "9--37% of what removing cell line costs at
%                               layers 2--9" -> layer 4 IS the 37%.        OK
%   04b-representation.tex:286-288 "No intervals are drawn ... normal-
%                               approximation standard errors over the 60
%                               prompts" -> inherited here.                OK
%   03-Apparatus.tex:249-252    "its expressed panel genes in descending order,
%                               terminated by \ENDCELL. The prompt supplies cell
%                               line, drug, dose, mechanism and the control
%                               cell's sentence" -> the five prompt fields, in
%                               that order, and the terminator.            OK
%   The layer-4 cell-line KL, 4.180, is printed in NO sentence of the thesis;
%   the section quotes only the ratios built from it (04b:258-262). It is a new
%   reading of a stored field and hook.json flags it as such. Precedent:
%   figs/slab.json does the same for variance_share.context.
%
% DESIGN DECISIONS, each with its reason
% --------------------------------------
%  1. BOTH STRIPS ARE DRAWN BY ONE MACRO, \tokstrip, taking only the y of the
%     box bottoms. Two hand-written strips would be two chances to differ, and
%     the single claim the figure has to carry is that they do NOT differ. The
%     macro makes "identical box for box" a property of the source, not of my
%     care.
%  2. THE FIVE PROMPT FIELDS AND THEIR ORDER ARE NOT MINE. They are
%     03-Apparatus.tex:250-252 verbatim, and the same five in the same order as
%     figs/licensing.tex's strip, so the two figures read as one object seen
%     twice. A sixth field would be unsourced and would break that.
%     PASS 9: the two strips now also share a GLYPH. A reviewer found that the
%     words agreed and the drawing did not -- hook drew rounded panelbg boxes
%     with rounded gutters and licensing drew square-cornered WHITE cells
%     sharing edges, so a reader met the same object twice and was shown two
%     objects. hook's treatment was judged the better mark and licensing.tex
%     adopted it (its cell/.style now carries fill=panelbg, rounded
%     corners=1pt, and its drug cell \tokdrug's accent 0.4pt border). Nothing
%     in THIS file changed for it; the cell edges still differ, because the two
%     strips are 8.20cm and 7.20cm and each distributes its own slack.
%  3. THE ELLIPSIS BOX is the mitigation for this figure's largest honesty
%     exposure. The real target is a cell sentence of hundreds of gene tokens;
%     four boxes is a schematic. "no token count is implied" appears in the
%     caption and in the ledger. Rejected alternative: drawing twenty boxes,
%     which implies a count just as firmly and only moves the lie.
%  4. THE DRUG BOX IS THE ONE TINTED PROMPT BOX (accent!12, accent rule). It is
%     the one place the prompt names the drug in plain text -- the fact that
%     makes sec:methods-probe's positive nearly free -- carried forward here so
%     the reader sees the hook deleting downstream of it. It is also the only
%     accent inside either strip, so the strips themselves stay furniture.
%  5. THE HOOK RULE SITS ABOVE THE PASS-TWO STRIP, not below it, so the eye
%     meets the intervention before it meets the boxes it acts on. The RULE is
%     what says "everywhere along here": it spans the strip's full width,
%     prompt half included, which is the correction to the reader's default
%     assumption that an intervention on generation touches only what is
%     generated.
%     REPAIRED, PASS 9: the nine descending TICKS are gone, replaced by one
%     Latex arrow into each half at 0.5pt, stopping 0.06cm above the box tops.
%     Two faults, found independently by both reviewers. (a) A tick per BOX is
%     not a tick per POSITION. "control cell sentence" is one box and hundreds
%     of positions; the ellipsis box is an unknown number. Counting a comb and
%     mapping it to positions is the first thing a reader does with a comb, and
%     it gave the wrong answer -- while the caption simultaneously had to say
%     "the token strip is schematic and implies no token count". (b) The ticks
%     were accent at nearly the rule's own weight and each terminated exactly
%     on its box's top border with no gap; for the DRUG box, whose border is
%     also accent, tick and border fused into one continuous stroke, so the one
%     cell that must be found first sprouted a flagpole and stopped reading as
%     a highlighted field. Marks running into rules uninterrupted is precisely
%     what the house rule forbids. The comb also read as a categorical axis,
%     which is exactly what the WHAT COMES OUT panel two bands below is. Two
%     arrows say "down into both halves" and assert no count; the words "at
%     every position, prompt and target alike" under the strip carry the claim.
%  5b.NO SEPARATE TAG ON THE RULE. One reading "the hook" was drawn at
%     (8.70,7.40) and taken out on the next build: PASS TWO's subtitle now ends
%     "...the hook is on one layer's residual stream" on baseline 7.58 and runs
%     to x = 13.5, so a tag 0.18cm beneath it collided with it and named the
%     same object a third time. The rule is named above by that subtitle and
%     below by the accent formula.
%  6. THE RULE RUNS UNBROKEN ACROSS THE 1.20cm GAP between the prompt strip and
%     the target strip. The gap is a reading device (it is where the divider
%     is); the hook does not know about it. Breaking the rule there would draw
%     the exact exemption the figure denies.
%  7. THE TWO BRACES SPAN THE TARGET SIDE ONLY. That the brace covers the
%     target and no more is the whole answer to "why is the KL taken over the
%     target tokens", and it is drawn rather than said.
%     PASS 9: they now start at the first target box's own left edge, 9.20, and
%     not at the prompt/target divider, 8.70. The caption says "the braces span
%     the target tokens only" and a reader is invited to test that against
%     "prompt boxes included" for the hook rule immediately above, so a brace
%     that began 0.50cm before the first token was contradicting its own label.
%     They are rulegrey, not accent:
%     they are furniture that delimits, and the colour budget is spent on the
%     hook, the KL arrow and the in-slab component.
%  8. THE KL ARROW IS THE ONLY STROKE ABOVE 1.0pt IN THE FIGURE (1.1pt).
%     Nothing else exceeds 1.0pt, so weight alone says which mark is the claim.
%     It runs between the two brace LABELS' baselines, so it visibly joins two
%     readings of the same tokens rather than two sentences.
%     REPAIRED, PASS 9: it was DOUBLE-headed and is now single-headed, tail at
%     the pass-one readout and head at the pass-two readout, matching the
%     direction its own label carries. KL(clean || ablated) is asymmetric; a
%     two-way head says "the difference between", i.e. that the two passes are
%     interchangeable, which is the one thing the measure is not. It was also
%     the figure's heaviest stroke saying the opposite of its label. No
%     double-headed arrow exists anywhere in the corpus: arrows.meta is used
%     only as -{Latex} for a displacement and -{Stealth} for a pipeline edge.
%  9. THE VECTOR PANEL DRAWS THE SLAB AS A BAND WITH THICKNESS, not as a line.
%     A line would say "one direction"; the slab is ten. Nothing in that panel
%     is to scale (a true 10-of-2048 slab is 0.5% of the squared norm and
%     cannot hold a labelled decomposition -- figs/slab.tex draws that true
%     angle, once, and this panel does not pretend to repeat it).
%     REPAIRED, PASS 9, AND THIS WAS THE PANEL'S WORST FAULT: THE RATIO WAS
%     THE WRONG WAY ROUND. h was drawn at 30 degrees above the band, so the
%     REMOVED component (hV)V^T measured 2.90cm and the SURVIVING remainder
%     h-(hV)V^T measured 1.70cm -- on the render, 228px against 130px, the
%     ten-dimensional slab drawn 1.75x LONGER than the 2038-dimensional
%     remainder, under a subtitle reading "ten of the 2048 directions". Cold,
%     the slab read as the dominant subspace and the residual as a leftover,
%     the exact inversion of the section's finding. h is now steep: component
%     0.48cm, complement 1.90cm, 4.0:1 the other way. It is deliberately NOT
%     the true 14.3:1 (86 degrees, cos^2(86) = 0.0049 = 10/2048): at that angle
%     the component is 0.13cm, shorter than its own arrowhead. slab.tex inset
%     one draws that geometry true and this panel defers to it in the caption.
%     Three smaller repairs in the same pass. (a) (hV)V^T was the only one of
%     the three vectors drawn HEADLESS, so in a decomposition the component
%     that is removed was the one not drawn as a vector and read as a baseline;
%     it now carries the same Latex head as its siblings, and its label carries
%     its fate, "-> 0", so the zeroing is a mark and not only the sentence
%     underneath. (b) Its label's 0.1cm rulegrey stub leader is deleted: it was
%     shorter than the 0.3cm threshold that triggers a leader and in the wrong
%     colour. (c) The construction is centred on the band's midpoint at 3.40,
%     because at 2.30 it crowded into the band's left third with 3.1cm of empty
%     band to its right.
%  9b.THE BAND IS panelbg WITH A rulegrey OUTLINE, not accent!15.
%     PASS 9, a cross-figure fix. Across this slate a PALE accent tint meant
%     "raw, still confounded" in fig:purify and "purified" here, while SOLID
%     accent meant V_pure in fig:purify and in the published slab.tex but "the
%     component being removed" here: the tint semantics were inverted between
%     two figures a few pages apart. And inside this figure alone, accent!12
%     (the drug field, both strips) and accent!15 (this band) sampled as
%     (225,234,242) and (218,229,238) -- a 3% tint difference that will be
%     indistinguishable on press, two meanings in one apparent colour. panelbg
%     with a rulegrey outline is slab.tex's own treatment for "a region of the
%     stream" (its \slabdisc, slab.tex:83-90), it leaves accent!12 uniquely
%     "the drug field" across licensing and hook, and it keeps solid accent for
%     the slab itself across the slate.
%     ONE RESIDUAL DIFFERENCE, AND IT IS DELIBERATE: fig:purify draws
%     V_pure as a SOLID accent bar and this figure draws it as a cream region.
%     The two are doing different jobs. In fig:purify V_pure is one subspace
%     among three and is the object every causal claim is made on, so it takes
%     the accent, exactly as slab.tex's published wedge does; here it is the
%     backdrop a vector is decomposed against, and an accent component has to
%     be legible ON it, which a solid accent fill would make impossible. The
%     tie between the two figures is the direct label, which is accent
%     $V_{\textrm{pure}}$ in both. Giving this band an accent OUTLINE instead
%     was considered and rejected: at panelbg fill with an accent 0.35pt border
%     it is the drug box's own treatment (accent!12 fill, accent 0.4pt border)
%     to within a tint, and one glyph would then mean two things again.
%  9c.ONE NAME FOR ONE OBJECT. The panel said V, the arrows said (hV)V^T and
%     the band said V_pure -- three names for the same ten directions, in a
%     figure that sits a page from fig:purify where V_pure is built. The
%     subtitle now states the identity once: "V = V_pure: ten of the 2048
%     directions".
% 10. THE RIGHT-ANGLE MARK AT THE ORIGIN is load-bearing, not decoration. It is
%     what distinguishes "the slab component is set to ZERO" from "the slab
%     component is reduced". Without it the vertical arrow is just a shorter
%     arrow.
% 11. THE READOUT AXIS IS LINEAR AND CARRIES NO BARS. It is one layer, so the
%     reason mechanism.py forbids a bar on a log axis (a bar's length is set by
%     wherever the axis is cut) does not arise -- and there are no bars in any
%     case, only three dots. Rejected alternative: seven layers on this axis,
%     which would rebuild fig:mechanism panel (b) inside a mechanism figure and
%     make the readout compete with the mechanism for the reader.
%     REPAIRED, PASS 9: EVERY DOT NOW CARRIES ITS VALUE. None of 1.563, 0.084
%     and 4.180 was printed on the drawing, and \kx{0.084} = 8.362 against a
%     zero tick at 8.20, so the null sat 1.6mm from zero and, under a 2.1pt
%     halo and a 1.4pt fill, was visually AT zero -- while the strip below
%     asserted "18.6x the null". A reader who takes the null for zero takes
%     18.6x for a claim they cannot check. Both reviewers proposed the same two
%     remedies, printing the values or moving to a log axis with a joining
%     stem; the first is taken, because it is the document's own no-legend
%     rule, it makes the ratio checkable from the marks, and it keeps the axis
%     comparable in kind with tab:workspace's nats column. The three
%     descriptions were also three near-full-width stacked lines in an order
%     that was neither the axis order nor a reading order -- the top line named
%     the rightmost dot, and the middle line, spanning crop x 275-1630, sat
%     directly over the accent dot at 770, so "the null" attached to the drug
%     slab on first pass. They are now short and value-first, each near its own
%     dot; the long descriptions live in the caption, which already carried
%     them. The three dots already sat over 2.1pt white halos in ladder.tex's
%     order (spine, halos, fills, labels) and still do; a reviewer reported
%     them as halo-less and that finding was overruled on the source.
% 12. NO MARK JOINS THE KL AXIS TO ANY ACCURACY OR SUBSPACE-FRACTION SCALE, and
%     no double-headed arrow spans two rules anywhere below the divider. Those
%     are different instruments with different comparators; fig:comparators is
%     the figure that owns that argument.
% 13. LAYER 4 IS THE ONE LAYER DRAWN because it is the slab's strongest showing
%     of the seven against the null (18.6x, the top of the thesis's 3.6-18.6x)
%     and the top of the 9-37% band against cell line. Drawing the best case
%     and still showing the accent dot sitting at a third of the ink dot is the
%     honest form of the claim. Layer 12 was the alternative and is not drawn:
%     there the drug and cell-line costs are nearly equal (0.133 against 0.135,
%     98%), which is the section's own named exception and would read as the
%     rule if it were the only layer shown. The caption says both, and points
%     at tab:workspace for all seven. hook.json records layer 12's three values.
% 14. THE TOP-RIGHT NOTE SAYS "AT ONE LAYER". The word "every" in "at every
%     position" ranges over the SEQUENCE, not over depth, and a reader who
%     meets "layer 4" only in the readout panel has no way to tell which. The
%     misreading -- that the slab is projected out at every layer as well as at
%     every position -- is exactly the kind this figure exists to prevent, and
%     no mark can carry it, so it is stated in words. It is a design fact from
%     the method (04b:164-171, one measurement per layer, one row per layer in
%     tab:workspace; 04b:191-195, the hook itself), not a stored quantity, and
%     hook.json records it as such. Set quietink at \scriptsize, because it is
%     furniture; that it also fills the figure's one empty quadrant is a
%     consequence and not the reason.
%
% WHAT SIX BUILDS CHANGED, with the measurement that forced each
% -------------------------------------------------------------
% (a) EVERY TEXT NODE IS anchor=base, so each y below IS a baseline. Build 1
%     used anchor=north west with an assumed 0.34cm node height. The real
%     height of a \footnotesize node at TikZ's default inner sep (0.3333em,
%     2.7pt a side) is 12 + 5.4 = 17.4pt = 0.61cm, nearly twice that, and the
%     consequences were visible: "the true sentence is teacher-forced, never
%     generated" put its descenders through the pass-one box tops, and the hook
%     rule at y = 7.52 ran clean through the middle of "the same prompt, the
%     same sentence, the same weights". Baseline anchoring removes the guess:
%     ink sits about 0.20cm above and 0.06cm below a \scriptsize baseline, and
%     0.25/0.08 at \footnotesize. frames.tex records the same fix for the same
%     reason, plus a second one -- a node with align= or text width builds a
%     minipage whose first line follows the AMBIENT \baselineskip, so its ink
%     moves with the prose before the float. No node here uses either.
% (b) THE PROMPT/TARGET DIVIDER MOVED OFF THE BOX EDGE. Build 1 put it at
%     x = 7.80, which is exactly the right-hand border of the "control cell
%     sentence" box: a 0.3pt rulegrey line drawn on top of a 0.35pt rulegrey
%     box border is invisible as a separate mark, and the requirement is that
%     the braces stop VISIBLY at it. It now sits at x = 8.40, the midpoint of
%     the 1.20cm gap between the strips, where nothing else is drawn, and it is
%     set as `ink, 0.5pt, densely dotted' -- this document's mark for a mere
%     coordinate worth naming, which is exactly what the boundary is. Both
%     braces now begin at 8.40 and their left tips land on it.
% (c) THE BRACE AMPLITUDE WENT FROM 2.5pt TO 3.5pt AT 0.35pt. At 2.5pt and
%     0.3pt, under a rounded-corner box border 0.10cm above it, the brace read
%     as a second box outline rather than as a brace; slab.tex and ladder.tex
%     both use 3.0-3.5pt for braces that have to be seen.
% (d) THE THREE READOUT LABELS WERE STAGGERED BY 0.50cm, not 0.33cm. At 0.33cm
%     "the purified drug slab" and "a matched-dimension random subspace --- the
%     null" overlapped by about 0.08cm of ink, because both labels are wide
%     enough to share x and the node boxes are taller than a line. 0.50cm gives
%     0.24cm of clear space between successive label blocks, which is also what
%     lets the two leaders leave their dots at distinguishable angles.
% (e) "ten of the 2048 directions" AND THE VECTOR PANEL'S CLOSING SENTENCE
%     overlapped outright in build 1 (the first sat at y 1.05-1.35, the second
%     at 0.95-1.15). The label was DELETED rather than moved: the panel's own
%     subtitle already reads "$V$ holds ten of the 2048 directions", so the
%     node under the band was saying a second time what the subtitle says
%     first, and deleting it also freed the space the closing pair needed.
% (f) THE STRIP WAS WIDENED FROM 13.20 TO 13.55 and the KL arrow moved from
%     13.75 to 14.05. At the narrower width the top-right quadrant, roughly
%     4.1 x 1.7cm, held nothing at all, and the KL label was the only mark
%     right of the strip. The right column now carries the "one layer" note as
%     well (decision 14) and band one reaches x = 16.5, which is where band
%     two's readout already ended.
% (g) THE IN-SLAB COMPONENT'S LABEL GOT A LEADER. $(hV)V^{\top}$ names a
%     segment that lies INSIDE the slab band, so the nearest place its label
%     can sit is 0.28cm above the band's top rule and 0.44cm above the mark
%     itself -- well past the ~0.3cm at which this document's rule requires a
%     leader. Built without one it read as floating in the interior of the
%     triangle. It now has a 0.14cm accent!45 hairline from just above the
%     segment to just under the label, which is the same device comparators.tex
%     uses for its 0.121 mark.
% (h) THE 18.6x / 37% LINE MOVED FROM quietink TO ink. Under the axis it was
%     the third of four grey centred lines and read as a fourth piece of axis
%     apparatus. It is a reading of the three marks, so it takes the register
%     this document gives an estimate rather than the one it gives furniture.
% (i) THE AXIS SPINE WAS CUT FROM 16.90 TO 16.60. The map's range top is 4.5
%     nats, but the largest mark is 4.180 and the last tick is 4, so a full-
%     range spine left 0.97cm of empty rule past the last tick and read as an
%     axis that had run out rather than as headroom. No coordinate moved.
% (j) THE PASS-ONE BOUNDARY RULE GREW A TAIL, from 8.72 down to 8.55. Its
%     label sits to its left on the brace label's baseline, and with the rule
%     stopping 0.14cm under the boxes the label read as a caption of the prompt
%     half rather than as a pointer at the boundary. The pass-two copy keeps
%     the symmetric 0.14/0.14, because it carries no label.
%
% LAYOUT NOTES, all forced by a build
% -----------------------------------
%   * fullwidth = textwidth 142mm + marginparwidth 28mm + marginparsep 4mm
%     = 174mm. (preamble.tex's own comment beside \newenvironment{fullwidth}
%     says 171mm and is stale arithmetic; the geometry keys give 174mm.) The
%     box is PINNED at 17.30 x 10.00cm = 173.0 x 100.0mm, 1.0mm inside the
%     measure. Rightmost ink is the axis spine at 16.90.
%   * the spec's token x-edges (0.00/1.20/2.20/3.00/4.30/7.30) do not hold
%     their labels at \scriptsize: "cell line" measures about 1.27cm of ink and
%     the box was 1.20cm, and "mechanism" the same against 1.30cm. The five
%     boxes were rewidened to 1.55/0.85/0.85/1.70/3.25, keeping the order and
%     keeping "control cell sentence" the widest field, and the prompt strip
%     now ends at 8.20 rather than 7.30. Everything keyed to the boundary (both
%     dividers, both braces, the target strip, the KL arrow) moved with it.
%     The five widths are set against figs/licensing.tex's measured ink for the
%     SAME five labels at the same size (cell line 0.954, drug 0.616 bold, dose
%     0.579, mechanism 1.424, control cell sentence 2.511; total 6.084cm), so
%     the clear space each side runs 0.117 to 0.370cm and no cell is near the
%     0.06cm floor. A first pass gave mechanism only 1.55cm, i.e. 0.063 a side,
%     while control cell sentence had 0.445; 0.15cm moved from the last cell to
%     the fourth evens that out. \ENDCELL's cell keeps 0.22 a side.
%   * the target strip is four boxes, so the strip is nine boxes and the hook
%     carries nine ticks. \ENDCELL is \texttt{[END\_CELL]}: ten inconsolata
%     characters, about 1.41cm at \scriptsize, so its box is 1.85cm.
%   * the prompt/target GAP is 1.00cm (8.20 to 9.20) and the boundary rule sits
%     at its midpoint, 8.70. Both braces begin there and their left tips land
%     on the rule, which is what "the brace stops dead at the divider" means
%     when it is drawn rather than described.
%   * nothing is scaled. 1cm here is 1cm on the page; there is no scale= key
%     and no \resizebox. style_v2.py records that v1's figures were silently
%     scaled to 0.887 by LaTeX and that this, not the drawing, read as "off".
%   * the pgfmath trap: \kx expands to a BRACED group, so \kx{1.563}+0.05 is
%     not a coordinate. Where a label needs a shim it is applied as an xshift
%     on the node, or the shim lives inside the braces. (gate.tex's header
%     records the same trap costing a compile.)
%   * \foreach prints its raw token, so the tick labels are given in the
%     \v/\t value/label pair form. Here every tick is a non-negative integer so
%     the pair form is not strictly needed, but it is the corpus idiom and it
%     keeps the axis editable without a surprise.
%
% TYPE
% ----
%   Ambient font=\small on the picture; every node sets its own size. The
%   dominant size is \scriptsize (8pt), with \footnotesize (10pt) only for the
%   four panel subtitles, the band-one sentence and the vector panel's two
%   closing lines. Nothing is set at body size. Registers: \sffamily bold
%   accent = a panel tag; roman accent = the intervention; roman ink = the
%   object being scored or the comparator; roman quietink = furniture and
%   annotation; \ttfamily = only \ENDCELL, which names a real token.
%
% ===========================================================================
% PASS 10 -- 2026-08-18 -- THE STYLE-CONTRACT REBUILD
% ===========================================================================
% EVERYTHING ABOVE THIS LINE IS LEFT INTACT and is the design record of passes
% 1-9. This section supersedes, BY NAME, only the claims above that this pass
% made false; a decision above that is not named here still stands, and the
% reasoning that produced it is still the reasoning this figure is audited
% against. NO NUMBER CHANGED IN THIS PASS: 1.563, 0.084, 4.180, n = 60, rank
% 10, residual width 2048 and layer 4 are byte-identical to pass 9, and the two
% ratios 18.6x and 37% were not altered but MOVED -- off the drawing and into
% the caption, which is where this document's contract puts a result.
%
% WHY THIS PASS EXISTS
% --------------------
% The supervisor's verdict on the 4.7-4.12 slate: the figures look "sloppy and
% childish", they must be CORRECT, and every axis must define what it measures.
% A binding style contract (S1-S25) and a cross-figure code (CF1-CF14) were then
% written against measurements of this document's own two reference figures --
% 4.11, figs/mechanism.py, the BODY standard, and A.3, figs/slab.tex, the
% appendix plate -- and the ruling was: in the body, take 4.11. Measured against
% it, this figure was the worst offender of the five.
%
%                          pass 9            4.11        pass 10 (measured)
%   type sizes >7pt        9                 2           2 chosen + 1 artefact
%   embedded faces         9                 1           6, four of them math/tt
%   italic characters      34%               0%          1.8%
%   stroke weights         9 (0.17, 0.68     --          5
%                          used once each)
%   all-caps sans banners  6                 0           0
%   ink at 200dpi          8.55%             5.3%        5.98%   (cap 6.0)
%   drawn box              173.1 x 97.0mm    --          173.99 x 99.95mm
%
% THE CORRECTNESS DEFECT, WHICH IS WHY THE FIGURE WAS REOPENED AT ALL
% -------------------------------------------------------------------
% THE READOUT AXIS IS NOW LOGARITHMIC. Design decision 11 above chose a linear
% nats axis, and pass 9 mitigated its consequence by printing every dot's value.
% BOTH ARE SUPERSEDED. The mitigation was not enough, and the header itself said
% so in as many words: on the linear map \kx = 8.20 + v/4.5*8.70 the null at
% 0.084 landed 0.162cm -- 1.6mm -- from the zero tick, under a 2.1pt halo and a
% 1.4pt fill, so the null dot WAS the origin at print size while the strip below
% asserted "18.6x the null". A ratio the drawing renders as infinite is not a
% mark a reader can check, whatever is printed beside it. Both quantities this
% figure quotes are RATIOS; on a log axis a horizontal gap IS a ratio, so the
% mark now means what the sentence means. 4.11 panel (b) published exactly this
% solution for exactly these numbers -- log axis, open and filled dots, the
% ratio read as a gap -- and this figure now uses it.
%   NEW MAP  \kx#1 = 8.40 + (ln(#1) - ln(0.05))/(ln(6.0) - ln(0.05)) * 8.98
%            domain [0.05, 6.0] nats, spine 8.40 to 17.38, labelled ticks at
%            0.1, 0.3, 1, 3
%   0.084 -> 9.3731 , 1.563 -> 14.8569 , 4.180 -> 16.7020
%   The null-to-slab gap is 5.4838cm and IS 18.6x. The slab-to-cell-line gap is
%   1.8452cm and IS 37% read the other way. Neither is annotated: the caption
%   carries both, which is where the contract puts a result.
%   The domain is deliberately wider than the data at BOTH ends, so no mark sits
%   on the cut and no bar is possible; 4.11(b) does the same, its ticks running
%   0.01-1.00 inside a taller axis. There are still no bars and there must not
%   be: on a log axis a bar's length is set by wherever the axis is cut.
%   SUPERSEDED: the linear map (0.084 -> 8.3624, 1.563 -> 11.2218,
%   4.180 -> 16.2813); decision 11's argument that "one layer, so a log axis is
%   not needed"; and note (i)'s spine crop, which no longer applies.
%
% THE FOUR MARKS AND LINES THAT WERE DELETED, EACH WITH ITS REASON
% ----------------------------------------------------------------
% (i)   THE TWO SHORT HOOK ARROWS (pass 9's replacement for the nine-tick comb)
%       are gone. They implied a count of two and the caption had to deny it in
%       a sentence of its own ("The two arrows off the rule mark that it acts in
%       both halves and imply no count of positions"). A drawing that needs its
%       caption to unsay it is drawing the wrong mark. The RULE's span is the
%       "everywhere" claim -- decision 5 above is right and survives -- and the
%       accent formula now sits ON the rule, 0.20cm above it, instead of 1.1cm
%       below the strip. This also empties the figure of every arrowhead except
%       the comparison arrow's two and the vector panel's three, which is what
%       S11 asks for: one arrowhead, one meaning.
% (ii)  THE ITALIC EPIGRAM "nothing is generated; the two passes differ only in
%       the hook" and the 17.30cm rulegrey divider that existed only to carry it
%       are deleted. It was 9.96pt Pagella-Italic quietink centred on the
%       figure's own x-centre -- the LARGEST type on the drawing, 25% larger
%       than the 7.97pt at which every measured number was set -- and its first
%       clause restates the caption's opening words, "Nothing is generated:".
%       It is the caption's sentence and the caption already had it.
% (iii) THE PRINTED RESULT "18.6x the null; 37% of what removing cell line
%       costs" is deleted from under the axis. A result may not sit in the
%       axis-label lane, and it was the third of four stacked centred lines
%       there. NOTE (h) ABOVE IS SUPERSEDED: moving that line from quietink to
%       ink so it would stop reading as axis apparatus was the wrong repair --
%       the line does not belong under the axis in any colour. Both ratios are
%       in the caption, which now also carries the qualifier the old caption
%       dropped: "9-37%" is the range over layers 2-9 ONLY; the true seven-layer
%       range of kl.pure_drug[0]/kl.cell_line[0] is 8.90% (layer 9) to 98.45%
%       (layer 12).
% (iv)  "the marks carry no interval; the caption says why" is deleted. It
%       restated the caption's own bold "no interval is drawn" and it was
%       italic. It is replaced by the single figure-foot declaration S17
%       requires, set \scriptsize quietink on the figure's x origin, which NAMES
%       the construction instead of pointing at the caption.
% (v)   "generation would begin here", the prompt/target divider's label, is
%       deleted, and this was the one deletion made for the ink budget rather
%       than for a defect. Recorded here because it was load-bearing and a
%       future pass may want it back. Its claim is drawn three other ways --
%       panel (a)'s head is "pass one, clean", both braces read "scored tokens"
%       so the prompt half is visibly not scored, and the caption's lede opens
%       "Nothing is generated" -- and the divider itself survives as what it
%       always was, a piece of the strip's structure. Cost measured: 27mm2 of
%       ink, against a 6.0% ceiling the figure clears at 5.98%.
%
% WHAT REPLACED THE SIX BANNERS AND THE FOUR ITALIC SUBTITLES
% -----------------------------------------------------------
% Every \sffamily\scriptsize\bfseries accent banner is gone, and with it the
% only sans face in the figure. The apparatus face belongs to the margin column,
% not to a figure read inside a running argument; A.3 may use it because A.3 is
% a reference plate a reader stops at. Panel heads are now 4.11's form exactly
% -- a lowercase parenthesised letter and a sentence-case descriptive title,
% \footnotesize roman ink, on the panel's own x origin, no colon, no bold, no
% caps:
%     (a) pass one, clean
%     (b) pass two, the slab projected out
%     (c) what the hook does to one vector
%     (d) what comes out
% The four \footnotesize\itshape quietink subtitles under them are deleted.
% Their content was distributed, not lost:
%   "the true sentence is teacher-forced, never generated"  -> already the
%       caption's opening clause; it was duplicating it (S18).
%   "the same prompt, the same sentence, the same weights"  -> now the drawn
%       tie between the strips (below).
%   "the hook is on ONE layer's residual stream"            -> now the axis
%       block's scope line, where the layer is named anyway.
%   "V = V_pure: ten of the 2048 directions"                -> the identity is
%       now the band's own direct label; the rank moved to the caption.
%   "layer 4: mean KL from the clean pass, over the target tokens, 60 prompts"
%       -> now the axis block itself, split into a quantity line and a scope
%       line, which is the position S5 requires.
% DECISION 14 ABOVE IS PARTLY SUPERSEDED: the "at one layer" fact is still
% stated in words, because no mark can carry it, but it is no longer a clause of
% PASS TWO's subtitle -- there is no subtitle -- and lives in the axis block's
% second line, the place S5 reserves for scope.
%
% THE TIE BETWEEN THE STRIPS, WHICH IS THE ONLY NEW TEXT IN THE FIGURE
% ---------------------------------------------------------------------
% The two strips are emitted by one macro and are therefore provably identical
% box for box -- the figure's central claim, and its best-executed feature.
% Nothing on the drawing said so. Worse, the two braces beneath them carried
% DIFFERENT labels ("per-token log p, read here only" against "the same tokens,
% read again"), which sends a reader looking for a difference between two
% objects whose sameness is the point. Both braces now carry the SAME label,
% "scored tokens", and one \scriptsize roman ink line reads "the same sentence,
% box for box". It sits UNDER both strips, on band one's own x origin and on the
% same baseline as pass two's brace label, so it is a statement about a pair the
% reader has just seen and it costs the drawing no extra row of ink.
%
% THE COMPARISON IS NO LONGER AN ARROW WITH ONE HEAD
% ---------------------------------------------------
% DECISION 8 ABOVE AND ITS PASS-9 REPAIR ARE SUPERSEDED. Pass 9 made the KL mark
% single-headed, reasoning that KL(clean || ablated) is directed and asymmetric
% and that a two-way head would say the passes are interchangeable. That
% reasoning is about the MEASURE; the MARK is about the two passes, and a naive
% reader read the single head as pass one FEEDING pass two -- generation, the
% precise misreading this figure exists to kill, drawn by the figure itself in
% its heaviest stroke. The two passes are independent scorings of one fixed
% sentence by one set of weights: that is a comparison, and S11 gives a
% comparison a double head or a brace. It is now a DOUBLE-headed ink arrow at
% 1.0pt, and the single accent arrowhead is reserved for an operation acting in
% a direction. It is no longer the figure's heaviest stroke either -- the 1.1pt
% was a tenth-of-a-point singleton in a nine-weight vocabulary -- and its label
% is the registry name, "KL from the clean forward pass", the same words the
% axis in panel (d) carries, because it is the same quantity and S19 gives an
% object exactly one name.
%
% THE AXIS DEFINITION, WHICH IS THE SUPERVISOR'S FIRST COMPLAINT
% --------------------------------------------------------------
% This figure was the only one of the five that already named quantity and unit
% in axis-label position, and note (h) plus decision 11 built that. The WORDS
% survive; the position, the rank and the stack around them change.
%   LINE 1  \footnotesize roman ink, on the axis's own left end, x = 8.40:
%           "KL from the clean forward pass (nats), logarithmic scale"
%   LINE 2  \scriptsize quietink, same x:
%           "mean over the target tokens of 60 teacher-forced prompts, layer 4"
% "Mean KL divergence, clean -> ablated (nats)" is RETIRED. It was centred and
% bold under the rule; more importantly it was a SECOND name for 4.11(b)'s "KL
% from the clean forward pass", and CF3 gives the quantity one name across the
% slate. The old panel subtitle's "layer 4 ... over the target tokens, 60
% prompts" is folded into line 2, so the scope is stated once instead of twice.
%
% THE READOUT IS THREE ROWS, NOT ONE, AND WHY
% -------------------------------------------
% S22 requires that within a row every direct label share one baseline and that
% no label cross the drawn measure. Three labels naming three objects at one
% layer -- "0.084 - matched-dimension random slab" (5.2cm), "1.563 - the
% purified drug slab" (4.2cm) and "4.180 - cell line, the positive control"
% (5.4cm) -- cannot share one baseline: the last two marks sit 1.85cm apart on
% the log map. Staggering labels off one rule is what S22 forbids; giving each
% object its OWN row is not. Each row therefore carries exactly one mark and one
% label on one baseline; the three rows share one log x axis; and they are
% ordered by value with the largest on top, so vertical position carries the
% same order as horizontal position and can assert nothing the axis does not.
% That ordering is 4.11(b)'s own, the filled slab dot above the open null dot.
% The rows sit 0.38, 0.82 and 1.26cm above the spine so the three read as marks
% over one scale. Each label is 0.26cm from its own mark, just inside the 3mm at
% which this document requires a leader, so no leader is drawn. THE LABEL SIDE
% DIFFERS BY ROW AND IT HAS TO: the two large marks are 0.7 and 2.5cm from the
% right end of the measure and their labels can only go left, while the null at
% 9.37 has 0.97cm to its left and its label can only go right. NO DROP LINES
% JOIN THE MARKS TO THE SPINE: a hairline from a dot down to a horizontal axis
% reads as a bar, and a bar on a log axis is forbidden in this document.
%
% MARKS, WEIGHTS, TYPE AND PROSE
% -------------------------------
% MARKS, S9/CF5, three glyphs and no more:
%   filled disc, r = 1.4pt, accent   the purified drug slab -- the object under
%                                    test, and the only accent mark in panel (d)
%   filled disc, r = 1.4pt, ink      cell line, the positive control
%   open disc,  r = 1.4pt, white     matched-dimension random slab, the null.
%               fill, 0.8pt ink edge This is 4.11(b)'s own encoding for this
%                                    same object (mechanism.py: mfc="white",
%                                    mew=1.1, color=style.INK), and adopting it
%                                    is CF5's stated requirement. The 0.8pt edge
%                                    is the mark vocabulary's own weight, not a
%                                    path weight, and it is the only 0.8pt here.
%   The three 2.1pt white halos are DELETED together with the two leaders they
%   existed to cover. Pass 9's overruling of the "no halo" finding stands as a
%   fact about pass 9; the halos are simply no longer needed, because no leader
%   now runs into a dot.
% WEIGHTS, S10. Measured in the built PDF: 0.30 x26, 0.40 x5, 0.80 x1, 0.90 x1,
% 1.00 x4. Five weights, against pass 9's nine.
%   0.30pt  furniture: every token-box border, both braces, the slab band's
%           outline, the right-angle mark, the two dashed construction lines,
%           the four axis ticks
%   0.40pt  the drug cell's accent border, both prompt/target dividers, the
%           axis rule
%   0.80pt  the open disc's edge (S9's own weight for that glyph; see above)
%   0.90pt  the hook rule -- USED ONCE, and justified here as S10 expressly
%           permits: it is a distinct kind of object, the intervention, the one
%           thing that differs between the two passes. It is accent and not
%           black, so it cannot be confused with the chance/gate/margin register
%           S6 reserves for a black solid 0.9pt rule, and this figure draws no
%           threshold at all.
%   1.00pt  panel (c)'s three vectors and the comparison arrow
%   RETIRED: 0.35pt (box borders and braces, folded into 0.30), 0.50pt (the
%   dotted dividers, see below), 1.10pt (the KL arrow), and the 0.17pt and
%   0.68pt singletons, which were never chosen at all -- they were the Latex
%   arrow tips' own stroke, and the tips are now declared line width=0pt so they
%   are pure fills and contribute no stroke weight.
% THE PROMPT/TARGET DIVIDERS are no longer "ink, 0.5pt, densely dotted".
%   NOTE (b) ABOVE IS SUPERSEDED. S7 reserves the densely-dotted 0.5pt stroke
%   for a zero rule, a magnification leader or a projection line, and requires a
%   dotted rule's label to be "0", to be a leader, or to be absent; this one
%   carried "generation would begin here", which is none of those. The line is
%   what it always was -- a piece of the strip's structure that delimits two
%   halves -- so it is now rulegrey 0.40pt solid, a strip weight, still at 7.94
%   in the middle of the empty 1.00cm gap where nothing else is drawn. Note
%   (b)'s actual finding, that a rulegrey line drawn ON a rulegrey box border is
%   invisible as a separate mark, is why it stays at mid-gap.
% TYPE, S1. Measured in the built PDF: 7.97pt x587 characters and 9.96pt x165,
% plus three exempt classes -- 6.85pt x20 (\ENDCELL, Inconsolata, the \texttt
% token name S1 exempts), 5.98pt x8 and 6.23pt x3 (subscripts and the \top
% superscript, which S1 exempts), and 8.30pt x10.
%   THE 8.30pt SPANS ARE NOT A THIRD CHOSEN SIZE and are recorded here so an
%   auditor need not rediscover them: they are ten glyphs -- \leftarrow, two
%   minus signs, six parentheses and \to -- set from CMSY10 and CMR10 inside the
%   two formulas $h \leftarrow h-(hV)V^{\top}$ and $(hV)V^{\top} \to 0$.
%   mathpazo substitutes Palatino for math letters but leaves the CM symbol
%   fonts in place and scales them by 1.041 to match Palatino's x-height, so a
%   \scriptsize (7.97pt) math operator prints at 7.97 x 1.041 = 8.30pt. It is
%   the same nominal size in a differently scaled font, it is a property of this
%   document's preamble and not of this figure, and every figure in the thesis
%   that sets inline math shows it. The one instance that COULD be removed was:
%   the ellipsis token box was $\cdots$ (CMSY10 at 8.30) and is now \dots, set
%   from Pagella at 7.97 like every other token label.
%   Registers: \footnotesize roman ink = the four panel heads and the axis's
%   quantity line, and nothing else, so the largest type on the drawing always
%   names a panel or a measured quantity and never narrates. \scriptsize carries
%   every label, tick, formula, tie and foot line. Zero \sffamily, zero
%   \bfseries, zero \itshape. Italic characters: 14 of 793, 1.8%, and all
%   fourteen are the math variables h and V, which are a symbol's own face and
%   not a voice. Pass 9 measured 34%.
% PROSE, S4: THE FIGURE NOW CARRIES NONE, which is 4.11's own austerity -- 4.11
%   has axis labels and direct labels and nothing else. Pass 9's two closing
%   lines, "the slab component is set to zero; every other direction is
%   untouched", are deleted because the drawing already carries both claims as
%   MARKS: the accent component's label ends "-> 0", the right angle says the
%   component is zeroed and not merely shortened, and the surviving vector is
%   labelled $h-(hV)V^{\top}$. A sentence whose deletion leaves no drawn object
%   unnamed is not definitional and S4 does not permit it. Decision 9c's "one
%   name for one object" survives, carried now by the band's own direct label,
%   "$V_{\textrm{pure}}$, the purified drug slab", which is also the gloss S19
%   requires at the symbol's first use.
%
% THE RIGHT ANGLE, AND THE ONE GEOMETRIC RATIO THAT MOVED
% --------------------------------------------------------
% The right-angle mark at the foot of the surviving vector was 0.16cm (1.6mm)
% and a naive reader did not see it until zooming. It is now 0.40cm (4.0mm) on
% both legs. Decision 10 above is why it may not simply be dropped: it is what
% distinguishes "the slab component is set to ZERO" from "the slab component is
% reduced". A 4mm glyph cannot sit on a 0.48cm arm without swallowing it, so the
% in-slab component grew from 0.48cm to 0.90cm and the surviving complement from
% 1.90cm to 3.14cm. THE RATIO IS UNCHANGED IN DIRECTION AND BARELY CHANGED IN
% SIZE: 3.49:1 in favour of the survivor against pass 9's 4.0:1, still the right
% way round, and still deliberately NOT the true 14.3:1 (86 degrees,
% cos^2 86 = 0.0049 = 10/2048), at which the component would be 0.13cm, shorter
% than its own arrowhead. figs/slab.tex inset one draws that geometry true, and
% the caption still sends a reader there. Nothing in this panel is a coordinate
% and nothing in it is measured.
%
% THE STRIP: SAME MACRO, SAME FIELDS, EVEN CLEAR SPACE
% -----------------------------------------------------
% The nine box edges were reset so that EVERY cell carries the same 0.13cm of
% clear space each side of its label's measured ink, against pass 9's 0.117 to
% 0.370cm. The strip therefore runs 0.00 to 12.65 instead of 0.00 to 13.55, the
% prompt half ends at 7.44 and the target half begins at 8.44, and everything
% keyed to the boundary (both dividers at \bnd = 7.94, both braces, the brace
% labels at 10.545, the hook rule, the comparison arrow at 13.40) moved with it.
% Two reasons, and both matter: the clear space is now a constant rather than a
% number that varied by a factor of three between cells, and the strip's own
% rules are 3.6cm shorter, which is 5mm2 of the 6.0% ink ceiling. The five
% prompt fields, their order, the [END_CELL] terminator and the ONE-MACRO
% construction are untouched. Ink widths the edges are set from, measured in
% figs/licensing.tex at the same size: cell line 0.954, drug 0.616, dose 0.579,
% mechanism 1.424, control cell sentence 2.511, [END_CELL] 1.41.
%
% GEOMETRY, MEASURED IN THE 200dpi RENDER
% ----------------------------------------
% Content box 173.99 x 99.95mm, against a 170-174mm fullwidth measure and a
% 100mm height cap. Ink 5.98% of the content box, against a 6.0% ceiling and
% pass 9's 8.55%; 793 non-space characters over 17393mm2 is 4.6 per 100mm2,
% against a 5.0 ceiling. Picture bbox 17.38 x 10.04cm (y 0.58 to 10.62).
% Leftmost ink is x = 0.00 -- the four panel heads, both strips, the hook rule
% and its formula, the tie, panel (c)'s band and label, and the foot line, all
% on it. Rightmost ink is the axis spine at 17.38. PANEL X ORIGINS, S13: band
% one and panel (c) at 0.00, panel (d) at 8.40. Measured in the built PDF, the
% left edges of the panel-level text lines are 58.34, 58.34, 58.48, 58.34,
% 58.34, 58.36, 58.34pt for the 0.00 panels -- a spread of 0.14pt -- and
% 296.45pt three times over for panel (d), a spread of 0.00pt.
%
% ===========================================================================
% PASS 11 -- 2026-08-18 -- THE ADVERSARIAL-REFUTATION REPAIR
% ===========================================================================
% EVERYTHING ABOVE THIS LINE IS LEFT INTACT. Three independent refuters read the
% pass-10 build and returned 33 findings. This section supersedes BY NAME only
% the claims above that this pass made false, records every finding APPLIED and
% every finding OVERRULED with its reason, and re-records the measurements.
% NO NUMBER CHANGED IN THIS PASS. 1.563, 0.084, 4.180, n = 60, rank 10, residual
% width 2048 and layer 4 are byte-identical to passes 9 and 10, the log map
% \kx is byte-identical to pass 10, and the three marks stand at the same three
% x to within the 0.02pt pgfmath residual recorded below. Nothing measured moved.
%
% TWO CLAIMS ABOVE ARE FALSE AND ARE SUPERSEDED HERE
% ---------------------------------------------------
% (S-1) The pass-10 note at the foot of this file said the replacement caption
%       "is also staged, keyed by \label, in figs/_PENDING_CAPTIONS.tex for a
%       human to integrate". THAT FILE DID NOT EXIST. It does now, this pass
%       created it, and it carries the caption plus an explicit list of the
%       clauses the integrator must DELETE from the live caption -- above all
%       Sections/04b-representation.tex:289-290, "The two arrows off the rule
%       mark that it acts in both halves and imply no count of positions",
%       which names two marks pass 10 deleted. Until that deletion is made the
%       figure ships contradicting its own caption. hook.json's .shape carried
%       the same false claim and is corrected.
% (S-2) The pass 1-9 block's AGREEMENT WITH THE TEXT cites 04b-representation
%       lines 229, 254, 260 and 286-288, and note 5/decision 5 cites 193, 195
%       and 200-202. Four of those are stale: tab:workspace's layer-4 row is now
%       line 338 (229 is fig:purify's caption), "3.6--18.6x" is 363, the
%       "9--37% ... at layers 2--9" sentence is 368-370, and "No intervals are
%       drawn in either panel" is 395. The 193/195/200-202 set is now 263, 264,
%       265 and 270-272 and pass 10 had already refreshed it in hook.json. The
%       "2048 is prose only at 04b:383" citation is also stale: 383 is
%       fig:mechanism's caption and the residual width is stated at 04b:648.
%       hook.json's citations are repaired; the stale text above is left in
%       place, superseded here, because S25 forbids rewriting a header.
%
% THE FIFTEEN REPAIRS, EACH WITH THE MEASUREMENT THAT FORCED IT
% --------------------------------------------------------------
% (1) PANEL (d), THE HEADLINE GAP WAS COVERED BY ITS OWN LABEL. This is the
%     serious one, because the log axis exists for one reason -- on a log scale
%     a horizontal gap IS a ratio -- and the gap was drawn to the pixel and then
%     papered over. The null's label ran RIGHT from its mark and occupied 100%
%     of the 5.48cm between the open null disc and the filled accent disc. It is
%     now two right-aligned lines running LEFT into the gutter, so the null's own
%     row is clear from its mark rightward and the ratio is a visible span.
% (2) PANEL (d), EVERY VALUE NOW ABUTS ITS OWN MARK. "4.180 --- cell line, the
%     positive control" put its numeral at the far end of the label from the dot
%     it names: measured in the pass-10 PDF, "4.180" began at x = 393.28, the
%     axis interval 0.32-0.43 nats, with its leading digit overlapping the "0.3"
%     tick numeral's column 12.5mm below it, while its own mark stood 48mm away.
%     On a value axis a numeral is read at the position it occupies (S21). Each
%     label is reversed to "name --- value"; no character is added.
% (3) PANEL (d), ALL THREE LABELS ARE NOW ON THE SAME SIDE OF THEIR MARKS. The
%     side flipped between rows two and three, so the disc trailed the text
%     twice and led it once, and at print a trailing 1.4pt disc reads as a
%     bullet ending a line -- three key entries, i.e. a legend, which is the one
%     construction this document forbids outright. All three labels now sit
%     left, all three discs right, and (1) is what made that possible.
% (4) PANEL (d), THE RULE NOW BEGINS AND ENDS ON A DRAWN TICK. At four labelled
%     ticks (0.1, 0.3, 1, 3) the rule overran its tick lane by 36.87pt (13.0mm)
%     at EACH end -- 29% of the axis carrying no numeral, which is the exact
%     defect S13 measures and condemns in fig:comparators row (d) -- and BOTH
%     extreme marks fell in those bare zones, so neither the null at 0.084 nor
%     the positive control at 4.180 could be placed against a printed number.
%     The domain's own ends are now ticks: 0.05, 0.1, 0.3, 1, 3, 6. The two end
%     numerals are set flush INSIDE the rule's ends rather than centred on their
%     ticks, because centred they overhung x = 17.38 by 0.09cm and made the
%     float overfull by 1.88pt; every other numeral is centred on its tick.
% (5) PANEL (d), THE TICK NUMERALS ARE INK, NOT quietink. quietink is this
%     document's comparator-and-null colour (S8/CF4). Spending it on the scale
%     while the figure's actual null was labelled in ink inverted the code.
%     fig:mechanism sets every tick numeral in ink. THE NULL'S OWN INK LABEL AND
%     ITS ink-EDGED OPEN DISC ARE DELIBERATE AND OVERRIDE CF4's quietink slot:
%     CF5 instructs this figure by name to adopt 4.11(b)'s encoding for this
%     object, and 4.11(b) draws it mfc="white", mew=1.1, color=style.INK. The
%     distinction between the null and the positive control is carried by the
%     GLYPH -- open against filled -- which is S9's own device, not by hue.
% (6) THE COMPARISON ARROW'S HEADS POINTED AT NOTHING. It ran between the two
%     brace LABELS' baselines, 9.04 to 6.88, so its upper head sat 0.50cm BELOW
%     the strip it was meant to span and both heads terminated in white space.
%     It now runs 9.75 to 7.59, the two strips' own vertical centres, at x =
%     13.10 rather than 13.40.
% (7) THE COMPARISON ARROW'S LABEL IS NOW ONE WORD, "compared". It had set
%     panel (d)'s axis quantity string a second time 130pt away, so a reader
%     could take this vertical 1.0pt rule for an axis OF that quantity; the axis
%     names the quantity once and this mark's job is to say that the two passes
%     are compared and not chained. S20 is satisfied -- the role is stated in a
%     verb of comparison beside the double head S11 gives a comparison. The 17
%     characters this frees are what pay for (8), (9) and (4) inside the 6.0%
%     ink ceiling, which the pass-10 build cleared by 2.8mm2.
% (8) THE TIE MOVED AND GAINED THE WEIGHTS. It sat at 6.88 -- the EXACT baseline
%     of pass two's brace label, measured to 0.00pt -- and directly under pass
%     two's prompt half, the one half of the strip that carries no brace, so it
%     read as that half's missing brace label, i.e. as a statement about ONE
%     strip when its whole job is to be a statement about TWO. It is now band
%     one's foot line on its own baseline at 6.44, 0.44cm below the brace label
%     and 0.76cm above band two's heads. And it read "the same sentence, box for
%     box", which carries only half of the question this figure must answer:
%     nothing anywhere on the drawing named the WEIGHTS, and the pass-10 note
%     claiming the deleted subtitle "the same prompt, the same sentence, the
%     same weights" had been inherited by the tie was wrong. It now reads "one
%     sentence, one set of weights, scored twice".
% (9) THE LAYER IS NAMED AT THE INTERVENTION. "at every position" ranges over
%     the SEQUENCE, not over depth, and the only "layer 4" on the drawing was in
%     the axis block 40mm away, attached to the READOUT. Panel (b)'s head now
%     reads "(b) pass two, the slab projected out at layer 4". DECISION 14 above
%     and pass 10's partial supersession of it are superseded again: the fact is
%     still stated in words because no mark can carry it, but it is stated where
%     the hook is drawn, and the axis block keeps its own "layer 4" because CF10
%     requires the layer named where the drawn layer's marks are.
% (10) THE HOOK RULE NO LONGER READS AS AN UNDERLINE. At y = 7.92 it sat 3.42pt
%     under an accent line of type and 3.40pt over the box tops -- equidistant,
%     same colour above and below, which is the heading-and-underline idiom. It
%     is at 7.88: 0.18cm of clear space above, 0.08cm below, so it binds to the
%     strip it acts on. It is NOT dropped onto the boxes, as one refuter
%     proposed: a mark running into a rule uninterrupted is what this file's own
%     note 5(b) forbids, and the drug cell's border is accent too.
% (11) PANEL (c), A PRINT COLLISION. The label $(hV)V^{\top} \to 0$ cleared the
%     top edge of the V_pure band by 0.72pt -- two pixels at 200dpi -- so its
%     parenthesis descenders would have closed on press. The band thins from
%     0.32 to 0.24cm and the label rises 0.02cm; the label is now 0.286cm from
%     its own mark, still inside the 3mm at which this document requires a
%     leader, and clears the band by 0.08cm.
% (12) PANEL (c), THE RIGHT ANGLE IS MIRRORED TO THE LEFT. On the right it was
%     unreadable: the 1.0pt black h vector crossed its 0.30pt horizontal leg
%     0.12cm from the corner, and its vertical leg terminated exactly on the
%     accent component's shaft. h rises to the RIGHT of the surviving vector, so
%     the corner between the survivor and the slab on the LEFT is empty. Same 90
%     degrees, same two arms, nothing crossing it. DECISION 10 above stands: the
%     mark is what distinguishes "set to ZERO" from "reduced".
% (13) PANEL (c), THE VERTICAL CONSTRUCTION GUIDE STOPS 0.20cm SHORT of h's tip.
%     It used to end inside the arrowhead. The HORIZONTAL guide still touches
%     both tips and must: it is the closing side of the parallelogram.
% (14) PANEL (c), THE BAND RUNS 0.00 TO 5.80 AND HAS LOST ITS OUTLINE. Two
%     findings, one repair. The lower band held a rectangle 38 x 46mm between
%     this panel and panel (d) with literally no ink in it -- 10% of the content
%     area and the largest void on the drawing -- and the band was byte-
%     identical in stroke and fill to the eighteen token cells above it
%     (rulegrey 0.30pt on panelbg), so ONE rectangle treatment carried two
%     meanings: "a field of the prompt", eighteen times, and "the purified drug
%     slab", once. Lengthening the band closes half the void at no ink cost
%     (panelbg renders above 230 grey and S15 does not count it); dropping the
%     outline makes the band the only fill-without-stroke region in the figure
%     and frees 8.7mm2. AN ACCENT OUTLINE WAS CONSIDERED AND REJECTED AGAIN, for
%     note 9b's reason: accent fill plus accent border is the drug cell's own
%     treatment to within a tint, and one glyph would mean two things again.
%     NOTE 9b IS THEREFORE AMENDED, NOT SUPERSEDED: the band is still panelbg
%     and still not accent; it simply no longer carries slab.tex's rulegrey
%     outline, because in THIS figure that outline is the token cell's own.
% (15) EVERY PANEL HAD TWO LEFT EDGES 1.20pt APART. TikZ's default 1pt inner
%     xsep meant a base-west node's ink began 1.20pt right of the coordinate it
%     was placed on, so the drawn rules started at one x and the panel's text at
%     another, visible under a flat-left glyph like the "K" of "KL". inner
%     xsep=0pt now. MEASURED IN THE BUILT PDF: the panel-level text lines of the
%     0.00 panels start at 57.3, 57.3, 57.5, 57.3, 57.3, 57.4, 57.3pt -- spread
%     0.2pt, and the residual is glyph side bearing, not node placement -- and
%     panel (d)'s head, its first tick numeral and both axis lines start at
%     295.5pt against an axis rule whose left end is 295.25pt.
%
% WHAT WAS PROPOSED AND OVERRULED, AND WHY
% -----------------------------------------
% (O-1) DELETE PANEL (c) AND RE-LETTER (d) AS (c). Refused. The spec's keep-list
%     protects the panel by name -- "h, h - (hV)V^T, the removed component going
%     to zero, and the deferral of the true 10/2048 angle to fig:slab" -- and
%     CF14 forbids losing load-bearing work. The finding's real content, the
%     void, is answered by (14) without deleting anything.
% (O-2) PUT ALL THREE MARKS ON ONE BASELINE AND STAGGER THE LABELS ABOVE WITH
%     LEADERS. Refused: staggering labels off one rule is exactly what S22
%     forbids, and the three names measure about 11cm against a 9cm axis, so one
%     baseline is not available. The three-row construction, one mark and one
%     label per row, is S22's own reading and it survives; (1), (2) and (3) fix
%     what was actually wrong with it.
% (O-3) DRAW DROP LINES OR PROJECTION LINES FROM THE MARKS TO THE SPINE.
%     Refused, and not for the reason pass 10 gave. S7 does sanction a densely
%     dotted 0.5pt quietink projection line, so pass 10's blanket objection was
%     too strong. The decisive objection is arithmetic: the three rows sit 0.38,
%     0.82 and 1.26cm above the spine in VALUE order, so three stems dropped to
%     the axis would have lengths that rise with the value they sit over, and
%     the panel would read as a bar chart of the three numbers against a
%     vertical scale that does not exist. A mark vocabulary must not encode a
%     quantity twice, once truly and once by accident.
% (O-4) ANNOTATE THE TWO GAPS WITH "18.6x" AND "37%" AS 4.11(b) DOES. Refused.
%     The spec deleted both printed ratios from this figure to the caption, and
%     a result is caption material here. 4.11(b) owns the annotated ratios --
%     CF1 says so -- and this figure now shows the gap and lets the caption
%     name it.
% (O-5) SET THE HOOK'S FORMULA IN ink AS DEFINITIONAL PROSE. Refused. It is the
%     direct label of an accent object, and S8 gives accent to "the drug-side
%     object under test AND ITS DIRECT LABELS". It is also \scriptsize, and S4's
%     prose is \footnotesize; calling it prose would break S4's size rule to fix
%     a colour it does not break. (10) fixes the underline reading instead.
% (O-6) DROP THE ACCENT FROM THE "drug" PROMPT CELL. Refused. It is a cross-
%     figure tie recorded in two files: decision 2 above, and figs/licensing.tex,
%     which ADOPTED this figure's cell treatment in pass 9 so that the two
%     strips read as one object seen twice. Removing it here would desynchronise
%     the pair. S8's own check lists what may not be accent -- a comparator, a
%     null, a discounted arm, a struck branch -- and the drug field is none of
%     them.
% (O-7) DRAW THE SECOND STORED RANDOM DRAW AS A FOURTH MARK. Refused as a
%     DRAWING change and ACCEPTED as a disclosure. workspace_probe.json really
%     does store layers.4.kl.random = 0.19796918605764707 beside the drawn
%     layers.4.kl.random_matched = 0.08388670384883881, and with
%     layers.4.raw_drug_dims = layers.4.pure_drug_dims = 10 the two are draws of
%     the same rank-10 construction, so the headline 18.6x would be 7.9x on the
%     other draw. That is a real and undisclosed exposure and the caption now
%     carries it. But the mark itself would be a fourth dot with a role no
%     section of the thesis states, drawn from a numeral that appears in no
%     table and in no sentence -- S20's own failure condition -- and tab:workspace
%     inherits the same single draw, so the figure would then disagree with the
%     table it is read against. hook.json records the second draw beside marks[1].
% (O-8) RENAME THE MARKS TO "ablating the purified drug slab" AND THE AXIS TO
%     "KL from the clean forward pass when a slab is ablated". Refused: S19's
%     registry and the spec both fix these strings, and line 1 is 4.11(b)'s
%     exact quantity name, which CF3 requires the slate to share. The verb the
%     finding wants is in the caption's lede, where a result belongs.
% (O-9) REVERSE OR BEHEAD THE (hV)V^T ARROW BECAUSE IT "POINTS AWAY FROM THE
%     ZERO IT GOES TO". Refused. It is a component of a vector decomposition and
%     points along the slab, which is its direction; "-> 0" is its fate, carried
%     by the label and by the right angle. Pass 9 changed it FROM headless TO
%     headed for a recorded reason -- the removed component was the only one of
%     three not drawn as a vector and read as a baseline -- and reversing the
%     head would draw a vector backwards.
%
% GEOMETRY AND STYLE, RE-MEASURED AT PASS 11
% -------------------------------------------
% Content box 173.99 x 99.95mm (fullwidth measure 170-174mm, height cap 100mm),
% unchanged from pass 10 to the hundredth of a millimetre. Ink 5.970% of the
% content box against a 6.0% ceiling, down from pass 10's 5.984% and pass 9's
% 8.55%. 709 non-space characters over 17390mm2 is 4.08 per 100mm2 against a 5.0
% ceiling. (The 793 recorded at pass 10 was counted by a different method; 709
% is the count of non-space characters in the built PDF's own text spans, and
% that method is now the one hook.json records.)
% TYPE, measured in the built PDF: 7.97pt x522 and 9.96pt x146 -- two chosen
% sizes -- plus the exempt classes 6.85 x20 (\ENDCELL, Inconsolata), 5.98 x8 and
% 6.23 x3 (sub- and superscripts) and 8.30 x10 (the CMSY10/CMR10 math symbols
% pass 10 explains and mathpazo causes). Faces: Pagella-Regular 653,
% PalladioL-Roma 9, Inconsolata 20, PalladioL-Ital 14, CMSY10 7, CMR10 6. Zero
% sans, zero bold, zero all-caps banners. Italic 14 of 709 = 2.0%, and all
% fourteen are the math variables h and V.
% STROKES, measured: 0.30 x27, 0.40 x5, 0.80 x1, 0.90 x1, 1.00 x4 -- the same
% five as pass 10. RECTANGLES: nineteen, and pass 10's count of them was wrong
% in its own arithmetic. Sixteen panelbg token cells, TWO accent!12 drug cells,
% eighteen cells in all, plus the V_pure band -- which is now the only one drawn
% without a stroke. ARROWHEADS: five, three on panel (c)'s vectors and two on
% the comparison arrow.
% THE AXIS, measured: spine x = 295.25 to 549.81pt, ticks at 295.25 (0.05),
% 332.10, 390.48, 454.51, 512.94 and 549.81 (6), i.e. the first and last drawn
% tick sit exactly on the rule's endpoints and there is no unnumbered overrun.
% Marks at 323.06 (null, between the 0.05 and 0.1 ticks), 478.26 (purified slab,
% between 1 and 3) and 530.58 (cell line, between 3 and 6). Each label ends
% 7.56-7.58pt from its own mark. The null-to-slab gap is 155.20pt = 5.474cm and
% carries no text on the null's own row.
%
% ===========================================================================
% PASS 12 -- 2026-08-19 -- THE SECOND ADVERSARIAL-REFUTATION REPAIR
% ===========================================================================
% EVERYTHING ABOVE THIS LINE IS LEFT INTACT. Three further refuters read the
% pass-11 build and returned 30 findings, two of them marked blocking and both
% blocking findings naming the same object. This section supersedes BY NAME the
% claims above that this pass made false, records every finding APPLIED and
% every finding OVERRULED with its reason, and re-measures the drawing.
%
% NO NUMBER CHANGED IN THIS PASS EITHER. 1.563, 0.084, 4.180, n = 60, rank 10,
% residual width 2048 and layer 4 are byte-identical to passes 9, 10 and 11.
% What changed is the value-to-POSITION map (\kx's domain, recorded below and
% in hook.json beside the two maps it supersedes) and where labels and marks
% sit. No mark's VALUE moved; every mark's x moved because the scale did.
%
% CLAIMS ABOVE THAT ARE NOW FALSE, SUPERSEDED HERE BY NAME
% ---------------------------------------------------------
% (T-1) PASS 11's O-2 ("put all three marks on one baseline" -- refused) and
%       O-3 ("drop lines from the marks" -- refused). Both are superseded; see
%       repair (1). O-2's arithmetic was about the LABELS, which still stagger;
%       it wrongly carried that conclusion over to the MARKS, which are 4.52cm
%       and 1.52cm apart and cannot collide. O-3's bar-chart objection is
%       answered by de-correlating leader length from value.
% (T-2) PASS 11's repair (2), "puts every numeral within a few per cent of its
%       own axis position". Measured, it did not: it moved each numeral one
%       decade along and left "4.180" sharing a column with the "3" tick, "1.563"
%       with the "1" tick and "0.084" with the "0.05" tick. Repair (2) below
%       replaces the approximation with an exact construction.
% (T-3) PASS 11's repair (3), "all three labels now sit left of their marks ...
%       a trailing 1.4pt disc reads as a bullet ending a line". The note
%       diagnosed the legend hazard correctly and then made ALL THREE rows
%       trail into a disc, which is the hazard. Superseded by repair (1).
% (T-4) PASS 11's repair (4), which set the domain-end tick numerals at
%       0.05 and 6. Superseded by repair (3): the ladder is now 0.03 / 0.1 /
%       0.3 / 1 / 3 / 10. The REASON recorded in repair (4) -- that the rule
%       must begin and end on a drawn tick and that every mark must have a
%       printed numeral either side of it -- stands and is preserved.
% (T-5) PASS 11's repair (13), which stopped the vertical construction guide
%       0.20cm short of h's tip "at the top end only". Its foot is repaired in
%       this pass; see (7).
% (T-6) PASS 11's repair (12), which mirrored panel (c)'s right angle to the
%       LEFT of the surviving vector. Superseded by (6): the origin cannot
%       carry that mark at all, whichever way it is mirrored.
% (T-7) LAYOUT NOTE (j) in the pass 1-9 block, "THE PASS-ONE BOUNDARY RULE GREW
%       A TAIL, from 8.72 down to 8.55", is superseded by (9): the label that
%       tail was grown to carry was deleted in pass 10, and the two dividers
%       are now identical.
% (T-8) PASS 11's claim at the head of its S-2 that "hook.json's citations are
%       repaired". They were not. Twelve of eighteen line citations in the
%       ledger did not resolve against the current Sections/04b-representation
%       .tex, and five ledger fields still described the pass-10 drawing. All
%       are repaired in this pass and hook.json now carries a pass-12
%       measurement block that supersedes the pass-10 style block by name.
%
% THE ELEVEN REPAIRS, EACH WITH THE MEASUREMENT THAT FORCED IT
% -------------------------------------------------------------
% (1) PANEL (d) IS REBUILT: THE THREE MARKS NOW SHARE ONE BASELINE AND EACH
%     NAME IS LEADERED TO ITS OWN MARK. This is the pass's one structural
%     change and it answers eight findings at once, including both blocking
%     ones. Measured in the pass-11 PDF: the three dot centres stood at
%     (323.06, 268.30), (478.26, 255.83) and (530.58, 243.36)pt -- three
%     baselines 0.44cm apart -- so the 155.20pt "gap" the log axis exists to
%     show was not a distance any reader could lay a ruler along, and the two
%     labels of the rows above ran x 363.68-470.69 and 394.28-523.00, i.e.
%     wholly inside it. Each row was also a text line ENDING in its own disc,
%     so at print three 1.4pt marks read as three bullets and the block read as
%     a key. And the null's two label lines began at x = 245.04 and 241.61
%     against panel (d)'s origin at 295.45 -- 19.0mm outside their own panel,
%     in panel (c)'s column, with "random slab" set 71pt above a band labelled
%     "the purified drug slab" in the same size and face.
%     Now: marks at y = 4.02 at their own three x, nothing else on that line;
%     names on rows 5.26 / 4.86 / 4.46; a 0.30pt leader in each mark's own
%     colour from just under its name to just above its disc. The row order is
%     null / cell line / purified slab, chosen for two reasons at once -- it is
%     the only order in which no leader crosses another row's label, and it
%     makes leader length (1.07 / 0.67 / 0.27cm) non-monotone in value (0.084 /
%     4.180 / 1.563), which is what kills O-3's bar-chart reading. The null's
%     name runs RIGHT from its mark because its mark is 1.59cm from the axis's
%     left end and the name measures 5.17cm; every text line in the panel now
%     lies inside 8.40-17.38.
% (2) EVERY PRINTED VALUE NOW STANDS EXACTLY AT ITS OWN AXIS POSITION. Each
%     label is aligned so its numeral's INNER edge sits on its mark's x:
%     "0.084" begins at \kx{0.084} = 9.9916, "4.180" ends at \kx{4.180} =
%     16.0316, "1.563" ends at \kx{1.563} = 14.5110. That is why one row is
%     value-first and two are value-last -- the value always abuts the mark and
%     the side is set by the measure, not by taste. It also lifts every printed
%     value 1.36cm or more clear of the spine, so the rule no longer carries
%     numerals within a 15pt band on both sides (S21).
% (3) THE TICK LADDER IS UNIFORM. 0.05 / 0.1 / 0.3 / 1 / 3 / 6 stepped x2, x3,
%     x3.33, x3, x2 and printed at 1.300 / 2.061 / 2.258 / 2.061 / 1.300cm -- a
%     74% spread, crowding a pair of numerals at each end of the rule, with the
%     null inside the crowded left pair. 0.03 / 0.1 / 0.3 / 1 / 3 / 10 is the
%     textbook 1-3 ladder and prints at 1.861 / 1.698 / 1.861 / 1.698 / 1.861,
%     a 10% spread. The rule still begins and ends on a drawn tick, no mark
%     sits on a tick, and every mark keeps a printed numeral on each side.
% (4) THE HOOK RULE NO LONGER HEADS A FLUSH-LEFT STACK. Title (9.96pt ink,
%     x0 57.34pt), formula (7.97pt accent, x0 57.48), accent 0.90pt rule (x0
%     57.14), boxes -- four decks on one left edge with 4.5pt of clear space
%     above the rule and 2.8pt below is the heading-and-underline idiom, and
%     the one mark whose whole job is "this operation runs the length of the
%     stream" was reading as "panel (b) starts here". The formula is now
%     right-aligned to the rule's own right end. It is a DIRECT LABEL of a mark
%     (S22), not a panel-level element (S13), and a rule's direct label belongs
%     at one of its ends.
% (5) "compared" IS ROTATED 90 DEGREES, READING UPWARD, BESIDE THE ARROW. Set
%     horizontally its baseline was 8.62 at \scriptsize while panel (b)'s head
%     was 8.54 at \footnotesize, 13.3cm away: two lines 2.27pt apart in
%     baseline, at two sizes, neither aligned nor separated. Rotated there is
%     no second baseline to near-miss and the label sits on the mark it names.
% (6) PANEL (c)'S RIGHT ANGLE MOVES TO THE FOOT OF THE PERPENDICULAR. See the
%     note at the mark itself. The short form: h = survivor + component with
%     both terms positive, so h always lies inside the quadrant the ORIGIN's
%     right angle would enclose, and no mirroring escapes that; the
%     decomposition is a rectangle, so the foot of the perpendicular carries
%     the same 90 degrees and is the one corner that is empty and whose two
%     arms both end on drawn strokes.
% (7) THE VERTICAL CONSTRUCTION GUIDE'S FOOT STOPS 0.14cm SHORT of the accent
%     component's arrow tip, the same treatment pass 11 gave its top end.
% (8) THE SURVIVING VECTOR IS LABELLED AT MID-SHAFT. At the tip, its baseline
%     (5.16) and h's (5.20) sat either side of the horizontal closing guide
%     (5.24), so the row read as "h - (hV)V^T - - - - - h", a dashed RELATION
%     between two expressions -- an equality or comparison the panel does not
%     assert. The guide cannot move; it is the closing side of the
%     parallelogram. One of the two labels had to leave its line, and the
%     survivor's is the one that has a shaft to sit beside.
% (9) THE TWO PROMPT/TARGET DIVIDERS ARE NOW IDENTICAL. Measured: pass one
%     0.686cm, symmetric; pass two 0.572cm, asymmetric. In a figure whose
%     single claim is that the two strips are the same object, the one
%     structural mark that visibly differed between them was a vestige.
% (10) \ENDCELL'S CELL IS 1.47cm, NOT 1.70. Its Inconsolata ink measures
%     1.204cm and every other cell in the strip carries 0.131-0.150cm of clear
%     space a side; this one carried 0.248. The strip end, the braces, the hook
%     rule and the comparison arrow are keyed to \stripend and moved with it.
% (11) THE INTERVAL DECLARATION IS SET ON TWO LINES. At 15.89cm it was the
%     widest object in the figure -- wider than the axis and wider than band
%     one -- which S14 forbids, and an 8pt grey line that wide, stranded a
%     centimetre below everything else, reads as a second and smaller caption
%     above the real one. Broken at the colon it measures 8.68cm and 7.10cm.
%
% WHAT WAS PROPOSED AND OVERRULED, AND WHY
% -----------------------------------------
% (P-1) DELETE PANEL (c) AND RE-LETTER (d) AS (c), giving the readout the full
%     17.38cm. Refused again, for O-1's reason: the spec's keep-list protects
%     the panel by name and CF14 forbids losing load-bearing work. The
%     finding's substantive point -- that the panel overstates the deletion --
%     is conceded in the caption, which also defers the true 10/2048 geometry
%     to fig:slab, where it is drawn at 86 degrees.
% (P-2) WIDEN THE IN-SLAB COMPONENT FROM 0.90 TO 1.35cm so the two black
%     vectors separate at the origin; and, from a different refuter, NARROW IT
%     TO 0.30cm so the panel stops overstating the deletion. Both refused,
%     because they are contradictory demands on one parameter: the angle
%     between h and the survivor IS the drawn ratio, so every millimetre of
%     separation at the origin is bought by overstating what the hook removes.
%     Pass 9 chose that trade-off deliberately at 4.0:1 and recorded it; the
%     built ratio is 3.49:1 and the fusion is confined to the bottom 1.2mm of a
%     31.4mm vector. The relocated right angle (6) also moves the corner a
%     reader has to read away from the fused base.
% (P-3) GIVE THE V_pure BAND AN ACCENT OUTLINE, so the object under test is not
%     drawn in the furniture colour while its own label is accent. The
%     COMPLAINT IS CONCEDED and the repair is refused on arithmetic. The
%     drawing measures 5.972% non-white at 200dpi inside a 173.99 x 99.95mm
%     content box against S15's 6.0% ceiling: 4.9mm2 of headroom. A 0.30pt
%     outline round the 5.80 x 0.24cm band costs 12.8mm2, and even a band
%     shortened to the extent its construction occupies costs more than the
%     ceiling allows -- the longest band an outline could be bought for is
%     2.07cm, which is shorter than the construction standing on it. Pass 9's
%     note 9b and pass 11's repair (14) also both refused an accent border, on
%     the ground that accent-on-panelbg is the drug cell's treatment to within
%     a tint. The band keeps its accent DIRECT LABEL, the accent component is
%     drawn lying inside it, and the same object is an accent filled disc in
%     panel (d), which is where the colour code is carried. RECORDED AS AN
%     OWED REPAIR for any future pass that finds ink to spend.
% (P-4) REBUILD \tokstrip AS TWO OUTER RECTANGLES PLUS INTERNAL VERTICALS, to
%     close the ~0.4mm notches the abutting cells' rounded corners leave in the
%     strip's top and bottom edges. Refused, and it is the closest call in this
%     pass. The notches are real and measurable (13 missing pixels across four
%     boundaries at 200dpi), and the repair would also SAVE 6.2mm2. But the
%     nine-rounded-cells glyph is a recorded cross-figure tie: decision 2 above
%     records that figs/licensing.tex ADOPTED this figure's cell treatment in
%     pass 9 precisely so a reader who meets the prompt strip twice meets one
%     object, and licensing.tex is out of this task's scope, so the repair
%     could not be mirrored and the two figures would diverge again. RECORDED
%     AS AN OWED REPAIR to be made in hook.tex and licensing.tex together.
% (P-5) CENTRE ALL SIX TICK NUMERALS AND CUT THE RULE TO 17.10 TO PAY FOR IT.
%     Refused; the reason is at the ladder itself. In short: at the LEFT end
%     centring puts "0.03" 0.27cm outside panel (d)'s x origin, and S13's
%     one-grid rule, measured at a 0.2pt spread, outranks the optical centring
%     of a domain-end numeral that labels no mark.
% (P-6) EXTEND THE HOOK RULE 0.25cm PAST BOTH STRIP ENDS. Refused: the
%     picture's bounding path is pinned at x = 0.00 and a rule at -0.25 would
%     put the content box at 176.3mm, past S14's 174mm ceiling. (4) breaks the
%     heading reading without moving the rule.
% (P-7) MOVE THE COMPARISON ARROW OUT TO x ~ 15.2 WITH LEADERS RUNNING BACK TO
%     THE STRIPS, to close the empty top-right quadrant. Refused. It would
%     detach the figure's comparison mark from the two objects it compares,
%     buy the composition with two 2.1cm horizontal leaders (4.5mm2, more than
%     the whole ink headroom) and re-open exactly the flow reading -- one strip
%     feeding a mark feeding the other -- that the double head exists to kill.
%     The empty quadrant is accepted and recorded.
% (P-8) RESTORE "generation would begin here" AT THE PROMPT/TARGET DIVIDER, or
%     delete the divider as an unlabelled duplicate of the brace's boundary.
%     Refused both ways. The spec's keep-list protects the divider by name; a
%     label there would restate what the 1.00cm gap already shows and cost ink
%     the ceiling does not have; and the divider is named where an inherited
%     phrase belongs, in the caption, which now reads "the divider marks only
%     where a generated continuation would have begun".
% (P-9) MIRROR PANEL (c)'S CONSTRUCTION so the removed component points LEFT
%     and h tips right. Refused: that is not a mirror, it is an error. With h
%     tipping right and the component pointing left, h is no longer the sum of
%     the two drawn vectors. The finding's DIAGNOSIS -- that the pass-11 mark
%     did not enclose the angle it claimed -- is correct and is applied as (6).
% (P-10) SET THE RIGHT ANGLE IN ink 0.40pt. Refused: the two construction
%     guides beside it are rulegrey 0.30pt, the mark is the third element of
%     that one apparatus, and splitting it off would give 0.40pt a sixth
%     meaning inside S10's five-weight table.
%
% GEOMETRY AND STYLE, RE-MEASURED AT PASS 12
% -------------------------------------------
% Content box 173.99 x 99.95mm -- fullwidth measure 170-174mm, height cap
% 100mm -- unchanged to the hundredth of a millimetre from passes 10 and 11.
% Ink 5.972% at 200dpi against a 6.0% ceiling. 710 non-space characters in the
% built PDF's own text spans over 17390mm2 is 4.08 per 100mm2 against 5.0.
% TYPE, measured in the built PDF: 7.97pt x523 and 9.96pt x146 -- two chosen
% sizes -- plus the exempt classes 6.85 x20 (\ENDCELL, Inconsolata), 5.98 x8 and
% 6.23 x3 (sub- and superscripts) and 8.30 x10 (CMSY10/CMR10 math symbols).
% Faces: Pagella-Regular 654, PalladioL-Roma 9, Inconsolata 20, PalladioL-Ital
% 14, CMSY10 7, CMR10 6. Zero sans, zero bold, zero all-caps banners. Italic
% 14 of 710 = 2.0%, and all fourteen are the math variables h and V.
% STROKES: 0.30 (furniture -- 18 cell borders' worth of rulegrey, both braces,
% six axis ticks, two construction guides, the right angle, three label
% leaders), 0.40 (two prompt/target dividers, the axis rule), 0.80 (ONE, the
% null's open-disc edge: CF5 instructs this figure by name to take 4.11(b)'s
% encoding for this object and mechanism.py draws it mew=1.1, so the singleton
% is a mandated distinct kind of object and not a stray weight), 0.90 (ONE, the
% hook rule: it is the figure's intervention and the single heaviest non-data
% stroke; at 1.00pt it would merge with the vectors and the comparison arrow,
% which are data paths, and at 0.40pt with the axis), 1.00 (three vectors and
% the comparison arrow). Five weights, both singletons justified.
% RECTANGLES: nineteen -- sixteen panelbg token cells, two accent!12 drug
% cells, and the V_pure band, which is still the only one drawn without a
% stroke. ARROWHEADS: five, three on panel (c)'s vectors and two on the
% comparison arrow; no other mark in the figure carries one.
% THE AXIS, measured in the built PDF: spine 295.25 to 549.81pt; ticks at
% 295.25 (0.03), 348.07 (0.1), 396.15 (0.3), 448.91 (1), 497.05 (3) and 549.81
% (10), i.e. tick spacings 52.82 / 48.08 / 52.76 / 48.14 / 52.76pt, a 10%
% spread against pass 11's 74%, and the first and last drawn tick still sit
% exactly on the rule's endpoints. The four interior numerals are centred on
% their ticks; the two domain-end numerals are set flush inside the rule's ends
% for the reason recorded at the ladder. THE THREE MARKS ARE AT 340.61 (null,
% white-filled open disc, 0.797pt edge), 468.48 (purified slab, accent filled)
% and 511.58 (cell line, ink filled) -- ALL THREE AT cy = 262.63, one baseline
% to the hundredth of a point. The null-to-slab gap is 127.87pt = 4.510cm and
% carries no ink at all; the gap the map predicts for a ratio of
% 18.634044258027835 is 4.521cm, so the drawn span IS the ratio to within
% 0.011cm, which is a third of the width of the stroke that draws the discs.
% The slab-to-cell-line gap is 43.10pt = 1.520cm against a predicted 1.5202cm.
% THE THREE LEADERS: 0.299pt, at x = 340.61 (ink, 30.33pt), 511.58 (ink,
% 19.00pt) and 468.48 (accent, 7.65pt) -- lengths 1.07 / 0.67 / 0.27cm against
% values 0.084 / 4.180 / 1.563, deliberately not monotone.
% ===========================================================================
%
% ===========================================================================
% v-BLOCK, 2026-08-19: WHOLE-SLATE PASS (fig:comparators, fig:licensing,
% fig:purify, fig:hook, fig:materiality reconciled against one another and
% against fig:mechanism, which is the document's reference standard and was
% NOT rebuilt). Nothing above this line is deleted; where an earlier claim is
% now false it is superseded EXPLICITLY below.
%
% WHY. The five figures were rebuilt independently and drifted. This pass owns
% only the differences BETWEEN them. No value changed in any of the five; the
% sibling .json ledgers carry the value-to-position maps, old and new.
%
% THE SLATE RULES THAT NOW HOLD ACROSS ALL FIVE (each stated once here, in
% every file, so any future single-figure edit can see what it would break):
%
%  R1 TICK NUMERALS are quietink, \scriptsize, inner sep 0, anchored north on
%     the tick and set 0.16cm below the axis rule, over a 0.10cm rulegrey
%     0.30pt tick. Before: comparators/licensing/purify quietink, hook and
%     materiality ink; ticks 0.10 / 0.12 / 0.14cm; three different gaps.
%  R2 A THRESHOLD'S OWN LABEL IS BLACK, matching the black solid 0.90pt rule
%     it names (S6). Before: comparators black, purify's "gate, 0.8" and
%     materiality's two margin labels ink.
%  R3 ONE NAME PER OBJECT (S19). The pre-registered 10% rule is "10%
%     materiality margin" in fig:comparators row (c), fig:materiality (a) and
%     (b), and in figs/family.py, which is where the name comes from.
%     fig:comparators' "pre-registered margin" is superseded; "pre-registered"
%     is now the caption's word, not the drawing's. Zero, where it is the null
%     a bar is read against, is "zero: no identity effect" in BOTH
%     fig:comparators row (c) and fig:materiality (b) -- one quantity, one
%     wording. V_pure is spelt with \textrm everywhere, as figs/slab.tex does.
%  R4 COLOUR KEY, one key for five figures (S8). accent = the drug-side object
%     under test. quietink = every comparator, null, control, replication or
%     discounted arm, AND its label, AND its leader. ink/black = thresholds,
%     axis rules, axis labels, panel heads, definitional prose. rulegrey =
%     furniture only. Within quietink the GLYPH separates two kinds: an OPEN
%     disc is a construction built to carry no signal (a null, a control, a
%     raw/before series); a FILLED disc is a real measurement that is not the
%     headline. fig:hook (d) drew both its null and its cell-line control in
%     ink and is corrected.
%  R5 PROJECTION AND ZERO RULES are quietink 0.50pt densely dotted (S7).
%     fig:hook's two panel-(c) projection lines were rulegrey 0.30pt densely
%     dashed and are corrected.
%  R6 DEFINITIONAL PROSE is \footnotesize ink and OUTRANKS every label (S4);
%     at most one block per figure, on the panel's own x origin. fig:hook's
%     one prose line was \scriptsize and is raised. fig:purify's panel-(b)
%     block carried an argument clause ("so causal claims use the purified
%     slab, not the raw") that its caption already carries, and is now
%     definitional only.
%  R7 ONE INTERVAL DECLARATION per figure, \scriptsize quietink, at the foot,
%     on the figure's x origin, opening with the word "Intervals:" and naming
%     EVERY panel including the ones with no measured mark (S17).
%     fig:comparators already did this and is the model; the other four now
%     match its form.
%  R8 THE FIVE-FIELD PROMPT STRIP is ONE object drawn in two figures, and the
%     two drawings are now congruent: cell boundaries 0 / 1.12 / 1.90 / 2.64 /
%     4.23 / 6.90 cm and cell height 0.42cm in both fig:licensing (a) and
%     fig:hook (a) and (b), labels centred in their cells, the drug cell
%     accent-bordered at 0.40pt on accent!12 and its label NOT bold. Before:
%     licensing 6.90cm wide with 0.62cm cells and a bold label, hook 7.44cm
%     wide with 0.42cm cells and a plain label -- the same five fields at two
%     sizes, four pages apart, with each caption pointing at the other.
%  R9 MEASURE. Every fullwidth figure draws 172.0-174.0mm; the one textwidth
%     figure draws 140.3mm. Ink at 200dpi inside the content bbox, counted as
%     greyscale luminance below 230 (the contract's own method), is now
%     5.63-5.99% on all five, under the 6.0% body ceiling, where before this
%     pass it ran 5.36-6.54%.
%
% WHAT THIS PASS DELIBERATELY DID NOT UNIFY, so that a later reader does not
% think it was missed:
%  * PANEL-HEAD CASE. fig:comparators sets "(a) Is it there?" and the other
%    four set "(a) where the drug name enters". comparators' heads are
%    interrogative SENTENCES and sentence case is correct for a sentence; the
%    v4 header argued this and the argument stands. fig:mechanism's lowercase
%    heads are noun phrases, not questions. One capital letter, four times.
%  * PANEL-HEAD POSITION. fig:comparators puts its heads in a right-aligned
%    left stub, the other four above the panel at its x origin. Moving them
%    above the rules costs about 20mm of height against a 118mm cap that the
%    figure already meets at 117.60mm. Not available.
%  * TWO-COLUMN GUTTER. purify and hook split at x = 8.40cm, licensing at
%    7.75cm. licensing cannot move right: "(d) what a negative would buy" is
%    4.85cm wide and its origin is already at 12.40 of a 17.30cm measure, with
%    1.5pt of slack. Moving purify and hook LEFT to 7.75 would be arbitrary.
% ===========================================================================
%
% CHANGED IN THIS FILE. (i) The prompt strip adopts fig:licensing's cell
% boundaries, 0 / 1.12 / 1.90 / 2.64 / 4.23 / 6.90cm, so the one object is one
% width in both figures (R8). The divider \bnd follows it from 7.94 to 7.40,
% the target strip from 8.44-12.42 to \tgt = 7.90 - \stripend = 11.88 with its
% four cell widths unchanged, the two braces and their "scored tokens" labels
% with it, and the comparison arrow from x 12.87 to 12.33 with its rotated
% label from 13.27 to 12.73. Panel (d), the axis and the foot do not move.
% (ii) Panel (d): the matched-dimension random slab and the cell-line positive
% control are comparators, so their marks, labels and leaders go ink ->
% quietink (R4). The null keeps the open disc and the control the filled one,
% which is the slate's glyph distinction between a construction built to carry
% no signal and a real measurement that is not the headline. accent now
% appears exactly once in panel (d), on the object under test.
% (iii) The six x-tick numerals go ink -> quietink, anchor=base -> anchor=north
% at a 0.06cm gap (R1); the end pair keep their north-west / north-east pull.
% (iv) The two panel-(c) projection lines go from rulegrey 0.30pt densely
% dashed to quietink 0.50pt densely dotted, which is the register
% fig:comparators, fig:purify and fig:materiality all use for a projection and
% which S7 reserves for it (R5).
% (v) "one sentence, one set of weights, scored twice" goes \scriptsize ->
% \footnotesize ink, so this figure's one prose block outranks its labels as
% fig:purify's and fig:licensing's do (R6).
% (vi) The hook formula's tail, "prompt and target alike", is deleted; the
% accent rule visibly spans the prompt half and runs across the divider, and
% the caption states it in full. This is the one deletion made for the ink
% ceiling rather than for a defect, and it is the first thing to restore.
% (vii) V_pure is respelt with \textrm, as figs/slab.tex and figs/purify.tex
% spell it. (viii) The foot block takes the slate's declaration form (R7).
% STILL TRUE AND STILL DECLARED: the 0.90pt accent hook rule is a singleton
% weight, kept under S10's own permission; it is accent, not black, so it
% cannot be read into the threshold register, and this figure draws no
% threshold. The ten 8.30pt glyphs are CMR10/CMSY10 math signs inside the two
% formulas, scaled 1.041 by mathpazo; they are not a third chosen type size
% and cannot be removed without deleting the definition of the hook.
% RE-MEASURED: 173.99 x 99.95mm; ink 5.955%; 806 characters, 4.63 per 100mm^2.
% One page, zero Overfull, zero Underfull.
% ===========================================================================
\begin{tikzpicture}[x=1cm, y=1cm, font=\small,
                    every node/.style={inner xsep=0pt, inner ysep=1pt}]
  % TWO SIZES, ONE FACE. \footnotesize (10.0pt) is reserved for the four panel
  % heads and the axis's quantity line; \scriptsize (8.0pt) carries everything
  % else. No \sffamily, no \bfseries, no \itshape anywhere in the picture.
  % inner xsep=0pt, PASS 11: with TikZ's 1pt x padding a base-west node's ink
  % began 1.20pt right of the coordinate it was placed on, so in every panel the
  % drawn rules started at one x and the text at another (S5, S13). The padding
  % is now zero in x and the text sits ON the panel's origin.
  \tikzset{
    ttl/.style  ={font=\footnotesize, text=ink,      anchor=base west},
    lab/.style  ={font=\scriptsize,   text=ink,      anchor=base west},
    labe/.style ={font=\scriptsize,   text=ink,      anchor=base east},
    qlab/.style ={font=\scriptsize,   text=quietink, anchor=base west},
    qlabe/.style={font=\scriptsize,   text=quietink, anchor=base east},
    tick/.style ={font=\scriptsize,   text=quietink, anchor=north,
                  inner sep=0pt},
    def/.style  ={font=\footnotesize, text=ink,      anchor=base west},
    alab/.style ={font=\scriptsize,   text=accent,   anchor=base west},
    brc/.style  ={draw=rulegrey, line width=0.30pt, decorate,
                  decoration={brace, amplitude=3.5pt, mirror}},
    arw/.style  ={-{Latex[length=3.4pt,width=2.8pt,line width=0pt]}},
    cmp/.style  ={{Latex[length=3.4pt,width=2.8pt,line width=0pt]}-%
                  {Latex[length=3.4pt,width=2.8pt,line width=0pt]}}}

  % --- one token box, and the drug box, so the two strips cannot differ ------
  % Border 0.30pt: a cell is a container, and containers are furniture here.
  % The drug cell keeps its 0.40pt accent border, because it is the one field
  % of the five that is an object in the argument -- the place the prompt names
  % the drug in plain text -- and it is the only accent inside either strip.
  \newcommand{\tokbox}[4]{%
    \draw[rulegrey, line width=0.30pt, fill=panelbg, rounded corners=1pt]
      (#1,#3) rectangle (#2,{#3+0.42});
    \node[font=\scriptsize, text=ink, inner sep=0pt]
      at ({(#1+#2)/2},{#3+0.20}) {#4};}
  \newcommand{\tokdrug}[3]{%
    \draw[accent, line width=0.40pt, fill=accent!12, rounded corners=1pt]
      (#1,#3) rectangle (#2,{#3+0.42});
    \node[font=\scriptsize, text=accent, inner sep=0pt]
      at ({(#1+#2)/2},{#3+0.20}) {drug};}
  % THE strip. One argument: the y of the box bottoms. Called twice, and that
  % is the whole guarantee that pass one and pass two are the same sentence.
  \newcommand{\tokstrip}[1]{%
    \tokbox{0.00}{1.12}{#1}{cell line}%
    \tokdrug{1.12}{1.90}{#1}%
    \tokbox{1.90}{2.64}{#1}{dose}%
    \tokbox{2.64}{4.23}{#1}{mechanism}%
    \tokbox{4.23}{6.90}{#1}{control cell sentence}%
    \tokbox{\tgt}{8.88}{#1}{gene$_1$}%
    \tokbox{8.88}{9.86}{#1}{gene$_2$}%
    \tokbox{9.86}{10.41}{#1}{\dots}%
    \tokbox{10.41}{\stripend}{#1}{\ENDCELL}}
  % the prompt/target boundary: mid-gap, where nothing else is drawn.
  \def\bnd{7.40}
  \def\tgt{7.90}
  % PASS 12: 12.42, not 12.65. \ENDCELL's cell was the one field of nine whose
  % clear space did not match the other eight: measured ink 1.204cm (Inconsolata
  % at \scriptsize, built PDF x 374.49-408.76pt) inside a 1.70cm cell is 0.248cm
  % a side against 0.131-0.150cm everywhere else, so at print it was the one
  % roomy cell in a strip whose whole claim is regularity. 1.204 + 2 x 0.133 =
  % 1.470. The strip end is a \def and the brace, the hook rule and the
  % comparison arrow are all keyed to it, so they move with it.
  \def\stripend{11.88}

% ###########################################################################
% # BAND ONE -- the two passes.  x origin 0.00
% ###########################################################################

  % ---------------- (a) pass one, clean -------------------------------------
  \node[ttl] at (0.00,10.34) {(a) pass one, clean};

  \tokstrip{9.54}

  % The divider is rulegrey 0.40pt SOLID, not ink 0.5pt densely dotted: S7
  % reserves the dotted stroke for a zero rule, a leader or a projection line
  % and forbids it a label of its own, and this one is labelled. It is a piece
  % of the strip's own structure, so it takes a strip weight, and it sits at
  % mid-gap because a rulegrey line on a rulegrey box border is invisible.
  % PASS 12: 9.40 to 9.98, not 9.40 to 10.10. The two dividers were drawn at two
  % different lengths and two different symmetries -- pass one 0.686cm with
  % 0.133cm clear above and below its boxes, pass two 0.572cm with 0.134 below
  % and 0.018 above (clipped to clear the hook rule at 7.88). The one structural
  % mark that visibly DIFFERED between two strips whose whole claim is that they
  % are identical was a leftover: layout note (j) below grew that tail to carry
  % the label "generation would begin here", which pass 10 deleted. NOTE (j) IS
  % THEREFORE SUPERSEDED -- its reason no longer exists. Both dividers now run
  % 0.58cm with the same 0.14/0.02 offsets from their own boxes.
  \draw[rulegrey, line width=0.40pt] (\bnd,9.40) -- (\bnd,9.98);

  % The brace spans the TARGET side only and starts at the first target box's
  % own left edge, 9.20. Both braces carry the SAME label: the two spans are
  % the same tokens read twice, and two different labels sent a reader looking
  % for a difference between them.
  \draw[brc] (\tgt,9.44) -- (\stripend,9.44);
  \node[font=\scriptsize, text=ink, anchor=base] at (9.89,9.04)
    {scored tokens};

  % ---------------- (b) pass two, the slab projected out --------------------
  % PASS 11: "at layer 4" is added to the head. The word "every" in "at every
  % position" ranges over the SEQUENCE, not over depth, and the only place the
  % depth was named was the axis block 40mm away, attached to the READOUT rather
  % than to the intervention -- so a reader could take the hook to run at every
  % layer as well as at every position, which is exactly the misreading decision
  % 14 above says the figure exists to prevent. The axis block keeps its own
  % "layer 4" because CF10 requires the layer named where the one drawn layer's
  % marks are.
  \node[ttl] at (0.00,8.54) {(b) pass two, the slab projected out at layer 4};

  % THE HOOK. The RULE is what says "everywhere along here": it spans the whole
  % strip, prompt half and target half, and runs straight across the
  % prompt/target boundary, which is the point. Its formula sits ON the rule
  % rather than 1.1cm under the strip, so the label names the mark it is beside.
  % The two short arrows of pass 9 are gone: they implied a count of two and
  % the caption had to deny it.
  % PASS 11: the rule drops from 7.92 to 7.88, so its leading is ASYMMETRIC --
  % 0.18cm of clear space above it to the formula's descenders, 0.08cm below it
  % to the box tops. At 7.92 it was equidistant between an accent line of type
  % above and the boxes below, and accent type with an accent rule tight beneath
  % it is the heading-and-underline idiom; the rule now binds to the strip it
  % acts on. It is NOT dropped onto the boxes: a mark running into a rule
  % uninterrupted is what the house rule forbids, and the drug cell's border is
  % accent too.
  % PASS 12: THE FORMULA IS RIGHT-ALIGNED TO THE RULE'S OWN RIGHT END. Flush
  % left it made a three-deck flush-left stack -- title (9.96pt ink, x0 57.34pt),
  % formula (7.97pt accent, x0 57.48pt), accent 0.90pt rule (x0 57.14pt), boxes
  % -- with 4.5pt of clear space above the rule and 2.8pt below, and an accent
  % rule tight under an accent line of type sharing its left edge is the
  % heading-and-underline idiom. The mark that has to say "this operation runs
  % along the stream" was reading as "panel (b) starts here". The formula is a
  % DIRECT LABEL of the rule, not a panel-level element, so S22 governs it and
  % not S13: a rule's direct label sits at one of its ends, and putting it at
  % the far end from the panel's title breaks the stack without moving the rule.
  % The rule is NOT extended past the strip ends, as one refuter proposed: the
  % picture's bounding path is pinned at x = 0.00 and a rule at -0.25 would put
  % the content box at 176.3mm, past S14's 174mm ceiling.
  \node[font=\scriptsize, text=accent, anchor=base east] at (\stripend,8.12)
    {$h \leftarrow h-(hV)V^{\top}$ at every position};
  \draw[accent, line width=0.90pt] (0.00,7.88) -- (\stripend,7.88);

  \tokstrip{7.38}

  \draw[rulegrey, line width=0.40pt] (\bnd,7.24) -- (\bnd,7.82);
  \draw[brc] (\tgt,7.28) -- (\stripend,7.28);
  \node[font=\scriptsize, text=ink, anchor=base] at (9.89,6.88)
    {scored tokens};

  % ---------------- the tie: the figure's central claim, drawn --------------
  % It sits UNDER both strips, on band one's own x origin, and it is now band
  % one's FOOT LINE: on its own baseline, 0.44cm below pass two's brace label
  % and 0.76cm above the two panel heads of band two, so it belongs to the band
  % it closes.
  % PASS 11, TWO REPAIRS. (i) POSITION. It was at 6.88, the EXACT baseline of
  % pass two's brace label, and directly under pass two's prompt half -- the one
  % half of the strip that carries no brace. A reader therefore read it as the
  % missing brace label of that half ("prompt half: the same sentence, box for
  % box | target half: scored tokens"), i.e. as a statement about ONE strip,
  % which is the opposite of what it says. Nothing else is now set on 6.44.
  % (ii) WORDING. "the same sentence, box for box" carried only half of the
  % question this figure has to answer -- are the two passes scoring the same
  % sentence WITH THE SAME WEIGHTS -- and nothing anywhere on the drawing named
  % the weights at all; a reader could infer them only from the formula being an
  % activation edit. The deleted pass-9 subtitle "the same prompt, the same
  % sentence, the same weights" did carry it and pass 10's note claimed the tie
  % had inherited it, which was not true. It now reads "one sentence, one set of
  % weights, scored twice". This also clears the S18 duplication with the
  % caption's "are the same sentence box for box", which is reworded in the
  % staged caption to "emitted by one macro, so no difference between them is
  % possible except the hook".
  \node[def] at (0.00,6.36) {one sentence, one set of weights, scored twice};

  % ---------------- the comparison -----------------------------------------
  % DOUBLE-headed, and ink. The two passes are independent scorings of one
  % fixed sentence by one set of weights; that is a comparison, not a flow, and
  % the single head read as pass one feeding pass two. The single accent
  % arrowhead is reserved for an operation acting in a direction.
  % PASS 11, TWO REPAIRS. (i) GEOMETRY. Its two heads pointed at nothing: it ran
  % between the two brace LABELS' baselines (9.04 to 6.88), so the upper head sat
  % 0.50cm BELOW the strip it was meant to span and both heads terminated in
  % white space. It now runs between the two strips' own vertical centres, 9.75
  % to 7.59, and sits 0.45cm off the strip end instead of 0.75cm, so each head is
  % level with the band it names.
  % (ii) LABEL. "KL from the clean / forward pass" is deleted and replaced by one
  % word. It set panel (d)'s axis quantity string a second time, 130pt away, so a
  % reader could take this vertical rule for an axis OF that quantity; the axis
  % names the quantity, and this mark's job is to say that the two passes are
  % compared and not chained. S20 is satisfied -- the role is stated in one verb
  % of comparison, beside a double head that S11 gives a comparison. The 17
  % characters this frees are what pay for "at layer 4", the reworded tie and the
  % two domain-end tick numerals inside the 6.0% ink ceiling.
  % PASS 12, TWO MOVES. (i) x = 12.87, not 13.10: the strip end moved from 12.65
  % to 12.42 and the arrow keeps pass 11's 0.45cm standoff from it.
  % (ii) THE LABEL IS ROTATED 90 DEGREES, READING UPWARD, BESIDE THE ARROW'S
  % SHAFT AND CENTRED ON ITS MIDPOINT. Set horizontally it sat on baseline 8.62
  % at \scriptsize while panel (b)'s head sat on 8.54 at \footnotesize, 13.3cm
  % away: two lines of type 2.27pt apart in baseline, at two sizes, neither
  % sharing a baseline nor clearly separated, which is the near-alignment a
  % reader reads as carelessness. Rotated, there is no second baseline to
  % near-miss, the label sits ON the mark it names instead of beside it, and it
  % adopts fig:mechanism's own rotated-label idiom. It reads upward, as 4.11's
  % y-labels do, so its ascenders point left and its ink clears the 1.00pt shaft
  % by 0.18cm.
  \draw[ink, line width=1.00pt, cmp] (12.33,9.75) -- (12.33,7.59);
  \node[font=\scriptsize, text=ink, anchor=base, rotate=90]
    at (12.73,8.67) {compared};

% ###########################################################################
% # BAND TWO, LEFT -- (c) what the hook does to one vector.  x origin 0.00
% ###########################################################################
  \node[ttl] at (0.00,5.68) {(c) what the hook does to one vector};

  % THE SLAB. PASS 11, TWO CHANGES, AND THE SECOND PAYS FOR MOST OF THIS PASS.
  % (i) IT RUNS 0.00 TO 5.80, not 0.00 to 3.80. The lower band held a rectangle
  % 38 x 46mm between this panel and panel (d) containing literally no ink --
  % 10% of the content area, and the largest single void on the drawing. The
  % band is the panel's own object and lengthening it is what a subspace does;
  % the construction moves with it to stay on the band's midpoint.
  % (ii) THE rulegrey OUTLINE IS GONE; the band is now a fill alone. It was
  % byte-identical in stroke and fill to the eighteen token cells above it
  % (rulegrey 0.30pt on panelbg), so one rectangle treatment carried two
  % meanings -- "a field of the prompt", eighteen times, and "the purified drug
  % slab, a ten-dimensional subspace", once -- and the only difference was
  % square against rounded corners on a 9pt-tall shape. An ACCENT outline was
  % considered and rejected again for pass 9's reason (accent fill plus accent
  % border is the drug cell's own treatment to within a tint), so the border
  % goes rather than changes colour. The band keeps its accent direct label, is
  % 24 times as wide as it is tall, and is the only fill-without-stroke region
  % in the figure. It also frees 8.7mm2 of the 6.0% ink ceiling, which is what
  % buys "at layer 4", the reworded tie and the two new domain-end ticks.
  \fill[panelbg] (0.00,1.98) rectangle (5.80,2.22);

  % construction lines first, so the arrows and the right angle cover them.
  % PASS 11: the vertical guide stops 0.20cm short of h's tip. It used to end
  % ON the tip, inside the arrowhead, so a leader ran into the mark it points
  % from. The horizontal guide still touches both tips: it IS the closing side
  % of the parallelogram and a gap there would break the construction.
  % PASS 12: the vertical guide now stops at 2.24, 0.14cm above the component's
  % line and just clear of the V_pure band's top edge. It used to run to 2.10,
  % which is exactly the accent component's arrow TIP, so a dashed leader ended
  % inside a mark and at 6x zoom the guide read as a continuation of the
  % component rather than as the side of a rectangle. The same 0.20cm-standoff
  % treatment pass 11 gave the guide's top end, applied to its foot.
  \draw[quietink, line width=0.50pt, densely dotted] (3.40,5.04) -- (3.40,2.24);
  \draw[quietink, line width=0.50pt, densely dotted] (2.50,5.24) -- (3.40,5.24);

  \draw[ink, line width=1.00pt, arw] (2.50,2.10) -- (3.40,5.24);
  \node[lab] at (3.50,5.20) {$h$};

  % PASS 12: THE SURVIVOR IS LABELLED AT MID-SHAFT, NOT AT ITS TIP. At the tip
  % its baseline was 5.16 and $h$'s was 5.20, with the horizontal closing guide
  % between them on 5.24 -- so the row read left to right as
  % "h - (hV)V^T - - - - - h", a dashed RELATION joining two expressions, i.e.
  % an assertion of equality or comparison the panel does not make (the two are
  % a vector and its component, not two sides of anything). Moving one of the
  % two labels off the guide's line is what kills that reading; the guide itself
  % cannot move, because it IS the closing side of the parallelogram and must
  % touch both tips. The label now sits beside the shaft it names, 0.12cm from
  % it and 0.56cm from $h$'s shaft at the same height, so the attachment is
  % unambiguous and no leader is required.
  \draw[ink, line width=1.00pt, arw] (2.50,2.10) -- (2.50,5.24);
  \node[labe] at (2.38,3.62) {$h-(hV)V^{\top}$};

  % the in-slab component, and its fate carried as a mark, not only in words.
  % PASS 11: its label rises from +0.26cm to +0.28cm above the segment and the
  % band thins from 0.32 to 0.24cm, because at the old geometry the label's
  % parenthesis descenders cleared the band's top rule by 0.72pt -- two pixels
  % at 200dpi -- and would have closed on press. The label is still 0.286cm
  % from its own mark, inside the 3mm at which this document requires a leader.
  % PASS 12: the label's x origin moves 3.46 -> 3.58. At 3.46 its opening
  % parenthesis began 0.06cm (2.0pt) right of the vertical guide at 3.40, so at
  % print the dashed guide and the "(" read as one glyph cluster. 0.18cm now.
  \draw[accent, line width=1.00pt, arw] (2.50,2.10) -- (3.40,2.10);
  \node[alab] at (3.58,2.38) {$(hV)V^{\top} \to 0$};

  % the right angle: the component is set to ZERO, not merely reduced. 0.40cm
  % on both legs, against pass 9's 0.16cm, which a reader did not see.
  % PASS 11: IT IS MIRRORED TO THE LEFT OF THE SURVIVING VECTOR. On the right it
  % was unreadable: the 1.0pt black h vector crossed its horizontal leg 0.12cm
  % from the corner, and its vertical leg terminated exactly on the accent
  % component's shaft, so a 0.30pt rulegrey glyph sat under two heavier strokes
  % and read as an artefact. h rises to the RIGHT of the survivor, so the corner
  % between the survivor and the slab on the LEFT is empty; the same 90 degrees,
  % the same two arms, nothing crossing it.
  % PASS 12: THE MARK MOVES TO THE FOOT OF THE PERPENDICULAR, (3.40,2.10).
  % Pass 11 mirrored it to the LEFT of the survivor to escape a collision, and
  % in doing so put it in a quadrant where NEITHER of its arms lies on a drawn
  % ray: the two rays that leave the shared origin (2.50,2.10) are the survivor
  % (up) and the accent component (right), so they bound the up-RIGHT quadrant,
  % while the mark enclosed the up-LEFT one -- its upper arm dead-ended on the
  % survivor's 1.0pt shaft and its lower arm ended at (2.10,2.10) inside the
  % V_pure band's fill, on the backward extension of a ray that is not drawn.
  % At 6x zoom it read as a detached grey bracket, and three refuters read it
  % that way independently.
  % THE ORIGIN CANNOT CARRY THE MARK AT ALL, and that is why this is a move and
  % not a re-mirror: h = survivor + component with both terms positive, so h
  % ALWAYS lies inside the very quadrant the origin's right angle encloses, and
  % at 0.40cm legs h crosses the far arm 0.115cm from its corner whichever way
  % the construction is mirrored. The decomposition is a RECTANGLE -- origin,
  % (3.40,2.10), (3.40,5.24), (2.50,5.24) -- so all four corners carry the same
  % 90 degrees, and the foot of the perpendicular is the one corner that is
  % empty. Its two arms end on drawn strokes: (3.40,2.50) on the vertical
  % construction guide, (3.00,2.10) on the component's own shaft. It asserts
  % what decision 10 says it must -- the surviving vector is orthogonal to the
  % slab, so the slab component is set to ZERO and not merely reduced.
  % IT STAYS rulegrey 0.30pt. One refuter asked for ink 0.40pt on the ground
  % that a load-bearing assertion should not share a stroke with a box border.
  % Overruled: the two construction guides in this same panel are rulegrey
  % 0.30pt, this mark is the third element of that one apparatus, and splitting
  % it off would give 0.40pt a sixth meaning inside S10's five-weight table.
  \draw[rulegrey, line width=0.30pt]
    (3.40,2.50) -- (3.00,2.50) -- (3.00,2.10);

  % ONE label for the band, 0.32cm under its own left edge: the symbol and
  % its gloss in one place, which is the only licence S19 gives the symbol.
  \node[alab] at (0.00,1.66) {$V_{\textrm{pure}}$, the purified drug slab};

% ###########################################################################
% # BAND TWO, RIGHT -- (d) what comes out.  x origin 8.40
% ###########################################################################
  \node[ttl] at (8.40,5.68) {(d) what comes out};

  % VALUE-TO-POSITION MAP, LOGARITHMIC. Braced: it is a pgfmath group, and a
  % shim added outside the braces is not a coordinate, so a label that needs a
  % shim takes it as an xshift on the node.
  %   \kx{v} = 8.40 + (ln v - ln 0.03)/(ln 10 - ln 0.03) * 8.98
  %   domain [0.03, 10] nats; spine 8.40 to 17.38; ticks 0.03 0.1 0.3 1 3 10
  %   0.084 -> 9.9916   1.563 -> 14.5110   4.180 -> 16.0316
  % A horizontal gap is a ratio: null to slab is 4.5194cm and IS 18.6x; slab to
  % cell line is 1.5206cm and IS 37% read the other way.
  %
  % PASS 12, THE DOMAIN CHANGES FROM [0.05, 6.0] TO [0.03, 10] AND NO VALUE
  % MOVES. The pass-11 ladder was 0.05 / 0.1 / 0.3 / 1 / 3 / 6: steps of x2, x3,
  % x3.33, x3, x2, which printed as spacings 1.300 / 2.061 / 2.258 / 2.061 /
  % 1.300cm -- a 74% spread, with a crowded pair of numerals jammed at each end
  % of the rule and the null sitting inside the crowded left pair. A reader
  % cannot infer the scale's rule from a ladder whose steps are not constant,
  % and fig:mechanism(b), the figure this axis was built to imitate, uses a
  % clean decade ladder at uniform spacing. 0.03 / 0.1 / 0.3 / 1 / 3 / 10 is the
  % textbook 1-3 ladder: steps x3.33, x3, x3.33, x3, x3.33, printing as
  % 1.861 / 1.698 / 1.861 / 1.698 / 1.861cm, a 10% spread. The rule still begins
  % and ends on a drawn tick (S13), every mark still has a printed numeral on
  % each side of it, and no mark sits on a tick. Only the value-to-position map
  % changes; 0.084, 1.563 and 4.180 are byte-identical to passes 9, 10 and 11.
  \def\kx#1{{8.40+(ln(#1)-ln(0.03))/(ln(10)-ln(0.03))*8.98}}

  % PASS 12 REBUILDS THIS PANEL, AND IT IS THE PASS'S ONE STRUCTURAL CHANGE.
  % ------------------------------------------------------------------------
  % WHAT WAS WRONG, measured in the pass-11 PDF and reported independently by
  % all three refuters, two of them as blocking:
  %  (a) THE GAP WAS NOT A GAP. The log axis exists for exactly one reason --
  %      on a log scale a horizontal distance IS a ratio -- and the three marks
  %      stood on three different baselines 0.44cm apart (268.30, 255.83,
  %      243.36pt), so there was no single line to lay a ruler along. Worse, the
  %      155.20pt between the null and the purified slab was filled, one row up,
  %      by "the purified drug slab --- 1.563" (x 363.68-470.69) and, two rows
  %      up, by "cell line, the positive control --- 4.180" (394.28-523.00).
  %      The caption asserted a span the drawing did not draw.
  %  (b) IT READ AS A LEGEND. Each row was a text line TERMINATING in its own
  %      disc, so at print each 1.4pt mark read as a bullet ending a line: three
  %      key entries, stacked in descending value order, i.e. the one
  %      construction S22 forbids outright. Pass 11's own note (3) names this
  %      hazard verbatim and then made all three rows trail instead of one.
  %  (c) THE NULL'S LABEL WAS OUTSIDE ITS OWN PANEL. Its two lines began at
  %      x = 245.04 and 241.61pt against a panel origin of 295.45 -- 19.0mm to
  %      the LEFT of the axis's own left end, in panel (c)'s column, 20pt from
  %      the end of panel (c)'s V_pure band, so "random slab" sat above a band
  %      labelled "the purified drug slab" in the same size and face.
  %  (d) THE ROWS WERE MIS-SET VERTICALLY. The null's first line (baseline 4.05)
  %      and the whole slab row (4.19) fell in one connected ink row; and panel
  %      (d)'s head stood 0.98cm above its first row where panel (c)'s stood
  %      0.48cm above its first content, so two side-by-side panels began at
  %      two depths.
  % ------------------------------------------------------------------------
  % WHAT IS DRAWN NOW. THE THREE MARKS SHARE ONE BASELINE, y = 4.02, at their
  % own three x. Nothing else is set on that line, so the null-to-slab distance
  % is a single uninterrupted 4.52cm horizontal span between two discs and the
  % caption's "a horizontal gap IS their ratio" is true of the drawing. Each
  % NAME sits on its own row above, ending (or, for the null, beginning) exactly
  % at its mark's x, and is joined to its mark by a 0.30pt leader in the mark's
  % own colour -- which is S22's own prescription for a label more than 3mm from
  % its mark, and which is what stops any row reading as a key entry, because a
  % mark tied to its name by a drawn leader is not a bullet.
  % PASS 11'S O-2 IS SUPERSEDED, NOT REVERSED. O-2 refused "one baseline with
  % the labels staggered above" on the ground that the three names measure about
  % 11cm against a 9cm axis. That is still true and the labels are still
  % staggered on three rows; what O-2 got wrong is that the MARKS have to
  % stagger with them. They do not: they are 4.52cm and 1.52cm apart and cannot
  % collide.
  % PASS 11'S O-3 IS SUPERSEDED. O-3 refused drop lines because three stems to
  % the spine would have lengths rising with value and would read as a bar
  % chart. These leaders do not touch the spine and their lengths are 1.07,
  % 0.67 and 0.27cm against values 0.084, 4.180 and 1.563 -- deliberately NOT
  % monotone in value, which is why the row order is null / cell line /
  % purified slab and not the pass-11 descending-value order. The row order is
  % also the only one in which no leader crosses another row's label: the
  % null's leader at 9.99 passes left of both rows below it, and the cell line's
  % at 16.03 passes right of the slab row's right end at 14.51.
  % S21, ANSWERED EXACTLY RATHER THAN APPROXIMATELY. A numeral printed near a
  % value axis is read at the position it occupies. Because each label is
  % aligned so that its numeral's INNER edge sits on its own mark's x, each
  % printed value now stands at its own axis position to the character:
  % "0.084" begins at \kx{0.084}, "4.180" ends at \kx{4.180}, "1.563" ends at
  % \kx{1.563}. That is why the null's row is value-first and the other two are
  % value-last -- the value always abuts the mark, and which side the name falls
  % on is set by the measure, the null's mark being only 1.59cm from the axis's
  % left end. No printed value is now within 1.36cm of the spine, so the rule
  % does not carry numerals on both sides.
  \def\my{4.02}
  % row one: the null.  Label runs RIGHT from its mark; the value leads.
  \node[qlab] at (\kx{0.084},5.26) {0.084 --- matched-dimension random slab};
  \draw[quietink, line width=0.30pt] (\kx{0.084},5.16) -- (\kx{0.084},4.09);
  % the null is an OPEN disc, white fill, 0.8pt edge -- 4.11(b)'s own encoding
  % for this same object (mechanism.py: mfc="white", mew=1.1, color=INK).
  \draw[quietink, line width=0.8pt, fill=white] (\kx{0.084},\my) circle (1.4pt);

  % row two: the positive control.  Label runs LEFT; the value trails.
  \node[qlabe] at (\kx{4.180},4.86) {cell line, the positive control --- 4.180};
  \draw[quietink, line width=0.30pt] (\kx{4.180},4.76) -- (\kx{4.180},4.09);
  \fill[quietink] (\kx{4.180},\my) circle (1.4pt);

  % row three: the object under test.
  \node[font=\scriptsize, text=accent, anchor=base east]
    at (\kx{1.563},4.46) {the purified drug slab --- 1.563};
  \draw[accent, line width=0.30pt] (\kx{1.563},4.36) -- (\kx{1.563},4.09);
  \fill[accent] (\kx{1.563},\my) circle (1.4pt);

  % THE TICK LADDER. PASS 11: six numerals, not four, and the outer two are the
  % domain's own ends, so the rule now BEGINS and ENDS on a drawn tick. At four
  % ticks the rule overran its tick lane by 36.87pt (13.0mm) at each end -- 29%
  % of the axis carrying no numeral, the exact defect S13 measures and condemns
  % in fig:comparators row (d) -- and BOTH extreme marks fell in those bare
  % zones, so neither the null nor the positive control could be placed against
  % a printed number. 0.05, 0.1, 0.3, 1, 3, 6 still reads as a log ladder and
  % puts a numeral either side of every mark.
  % The numerals are INK, not quietink. quietink is this document's comparator
  % and null colour (S8/CF4); spending it on the scale while the figure's actual
  % null is labelled in ink inverted the code. fig:mechanism sets every tick
  % numeral in ink and reserves quietink for its annotations.
  \draw[ink, line width=0.40pt] (8.40,3.44) -- (17.38,3.44);
  \foreach \v in {0.03, 0.1, 0.3, 1, 3, 10}{
    \draw[rulegrey, line width=0.30pt] (\kx{\v},3.44) -- (\kx{\v},3.34);}
  \foreach \v/\t in {0.1/0.1, 0.3/0.3, 1/1, 3/3}{
    \node[tick] at (\kx{\v},3.28) {\t};}
  % the two domain-end numerals are pulled flush INSIDE the rule's own ends, so
  % no ink crosses the drawn measure (S14). Centred on their ticks they would
  % have overhung x = 17.38 by 0.09cm and the float overfull by 1.88pt.
  % PASS 12 CONSIDERED AND OVERRULED centring all six and cutting the rule to
  % 17.10 to pay for it. At the LEFT end that is not available: "0.03" centred
  % on 8.40 begins at 8.13, which is 0.27cm outside panel (d)'s x origin, and
  % S13's one-grid rule -- measured at a 0.2pt spread by pass 11's repair (15)
  % -- outranks the optical centring of a domain-end numeral that labels no
  % mark. The residual is 0.275cm on "0.03" and 0.16cm on "10", which on this
  % ladder is 7% and 4% of a decade.
  \node[tick, anchor=north west] at (8.40,3.28) {0.03};
  \node[tick, anchor=north east] at (17.38,3.28) {10};

  % THE AXIS BLOCK. Quantity and unit on line one, scope on line two, both
  % left-aligned to the axis's own left end, directly under the tick numerals.
  % No result, no interval disclaimer and no cross-reference may sit here.
  \node[ttl]  at (8.40,2.68)
    {KL from the clean forward pass (nats), logarithmic scale};
  \node[qlab] at (8.40,2.34)
    {mean over the target tokens of 60 teacher-forced prompts, layer 4};

% ###########################################################################
% # THE FIGURE'S ONE INTERVAL DECLARATION.  x origin 0.00
% ###########################################################################
  % PASS 12: SET ON TWO LINES, NOT ONE. As one line it measured 15.89cm and was
  % the WIDEST OBJECT IN THE FIGURE -- wider than the axis (8.98), wider than
  % band one (14.54) -- which S14's check forbids outright ("the widest object
  % must be an axis or a panel, never a label"), and a 15.89cm line of 8pt grey
  % stranded a centimetre below everything else reads as a second, smaller
  % caption set above the real one. S17 expressly allows two lines. Broken at
  % the colon, the block measures about 7.1cm, and its two baselines are set so
  % the picture's lowest ink is unchanged and the height cap is not touched.
  \node[qlab] at (0.00,1.04)
    {Intervals: none in (d) --- the stored dispersion is a normal-approximation};
  \node[qlab] at (0.00,0.70)
    {standard error, not this document's convention. (a) to (c) carry no measured mark.};

  % bounding box, so the picture is exactly the measure it claims
  \path (0,0.58) rectangle (17.38,10.62);
\end{tikzpicture}
%
% ---------------------------------------------------------------------------
% READY TO PASTE into Sections/04b-representation.tex, at sec:res-used,
% immediately after the paragraph "Step four, the causal measurement in logit
% space" (04b:261-275). The caption below is the PASS 11 replacement AS REVISED
% AT PASS 12. It is also staged, keyed by \label{fig:hook}, in
% figs/_PENDING_CAPTIONS.tex -- a file PASS 11 CREATED, because the pass-10 note
% in this position asserted it existed and it did not. NO BUILDER EDITS
% Sections/. The staged copy carries the pass-12 revision note; the two copies
% are kept byte-equal in the \caption itself and figs/_PENDING_CAPTIONS.tex is
% the authoritative one.
% PASS 12 REVISED THREE THINGS IN IT, none of them a numeral already there:
%   * panel (d)'s marks are described as standing on ONE LINE, tied to their
%     names by leaders, because pass 11's three baselines meant the "horizontal
%     gap IS the ratio" sentence described a span that was not on the page;
%   * the POSITIVE CONTROL's size confound is disclosed for the first time
%     (cell line's subspace fraction is 20.9x the purified slab's at layer 4,
%     15-22x over the seven -- 04b:370), beside the null's 5.6x, because the
%     37% was being quoted against a comparator matched on nothing at all;
%   * the random-draw disclosure is strengthened from "one random draw" to the
%     fact it understated: with two stored draws at seven layers the file holds
%     fourteen ratios and the drawn 18.634 is the largest of them, and the two
%     draws swap order with depth (at layer 6, 14.7x undrawn against 5.7x
%     drawn).
% Two figures inside the DELETIONS block of the staged copy are superseded by
% the new log domain: the null-to-slab gap is 4.510cm, not 5.474, and the
% slab-to-cell-line gap is 1.520cm, not 1.845. No value moved; the map did.
%
% \begin{figure}[htbp]
% \begin{fullwidth}
% \centering
% \input{figs/hook}
% \caption[The causal measurement]{\textbf{Nothing is generated: one true
% sentence is teacher-forced and scored, a hook projects the purified drug slab
% out at every position, and the identical tokens are scored again --- and
% removing the slab costs $18.6\times$ what removing a matched-dimension random
% slab costs, which is still only $37\%$ of what removing cell line costs.}
% Panels (a) and (b) are emitted by one macro, so no difference between them is
% possible except the hook. Its accent rule spans the whole strip, prompt half
% included, and runs straight across the prompt/target divider, because the
% projection is applied at every position of layer 4 --- an intervention on the
% computation, not on generation, and the divider marks only where a generated
% continuation would have begun. The braces mark the tokens the mean KL is read
% over, identically in both passes, and the double-headed arrow is that
% comparison: two independent scorings of one fixed sentence by one set of
% weights, not a flow from the first into the second. Panel (d) is logarithmic
% and its three marks stand on one line above the scale, each tied to its name
% by a leader, so a horizontal gap between two marks \emph{is} their ratio ---
% null to purified slab is the $18.6\times$, purified slab to cell line is the
% $37\%$ read the other way --- and the filled disc is the object under test
% while the open disc is the null, as in \cref{fig:mechanism}(b). Provenance:
% \texttt{workspace\_probe.json}, \texttt{layers.4.kl}, $n = 60$ prompts each
% --- purified drug slab $1.563$, matched-dimension random slab $0.084$,
% cell-line positive control $4.180$ nats, each printed at its own mark.
% \emph{What is schematic}: the token strip implies no token count, since a real
% target is a cell sentence of hundreds of gene tokens which the ellipsis box
% stands for, and the five prompt fields are those of \cref{sec:methods-data},
% drawn as they are in \cref{fig:licensing}; and nothing in panel (c) is to
% scale --- $V_{\textrm{pure}}$ is ten directions of the $2048$ the residual
% stream carries, but the panel draws the surviving component about three and a
% half times the deleted one, which is the right direction and not the right
% ratio, and $10/2048$ is drawn true once, in \cref{fig:slab}. What may not be
% read across: \textbf{no interval is drawn}; the line under the drawing states
% the reason, and as in \cref{tab:workspace} these are sign-and-magnitude
% readings. Neither gap is matched on size. The null matches the slab's
% \emph{rank} and not the variance it removes --- the purified slab's subspace
% fraction is $5.6\times$ the null's here --- and the positive control is matched
% on neither: cell line's subspace fraction is $20.9\times$ the purified slab's
% at this layer and $15$--$22\times$ across the seven, so the $18.6\times$ and
% the $37\%$ each mix drug content with a size difference, and in opposite
% directions. Nor is the null one thing: the same file stores a second draw of
% the identical rank-$10$ construction at every layer
% (\texttt{layers.4.kl.random}, $0.198$ nats here) on which this ratio would read
% $7.9\times$, so the drawn $18.6\times$ is the largest of the fourteen
% (layer, draw) combinations the file holds, and the two draws disagree in both
% directions with depth --- at layer 6 the undrawn draw gives $14.7\times$
% against the drawn $5.7\times$. \cref{tab:workspace}'s Ratio column inherits the
% same single draw. Layer 4 is likewise one layer of seven, chosen because it is
% the slab's strongest showing against the null ($18.6\times$, the top of
% $3.6$--$18.6\times$) and against cell line. That $37\%$ is the top of the
% $9$--$37\%$ band \emph{at layers $2$--$9$ only}: over all seven layers the drug
% slab's ablation cost runs from $8.90\%$ of the cell-line slab's at layer 9 to
% $98.45\%$ at layer 12, and it is because layer 12 is that near-equality that
% layer 12 is not the layer drawn. \cref{tab:workspace} carries all seven and
% \cref{fig:mechanism}(b) plots them.}
% \label{fig:hook}
% \end{fullwidth}
% \end{figure}
%
% CAPTION NOTE, PASS 10. Six changes, each forced by a change to the drawing or
% by the contract that a result lives in the caption:
%   * THE LEDE NOW STATES THE FINDING, with both ratios, because both were
%     deleted from under the axis where they had been sitting in the axis-label
%     lane. The old lede described the apparatus only.
%   * "The two arrows off the rule mark that it acts in both halves and imply no
%     count of positions" is deleted: the arrows are gone, and a caption that
%     has to unsay a mark was the reason to delete them.
%   * THE DOUBLE-HEADED ARROW IS NAMED, and named as a COMPARISON, because the
%     old single head was read as pass one feeding pass two -- the misreading
%     this figure exists to kill.
%   * THE LOG AXIS IS NAMED, and the sentence that a gap IS a ratio is what
%     licenses a reader to check $18.6\times$ off the drawing. The open/filled
%     disc code is named too, and tied to fig:mechanism(b), which is where it
%     comes from.
%   * "the divider marks only where a generated continuation would have begun"
%     is new, and carries what the drawing lost when the divider's own label,
%     "generation would begin here", was deleted for the ink ceiling.
%   * "$V_{\textrm{pure}}$ is ten directions of the 2048 the residual stream
%     carries" is new, and carries what panel (c) lost when its subtitle was
%     replaced by the band's one-line direct label. The panel's true ratio is
%     also restated: 3.49:1 as drawn, against 14.3:1 true.
%   * THE QUALIFIER THE OLD CAPTION DROPPED IS RESTORED. "$9$--$37\%$" is the
%     range over layers 2--9 only. Recomputed over all seven,
%     kl.pure_drug[0]/kl.cell_line[0] is 0.1388, 0.3740, 0.3331, 0.1100, 0.0890,
%     0.9845, 0.2501 at layers 2/4/6/8/9/12/16, i.e. 8.90% to 98.45%. The old
%     caption gave "9--37%" unqualified and then restored the information one
%     sentence later as "at layer 12 the two costs are nearly equal (98%)"; it
%     is now stated once, correctly, in one place.
% ---------------------------------------------------------------------------

% \caption[The causal measurement]{\textbf{Nothing is generated: one true
% sentence is teacher-forced and scored, a hook projects the purified drug slab
% out at every position, and the identical tokens are scored again --- and
% removing the slab costs $18.6\times$ what removing a matched-dimension random
% slab costs, which is still only $37\%$ of what removing cell line costs.}
% Panels (a) and (b) are emitted by one macro, so no difference between them is
% possible except the hook. Its accent rule spans the whole strip, prompt half
% included, and runs straight across the prompt/target divider, because the
% projection is applied at every position of layer 4 --- an intervention on the
% computation, not on generation, and the divider marks only where a generated
% continuation would have begun. The braces mark the tokens the mean KL is read
% over, identically in both passes, and the double-headed arrow is that
% comparison: two independent scorings of one fixed sentence by one set of
% weights, not a flow from the first into the second. Panel (d) is logarithmic
% and its three marks stand on one line above the scale, each tied to its name
% by a leader, so a horizontal gap between two marks \emph{is} their ratio ---
% null to purified slab is the $18.6\times$, purified slab to cell line is the
% $37\%$ read the other way --- and the filled disc is the object under test
% while the open disc is the null, as in \cref{fig:mechanism}(b). Provenance:
% \texttt{workspace\_probe.json}, \texttt{layers.4.kl}, $n = 60$ prompts each
% --- purified drug slab $1.563$, matched-dimension random slab $0.084$,
% cell-line positive control $4.180$ nats, each printed at its own mark.
% \emph{What is schematic}: the token strip implies no token count, since a real
% target is a cell sentence of hundreds of gene tokens which the ellipsis box
% stands for, and the five prompt fields are those of \cref{sec:methods-data},
% drawn as they are in \cref{fig:licensing}; and nothing in panel (c) is to
% scale --- $V_{\textrm{pure}}$ is ten directions of the $2048$ the residual
% stream carries, but the panel draws the surviving component about three and a
% half times the deleted one, which is the right direction and not the right
% ratio, and $10/2048$ is drawn true once, in \cref{fig:slab}. What may not be
% read across: \textbf{no interval is drawn}; the line under the drawing states
% the reason, and as in \cref{tab:workspace} these are sign-and-magnitude
% readings. Neither gap is matched on size. The null matches the slab's
% \emph{rank} and not the variance it removes --- the purified slab's subspace
% fraction is $5.6\times$ the null's here --- and the positive control is matched
% on neither: cell line's subspace fraction is $20.9\times$ the purified slab's
% at this layer and $15$--$22\times$ across the seven, so the $18.6\times$ and
% the $37\%$ each mix drug content with a size difference, and in opposite
% directions. Nor is the null one thing: the same file stores a second draw of
% the identical rank-$10$ construction at every layer
% (\texttt{layers.4.kl.random}, $0.198$ nats here) on which this ratio would read
% $7.9\times$, so the drawn $18.6\times$ is the largest of the fourteen
% (layer, draw) combinations the file holds, and the two draws disagree in both
% directions with depth --- at layer 6 the undrawn draw gives $14.7\times$
% against the drawn $5.7\times$. \cref{tab:workspace}'s Ratio column inherits the
% same single draw. Layer 4 is likewise one layer of seven, chosen because it is
% the slab's strongest showing against the null ($18.6\times$, the top of
% $3.6$--$18.6\times$) and against cell line. That $37\%$ is the top of the
% $9$--$37\%$ band \emph{at layers $2$--$9$ only}: over all seven layers the drug
% slab's ablation cost runs from $8.90\%$ of the cell-line slab's at layer 9 to
% $98.45\%$ at layer 12, and it is because layer 12 is that near-equality that
% layer 12 is not the layer drawn. \cref{tab:workspace} carries all seven and
% \cref{fig:mechanism}(b) plots them.}
% \label{fig:hook}
% \end{fullwidth}
% \end{figure}
%
% CAPTION NOTE, PASS 10. Six changes, each forced by a change to the drawing or
% by the contract that a result lives in the caption:
%   * THE LEDE NOW STATES THE FINDING, with both ratios, because both were
%     deleted from under the axis where they had been sitting in the axis-label
%     lane. The old lede described the apparatus only.
%   * "The two arrows off the rule mark that it acts in both halves and imply no
%     count of positions" is deleted: the arrows are gone, and a caption that
%     has to unsay a mark was the reason to delete them.
%   * THE DOUBLE-HEADED ARROW IS NAMED, and named as a COMPARISON, because the
%     old single head was read as pass one feeding pass two -- the misreading
%     this figure exists to kill.
%   * THE LOG AXIS IS NAMED, and the sentence that a gap IS a ratio is what
%     licenses a reader to check $18.6\times$ off the drawing. The open/filled
%     disc code is named too, and tied to fig:mechanism(b), which is where it
%     comes from.
%   * "the divider marks only where a generated continuation would have begun"
%     is new, and carries what the drawing lost when the divider's own label,
%     "generation would begin here", was deleted for the ink ceiling.
%   * "$V_{\textrm{pure}}$ is ten directions of the 2048 the residual stream
%     carries" is new, and carries what panel (c) lost when its subtitle was
%     replaced by the band's one-line direct label. The panel's true ratio is
%     also restated: 3.49:1 as drawn, against 14.3:1 true.
%   * THE QUALIFIER THE OLD CAPTION DROPPED IS RESTORED. "$9$--$37\%$" is the
%     range over layers 2--9 only. Recomputed over all seven,
%     kl.pure_drug[0]/kl.cell_line[0] is 0.1388, 0.3740, 0.3331, 0.1100, 0.0890,
%     0.9845, 0.2501 at layers 2/4/6/8/9/12/16, i.e. 8.90% to 98.45%. The old
%     caption gave "9--37%" unqualified and then restored the information one
%     sentence later as "at layer 12 the two costs are nearly equal (98%)"; it
%     is now stated once, correctly, in one place.
% ---------------------------------------------------------------------------

% \caption[The causal measurement]{\textbf{Nothing is generated: one true
% sentence is teacher-forced and scored, a hook projects the purified drug slab
% out at every position, and the identical tokens are scored again --- and
% removing the slab costs $18.6\times$ what removing a matched-dimension random
% slab costs, which is still only $37\%$ of what removing cell line costs.}
% Panels (a) and (b) are emitted by one macro, so no difference between them is
% possible except the hook. Its accent rule spans the whole strip, prompt half
% included, and runs straight across the prompt/target divider, because the
% projection is applied at every position of layer 4 --- an intervention on the
% computation, not on generation, and the divider marks only where a generated
% continuation would have begun. The braces mark the tokens the mean KL is read
% over, identically in both passes, and the double-headed arrow is that
% comparison: two independent scorings of one fixed sentence by one set of
% weights, not a flow from the first into the second. Panel (d) is logarithmic
% and its three marks stand on one line above the scale, each tied to its name
% by a leader, so a horizontal gap between two marks \emph{is} their ratio ---
% null to purified slab is the $18.6\times$, purified slab to cell line is the
% $37\%$ read the other way --- and the filled disc is the object under test
% while the open disc is the null, as in \cref{fig:mechanism}(b). Provenance:
% \texttt{workspace\_probe.json}, \texttt{layers.4.kl}, $n = 60$ prompts each
% --- purified drug slab $1.563$, matched-dimension random slab $0.084$,
% cell-line positive control $4.180$ nats, each printed at its own mark.
% \emph{What is schematic}: the token strip implies no token count, since a real
% target is a cell sentence of hundreds of gene tokens which the ellipsis box
% stands for, and the five prompt fields are those of \cref{sec:methods-data},
% drawn as they are in \cref{fig:licensing}; and nothing in panel (c) is to
% scale --- $V_{\textrm{pure}}$ is ten directions of the $2048$ the residual
% stream carries, but the panel draws the surviving component about three and a
% half times the deleted one, which is the right direction and not the right
% ratio, and $10/2048$ is drawn true once, in \cref{fig:slab}. What may not be
% read across: \textbf{no interval is drawn}; the line under the drawing states
% the reason, and as in \cref{tab:workspace} these are sign-and-magnitude
% readings. Neither gap is matched on size. The null matches the slab's
% \emph{rank} and not the variance it removes --- the purified slab's subspace
% fraction is $5.6\times$ the null's here --- and the positive control is matched
% on neither: cell line's subspace fraction is $20.9\times$ the purified slab's
% at this layer and $15$--$22\times$ across the seven, so the $18.6\times$ and
% the $37\%$ each mix drug content with a size difference, and in opposite
% directions. Nor is the null one thing: the same file stores a second draw of
% the identical rank-$10$ construction at every layer
% (\texttt{layers.4.kl.random}, $0.198$ nats here) on which this ratio would read
% $7.9\times$, so the drawn $18.6\times$ is the largest of the fourteen
% (layer, draw) combinations the file holds, and the two draws disagree in both
% directions with depth --- at layer 6 the undrawn draw gives $14.7\times$
% against the drawn $5.7\times$. \cref{tab:workspace}'s Ratio column inherits the
% same single draw. Layer 4 is likewise one layer of seven, chosen because it is
% the slab's strongest showing against the null ($18.6\times$, the top of
% $3.6$--$18.6\times$) and against cell line. That $37\%$ is the top of the
% $9$--$37\%$ band \emph{at layers $2$--$9$ only}: over all seven layers the drug
% slab's ablation cost runs from $8.90\%$ of the cell-line slab's at layer 9 to
% $98.45\%$ at layer 12, and it is because layer 12 is that near-equality that
% layer 12 is not the layer drawn. \cref{tab:workspace} carries all seven and
% \cref{fig:mechanism}(b) plots them.}
% \label{fig:hook}
% \end{fullwidth}
% \end{figure}
%
% CAPTION NOTE, PASS 10. Six changes, each forced by a change to the drawing or
% by the contract that a result lives in the caption:
%   * THE LEDE NOW STATES THE FINDING, with both ratios, because both were
%     deleted from under the axis where they had been sitting in the axis-label
%     lane. The old lede described the apparatus only.
%   * "The two arrows off the rule mark that it acts in both halves and imply no
%     count of positions" is deleted: the arrows are gone, and a caption that
%     has to unsay a mark was the reason to delete them.
%   * THE DOUBLE-HEADED ARROW IS NAMED, and named as a COMPARISON, because the
%     old single head was read as pass one feeding pass two -- the misreading
%     this figure exists to kill.
%   * THE LOG AXIS IS NAMED, and the sentence that a gap IS a ratio is what
%     licenses a reader to check $18.6\times$ off the drawing. The open/filled
%     disc code is named too, and tied to fig:mechanism(b), which is where it
%     comes from.
%   * "the divider marks only where a generated continuation would have begun"
%     is new, and carries what the drawing lost when the divider's own label,
%     "generation would begin here", was deleted for the ink ceiling.
%   * "$V_{\textrm{pure}}$ is ten directions of the 2048 the residual stream
%     carries" is new, and carries what panel (c) lost when its subtitle was
%     replaced by the band's one-line direct label. The panel's true ratio is
%     also restated: 3.49:1 as drawn, against 14.3:1 true.
%   * THE QUALIFIER THE OLD CAPTION DROPPED IS RESTORED. "$9$--$37\%$" is the
%     range over layers 2--9 only. Recomputed over all seven,
%     kl.pure_drug[0]/kl.cell_line[0] is 0.1388, 0.3740, 0.3331, 0.1100, 0.0890,
%     0.9845, 0.2501 at layers 2/4/6/8/9/12/16, i.e. 8.90% to 98.45%. The old
%     caption gave "9--37%" unqualified and then restored the information one
%     sentence later as "at layer 12 the two costs are nearly equal (98%)"; it
%     is now stated once, correctly, in one place.
% ---------------------------------------------------------------------------

\caption[The causal measurement]{\textbf{Nothing is generated: one true
sentence is teacher-forced and scored, a hook projects the purified drug slab
out at every position, and the identical tokens are scored again --- and
removing the slab costs $18.6\times$ what removing a matched-dimension random
slab costs, which is still only $37\%$ of what removing cell line costs.}
Panels~(a) and~(b) are emitted by one macro and are therefore the same sentence
box for box; the only difference between them is the hook. Its accent rule
spans the whole strip, prompt half included, and runs straight across the
prompt/target divider, because $h \leftarrow h - (hV)V^{\top}$ is applied at
every position --- an intervention on the computation, not on generation, and
the divider marks only where a generated continuation would have begun. The
braces mark the tokens the mean KL is read over, identically in both passes,
and the double-headed arrow is that comparison: two independent scorings of one
fixed sentence by one set of weights, not a flow from the first into the
second. Panel~(d) is logarithmic, so a horizontal gap between two marks
\emph{is} their ratio --- null to purified slab is the $18.6\times$, purified
slab to cell line is the $37\%$ read the other way. Three marks, three glyphs:
the accent disc is the object under test, the open disc is the null, and the
grey disc is the cell-line positive control, a real measurement that is not the
headline. Provenance: \texttt{workspace\_probe.json}, \texttt{layers.4.kl}, $n
= 60$ prompts each --- purified drug slab $1.563$, matched-dimension random
slab $0.084$, cell-line positive control $4.180$ nats, each printed at its own
mark. \emph{What is schematic}: the token strip implies no token count, since a
real target is a cell sentence of hundreds of gene tokens which the ellipsis
box stands for, and the five prompt fields are those of
\cref{sec:methods-data}, drawn at the same widths as in \cref{fig:licensing};
and nothing in panel~(c) is to scale --- $V_{\textrm{pure}}$ is ten directions
of the $2048$ the residual stream carries, but the panel draws the surviving
component about three and a half times the deleted one, which is the right
direction and not the right ratio. $10/2048$ of the squared norm is drawn true
once, in \cref{fig:slab}, and is not repeated here. What may not be read
across: \textbf{no interval is drawn}, because the stored dispersion is a
normal-approximation standard error over the $60$ prompts and not this
document's interval convention (95\% bootstrap clustered over cell lines), so
these are sign-and-magnitude readings as in \cref{tab:workspace}; and layer $4$
is one layer of seven, chosen because it is the slab's strongest showing
against the null ($18.6\times$, the top of $3.6$--$18.6\times$) and against
cell line. That $37\%$ is the top of the $9$--$37\%$ band \emph{at layers $2$
to $9$ only}: over all seven layers the drug slab's ablation cost runs from
$8.90\%$ of the cell-line slab's at layer $9$ to $98.45\%$ at layer $12$, and
it is because layer $12$ is that near-equality that layer $12$ is not the layer
drawn. \cref{tab:workspace} carries all seven and \cref{fig:mechanism}(b) plots
them.}
\label{fig:hook}
\end{fullwidth}
\end{figure}

% ===========================================================================
\begin{table}[htbp]
\begin{fullwidth}
\begin{threeparttable}
\caption{\textbf{Drug-label information becomes more decodable with depth while
the same slab carries little functional weight.} All seven layers pass the
removal gate (88--95\% of above-chance drug signal removed); chance is $0.083$.
The last column is ablation cost divided by the matched-dimension random null on
the same layer, a sign-and-magnitude reading and not an estimate with an
interval. The null matches the slab's rank, not the variance it removes --- the
slab's subspace fraction is roughly six times the null's --- so the ratio mixes
drug content with that difference.}
\label{tab:workspace}
\small
\begin{tabular}{@{}lrrrrr@{}}
\toprule
\thead{Layer} & \thead{Decode} & \thead{Decode} & \thead{Ablate pure}
  & \thead{Random} & \thead{Ratio} \\
 & \thead{(raw)} & \thead{(context removed)} & \thead{drug (KL)}
  & \thead{matched dim} & \\
\midrule
2 & 0.408 & 0.354 & 0.155 & 0.012 & $12.7\times$ \\
4 & 0.585 & 0.429 & 1.563 & 0.084 & $18.6\times$ \\
6 & 0.692 & 0.529 & 1.435 & 0.251 & $5.7\times$ \\
8 & 0.779 & 0.569 & 0.596 & 0.042 & $14.1\times$ \\
9 & \textbf{0.821} & 0.602 & 0.418 & 0.100 & $4.2\times$ \\
12 & 0.754 & 0.873 & 0.133 & 0.035 & $3.8\times$ \\
16 & 0.758 & \textbf{0.883} & 0.015 & 0.004 & $3.6\times$ \\
\bottomrule
\end{tabular}
\begin{tablenotes}[flushleft]\footnotesize
\item Decode columns are accuracies over twelve drugs, answering different
questions and not two estimates of one quantity. The raw column includes whatever
cell line, plate and dose contribute; the other is the same probe after those
directions are removed.
\end{tablenotes}
\end{threeparttable}
\end{fullwidth}
\end{table}

% ===========================================================================
\paragraph{What the ablation shows is a dissociation} With cell line, plate and
dose projected out, drug-label decodability rises with depth, from $0.354$ at
layer 2 to $0.883$ at layer 16. Yet the purified slab's descriptive subspace
fraction stays flat at $\approx 0.03$ against $\approx 0.6$ for cell line and
$\approx 0.005$ for the matched-dimension null, an
activation-space summary and not a decomposition of biological drug-response
variance, and ablating it costs $3.6$--$18.6\times$ the matched-dimension random
null. \Cref{fig:mechanism} draws the dissociation whole: the rise, the flat
fraction beneath it, and each layer's ablation cost beside its random null.

That is small beside the cell-line positive control, which the same ablation
reads at $1.0$--$11.2\times$ the drug's. Removing the drug subspace costs
$9$--$37\%$ of what removing cell line costs at layers 2--9, with layer 12 the
single exception at $98\%$ before layer 16 falls back to $25\%$. On subspace
fraction cell line measures $15$--$22\times$ the purified slab. Sensitivity is
established at a bit player's weight, across all seven layers.

\begin{figure}[htbp]
\begin{fullwidth}
\centering
\includegraphics{fig-mechanism}
\caption[Decodable with depth, at a bit player's weight]{\textbf{Drug-label
information becomes more decodable with depth while interventions on the
purified slab have limited effect.} Layer is a true numeric axis in both
panels, so the $9$--$12$--$16$ spacing is drawn as it is; the two dots at the
far left of~(b), outside the tick range, are the two direct-label keys and not
data. \emph{(a) Top:} drug-label probe accuracy over \emph{twelve} drug
classes, so chance is $1/12 = 0.083$, on \num{480} held-out predictions. Two
series: the dotted open one is the raw probe, and the solid filled one is the
same probe once cell line, plate and dose are projected out. They \emph{cross}
between layers $9$ and $12$ --- removing context lowers decodability at the
five shallowest depths and raises it at the two deepest, to $0.873$ and
$0.883$, and this panel is the only place in the document where that is
visible. The flat grey line is the same probe after the drug slab is projected
out, measured at all seven depths and identical at each. \emph{Bottom:} the
descriptive subspace fraction of the purified drug slab against the cell-line
positive control, on its own separate scale --- an activation-space summary,
not a decomposition of biological drug-response variance. The double-headed
arrow at layer $16$ joins one point in each half; it marks that the two
quantities move apart, and its \emph{length} measures nothing, the two scales
being incommensurable. \emph{Encoded, not read} names that contrast, but the
evidence for \emph{not read} is panel~(b)'s ablation cost and not this panel,
which is descriptive throughout. \emph{(b)} Mean KL from the clean forward pass
in \emph{nats} when the purified drug slab is ablated, against a
matched-dimension random slab, $n = 60$ prompts per entry. Log axis because
late-layer divergences are small for everything, and paired dots rather than
bars: on a log axis a bar's length is measured from an arbitrary axis floor,
whereas the vertical gap between a pair \emph{is} their ratio, which is the
number annotated beside each pair and the number in \cref{tab:workspace}. That
table's caveat holds here: the null is matched to the slab's rank and not to
the variance the slab removes --- the slab's subspace fraction runs $4.9$ to
$7.0$ times the null's, plotted in the lower half of~(a) --- so the ratio mixes
drug content with that difference. The cell-line positive control, against
which the ablation cost is $9$ to $37\%$ at layers $2$ to $9$, is not drawn
here; \cref{fig:hook}(d) draws it at layer $4$ and \cref{tab:workspace} carries
all seven. \textbf{No intervals are drawn in either panel, for two different
reasons.} In~(b) the stored dispersions are normal-approximation standard
errors over the $60$ prompts, not this document's interval convention (95\%
bootstrap clustered over cell lines), so as in \cref{tab:workspace} this is a
sign-and-magnitude reading. In~(a) no dispersion was computed at all: the
accuracies are stored as bare floats. These accuracies come from a different
instrument from the twenty-four-class probe sweep run over the five arms of
\cref{sec:res-mechanism}, whose chance is $1/24$ on a different label set, and
the two must not be read against each other. Whether drug \emph{identity} is
read requires \cref{sec:res-mechanism}.}
\label{fig:mechanism}
\end{fullwidth}
\end{figure}

% The direct-arm fold sits after fig:mechanism and before "Where ablation stops",
% so that paragraph's last sentence still hands off directly to sec:res-mechanism.
% The pre-registered joint reading rule is not cleared: rel_B is below the 0.40
% gate, so the "only ever a label" reading is withheld and the arm is reported as
% a bound on a fitted map. "leave-one-drug-out" is avoided; the glossary reserves
% it for the generic construction. Frame changes to residual here, marked with
% \framecue on the precedent of 04a line 465, not by changing \beatsection{E}.
% The drafted phrase "at one dose" is dropped: the roster's doses are not uniform
% across drugs (biology/dose_frame/per_drug), and this arm partials nothing out,
% so naming a dose would assert a control the measurement does not carry.
% The drafted \gloss{dual form} is demoted into the sentence: as a margin note it
% collided in this float-dense stretch and cost a "Marginpar on page 48 moved".
% \framecue is the one margin note here, and it stays.
%
% REPAIR PASS, after three refuters read the draft.
% (1) X is the K-averaged residual stream at one prompt position and is never
% projected onto V_pure, so the slab does not enter this measurement at all.
% The body now says "residual stream", the verdict says "this model's
% activations", and the roster fence is paired with a stream/slab fence. The
% draft verdict read "what a fitted map recovers from the slab is bounded below
% the name string", which asserted the slab, generalised one ridge to fitted
% maps at large, and used "bounded below" for its opposite. It is gone.
% (2) The cell-line positive control was run with the RSA statistic against a
% 2000-draw null and never with this ridge: no positive_control cell in either
% file carries a direct_arm block. So the joint rule's control half is met only
% by borrowing, and the draft's flat "the control is positive" contradicted the
% concession three paragraphs above it.
% (3) Layer attribution. The activations median and every permutation figure
% are layer 12. The layer-16 median is 0.999916, and there the null's q95
% exceeds its own observed value, so neither travels.
% (4) alphas_chosen_median sits at the grid floor, but alphas_chosen_max is 100
% at layer 12 and 1e5, the grid ceiling, at layer 16. The floor claim is the
% median fold's, and it is the stated premise for "interpolates".
% (5) The mean-r increment is negative at all 24 grid cells in each stratum,
% but the two prompt variants are byte-identical at the drug-name position, so
% 24 cells are 20 distinct measurements. The count is dropped; "every setting"
% carries the breadth without asserting 24 independent confirmations.
% (6) Correlations are printed to five places so that the printed increment is
% the difference of the printed operands. At four places 0.9767 - 0.9323 reads
% 0.0444 against a stated -0.0445, and a reader who subtracts gets a third
% number.
\paragraph{The ridge does better from the name than from the activations} On
\num{50} tier-2 drugs, none of them seen in fine-tuning, a ridge is fitted from
the fine-tuned checkpoint's residual stream at the last prompt position to each
drug's measured residual profile, in one cell line.\framecue{\Rframe}{this
result and its scope limit are residual-frame, \cref{sec:methods-residual}} It
is fitted in the dual form, through the drug-by-drug kernel rather than the far
larger activation covariance, one drug held out at a time. The penalty is
chosen inside each training fold by holding out one more drug, never fixed in
advance. The activations are averaged over sixteen control sentences; the target
is that drug's pseudobulk minus the roster mean, signed-rank encoded as in
\cref{sec:methods-residual}. Both sides are centred across drugs, not referenced
to a plate as that section's target is. No score here belongs beside a
residual-frame $\NIR$. The identical ridge is then fitted from the drug's name
alone, character trigrams plus indicators for the tokens the model splits that
name into. No confound is partialled out; the name features are the control,
fixed across layer, position and checkpoint, and the quantity reported is the
increment of activations over names in held-out mean correlation. The roster is
\num{50} drugs and not the twelve of the steps above, and the two may not be
read across. The ridge reads the whole stream, not the purified slab.

% KNOWN RESIDUAL. This paragraph sets fourteen lines and its last two fall past
% fig:mechanism onto the next page, so "the fifth | decimal place" breaks across
% a page turn and a twenty-line caption. \looseness=-1 and -2 were both tried
% and TeX finds no shorter setting. Closing it needs about twenty-six words out
% of this paragraph and the one above, which costs either the dual form's
% rationale or the fixed-comparator fact, and a refuter put each of those here.
% The claim box, which preamble.tex 882-897 names as the break these hairlines
% exist to prevent, does close whole; that was the break worth buying.
Held-out mean correlation at layer 12 reads $0.93225$ from the activations
against $0.97672$ from the name features, an increment of $-0.0445$ that layer
16 matches to four places. The sign is negative at every setting of a grid
crossing two checkpoints, two prompt variants, two layers and three positions,
and again at every setting in a second cell line, where a shared well structure
makes that agreement weaker than it looks. The pre-registration licenses the
increment at none of them. Neither mean describes the typical drug. Median
held-out correlation at layer 12 is $0.999953$ from the activations and
$0.999997$ from the name features, so on the typical drug the two fits differ in
the fifth decimal place, and each mean is a tail statistic over a few badly fit
drugs. No null is reported for the median; it is described, not tested.

None of these correlations is a reconstruction. At layer 12, permuting the drugs
of the target and running the identical pipeline still gives a mean correlation
of $0.9220$, with $0.9318$ at the ninety-fifth percentile of that null against
the $0.93225$ observed. Under the same permutation the held-out drug is still
retrieved $0.9108$ of the time on average, with $0.94$ at that percentile
against an observed $0.92$. This comparison's analytic retrieval null of $1/50$
is not the operative one here. What the increment carries is a comparison
between two fits and not the height of either.

\Cref{sec:methods-probe} discounts its own decode positive because the prompt
names the drug in plain text. The name fit measures that discount, and the
activations do not beat it. The repair this invites, a better linear readout,
was tried in its simplest form and did not arrive. That locates one ridge and no
more. It is not scored against the baselines of \cref{sec:res-lookup}, and no
positive control was run on it, so it has never been shown to register a
positive.

\begin{scopelimit}[carried into \cref{sec:res-mechanism} and \cref{sec:limitations}]
Three things bound what that increment carries. The fold is not fully held out
in its centring. Activations and target are both centred over all \num{50} drugs
including the held-out one, and the kernel is rescaled by the mean of its own
full diagonal; the penalty the inner hold-out selects sits at the floor of its
grid in the median fold at both layers, so the fit interpolates its training
drugs and the absolute correlations are inflated by construction. The
permutation null runs the identical pipeline, so the sign is tested against it;
no corrected correlation is estimated here. The increment's own permutation null
is not centred at zero, so a $p$ on it is not a test of whether the activations
add, and none is quoted above. And the stronger reading, that the drug region
was only ever a label, is \notestablished here. Its pre-registered condition
covers a joint result of which this section reports one half, and pairs a
cell-line positive control on the identical code path with a split-half
reliability of at least $0.40$ on the target. The control is positive at both
layers at this position in both cell lines, but with a different statistic
against a different null, never with this ridge. The target's reliability sits
at $0.385$ in the cell line the scores come from, and lower in the second. The
threshold is read on the bootstrap lower bound, which is $0.244$. That
reliability is biased upward: each drug nests within a treated well, and both
halves of the split inherit that well's technical structure. Whether more cells
per drug, or a plating design separating drug from well, would reach the
threshold is untested: the well-split companion covers one drug and returns
nothing. A different fit will not reach it.
\end{scopelimit}

\paragraph{Where ablation stops} It cannot show that the slab's drug-specific
content is what is read. The purified directions could carry generic
strong-programme structure that merely co-varies with drug across the twelve.
That needs a test which replaces what the slab says.

\begin{claim}
The purified drug-label slab carries decodable information and generation is
sensitive to its removal, at a bit player's weight. Once cell line, plate and
dose are projected out, decodability rises monotonically with depth, which the
raw probe, peaking at layer 9, does not, while the descriptive subspace fraction
stays flat at $\approx 0.03$, roughly $5\%$ of the cell line's $\approx 0.6$, and
is not a biological variance decomposition. Ablating the slab costs $9$--$37\%$
of what ablating cell line costs at layers $2$--$9$, with layer $12$ the
exception at $98\%$. Whether the slab's content is read, ablation cannot say.
On a separate roster of \num{50} drugs, a ridge from this model's activations
recovers each drug's measured residual profile less well, on held-out mean
correlation, than the same ridge given only that drug's name.
\end{claim}

% ===========================================================================
\beatsection{E}{3}{Identity is read, at very low gain}{sec:res-mechanism}
% \claimlede must follow \beatsection: \beatsection clears the stored lede.
\claimlede{The generation decoder is sensitive to the drug-label activation
signal at very low gain, in one training arm and at one rung of the displacement
ladder.}

\begin{question}[3]
Removal cannot distinguish a region whose \emph{content} is read from one whose
mere presence matters. The last question needs an instrument that replaces what
the region says. Push the activations from the drug the prompt names towards
another drug, and ask whether the model moves its preference towards that other
drug's own held-out sentence. It moves, by a sliver, in one arm at one depth. The
design is given at length because three earlier versions could have manufactured
an effect.
\end{question}

\noindent
The verdict first, because the instrument takes several pages to describe and the
verdict is small. In one of the five training arms of \cref{tab:trainconfig}, at
one depth of seven, injecting half the class-mean displacement from drug A to
drug B inside
the purified subspace ($\alpha = 0.5$, the rung the multiplicity correction is
applied to) moves the model's preference towards drug B's own held-out sentence,
against an identically constructed injection naming a third drug, by
\ci{+0.0197}{+0.0063}{+0.0335} of the preference gap the model already has. One
to two per cent, against a $10\%$ margin pre-registered for calling an effect
material. Two further seeds reproduce it, and of twenty-six headline tests it is
the only survivor of multiplicity correction.

Read at face value that says the generation decoder does consult which drug the
region names, at a gain far too small to account for anything in
\cref{sec:res-blind}. Read with the multiplicity discipline below it says less.

\begin{scopelimit}[carried into \cref{sec:res-epsilon} and \cref{sec:limitations}]
This is the weakest load-bearing measurement in the chapter and it is carried as
such. No argument downstream rests on it alone, and a reader who discards it
entirely loses nothing that the drug-blindness result depends on. What the
following section tests is a prediction of the \emph{interpretation} of this
result, not the result itself, and the inference there is written to run one way
only for that reason.
\end{scopelimit}

% ===========================================================================
\begin{figure}[htbp]
\centering
% ===========================================================================
% materiality.tex -- fig:materiality. The one surviving identity effect drawn
% twice on one nats ruler: once against the whole preference gap it is divided
% by, where it is a 3.6mm smudge on a 133mm axis, and once at exactly eight
% times, where the swap arm, the permuted control and their difference separate.
%
% FIGURE BODY ONLY. No preamble, no document environment, no float, no caption.
% \input it from inside a figure float in Sections/04b-representation.tex at
% sec:res-mechanism. The float wrapper and caption sit at the foot of this file,
% commented out, ready to paste -- the staging convention of figs/two-scales.tex.
% Compiles against thesis_v2/preamble.tex and nothing else: no \includegraphics,
% no external data at draw time, no TikZ library beyond the ones preamble.tex
% already loads (this file uses decorations.pathreplacing and arrows.meta only).
%
% WHY THIS FIGURE EXISTS. Section 4.5 reports every effect "as a fraction of
% that model's own clean preference gap" (04b:396-398) and NO figure anywhere in
% the document shows that gap. fig:family draws the numerator on a diverging
% colour scale; fig:comparators row C draws it as a fraction. The denominator is
% never drawn, and neither is the subtraction that produces the numerator. This
% figure's subject is the denominator, and the subtraction.
%
% EVERY NUMBER DRAWN, AND WHERE IT COMES FROM
% -------------------------------------------
% ONE arm, ONE layer, ONE rung: the residual arm, layer 12, alpha = 0.5. Source
% for every coordinate: RESULTS_cluster/probe_arm_residual.json, all under the
% key path layers.12.swap. Sibling ledger of the exact values drawn, at full
% machine precision, with key paths: thesis_v2/figs/materiality.json.
%
%   quantity                key path              value drawn (full precision)
%   clean preference gap    .clean_ab_gap         0.06420005919957501
%   10% materiality margin  .equiv_margin         0.006420005919957502
%   permuted arm            .effect_perm          0.0006409306049362111
%   swap arm                .effect_swap          0.0019027651015006252
%   swap - permuted         .diff_perm            0.0012618344965644142
%   its 95% interval        .ci_perm     [0.0004024989353392094,
%                                         0.002150955667015771]
%
% One further quantity is COMPUTED, and is declared as computed in the ledger:
%   one eighth of the gap   .clean_ab_gap / 8     0.008025007399946877
% It is band two's full scale and band one's magnifier bracket, and it is the
% only reason the word "eighth" appears anywhere on the figure.
%
% VERIFIED ON BUILD, not assumed:
%   * .equiv_margin / .clean_ab_gap = 0.1 exactly (config.equiv_frac = 0.1), so
%     the two dashed rules stand at one tenth of band one's full scale.
%   * .effect_swap - .effect_perm = 0.001261834496564414 against .diff_perm
%     = 0.0012618344965644142. They agree to fifteen significant figures and
%     differ in the last bit only, from summation order; the stored .diff_perm
%     is what is drawn. The subtraction the figure draws is therefore the
%     subtraction the pipeline performed, and this was checked, not asserted.
%   * .diff_perm / .clean_ab_gap = 0.019654724813287512 and
%     .ci_perm / .clean_ab_gap = [0.0062694480403512436, 0.03350395145788292],
%     reproducing .diff_perm_norm and .ci_perm_norm to sixteen digits. That is
%     the cross-check that this figure's nats and the thesis's fractions are the
%     same measurement; the normalised triple is printed in the caption and is
%     NOT drawn.
%   * 0.06420005919957501 / 0.0012618344965644142 = 50.87835161772162 and
%     0.006420005919957502 / 0.0012618344965644142 = 5.087835161772163. Both
%     COMPUTED, declared as computed in the ledger; they are the "one fiftieth"
%     and "one fifth" of the closing line.
%
% THE ci_perm TRAP, and why the header names both draws
% -----------------------------------------------------
% The interval drawn is layers.12.swap.ci_perm, the TOP-LEVEL key:
%   [0.0004024989353392094, 0.002150955667015771], p_perm = 0.001,
%   normalising to the thesis's [+0.0063, +0.0335].
% layers.12.swap.ladder[alpha = 0.5].ci_perm stores a SECOND bootstrap draw of
% the same effect,
%   [0.0003982749538768512, 0.002137823621073976], p_perm = 0.004,
% with the point estimate identical to sixteen digits. It normalises to
% [+0.0062, +0.0333] and would print one digit off tab:identity while silently
% swapping the surviving p for the non-surviving one -- 0.001 and 0.004 lie
% either side of the 0.05/26 admitted at rank one (04b:526-531). A builder who
% reads the interval from the ladder because that is where the alphas live makes
% exactly this error. The ladder draw is recorded in materiality.json under
% not_drawn_second_bootstrap_draw and is not drawn here.
%
% AGREEMENT WITH THE TEXT (checked line by line, not assumed):
%   04b:340-341  "\ci{+0.0197}{+0.0063}{+0.0335} of the preference gap the model
%                already has. One to two per cent, against a $10\%$ margin
%                pre-registered for calling an effect material."
%                -> the drawn nats normalise to exactly that triple; the margin
%                   is drawn twice, once per band.                          OK
%   04b:443      "In raw nats the rungs read $+0.0013$" (first rung = the
%                headline) -> row three's label prints $+0.0013$.           OK
%   04b:500-502  "60 rows in 58 pair clusters over 24 anchor drugs"
%                -> row three's second label line prints 58 pair clusters;
%                   60 and 24 are in the caption.                           OK
%   04b:507-513  "The interval must lie inside a pre-registered margin ... and
%                contain zero. ... Under textbook TOST the headline ... would be
%                returned as equivalent, read as a null; this rule returns a
%                directional effect smaller than the margin."
%                -> the two closing ink lines track this sentence, which is the
%                   section's subtlest point and is drawn nowhere else.     OK
%   04b:568      "a fifth or less of the $10\%$ margin"
%                -> the closing accent line prints "one fifth of the margin".OK
%
% DESIGN DECISIONS, each with its reason
% --------------------------------------
%  1. TEXTWIDTH, NOT FULLWIDTH. sec:res-mechanism already carries two fullwidth
%     objects, fig:family and tab:identity, and a fullwidth float forbids a
%     margin note on its page (preamble.tex:522-523). Dropping fullwidth buys
%     this page's margin notes back -- the trade ladder.tex and two-scales.tex
%     both made deliberately. Drawn to 141.9mm against the 142mm measure; the
%     observed band in the corpus is 139.6-141.9mm.
%  2. A RAW-NATS AXIS, WHICH IS LEGITIMATE HERE AND NOWHERE ACROSS ARMS. The
%     five arms emit different target formats and their clean gaps span nearly
%     fivefold, so a single nats rule laid across arms would be one
%     arm's ruler drawn over five arms' data. RECOMPUTED, because the figure
%     brief quoted "4.4x" for this and 4.4 is not the span: over the five
%     probe_arm_*.json files clean_ab_gap runs 0.06420005919957501 (this arm,
%     this layer) to 0.3138025708635225 (consensus, layer 4), a ratio of
%     4.887886, and 4.410308 is the narrower consensus-layer-2 to
%     residual-layer-12 ratio. The caption prints the two endpoint values and
%     the word "nearly fivefold" so a reader can check it. Everything on this
%     figure is one
%     arm at one depth at one rung, so the raw scale is the arm's own and the
%     10% margin is a single rule with a single meaning. fig:family's lower
%     strip is a single-arm object for the same reason.
%  3. THE SAME QUANTITY DRAWN TWICE, NOT TWO QUANTITIES, AND THE MAGNIFICATION
%     IS EXACTLY EIGHT. Band two's full scale is .clean_ab_gap / 8, so the panel
%     is exactly the first eighth of the axis above it and the factor printed on
%     its axis is exactly 8 and not a rounded number. This was arrived at by
%     correction, and the arithmetic is worth recording. The first draft ran
%     band two to 0.0075 nats and called the panel "the first twelfth" at
%     "8.56x": 0.0075 / 0.06420006 = 0.11682, which is one 8.56th and not one
%     twelfth, so the panel's own title was false by a factor of 1.4 and the
%     magnification was a number no reader could check. A true twelfth cannot be
%     drawn at all -- .clean_ab_gap / 12 = 0.005350 nats is SMALLER than the
%     0.006420 margin, so the margin would fall off the right-hand end of the
%     panel and the figure would lose its subject. One eighth is the nearest
%     nameable fraction that still contains the margin, and at that scale the
%     margin sits at 80% of the panel's width, which is where it does the most
%     work. Nothing else changed: every drawn quantity is the same number.
%  4. BAND ONE IS SPARE, BUT ITS MEASUREMENT IS NOW NAMED. Four tick numerals,
%     not seven: the band's job is one impression -- 3.6mm of ink 0.8mm from the
%     origin on a 133mm axis -- and a crowded band invites reading rather than
%     seeing. Every other label lives in band two.
%     REPAIRED, PASS 8, and both reviewers called the old state blocking. The
%     measurement carried NO label at all, on the theory that the magnifier
%     would carry it; its only pointer was the italic "the effect and its
%     interval, magnified below" set 11cm to its right at the far end of a
%     nearly horizontal grey leader, against the document's own rule that a
%     label more than ~0.3cm from its mark gets a leader. One reviewer could not
%     tell whether the mark was a mark, a tick or a printing artefact; the other
%     read it as a stray em-dash. It now carries an accent label, "the effect",
%     centred on the estimate dot 0.22cm above it, and a 0.25pt stem from the
%     dot down to the spine so the mark is anchored to a coordinate rather than
%     floating. The label is two words and not "the effect, +0.0013 nats"
%     because the measure between the axis origin at 0.55 and the margin rule at
%     1.88 holds 1.33cm at \scriptsize while the longer form is 2.55cm and would
%     have crossed the rule; row three of band two prints the value in full,
%     3.5cm below. The italic pointer is deleted with the leader it rode on.
%  5. THE MEASUREMENT SITS 0.22cm ABOVE BAND ONE'S SPINE, not on it. On the
%     spine, a 3.6mm accent bar 0.8mm from the origin reads as a thickening of
%     the axis or as a printing defect. Lifted clear, with its dot over a 2.1pt
%     white halo, it reads as a mark. Checked in the 200dpi render and again in
%     a 550dpi crop: the bar, the dot and the spine are three separable things.
%  6. THE MAGNIFIER'S SOURCE IS A NAMED SPAN, NOT TWO BARE COORDINATES.
%     REPAIRED, PASS 8, reported independently by both reviewers. It used to be
%     a vertical drop at x = 0.55, a second vertical at x = 2.2125, and a line
%     from (2.2125,7.20) to (13.85,6.75) -- 11.6cm horizontal over 0.45cm
%     vertical, about 2 degrees. Read cold, the vertical was a drop from the
%     axis origin and the long line was a rule under band one, and the italic
%     that rode on it annotated band one's axis rather than announcing a
%     magnification. slab.tex's version works because both its leaders leave the
%     magnified REGION at a visible angle, and the fix is to make the region an
%     object first. A rulegrey brace, the corpus's own idiom, now spans exactly
%     the magnified eighth on band one -- \ga{0} = 0.55 to
%     \ga{clean_ab_gap/8} = 2.2125 -- and both leaders leave its two ends, one
%     at 66 degrees and one shallow. THE SHALLOW ONE CANNOT BE FIXED and should
%     not be pretended away: a 1.66cm span is being opened into a 13.90cm panel,
%     so one of the two leaders must run nearly flat whatever is done. What
%     changed is that it now leaves a bracket rather than a bare coordinate, so
%     the assembly reads as source-and-enlargement rather than as a crooked
%     rule. The leaders now land on the panel's own top corners (0.25 and
%     14.15) rather than on band two's data extremes: with the brace naming the
%     source, the leaders' job is to point at the PANEL, and corner-to-corner is
%     what a magnifier looks like. The bracket contains both the effect AND the
%     margin rule, which is itself the point: the whole of section 4.5's
%     argument happens inside the first eighth of the gap.
%  7. THREE ROWS IN THE ORDER PERMUTED, SWAP, DIFFERENCE. Control first,
%     treatment second, measurement third, so the subtraction reads downward in
%     the order the sentence "swap minus permuted" is spoken.
%  8. AN INTERVAL IS DRAWN ON ROW THREE AND ONLY THERE. No interval is stored
%     for either arm alone -- .ci_perm is an interval on the DIFFERENCE -- so
%     rows one and two are bare solid dots. In this document a bare solid dot
%     with no bar is how the absence of an interval is said (comparators.tex
%     changed exactly that mark from a hollow ring to a solid fill, because the
%     absent bar is what carries the meaning). Row three's bar is the document's
%     own convention: a 95% cluster bootstrap over ordered (A,B) pairs.
%  9. THE SIX-MARK TRANSLATION DEVICE IS DELETED, AND THE ROWS CARRY THE
%     CORPUS'S ROW APPARATUS INSTEAD.
%     PASS 8, and the spec pre-authorised this deletion. The device was a
%     hairline from each of the two dots down to the ends of a brace, that
%     brace, a leftward arrow sliding it to the origin, a second brace of equal
%     length from the origin, and a hairline dropping onto row three's dot: five
%     grey structural marks against three measured ones, carrying the identity
%     swap - permuted = the interval's point estimate. It is arithmetically
%     exact -- \gb{effect_swap} - \gb{effect_perm} = \gb{diff_perm} - 0.55 =
%     2.09126 -- and it did not read. Neither brace was labelled with a value,
%     so 0.00190 - 0.00064 = 0.00126 was not recoverable from the marks; the
%     arrow's head landed in white space with no mark at it; and the one
%     sentence that explained the chain, "the same span, measured from zero",
%     sat about 4cm to the right of the arrow it named, at a height between the
%     two braces, attached to nothing. Both reviewers reported having to
%     magnify 3x and reverse-engineer the arithmetic, and one noted that nothing
%     in the construction was accent, so the eye never entered it. Row three's
%     own label already STATES the identity in words -- "swap - permuted" -- so
%     the deletion costs nothing and removes five hairlines. The house rule it
%     follows: an identity must be readable from the marks or not drawn.
%     What replaces it is what fig:comparators gives every row: a rulegrey!55
%     0.3pt divider under each of the three rows, spanning the panel's inner
%     width. Before, the three quantities sat at three heights with no dividers,
%     no stubs and no alignment cues, so the vertical read as a second dimension
%     when it is only a stacking order -- and the higher dot happens to be the
%     SMALLER value, which invites a reader to look for an ordering that does
%     not exist. Rows moved from 5.52/4.94/3.66 to 5.40/4.72/4.04 and the
%     margin brace from 3.10 to 3.30 to take up the slack the deletion left.
% 10. THE MARGIN IS DRAWN TWICE AND BRACED ONCE. On band one it stands 1.33cm
%     from the origin; on band two, 10.64cm from it, and the 7.08cm of clear
%     space between the interval's upper end and the margin rule is the figure's
%     answer to "is it material". The brace under band two names what the margin
%     is a tenth OF -- 0.00642 nats out of 0.0642 -- which no existing figure
%     does.
% 11. THE FRAME IS NAMED ON THE DRAWING and not only in the caption: a quietink
%     line at the foot says which arm, which layer and which rung. This document
%     runs five arms whose clean preference gaps differ by 4.89x, so a raw-nats
%     axis that does not say whose nats they are is under-labelled in exactly
%     the way gate.tex's fifth fix names. \ttfamily on \texttt{residual} because
%     that is how the results files name the arm (ladder.tex's register key).
% 12. WHAT IS NOT DRAWN. The isotropic arm (.effect_iso = 0.0010863731109609438)
%     is stored at this cell and is omitted: 04b:377-380 keeps it "only as a
%     damage diagnostic", it is not the control the headline is measured
%     against, and a third bare dot between rows one and two would read as a
%     second comparison. The alpha ladder is omitted because fig:family's lower
%     strip already draws it and this figure is one rung. p, q, n and the
%     normalised triple are in the caption, not on the drawing.
% 13. REJECTED ALTERNATIVE: a broken axis. Drawing the whole gap and the first
%     eighth on one axis with a break would have saved 3.5cm of height, and it
%     would have destroyed the one thing the figure is for -- the reader has to
%     see the unbroken 133mm and the 3.6mm on it in the same glance before the
%     magnification means anything. two-scales.tex breaks an axis and says in
%     its own header that the break is admissible only because the off-scale
%     value is named in words; here the off-scale value IS the subject.
% 14. REJECTED ALTERNATIVE: bars from zero for the two arm effects. A bar length
%     would have made effect_swap look three times effect_perm, which is true of
%     the numbers and is not the measurement: the measurement is their
%     difference with an interval, and neither arm alone is quoted anywhere in
%     the thesis. Dots keep the arms as positions and the difference as the only
%     thing given a bar.
% 15. REJECTED ALTERNATIVE: naming the 0.0642 axis "chance" or drawing a 0.50
%     rule. There is no chance value on this figure. The reference here is the
%     model's own clean preference gap, an arm-specific measured quantity, so it
%     is drawn as the axis itself and named at the axis's right end rather than
%     as a rule inside it. Noted because the document's two conflicting habits
%     for a chance rule -- solid black in style_v2.chance_line(), 0.5pt densely
%     dotted ink in ladder.tex -- do not apply here, and their absence should be
%     deliberate rather than accidental.
%
% LAYOUT NOTES, all forced by a build
% -----------------------------------
%   * TWO MAPS, both braced pgfmath groups. A shim may not be added outside the
%     braces -- \gb{v}+0.05 is not a coordinate -- so every label offset here is
%     an explicit literal or an xshift on the node (the gate.tex trap). Both
%     maps are used directly as coordinate components, (\gb{v},y), which is
%     comparators.tex's idiom and keeps every data x exact rather than rounded
%     into the source.
%       \ga{0.06420005919957501}  = 13.850005   (band one, the whole gap)
%       \gb{0.008025007399946877} = 13.850005   (band two, one eighth of it)
%     The two right-hand ends coincide to 5 decimal places by construction:
%     1657.32 = 8 x 207.1650 exactly.
%   * Tick numerals use the \foreach \v/\t value/label pair form. \foreach \v
%     alone prints the raw token, so 0.020 would set as 0.02 in one place and
%     0.0 in another; the pair form pins the printed string.
%   * inner sep = 1pt ON THE WHOLE PICTURE, as gate.tex sets it. The default
%     0.3333em is 3.65pt at the ambient \small, and at pass 2 that invisible
%     padding -- not the ink -- was what collided: the band-one word tag's box
%     and the margin label's box overlapped while their glyphs were 1.2mm
%     apart. With inner sep pinned, every clearance below is box arithmetic
%     rather than judgement. The consequence is that a label placed at
%     mark + 0.30 would sit 0.10cm closer to its mark than comparators.tex's,
%     so every row label here is placed at mark + 0.40, which puts the glyph
%     0.39cm clear of the marker edge -- comparators' measured value.
%   * TICK NUMERALS CARRY fill=white AND ARE DRAWN LAST IN THEIR BAND. Pass 1
%     ran the magnifier's left leader straight through band one's "0", because
%     band one's origin and the magnifier's left edge are the same x = 0.55.
%     The white box interrupts the rule across the numeral instead --
%     comparators.tex's idiom for its zero rule, done with a box rather than by
%     hand because the numerals are drawn in a \foreach. Drawing order matters:
%     spine, then magnifier, then numerals.
%   * THE PANEL IS 3mm WIDER THAN ITS DATA ON EACH SIDE, (0.25,1.50) to
%     (14.15,6.75) around data running 0.55 to 13.85. At pass 2 the panel ran
%     exactly 0.55 to 13.85 and its left dashed border was drawn at the same x
%     as band two's dotted zero rule; the two were indistinguishable in the
%     600dpi crop, so the figure's zero coordinate was reading as the frame of
%     its own panel. Widening the frame separates them by 3mm. PASS 8 re-aimed
%     the magnifier's leaders at the panel's own top corners (0.25 and 14.15),
%     which is the opposite of what pass 2 did and is right for the same reason
%     the brace is right: with the brace naming the SOURCE span exactly, the
%     leaders' job is to point at the PANEL, and corner-to-corner is what a
%     magnifier looks like.
%   * BAND TWO'S WORD TAG AND ITS MARGIN LABEL SIT INSIDE THE PANEL, on one
%     baseline at y = 6.44, and the axis title sits inside it too. At pass 2 the
%     tag, the margin label, "no movement" and the axis title all straddled the
%     dashed border. Anything that names something inside the panel is now
%     inside the panel.
%   * The band-two tag is anchored at x = 0.75 and not at x = 0.00, because the
%     magnifier's left leader is a vertical at x = 0.55.
%   * PASS 8: the translation device and its label are deleted (decision 9),
%     which took five hairlines and one floating sentence out of the panel.
%     Rows moved up to take the slack -- 5.52/4.94/3.66 became 5.40/4.72/4.04,
%     and the margin brace 3.10 became 3.30 -- and each row gained a rulegrey!55
%     0.3pt divider under it, spanning 0.35 to 14.05, which is 0.10cm inside the
%     panel's dashed frame at both ends.
%   * PASS 8, THE MARGIN LABEL STAYS ON THE 8.44 ROW. The build that moved it up
%     to the tag's row put it at x = 1.96 and "THE WHOLE GAP" measures 1.95cm of
%     \sffamily\scriptsize\bfseries from x = 0, so the two ran together as
%     "THE WHOLE GAP0% materiality margin" in the render. Both labels sit on the
%     8.44 row instead and the dashed margin rule at x = 1.88 separates them:
%     "the effect" is centred on the estimate dot and ends at 1.29, the rule is
%     at 1.88, the margin label starts at 1.98. The rule stops at 8.56, which is
%     0.09cm under the tag's descender line and is why it may not be carried
%     higher to meet a label on the tag's row.
%   * BAND ONE'S RIGHT END is a terminal tick RISING from the spine, 0.30cm
%     tall against the 0.10cm hanging axis ticks, with its two label lines
%     right-aligned on the same x = 13.85 so the tick reads as their anchor. At
%     pass 3 that rule ran from the spine up to the second label's baseline and
%     read as a bracket enclosing the label block.
%   * The closing prose was three lines at pass 2 and is two at pass 3; the
%     28pt that bought went into the clearances above.
%   * Every prose node is a single \scriptsize line anchored base west. No
%     `align=` and no `text width=` anywhere: those build a minipage whose first
%     line follows the AMBIENT \baselineskip, so the node's ink moves with the
%     prose before the float (frames.tex measured a 3.7pt drift).
%   * The bounding box is pinned with a bare \path so the measure is exact and
%     nothing is clipped: (0.0000,0.30) to (14.1900,8.98) = 141.9 x 86.8mm.
%
% TYPE. Ambient \small on the picture; every node sets \scriptsize = 8pt, which
% is 0.67 of the 12pt body and below the 11pt caption. \ttfamily on the two arm
% names and on the arm named at the foot, because that is how the results files
% name them. Nothing is scaled: 1cm here is 1cm on the page.
% ===========================================================================
%
% ===========================================================================
% PASS 9 -- 2026-08-18. REBUILD AGAINST THE STYLE CONTRACT AND THE SUPERVISOR'S
% VERDICT ("no definition on the x axis for what we are measuring"; "make it
% look more professional, right now it looks a bit sloppy and childish").
% NOTHING ABOVE IS DELETED. The design rationale of passes 1-8 is the audit
% trail and stays; what follows supersedes, EXPLICITLY AND BY NAME, each clause
% of it that this pass made false. The reference standard adopted is fig 4.11
% (figs/fig-mechanism.pdf, printed page 51) and, for the threshold, fig 4.13
% (figs/family.pdf), both of which were rendered and read side by side with this
% figure during the build.
%
% NO NUMERIC VALUE CHANGED IN THIS PASS. Every quantity drawn or printed is the
% same float from the same key path as in pass 8: clean_ab_gap 0.06420005919957501,
% equiv_margin 0.006420005919957502, effect_perm 0.0006409306049362111,
% effect_swap 0.0019027651015006252, diff_perm 0.0012618344965644142,
% ci_perm [0.0004024989353392094, 0.002150955667015771], n_pair_clusters 58.
% The one printed numeral that left the drawing is "one fiftieth / one fifth"
% (a ratio in words), and it went to the caption, where it already was.
%
% W1. SUPERSEDES LAYOUT NOTE "TWO MAPS" AND THE MAP CONSTANTS RECORDED THERE.
%     Both maps lose their 0.55cm left inset and gain a longer barrel:
%         old   \ga#1 = {0.55+(#1)*207.1650}   \gb#1 = {0.55+(#1)*1657.3200}
%         new   \ga#1 = {(#1)*215.7500}        \gb#1 = {(#1)*1726.0000}
%     Two things are preserved exactly and were recomputed on the build rather
%     than assumed: 1726.0000 = 8 x 215.7500 EXACTLY, so the magnification is
%     still exactly eight and not a rounded number (this is the correction made
%     at pass 4 and it survives untouched); and both bands' right-hand ends
%     still coincide, now at \ga{0.06420005919957501} = \gb{0.008025007399946877}
%     = 13.851162772308308cm. Every data x is still computed inside pgfmath from
%     the full-precision float, so no value is rounded into the source.
%     WHY: (i) S13, one x origin per panel. With the old map the axis began at
%     0.55 while every text block began at 0.00, so the panel had two left edges
%     0.55cm apart and the left column measured three (58.0, 63.9, 70.4pt).
%     Now the axis's left end, the panel titles, both axis-label blocks, the
%     span bracket and the foot line all sit on x = 0.0000. (ii) S5 puts the
%     axis-label line under the tick numerals and LEFT-ALIGNED TO THE AXIS'S
%     LEFT END; that is only the same thing as the panel origin if the two
%     coincide. (iii) S14, measure. Deleting the dashed enclosure (W4) took away
%     the object that used to set the width, and a 133.0mm barrel inside a
%     textwidth figure measures about 134mm of content, under the 138mm floor.
%     215.75 puts the barrel at 138.51mm and the content bbox at 139.6mm, inside
%     138-142, with the widest object the axis itself and the data region 100%
%     of the drawn width.
%     The old map is kept in materiality.json beside the new one, per S25.
%
% W2. AXIS DEFINITION, WHICH IS THE SUPERVISOR'S FIRST COMPLAINT AND THE REASON
%     THIS PASS EXISTS. Pass 8 named neither axis. Band one carried bare 0 /
%     0.02 / 0.04 / 0.06 with no quantity and no unit anywhere on it -- the unit
%     reached the reader only through the terminal label -- and band two carried
%     "nats, magnified 8x from the axis above", which names a unit and a
%     transformation but no quantity. 4.11 sets "probe accuracy", "descriptive
%     subspace fraction", "KL from the clean forward pass" and "layer" in
%     axis-label position five times over. Each band now carries, directly under
%     its tick numerals and left-aligned to its axis's left end:
%         line 1, \footnotesize roman ink:  Preference shift, nats
%         line 2, \scriptsize quietink:     (a) residual arm, layer 12, alpha =
%                                               0.5; the full width of this rule
%                                               is the model's clean preference
%                                               gap
%                                           (b) the first eighth of the axis
%                                               above, magnified 8x
%     The quantity is repeated deliberately: two scales, each self-describing.
%     "Preference shift" is the section-4.5 registry name (S19), so the axis, the
%     caption and fig:family's lower strip now use one string for one quantity.
%
% W3. THE THRESHOLD. SUPERSEDES DESIGN DECISION 10 AND THE marginrule STYLE.
%     The 10% pre-registered materiality margin was accent 0.60pt dashed here,
%     ink 0.50pt densely dotted in fig:comparators row C, and solid black 0.90pt
%     in fig:family's lower strip: one pre-registered threshold, three renderings
%     inside nine pages. It is now BLACK, SOLID, 0.90pt in both bands, labelled
%     \scriptsize roman black, which is style_v2.chance_line()'s own hard-coded
%     register and fig:family's. This is the highest-priority change in the set
%     and it is a correctness repair, not a preference: accent means the object
%     under test (S8), and a pre-registered bound is not under test.
%     Design decision 15's note -- that the document's two habits for a chance
%     rule "do not apply here" -- is now false as written and is superseded:
%     there is still no chance value on this figure, but the black-solid-0.9pt
%     register does apply, to the margin.
%
% W4. DELETED: THE DASHED ROUNDED ENCLOSURE round band two (390 x 149pt), with
%     the layout note that sized it 3mm wider than its data on each side. S12
%     forbids a dashed rectangle drawn round a panel outright, and it was this
%     figure's only rectangle. Band two now sits directly under band one in
%     4.13's idiom -- the magnified axis under the full axis, with the
%     magnification named in axis-label position -- which is also what fixes
%     band two's axis under W2. Two consequences recorded so they are not read
%     as accidents: the layout note about the panel's left border colliding with
%     band two's zero rule is moot, since there is no border; and the note that
%     band two's tag and margin label "sit inside the panel" is moot for the same
%     reason.
%
% W5. DELETED: THE FOUR PALE ROW SEPARATORS AND THE MARGIN BRACE. SUPERSEDES
%     DESIGN DECISION 9's replacement clause and the pass-8 note that added them.
%     Five near-identical thin horizontals stood inside 30mm of vertical space --
%     three rulegrey!55 row dividers, one rulegrey brace carrying "10% of the gap
%     = 0.00642 nats", plus the axis rule -- and the one that carried meaning was
%     indistinguishable from the ones that did not. The dividers go because the
%     rows no longer need them: with the enclosure gone and one x origin, the
%     three rows read as a stack. The brace goes because the same fact is now
%     carried without any stroke at all, by printing the margin's value in band
%     one's threshold label ("10% materiality margin, 0.00642 nats") beside band
%     one's terminal label ("the clean preference gap, 0.0642 nats"): the tenth
%     is then a division a reader performs on two printed numbers 8cm apart
%     rather than a brace they must trace.
%
% W6. THE INTERVAL IS DRAWN ONCE. SUPERSEDES DESIGN DECISION 4 and the pass-8
%     repair that labelled band one's bar. Pass 8 drew the same estimate as a
%     3.6mm bar-with-dot in band one AND as a 30mm bar-with-dot in band two,
%     while the caption said "it is the one bar on the figure"; a naive reader
%     counted two and spent time working out which two quantities they were.
%     Rule applied: an interval is drawn at the scale where it is legible and
%     nowhere else. Band one keeps a single filled accent disc for the estimate,
%     with no bar; band two carries the bar. Band one's dot keeps its label and
%     its leader, and the label is now the SAME STRING as band two's row three,
%     "swap - permuted", because S19 forbids two names for one JSON key and pass
%     8 had "the effect" here and "swap - permuted" there.
%
% W7. DELETED: THE TWO BANNERS, replaced by 4.11's panel heads. "THE WHOLE GAP"
%     and "THE FIRST EIGHTH, MAGNIFIED" were \sffamily\scriptsize\bfseries accent
%     caps -- the preamble's margin-apparatus face, which belongs to the margin
%     and not to a body figure (S2). They are now "(a) the whole preference gap"
%     and "(b) the first eighth, magnified 8x", \footnotesize roman ink,
%     lowercase, left-aligned on the panel origin, no colon and no bold. The
%     tag/.style is deleted rather than left unused.
%
% W8. DELETED: THE ITALIC "no movement". It sat unattached above the permuted
%     row and a naive reader took it for an assertion about that arm rather than
%     a name for the zero rule. The zero rule survives -- densely dotted quietink
%     0.5pt, which is exactly S7's reserved register for a coordinate that
%     decides nothing -- and carries no label, because its foot is the axis's own
%     "0" numeral. With it goes the last italic on the drawing: italic share is
%     now 0%, against 10% at pass 8 and 0% in 4.11.
%
% W9. DELETED: THE CLOSING PROSE BLOCK AND THE ACCENT APHORISM. The two ink lines
%     ("The interval excludes zero and lies inside the margin. Textbook TOST
%     returns that as equivalence and reads it as a null; this rule returns a
%     directional effect smaller than the margin.") and the accent italic ("one
%     fiftieth of the gap; one fifth of the margin.") are argument, and in this
%     document the argument lives in the caption (S3, S4, S18). At print the
%     reader met an 8pt caption, a blank, then the real 10.9pt caption -- two
%     captions, the smaller one first. Both lines are written into
%     figs/_PENDING_CAPTIONS.tex, keyed by \label, for a human to integrate;
%     the aphorism is already the caption's lede. The TOST asymmetry is the
%     section's subtlest point and it is NOT lost: it is caption material, and
%     the configuration it describes -- an interval that excludes zero and lies
%     wholly inside the margin -- is still visible in band two, which is the only
%     part of that sentence a drawing can carry.
%     SUPERSEDES the text_agreement_check on 04b:507-513 in the ledger, which
%     asserted the drawing restates that sentence. It no longer does.
%
% W10. THE MAGNIFIER. SUPERSEDES DESIGN DECISION 6 AND THE PASS-8 BRACE-AND-TWO-
%     LEADERS CONSTRUCTION. What is kept: the source span is a NAMED OBJECT and
%     not two bare coordinates. What is deleted: both leaders, including the
%     shallow one that pass 8's own header admitted "CANNOT BE FIXED" and that
%     the naive reader read as a trend line before reading it as a leader.
%     The reasoning, recorded because a later builder will be tempted to put them
%     back. A leader pair from the ends of band one's [0, gap/8] segment to the
%     ends of band two's rule is geometrically forced to be asymmetric -- a
%     17.3mm span opening into a 138.5mm rule -- so the right-hand leader runs at
%     under 4 degrees whatever is done, and 4 degrees of accent-free hairline
%     running the full measure is a trend line to every eye that meets it first.
%     Worse, it cannot be routed: S5 now requires band one's two-line axis-label
%     block directly under band one's tick numerals, which puts full-width text
%     in the only lane the leaders could cross. The magnification is instead
%     carried by two objects that cannot be misread: a 0.30pt rulegrey bracket on
%     band one spanning exactly \ga{0} = 0.0000 to \ga{clean_ab_gap/8} = 1.731395,
%     labelled at its own end ("the eighth magnified below"), and band two's
%     axis-label line 2, which names the magnification in axis-label position
%     where S5 puts it. The bracket's two ends stand at the same x as band two's
%     rule origin, so the vertical alignment does the joining that a hairline was
%     doing badly.
%
% W11. MARK VOCABULARY (S9, CF5). SUPERSEDES DESIGN DECISION 8's claim that a
%     bare SOLID dot is how this document says "no interval". It is not the only
%     thing a dot has to say. The permuted arm is the CONTROL, and 4.11's own
%     encoding for a comparator or a null is an open disc, white fill, 0.8pt edge
%     in its colour (mfc='white' on the raw probe and on the random null). Row one
%     is therefore an open quietink disc; rows two and three and band one's
%     estimate are filled accent discs, all at radius 1.4pt, one radius on the
%     figure. The absence of an interval is now said where the contract puts it,
%     in the single \scriptsize interval declaration at the foot (S17), which
%     names the bootstrap construction for the panel that has a bar and says in
%     as many words that the two arms and band one carry none. The 0.8pt disc
%     edge is a MARK edge, not a rule weight; the figure's rules use exactly five
%     weights, 0.30 / 0.40 / 0.50 / 0.90 / 1.00pt.
%     TWO SINGLETONS ARE DECLARED HERE RATHER THAN REMOVED, as S10 permits. The
%     0.50pt densely dotted stroke is used once, for band two's zero rule, and is
%     a distinct kind of object -- S7 reserves that stroke for a coordinate that
%     decides nothing, and this figure has exactly one such coordinate. The
%     1.00pt stroke is used once, for the interval bar, and is a distinct kind of
%     object for the same reason: this figure stores exactly one interval.
%
% W12. TYPE (S1, S16). Two sizes above 7pt and one face: \footnotesize (10.0pt)
%     for the two panel heads, both axis-label line ones, and nothing else;
%     \scriptsize (8.0pt) for tick numerals, direct labels, scope lines, the
%     threshold labels and the interval declaration. No \sffamily, no \bfseries,
%     no \itshape anywhere. \ttfamily survives on the four arm tokens
%     (residual, permuted, swap) because that is how the results files name them
%     and S1 exempts the \texttt token name. Pass 8 printed four sizes and set
%     its largest type on a grey italic epigram; the largest type now belongs to
%     a panel head or an axis label, which is the whole of CF6.
%
% W13. TICK LANES MADE CONGRUENT (S13). Band one ticks at 0 / 0.02 / 0.04 / 0.06
%     land on x = 0.0000 / 4.3150 / 8.6300 / 12.9450. Band two's ticks were
%     chosen so that they land on THE SAME FOUR x: 0 / 0.0025 / 0.005 / 0.0075,
%     which is possible only because band two is exactly one eighth of band one
%     and is a visible consequence of the exact 8. Both rules share both
%     endpoints (0.0000 and 13.851163) and both first ticks sit on the left
%     endpoint. Band one alone carries a terminal tick at the right endpoint,
%     rising 0.20cm from the spine, because that endpoint is the clean preference
%     gap and is the figure's denominator; band two's right endpoint is one
%     eighth of it and is named by band two's axis label rather than by a mark,
%     so it takes no tick and no numeral -- no unlabelled mark is left on the
%     drawing (S20).
%
% W14. MEASURE AND HEIGHT, MEASURED ON THE FINAL BUILD, not estimated: content
%     bbox 140.34 x 83.82mm, ink 5.14% of pixels below 230 grey at 200dpi, 4.88
%     characters of figure text per 100mm^2, one page, zero Overfull and zero
%     Underfull boxes. Against pass 8's 141.2 x 85.9mm at 7.03% ink and against
%     the contract's 138-142mm measure, 85mm height cap, 6.0% ink ceiling and
%     5.0 characters/100mm^2 ceiling. Every millimetre of the height and every
%     point of the ink came from W4, W5, W8, W9 and W10 -- deletions -- and
%     nothing measured was removed to get them.
%
% W15. STALE CLAIMS ELSEWHERE IN THIS HEADER, SUPERSEDED BY NAME so that an
%     auditor reading top to bottom is not misled.
%     (i) The opening paragraph's "where it is a 3.6mm smudge on a 133mm axis"
%         is now false twice over: the axis is 139.96mm, and band (a) no longer
%         draws the interval at all, so what stands 2.75mm from the origin is a
%         single 1.4pt disc and not a 3.6mm bar (W6).
%     (ii) The file-header line "this file uses decorations.pathreplacing and
%         arrows.meta only" is now false in the harmless direction: this figure
%         uses NEITHER. The braces went with W5 and W10 and there was never an
%         arrowhead. It still loads no library preamble.tex does not.
%     (iii) DESIGN DECISION 1's "Drawn to 141.9mm against the 142mm measure" is
%         superseded by W14's 140.34mm. The decision itself -- textwidth and not
%         fullwidth, to buy the page's margin notes back -- stands unchanged.
%     (iv) DESIGN DECISION 3's arithmetic is intact and was recomputed on this
%         build; only the cm-per-nat constants moved (W1). The first draft's
%         "first twelfth at 8.56x" correction, and the reason a true twelfth
%         cannot be drawn at all, remain exactly as recorded.
%     (v) The AGREEMENT WITH THE TEXT entry for 04b:507-513 ("the two closing
%         ink lines track this sentence") is superseded by W9: those lines are
%         deleted from the drawing and staged for the caption. The entry for
%         04b:568 ("the closing accent line prints one fifth of the margin") is
%         superseded for the same reason. Both sentences remain TRUE of the
%         figure's subject; they are no longer drawn on it.
%
% W16. TWO DEVIATIONS FROM THE SPEC'S LITERAL AXIS-LABEL STRINGS, both forced by
%     other clauses of the same contract, both recorded rather than quietly
%     taken.
%     (i) Band (a)'s scope line is set as "residual arm, layer 12, alpha = 0.5"
%         and NOT as "residual arm, layer 12, alpha = 0.5; the full width of
%         this rule is the model's clean preference gap". The deleted clause
%         restates, in twelve words, the terminal label "the clean preference
%         gap, 0.0642 nats" that the spec's own keep-list requires at the axis's
%         right end, and S18 forbids a drawing from carrying one clause twice.
%         It also costs 70 characters, and at 70 more characters the figure
%         needs 89.8mm of height to stay under the 5.0 characters/100mm^2
%         ceiling, which is past the 85mm cap. The quantity and its unit are on
%         line one, which is what S5 actually requires; what came off is a
%         second statement of what the axis's right end already says.
%     (ii) Band (b)'s scope line is set verbatim as the spec gives it, "the
%         first eighth of the axis above, magnified 8x", and the 14 characters
%         that costs were bought by opening 3mm more air between the bands
%         rather than by cutting anything else.
%     Line one of both blocks is verbatim: "Preference shift, nats". Recorded
%     because 4.11 and 4.13 both set their axis labels entirely lowercase and
%     4.10 sets its capitalised, so the slate is split on case and the spec's
%     capitalisation was followed rather than guessed at.
% ===========================================================================
% ===========================================================================
% PASS 10 -- 2026-08-18. REPAIR PASS on the pass-9 rebuild. Three adversarial
% readers went over the built figure independently and returned thirty-three
% findings. NOTHING ABOVE IS DELETED: passes 1-9 are the audit trail and stay
% as written. Every clause of pass 9 that this pass made false, and every
% number pass 9 recorded that this file has never drawn, is superseded here BY
% NAME (S25). The findings that were OVERRULED are recorded under X16 with the
% reason, so that a later builder does not re-apply them blind.
%
% NO NUMERIC VALUE CHANGED IN THIS PASS EITHER. Every quantity drawn or printed
% is the same float from the same key path as at passes 8 and 9:
% clean_ab_gap 0.06420005919957501, equiv_margin 0.006420005919957502,
% effect_perm 0.0006409306049362111, effect_swap 0.0019027651015006252,
% diff_perm 0.0012618344965644142, ci_perm [0.0004024989353392094,
% 0.002150955667015771], n_pair_clusters 58, and the exact magnification 8.
% Both maps are untouched: \ga = (#1)*218.0000 and \gb = (#1)*1744.0000, with
% 1744.0000 = 8 x 218.0000 exactly. What moved is where labels sit, where two
% rules start, how long two others are, and what four sentences say.
%
% X1. STALE CONSTANTS IN THE PASS-9 BLOCK, CORRECTED BY NAME. W1, W10, W13 and
%     W14 record a map, a shared endpoint, a bracket end, four tick positions
%     and a measure that this file has never drawn. materiality.json has always
%     recorded the drawn ones correctly, and the drawing is not affected --
%     every data x is computed inside pgfmath from the full-precision float --
%     but S25 makes the header an audited artefact, so an auditor recomputing a
%     position from W1 must not be handed a figure that does not exist. The
%     corrections, each superseding the clause named:
%       W1  "new \ga#1 = {(#1)*215.7500}  \gb#1 = {(#1)*1726.0000}"
%           -> the drawn map is \ga#1 = {(#1)*218.0000} and
%              \gb#1 = {(#1)*1744.0000}. Both properties W1 claims survive and
%              were rechecked on this build: 1744.0000 = 8 x 218.0000 exactly
%              (8*218.0 == 1744.0 is True in IEEE double), and both bands' right
%              ends coincide.
%       W1  "= 13.851162772308308cm" for the shared right end
%           -> 13.995612905507352cm = 0.06420005919957501 x 218, which is what
%              is hard-coded at the spine, at the two end caps and at the
%              terminal label, and what materiality.json records.
%       W1  "215.75 puts the barrel at 138.51mm and the content bbox at 139.6mm"
%           -> the barrel is 139.956mm and the content bbox 139.96mm, measured
%              on this build. Both are inside S14's 138-142mm.
%       W10 "\ga{clean_ab_gap/8} = 1.731395" for the bracket's right end
%           -> 1.749451613188419cm = 0.008025007399946877 x 218.
%       W13 "land on x = 0.0000 / 4.3150 / 8.6300 / 12.9450"
%           -> 0.0000 / 4.3600 / 8.7200 / 13.0800, for both bands, which is
%              what makes the two tick lanes congruent.
%       W13 "Both rules share both endpoints (0.0000 and 13.851163)"
%           -> (0.0000 and 13.995613).
%       W14 "content bbox 140.34 x 83.82mm, ink 5.14%, 4.88 characters per
%           100mm^2" -> superseded by X17's measurement of THIS build.
%     No 215.75 build was ever produced; the pass-9 numbers are a draft's, not
%     a superseded state of the drawing, and they are recorded here as wrong
%     rather than as an earlier truth.
%
% X2. THE ONE INTERVAL DECLARATION SAID SOMETHING FALSE ABOUT THE INSTRUMENT,
%     AND SAID IT ABOUT THE ONE MARK THAT HAS AN INTERVAL. SUPERSEDES W11's
%     description of the declaration and the ledger's intervals.declaration_on_
%     the_drawing. Pass 9's second line read "No interval is returned for
%     permuted, for swap, or for (a)." Band (a)'s single disc IS diff_perm, and
%     layers.12.swap.ci_perm is an interval on diff_perm; it is the very
%     interval drawn 8x larger in (b). "Returned" is this document's word for
%     what the instrument stores (S17), so the sentence asserted that the
%     band-(a) estimate has no uncertainty available and implied that (a) and
%     (b) draw different quantities. The declaration now separates
%     non-existence from non-drawing:
%       "95% bootstrap over 58 ordered pair clusters, on the difference; drawn
%        once, in (b)."
%       "No interval is stored for permuted or for swap alone."
%     Both lines are true of the file, they still fit S17's one block of at
%     most two lines, and the reason band (a) carries no bar is now the same
%     reason the ledger has always given -- legibility -- rather than a claim
%     about the instrument.
%
% X3. THE THRESHOLD IN (a) STANDS ON ITS SPINE. SUPERSEDES the pass-9 y extent
%     8.38 .. 9.04. Drawn from 8.38 it crossed the axis rule and hung 0.16cm
%     below it, past the ends of the hanging ticks and into the tick-numeral
%     lane, so the heaviest stroke in the panel doubled as an outsized
%     unnumbered tick standing between the numerals "0" and "0.02" -- exactly
%     the "the threshold merges with the tick furniture" defect S7 indicts in
%     fig:comparators row (a). It now runs 8.54 .. 9.04, from the spine upward,
%     which is what (b)'s half of the same threshold has always done. One
%     pre-registered threshold, one stroke AND one geometry, in both bands.
%
% X4. BAND (a)'s ONE MARK IS NOW LABELLED AT ITSELF. SUPERSEDES W6's placement.
%     Pass 9 anchored "swap - permuted" base west on x = 0.00, the panel
%     origin, while the leader fell from \ga{diff_perm} = 0.2751 -- so the
%     leader descended from inside the word "swap", and the mark it landed on
%     read as the swap arm, whose value is 0.0019 and not the 0.0013 drawn.
%     That is the only mark in band (a). The label now begins on the mark's own
%     x and the leader falls from its first character.
%
% X5. THE LEADER AND THE PROJECTION TOUCH WHAT THEY POINT AT. The accent leader
%     stopped at 9.12 against a disc whose top edge is 8.9493 -- 4.84pt of
%     white -- so at print it was a detached accent tick under the label, of
%     the same length, weight and orientation as an axis tick. The dotted
%     projection was detached at both ends, 2.0pt below the disc and 0.57pt
%     above the spine. Leader 9.46 -> 8.95 lands on the disc; projection
%     8.85 -> 8.54 leaves the disc's lower edge and reaches the rule.
%
% X6. THREE LABELS THAT READ AS ONE BLOCK ARE SEPARATED, and no label sits on a
%     panel head's baseline any more. SUPERSEDES the pass-9 placements.
%     (i) Band (a)'s accent label moved 9.44 -> 9.56. At 9.44 it stood 7.94pt
%         above the threshold label -- set solid for 7.97pt type -- and the two
%         overlap horizontally by 0.6cm, so they read as one ragged two-line
%         paragraph headed by an accent line. The gap is now 11.3pt.
%     (ii) Band (a)'s terminal label moved off the 9.16 baseline onto 8.86,
%         0.10cm above the end cap it names. At 9.16 it shared a baseline AND
%         the "<name>, <value> nats" construction with the threshold label 8cm
%         to its left, so "the clean preference gap, 0.0642 nats" read as the
%         label of a RULE and a reader looked for a rule at 0.0642 and found a
%         0.22cm cap. It is now unmistakably the label of the axis's own end.
%     (iii) Band (b)'s threshold label moved 5.70 -> 5.44, off the panel head's
%         exact baseline, and its rule's top came down 5.56 -> 5.32 to meet it.
%         On the pass-9 build, scanning the two panel heads read "(b) the first
%         eighth, magnified 8x  ...  10% materiality margin" as one title line.
%         Both threshold labels now stand 0.12cm above their own rule's top.
%
% X7. THE DENOMINATOR IS DELIMITED AT BOTH ENDS. The gap -- the span whose full
%     width is the quantity every effect in section 4.5 is divided by -- was
%     capped at the right by a rising 0.40pt ink stub and at the left by
%     nothing but the same rulegrey hanging tick that 0.02, 0.04 and 0.06 carry,
%     so it did not read as a bounded span at all. A matching ink stub now
%     rises at x = 0. No new weight, no new colour, no height.
%
% X8. THE MAGNIFIER BRACKET: RAISED, AND ITS LABEL NAMES THE OBJECT. It moved
%     7.14 -> 7.22 (rail), tight under the scope line at the clearance the
%     scope line's descenders allow, and its label became "the eighth magnified
%     in (b)" from "magnified in (b)", so the mark is named as well as pointed
%     (S20). TWO HALVES OF WHAT WAS ASKED WERE OVERRULED; see X16 (b) and (c).
%
% X9. BAND (b)'s ZERO RULE NO LONGER READS AS A PANEL EDGE. SUPERSEDES the
%     pass-9 extent 3.58 .. 5.56. Drawn from the spine it was collinear and
%     continuous with the axis's left terminus, that terminus's hanging tick
%     and the panel's x origin, ran the full height of the data region, and
%     matched the threshold rule on the right for length to the point -- so all
%     three readers took the pair for the two sides of a frame round the rows,
%     which is the object S12 exists to forbid, and one noted that band (a)
%     draws no such rule at its own zero. It now runs 3.86 .. 5.28: it spans
%     the three rows only, stands 0.28cm clear of the spine and 0.42cm clear of
%     the panel head, and reads as the coordinate the interval bar's left end
%     clears by 0.70cm. It is still S7's register and still unlabelled, because
%     the axis's own "0" numeral stands at its foot.
%
% X10. ONE DASH, ONE MEANING. Rows one and two glossed their arm with an em
%     dash in the same position after the token where row three and band (a)
%     set a math minus, so one glyph class carried "namely" twice and
%     "subtract" once on a figure whose entire subject is a subtraction. The
%     two glosses now take a comma. The minus in "swap - permuted" is the only
%     dash inside any label on the drawing.
%
% X11. THE FOOT BLOCK'S LEADING. The two lines were set 0.24cm apart, which is
%     6.81pt of leading for 7.97pt type -- tighter than solid, tighter than
%     every other stacked pair in this figure, and tight enough that the
%     descender of "bootstrap" met the dot of the "i" below it. They are now
%     0.28cm apart, the figure's own value everywhere else, and the pinned
%     bounding box floor went 1.76 -> 1.72 so descenders stay inside it.
%
% X12. ONE X ORIGIN, MEASURED RATHER THAN ASSERTED. SUPERSEDES the layout note
%     "inner sep = 1pt ON THE WHOLE PICTURE" in its horizontal half only. With
%     inner sep = 1pt every left-anchored text box stood 1.19pt (0.42mm) inside
%     the axis's left end, so both rules and both left end ticks protruded to
%     the left of the entire text column and S5's "left-aligned to the axis's
%     left end" was met to 0.42mm rather than exactly. inner ysep stays 1pt, so
%     every vertical clearance recorded above is still box arithmetic; inner
%     xsep is 0pt. Measured on this build, the ten panel-level text lines now
%     share ONE ink x0, 59.81pt, against an axis rule beginning at 59.615pt --
%     the residue is the first glyph's own side bearing and nothing else.
%     Band (b)'s three row labels are now placed with an xshift on the MARK's
%     own coordinate rather than at a literal x, so each is tied to the thing
%     it names: 5pt right of the disc in rows one and two, 5pt right of the
%     interval's upper end in row three, which is that row's rightmost ink.
%
% X13. THE AXIS LABELS ARE SET LOWERCASE. SUPERSEDES W16's closing paragraph and
%     repair_pass_9's not_applied entry, both of which recorded that the spec's
%     capitalised string "Preference shift, nats" was followed deliberately
%     because the slate is split on case. It is now "preference shift, nats" in
%     both bands. The reason for changing a spec-verbatim string, recorded so a
%     human can put it back in one edit if the slate is unified the other way:
%     the contract's own ruling is "between 4.11 and A.3: take 4.11", 4.11 sets
%     all four of its axis labels lowercase and 4.13 sets both of its, S19's
%     registry gives the quantity's name lowercase ("preference shift"), and
%     this figure's own panel heads are lowercase after S2 -- so capitalising
%     the axis label alone left one string in the figure with a case of its own.
%     Only the case changed; the quantity and the unit are the spec's words.
%
% X14. THE MAGNIFICATION IS NOW STATED TWICE, NOT THREE TIMES. SUPERSEDES
%     W16(ii), which recorded band (b)'s scope line as set verbatim from the
%     spec. It read "the first eighth of the axis above, magnified 8x" and sat
%     3.2cm below a panel head reading "(b) the first eighth, magnified 8x".
%     The head keeps the factor; the axis-label line keeps the relation to the
%     axis above, which is the fact the head does not carry and the fact an
%     axis-label scope line is for. The exactness of the eight is caption
%     material and the caption states it. Band (a)'s bracket label is the third
%     appearance and it stays: it is the only one that is geometric, it names
%     the bracket, and it is 4cm from the other two.
%
% X15. figs/_PENDING_CAPTIONS.tex NOW EXISTS. SUPERSEDES W9's assertion, and
%     the four assertions in materiality.json, that the sentences cut from the
%     drawing had been written into it: at pass 9 the file did not exist, so
%     the two deleted sentences and the whole proposed caption survived only in
%     the commented block at the foot of this file, which no build reads and no
%     integrator was pointed at. The file is created at this pass, keyed by
%     \label{fig:materiality}, and carries the proposed caption in full, the
%     two sentences deleted under W9 verbatim, and a list of the clauses of the
%     PRINTED caption (Sections/04b-representation.tex:597-620) that the
%     rebuild made false. No builder on this pass may edit Sections/.
%
% X16. FINDINGS OVERRULED, WITH THE REASON, so they are not re-applied blind.
%     (a) "Drop a quietink dotted projection from each of band (b)'s three
%         marks to the spine, as band (a) does for its one mark." OVERRULED.
%         Two of the three would cross the interval bar: the permuted mark
%         stands at x = 1.118 and the swap mark at x = 3.318, and the bar runs
%         0.702 to 3.751 on the row beneath them, so the projections would
%         strike it at two interior points and read as endpoints of something.
%         Band (a)'s projection is not a convention this figure owes band (b);
%         it is a rescue for a single disc standing 2.75mm from the origin on a
%         139.96mm rule with no tick within 4cm. In (b) every mark is on a row,
%         direct-labelled at itself, 4.8 to 15.6mm above a rule whose ticks
%         stand every 4.36cm, and the two arm values the projections would help
%         a reader read off are now PRINTED, in the caption, which is where the
%         same reviewer asked for them.
%     (b) "Raise the bracket into the lane immediately under the tick numerals."
%         OVERRULED. S5 puts the axis-label block directly under the tick
%         numerals; the bracket carries a \scriptsize label of its own, so
%         moving it there interposes a second text block between the numerals
%         and the line that names the quantity, which is the supervisor's first
%         complaint reintroduced by the back door. It was raised as far as the
%         scope line's descenders allow instead (X8).
%     (c) "Reverse the bracket's stubs so they turn UP toward the rule."
%         OVERRULED. With the bracket necessarily BELOW the two-line axis-label
%         block, up-turned stubs bracket the scope line's text -- "residual
%         arm, layer 12, alpha = 0.5" runs 0 to 3.66cm and the bracket spans 0
%         to 1.749cm, so they would enclose its first half. Down-turned stubs
%         agree with the label's own pointer, "the eighth magnified in (b)",
%         and point at the panel that opens the span. This was pass 9's stated
%         intent and it survives the pass.
%     (d) "Drop \texttt on the arm names, or scale the mono, because Inconsolata
%         sets at 6.85pt inside 7.97pt roman." OVERRULED. 6.85pt is the
%         preamble's own scaling of the mono face to Pagella, not a second type
%         size chosen here, and S1 exempts the \texttt token name from the size
%         count for exactly this reason. The tokens are JSON key names --
%         swap, permuted, residual -- and setting them roman would turn
%         "swap, points at the drug scored" into a sentence about an English
%         verb. The register is the corpus's (ladder.tex, comparators.tex).
%     (e) "The threshold carries two label strings, '10% materiality margin,
%         0.00642 nats' in (a) and '10% materiality margin' in (b), which fails
%         S19's mechanical check." OVERRULED on the ledger's own standing
%         principle, applied to diff_perm at W6: the NAME is "10% materiality
%         margin" in both bands, and (a) additionally prints the VALUE. A value
%         is not a second name. It is printed in (a) and not in (b) because (a)
%         is the band that also prints the quantity it is a tenth of, so the
%         tenth is a division a reader performs on two printed numbers 8cm
%         apart -- which is what W5 traded the deleted brace for.
%     (f) "Restore the spec's clause 'the full width of this rule is the
%         model's clean preference gap' to band (a)'s scope line." OVERRULED
%         again, for W16(i)'s reason and one more: with the terminal label now
%         standing on its own end cap (X6 ii) and both ends of the span capped
%         in ink (X7), the drawing says geometrically what that clause says in
%         twelve words, and the panel head says it in five.
%
% X17. MEASURED ON THE PASS-10 BUILD, not estimated. figs/_preview.py --width
%     text materiality reports ok, pages = 1, Overfull 0, Underfull 0. Content
%     bounding box 139.96 x 84.12mm (S14's 138-142mm; the 85mm height cap).
%     Ink 5.07% of pixels below 230 grey at 200dpi inside that box (ceiling
%     6.0%). 572 characters of figure text, 4.86 per 100mm^2 (ceiling 5.0).
%     Rule weights, counted off the PDF: 0.30pt x10 (nine rulegrey furniture
%     strokes and one accent leader), 0.40pt x4 (the two spines and the two end
%     caps), 0.50pt densely dotted x2, 0.90pt x2, 1.00pt x1 -- five weights,
%     each with more than one instance or declared as a distinct kind of object
%     at W11 -- plus the 0.80pt open-disc edge, which is a mark edge and not a
%     rule. Type: 7.97pt (400 characters) and 9.96pt (105) in TeXGyrePagella,
%     which are \scriptsize and \footnotesize; 6.85pt (58) in Inconsolata,
%     which S1 exempts; and 8.30pt (8) and 10.38pt (1) in the math companions,
%     which are the SAME two nominal sizes set from the math families -- the
%     minus in "swap - permuted" and the \times in "magnified 8x". Italic: the
%     render carries exactly one span in an italic-named face, PazoMath-Italic
%     at 7.97pt, and it is the mathematical $\alpha$ of "alpha = 0.5" -- 2 of
%     572 characters, 0.35%, which is a symbol's own face and not an italic
%     voice, the use S16 allows. There is no \itshape anywhere in the source.
%     PASS 9's W12 and the ledger's S16 line both said "0%"; the figure of
%     record is 0.35%, and it is the alpha. Sans: 0%. Bold: 0%. Rectangles: 0.
%     Arrowheads: 0.
% ===========================================================================
% ===========================================================================
% ===========================================================================
% PASS 11 -- 2026-08-19. REPAIR PASS on the pass-10 repair. Three adversarial
% readers went over the built figure independently and returned twenty-three
% findings. NOTHING ABOVE IS DELETED: passes 1-10 are the audit trail and stay
% as written. Every clause of passes 1-10 that this pass made false, and every
% clause those passes left false, is superseded here BY NAME (S25). The
% findings that were OVERRULED are recorded under Y15 with the reason, so that
% a later builder does not re-apply them blind.
%
% NO NUMERIC VALUE CHANGED IN THIS PASS EITHER. Every quantity drawn or
% printed is the same float from the same key path as at passes 8, 9 and 10:
% clean_ab_gap 0.06420005919957501, equiv_margin 0.006420005919957502,
% effect_perm 0.0006409306049362111, effect_swap 0.0019027651015006252,
% diff_perm 0.0012618344965644142, ci_perm [0.0004024989353392094,
% 0.002150955667015771], n_pair_clusters 58, and the exact magnification 8.
% Both maps are untouched: \ga = (#1)*218.0000 and \gb = (#1)*1744.0000, with
% 1744.0000 = 8 x 218.0000 exactly (rechecked on this build: 8*218.0 == 1744.0
% is True in IEEE double, and 0.06420005919957501*218.0 ==
% 0.008025007399946877*1744.0 is True, both ends coinciding at
% 13.995612905507352cm). ONE PRINTED STRING changed its number of decimal
% places without changing its value -- band (b)'s middle tick numeral, 0.005
% -> 0.0050 (Y11) -- and one value that was printed on a row of (b) is now
% printed in (a) instead (Y2, Y8). Nothing was rounded differently, and nothing
% new was computed.
%
% WHAT MOVED, IN ONE LINE EACH. The changes are commented at the point of edit
% as Y1..Y13 and summarised here.
%   Y1  band (a)'s panel head names the object by the registry string.
%   Y2  band (a)'s one label prints the estimate and drops 0.06cm.
%   Y3  the axis's right end is computed by \ga, not transcribed.
%   Y4  band (a)'s scope line drops 0.12cm and says "layer 12 of seven".
%   Y5  the magnifier bracket's stubs are visible and it stands clear.
%   Y6  band (b)'s zero rule is named "zero".
%   Y7  the mark-to-label gap in all three rows of (b) goes 5pt -> 9pt.
%   Y8  the estimate's value comes off (b)'s row-three baseline label.
%   Y9  0.06cm more air under row two.
%   Y10 band (b)'s rule is capped at its right end, in (a)'s cap.
%   Y11 band (b)'s middle tick numeral prints 0.0050.
%   Y12 the foot block's leading and its stand-off both open.
%   Y13 the pinned bounding box floor goes 1.72 -> 1.68.
%
% Y14. CLAUSES OF PASSES 1-10 THAT WERE ALREADY FALSE OF THE BUILT FIGURE AND
%     WERE SUPERSEDED NOWHERE. W15 and X1 each did this job for the pass they
%     followed; neither went back over the pass-1 to pass-8 LAYOUT NOTES and
%     DESIGN DECISIONS, so an auditor reading this header top to bottom was
%     handed five descriptions of a figure that does not exist -- the exact
%     hazard X1 was written to prevent. Each is named and given its drawn value.
%     (i) DESIGN DECISION 5, "THE MEASUREMENT SITS 0.22cm ABOVE BAND ONE'S
%         SPINE ... with its dot over a 2.1pt white halo ... the bar, the dot
%         and the spine are three separable things."
%         -> The disc is drawn at y = 8.90 against a spine at 8.54, which is
%            0.36cm and not 0.22cm. There is NO white halo under it: the only
%            \fill[white] in the file is band (b) row one's open control disc.
%            And there is no bar in (a) at all, since W6 drew the interval once
%            and gave (a) a single filled disc. The decision's REASON -- that a
%            mark on the spine reads as a thickening of the axis, so it is
%            lifted clear -- survives; its three measurements do not.
%     (ii) LAYOUT NOTE "TICK NUMERALS CARRY fill=white AND ARE DRAWN LAST IN
%         THEIR BAND."
%         -> num/.style and numL/.style carry no fill. The white box existed to
%            interrupt the magnifier's left leader where it crossed band one's
%            "0"; W10 deleted that leader, so nothing crosses the numerals and
%            drawing order no longer matters.
%     (iii) LAYOUT NOTE "The band-two tag is anchored at x = 0.75 and not at
%         x = 0.00, because the magnifier's left leader is a vertical at
%         x = 0.55."
%         -> Neither object exists: the tag went at W7 (replaced by 4.11's
%            panel head) and the leader at W10.
%     (iv) LAYOUT NOTE "BAND ONE'S RIGHT END is a terminal tick RISING from the
%         spine, 0.30cm tall against the 0.10cm hanging axis ticks, with its
%         two label lines right-aligned on the same x = 13.85."
%         -> It rises 0.22cm, from 8.54 to 8.76, at x =
%            \ga{0.06420005919957501} = 13.995612905507352, and it carries ONE
%            label line, not two. Band (b)'s right end now carries the same cap
%            at the same x (Y10).
%     (v) LAYOUT NOTE "The bounding box is pinned with a bare \path so the
%         measure is exact and nothing is clipped: (0.0000,0.30) to
%         (14.1900,8.98) = 141.9 x 86.8mm."
%         -> It is pinned (0.00,1.68) to (13.9957,10.18) after Y13, and the
%            measured content box is at Y16. The 141.9 x 86.8mm figure is a
%            pass-2 measurement and is over both the 142mm measure and the 85mm
%            height cap.
%     (vi) W11's "The 0.50pt densely dotted stroke is used once, for band two's
%         zero rule, and is a distinct kind of object."
%         -> It has two instances in this build and had two in pass 10, as
%            X17's own weight count records: band (a)'s projection from the disc
%            to the spine, and band (b)'s zero rule. It is not a singleton and
%            needs no singleton declaration under S10. The 1.00pt interval bar
%            IS a singleton and its declaration stands.
%
% Y15. FINDINGS OVERRULED, WITH THE REASON, so they are not re-applied blind.
%     (a) "Take the open control disc's edge to ink rather than quietink, as
%         fig 4.11(b) does for its matched-dimension random slab." OVERRULED.
%         S8 gives quietink to a comparator, a null or a control and S9 says
%         the open disc's 0.8pt edge is drawn "in its colour"; ink is the rank
%         S8 reserves for thresholds, axes, axis labels, titles and definitional
%         prose, so an ink ring would put the control in the apparatus rank and
%         make the comparator heavier than the accent mark under test. 4.11(b)
%         can draw its null black because its palette has no quietink rank at
%         all; this figure has four ranks and the control's is quietink. The
%         defect the finding reports is real and blocking and is fixed by the
%         gap alone (Y7): 7.2pt of white against the mono face's own 3.43pt
%         character advance, where 3.28pt was letter-spacing.
%         The companion finding, "quietink carries two ranks -- the one
%         measured comparator's label and all four apparatus lines", is
%         OVERRULED for the same reason from the other side: S5 mandates
%         quietink for both scope lines and S17 for the foot block, so the
%         collision cannot be broken by recolouring either party. It is broken
%         by position, which is how this figure has always broken it: the
%         control's label is anchored on a mark, on a row, inside the data
%         field; every apparatus line is panel-level text on the figure's one
%         x origin, under an axis.
%     (b) "Drop band (b)'s two arm dots. Their values are quoted in no sentence
%         of the thesis and the proposed caption already prints them, and the
%         one interval bar horizontally contains both of them." OVERRULED.
%         The spec's keep-list makes the three-row structure -- permuted, swap,
%         and their difference -- load-bearing, and CF14 fails a rebuild that
%         loses a load-bearing piece even if it passes every typographic check.
%         Without the two arms the difference is a bare number with a bar and
%         the figure stops drawing the subtraction that is its subject; design
%         decisions 7 and 14 record that reasoning and it stands.
%         THE DIAGNOSIS IS ACCEPTED AND IS REAL, and it is answered at Y9
%         rather than by deletion. The arithmetic is an accident worth stating
%         plainly: the bar runs \gb{0.0004024989353392094} = 0.702cm to
%         \gb{0.002150955667015771} = 3.751cm, the control stands at 1.118 and
%         the arm at 3.318, so the interval on the DIFFERENCE brackets both of
%         its own operands, and \gb{diff_perm} = 2.201 falls within 0.02cm of
%         their midpoint. What is done about it: the row pitch is no longer
%         flat (0.58 / 0.64 rather than 0.54 / 0.54), so the derived row sits
%         off the pair it is derived from; the row's name states the derivation
%         at the bar's own end; and rows one and two name their arms
%         individually immediately above, so "swap $-$ permuted" is read after
%         "permuted" and "swap" and not before them.
%         The finding's own fallback -- "move rows one and two onto their own
%         short strip with their own zero" -- is also OVERRULED: a second rule
%         for one quantity at one scale inside one panel is a second axis a
%         reader has to reconcile with the first, against S13's one grid, and
%         it would put two zeros 5mm apart on a figure whose subject is a
%         distance from zero.
%     (c) "Move the magnifier bracket into band (a)'s tick lane, rail at the
%         hanging ticks' depth with its stubs turned UP to the spine; or, if
%         that lane cannot be had, delete the bracket and tint the magnified
%         span in (a) and the whole of (b)'s data field instead." OVERRULED,
%         both halves. X16 (b) and (c) already overruled the first and the
%         reason holds: a bracket in the tick lane needs its label in the
%         tick-numeral lane, which is the defect S24 indicts in fig:purify,
%         where the word "four" sets on the x-tick baseline and reads as a
%         leftmost tick. The tint is worse, not better: S12 permits a filled
%         rectangle only where the rectangle IS an object in the argument and
%         forbids outright a rectangle drawn round a panel, and a tint over the
%         whole of (b)'s data field is that rectangle by another means -- the
%         object W4 deleted. The finding's real complaint, that the assembly did
%         not read as a bracket at all, is accepted and fixed at Y5.
%     (d) "Raise the bracket's rail onto its label's own baseline, so the rail
%         and the label read as one object rather than the rail reading as a
%         strike-through beside the words." OVERRULED as already true. The rail
%         at 7.02 stands 0.06cm above a label baseline at 6.96, which is that
%         label's own x-height centre, so rail and label are optically on one
%         line. What made it read as a rule was the 0.51mm stubs and the scope
%         line overhanging both its ends; both are fixed at Y5.
%     (e) "Label band (b)'s right cap 'one eighth of it, 0.008025 nats', so
%         that the exact eight becomes a division a reader can perform on two
%         printed numbers." OVERRULED on geometry, not on principle. Band (b)'s
%         threshold rule stands at \gb{0.006420005919957502} = 11.196cm and
%         runs the full height of the data field; a label right-aligned on the
%         axis's end at 13.996cm measures 3.4 to 4.7cm at \scriptsize and
%         crosses it, and the only string short enough to clear the rule is a
%         bare numeral with no name, which S20 forbids. There is no lane above
%         the rule either: (b)'s threshold label already occupies it, ending at
%         11.10cm. The magnification is stated by (b)'s panel head, by (b)'s
%         axis-label scope line and by (a)'s bracket label, and its EXACTNESS
%         is caption material (X14, and the caption states it). The cap added
%         at Y10 is axis furniture -- the rule's own end -- and not a mark, so
%         it needs no name under S20; its end is named by (b)'s axis label,
%         which is W13's reasoning and survives.
%     (f) "Restore the spec's clause 'the full width of this rule is the
%         model's clean preference gap' to band (a)'s scope line: line one
%         reads 'preference shift, nats', but band (a)'s full width is
%         clean_ab_gap, the model's contrast score with NO intervention -- a
%         level, not a shift -- so a reader who takes the axis label at face
%         value reads the right terminus as a preference shift of 0.0642 nats."
%         OVERRULED a third time. The REASON RECORDED AT W16(i) AND X16(f) IS
%         WITHDRAWN AND REPLACED, because it was wrong: both cited S18, and
%         S18 forbids a drawing from restating a clause of ITS OWN CAPTION or
%         of the paragraph that introduces the float. It says nothing about a
%         drawing carrying two related statements of its own. The clause is
%         refused on S5 and S15 instead. S5 caps the axis-label block at two
%         lines and gives line two to scope or a scale note; a thirty-character
%         clause about what the axis's other end IS is neither. And the fact is
%         already stated at the two places a reader looks for it: the panel head
%         reads "(a) the whole clean preference gap" after Y1, and the terminal
%         label at the axis's own right end reads "the clean preference gap,
%         0.0642 nats". THE FINDING'S UNDERLYING OBSERVATION IS RIGHT AND IS
%         RECORDED HERE so that no later reader has to rediscover it: this
%         axis's quantity name is the quantity its MARKS carry, and its right
%         terminus is a level of that same contrast score measured with no
%         intervention, named as such at the terminus. Laying the two on one
%         rule is the figure's whole construction and the thesis's own -- 04b:
%         663 divides every effect by "that model's own clean preference gap,
%         measured with no intervention".
%     (g) "Scale the mono up so that \texttt sets at the roman's x-height:
%         6.85 x (3.96/3.24) = 8.37pt." OVERRULED again, as at X16(d), and for
%         a reason that finding does not meet. 6.85pt is preamble.tex's global
%         scaling of Inconsolata to Pagella and it applies to every \texttt in
%         the thesis. Overriding it inside one figure would make this figure's
%         JSON key names set differently from comparators.tex's, ladder.tex's
%         and the body text's, which is a worse inconsistency than an x-height
%         step inside a label; and S1 exempts the \texttt token name from the
%         size count precisely so that the document's global scaling can stand
%         without a figure having to fight it.
%     (h) "In band (a) the two black labels that share the '<name>, <value>
%         nats' construction are the CLOSER pair (0.30cm) while the accent
%         label, which differs from both in face, colour and construction, is
%         the further (0.40cm); drop the terminal label 8.86 -> 8.80 and lift
%         the threshold label 9.16 -> 9.22." OVERRULED as proposed, accepted in
%         substance. The two black labels are 34mm apart horizontally, sit on
%         different baselines, and one is right-aligned onto the axis's own end
%         cap while the other is left-set beside a rule 8cm away; X6(ii)
%         separated them for the stronger reason that they had SHARED a
%         baseline, and the proposed 8.80 would set a 7.97pt baseline 1.1pt
%         above the 0.22cm cap it labels, re-creating the collision X6(ii)
%         removed. After Y2 the stack reads 0.40 / 0.34 / 0.30 from the head
%         down -- monotone, so nothing alternates and no pair is singled out by
%         its spacing.
%
% Y16. MEASURED ON THE PASS-11 BUILD, not estimated, and SUPERSEDING X17's
%     bounding box by name. figs/_preview.py --width text materiality reports
%     ok, pages = 1, Overfull 0, Underfull 0.
%       content bounding box   140.34 x 84.33mm   (S14's 138-142mm; 85mm cap)
%       ink                    5.36% of pixels below 230 grey at 200dpi (6.0%)
%       figure text            587 characters, 4.96 per 100mm^2   (ceiling 5.0)
%     X17 recorded 139.96 x 84.12mm for the pass-10 build; re-measured with the
%     contract's own threshold (pixels below 230 grey at 200dpi, clipped to the
%     content box) that build is 140.21 x 83.82mm. The +0.25mm and -0.30mm are
%     a method difference and not a drawing change, and they are recorded here
%     rather than left as an unexplained discrepancy: X17's number came from a
%     box measured before the antialiased edge columns of the axis caps were
%     counted. Both builds are inside every limit; this one is 0.67mm under the
%     height cap and 0.04 characters per 100mm^2 under the density ceiling,
%     which is why six characters came out of the foot block when "of seven",
%     "clean" and "zero" went in.
%     Rule weights, counted off the PDF: 0.30pt x10 (nine rulegrey furniture
%     strokes -- eight ticks and the bracket polyline -- plus one accent
%     leader), 0.40pt x5 (two spines and three end caps), 0.50pt densely dotted
%     x2, 0.90pt x2, 1.00pt x1. Five weights, no singleton except the interval
%     bar, which W11 declares as a distinct kind of object. Plus the 0.80pt
%     open-disc edge, which is a mark edge and not a rule.
%     Type: 7.97pt (417 characters) and 9.96pt (111) in TeXGyrePagella, which
%     are \scriptsize and \footnotesize; 6.85pt (58) in Inconsolata, which S1
%     exempts; 8.30pt (8) and 10.38pt (1) in the math companions, which are the
%     same two nominal sizes set from the math families. Italic: one span,
%     PazoMath-Italic at 7.97pt, the mathematical $\alpha$ -- 2 of 587
%     characters, 0.34%, a symbol's own face and not an italic voice (S16).
%     Sans 0%. Bold 0%. Rectangles 0. Arrowheads 0.
%
% Y17. THE CAPTION. figs/_PENDING_CAPTIONS.tex is updated in the same pass and
%     three of its clauses changed with the drawing: band (a)'s mark now prints
%     its value, band (b)'s row three no longer does, and (b)'s rule is capped
%     at both -- at its right -- end. ONE ITEM IN THAT FILE IS BLOCKING and is
%     flagged there as such: the printed caption (04b:617-618) says only "this
%     is one arm at one depth at one rung", which is a scope statement, while
%     layer 12 is the argmax of the anchor seed's seven-depth sweep and the
%     only one of the seven whose interval excludes zero. The drawing now says
%     "layer 12 of seven" (Y4), which discloses that the depth was chosen from
%     a set but not that it was the favourable one; only the caption can carry
%     that, and no builder on this pass may edit Sections/.
% ===========================================================================
%
% ===========================================================================
% v-BLOCK, 2026-08-19: WHOLE-SLATE PASS (fig:comparators, fig:licensing,
% fig:purify, fig:hook, fig:materiality reconciled against one another and
% against fig:mechanism, which is the document's reference standard and was
% NOT rebuilt). Nothing above this line is deleted; where an earlier claim is
% now false it is superseded EXPLICITLY below.
%
% WHY. The five figures were rebuilt independently and drifted. This pass owns
% only the differences BETWEEN them. No value changed in any of the five; the
% sibling .json ledgers carry the value-to-position maps, old and new.
%
% THE SLATE RULES THAT NOW HOLD ACROSS ALL FIVE (each stated once here, in
% every file, so any future single-figure edit can see what it would break):
%
%  R1 TICK NUMERALS are quietink, \scriptsize, inner sep 0, anchored north on
%     the tick and set 0.16cm below the axis rule, over a 0.10cm rulegrey
%     0.30pt tick. Before: comparators/licensing/purify quietink, hook and
%     materiality ink; ticks 0.10 / 0.12 / 0.14cm; three different gaps.
%  R2 A THRESHOLD'S OWN LABEL IS BLACK, matching the black solid 0.90pt rule
%     it names (S6). Before: comparators black, purify's "gate, 0.8" and
%     materiality's two margin labels ink.
%  R3 ONE NAME PER OBJECT (S19). The pre-registered 10% rule is "10%
%     materiality margin" in fig:comparators row (c), fig:materiality (a) and
%     (b), and in figs/family.py, which is where the name comes from.
%     fig:comparators' "pre-registered margin" is superseded; "pre-registered"
%     is now the caption's word, not the drawing's. Zero, where it is the null
%     a bar is read against, is "zero: no identity effect" in BOTH
%     fig:comparators row (c) and fig:materiality (b) -- one quantity, one
%     wording. V_pure is spelt with \textrm everywhere, as figs/slab.tex does.
%  R4 COLOUR KEY, one key for five figures (S8). accent = the drug-side object
%     under test. quietink = every comparator, null, control, replication or
%     discounted arm, AND its label, AND its leader. ink/black = thresholds,
%     axis rules, axis labels, panel heads, definitional prose. rulegrey =
%     furniture only. Within quietink the GLYPH separates two kinds: an OPEN
%     disc is a construction built to carry no signal (a null, a control, a
%     raw/before series); a FILLED disc is a real measurement that is not the
%     headline. fig:hook (d) drew both its null and its cell-line control in
%     ink and is corrected.
%  R5 PROJECTION AND ZERO RULES are quietink 0.50pt densely dotted (S7).
%     fig:hook's two panel-(c) projection lines were rulegrey 0.30pt densely
%     dashed and are corrected.
%  R6 DEFINITIONAL PROSE is \footnotesize ink and OUTRANKS every label (S4);
%     at most one block per figure, on the panel's own x origin. fig:hook's
%     one prose line was \scriptsize and is raised. fig:purify's panel-(b)
%     block carried an argument clause ("so causal claims use the purified
%     slab, not the raw") that its caption already carries, and is now
%     definitional only.
%  R7 ONE INTERVAL DECLARATION per figure, \scriptsize quietink, at the foot,
%     on the figure's x origin, opening with the word "Intervals:" and naming
%     EVERY panel including the ones with no measured mark (S17).
%     fig:comparators already did this and is the model; the other four now
%     match its form.
%  R8 THE FIVE-FIELD PROMPT STRIP is ONE object drawn in two figures, and the
%     two drawings are now congruent: cell boundaries 0 / 1.12 / 1.90 / 2.64 /
%     4.23 / 6.90 cm and cell height 0.42cm in both fig:licensing (a) and
%     fig:hook (a) and (b), labels centred in their cells, the drug cell
%     accent-bordered at 0.40pt on accent!12 and its label NOT bold. Before:
%     licensing 6.90cm wide with 0.62cm cells and a bold label, hook 7.44cm
%     wide with 0.42cm cells and a plain label -- the same five fields at two
%     sizes, four pages apart, with each caption pointing at the other.
%  R9 MEASURE. Every fullwidth figure draws 172.0-174.0mm; the one textwidth
%     figure draws 140.3mm. Ink at 200dpi inside the content bbox, counted as
%     greyscale luminance below 230 (the contract's own method), is now
%     5.63-5.99% on all five, under the 6.0% body ceiling, where before this
%     pass it ran 5.36-6.54%.
%
% WHAT THIS PASS DELIBERATELY DID NOT UNIFY, so that a later reader does not
% think it was missed:
%  * PANEL-HEAD CASE. fig:comparators sets "(a) Is it there?" and the other
%    four set "(a) where the drug name enters". comparators' heads are
%    interrogative SENTENCES and sentence case is correct for a sentence; the
%    v4 header argued this and the argument stands. fig:mechanism's lowercase
%    heads are noun phrases, not questions. One capital letter, four times.
%  * PANEL-HEAD POSITION. fig:comparators puts its heads in a right-aligned
%    left stub, the other four above the panel at its x origin. Moving them
%    above the rules costs about 20mm of height against a 118mm cap that the
%    figure already meets at 117.60mm. Not available.
%  * TWO-COLUMN GUTTER. purify and hook split at x = 8.40cm, licensing at
%    7.75cm. licensing cannot move right: "(d) what a negative would buy" is
%    4.85cm wide and its origin is already at 12.40 of a 17.30cm measure, with
%    1.5pt of slack. Moving purify and hook LEFT to 7.75 would be arbitrary.
% ===========================================================================
%
% CHANGED IN THIS FILE. (i) No math sign is set from a CM font any more. The
% panel-(b) head's 8x (CMSY10 at 10.38pt, the largest glyph on the figure and
% a THIRD type size above 7pt) is now 8\texttimes{}; the two minus signs of
% "swap - permuted" are \textminus{}; +0.0013, 10%, 0.00642, 0.0642, 12, 58
% and 95% lose their math mode, and $\alpha = 0.5$ becomes $\alpha$ = 0.5,
% which is fig:comparators' idiom. The figure now prints exactly 9.96 and
% 7.97pt in one face, plus the \texttt arm names. (ii) Tick numerals go ink ->
% quietink with inner sep 0, set 0.16cm below their rule (R1). (iii) Both
% "10% materiality margin" labels go black, matching the black solid 0.90pt
% rule they name (R2). (iv) Band (b)'s zero label "zero" -> "zero: no identity
% effect", which is fig:comparators row (c)'s wording for the same null under
% the same quantity (R3), and goes ink -> quietink so that the label matches
% the dotted quietink rule it names. (v) \texttt{residual} -> residual in the
% scope line: fig:comparators names the same arm "residual seeds at layer 12"
% in plain roman, and an arm name in prose is not a code token. \texttt{swap}
% and \texttt{permuted} stay, because they are the two stored conditions and
% tab:identity sets them so. (vi) The foot block takes the slate's declaration
% form (R7).
% NOT CHANGED, and still declared: the 1.00pt interval bar is a singleton
% weight because this figure stores exactly one interval (S10's permission);
% the 0.80pt open-disc edge is a MARK edge, not a rule.
% RE-MEASURED: 140.34 x 84.45mm (cap 85, textwidth band 138-142); ink 5.630%;
% 620 characters, 5.23 per 100mm^2. One page, zero Overfull, zero Underfull.
% ===========================================================================
\begin{tikzpicture}[x=1cm, y=1cm, font=\small, inner ysep=1pt, inner xsep=0pt,
  %% TYPE. Two sizes above 7pt and one face (S1, W12). \footnotesize = 10.0pt
  %% carries the two panel heads and the two axis-label line ones and nothing
  %% else; \scriptsize = 8.0pt carries ticks, direct labels, scope lines, the
  %% threshold labels and the interval declaration. No sans, no bold, no italic.
  %% INNER SEP, PASS 10 (X12). It was inner sep = 1pt on the whole picture; the
  %% vertical half is unchanged, so every clearance below is still box
  %% arithmetic, but the horizontal half is now 0pt. With xsep = 1pt every
  %% left-anchored text box stood 1.19pt (0.42mm) inside the axis's left end,
  %% so both rules protruded to the left of the entire text column and S5's
  %% "left-aligned to the axis's left end" was met to 0.42mm and not exactly.
  head/.style  ={font=\footnotesize, text=ink,      anchor=base west},
  axlab/.style ={font=\footnotesize, text=ink,      anchor=base west},
  scope/.style ={font=\scriptsize,   text=quietink, anchor=base west},
  qlabz/.style ={font=\scriptsize,   text=quietink, anchor=base west},
  lab/.style   ={font=\scriptsize,   text=ink,      anchor=base west},
  labR/.style  ={font=\scriptsize,   text=ink,      anchor=base east},
  acc/.style   ={font=\scriptsize,   text=accent,   anchor=base west},
  rowq/.style  ={font=\scriptsize,   text=quietink, anchor=west},
  rowa/.style  ={font=\scriptsize,   text=accent,   anchor=west},
  %% Tick numerals are INK, not quietink: they are part of the axis, and S8
  %% gives ink to axes and axis labels while quietink means a comparator or a
  %% control. 4.11 and 4.13 both set their tick numerals dark for this reason.
  num/.style   ={font=\scriptsize,   text=quietink, anchor=north,
                 inner sep=0pt},
  numL/.style  ={font=\scriptsize,   text=quietink, anchor=north west,
                 inner sep=0pt},
  blab/.style  ={font=\scriptsize,   text=black,    anchor=base west},
  blabR/.style ={font=\scriptsize,   text=black,    anchor=base east},
  %% FIVE STROKE WEIGHTS, each with one meaning (S10, W11).
  tick/.style  ={rulegrey, line width=0.3pt},                  % furniture
  lead/.style  ={rulegrey, line width=0.3pt},                  % furniture
  accled/.style={accent,   line width=0.3pt},                  % furniture, in the mark's colour (S22)
  spine/.style ={ink,      line width=0.4pt},                  % axis rule and its two end caps
  proj/.style  ={quietink, line width=0.5pt, densely dotted},  % zero rule / projection (S7)
  thresh/.style={black,    line width=0.9pt},                  % the one threshold (S6, W3)
  bar/.style   ={accent,   line width=1.0pt},                  % interval
  ctrl/.style  ={quietink, line width=0.8pt}]                  % open-disc edge: a MARK edge, not a rule (S9, W11)

  %% ==== THE TWO MAPS ====================================================
  %% Superseded constants, kept in materiality.json: 0.55+(#1)*207.1650 and
  %% 0.55+(#1)*1657.3200. The 0.55cm inset is gone so that the axis's left end
  %% and the panel's x origin are one coordinate (S13, W1).
  %% 1744.0000 = 8 x 218.0000 EXACTLY, so the magnification is exactly eight,
  %% and \ga{clean_ab_gap} = \gb{clean_ab_gap/8} = 13.995613cm exactly.
  %% PASS 10, X1: the pass-9 header block records these two constants as
  %% 215.7500 and 1726.0000 and the shared right end as 13.851163cm. Those
  %% numbers were never drawn by this file. The two lines below are the map,
  %% and materiality.json has always recorded them correctly.
  \def\ga#1{{(#1)*218.0000}}   % band (a): 0 .. 0.06420005919957501 nats
  \def\gb#1{{(#1)*1744.0000}}  % band (b): 0 .. 0.008025007399946877 nats

  %% ==================== BAND (a) -- THE WHOLE GAP ========================
  %% PASS 11, Y1: the head names the object by the registry string. It read
  %% "(a) the whole preference gap" while the terminal label 8cm to its right
  %% read "the clean preference gap", so one JSON key (clean_ab_gap) carried
  %% two names inside one panel, which is what S19's mechanical check forbids.
  %% "whole" is kept because it is the contrast with (b)'s "first eighth".
  \node[head] at (0.00,9.90) {(a) the whole clean preference gap};

  % The estimate, drawn once and without a bar (W6). At true scale it stands
  % 2.75mm from the origin on a 139.96mm rule, which is the panel's whole point.
  % Its name is band (b) row three's name, because one JSON key takes one name.
  % PASS 10, X4: THE LABEL IS ANCHORED AT THE MARK, not at the panel origin.
  % Pass 9 set it base west on x = 0.00 while its leader fell from
  % \ga{diff_perm} = 0.2751, so the leader descended from inside the word
  % "swap" and the mark read as the SWAP arm -- whose value is 0.0019, not the
  % 0.0013 drawn. That is the one mark in band (a); mis-naming it loses the
  % panel. The label now begins on the mark's own x, and the leader falls from
  % its first character.
  % PASS 11, Y2: THE LABEL PRINTS THE VALUE, AND IT DROPPED 0.06cm.
  % Two defects, one edit. (i) At 9.56 the label stood 0.34cm under the panel
  % head, near-left-aligned with it (7.8pt of indent), so head and label read
  % as a two-line title and the panel's single mark lost its name; the head's
  % descenders cleared the label's ascenders by 2.6pt. At 9.50 that clearance
  % is 4.3pt and the gap to the threshold label below is 0.34cm = 9.7pt, which
  % is still wider than the 7.94pt X6(i) condemned. (ii) Without a value the
  % math minus between two mono tokens parses as easily as a gloss dash --
  % "swap, that is, permuted" -- so the one mark in (a) could be read as the
  % SWAP arm, whose value is 0.0019 and not the 0.0013 drawn. The value now
  % settles it, and it also makes (a) self-sufficient for the figure's
  % question: 0.0013 against the 0.0642 printed at the axis's right end, and
  % against the 0.00642 printed at the threshold.
  % WHY A VALUE MAY BE PRINTED HERE AND NOT ON A ROW OF (b). This label sits
  % 0.96cm above the spine and is leadered to its disc; positions in that lane
  % are not read as values, which is why (a)'s threshold label already prints
  % 0.00642 three centimetres right of the rule it names and its terminal
  % label prints 0.0642 four centimetres left of the cap it names. A label set
  % on a data row's OWN baseline is read positionally, which is why the value
  % came off (b)'s row three (Y8).
  \node[acc] at (\ga{0.0012618344965644142},9.50)
    {\texttt{swap} \textminus{} \texttt{permuted}, +0.0013 nats};
  % The label's leader is accent 0.30pt, the mark's own colour (S22). PASS 10,
  % X5: it now LANDS ON THE DISC. It used to stop at 9.12 against a disc whose
  % top edge is 8.9493, leaving 4.84pt of white, so the leader read as a stray
  % accent tick under the label and the label read as unattached.
  \draw[accled] (\ga{0.0012618344965644142},9.40)
             -- (\ga{0.0012618344965644142},8.95);
  % The stem down onto the spine anchors the mark to a coordinate rather than
  % leaving it floating above the rule; pass 8 added it after both reviewers
  % reported they could not tell whether the mark was a mark, a tick or a
  % printing artefact, and it survives the rebuild. It is drawn in the PROJECTION
  % register -- quietink 0.5pt densely dotted, which is what S7 reserves that
  % stroke for -- and not in the leader's accent, because on the build that drew
  % leader and stem at one weight and one colour the assembly read as a VERTICAL
  % interval bar with a dot on it, which is the one thing this figure must not
  % appear to draw. PASS 10, X5: it now touches the disc's lower edge (8.85
  % against the disc's 8.8507) and the spine (8.54), instead of floating 2.0pt
  % clear of the one and 0.57pt clear of the other.
  \draw[proj] (\ga{0.0012618344965644142},8.85)
           -- (\ga{0.0012618344965644142},8.54);
  \fill[accent] (\ga{0.0012618344965644142},8.90) circle (1.4pt);

  % The one threshold, black and solid, in both bands (S6, CF2, W3). Its value
  % is printed here and the quantity it is a tenth of is printed at the axis's
  % right end, so the tenth is a division on two printed numbers and needs no
  % brace (W5).
  % PASS 10, X3: IT STANDS ON THE SPINE AND NO LONGER CROSSES IT. Pass 9 drew
  % it from 8.38, which is 0.06cm BELOW the hanging ticks' ends and 0.02cm above
  % the tick numerals, so the heaviest stroke in the panel doubled as an
  % outsized unnumbered tick sitting between the numerals "0" and "0.02" -- the
  % same "the threshold merges with the tick furniture" defect S7 indicts in
  % fig:comparators row (a). Band (b) already drew its half of the same
  % threshold from its own spine upward; the two now agree.
  \draw[thresh] (\ga{0.006420005919957502},8.54)
             -- (\ga{0.006420005919957502},9.04);
  \node[blab] at (1.50,9.16) {10\% materiality margin, 0.00642 nats};

  % The axis. PASS 10, X7: BOTH ENDS NOW CARRY THE SAME CAP. The right end is
  % the clean preference gap and has always carried a rising 0.22cm ink stub;
  % the left end carried only the same rulegrey hanging tick as 0.02, 0.04 and
  % 0.06, so the span whose full width IS the figure's denominator was
  % delimited at one end by an ink stub and at the other by an ordinary tick,
  % and did not read as a bounded span at all. Two matching caps, no new weight
  % and no new colour.
  % PASS 11, Y3: THE AXIS'S RIGHT END IS COMPUTED, NOT TRANSCRIBED. It was
  % hard-coded 13.995613 here, at the end cap and at the terminal label, while
  % W1 and the layout note both assert that "every data x is still computed
  % inside pgfmath from the full-precision float, so no value is rounded into
  % the source". \ga{0.06420005919957501} = 13.995612905507352, so the literal
  % was short by 3.5e-7cm -- invisible, but it is the ONE coordinate this
  % figure exists to draw, and the claim about it was false. Only the pinned
  % bounding box below still carries a rounded literal, and it carries no
  % value: it is a box corner, declared as such at Y13.
  \draw[spine] (0.00,8.54) -- (\ga{0.06420005919957501},8.54);
  \draw[spine] (0.00,8.54) -- (0.00,8.76);
  \draw[spine] (\ga{0.06420005919957501},8.54) -- (\ga{0.06420005919957501},8.76);
  % PASS 10, X6: the terminal label sits ON ITS OWN CAP, at y = 8.86, and no
  % longer shares the 9.16 baseline and the "<name>, <value> nats" construction
  % with the threshold label 8cm to its left. Read cold on the pass-9 build the
  % two were one row of two labels in one style, so "the clean preference gap,
  % 0.0642 nats" read as the label of a RULE and a reader looked for a rule at
  % 0.0642 and found a 0.22cm end cap. It is now 0.10cm above that cap.
  \node[labR] at (\ga{0.06420005919957501},8.86)
    {the clean preference gap, 0.0642 nats};
  \foreach \v/\t in {0.02/0.02, 0.04/0.04, 0.06/0.06}{
    \draw[tick] (\ga{\v},8.54) -- (\ga{\v},8.44);
    \node[num] at (\ga{\v},8.38) {\t};}
  % Zero is both the axis's left end and the panel's x origin, so its numeral is
  % set flush with that origin rather than centred on it: a centred "0" is the
  % only object that would break the figure's single left edge (S13).
  \draw[tick]  (0.00,8.54) -- (0.00,8.44);
  \node[numL] at (0.00,8.38) {0};

  % AXIS LABEL, in axis-label position: directly under the tick numerals, left
  % edge on the axis's left end, quantity and unit on line one (S5, W2).
  % PASS 10, X13: set lowercase, as 4.11 and 4.13 set every axis label of
  % theirs and as S19's registry gives the name.
  % PASS 11, Y4: THE SCOPE LINE DROPS 0.12cm AND DISCLOSES THE SELECTION.
  % (i) LEADING. It was set 0.28cm under a 9.96pt line, which is 7.94pt of
  % baseline-to-baseline for type 2pt larger than the lead -- negative leading,
  % and the tightest of the four stacked pairs in this figure. Four such pairs
  % were set at or under solid and it is the single thing that made the figure
  % look crowded beside 4.11. Both axis-label blocks and the foot block now
  % open to 0.38cm (10.8pt), and the whole of the added height was paid for out
  % of the white between the two bands, which fell 1.14cm -> 0.77cm.
  % (ii) SELECTION. "layer 12" alone is a scope statement. Layer 12 is the
  % argmax of the anchor seed's seven-depth sweep and the only one of the seven
  % whose interval excludes zero (diff_perm_norm with p_perm: +0.0022 p 0.583,
  % -0.0050 p 0.139, -0.0025 p 0.4495, +0.0041 p 0.3955, -0.0019 p 0.635,
  % +0.0197 p 0.001, +0.0059 p 0.4885 at layers 2/4/6/8/9/12/16). "of seven"
  % is nine characters and it tells a reader on the drawing that the depth was
  % one of a set; WHICH one, and why that matters, is caption material and the
  % proposed caption in figs/_PENDING_CAPTIONS.tex carries it. Until a human
  % integrates that caption the printed one (04b:617-618) says only "one arm at
  % one depth at one rung", which is scope and not selection.
  \node[axlab] at (0.00,7.78) {preference shift, nats};
  \node[scope] at (0.00,7.40)
    {residual arm, layer 12 of seven, $\alpha$ = 0.5};

  % The magnified span, named as an object and not traced by a leader (W10).
  % Its two stubs turn DOWN, towards the panel that opens the span, so the
  % bracket points at (b) instead of back at an axis it already stands under.
  % PASS 10, X8: raised 0.16cm, tight under the scope line, and its label now
  % NAMES the object as well as pointing at (b). Two reviewers asked for it to
  % be moved into the lane immediately under the tick numerals with its stubs
  % reversed; both halves of that were overruled, and the reason is recorded in
  % the pass-10 block under X8.
  % PASS 11, Y5: THE STUBS ARE LONG ENOUGH TO SEE, AND THE BRACKET STANDS
  % CLEAR OF THE AXIS-LABEL BLOCK. Measured on the pass-10 render, each stub
  % was 4 pixels at 200dpi -- 1.44pt, 0.51mm -- of 0.30pt rulegrey, which does
  % not survive print, so the assembly was a bare grey horizontal sitting
  % 0.34cm under a scope line that runs past both of its ends: it read as an
  % underline of "residual arm, layer 12," or as an editorial divider, and not
  % as a bracket. Two changes, no new weight and no new colour: the stubs go
  % 0.06cm -> 0.16cm, and the gap from the scope line goes 0.34cm -> 0.44cm,
  % which is larger than the 0.38cm internal leading of the axis-label block
  % above it, so the bracket stops reading as that block's third line.
  % Two harder proposals were overruled and the reasons are at Y15 (c) and (d).
  \draw[lead] (0.00,6.86) -- (0.00,7.02)
           -- (\ga{0.008025007399946877},7.02)
           -- (\ga{0.008025007399946877},6.86);
  \node[scope] at (1.85,6.96) {the eighth magnified in (b)};

  %% ============= BAND (b) -- THE FIRST EIGHTH, MAGNIFIED =================
  %% PASS 11: the whole of band (b) is raised 0.14cm and its axis lowered
  %% 0.10cm, which is where the leading opened at Y4 and Y11 was paid from.
  \node[head] at (0.00,5.84) {(b) the first eighth, magnified 8\texttimes{}};
  % The threshold label sits against the rule it names and takes the LEFT of it
  % here and the RIGHT of it in (a): in each band the label goes on the side
  % where the rule has room, which in (a) -- where the rule stands at a tenth of
  % the measure -- is the right, and in (b) -- where it stands at four fifths --
  % is the left. Set right of the rule here it measures 2.98cm from 11.30 and
  % crosses the right end of the axis at 13.996 (checked on a build, which
  % reported Overfull \hbox 2.84pt).
  % PASS 10, X6: it has come off the panel head's baseline (5.70) and down onto
  % its own, 0.12cm above the rule's top, in the same relation to its rule as
  % (a)'s label is to (a)'s. On the pass-9 build it sat at exactly the head's
  % baseline, so scanning the two panel heads read "(b) the first eighth,
  % magnified 8x ... 10% materiality margin" as one title line.
  \node[blabR] at (11.10,5.46) {10\% materiality margin};
  \draw[thresh] (\gb{0.006420005919957502},3.68)
             -- (\gb{0.006420005919957502},5.34);

  % Zero. Densely dotted quietink 0.5pt is this document's stroke for a
  % coordinate that decides nothing (S7).
  % PASS 11, Y6: IT IS NAMED, ON THE SAME BASELINE AS THE THRESHOLD LABEL AND
  % AT THE SAME 0.12cm ABOVE ITS OWN RULE'S TOP. SUPERSEDES W8's closing
  % clause, that the rule "carries no label, because its foot is the axis's own
  % 0 numeral". The equivalence criterion of 04b:507-513 is a CONJUNCTION --
  % the interval must lie inside the margin AND contain zero -- and this
  % figure's result is that it does the first and not the second. Pass 10 drew
  % the margin conjunct as the heaviest stroke on the page and the zero
  % conjunct as an unlabelled 0.5pt dotted rule standing on the panel's own x
  % origin, where it has no figure/ground of its own; at a 4x squint it
  % disappeared entirely. Naming it is what the figure's spec authorises ("if
  % the zero rule needs a name it takes 'zero' in \scriptsize roman ink at the
  % rule") and it is the whole of the fix: the rule is NOT promoted to the
  % threshold register, because S6 reserves black solid 0.90pt for a threshold
  % and S7 reserves this stroke for a coordinate that decides nothing -- which
  % zero, as a coordinate, is. What decides is whether the bar clears it, and
  % the bar clears it by 0.70cm.
  % The label is placed with an xshift on the rule's own x, as the row labels
  % are placed on their marks', so it is tied to the thing it names; the two
  % vertical rules of (b) now carry one \scriptsize label each, on one
  % baseline, one at each rule's top (S22).
  % PASS 10, X9: it no longer starts on the spine. Drawn 3.58 to 5.56 it was
  % collinear and continuous with the axis's left terminus, its hanging tick and
  % the panel's x origin, ran the full height of the data region, and matched
  % the threshold rule on the right for length -- so all three reviewers read
  % the pair as the two sides of a frame round the rows, which is the object
  % S12 exists to forbid. It now spans the three rows only, 0.28cm clear of the
  % spine and 0.42cm clear of the panel head, and reads as what it is: the
  % coordinate the interval bar's left end clears by 0.70cm.
  \draw[proj] (0.00,3.94) -- (0.00,5.34);
  \node[qlabz, xshift=3pt] at (0.00,5.46) {zero: no identity effect};

  % ROW ONE, the control: an open disc, white fill, 0.8pt quietink edge, which
  % is 4.11's own encoding for a comparator or a null (S9, CF5, W11).
  % PASS 10, X10: the gloss separator is a comma. It was an em dash, in the same
  % position after the token as row three's math minus, so one glyph class
  % carried "namely" in two rows and "subtract" in the third on a figure whose
  % whole subject is a subtraction. The minus in row three and in band (a) is
  % now the only dash inside a label anywhere on the drawing.
  % X12: the three row labels are placed with an xshift on the mark's own
  % coordinate rather than at a literal x, so each label is tied to its mark.
  % PASS 11, Y7: THE MARK-TO-LABEL GAP GOES 5pt -> 9pt IN ALL THREE ROWS.
  % SUPERSEDES X12's closing clause, which set it at 5pt. At 5pt the open
  % control disc stood 3.28pt from the first glyph of \texttt{permuted}, which
  % is LESS than Inconsolata's own character advance at this size (3.43pt),
  % while its outer diameter (3.6pt) is 11% larger than that face's x-height
  % (3.24pt measured at 200dpi) and its 0.80pt edge is of comparable weight to
  % the stems beside it. The one comparator mark on the drawing therefore set
  % as a lowercase letter o prefixing the word: at 12x it reads "o permuted".
  % Row two escaped only because its disc is filled and accent. At 9pt the gap
  % is 7.2pt, twice the mono advance, and the ring is a mark again.
  % A second fix was proposed for the same defect -- take the ring's edge to
  % ink, as 4.11(b) does -- and it is overruled at Y15 (a).
  \fill[white]  (\gb{0.0006409306049362111},5.24) circle (1.4pt);
  \draw[ctrl]   (\gb{0.0006409306049362111},5.24) circle (1.4pt);
  \node[rowq, xshift=9pt] at (\gb{0.0006409306049362111},5.24)
    {\texttt{permuted}, same construction, wrong drug};

  % ROW TWO, the arm under test.
  \fill[accent] (\gb{0.0019027651015006252},4.66) circle (1.4pt);
  \node[rowa, xshift=9pt] at (\gb{0.0019027651015006252},4.66)
    {\texttt{swap}, points at the drug scored};

  % ROW THREE, the measurement, and the only interval on the figure. The label
  % is anchored on the interval's upper end, which is the rightmost ink in the
  % row, so it clears the bar by the same 9pt the other two clear their discs.
  % PASS 11, Y8: THE VALUE COMES OFF THIS LABEL. It printed "$+0.0013$ nats"
  % 48.9pt (17.2mm) to the right of the disc whose value it is, ON THE ROW'S
  % OWN BASELINE, where position is value: the numeral stood at x = 0.00339 on
  % (b)'s axis, 2.7 times the number it printed and past the 0.0025 tick, on
  % the one panel whose job is to let a reader place the estimate between zero
  % and the 0.00642 margin. Rows one and two anchor at their marks and this
  % row anchored at its bar's end, so one panel carried two anchoring rules and
  % the one that misplaced its label was the one that printed a number. The
  % label is now the row's NAME only, which is the same string band (a) uses
  % for the same JSON key (S19), and the value is printed once, in (a), at a
  % label lifted off the data lane and leadered to its disc (Y2). After this
  % edit no numeral anywhere on the drawing stands on a data row's baseline.
  % PASS 11, Y9: 0.06cm MORE AIR UNDER ROW TWO. The bar on this row spans
  % 0.702 to 3.751cm and the two arm marks stand at 1.118 and 3.318, so the
  % interval on the DIFFERENCE horizontally contains both of its own operands
  % and the difference disc falls within 0.02cm of their midpoint by
  % arithmetic accident. A reader can be offered "a 95% interval that brackets
  % both arms", which is the opposite of what the row says. The row pitch is
  % now 0.58 / 0.64 rather than a flat 0.54, so the derived row sits off the
  % pair it is derived from; the label states the derivation; and the proposal
  % to delete the two arm marks outright is overruled at Y15 (b).
  \draw[bar] (\gb{0.0004024989353392094},4.02)
          -- (\gb{0.002150955667015771},4.02);
  \fill[accent] (\gb{0.0012618344965644142},4.02) circle (1.4pt);
  \node[rowa, xshift=9pt] at (\gb{0.002150955667015771},4.02)
    {\texttt{swap} \textminus{} \texttt{permuted}};

  % The magnified axis, its ticks landing on the same four x as band (a)'s so
  % that the two rules read as one rule at two scales (W13).
  % PASS 11, Y10: BAND (b) IS CAPPED AT ITS RIGHT END, AND ITS END IS COMPUTED.
  % SUPERSEDES W13's closing clause ("band two's right endpoint ... takes no
  % tick and no numeral"). W13 declined a terminal MARK on the ground that
  % (b)'s right end is named by (b)'s axis label rather than by a mark, and
  % that reasoning stands -- but what it left was 26.01pt (9.17mm) of bare
  % rule running past the last drawn tick to an end with nothing on it, which
  % is the same rule-end-to-last-tick defect S13 measures at 8.1mm in
  % fig:comparators row (c). Band (a) has the identical overhang and closes it
  % with a 0.22cm ink cap; band (b) now closes it with the same cap, at the
  % same x, in the same weight and colour, so the two rules terminate
  % identically at the end where the overhang is and read as one rule at two
  % scales. The cap is axis furniture -- the rule's own end -- and not a data
  % mark, so S20 is untouched; the proposal to hang a value on it is overruled
  % at Y15 (e). NO LEFT cap is added in (b): a 0.22cm ink stub rising from the
  % spine at x = 0 would stand 0.14cm under the dotted zero rule and collinear
  % with it, which is exactly the frame-edge reading X9 removed.
  \draw[spine] (0.00,3.68) -- (\gb{0.008025007399946877},3.68);
  \draw[spine] (\gb{0.008025007399946877},3.68)
            -- (\gb{0.008025007399946877},3.90);
  % PASS 11, Y11: the middle numeral prints 0.0050 and not 0.005. Band (a)'s
  % three numerals carry two decimals each and band (b)'s carry four, and the
  % one visible proof that (b) is (a) divided by eight is that the two lanes
  % put congruent numerals on the same four x. A dropped place broke the
  % column of trailing digits that carries it. The VALUE is unchanged; this is
  % a print string, which is why the \foreach takes the value/label pair form.
  \foreach \v/\t in {0.0025/0.0025, 0.005/0.0050, 0.0075/0.0075}{
    \draw[tick] (\gb{\v},3.68) -- (\gb{\v},3.58);
    \node[num] at (\gb{\v},3.52) {\t};}
  \draw[tick]  (0.00,3.68) -- (0.00,3.58);
  \node[numL] at (0.00,3.52) {0};

  \node[axlab] at (0.00,2.92) {preference shift, nats};
  % PASS 10, X14: "magnified 8x" comes off this line. The panel head 3.2cm
  % above already carries it, and the exactness of the eight is caption
  % material; what the axis-label line has to carry is the relation to the axis
  % above, which is what is left.
  \node[scope] at (0.00,2.54) {the first eighth of the axis above};

  %% ==== THE ONE INTERVAL DECLARATION, at the foot (S17, CF8) =============
  %% PASS 10, X2. The pass-9 second line read "No interval is returned for
  %% permuted, for swap, or for (a)." That was false, and false about the one
  %% mark on the figure that does have an interval: band (a)'s disc IS
  %% diff_perm, and layers.12.swap.ci_perm is an interval on diff_perm. In this
  %% document "returned" is what the instrument does, so the sentence told a
  %% reader that band (a)'s estimate has no uncertainty available and, by
  %% implication, that (a) and (b) draw different quantities. The declaration
  %% now separates non-existence from non-drawing: the arms have no stored
  %% interval, the difference has one, and it is drawn once.
  %% Leading is 0.28cm here, as everywhere else in this figure. At pass 9 the
  %% two lines were set 0.24cm apart, tighter than solid for 8pt type, and the
  %% descender of "bootstrap" met the dot of the "i" below it.
  %% PASS 11, Y12: 0.28cm -> 0.30cm between the two lines, and 0.46cm from the
  %% axis-label block above rather than 0.42cm. 0.28cm is exactly solid for
  %% 7.97pt type and the two lines' boxes met; the block also has to stand off
  %% the axis-label block by more than that block's own 0.38cm internal
  %% leading, or the four lines read as one four-line paragraph.
  %% PASS 11: six characters come out of these two lines, and nothing else, to
  %% keep the figure under S15's 5.0 characters per 100mm^2 after "of seven",
  %% "clean" and "zero" went in. Both sentences say exactly what they said.
  \node[scope] at (0.00,2.08) {Intervals: (b) 95\% bootstrap over 58 ordered pair clusters, on the difference, drawn once.};
  \node[scope] at (0.00,1.78) {None is stored for \texttt{permuted} or \texttt{swap} alone, and (a) carries none.};

  % The bounding box IS the body block; nothing may stick out of it.
  % PASS 11, Y13: the floor goes 1.72 -> 1.68 to keep the foot block's
  % descenders inside. This \path is the one rounded literal left in the file
  % and it carries no value: 13.9957 is a box corner, not a coordinate any
  % mark is drawn at, and every data x is computed by \ga or \gb from the
  % full-precision float (Y3).
  \path (0.00,1.68) rectangle (13.9957,10.18);
\end{tikzpicture}
%
% ---------------------------------------------------------------------------
% READY TO PASTE into Sections/04b-representation.tex, at sec:res-mechanism,
% immediately after the verdict paragraph (04b:332-342) or, if the page breaks
% badly, after the equivalence-criterion paragraph (04b:412-419). The same block
% was staged, at pass 9, in a figs/_PENDING_FLOATS.tex that does not exist;
% PASS 10 supersedes that: the float wrapper is HERE and the caption is in
% figs/_PENDING_CAPTIONS.tex, which does exist as of this pass (X15).
%
% PASS 9 NOTE. The caption below is the PROPOSED replacement and it is not yet
% in Sections/. No builder on this pass may edit Sections/04b-representation.tex,
% so the same text is written into figs/_PENDING_CAPTIONS.tex keyed by
% \label{fig:materiality} for a human to integrate. Until that happens the
% caption printed in main.pdf is the pass-8 one, and three of its clauses are
% false of the rebuilt drawing -- see CAPTION NOTE, PASS 9 at the foot.
%
% \begin{figure}[htbp]
% \centering
% % ===========================================================================
% materiality.tex -- fig:materiality. The one surviving identity effect drawn
% twice on one nats ruler: once against the whole preference gap it is divided
% by, where it is a 3.6mm smudge on a 133mm axis, and once at exactly eight
% times, where the swap arm, the permuted control and their difference separate.
%
% FIGURE BODY ONLY. No preamble, no document environment, no float, no caption.
% \input it from inside a figure float in Sections/04b-representation.tex at
% sec:res-mechanism. The float wrapper and caption sit at the foot of this file,
% commented out, ready to paste -- the staging convention of figs/two-scales.tex.
% Compiles against thesis_v2/preamble.tex and nothing else: no \includegraphics,
% no external data at draw time, no TikZ library beyond the ones preamble.tex
% already loads (this file uses decorations.pathreplacing and arrows.meta only).
%
% WHY THIS FIGURE EXISTS. Section 4.5 reports every effect "as a fraction of
% that model's own clean preference gap" (04b:396-398) and NO figure anywhere in
% the document shows that gap. fig:family draws the numerator on a diverging
% colour scale; fig:comparators row C draws it as a fraction. The denominator is
% never drawn, and neither is the subtraction that produces the numerator. This
% figure's subject is the denominator, and the subtraction.
%
% EVERY NUMBER DRAWN, AND WHERE IT COMES FROM
% -------------------------------------------
% ONE arm, ONE layer, ONE rung: the residual arm, layer 12, alpha = 0.5. Source
% for every coordinate: RESULTS_cluster/probe_arm_residual.json, all under the
% key path layers.12.swap. Sibling ledger of the exact values drawn, at full
% machine precision, with key paths: thesis_v2/figs/materiality.json.
%
%   quantity                key path              value drawn (full precision)
%   clean preference gap    .clean_ab_gap         0.06420005919957501
%   10% materiality margin  .equiv_margin         0.006420005919957502
%   permuted arm            .effect_perm          0.0006409306049362111
%   swap arm                .effect_swap          0.0019027651015006252
%   swap - permuted         .diff_perm            0.0012618344965644142
%   its 95% interval        .ci_perm     [0.0004024989353392094,
%                                         0.002150955667015771]
%
% One further quantity is COMPUTED, and is declared as computed in the ledger:
%   one eighth of the gap   .clean_ab_gap / 8     0.008025007399946877
% It is band two's full scale and band one's magnifier bracket, and it is the
% only reason the word "eighth" appears anywhere on the figure.
%
% VERIFIED ON BUILD, not assumed:
%   * .equiv_margin / .clean_ab_gap = 0.1 exactly (config.equiv_frac = 0.1), so
%     the two dashed rules stand at one tenth of band one's full scale.
%   * .effect_swap - .effect_perm = 0.001261834496564414 against .diff_perm
%     = 0.0012618344965644142. They agree to fifteen significant figures and
%     differ in the last bit only, from summation order; the stored .diff_perm
%     is what is drawn. The subtraction the figure draws is therefore the
%     subtraction the pipeline performed, and this was checked, not asserted.
%   * .diff_perm / .clean_ab_gap = 0.019654724813287512 and
%     .ci_perm / .clean_ab_gap = [0.0062694480403512436, 0.03350395145788292],
%     reproducing .diff_perm_norm and .ci_perm_norm to sixteen digits. That is
%     the cross-check that this figure's nats and the thesis's fractions are the
%     same measurement; the normalised triple is printed in the caption and is
%     NOT drawn.
%   * 0.06420005919957501 / 0.0012618344965644142 = 50.87835161772162 and
%     0.006420005919957502 / 0.0012618344965644142 = 5.087835161772163. Both
%     COMPUTED, declared as computed in the ledger; they are the "one fiftieth"
%     and "one fifth" of the closing line.
%
% THE ci_perm TRAP, and why the header names both draws
% -----------------------------------------------------
% The interval drawn is layers.12.swap.ci_perm, the TOP-LEVEL key:
%   [0.0004024989353392094, 0.002150955667015771], p_perm = 0.001,
%   normalising to the thesis's [+0.0063, +0.0335].
% layers.12.swap.ladder[alpha = 0.5].ci_perm stores a SECOND bootstrap draw of
% the same effect,
%   [0.0003982749538768512, 0.002137823621073976], p_perm = 0.004,
% with the point estimate identical to sixteen digits. It normalises to
% [+0.0062, +0.0333] and would print one digit off tab:identity while silently
% swapping the surviving p for the non-surviving one -- 0.001 and 0.004 lie
% either side of the 0.05/26 admitted at rank one (04b:526-531). A builder who
% reads the interval from the ladder because that is where the alphas live makes
% exactly this error. The ladder draw is recorded in materiality.json under
% not_drawn_second_bootstrap_draw and is not drawn here.
%
% AGREEMENT WITH THE TEXT (checked line by line, not assumed):
%   04b:340-341  "\ci{+0.0197}{+0.0063}{+0.0335} of the preference gap the model
%                already has. One to two per cent, against a $10\%$ margin
%                pre-registered for calling an effect material."
%                -> the drawn nats normalise to exactly that triple; the margin
%                   is drawn twice, once per band.                          OK
%   04b:443      "In raw nats the rungs read $+0.0013$" (first rung = the
%                headline) -> row three's label prints $+0.0013$.           OK
%   04b:500-502  "60 rows in 58 pair clusters over 24 anchor drugs"
%                -> row three's second label line prints 58 pair clusters;
%                   60 and 24 are in the caption.                           OK
%   04b:507-513  "The interval must lie inside a pre-registered margin ... and
%                contain zero. ... Under textbook TOST the headline ... would be
%                returned as equivalent, read as a null; this rule returns a
%                directional effect smaller than the margin."
%                -> the two closing ink lines track this sentence, which is the
%                   section's subtlest point and is drawn nowhere else.     OK
%   04b:568      "a fifth or less of the $10\%$ margin"
%                -> the closing accent line prints "one fifth of the margin".OK
%
% DESIGN DECISIONS, each with its reason
% --------------------------------------
%  1. TEXTWIDTH, NOT FULLWIDTH. sec:res-mechanism already carries two fullwidth
%     objects, fig:family and tab:identity, and a fullwidth float forbids a
%     margin note on its page (preamble.tex:522-523). Dropping fullwidth buys
%     this page's margin notes back -- the trade ladder.tex and two-scales.tex
%     both made deliberately. Drawn to 141.9mm against the 142mm measure; the
%     observed band in the corpus is 139.6-141.9mm.
%  2. A RAW-NATS AXIS, WHICH IS LEGITIMATE HERE AND NOWHERE ACROSS ARMS. The
%     five arms emit different target formats and their clean gaps span nearly
%     fivefold, so a single nats rule laid across arms would be one
%     arm's ruler drawn over five arms' data. RECOMPUTED, because the figure
%     brief quoted "4.4x" for this and 4.4 is not the span: over the five
%     probe_arm_*.json files clean_ab_gap runs 0.06420005919957501 (this arm,
%     this layer) to 0.3138025708635225 (consensus, layer 4), a ratio of
%     4.887886, and 4.410308 is the narrower consensus-layer-2 to
%     residual-layer-12 ratio. The caption prints the two endpoint values and
%     the word "nearly fivefold" so a reader can check it. Everything on this
%     figure is one
%     arm at one depth at one rung, so the raw scale is the arm's own and the
%     10% margin is a single rule with a single meaning. fig:family's lower
%     strip is a single-arm object for the same reason.
%  3. THE SAME QUANTITY DRAWN TWICE, NOT TWO QUANTITIES, AND THE MAGNIFICATION
%     IS EXACTLY EIGHT. Band two's full scale is .clean_ab_gap / 8, so the panel
%     is exactly the first eighth of the axis above it and the factor printed on
%     its axis is exactly 8 and not a rounded number. This was arrived at by
%     correction, and the arithmetic is worth recording. The first draft ran
%     band two to 0.0075 nats and called the panel "the first twelfth" at
%     "8.56x": 0.0075 / 0.06420006 = 0.11682, which is one 8.56th and not one
%     twelfth, so the panel's own title was false by a factor of 1.4 and the
%     magnification was a number no reader could check. A true twelfth cannot be
%     drawn at all -- .clean_ab_gap / 12 = 0.005350 nats is SMALLER than the
%     0.006420 margin, so the margin would fall off the right-hand end of the
%     panel and the figure would lose its subject. One eighth is the nearest
%     nameable fraction that still contains the margin, and at that scale the
%     margin sits at 80% of the panel's width, which is where it does the most
%     work. Nothing else changed: every drawn quantity is the same number.
%  4. BAND ONE IS SPARE, BUT ITS MEASUREMENT IS NOW NAMED. Four tick numerals,
%     not seven: the band's job is one impression -- 3.6mm of ink 0.8mm from the
%     origin on a 133mm axis -- and a crowded band invites reading rather than
%     seeing. Every other label lives in band two.
%     REPAIRED, PASS 8, and both reviewers called the old state blocking. The
%     measurement carried NO label at all, on the theory that the magnifier
%     would carry it; its only pointer was the italic "the effect and its
%     interval, magnified below" set 11cm to its right at the far end of a
%     nearly horizontal grey leader, against the document's own rule that a
%     label more than ~0.3cm from its mark gets a leader. One reviewer could not
%     tell whether the mark was a mark, a tick or a printing artefact; the other
%     read it as a stray em-dash. It now carries an accent label, "the effect",
%     centred on the estimate dot 0.22cm above it, and a 0.25pt stem from the
%     dot down to the spine so the mark is anchored to a coordinate rather than
%     floating. The label is two words and not "the effect, +0.0013 nats"
%     because the measure between the axis origin at 0.55 and the margin rule at
%     1.88 holds 1.33cm at \scriptsize while the longer form is 2.55cm and would
%     have crossed the rule; row three of band two prints the value in full,
%     3.5cm below. The italic pointer is deleted with the leader it rode on.
%  5. THE MEASUREMENT SITS 0.22cm ABOVE BAND ONE'S SPINE, not on it. On the
%     spine, a 3.6mm accent bar 0.8mm from the origin reads as a thickening of
%     the axis or as a printing defect. Lifted clear, with its dot over a 2.1pt
%     white halo, it reads as a mark. Checked in the 200dpi render and again in
%     a 550dpi crop: the bar, the dot and the spine are three separable things.
%  6. THE MAGNIFIER'S SOURCE IS A NAMED SPAN, NOT TWO BARE COORDINATES.
%     REPAIRED, PASS 8, reported independently by both reviewers. It used to be
%     a vertical drop at x = 0.55, a second vertical at x = 2.2125, and a line
%     from (2.2125,7.20) to (13.85,6.75) -- 11.6cm horizontal over 0.45cm
%     vertical, about 2 degrees. Read cold, the vertical was a drop from the
%     axis origin and the long line was a rule under band one, and the italic
%     that rode on it annotated band one's axis rather than announcing a
%     magnification. slab.tex's version works because both its leaders leave the
%     magnified REGION at a visible angle, and the fix is to make the region an
%     object first. A rulegrey brace, the corpus's own idiom, now spans exactly
%     the magnified eighth on band one -- \ga{0} = 0.55 to
%     \ga{clean_ab_gap/8} = 2.2125 -- and both leaders leave its two ends, one
%     at 66 degrees and one shallow. THE SHALLOW ONE CANNOT BE FIXED and should
%     not be pretended away: a 1.66cm span is being opened into a 13.90cm panel,
%     so one of the two leaders must run nearly flat whatever is done. What
%     changed is that it now leaves a bracket rather than a bare coordinate, so
%     the assembly reads as source-and-enlargement rather than as a crooked
%     rule. The leaders now land on the panel's own top corners (0.25 and
%     14.15) rather than on band two's data extremes: with the brace naming the
%     source, the leaders' job is to point at the PANEL, and corner-to-corner is
%     what a magnifier looks like. The bracket contains both the effect AND the
%     margin rule, which is itself the point: the whole of section 4.5's
%     argument happens inside the first eighth of the gap.
%  7. THREE ROWS IN THE ORDER PERMUTED, SWAP, DIFFERENCE. Control first,
%     treatment second, measurement third, so the subtraction reads downward in
%     the order the sentence "swap minus permuted" is spoken.
%  8. AN INTERVAL IS DRAWN ON ROW THREE AND ONLY THERE. No interval is stored
%     for either arm alone -- .ci_perm is an interval on the DIFFERENCE -- so
%     rows one and two are bare solid dots. In this document a bare solid dot
%     with no bar is how the absence of an interval is said (comparators.tex
%     changed exactly that mark from a hollow ring to a solid fill, because the
%     absent bar is what carries the meaning). Row three's bar is the document's
%     own convention: a 95% cluster bootstrap over ordered (A,B) pairs.
%  9. THE SIX-MARK TRANSLATION DEVICE IS DELETED, AND THE ROWS CARRY THE
%     CORPUS'S ROW APPARATUS INSTEAD.
%     PASS 8, and the spec pre-authorised this deletion. The device was a
%     hairline from each of the two dots down to the ends of a brace, that
%     brace, a leftward arrow sliding it to the origin, a second brace of equal
%     length from the origin, and a hairline dropping onto row three's dot: five
%     grey structural marks against three measured ones, carrying the identity
%     swap - permuted = the interval's point estimate. It is arithmetically
%     exact -- \gb{effect_swap} - \gb{effect_perm} = \gb{diff_perm} - 0.55 =
%     2.09126 -- and it did not read. Neither brace was labelled with a value,
%     so 0.00190 - 0.00064 = 0.00126 was not recoverable from the marks; the
%     arrow's head landed in white space with no mark at it; and the one
%     sentence that explained the chain, "the same span, measured from zero",
%     sat about 4cm to the right of the arrow it named, at a height between the
%     two braces, attached to nothing. Both reviewers reported having to
%     magnify 3x and reverse-engineer the arithmetic, and one noted that nothing
%     in the construction was accent, so the eye never entered it. Row three's
%     own label already STATES the identity in words -- "swap - permuted" -- so
%     the deletion costs nothing and removes five hairlines. The house rule it
%     follows: an identity must be readable from the marks or not drawn.
%     What replaces it is what fig:comparators gives every row: a rulegrey!55
%     0.3pt divider under each of the three rows, spanning the panel's inner
%     width. Before, the three quantities sat at three heights with no dividers,
%     no stubs and no alignment cues, so the vertical read as a second dimension
%     when it is only a stacking order -- and the higher dot happens to be the
%     SMALLER value, which invites a reader to look for an ordering that does
%     not exist. Rows moved from 5.52/4.94/3.66 to 5.40/4.72/4.04 and the
%     margin brace from 3.10 to 3.30 to take up the slack the deletion left.
% 10. THE MARGIN IS DRAWN TWICE AND BRACED ONCE. On band one it stands 1.33cm
%     from the origin; on band two, 10.64cm from it, and the 7.08cm of clear
%     space between the interval's upper end and the margin rule is the figure's
%     answer to "is it material". The brace under band two names what the margin
%     is a tenth OF -- 0.00642 nats out of 0.0642 -- which no existing figure
%     does.
% 11. THE FRAME IS NAMED ON THE DRAWING and not only in the caption: a quietink
%     line at the foot says which arm, which layer and which rung. This document
%     runs five arms whose clean preference gaps differ by 4.89x, so a raw-nats
%     axis that does not say whose nats they are is under-labelled in exactly
%     the way gate.tex's fifth fix names. \ttfamily on \texttt{residual} because
%     that is how the results files name the arm (ladder.tex's register key).
% 12. WHAT IS NOT DRAWN. The isotropic arm (.effect_iso = 0.0010863731109609438)
%     is stored at this cell and is omitted: 04b:377-380 keeps it "only as a
%     damage diagnostic", it is not the control the headline is measured
%     against, and a third bare dot between rows one and two would read as a
%     second comparison. The alpha ladder is omitted because fig:family's lower
%     strip already draws it and this figure is one rung. p, q, n and the
%     normalised triple are in the caption, not on the drawing.
% 13. REJECTED ALTERNATIVE: a broken axis. Drawing the whole gap and the first
%     eighth on one axis with a break would have saved 3.5cm of height, and it
%     would have destroyed the one thing the figure is for -- the reader has to
%     see the unbroken 133mm and the 3.6mm on it in the same glance before the
%     magnification means anything. two-scales.tex breaks an axis and says in
%     its own header that the break is admissible only because the off-scale
%     value is named in words; here the off-scale value IS the subject.
% 14. REJECTED ALTERNATIVE: bars from zero for the two arm effects. A bar length
%     would have made effect_swap look three times effect_perm, which is true of
%     the numbers and is not the measurement: the measurement is their
%     difference with an interval, and neither arm alone is quoted anywhere in
%     the thesis. Dots keep the arms as positions and the difference as the only
%     thing given a bar.
% 15. REJECTED ALTERNATIVE: naming the 0.0642 axis "chance" or drawing a 0.50
%     rule. There is no chance value on this figure. The reference here is the
%     model's own clean preference gap, an arm-specific measured quantity, so it
%     is drawn as the axis itself and named at the axis's right end rather than
%     as a rule inside it. Noted because the document's two conflicting habits
%     for a chance rule -- solid black in style_v2.chance_line(), 0.5pt densely
%     dotted ink in ladder.tex -- do not apply here, and their absence should be
%     deliberate rather than accidental.
%
% LAYOUT NOTES, all forced by a build
% -----------------------------------
%   * TWO MAPS, both braced pgfmath groups. A shim may not be added outside the
%     braces -- \gb{v}+0.05 is not a coordinate -- so every label offset here is
%     an explicit literal or an xshift on the node (the gate.tex trap). Both
%     maps are used directly as coordinate components, (\gb{v},y), which is
%     comparators.tex's idiom and keeps every data x exact rather than rounded
%     into the source.
%       \ga{0.06420005919957501}  = 13.850005   (band one, the whole gap)
%       \gb{0.008025007399946877} = 13.850005   (band two, one eighth of it)
%     The two right-hand ends coincide to 5 decimal places by construction:
%     1657.32 = 8 x 207.1650 exactly.
%   * Tick numerals use the \foreach \v/\t value/label pair form. \foreach \v
%     alone prints the raw token, so 0.020 would set as 0.02 in one place and
%     0.0 in another; the pair form pins the printed string.
%   * inner sep = 1pt ON THE WHOLE PICTURE, as gate.tex sets it. The default
%     0.3333em is 3.65pt at the ambient \small, and at pass 2 that invisible
%     padding -- not the ink -- was what collided: the band-one word tag's box
%     and the margin label's box overlapped while their glyphs were 1.2mm
%     apart. With inner sep pinned, every clearance below is box arithmetic
%     rather than judgement. The consequence is that a label placed at
%     mark + 0.30 would sit 0.10cm closer to its mark than comparators.tex's,
%     so every row label here is placed at mark + 0.40, which puts the glyph
%     0.39cm clear of the marker edge -- comparators' measured value.
%   * TICK NUMERALS CARRY fill=white AND ARE DRAWN LAST IN THEIR BAND. Pass 1
%     ran the magnifier's left leader straight through band one's "0", because
%     band one's origin and the magnifier's left edge are the same x = 0.55.
%     The white box interrupts the rule across the numeral instead --
%     comparators.tex's idiom for its zero rule, done with a box rather than by
%     hand because the numerals are drawn in a \foreach. Drawing order matters:
%     spine, then magnifier, then numerals.
%   * THE PANEL IS 3mm WIDER THAN ITS DATA ON EACH SIDE, (0.25,1.50) to
%     (14.15,6.75) around data running 0.55 to 13.85. At pass 2 the panel ran
%     exactly 0.55 to 13.85 and its left dashed border was drawn at the same x
%     as band two's dotted zero rule; the two were indistinguishable in the
%     600dpi crop, so the figure's zero coordinate was reading as the frame of
%     its own panel. Widening the frame separates them by 3mm. PASS 8 re-aimed
%     the magnifier's leaders at the panel's own top corners (0.25 and 14.15),
%     which is the opposite of what pass 2 did and is right for the same reason
%     the brace is right: with the brace naming the SOURCE span exactly, the
%     leaders' job is to point at the PANEL, and corner-to-corner is what a
%     magnifier looks like.
%   * BAND TWO'S WORD TAG AND ITS MARGIN LABEL SIT INSIDE THE PANEL, on one
%     baseline at y = 6.44, and the axis title sits inside it too. At pass 2 the
%     tag, the margin label, "no movement" and the axis title all straddled the
%     dashed border. Anything that names something inside the panel is now
%     inside the panel.
%   * The band-two tag is anchored at x = 0.75 and not at x = 0.00, because the
%     magnifier's left leader is a vertical at x = 0.55.
%   * PASS 8: the translation device and its label are deleted (decision 9),
%     which took five hairlines and one floating sentence out of the panel.
%     Rows moved up to take the slack -- 5.52/4.94/3.66 became 5.40/4.72/4.04,
%     and the margin brace 3.10 became 3.30 -- and each row gained a rulegrey!55
%     0.3pt divider under it, spanning 0.35 to 14.05, which is 0.10cm inside the
%     panel's dashed frame at both ends.
%   * PASS 8, THE MARGIN LABEL STAYS ON THE 8.44 ROW. The build that moved it up
%     to the tag's row put it at x = 1.96 and "THE WHOLE GAP" measures 1.95cm of
%     \sffamily\scriptsize\bfseries from x = 0, so the two ran together as
%     "THE WHOLE GAP0% materiality margin" in the render. Both labels sit on the
%     8.44 row instead and the dashed margin rule at x = 1.88 separates them:
%     "the effect" is centred on the estimate dot and ends at 1.29, the rule is
%     at 1.88, the margin label starts at 1.98. The rule stops at 8.56, which is
%     0.09cm under the tag's descender line and is why it may not be carried
%     higher to meet a label on the tag's row.
%   * BAND ONE'S RIGHT END is a terminal tick RISING from the spine, 0.30cm
%     tall against the 0.10cm hanging axis ticks, with its two label lines
%     right-aligned on the same x = 13.85 so the tick reads as their anchor. At
%     pass 3 that rule ran from the spine up to the second label's baseline and
%     read as a bracket enclosing the label block.
%   * The closing prose was three lines at pass 2 and is two at pass 3; the
%     28pt that bought went into the clearances above.
%   * Every prose node is a single \scriptsize line anchored base west. No
%     `align=` and no `text width=` anywhere: those build a minipage whose first
%     line follows the AMBIENT \baselineskip, so the node's ink moves with the
%     prose before the float (frames.tex measured a 3.7pt drift).
%   * The bounding box is pinned with a bare \path so the measure is exact and
%     nothing is clipped: (0.0000,0.30) to (14.1900,8.98) = 141.9 x 86.8mm.
%
% TYPE. Ambient \small on the picture; every node sets \scriptsize = 8pt, which
% is 0.67 of the 12pt body and below the 11pt caption. \ttfamily on the two arm
% names and on the arm named at the foot, because that is how the results files
% name them. Nothing is scaled: 1cm here is 1cm on the page.
% ===========================================================================
%
% ===========================================================================
% PASS 9 -- 2026-08-18. REBUILD AGAINST THE STYLE CONTRACT AND THE SUPERVISOR'S
% VERDICT ("no definition on the x axis for what we are measuring"; "make it
% look more professional, right now it looks a bit sloppy and childish").
% NOTHING ABOVE IS DELETED. The design rationale of passes 1-8 is the audit
% trail and stays; what follows supersedes, EXPLICITLY AND BY NAME, each clause
% of it that this pass made false. The reference standard adopted is fig 4.11
% (figs/fig-mechanism.pdf, printed page 51) and, for the threshold, fig 4.13
% (figs/family.pdf), both of which were rendered and read side by side with this
% figure during the build.
%
% NO NUMERIC VALUE CHANGED IN THIS PASS. Every quantity drawn or printed is the
% same float from the same key path as in pass 8: clean_ab_gap 0.06420005919957501,
% equiv_margin 0.006420005919957502, effect_perm 0.0006409306049362111,
% effect_swap 0.0019027651015006252, diff_perm 0.0012618344965644142,
% ci_perm [0.0004024989353392094, 0.002150955667015771], n_pair_clusters 58.
% The one printed numeral that left the drawing is "one fiftieth / one fifth"
% (a ratio in words), and it went to the caption, where it already was.
%
% W1. SUPERSEDES LAYOUT NOTE "TWO MAPS" AND THE MAP CONSTANTS RECORDED THERE.
%     Both maps lose their 0.55cm left inset and gain a longer barrel:
%         old   \ga#1 = {0.55+(#1)*207.1650}   \gb#1 = {0.55+(#1)*1657.3200}
%         new   \ga#1 = {(#1)*215.7500}        \gb#1 = {(#1)*1726.0000}
%     Two things are preserved exactly and were recomputed on the build rather
%     than assumed: 1726.0000 = 8 x 215.7500 EXACTLY, so the magnification is
%     still exactly eight and not a rounded number (this is the correction made
%     at pass 4 and it survives untouched); and both bands' right-hand ends
%     still coincide, now at \ga{0.06420005919957501} = \gb{0.008025007399946877}
%     = 13.851162772308308cm. Every data x is still computed inside pgfmath from
%     the full-precision float, so no value is rounded into the source.
%     WHY: (i) S13, one x origin per panel. With the old map the axis began at
%     0.55 while every text block began at 0.00, so the panel had two left edges
%     0.55cm apart and the left column measured three (58.0, 63.9, 70.4pt).
%     Now the axis's left end, the panel titles, both axis-label blocks, the
%     span bracket and the foot line all sit on x = 0.0000. (ii) S5 puts the
%     axis-label line under the tick numerals and LEFT-ALIGNED TO THE AXIS'S
%     LEFT END; that is only the same thing as the panel origin if the two
%     coincide. (iii) S14, measure. Deleting the dashed enclosure (W4) took away
%     the object that used to set the width, and a 133.0mm barrel inside a
%     textwidth figure measures about 134mm of content, under the 138mm floor.
%     215.75 puts the barrel at 138.51mm and the content bbox at 139.6mm, inside
%     138-142, with the widest object the axis itself and the data region 100%
%     of the drawn width.
%     The old map is kept in materiality.json beside the new one, per S25.
%
% W2. AXIS DEFINITION, WHICH IS THE SUPERVISOR'S FIRST COMPLAINT AND THE REASON
%     THIS PASS EXISTS. Pass 8 named neither axis. Band one carried bare 0 /
%     0.02 / 0.04 / 0.06 with no quantity and no unit anywhere on it -- the unit
%     reached the reader only through the terminal label -- and band two carried
%     "nats, magnified 8x from the axis above", which names a unit and a
%     transformation but no quantity. 4.11 sets "probe accuracy", "descriptive
%     subspace fraction", "KL from the clean forward pass" and "layer" in
%     axis-label position five times over. Each band now carries, directly under
%     its tick numerals and left-aligned to its axis's left end:
%         line 1, \footnotesize roman ink:  Preference shift, nats
%         line 2, \scriptsize quietink:     (a) residual arm, layer 12, alpha =
%                                               0.5; the full width of this rule
%                                               is the model's clean preference
%                                               gap
%                                           (b) the first eighth of the axis
%                                               above, magnified 8x
%     The quantity is repeated deliberately: two scales, each self-describing.
%     "Preference shift" is the section-4.5 registry name (S19), so the axis, the
%     caption and fig:family's lower strip now use one string for one quantity.
%
% W3. THE THRESHOLD. SUPERSEDES DESIGN DECISION 10 AND THE marginrule STYLE.
%     The 10% pre-registered materiality margin was accent 0.60pt dashed here,
%     ink 0.50pt densely dotted in fig:comparators row C, and solid black 0.90pt
%     in fig:family's lower strip: one pre-registered threshold, three renderings
%     inside nine pages. It is now BLACK, SOLID, 0.90pt in both bands, labelled
%     \scriptsize roman black, which is style_v2.chance_line()'s own hard-coded
%     register and fig:family's. This is the highest-priority change in the set
%     and it is a correctness repair, not a preference: accent means the object
%     under test (S8), and a pre-registered bound is not under test.
%     Design decision 15's note -- that the document's two habits for a chance
%     rule "do not apply here" -- is now false as written and is superseded:
%     there is still no chance value on this figure, but the black-solid-0.9pt
%     register does apply, to the margin.
%
% W4. DELETED: THE DASHED ROUNDED ENCLOSURE round band two (390 x 149pt), with
%     the layout note that sized it 3mm wider than its data on each side. S12
%     forbids a dashed rectangle drawn round a panel outright, and it was this
%     figure's only rectangle. Band two now sits directly under band one in
%     4.13's idiom -- the magnified axis under the full axis, with the
%     magnification named in axis-label position -- which is also what fixes
%     band two's axis under W2. Two consequences recorded so they are not read
%     as accidents: the layout note about the panel's left border colliding with
%     band two's zero rule is moot, since there is no border; and the note that
%     band two's tag and margin label "sit inside the panel" is moot for the same
%     reason.
%
% W5. DELETED: THE FOUR PALE ROW SEPARATORS AND THE MARGIN BRACE. SUPERSEDES
%     DESIGN DECISION 9's replacement clause and the pass-8 note that added them.
%     Five near-identical thin horizontals stood inside 30mm of vertical space --
%     three rulegrey!55 row dividers, one rulegrey brace carrying "10% of the gap
%     = 0.00642 nats", plus the axis rule -- and the one that carried meaning was
%     indistinguishable from the ones that did not. The dividers go because the
%     rows no longer need them: with the enclosure gone and one x origin, the
%     three rows read as a stack. The brace goes because the same fact is now
%     carried without any stroke at all, by printing the margin's value in band
%     one's threshold label ("10% materiality margin, 0.00642 nats") beside band
%     one's terminal label ("the clean preference gap, 0.0642 nats"): the tenth
%     is then a division a reader performs on two printed numbers 8cm apart
%     rather than a brace they must trace.
%
% W6. THE INTERVAL IS DRAWN ONCE. SUPERSEDES DESIGN DECISION 4 and the pass-8
%     repair that labelled band one's bar. Pass 8 drew the same estimate as a
%     3.6mm bar-with-dot in band one AND as a 30mm bar-with-dot in band two,
%     while the caption said "it is the one bar on the figure"; a naive reader
%     counted two and spent time working out which two quantities they were.
%     Rule applied: an interval is drawn at the scale where it is legible and
%     nowhere else. Band one keeps a single filled accent disc for the estimate,
%     with no bar; band two carries the bar. Band one's dot keeps its label and
%     its leader, and the label is now the SAME STRING as band two's row three,
%     "swap - permuted", because S19 forbids two names for one JSON key and pass
%     8 had "the effect" here and "swap - permuted" there.
%
% W7. DELETED: THE TWO BANNERS, replaced by 4.11's panel heads. "THE WHOLE GAP"
%     and "THE FIRST EIGHTH, MAGNIFIED" were \sffamily\scriptsize\bfseries accent
%     caps -- the preamble's margin-apparatus face, which belongs to the margin
%     and not to a body figure (S2). They are now "(a) the whole preference gap"
%     and "(b) the first eighth, magnified 8x", \footnotesize roman ink,
%     lowercase, left-aligned on the panel origin, no colon and no bold. The
%     tag/.style is deleted rather than left unused.
%
% W8. DELETED: THE ITALIC "no movement". It sat unattached above the permuted
%     row and a naive reader took it for an assertion about that arm rather than
%     a name for the zero rule. The zero rule survives -- densely dotted quietink
%     0.5pt, which is exactly S7's reserved register for a coordinate that
%     decides nothing -- and carries no label, because its foot is the axis's own
%     "0" numeral. With it goes the last italic on the drawing: italic share is
%     now 0%, against 10% at pass 8 and 0% in 4.11.
%
% W9. DELETED: THE CLOSING PROSE BLOCK AND THE ACCENT APHORISM. The two ink lines
%     ("The interval excludes zero and lies inside the margin. Textbook TOST
%     returns that as equivalence and reads it as a null; this rule returns a
%     directional effect smaller than the margin.") and the accent italic ("one
%     fiftieth of the gap; one fifth of the margin.") are argument, and in this
%     document the argument lives in the caption (S3, S4, S18). At print the
%     reader met an 8pt caption, a blank, then the real 10.9pt caption -- two
%     captions, the smaller one first. Both lines are written into
%     figs/_PENDING_CAPTIONS.tex, keyed by \label, for a human to integrate;
%     the aphorism is already the caption's lede. The TOST asymmetry is the
%     section's subtlest point and it is NOT lost: it is caption material, and
%     the configuration it describes -- an interval that excludes zero and lies
%     wholly inside the margin -- is still visible in band two, which is the only
%     part of that sentence a drawing can carry.
%     SUPERSEDES the text_agreement_check on 04b:507-513 in the ledger, which
%     asserted the drawing restates that sentence. It no longer does.
%
% W10. THE MAGNIFIER. SUPERSEDES DESIGN DECISION 6 AND THE PASS-8 BRACE-AND-TWO-
%     LEADERS CONSTRUCTION. What is kept: the source span is a NAMED OBJECT and
%     not two bare coordinates. What is deleted: both leaders, including the
%     shallow one that pass 8's own header admitted "CANNOT BE FIXED" and that
%     the naive reader read as a trend line before reading it as a leader.
%     The reasoning, recorded because a later builder will be tempted to put them
%     back. A leader pair from the ends of band one's [0, gap/8] segment to the
%     ends of band two's rule is geometrically forced to be asymmetric -- a
%     17.3mm span opening into a 138.5mm rule -- so the right-hand leader runs at
%     under 4 degrees whatever is done, and 4 degrees of accent-free hairline
%     running the full measure is a trend line to every eye that meets it first.
%     Worse, it cannot be routed: S5 now requires band one's two-line axis-label
%     block directly under band one's tick numerals, which puts full-width text
%     in the only lane the leaders could cross. The magnification is instead
%     carried by two objects that cannot be misread: a 0.30pt rulegrey bracket on
%     band one spanning exactly \ga{0} = 0.0000 to \ga{clean_ab_gap/8} = 1.731395,
%     labelled at its own end ("the eighth magnified below"), and band two's
%     axis-label line 2, which names the magnification in axis-label position
%     where S5 puts it. The bracket's two ends stand at the same x as band two's
%     rule origin, so the vertical alignment does the joining that a hairline was
%     doing badly.
%
% W11. MARK VOCABULARY (S9, CF5). SUPERSEDES DESIGN DECISION 8's claim that a
%     bare SOLID dot is how this document says "no interval". It is not the only
%     thing a dot has to say. The permuted arm is the CONTROL, and 4.11's own
%     encoding for a comparator or a null is an open disc, white fill, 0.8pt edge
%     in its colour (mfc='white' on the raw probe and on the random null). Row one
%     is therefore an open quietink disc; rows two and three and band one's
%     estimate are filled accent discs, all at radius 1.4pt, one radius on the
%     figure. The absence of an interval is now said where the contract puts it,
%     in the single \scriptsize interval declaration at the foot (S17), which
%     names the bootstrap construction for the panel that has a bar and says in
%     as many words that the two arms and band one carry none. The 0.8pt disc
%     edge is a MARK edge, not a rule weight; the figure's rules use exactly five
%     weights, 0.30 / 0.40 / 0.50 / 0.90 / 1.00pt.
%     TWO SINGLETONS ARE DECLARED HERE RATHER THAN REMOVED, as S10 permits. The
%     0.50pt densely dotted stroke is used once, for band two's zero rule, and is
%     a distinct kind of object -- S7 reserves that stroke for a coordinate that
%     decides nothing, and this figure has exactly one such coordinate. The
%     1.00pt stroke is used once, for the interval bar, and is a distinct kind of
%     object for the same reason: this figure stores exactly one interval.
%
% W12. TYPE (S1, S16). Two sizes above 7pt and one face: \footnotesize (10.0pt)
%     for the two panel heads, both axis-label line ones, and nothing else;
%     \scriptsize (8.0pt) for tick numerals, direct labels, scope lines, the
%     threshold labels and the interval declaration. No \sffamily, no \bfseries,
%     no \itshape anywhere. \ttfamily survives on the four arm tokens
%     (residual, permuted, swap) because that is how the results files name them
%     and S1 exempts the \texttt token name. Pass 8 printed four sizes and set
%     its largest type on a grey italic epigram; the largest type now belongs to
%     a panel head or an axis label, which is the whole of CF6.
%
% W13. TICK LANES MADE CONGRUENT (S13). Band one ticks at 0 / 0.02 / 0.04 / 0.06
%     land on x = 0.0000 / 4.3150 / 8.6300 / 12.9450. Band two's ticks were
%     chosen so that they land on THE SAME FOUR x: 0 / 0.0025 / 0.005 / 0.0075,
%     which is possible only because band two is exactly one eighth of band one
%     and is a visible consequence of the exact 8. Both rules share both
%     endpoints (0.0000 and 13.851163) and both first ticks sit on the left
%     endpoint. Band one alone carries a terminal tick at the right endpoint,
%     rising 0.20cm from the spine, because that endpoint is the clean preference
%     gap and is the figure's denominator; band two's right endpoint is one
%     eighth of it and is named by band two's axis label rather than by a mark,
%     so it takes no tick and no numeral -- no unlabelled mark is left on the
%     drawing (S20).
%
% W14. MEASURE AND HEIGHT, MEASURED ON THE FINAL BUILD, not estimated: content
%     bbox 140.34 x 83.82mm, ink 5.14% of pixels below 230 grey at 200dpi, 4.88
%     characters of figure text per 100mm^2, one page, zero Overfull and zero
%     Underfull boxes. Against pass 8's 141.2 x 85.9mm at 7.03% ink and against
%     the contract's 138-142mm measure, 85mm height cap, 6.0% ink ceiling and
%     5.0 characters/100mm^2 ceiling. Every millimetre of the height and every
%     point of the ink came from W4, W5, W8, W9 and W10 -- deletions -- and
%     nothing measured was removed to get them.
%
% W15. STALE CLAIMS ELSEWHERE IN THIS HEADER, SUPERSEDED BY NAME so that an
%     auditor reading top to bottom is not misled.
%     (i) The opening paragraph's "where it is a 3.6mm smudge on a 133mm axis"
%         is now false twice over: the axis is 139.96mm, and band (a) no longer
%         draws the interval at all, so what stands 2.75mm from the origin is a
%         single 1.4pt disc and not a 3.6mm bar (W6).
%     (ii) The file-header line "this file uses decorations.pathreplacing and
%         arrows.meta only" is now false in the harmless direction: this figure
%         uses NEITHER. The braces went with W5 and W10 and there was never an
%         arrowhead. It still loads no library preamble.tex does not.
%     (iii) DESIGN DECISION 1's "Drawn to 141.9mm against the 142mm measure" is
%         superseded by W14's 140.34mm. The decision itself -- textwidth and not
%         fullwidth, to buy the page's margin notes back -- stands unchanged.
%     (iv) DESIGN DECISION 3's arithmetic is intact and was recomputed on this
%         build; only the cm-per-nat constants moved (W1). The first draft's
%         "first twelfth at 8.56x" correction, and the reason a true twelfth
%         cannot be drawn at all, remain exactly as recorded.
%     (v) The AGREEMENT WITH THE TEXT entry for 04b:507-513 ("the two closing
%         ink lines track this sentence") is superseded by W9: those lines are
%         deleted from the drawing and staged for the caption. The entry for
%         04b:568 ("the closing accent line prints one fifth of the margin") is
%         superseded for the same reason. Both sentences remain TRUE of the
%         figure's subject; they are no longer drawn on it.
%
% W16. TWO DEVIATIONS FROM THE SPEC'S LITERAL AXIS-LABEL STRINGS, both forced by
%     other clauses of the same contract, both recorded rather than quietly
%     taken.
%     (i) Band (a)'s scope line is set as "residual arm, layer 12, alpha = 0.5"
%         and NOT as "residual arm, layer 12, alpha = 0.5; the full width of
%         this rule is the model's clean preference gap". The deleted clause
%         restates, in twelve words, the terminal label "the clean preference
%         gap, 0.0642 nats" that the spec's own keep-list requires at the axis's
%         right end, and S18 forbids a drawing from carrying one clause twice.
%         It also costs 70 characters, and at 70 more characters the figure
%         needs 89.8mm of height to stay under the 5.0 characters/100mm^2
%         ceiling, which is past the 85mm cap. The quantity and its unit are on
%         line one, which is what S5 actually requires; what came off is a
%         second statement of what the axis's right end already says.
%     (ii) Band (b)'s scope line is set verbatim as the spec gives it, "the
%         first eighth of the axis above, magnified 8x", and the 14 characters
%         that costs were bought by opening 3mm more air between the bands
%         rather than by cutting anything else.
%     Line one of both blocks is verbatim: "Preference shift, nats". Recorded
%     because 4.11 and 4.13 both set their axis labels entirely lowercase and
%     4.10 sets its capitalised, so the slate is split on case and the spec's
%     capitalisation was followed rather than guessed at.
% ===========================================================================
% ===========================================================================
% PASS 10 -- 2026-08-18. REPAIR PASS on the pass-9 rebuild. Three adversarial
% readers went over the built figure independently and returned thirty-three
% findings. NOTHING ABOVE IS DELETED: passes 1-9 are the audit trail and stay
% as written. Every clause of pass 9 that this pass made false, and every
% number pass 9 recorded that this file has never drawn, is superseded here BY
% NAME (S25). The findings that were OVERRULED are recorded under X16 with the
% reason, so that a later builder does not re-apply them blind.
%
% NO NUMERIC VALUE CHANGED IN THIS PASS EITHER. Every quantity drawn or printed
% is the same float from the same key path as at passes 8 and 9:
% clean_ab_gap 0.06420005919957501, equiv_margin 0.006420005919957502,
% effect_perm 0.0006409306049362111, effect_swap 0.0019027651015006252,
% diff_perm 0.0012618344965644142, ci_perm [0.0004024989353392094,
% 0.002150955667015771], n_pair_clusters 58, and the exact magnification 8.
% Both maps are untouched: \ga = (#1)*218.0000 and \gb = (#1)*1744.0000, with
% 1744.0000 = 8 x 218.0000 exactly. What moved is where labels sit, where two
% rules start, how long two others are, and what four sentences say.
%
% X1. STALE CONSTANTS IN THE PASS-9 BLOCK, CORRECTED BY NAME. W1, W10, W13 and
%     W14 record a map, a shared endpoint, a bracket end, four tick positions
%     and a measure that this file has never drawn. materiality.json has always
%     recorded the drawn ones correctly, and the drawing is not affected --
%     every data x is computed inside pgfmath from the full-precision float --
%     but S25 makes the header an audited artefact, so an auditor recomputing a
%     position from W1 must not be handed a figure that does not exist. The
%     corrections, each superseding the clause named:
%       W1  "new \ga#1 = {(#1)*215.7500}  \gb#1 = {(#1)*1726.0000}"
%           -> the drawn map is \ga#1 = {(#1)*218.0000} and
%              \gb#1 = {(#1)*1744.0000}. Both properties W1 claims survive and
%              were rechecked on this build: 1744.0000 = 8 x 218.0000 exactly
%              (8*218.0 == 1744.0 is True in IEEE double), and both bands' right
%              ends coincide.
%       W1  "= 13.851162772308308cm" for the shared right end
%           -> 13.995612905507352cm = 0.06420005919957501 x 218, which is what
%              is hard-coded at the spine, at the two end caps and at the
%              terminal label, and what materiality.json records.
%       W1  "215.75 puts the barrel at 138.51mm and the content bbox at 139.6mm"
%           -> the barrel is 139.956mm and the content bbox 139.96mm, measured
%              on this build. Both are inside S14's 138-142mm.
%       W10 "\ga{clean_ab_gap/8} = 1.731395" for the bracket's right end
%           -> 1.749451613188419cm = 0.008025007399946877 x 218.
%       W13 "land on x = 0.0000 / 4.3150 / 8.6300 / 12.9450"
%           -> 0.0000 / 4.3600 / 8.7200 / 13.0800, for both bands, which is
%              what makes the two tick lanes congruent.
%       W13 "Both rules share both endpoints (0.0000 and 13.851163)"
%           -> (0.0000 and 13.995613).
%       W14 "content bbox 140.34 x 83.82mm, ink 5.14%, 4.88 characters per
%           100mm^2" -> superseded by X17's measurement of THIS build.
%     No 215.75 build was ever produced; the pass-9 numbers are a draft's, not
%     a superseded state of the drawing, and they are recorded here as wrong
%     rather than as an earlier truth.
%
% X2. THE ONE INTERVAL DECLARATION SAID SOMETHING FALSE ABOUT THE INSTRUMENT,
%     AND SAID IT ABOUT THE ONE MARK THAT HAS AN INTERVAL. SUPERSEDES W11's
%     description of the declaration and the ledger's intervals.declaration_on_
%     the_drawing. Pass 9's second line read "No interval is returned for
%     permuted, for swap, or for (a)." Band (a)'s single disc IS diff_perm, and
%     layers.12.swap.ci_perm is an interval on diff_perm; it is the very
%     interval drawn 8x larger in (b). "Returned" is this document's word for
%     what the instrument stores (S17), so the sentence asserted that the
%     band-(a) estimate has no uncertainty available and implied that (a) and
%     (b) draw different quantities. The declaration now separates
%     non-existence from non-drawing:
%       "95% bootstrap over 58 ordered pair clusters, on the difference; drawn
%        once, in (b)."
%       "No interval is stored for permuted or for swap alone."
%     Both lines are true of the file, they still fit S17's one block of at
%     most two lines, and the reason band (a) carries no bar is now the same
%     reason the ledger has always given -- legibility -- rather than a claim
%     about the instrument.
%
% X3. THE THRESHOLD IN (a) STANDS ON ITS SPINE. SUPERSEDES the pass-9 y extent
%     8.38 .. 9.04. Drawn from 8.38 it crossed the axis rule and hung 0.16cm
%     below it, past the ends of the hanging ticks and into the tick-numeral
%     lane, so the heaviest stroke in the panel doubled as an outsized
%     unnumbered tick standing between the numerals "0" and "0.02" -- exactly
%     the "the threshold merges with the tick furniture" defect S7 indicts in
%     fig:comparators row (a). It now runs 8.54 .. 9.04, from the spine upward,
%     which is what (b)'s half of the same threshold has always done. One
%     pre-registered threshold, one stroke AND one geometry, in both bands.
%
% X4. BAND (a)'s ONE MARK IS NOW LABELLED AT ITSELF. SUPERSEDES W6's placement.
%     Pass 9 anchored "swap - permuted" base west on x = 0.00, the panel
%     origin, while the leader fell from \ga{diff_perm} = 0.2751 -- so the
%     leader descended from inside the word "swap", and the mark it landed on
%     read as the swap arm, whose value is 0.0019 and not the 0.0013 drawn.
%     That is the only mark in band (a). The label now begins on the mark's own
%     x and the leader falls from its first character.
%
% X5. THE LEADER AND THE PROJECTION TOUCH WHAT THEY POINT AT. The accent leader
%     stopped at 9.12 against a disc whose top edge is 8.9493 -- 4.84pt of
%     white -- so at print it was a detached accent tick under the label, of
%     the same length, weight and orientation as an axis tick. The dotted
%     projection was detached at both ends, 2.0pt below the disc and 0.57pt
%     above the spine. Leader 9.46 -> 8.95 lands on the disc; projection
%     8.85 -> 8.54 leaves the disc's lower edge and reaches the rule.
%
% X6. THREE LABELS THAT READ AS ONE BLOCK ARE SEPARATED, and no label sits on a
%     panel head's baseline any more. SUPERSEDES the pass-9 placements.
%     (i) Band (a)'s accent label moved 9.44 -> 9.56. At 9.44 it stood 7.94pt
%         above the threshold label -- set solid for 7.97pt type -- and the two
%         overlap horizontally by 0.6cm, so they read as one ragged two-line
%         paragraph headed by an accent line. The gap is now 11.3pt.
%     (ii) Band (a)'s terminal label moved off the 9.16 baseline onto 8.86,
%         0.10cm above the end cap it names. At 9.16 it shared a baseline AND
%         the "<name>, <value> nats" construction with the threshold label 8cm
%         to its left, so "the clean preference gap, 0.0642 nats" read as the
%         label of a RULE and a reader looked for a rule at 0.0642 and found a
%         0.22cm cap. It is now unmistakably the label of the axis's own end.
%     (iii) Band (b)'s threshold label moved 5.70 -> 5.44, off the panel head's
%         exact baseline, and its rule's top came down 5.56 -> 5.32 to meet it.
%         On the pass-9 build, scanning the two panel heads read "(b) the first
%         eighth, magnified 8x  ...  10% materiality margin" as one title line.
%         Both threshold labels now stand 0.12cm above their own rule's top.
%
% X7. THE DENOMINATOR IS DELIMITED AT BOTH ENDS. The gap -- the span whose full
%     width is the quantity every effect in section 4.5 is divided by -- was
%     capped at the right by a rising 0.40pt ink stub and at the left by
%     nothing but the same rulegrey hanging tick that 0.02, 0.04 and 0.06 carry,
%     so it did not read as a bounded span at all. A matching ink stub now
%     rises at x = 0. No new weight, no new colour, no height.
%
% X8. THE MAGNIFIER BRACKET: RAISED, AND ITS LABEL NAMES THE OBJECT. It moved
%     7.14 -> 7.22 (rail), tight under the scope line at the clearance the
%     scope line's descenders allow, and its label became "the eighth magnified
%     in (b)" from "magnified in (b)", so the mark is named as well as pointed
%     (S20). TWO HALVES OF WHAT WAS ASKED WERE OVERRULED; see X16 (b) and (c).
%
% X9. BAND (b)'s ZERO RULE NO LONGER READS AS A PANEL EDGE. SUPERSEDES the
%     pass-9 extent 3.58 .. 5.56. Drawn from the spine it was collinear and
%     continuous with the axis's left terminus, that terminus's hanging tick
%     and the panel's x origin, ran the full height of the data region, and
%     matched the threshold rule on the right for length to the point -- so all
%     three readers took the pair for the two sides of a frame round the rows,
%     which is the object S12 exists to forbid, and one noted that band (a)
%     draws no such rule at its own zero. It now runs 3.86 .. 5.28: it spans
%     the three rows only, stands 0.28cm clear of the spine and 0.42cm clear of
%     the panel head, and reads as the coordinate the interval bar's left end
%     clears by 0.70cm. It is still S7's register and still unlabelled, because
%     the axis's own "0" numeral stands at its foot.
%
% X10. ONE DASH, ONE MEANING. Rows one and two glossed their arm with an em
%     dash in the same position after the token where row three and band (a)
%     set a math minus, so one glyph class carried "namely" twice and
%     "subtract" once on a figure whose entire subject is a subtraction. The
%     two glosses now take a comma. The minus in "swap - permuted" is the only
%     dash inside any label on the drawing.
%
% X11. THE FOOT BLOCK'S LEADING. The two lines were set 0.24cm apart, which is
%     6.81pt of leading for 7.97pt type -- tighter than solid, tighter than
%     every other stacked pair in this figure, and tight enough that the
%     descender of "bootstrap" met the dot of the "i" below it. They are now
%     0.28cm apart, the figure's own value everywhere else, and the pinned
%     bounding box floor went 1.76 -> 1.72 so descenders stay inside it.
%
% X12. ONE X ORIGIN, MEASURED RATHER THAN ASSERTED. SUPERSEDES the layout note
%     "inner sep = 1pt ON THE WHOLE PICTURE" in its horizontal half only. With
%     inner sep = 1pt every left-anchored text box stood 1.19pt (0.42mm) inside
%     the axis's left end, so both rules and both left end ticks protruded to
%     the left of the entire text column and S5's "left-aligned to the axis's
%     left end" was met to 0.42mm rather than exactly. inner ysep stays 1pt, so
%     every vertical clearance recorded above is still box arithmetic; inner
%     xsep is 0pt. Measured on this build, the ten panel-level text lines now
%     share ONE ink x0, 59.81pt, against an axis rule beginning at 59.615pt --
%     the residue is the first glyph's own side bearing and nothing else.
%     Band (b)'s three row labels are now placed with an xshift on the MARK's
%     own coordinate rather than at a literal x, so each is tied to the thing
%     it names: 5pt right of the disc in rows one and two, 5pt right of the
%     interval's upper end in row three, which is that row's rightmost ink.
%
% X13. THE AXIS LABELS ARE SET LOWERCASE. SUPERSEDES W16's closing paragraph and
%     repair_pass_9's not_applied entry, both of which recorded that the spec's
%     capitalised string "Preference shift, nats" was followed deliberately
%     because the slate is split on case. It is now "preference shift, nats" in
%     both bands. The reason for changing a spec-verbatim string, recorded so a
%     human can put it back in one edit if the slate is unified the other way:
%     the contract's own ruling is "between 4.11 and A.3: take 4.11", 4.11 sets
%     all four of its axis labels lowercase and 4.13 sets both of its, S19's
%     registry gives the quantity's name lowercase ("preference shift"), and
%     this figure's own panel heads are lowercase after S2 -- so capitalising
%     the axis label alone left one string in the figure with a case of its own.
%     Only the case changed; the quantity and the unit are the spec's words.
%
% X14. THE MAGNIFICATION IS NOW STATED TWICE, NOT THREE TIMES. SUPERSEDES
%     W16(ii), which recorded band (b)'s scope line as set verbatim from the
%     spec. It read "the first eighth of the axis above, magnified 8x" and sat
%     3.2cm below a panel head reading "(b) the first eighth, magnified 8x".
%     The head keeps the factor; the axis-label line keeps the relation to the
%     axis above, which is the fact the head does not carry and the fact an
%     axis-label scope line is for. The exactness of the eight is caption
%     material and the caption states it. Band (a)'s bracket label is the third
%     appearance and it stays: it is the only one that is geometric, it names
%     the bracket, and it is 4cm from the other two.
%
% X15. figs/_PENDING_CAPTIONS.tex NOW EXISTS. SUPERSEDES W9's assertion, and
%     the four assertions in materiality.json, that the sentences cut from the
%     drawing had been written into it: at pass 9 the file did not exist, so
%     the two deleted sentences and the whole proposed caption survived only in
%     the commented block at the foot of this file, which no build reads and no
%     integrator was pointed at. The file is created at this pass, keyed by
%     \label{fig:materiality}, and carries the proposed caption in full, the
%     two sentences deleted under W9 verbatim, and a list of the clauses of the
%     PRINTED caption (Sections/04b-representation.tex:597-620) that the
%     rebuild made false. No builder on this pass may edit Sections/.
%
% X16. FINDINGS OVERRULED, WITH THE REASON, so they are not re-applied blind.
%     (a) "Drop a quietink dotted projection from each of band (b)'s three
%         marks to the spine, as band (a) does for its one mark." OVERRULED.
%         Two of the three would cross the interval bar: the permuted mark
%         stands at x = 1.118 and the swap mark at x = 3.318, and the bar runs
%         0.702 to 3.751 on the row beneath them, so the projections would
%         strike it at two interior points and read as endpoints of something.
%         Band (a)'s projection is not a convention this figure owes band (b);
%         it is a rescue for a single disc standing 2.75mm from the origin on a
%         139.96mm rule with no tick within 4cm. In (b) every mark is on a row,
%         direct-labelled at itself, 4.8 to 15.6mm above a rule whose ticks
%         stand every 4.36cm, and the two arm values the projections would help
%         a reader read off are now PRINTED, in the caption, which is where the
%         same reviewer asked for them.
%     (b) "Raise the bracket into the lane immediately under the tick numerals."
%         OVERRULED. S5 puts the axis-label block directly under the tick
%         numerals; the bracket carries a \scriptsize label of its own, so
%         moving it there interposes a second text block between the numerals
%         and the line that names the quantity, which is the supervisor's first
%         complaint reintroduced by the back door. It was raised as far as the
%         scope line's descenders allow instead (X8).
%     (c) "Reverse the bracket's stubs so they turn UP toward the rule."
%         OVERRULED. With the bracket necessarily BELOW the two-line axis-label
%         block, up-turned stubs bracket the scope line's text -- "residual
%         arm, layer 12, alpha = 0.5" runs 0 to 3.66cm and the bracket spans 0
%         to 1.749cm, so they would enclose its first half. Down-turned stubs
%         agree with the label's own pointer, "the eighth magnified in (b)",
%         and point at the panel that opens the span. This was pass 9's stated
%         intent and it survives the pass.
%     (d) "Drop \texttt on the arm names, or scale the mono, because Inconsolata
%         sets at 6.85pt inside 7.97pt roman." OVERRULED. 6.85pt is the
%         preamble's own scaling of the mono face to Pagella, not a second type
%         size chosen here, and S1 exempts the \texttt token name from the size
%         count for exactly this reason. The tokens are JSON key names --
%         swap, permuted, residual -- and setting them roman would turn
%         "swap, points at the drug scored" into a sentence about an English
%         verb. The register is the corpus's (ladder.tex, comparators.tex).
%     (e) "The threshold carries two label strings, '10% materiality margin,
%         0.00642 nats' in (a) and '10% materiality margin' in (b), which fails
%         S19's mechanical check." OVERRULED on the ledger's own standing
%         principle, applied to diff_perm at W6: the NAME is "10% materiality
%         margin" in both bands, and (a) additionally prints the VALUE. A value
%         is not a second name. It is printed in (a) and not in (b) because (a)
%         is the band that also prints the quantity it is a tenth of, so the
%         tenth is a division a reader performs on two printed numbers 8cm
%         apart -- which is what W5 traded the deleted brace for.
%     (f) "Restore the spec's clause 'the full width of this rule is the
%         model's clean preference gap' to band (a)'s scope line." OVERRULED
%         again, for W16(i)'s reason and one more: with the terminal label now
%         standing on its own end cap (X6 ii) and both ends of the span capped
%         in ink (X7), the drawing says geometrically what that clause says in
%         twelve words, and the panel head says it in five.
%
% X17. MEASURED ON THE PASS-10 BUILD, not estimated. figs/_preview.py --width
%     text materiality reports ok, pages = 1, Overfull 0, Underfull 0. Content
%     bounding box 139.96 x 84.12mm (S14's 138-142mm; the 85mm height cap).
%     Ink 5.07% of pixels below 230 grey at 200dpi inside that box (ceiling
%     6.0%). 572 characters of figure text, 4.86 per 100mm^2 (ceiling 5.0).
%     Rule weights, counted off the PDF: 0.30pt x10 (nine rulegrey furniture
%     strokes and one accent leader), 0.40pt x4 (the two spines and the two end
%     caps), 0.50pt densely dotted x2, 0.90pt x2, 1.00pt x1 -- five weights,
%     each with more than one instance or declared as a distinct kind of object
%     at W11 -- plus the 0.80pt open-disc edge, which is a mark edge and not a
%     rule. Type: 7.97pt (400 characters) and 9.96pt (105) in TeXGyrePagella,
%     which are \scriptsize and \footnotesize; 6.85pt (58) in Inconsolata,
%     which S1 exempts; and 8.30pt (8) and 10.38pt (1) in the math companions,
%     which are the SAME two nominal sizes set from the math families -- the
%     minus in "swap - permuted" and the \times in "magnified 8x". Italic: the
%     render carries exactly one span in an italic-named face, PazoMath-Italic
%     at 7.97pt, and it is the mathematical $\alpha$ of "alpha = 0.5" -- 2 of
%     572 characters, 0.35%, which is a symbol's own face and not an italic
%     voice, the use S16 allows. There is no \itshape anywhere in the source.
%     PASS 9's W12 and the ledger's S16 line both said "0%"; the figure of
%     record is 0.35%, and it is the alpha. Sans: 0%. Bold: 0%. Rectangles: 0.
%     Arrowheads: 0.
% ===========================================================================
% ===========================================================================
% ===========================================================================
% PASS 11 -- 2026-08-19. REPAIR PASS on the pass-10 repair. Three adversarial
% readers went over the built figure independently and returned twenty-three
% findings. NOTHING ABOVE IS DELETED: passes 1-10 are the audit trail and stay
% as written. Every clause of passes 1-10 that this pass made false, and every
% clause those passes left false, is superseded here BY NAME (S25). The
% findings that were OVERRULED are recorded under Y15 with the reason, so that
% a later builder does not re-apply them blind.
%
% NO NUMERIC VALUE CHANGED IN THIS PASS EITHER. Every quantity drawn or
% printed is the same float from the same key path as at passes 8, 9 and 10:
% clean_ab_gap 0.06420005919957501, equiv_margin 0.006420005919957502,
% effect_perm 0.0006409306049362111, effect_swap 0.0019027651015006252,
% diff_perm 0.0012618344965644142, ci_perm [0.0004024989353392094,
% 0.002150955667015771], n_pair_clusters 58, and the exact magnification 8.
% Both maps are untouched: \ga = (#1)*218.0000 and \gb = (#1)*1744.0000, with
% 1744.0000 = 8 x 218.0000 exactly (rechecked on this build: 8*218.0 == 1744.0
% is True in IEEE double, and 0.06420005919957501*218.0 ==
% 0.008025007399946877*1744.0 is True, both ends coinciding at
% 13.995612905507352cm). ONE PRINTED STRING changed its number of decimal
% places without changing its value -- band (b)'s middle tick numeral, 0.005
% -> 0.0050 (Y11) -- and one value that was printed on a row of (b) is now
% printed in (a) instead (Y2, Y8). Nothing was rounded differently, and nothing
% new was computed.
%
% WHAT MOVED, IN ONE LINE EACH. The changes are commented at the point of edit
% as Y1..Y13 and summarised here.
%   Y1  band (a)'s panel head names the object by the registry string.
%   Y2  band (a)'s one label prints the estimate and drops 0.06cm.
%   Y3  the axis's right end is computed by \ga, not transcribed.
%   Y4  band (a)'s scope line drops 0.12cm and says "layer 12 of seven".
%   Y5  the magnifier bracket's stubs are visible and it stands clear.
%   Y6  band (b)'s zero rule is named "zero".
%   Y7  the mark-to-label gap in all three rows of (b) goes 5pt -> 9pt.
%   Y8  the estimate's value comes off (b)'s row-three baseline label.
%   Y9  0.06cm more air under row two.
%   Y10 band (b)'s rule is capped at its right end, in (a)'s cap.
%   Y11 band (b)'s middle tick numeral prints 0.0050.
%   Y12 the foot block's leading and its stand-off both open.
%   Y13 the pinned bounding box floor goes 1.72 -> 1.68.
%
% Y14. CLAUSES OF PASSES 1-10 THAT WERE ALREADY FALSE OF THE BUILT FIGURE AND
%     WERE SUPERSEDED NOWHERE. W15 and X1 each did this job for the pass they
%     followed; neither went back over the pass-1 to pass-8 LAYOUT NOTES and
%     DESIGN DECISIONS, so an auditor reading this header top to bottom was
%     handed five descriptions of a figure that does not exist -- the exact
%     hazard X1 was written to prevent. Each is named and given its drawn value.
%     (i) DESIGN DECISION 5, "THE MEASUREMENT SITS 0.22cm ABOVE BAND ONE'S
%         SPINE ... with its dot over a 2.1pt white halo ... the bar, the dot
%         and the spine are three separable things."
%         -> The disc is drawn at y = 8.90 against a spine at 8.54, which is
%            0.36cm and not 0.22cm. There is NO white halo under it: the only
%            \fill[white] in the file is band (b) row one's open control disc.
%            And there is no bar in (a) at all, since W6 drew the interval once
%            and gave (a) a single filled disc. The decision's REASON -- that a
%            mark on the spine reads as a thickening of the axis, so it is
%            lifted clear -- survives; its three measurements do not.
%     (ii) LAYOUT NOTE "TICK NUMERALS CARRY fill=white AND ARE DRAWN LAST IN
%         THEIR BAND."
%         -> num/.style and numL/.style carry no fill. The white box existed to
%            interrupt the magnifier's left leader where it crossed band one's
%            "0"; W10 deleted that leader, so nothing crosses the numerals and
%            drawing order no longer matters.
%     (iii) LAYOUT NOTE "The band-two tag is anchored at x = 0.75 and not at
%         x = 0.00, because the magnifier's left leader is a vertical at
%         x = 0.55."
%         -> Neither object exists: the tag went at W7 (replaced by 4.11's
%            panel head) and the leader at W10.
%     (iv) LAYOUT NOTE "BAND ONE'S RIGHT END is a terminal tick RISING from the
%         spine, 0.30cm tall against the 0.10cm hanging axis ticks, with its
%         two label lines right-aligned on the same x = 13.85."
%         -> It rises 0.22cm, from 8.54 to 8.76, at x =
%            \ga{0.06420005919957501} = 13.995612905507352, and it carries ONE
%            label line, not two. Band (b)'s right end now carries the same cap
%            at the same x (Y10).
%     (v) LAYOUT NOTE "The bounding box is pinned with a bare \path so the
%         measure is exact and nothing is clipped: (0.0000,0.30) to
%         (14.1900,8.98) = 141.9 x 86.8mm."
%         -> It is pinned (0.00,1.68) to (13.9957,10.18) after Y13, and the
%            measured content box is at Y16. The 141.9 x 86.8mm figure is a
%            pass-2 measurement and is over both the 142mm measure and the 85mm
%            height cap.
%     (vi) W11's "The 0.50pt densely dotted stroke is used once, for band two's
%         zero rule, and is a distinct kind of object."
%         -> It has two instances in this build and had two in pass 10, as
%            X17's own weight count records: band (a)'s projection from the disc
%            to the spine, and band (b)'s zero rule. It is not a singleton and
%            needs no singleton declaration under S10. The 1.00pt interval bar
%            IS a singleton and its declaration stands.
%
% Y15. FINDINGS OVERRULED, WITH THE REASON, so they are not re-applied blind.
%     (a) "Take the open control disc's edge to ink rather than quietink, as
%         fig 4.11(b) does for its matched-dimension random slab." OVERRULED.
%         S8 gives quietink to a comparator, a null or a control and S9 says
%         the open disc's 0.8pt edge is drawn "in its colour"; ink is the rank
%         S8 reserves for thresholds, axes, axis labels, titles and definitional
%         prose, so an ink ring would put the control in the apparatus rank and
%         make the comparator heavier than the accent mark under test. 4.11(b)
%         can draw its null black because its palette has no quietink rank at
%         all; this figure has four ranks and the control's is quietink. The
%         defect the finding reports is real and blocking and is fixed by the
%         gap alone (Y7): 7.2pt of white against the mono face's own 3.43pt
%         character advance, where 3.28pt was letter-spacing.
%         The companion finding, "quietink carries two ranks -- the one
%         measured comparator's label and all four apparatus lines", is
%         OVERRULED for the same reason from the other side: S5 mandates
%         quietink for both scope lines and S17 for the foot block, so the
%         collision cannot be broken by recolouring either party. It is broken
%         by position, which is how this figure has always broken it: the
%         control's label is anchored on a mark, on a row, inside the data
%         field; every apparatus line is panel-level text on the figure's one
%         x origin, under an axis.
%     (b) "Drop band (b)'s two arm dots. Their values are quoted in no sentence
%         of the thesis and the proposed caption already prints them, and the
%         one interval bar horizontally contains both of them." OVERRULED.
%         The spec's keep-list makes the three-row structure -- permuted, swap,
%         and their difference -- load-bearing, and CF14 fails a rebuild that
%         loses a load-bearing piece even if it passes every typographic check.
%         Without the two arms the difference is a bare number with a bar and
%         the figure stops drawing the subtraction that is its subject; design
%         decisions 7 and 14 record that reasoning and it stands.
%         THE DIAGNOSIS IS ACCEPTED AND IS REAL, and it is answered at Y9
%         rather than by deletion. The arithmetic is an accident worth stating
%         plainly: the bar runs \gb{0.0004024989353392094} = 0.702cm to
%         \gb{0.002150955667015771} = 3.751cm, the control stands at 1.118 and
%         the arm at 3.318, so the interval on the DIFFERENCE brackets both of
%         its own operands, and \gb{diff_perm} = 2.201 falls within 0.02cm of
%         their midpoint. What is done about it: the row pitch is no longer
%         flat (0.58 / 0.64 rather than 0.54 / 0.54), so the derived row sits
%         off the pair it is derived from; the row's name states the derivation
%         at the bar's own end; and rows one and two name their arms
%         individually immediately above, so "swap $-$ permuted" is read after
%         "permuted" and "swap" and not before them.
%         The finding's own fallback -- "move rows one and two onto their own
%         short strip with their own zero" -- is also OVERRULED: a second rule
%         for one quantity at one scale inside one panel is a second axis a
%         reader has to reconcile with the first, against S13's one grid, and
%         it would put two zeros 5mm apart on a figure whose subject is a
%         distance from zero.
%     (c) "Move the magnifier bracket into band (a)'s tick lane, rail at the
%         hanging ticks' depth with its stubs turned UP to the spine; or, if
%         that lane cannot be had, delete the bracket and tint the magnified
%         span in (a) and the whole of (b)'s data field instead." OVERRULED,
%         both halves. X16 (b) and (c) already overruled the first and the
%         reason holds: a bracket in the tick lane needs its label in the
%         tick-numeral lane, which is the defect S24 indicts in fig:purify,
%         where the word "four" sets on the x-tick baseline and reads as a
%         leftmost tick. The tint is worse, not better: S12 permits a filled
%         rectangle only where the rectangle IS an object in the argument and
%         forbids outright a rectangle drawn round a panel, and a tint over the
%         whole of (b)'s data field is that rectangle by another means -- the
%         object W4 deleted. The finding's real complaint, that the assembly did
%         not read as a bracket at all, is accepted and fixed at Y5.
%     (d) "Raise the bracket's rail onto its label's own baseline, so the rail
%         and the label read as one object rather than the rail reading as a
%         strike-through beside the words." OVERRULED as already true. The rail
%         at 7.02 stands 0.06cm above a label baseline at 6.96, which is that
%         label's own x-height centre, so rail and label are optically on one
%         line. What made it read as a rule was the 0.51mm stubs and the scope
%         line overhanging both its ends; both are fixed at Y5.
%     (e) "Label band (b)'s right cap 'one eighth of it, 0.008025 nats', so
%         that the exact eight becomes a division a reader can perform on two
%         printed numbers." OVERRULED on geometry, not on principle. Band (b)'s
%         threshold rule stands at \gb{0.006420005919957502} = 11.196cm and
%         runs the full height of the data field; a label right-aligned on the
%         axis's end at 13.996cm measures 3.4 to 4.7cm at \scriptsize and
%         crosses it, and the only string short enough to clear the rule is a
%         bare numeral with no name, which S20 forbids. There is no lane above
%         the rule either: (b)'s threshold label already occupies it, ending at
%         11.10cm. The magnification is stated by (b)'s panel head, by (b)'s
%         axis-label scope line and by (a)'s bracket label, and its EXACTNESS
%         is caption material (X14, and the caption states it). The cap added
%         at Y10 is axis furniture -- the rule's own end -- and not a mark, so
%         it needs no name under S20; its end is named by (b)'s axis label,
%         which is W13's reasoning and survives.
%     (f) "Restore the spec's clause 'the full width of this rule is the
%         model's clean preference gap' to band (a)'s scope line: line one
%         reads 'preference shift, nats', but band (a)'s full width is
%         clean_ab_gap, the model's contrast score with NO intervention -- a
%         level, not a shift -- so a reader who takes the axis label at face
%         value reads the right terminus as a preference shift of 0.0642 nats."
%         OVERRULED a third time. The REASON RECORDED AT W16(i) AND X16(f) IS
%         WITHDRAWN AND REPLACED, because it was wrong: both cited S18, and
%         S18 forbids a drawing from restating a clause of ITS OWN CAPTION or
%         of the paragraph that introduces the float. It says nothing about a
%         drawing carrying two related statements of its own. The clause is
%         refused on S5 and S15 instead. S5 caps the axis-label block at two
%         lines and gives line two to scope or a scale note; a thirty-character
%         clause about what the axis's other end IS is neither. And the fact is
%         already stated at the two places a reader looks for it: the panel head
%         reads "(a) the whole clean preference gap" after Y1, and the terminal
%         label at the axis's own right end reads "the clean preference gap,
%         0.0642 nats". THE FINDING'S UNDERLYING OBSERVATION IS RIGHT AND IS
%         RECORDED HERE so that no later reader has to rediscover it: this
%         axis's quantity name is the quantity its MARKS carry, and its right
%         terminus is a level of that same contrast score measured with no
%         intervention, named as such at the terminus. Laying the two on one
%         rule is the figure's whole construction and the thesis's own -- 04b:
%         663 divides every effect by "that model's own clean preference gap,
%         measured with no intervention".
%     (g) "Scale the mono up so that \texttt sets at the roman's x-height:
%         6.85 x (3.96/3.24) = 8.37pt." OVERRULED again, as at X16(d), and for
%         a reason that finding does not meet. 6.85pt is preamble.tex's global
%         scaling of Inconsolata to Pagella and it applies to every \texttt in
%         the thesis. Overriding it inside one figure would make this figure's
%         JSON key names set differently from comparators.tex's, ladder.tex's
%         and the body text's, which is a worse inconsistency than an x-height
%         step inside a label; and S1 exempts the \texttt token name from the
%         size count precisely so that the document's global scaling can stand
%         without a figure having to fight it.
%     (h) "In band (a) the two black labels that share the '<name>, <value>
%         nats' construction are the CLOSER pair (0.30cm) while the accent
%         label, which differs from both in face, colour and construction, is
%         the further (0.40cm); drop the terminal label 8.86 -> 8.80 and lift
%         the threshold label 9.16 -> 9.22." OVERRULED as proposed, accepted in
%         substance. The two black labels are 34mm apart horizontally, sit on
%         different baselines, and one is right-aligned onto the axis's own end
%         cap while the other is left-set beside a rule 8cm away; X6(ii)
%         separated them for the stronger reason that they had SHARED a
%         baseline, and the proposed 8.80 would set a 7.97pt baseline 1.1pt
%         above the 0.22cm cap it labels, re-creating the collision X6(ii)
%         removed. After Y2 the stack reads 0.40 / 0.34 / 0.30 from the head
%         down -- monotone, so nothing alternates and no pair is singled out by
%         its spacing.
%
% Y16. MEASURED ON THE PASS-11 BUILD, not estimated, and SUPERSEDING X17's
%     bounding box by name. figs/_preview.py --width text materiality reports
%     ok, pages = 1, Overfull 0, Underfull 0.
%       content bounding box   140.34 x 84.33mm   (S14's 138-142mm; 85mm cap)
%       ink                    5.36% of pixels below 230 grey at 200dpi (6.0%)
%       figure text            587 characters, 4.96 per 100mm^2   (ceiling 5.0)
%     X17 recorded 139.96 x 84.12mm for the pass-10 build; re-measured with the
%     contract's own threshold (pixels below 230 grey at 200dpi, clipped to the
%     content box) that build is 140.21 x 83.82mm. The +0.25mm and -0.30mm are
%     a method difference and not a drawing change, and they are recorded here
%     rather than left as an unexplained discrepancy: X17's number came from a
%     box measured before the antialiased edge columns of the axis caps were
%     counted. Both builds are inside every limit; this one is 0.67mm under the
%     height cap and 0.04 characters per 100mm^2 under the density ceiling,
%     which is why six characters came out of the foot block when "of seven",
%     "clean" and "zero" went in.
%     Rule weights, counted off the PDF: 0.30pt x10 (nine rulegrey furniture
%     strokes -- eight ticks and the bracket polyline -- plus one accent
%     leader), 0.40pt x5 (two spines and three end caps), 0.50pt densely dotted
%     x2, 0.90pt x2, 1.00pt x1. Five weights, no singleton except the interval
%     bar, which W11 declares as a distinct kind of object. Plus the 0.80pt
%     open-disc edge, which is a mark edge and not a rule.
%     Type: 7.97pt (417 characters) and 9.96pt (111) in TeXGyrePagella, which
%     are \scriptsize and \footnotesize; 6.85pt (58) in Inconsolata, which S1
%     exempts; 8.30pt (8) and 10.38pt (1) in the math companions, which are the
%     same two nominal sizes set from the math families. Italic: one span,
%     PazoMath-Italic at 7.97pt, the mathematical $\alpha$ -- 2 of 587
%     characters, 0.34%, a symbol's own face and not an italic voice (S16).
%     Sans 0%. Bold 0%. Rectangles 0. Arrowheads 0.
%
% Y17. THE CAPTION. figs/_PENDING_CAPTIONS.tex is updated in the same pass and
%     three of its clauses changed with the drawing: band (a)'s mark now prints
%     its value, band (b)'s row three no longer does, and (b)'s rule is capped
%     at both -- at its right -- end. ONE ITEM IN THAT FILE IS BLOCKING and is
%     flagged there as such: the printed caption (04b:617-618) says only "this
%     is one arm at one depth at one rung", which is a scope statement, while
%     layer 12 is the argmax of the anchor seed's seven-depth sweep and the
%     only one of the seven whose interval excludes zero. The drawing now says
%     "layer 12 of seven" (Y4), which discloses that the depth was chosen from
%     a set but not that it was the favourable one; only the caption can carry
%     that, and no builder on this pass may edit Sections/.
% ===========================================================================
%
% ===========================================================================
% v-BLOCK, 2026-08-19: WHOLE-SLATE PASS (fig:comparators, fig:licensing,
% fig:purify, fig:hook, fig:materiality reconciled against one another and
% against fig:mechanism, which is the document's reference standard and was
% NOT rebuilt). Nothing above this line is deleted; where an earlier claim is
% now false it is superseded EXPLICITLY below.
%
% WHY. The five figures were rebuilt independently and drifted. This pass owns
% only the differences BETWEEN them. No value changed in any of the five; the
% sibling .json ledgers carry the value-to-position maps, old and new.
%
% THE SLATE RULES THAT NOW HOLD ACROSS ALL FIVE (each stated once here, in
% every file, so any future single-figure edit can see what it would break):
%
%  R1 TICK NUMERALS are quietink, \scriptsize, inner sep 0, anchored north on
%     the tick and set 0.16cm below the axis rule, over a 0.10cm rulegrey
%     0.30pt tick. Before: comparators/licensing/purify quietink, hook and
%     materiality ink; ticks 0.10 / 0.12 / 0.14cm; three different gaps.
%  R2 A THRESHOLD'S OWN LABEL IS BLACK, matching the black solid 0.90pt rule
%     it names (S6). Before: comparators black, purify's "gate, 0.8" and
%     materiality's two margin labels ink.
%  R3 ONE NAME PER OBJECT (S19). The pre-registered 10% rule is "10%
%     materiality margin" in fig:comparators row (c), fig:materiality (a) and
%     (b), and in figs/family.py, which is where the name comes from.
%     fig:comparators' "pre-registered margin" is superseded; "pre-registered"
%     is now the caption's word, not the drawing's. Zero, where it is the null
%     a bar is read against, is "zero: no identity effect" in BOTH
%     fig:comparators row (c) and fig:materiality (b) -- one quantity, one
%     wording. V_pure is spelt with \textrm everywhere, as figs/slab.tex does.
%  R4 COLOUR KEY, one key for five figures (S8). accent = the drug-side object
%     under test. quietink = every comparator, null, control, replication or
%     discounted arm, AND its label, AND its leader. ink/black = thresholds,
%     axis rules, axis labels, panel heads, definitional prose. rulegrey =
%     furniture only. Within quietink the GLYPH separates two kinds: an OPEN
%     disc is a construction built to carry no signal (a null, a control, a
%     raw/before series); a FILLED disc is a real measurement that is not the
%     headline. fig:hook (d) drew both its null and its cell-line control in
%     ink and is corrected.
%  R5 PROJECTION AND ZERO RULES are quietink 0.50pt densely dotted (S7).
%     fig:hook's two panel-(c) projection lines were rulegrey 0.30pt densely
%     dashed and are corrected.
%  R6 DEFINITIONAL PROSE is \footnotesize ink and OUTRANKS every label (S4);
%     at most one block per figure, on the panel's own x origin. fig:hook's
%     one prose line was \scriptsize and is raised. fig:purify's panel-(b)
%     block carried an argument clause ("so causal claims use the purified
%     slab, not the raw") that its caption already carries, and is now
%     definitional only.
%  R7 ONE INTERVAL DECLARATION per figure, \scriptsize quietink, at the foot,
%     on the figure's x origin, opening with the word "Intervals:" and naming
%     EVERY panel including the ones with no measured mark (S17).
%     fig:comparators already did this and is the model; the other four now
%     match its form.
%  R8 THE FIVE-FIELD PROMPT STRIP is ONE object drawn in two figures, and the
%     two drawings are now congruent: cell boundaries 0 / 1.12 / 1.90 / 2.64 /
%     4.23 / 6.90 cm and cell height 0.42cm in both fig:licensing (a) and
%     fig:hook (a) and (b), labels centred in their cells, the drug cell
%     accent-bordered at 0.40pt on accent!12 and its label NOT bold. Before:
%     licensing 6.90cm wide with 0.62cm cells and a bold label, hook 7.44cm
%     wide with 0.42cm cells and a plain label -- the same five fields at two
%     sizes, four pages apart, with each caption pointing at the other.
%  R9 MEASURE. Every fullwidth figure draws 172.0-174.0mm; the one textwidth
%     figure draws 140.3mm. Ink at 200dpi inside the content bbox, counted as
%     greyscale luminance below 230 (the contract's own method), is now
%     5.63-5.99% on all five, under the 6.0% body ceiling, where before this
%     pass it ran 5.36-6.54%.
%
% WHAT THIS PASS DELIBERATELY DID NOT UNIFY, so that a later reader does not
% think it was missed:
%  * PANEL-HEAD CASE. fig:comparators sets "(a) Is it there?" and the other
%    four set "(a) where the drug name enters". comparators' heads are
%    interrogative SENTENCES and sentence case is correct for a sentence; the
%    v4 header argued this and the argument stands. fig:mechanism's lowercase
%    heads are noun phrases, not questions. One capital letter, four times.
%  * PANEL-HEAD POSITION. fig:comparators puts its heads in a right-aligned
%    left stub, the other four above the panel at its x origin. Moving them
%    above the rules costs about 20mm of height against a 118mm cap that the
%    figure already meets at 117.60mm. Not available.
%  * TWO-COLUMN GUTTER. purify and hook split at x = 8.40cm, licensing at
%    7.75cm. licensing cannot move right: "(d) what a negative would buy" is
%    4.85cm wide and its origin is already at 12.40 of a 17.30cm measure, with
%    1.5pt of slack. Moving purify and hook LEFT to 7.75 would be arbitrary.
% ===========================================================================
%
% CHANGED IN THIS FILE. (i) No math sign is set from a CM font any more. The
% panel-(b) head's 8x (CMSY10 at 10.38pt, the largest glyph on the figure and
% a THIRD type size above 7pt) is now 8\texttimes{}; the two minus signs of
% "swap - permuted" are \textminus{}; +0.0013, 10%, 0.00642, 0.0642, 12, 58
% and 95% lose their math mode, and $\alpha = 0.5$ becomes $\alpha$ = 0.5,
% which is fig:comparators' idiom. The figure now prints exactly 9.96 and
% 7.97pt in one face, plus the \texttt arm names. (ii) Tick numerals go ink ->
% quietink with inner sep 0, set 0.16cm below their rule (R1). (iii) Both
% "10% materiality margin" labels go black, matching the black solid 0.90pt
% rule they name (R2). (iv) Band (b)'s zero label "zero" -> "zero: no identity
% effect", which is fig:comparators row (c)'s wording for the same null under
% the same quantity (R3), and goes ink -> quietink so that the label matches
% the dotted quietink rule it names. (v) \texttt{residual} -> residual in the
% scope line: fig:comparators names the same arm "residual seeds at layer 12"
% in plain roman, and an arm name in prose is not a code token. \texttt{swap}
% and \texttt{permuted} stay, because they are the two stored conditions and
% tab:identity sets them so. (vi) The foot block takes the slate's declaration
% form (R7).
% NOT CHANGED, and still declared: the 1.00pt interval bar is a singleton
% weight because this figure stores exactly one interval (S10's permission);
% the 0.80pt open-disc edge is a MARK edge, not a rule.
% RE-MEASURED: 140.34 x 84.45mm (cap 85, textwidth band 138-142); ink 5.630%;
% 620 characters, 5.23 per 100mm^2. One page, zero Overfull, zero Underfull.
% ===========================================================================
\begin{tikzpicture}[x=1cm, y=1cm, font=\small, inner ysep=1pt, inner xsep=0pt,
  %% TYPE. Two sizes above 7pt and one face (S1, W12). \footnotesize = 10.0pt
  %% carries the two panel heads and the two axis-label line ones and nothing
  %% else; \scriptsize = 8.0pt carries ticks, direct labels, scope lines, the
  %% threshold labels and the interval declaration. No sans, no bold, no italic.
  %% INNER SEP, PASS 10 (X12). It was inner sep = 1pt on the whole picture; the
  %% vertical half is unchanged, so every clearance below is still box
  %% arithmetic, but the horizontal half is now 0pt. With xsep = 1pt every
  %% left-anchored text box stood 1.19pt (0.42mm) inside the axis's left end,
  %% so both rules protruded to the left of the entire text column and S5's
  %% "left-aligned to the axis's left end" was met to 0.42mm and not exactly.
  head/.style  ={font=\footnotesize, text=ink,      anchor=base west},
  axlab/.style ={font=\footnotesize, text=ink,      anchor=base west},
  scope/.style ={font=\scriptsize,   text=quietink, anchor=base west},
  qlabz/.style ={font=\scriptsize,   text=quietink, anchor=base west},
  lab/.style   ={font=\scriptsize,   text=ink,      anchor=base west},
  labR/.style  ={font=\scriptsize,   text=ink,      anchor=base east},
  acc/.style   ={font=\scriptsize,   text=accent,   anchor=base west},
  rowq/.style  ={font=\scriptsize,   text=quietink, anchor=west},
  rowa/.style  ={font=\scriptsize,   text=accent,   anchor=west},
  %% Tick numerals are INK, not quietink: they are part of the axis, and S8
  %% gives ink to axes and axis labels while quietink means a comparator or a
  %% control. 4.11 and 4.13 both set their tick numerals dark for this reason.
  num/.style   ={font=\scriptsize,   text=quietink, anchor=north,
                 inner sep=0pt},
  numL/.style  ={font=\scriptsize,   text=quietink, anchor=north west,
                 inner sep=0pt},
  blab/.style  ={font=\scriptsize,   text=black,    anchor=base west},
  blabR/.style ={font=\scriptsize,   text=black,    anchor=base east},
  %% FIVE STROKE WEIGHTS, each with one meaning (S10, W11).
  tick/.style  ={rulegrey, line width=0.3pt},                  % furniture
  lead/.style  ={rulegrey, line width=0.3pt},                  % furniture
  accled/.style={accent,   line width=0.3pt},                  % furniture, in the mark's colour (S22)
  spine/.style ={ink,      line width=0.4pt},                  % axis rule and its two end caps
  proj/.style  ={quietink, line width=0.5pt, densely dotted},  % zero rule / projection (S7)
  thresh/.style={black,    line width=0.9pt},                  % the one threshold (S6, W3)
  bar/.style   ={accent,   line width=1.0pt},                  % interval
  ctrl/.style  ={quietink, line width=0.8pt}]                  % open-disc edge: a MARK edge, not a rule (S9, W11)

  %% ==== THE TWO MAPS ====================================================
  %% Superseded constants, kept in materiality.json: 0.55+(#1)*207.1650 and
  %% 0.55+(#1)*1657.3200. The 0.55cm inset is gone so that the axis's left end
  %% and the panel's x origin are one coordinate (S13, W1).
  %% 1744.0000 = 8 x 218.0000 EXACTLY, so the magnification is exactly eight,
  %% and \ga{clean_ab_gap} = \gb{clean_ab_gap/8} = 13.995613cm exactly.
  %% PASS 10, X1: the pass-9 header block records these two constants as
  %% 215.7500 and 1726.0000 and the shared right end as 13.851163cm. Those
  %% numbers were never drawn by this file. The two lines below are the map,
  %% and materiality.json has always recorded them correctly.
  \def\ga#1{{(#1)*218.0000}}   % band (a): 0 .. 0.06420005919957501 nats
  \def\gb#1{{(#1)*1744.0000}}  % band (b): 0 .. 0.008025007399946877 nats

  %% ==================== BAND (a) -- THE WHOLE GAP ========================
  %% PASS 11, Y1: the head names the object by the registry string. It read
  %% "(a) the whole preference gap" while the terminal label 8cm to its right
  %% read "the clean preference gap", so one JSON key (clean_ab_gap) carried
  %% two names inside one panel, which is what S19's mechanical check forbids.
  %% "whole" is kept because it is the contrast with (b)'s "first eighth".
  \node[head] at (0.00,9.90) {(a) the whole clean preference gap};

  % The estimate, drawn once and without a bar (W6). At true scale it stands
  % 2.75mm from the origin on a 139.96mm rule, which is the panel's whole point.
  % Its name is band (b) row three's name, because one JSON key takes one name.
  % PASS 10, X4: THE LABEL IS ANCHORED AT THE MARK, not at the panel origin.
  % Pass 9 set it base west on x = 0.00 while its leader fell from
  % \ga{diff_perm} = 0.2751, so the leader descended from inside the word
  % "swap" and the mark read as the SWAP arm -- whose value is 0.0019, not the
  % 0.0013 drawn. That is the one mark in band (a); mis-naming it loses the
  % panel. The label now begins on the mark's own x, and the leader falls from
  % its first character.
  % PASS 11, Y2: THE LABEL PRINTS THE VALUE, AND IT DROPPED 0.06cm.
  % Two defects, one edit. (i) At 9.56 the label stood 0.34cm under the panel
  % head, near-left-aligned with it (7.8pt of indent), so head and label read
  % as a two-line title and the panel's single mark lost its name; the head's
  % descenders cleared the label's ascenders by 2.6pt. At 9.50 that clearance
  % is 4.3pt and the gap to the threshold label below is 0.34cm = 9.7pt, which
  % is still wider than the 7.94pt X6(i) condemned. (ii) Without a value the
  % math minus between two mono tokens parses as easily as a gloss dash --
  % "swap, that is, permuted" -- so the one mark in (a) could be read as the
  % SWAP arm, whose value is 0.0019 and not the 0.0013 drawn. The value now
  % settles it, and it also makes (a) self-sufficient for the figure's
  % question: 0.0013 against the 0.0642 printed at the axis's right end, and
  % against the 0.00642 printed at the threshold.
  % WHY A VALUE MAY BE PRINTED HERE AND NOT ON A ROW OF (b). This label sits
  % 0.96cm above the spine and is leadered to its disc; positions in that lane
  % are not read as values, which is why (a)'s threshold label already prints
  % 0.00642 three centimetres right of the rule it names and its terminal
  % label prints 0.0642 four centimetres left of the cap it names. A label set
  % on a data row's OWN baseline is read positionally, which is why the value
  % came off (b)'s row three (Y8).
  \node[acc] at (\ga{0.0012618344965644142},9.50)
    {\texttt{swap} \textminus{} \texttt{permuted}, +0.0013 nats};
  % The label's leader is accent 0.30pt, the mark's own colour (S22). PASS 10,
  % X5: it now LANDS ON THE DISC. It used to stop at 9.12 against a disc whose
  % top edge is 8.9493, leaving 4.84pt of white, so the leader read as a stray
  % accent tick under the label and the label read as unattached.
  \draw[accled] (\ga{0.0012618344965644142},9.40)
             -- (\ga{0.0012618344965644142},8.95);
  % The stem down onto the spine anchors the mark to a coordinate rather than
  % leaving it floating above the rule; pass 8 added it after both reviewers
  % reported they could not tell whether the mark was a mark, a tick or a
  % printing artefact, and it survives the rebuild. It is drawn in the PROJECTION
  % register -- quietink 0.5pt densely dotted, which is what S7 reserves that
  % stroke for -- and not in the leader's accent, because on the build that drew
  % leader and stem at one weight and one colour the assembly read as a VERTICAL
  % interval bar with a dot on it, which is the one thing this figure must not
  % appear to draw. PASS 10, X5: it now touches the disc's lower edge (8.85
  % against the disc's 8.8507) and the spine (8.54), instead of floating 2.0pt
  % clear of the one and 0.57pt clear of the other.
  \draw[proj] (\ga{0.0012618344965644142},8.85)
           -- (\ga{0.0012618344965644142},8.54);
  \fill[accent] (\ga{0.0012618344965644142},8.90) circle (1.4pt);

  % The one threshold, black and solid, in both bands (S6, CF2, W3). Its value
  % is printed here and the quantity it is a tenth of is printed at the axis's
  % right end, so the tenth is a division on two printed numbers and needs no
  % brace (W5).
  % PASS 10, X3: IT STANDS ON THE SPINE AND NO LONGER CROSSES IT. Pass 9 drew
  % it from 8.38, which is 0.06cm BELOW the hanging ticks' ends and 0.02cm above
  % the tick numerals, so the heaviest stroke in the panel doubled as an
  % outsized unnumbered tick sitting between the numerals "0" and "0.02" -- the
  % same "the threshold merges with the tick furniture" defect S7 indicts in
  % fig:comparators row (a). Band (b) already drew its half of the same
  % threshold from its own spine upward; the two now agree.
  \draw[thresh] (\ga{0.006420005919957502},8.54)
             -- (\ga{0.006420005919957502},9.04);
  \node[blab] at (1.50,9.16) {10\% materiality margin, 0.00642 nats};

  % The axis. PASS 10, X7: BOTH ENDS NOW CARRY THE SAME CAP. The right end is
  % the clean preference gap and has always carried a rising 0.22cm ink stub;
  % the left end carried only the same rulegrey hanging tick as 0.02, 0.04 and
  % 0.06, so the span whose full width IS the figure's denominator was
  % delimited at one end by an ink stub and at the other by an ordinary tick,
  % and did not read as a bounded span at all. Two matching caps, no new weight
  % and no new colour.
  % PASS 11, Y3: THE AXIS'S RIGHT END IS COMPUTED, NOT TRANSCRIBED. It was
  % hard-coded 13.995613 here, at the end cap and at the terminal label, while
  % W1 and the layout note both assert that "every data x is still computed
  % inside pgfmath from the full-precision float, so no value is rounded into
  % the source". \ga{0.06420005919957501} = 13.995612905507352, so the literal
  % was short by 3.5e-7cm -- invisible, but it is the ONE coordinate this
  % figure exists to draw, and the claim about it was false. Only the pinned
  % bounding box below still carries a rounded literal, and it carries no
  % value: it is a box corner, declared as such at Y13.
  \draw[spine] (0.00,8.54) -- (\ga{0.06420005919957501},8.54);
  \draw[spine] (0.00,8.54) -- (0.00,8.76);
  \draw[spine] (\ga{0.06420005919957501},8.54) -- (\ga{0.06420005919957501},8.76);
  % PASS 10, X6: the terminal label sits ON ITS OWN CAP, at y = 8.86, and no
  % longer shares the 9.16 baseline and the "<name>, <value> nats" construction
  % with the threshold label 8cm to its left. Read cold on the pass-9 build the
  % two were one row of two labels in one style, so "the clean preference gap,
  % 0.0642 nats" read as the label of a RULE and a reader looked for a rule at
  % 0.0642 and found a 0.22cm end cap. It is now 0.10cm above that cap.
  \node[labR] at (\ga{0.06420005919957501},8.86)
    {the clean preference gap, 0.0642 nats};
  \foreach \v/\t in {0.02/0.02, 0.04/0.04, 0.06/0.06}{
    \draw[tick] (\ga{\v},8.54) -- (\ga{\v},8.44);
    \node[num] at (\ga{\v},8.38) {\t};}
  % Zero is both the axis's left end and the panel's x origin, so its numeral is
  % set flush with that origin rather than centred on it: a centred "0" is the
  % only object that would break the figure's single left edge (S13).
  \draw[tick]  (0.00,8.54) -- (0.00,8.44);
  \node[numL] at (0.00,8.38) {0};

  % AXIS LABEL, in axis-label position: directly under the tick numerals, left
  % edge on the axis's left end, quantity and unit on line one (S5, W2).
  % PASS 10, X13: set lowercase, as 4.11 and 4.13 set every axis label of
  % theirs and as S19's registry gives the name.
  % PASS 11, Y4: THE SCOPE LINE DROPS 0.12cm AND DISCLOSES THE SELECTION.
  % (i) LEADING. It was set 0.28cm under a 9.96pt line, which is 7.94pt of
  % baseline-to-baseline for type 2pt larger than the lead -- negative leading,
  % and the tightest of the four stacked pairs in this figure. Four such pairs
  % were set at or under solid and it is the single thing that made the figure
  % look crowded beside 4.11. Both axis-label blocks and the foot block now
  % open to 0.38cm (10.8pt), and the whole of the added height was paid for out
  % of the white between the two bands, which fell 1.14cm -> 0.77cm.
  % (ii) SELECTION. "layer 12" alone is a scope statement. Layer 12 is the
  % argmax of the anchor seed's seven-depth sweep and the only one of the seven
  % whose interval excludes zero (diff_perm_norm with p_perm: +0.0022 p 0.583,
  % -0.0050 p 0.139, -0.0025 p 0.4495, +0.0041 p 0.3955, -0.0019 p 0.635,
  % +0.0197 p 0.001, +0.0059 p 0.4885 at layers 2/4/6/8/9/12/16). "of seven"
  % is nine characters and it tells a reader on the drawing that the depth was
  % one of a set; WHICH one, and why that matters, is caption material and the
  % proposed caption in figs/_PENDING_CAPTIONS.tex carries it. Until a human
  % integrates that caption the printed one (04b:617-618) says only "one arm at
  % one depth at one rung", which is scope and not selection.
  \node[axlab] at (0.00,7.78) {preference shift, nats};
  \node[scope] at (0.00,7.40)
    {residual arm, layer 12 of seven, $\alpha$ = 0.5};

  % The magnified span, named as an object and not traced by a leader (W10).
  % Its two stubs turn DOWN, towards the panel that opens the span, so the
  % bracket points at (b) instead of back at an axis it already stands under.
  % PASS 10, X8: raised 0.16cm, tight under the scope line, and its label now
  % NAMES the object as well as pointing at (b). Two reviewers asked for it to
  % be moved into the lane immediately under the tick numerals with its stubs
  % reversed; both halves of that were overruled, and the reason is recorded in
  % the pass-10 block under X8.
  % PASS 11, Y5: THE STUBS ARE LONG ENOUGH TO SEE, AND THE BRACKET STANDS
  % CLEAR OF THE AXIS-LABEL BLOCK. Measured on the pass-10 render, each stub
  % was 4 pixels at 200dpi -- 1.44pt, 0.51mm -- of 0.30pt rulegrey, which does
  % not survive print, so the assembly was a bare grey horizontal sitting
  % 0.34cm under a scope line that runs past both of its ends: it read as an
  % underline of "residual arm, layer 12," or as an editorial divider, and not
  % as a bracket. Two changes, no new weight and no new colour: the stubs go
  % 0.06cm -> 0.16cm, and the gap from the scope line goes 0.34cm -> 0.44cm,
  % which is larger than the 0.38cm internal leading of the axis-label block
  % above it, so the bracket stops reading as that block's third line.
  % Two harder proposals were overruled and the reasons are at Y15 (c) and (d).
  \draw[lead] (0.00,6.86) -- (0.00,7.02)
           -- (\ga{0.008025007399946877},7.02)
           -- (\ga{0.008025007399946877},6.86);
  \node[scope] at (1.85,6.96) {the eighth magnified in (b)};

  %% ============= BAND (b) -- THE FIRST EIGHTH, MAGNIFIED =================
  %% PASS 11: the whole of band (b) is raised 0.14cm and its axis lowered
  %% 0.10cm, which is where the leading opened at Y4 and Y11 was paid from.
  \node[head] at (0.00,5.84) {(b) the first eighth, magnified 8\texttimes{}};
  % The threshold label sits against the rule it names and takes the LEFT of it
  % here and the RIGHT of it in (a): in each band the label goes on the side
  % where the rule has room, which in (a) -- where the rule stands at a tenth of
  % the measure -- is the right, and in (b) -- where it stands at four fifths --
  % is the left. Set right of the rule here it measures 2.98cm from 11.30 and
  % crosses the right end of the axis at 13.996 (checked on a build, which
  % reported Overfull \hbox 2.84pt).
  % PASS 10, X6: it has come off the panel head's baseline (5.70) and down onto
  % its own, 0.12cm above the rule's top, in the same relation to its rule as
  % (a)'s label is to (a)'s. On the pass-9 build it sat at exactly the head's
  % baseline, so scanning the two panel heads read "(b) the first eighth,
  % magnified 8x ... 10% materiality margin" as one title line.
  \node[blabR] at (11.10,5.46) {10\% materiality margin};
  \draw[thresh] (\gb{0.006420005919957502},3.68)
             -- (\gb{0.006420005919957502},5.34);

  % Zero. Densely dotted quietink 0.5pt is this document's stroke for a
  % coordinate that decides nothing (S7).
  % PASS 11, Y6: IT IS NAMED, ON THE SAME BASELINE AS THE THRESHOLD LABEL AND
  % AT THE SAME 0.12cm ABOVE ITS OWN RULE'S TOP. SUPERSEDES W8's closing
  % clause, that the rule "carries no label, because its foot is the axis's own
  % 0 numeral". The equivalence criterion of 04b:507-513 is a CONJUNCTION --
  % the interval must lie inside the margin AND contain zero -- and this
  % figure's result is that it does the first and not the second. Pass 10 drew
  % the margin conjunct as the heaviest stroke on the page and the zero
  % conjunct as an unlabelled 0.5pt dotted rule standing on the panel's own x
  % origin, where it has no figure/ground of its own; at a 4x squint it
  % disappeared entirely. Naming it is what the figure's spec authorises ("if
  % the zero rule needs a name it takes 'zero' in \scriptsize roman ink at the
  % rule") and it is the whole of the fix: the rule is NOT promoted to the
  % threshold register, because S6 reserves black solid 0.90pt for a threshold
  % and S7 reserves this stroke for a coordinate that decides nothing -- which
  % zero, as a coordinate, is. What decides is whether the bar clears it, and
  % the bar clears it by 0.70cm.
  % The label is placed with an xshift on the rule's own x, as the row labels
  % are placed on their marks', so it is tied to the thing it names; the two
  % vertical rules of (b) now carry one \scriptsize label each, on one
  % baseline, one at each rule's top (S22).
  % PASS 10, X9: it no longer starts on the spine. Drawn 3.58 to 5.56 it was
  % collinear and continuous with the axis's left terminus, its hanging tick and
  % the panel's x origin, ran the full height of the data region, and matched
  % the threshold rule on the right for length -- so all three reviewers read
  % the pair as the two sides of a frame round the rows, which is the object
  % S12 exists to forbid. It now spans the three rows only, 0.28cm clear of the
  % spine and 0.42cm clear of the panel head, and reads as what it is: the
  % coordinate the interval bar's left end clears by 0.70cm.
  \draw[proj] (0.00,3.94) -- (0.00,5.34);
  \node[qlabz, xshift=3pt] at (0.00,5.46) {zero: no identity effect};

  % ROW ONE, the control: an open disc, white fill, 0.8pt quietink edge, which
  % is 4.11's own encoding for a comparator or a null (S9, CF5, W11).
  % PASS 10, X10: the gloss separator is a comma. It was an em dash, in the same
  % position after the token as row three's math minus, so one glyph class
  % carried "namely" in two rows and "subtract" in the third on a figure whose
  % whole subject is a subtraction. The minus in row three and in band (a) is
  % now the only dash inside a label anywhere on the drawing.
  % X12: the three row labels are placed with an xshift on the mark's own
  % coordinate rather than at a literal x, so each label is tied to its mark.
  % PASS 11, Y7: THE MARK-TO-LABEL GAP GOES 5pt -> 9pt IN ALL THREE ROWS.
  % SUPERSEDES X12's closing clause, which set it at 5pt. At 5pt the open
  % control disc stood 3.28pt from the first glyph of \texttt{permuted}, which
  % is LESS than Inconsolata's own character advance at this size (3.43pt),
  % while its outer diameter (3.6pt) is 11% larger than that face's x-height
  % (3.24pt measured at 200dpi) and its 0.80pt edge is of comparable weight to
  % the stems beside it. The one comparator mark on the drawing therefore set
  % as a lowercase letter o prefixing the word: at 12x it reads "o permuted".
  % Row two escaped only because its disc is filled and accent. At 9pt the gap
  % is 7.2pt, twice the mono advance, and the ring is a mark again.
  % A second fix was proposed for the same defect -- take the ring's edge to
  % ink, as 4.11(b) does -- and it is overruled at Y15 (a).
  \fill[white]  (\gb{0.0006409306049362111},5.24) circle (1.4pt);
  \draw[ctrl]   (\gb{0.0006409306049362111},5.24) circle (1.4pt);
  \node[rowq, xshift=9pt] at (\gb{0.0006409306049362111},5.24)
    {\texttt{permuted}, same construction, wrong drug};

  % ROW TWO, the arm under test.
  \fill[accent] (\gb{0.0019027651015006252},4.66) circle (1.4pt);
  \node[rowa, xshift=9pt] at (\gb{0.0019027651015006252},4.66)
    {\texttt{swap}, points at the drug scored};

  % ROW THREE, the measurement, and the only interval on the figure. The label
  % is anchored on the interval's upper end, which is the rightmost ink in the
  % row, so it clears the bar by the same 9pt the other two clear their discs.
  % PASS 11, Y8: THE VALUE COMES OFF THIS LABEL. It printed "$+0.0013$ nats"
  % 48.9pt (17.2mm) to the right of the disc whose value it is, ON THE ROW'S
  % OWN BASELINE, where position is value: the numeral stood at x = 0.00339 on
  % (b)'s axis, 2.7 times the number it printed and past the 0.0025 tick, on
  % the one panel whose job is to let a reader place the estimate between zero
  % and the 0.00642 margin. Rows one and two anchor at their marks and this
  % row anchored at its bar's end, so one panel carried two anchoring rules and
  % the one that misplaced its label was the one that printed a number. The
  % label is now the row's NAME only, which is the same string band (a) uses
  % for the same JSON key (S19), and the value is printed once, in (a), at a
  % label lifted off the data lane and leadered to its disc (Y2). After this
  % edit no numeral anywhere on the drawing stands on a data row's baseline.
  % PASS 11, Y9: 0.06cm MORE AIR UNDER ROW TWO. The bar on this row spans
  % 0.702 to 3.751cm and the two arm marks stand at 1.118 and 3.318, so the
  % interval on the DIFFERENCE horizontally contains both of its own operands
  % and the difference disc falls within 0.02cm of their midpoint by
  % arithmetic accident. A reader can be offered "a 95% interval that brackets
  % both arms", which is the opposite of what the row says. The row pitch is
  % now 0.58 / 0.64 rather than a flat 0.54, so the derived row sits off the
  % pair it is derived from; the label states the derivation; and the proposal
  % to delete the two arm marks outright is overruled at Y15 (b).
  \draw[bar] (\gb{0.0004024989353392094},4.02)
          -- (\gb{0.002150955667015771},4.02);
  \fill[accent] (\gb{0.0012618344965644142},4.02) circle (1.4pt);
  \node[rowa, xshift=9pt] at (\gb{0.002150955667015771},4.02)
    {\texttt{swap} \textminus{} \texttt{permuted}};

  % The magnified axis, its ticks landing on the same four x as band (a)'s so
  % that the two rules read as one rule at two scales (W13).
  % PASS 11, Y10: BAND (b) IS CAPPED AT ITS RIGHT END, AND ITS END IS COMPUTED.
  % SUPERSEDES W13's closing clause ("band two's right endpoint ... takes no
  % tick and no numeral"). W13 declined a terminal MARK on the ground that
  % (b)'s right end is named by (b)'s axis label rather than by a mark, and
  % that reasoning stands -- but what it left was 26.01pt (9.17mm) of bare
  % rule running past the last drawn tick to an end with nothing on it, which
  % is the same rule-end-to-last-tick defect S13 measures at 8.1mm in
  % fig:comparators row (c). Band (a) has the identical overhang and closes it
  % with a 0.22cm ink cap; band (b) now closes it with the same cap, at the
  % same x, in the same weight and colour, so the two rules terminate
  % identically at the end where the overhang is and read as one rule at two
  % scales. The cap is axis furniture -- the rule's own end -- and not a data
  % mark, so S20 is untouched; the proposal to hang a value on it is overruled
  % at Y15 (e). NO LEFT cap is added in (b): a 0.22cm ink stub rising from the
  % spine at x = 0 would stand 0.14cm under the dotted zero rule and collinear
  % with it, which is exactly the frame-edge reading X9 removed.
  \draw[spine] (0.00,3.68) -- (\gb{0.008025007399946877},3.68);
  \draw[spine] (\gb{0.008025007399946877},3.68)
            -- (\gb{0.008025007399946877},3.90);
  % PASS 11, Y11: the middle numeral prints 0.0050 and not 0.005. Band (a)'s
  % three numerals carry two decimals each and band (b)'s carry four, and the
  % one visible proof that (b) is (a) divided by eight is that the two lanes
  % put congruent numerals on the same four x. A dropped place broke the
  % column of trailing digits that carries it. The VALUE is unchanged; this is
  % a print string, which is why the \foreach takes the value/label pair form.
  \foreach \v/\t in {0.0025/0.0025, 0.005/0.0050, 0.0075/0.0075}{
    \draw[tick] (\gb{\v},3.68) -- (\gb{\v},3.58);
    \node[num] at (\gb{\v},3.52) {\t};}
  \draw[tick]  (0.00,3.68) -- (0.00,3.58);
  \node[numL] at (0.00,3.52) {0};

  \node[axlab] at (0.00,2.92) {preference shift, nats};
  % PASS 10, X14: "magnified 8x" comes off this line. The panel head 3.2cm
  % above already carries it, and the exactness of the eight is caption
  % material; what the axis-label line has to carry is the relation to the axis
  % above, which is what is left.
  \node[scope] at (0.00,2.54) {the first eighth of the axis above};

  %% ==== THE ONE INTERVAL DECLARATION, at the foot (S17, CF8) =============
  %% PASS 10, X2. The pass-9 second line read "No interval is returned for
  %% permuted, for swap, or for (a)." That was false, and false about the one
  %% mark on the figure that does have an interval: band (a)'s disc IS
  %% diff_perm, and layers.12.swap.ci_perm is an interval on diff_perm. In this
  %% document "returned" is what the instrument does, so the sentence told a
  %% reader that band (a)'s estimate has no uncertainty available and, by
  %% implication, that (a) and (b) draw different quantities. The declaration
  %% now separates non-existence from non-drawing: the arms have no stored
  %% interval, the difference has one, and it is drawn once.
  %% Leading is 0.28cm here, as everywhere else in this figure. At pass 9 the
  %% two lines were set 0.24cm apart, tighter than solid for 8pt type, and the
  %% descender of "bootstrap" met the dot of the "i" below it.
  %% PASS 11, Y12: 0.28cm -> 0.30cm between the two lines, and 0.46cm from the
  %% axis-label block above rather than 0.42cm. 0.28cm is exactly solid for
  %% 7.97pt type and the two lines' boxes met; the block also has to stand off
  %% the axis-label block by more than that block's own 0.38cm internal
  %% leading, or the four lines read as one four-line paragraph.
  %% PASS 11: six characters come out of these two lines, and nothing else, to
  %% keep the figure under S15's 5.0 characters per 100mm^2 after "of seven",
  %% "clean" and "zero" went in. Both sentences say exactly what they said.
  \node[scope] at (0.00,2.08) {Intervals: (b) 95\% bootstrap over 58 ordered pair clusters, on the difference, drawn once.};
  \node[scope] at (0.00,1.78) {None is stored for \texttt{permuted} or \texttt{swap} alone, and (a) carries none.};

  % The bounding box IS the body block; nothing may stick out of it.
  % PASS 11, Y13: the floor goes 1.72 -> 1.68 to keep the foot block's
  % descenders inside. This \path is the one rounded literal left in the file
  % and it carries no value: 13.9957 is a box corner, not a coordinate any
  % mark is drawn at, and every data x is computed by \ga or \gb from the
  % full-precision float (Y3).
  \path (0.00,1.68) rectangle (13.9957,10.18);
\end{tikzpicture}
%
% ---------------------------------------------------------------------------
% READY TO PASTE into Sections/04b-representation.tex, at sec:res-mechanism,
% immediately after the verdict paragraph (04b:332-342) or, if the page breaks
% badly, after the equivalence-criterion paragraph (04b:412-419). The same block
% was staged, at pass 9, in a figs/_PENDING_FLOATS.tex that does not exist;
% PASS 10 supersedes that: the float wrapper is HERE and the caption is in
% figs/_PENDING_CAPTIONS.tex, which does exist as of this pass (X15).
%
% PASS 9 NOTE. The caption below is the PROPOSED replacement and it is not yet
% in Sections/. No builder on this pass may edit Sections/04b-representation.tex,
% so the same text is written into figs/_PENDING_CAPTIONS.tex keyed by
% \label{fig:materiality} for a human to integrate. Until that happens the
% caption printed in main.pdf is the pass-8 one, and three of its clauses are
% false of the rebuilt drawing -- see CAPTION NOTE, PASS 9 at the foot.
%
% \begin{figure}[htbp]
% \centering
% % ===========================================================================
% materiality.tex -- fig:materiality. The one surviving identity effect drawn
% twice on one nats ruler: once against the whole preference gap it is divided
% by, where it is a 3.6mm smudge on a 133mm axis, and once at exactly eight
% times, where the swap arm, the permuted control and their difference separate.
%
% FIGURE BODY ONLY. No preamble, no document environment, no float, no caption.
% \input it from inside a figure float in Sections/04b-representation.tex at
% sec:res-mechanism. The float wrapper and caption sit at the foot of this file,
% commented out, ready to paste -- the staging convention of figs/two-scales.tex.
% Compiles against thesis_v2/preamble.tex and nothing else: no \includegraphics,
% no external data at draw time, no TikZ library beyond the ones preamble.tex
% already loads (this file uses decorations.pathreplacing and arrows.meta only).
%
% WHY THIS FIGURE EXISTS. Section 4.5 reports every effect "as a fraction of
% that model's own clean preference gap" (04b:396-398) and NO figure anywhere in
% the document shows that gap. fig:family draws the numerator on a diverging
% colour scale; fig:comparators row C draws it as a fraction. The denominator is
% never drawn, and neither is the subtraction that produces the numerator. This
% figure's subject is the denominator, and the subtraction.
%
% EVERY NUMBER DRAWN, AND WHERE IT COMES FROM
% -------------------------------------------
% ONE arm, ONE layer, ONE rung: the residual arm, layer 12, alpha = 0.5. Source
% for every coordinate: RESULTS_cluster/probe_arm_residual.json, all under the
% key path layers.12.swap. Sibling ledger of the exact values drawn, at full
% machine precision, with key paths: thesis_v2/figs/materiality.json.
%
%   quantity                key path              value drawn (full precision)
%   clean preference gap    .clean_ab_gap         0.06420005919957501
%   10% materiality margin  .equiv_margin         0.006420005919957502
%   permuted arm            .effect_perm          0.0006409306049362111
%   swap arm                .effect_swap          0.0019027651015006252
%   swap - permuted         .diff_perm            0.0012618344965644142
%   its 95% interval        .ci_perm     [0.0004024989353392094,
%                                         0.002150955667015771]
%
% One further quantity is COMPUTED, and is declared as computed in the ledger:
%   one eighth of the gap   .clean_ab_gap / 8     0.008025007399946877
% It is band two's full scale and band one's magnifier bracket, and it is the
% only reason the word "eighth" appears anywhere on the figure.
%
% VERIFIED ON BUILD, not assumed:
%   * .equiv_margin / .clean_ab_gap = 0.1 exactly (config.equiv_frac = 0.1), so
%     the two dashed rules stand at one tenth of band one's full scale.
%   * .effect_swap - .effect_perm = 0.001261834496564414 against .diff_perm
%     = 0.0012618344965644142. They agree to fifteen significant figures and
%     differ in the last bit only, from summation order; the stored .diff_perm
%     is what is drawn. The subtraction the figure draws is therefore the
%     subtraction the pipeline performed, and this was checked, not asserted.
%   * .diff_perm / .clean_ab_gap = 0.019654724813287512 and
%     .ci_perm / .clean_ab_gap = [0.0062694480403512436, 0.03350395145788292],
%     reproducing .diff_perm_norm and .ci_perm_norm to sixteen digits. That is
%     the cross-check that this figure's nats and the thesis's fractions are the
%     same measurement; the normalised triple is printed in the caption and is
%     NOT drawn.
%   * 0.06420005919957501 / 0.0012618344965644142 = 50.87835161772162 and
%     0.006420005919957502 / 0.0012618344965644142 = 5.087835161772163. Both
%     COMPUTED, declared as computed in the ledger; they are the "one fiftieth"
%     and "one fifth" of the closing line.
%
% THE ci_perm TRAP, and why the header names both draws
% -----------------------------------------------------
% The interval drawn is layers.12.swap.ci_perm, the TOP-LEVEL key:
%   [0.0004024989353392094, 0.002150955667015771], p_perm = 0.001,
%   normalising to the thesis's [+0.0063, +0.0335].
% layers.12.swap.ladder[alpha = 0.5].ci_perm stores a SECOND bootstrap draw of
% the same effect,
%   [0.0003982749538768512, 0.002137823621073976], p_perm = 0.004,
% with the point estimate identical to sixteen digits. It normalises to
% [+0.0062, +0.0333] and would print one digit off tab:identity while silently
% swapping the surviving p for the non-surviving one -- 0.001 and 0.004 lie
% either side of the 0.05/26 admitted at rank one (04b:526-531). A builder who
% reads the interval from the ladder because that is where the alphas live makes
% exactly this error. The ladder draw is recorded in materiality.json under
% not_drawn_second_bootstrap_draw and is not drawn here.
%
% AGREEMENT WITH THE TEXT (checked line by line, not assumed):
%   04b:340-341  "\ci{+0.0197}{+0.0063}{+0.0335} of the preference gap the model
%                already has. One to two per cent, against a $10\%$ margin
%                pre-registered for calling an effect material."
%                -> the drawn nats normalise to exactly that triple; the margin
%                   is drawn twice, once per band.                          OK
%   04b:443      "In raw nats the rungs read $+0.0013$" (first rung = the
%                headline) -> row three's label prints $+0.0013$.           OK
%   04b:500-502  "60 rows in 58 pair clusters over 24 anchor drugs"
%                -> row three's second label line prints 58 pair clusters;
%                   60 and 24 are in the caption.                           OK
%   04b:507-513  "The interval must lie inside a pre-registered margin ... and
%                contain zero. ... Under textbook TOST the headline ... would be
%                returned as equivalent, read as a null; this rule returns a
%                directional effect smaller than the margin."
%                -> the two closing ink lines track this sentence, which is the
%                   section's subtlest point and is drawn nowhere else.     OK
%   04b:568      "a fifth or less of the $10\%$ margin"
%                -> the closing accent line prints "one fifth of the margin".OK
%
% DESIGN DECISIONS, each with its reason
% --------------------------------------
%  1. TEXTWIDTH, NOT FULLWIDTH. sec:res-mechanism already carries two fullwidth
%     objects, fig:family and tab:identity, and a fullwidth float forbids a
%     margin note on its page (preamble.tex:522-523). Dropping fullwidth buys
%     this page's margin notes back -- the trade ladder.tex and two-scales.tex
%     both made deliberately. Drawn to 141.9mm against the 142mm measure; the
%     observed band in the corpus is 139.6-141.9mm.
%  2. A RAW-NATS AXIS, WHICH IS LEGITIMATE HERE AND NOWHERE ACROSS ARMS. The
%     five arms emit different target formats and their clean gaps span nearly
%     fivefold, so a single nats rule laid across arms would be one
%     arm's ruler drawn over five arms' data. RECOMPUTED, because the figure
%     brief quoted "4.4x" for this and 4.4 is not the span: over the five
%     probe_arm_*.json files clean_ab_gap runs 0.06420005919957501 (this arm,
%     this layer) to 0.3138025708635225 (consensus, layer 4), a ratio of
%     4.887886, and 4.410308 is the narrower consensus-layer-2 to
%     residual-layer-12 ratio. The caption prints the two endpoint values and
%     the word "nearly fivefold" so a reader can check it. Everything on this
%     figure is one
%     arm at one depth at one rung, so the raw scale is the arm's own and the
%     10% margin is a single rule with a single meaning. fig:family's lower
%     strip is a single-arm object for the same reason.
%  3. THE SAME QUANTITY DRAWN TWICE, NOT TWO QUANTITIES, AND THE MAGNIFICATION
%     IS EXACTLY EIGHT. Band two's full scale is .clean_ab_gap / 8, so the panel
%     is exactly the first eighth of the axis above it and the factor printed on
%     its axis is exactly 8 and not a rounded number. This was arrived at by
%     correction, and the arithmetic is worth recording. The first draft ran
%     band two to 0.0075 nats and called the panel "the first twelfth" at
%     "8.56x": 0.0075 / 0.06420006 = 0.11682, which is one 8.56th and not one
%     twelfth, so the panel's own title was false by a factor of 1.4 and the
%     magnification was a number no reader could check. A true twelfth cannot be
%     drawn at all -- .clean_ab_gap / 12 = 0.005350 nats is SMALLER than the
%     0.006420 margin, so the margin would fall off the right-hand end of the
%     panel and the figure would lose its subject. One eighth is the nearest
%     nameable fraction that still contains the margin, and at that scale the
%     margin sits at 80% of the panel's width, which is where it does the most
%     work. Nothing else changed: every drawn quantity is the same number.
%  4. BAND ONE IS SPARE, BUT ITS MEASUREMENT IS NOW NAMED. Four tick numerals,
%     not seven: the band's job is one impression -- 3.6mm of ink 0.8mm from the
%     origin on a 133mm axis -- and a crowded band invites reading rather than
%     seeing. Every other label lives in band two.
%     REPAIRED, PASS 8, and both reviewers called the old state blocking. The
%     measurement carried NO label at all, on the theory that the magnifier
%     would carry it; its only pointer was the italic "the effect and its
%     interval, magnified below" set 11cm to its right at the far end of a
%     nearly horizontal grey leader, against the document's own rule that a
%     label more than ~0.3cm from its mark gets a leader. One reviewer could not
%     tell whether the mark was a mark, a tick or a printing artefact; the other
%     read it as a stray em-dash. It now carries an accent label, "the effect",
%     centred on the estimate dot 0.22cm above it, and a 0.25pt stem from the
%     dot down to the spine so the mark is anchored to a coordinate rather than
%     floating. The label is two words and not "the effect, +0.0013 nats"
%     because the measure between the axis origin at 0.55 and the margin rule at
%     1.88 holds 1.33cm at \scriptsize while the longer form is 2.55cm and would
%     have crossed the rule; row three of band two prints the value in full,
%     3.5cm below. The italic pointer is deleted with the leader it rode on.
%  5. THE MEASUREMENT SITS 0.22cm ABOVE BAND ONE'S SPINE, not on it. On the
%     spine, a 3.6mm accent bar 0.8mm from the origin reads as a thickening of
%     the axis or as a printing defect. Lifted clear, with its dot over a 2.1pt
%     white halo, it reads as a mark. Checked in the 200dpi render and again in
%     a 550dpi crop: the bar, the dot and the spine are three separable things.
%  6. THE MAGNIFIER'S SOURCE IS A NAMED SPAN, NOT TWO BARE COORDINATES.
%     REPAIRED, PASS 8, reported independently by both reviewers. It used to be
%     a vertical drop at x = 0.55, a second vertical at x = 2.2125, and a line
%     from (2.2125,7.20) to (13.85,6.75) -- 11.6cm horizontal over 0.45cm
%     vertical, about 2 degrees. Read cold, the vertical was a drop from the
%     axis origin and the long line was a rule under band one, and the italic
%     that rode on it annotated band one's axis rather than announcing a
%     magnification. slab.tex's version works because both its leaders leave the
%     magnified REGION at a visible angle, and the fix is to make the region an
%     object first. A rulegrey brace, the corpus's own idiom, now spans exactly
%     the magnified eighth on band one -- \ga{0} = 0.55 to
%     \ga{clean_ab_gap/8} = 2.2125 -- and both leaders leave its two ends, one
%     at 66 degrees and one shallow. THE SHALLOW ONE CANNOT BE FIXED and should
%     not be pretended away: a 1.66cm span is being opened into a 13.90cm panel,
%     so one of the two leaders must run nearly flat whatever is done. What
%     changed is that it now leaves a bracket rather than a bare coordinate, so
%     the assembly reads as source-and-enlargement rather than as a crooked
%     rule. The leaders now land on the panel's own top corners (0.25 and
%     14.15) rather than on band two's data extremes: with the brace naming the
%     source, the leaders' job is to point at the PANEL, and corner-to-corner is
%     what a magnifier looks like. The bracket contains both the effect AND the
%     margin rule, which is itself the point: the whole of section 4.5's
%     argument happens inside the first eighth of the gap.
%  7. THREE ROWS IN THE ORDER PERMUTED, SWAP, DIFFERENCE. Control first,
%     treatment second, measurement third, so the subtraction reads downward in
%     the order the sentence "swap minus permuted" is spoken.
%  8. AN INTERVAL IS DRAWN ON ROW THREE AND ONLY THERE. No interval is stored
%     for either arm alone -- .ci_perm is an interval on the DIFFERENCE -- so
%     rows one and two are bare solid dots. In this document a bare solid dot
%     with no bar is how the absence of an interval is said (comparators.tex
%     changed exactly that mark from a hollow ring to a solid fill, because the
%     absent bar is what carries the meaning). Row three's bar is the document's
%     own convention: a 95% cluster bootstrap over ordered (A,B) pairs.
%  9. THE SIX-MARK TRANSLATION DEVICE IS DELETED, AND THE ROWS CARRY THE
%     CORPUS'S ROW APPARATUS INSTEAD.
%     PASS 8, and the spec pre-authorised this deletion. The device was a
%     hairline from each of the two dots down to the ends of a brace, that
%     brace, a leftward arrow sliding it to the origin, a second brace of equal
%     length from the origin, and a hairline dropping onto row three's dot: five
%     grey structural marks against three measured ones, carrying the identity
%     swap - permuted = the interval's point estimate. It is arithmetically
%     exact -- \gb{effect_swap} - \gb{effect_perm} = \gb{diff_perm} - 0.55 =
%     2.09126 -- and it did not read. Neither brace was labelled with a value,
%     so 0.00190 - 0.00064 = 0.00126 was not recoverable from the marks; the
%     arrow's head landed in white space with no mark at it; and the one
%     sentence that explained the chain, "the same span, measured from zero",
%     sat about 4cm to the right of the arrow it named, at a height between the
%     two braces, attached to nothing. Both reviewers reported having to
%     magnify 3x and reverse-engineer the arithmetic, and one noted that nothing
%     in the construction was accent, so the eye never entered it. Row three's
%     own label already STATES the identity in words -- "swap - permuted" -- so
%     the deletion costs nothing and removes five hairlines. The house rule it
%     follows: an identity must be readable from the marks or not drawn.
%     What replaces it is what fig:comparators gives every row: a rulegrey!55
%     0.3pt divider under each of the three rows, spanning the panel's inner
%     width. Before, the three quantities sat at three heights with no dividers,
%     no stubs and no alignment cues, so the vertical read as a second dimension
%     when it is only a stacking order -- and the higher dot happens to be the
%     SMALLER value, which invites a reader to look for an ordering that does
%     not exist. Rows moved from 5.52/4.94/3.66 to 5.40/4.72/4.04 and the
%     margin brace from 3.10 to 3.30 to take up the slack the deletion left.
% 10. THE MARGIN IS DRAWN TWICE AND BRACED ONCE. On band one it stands 1.33cm
%     from the origin; on band two, 10.64cm from it, and the 7.08cm of clear
%     space between the interval's upper end and the margin rule is the figure's
%     answer to "is it material". The brace under band two names what the margin
%     is a tenth OF -- 0.00642 nats out of 0.0642 -- which no existing figure
%     does.
% 11. THE FRAME IS NAMED ON THE DRAWING and not only in the caption: a quietink
%     line at the foot says which arm, which layer and which rung. This document
%     runs five arms whose clean preference gaps differ by 4.89x, so a raw-nats
%     axis that does not say whose nats they are is under-labelled in exactly
%     the way gate.tex's fifth fix names. \ttfamily on \texttt{residual} because
%     that is how the results files name the arm (ladder.tex's register key).
% 12. WHAT IS NOT DRAWN. The isotropic arm (.effect_iso = 0.0010863731109609438)
%     is stored at this cell and is omitted: 04b:377-380 keeps it "only as a
%     damage diagnostic", it is not the control the headline is measured
%     against, and a third bare dot between rows one and two would read as a
%     second comparison. The alpha ladder is omitted because fig:family's lower
%     strip already draws it and this figure is one rung. p, q, n and the
%     normalised triple are in the caption, not on the drawing.
% 13. REJECTED ALTERNATIVE: a broken axis. Drawing the whole gap and the first
%     eighth on one axis with a break would have saved 3.5cm of height, and it
%     would have destroyed the one thing the figure is for -- the reader has to
%     see the unbroken 133mm and the 3.6mm on it in the same glance before the
%     magnification means anything. two-scales.tex breaks an axis and says in
%     its own header that the break is admissible only because the off-scale
%     value is named in words; here the off-scale value IS the subject.
% 14. REJECTED ALTERNATIVE: bars from zero for the two arm effects. A bar length
%     would have made effect_swap look three times effect_perm, which is true of
%     the numbers and is not the measurement: the measurement is their
%     difference with an interval, and neither arm alone is quoted anywhere in
%     the thesis. Dots keep the arms as positions and the difference as the only
%     thing given a bar.
% 15. REJECTED ALTERNATIVE: naming the 0.0642 axis "chance" or drawing a 0.50
%     rule. There is no chance value on this figure. The reference here is the
%     model's own clean preference gap, an arm-specific measured quantity, so it
%     is drawn as the axis itself and named at the axis's right end rather than
%     as a rule inside it. Noted because the document's two conflicting habits
%     for a chance rule -- solid black in style_v2.chance_line(), 0.5pt densely
%     dotted ink in ladder.tex -- do not apply here, and their absence should be
%     deliberate rather than accidental.
%
% LAYOUT NOTES, all forced by a build
% -----------------------------------
%   * TWO MAPS, both braced pgfmath groups. A shim may not be added outside the
%     braces -- \gb{v}+0.05 is not a coordinate -- so every label offset here is
%     an explicit literal or an xshift on the node (the gate.tex trap). Both
%     maps are used directly as coordinate components, (\gb{v},y), which is
%     comparators.tex's idiom and keeps every data x exact rather than rounded
%     into the source.
%       \ga{0.06420005919957501}  = 13.850005   (band one, the whole gap)
%       \gb{0.008025007399946877} = 13.850005   (band two, one eighth of it)
%     The two right-hand ends coincide to 5 decimal places by construction:
%     1657.32 = 8 x 207.1650 exactly.
%   * Tick numerals use the \foreach \v/\t value/label pair form. \foreach \v
%     alone prints the raw token, so 0.020 would set as 0.02 in one place and
%     0.0 in another; the pair form pins the printed string.
%   * inner sep = 1pt ON THE WHOLE PICTURE, as gate.tex sets it. The default
%     0.3333em is 3.65pt at the ambient \small, and at pass 2 that invisible
%     padding -- not the ink -- was what collided: the band-one word tag's box
%     and the margin label's box overlapped while their glyphs were 1.2mm
%     apart. With inner sep pinned, every clearance below is box arithmetic
%     rather than judgement. The consequence is that a label placed at
%     mark + 0.30 would sit 0.10cm closer to its mark than comparators.tex's,
%     so every row label here is placed at mark + 0.40, which puts the glyph
%     0.39cm clear of the marker edge -- comparators' measured value.
%   * TICK NUMERALS CARRY fill=white AND ARE DRAWN LAST IN THEIR BAND. Pass 1
%     ran the magnifier's left leader straight through band one's "0", because
%     band one's origin and the magnifier's left edge are the same x = 0.55.
%     The white box interrupts the rule across the numeral instead --
%     comparators.tex's idiom for its zero rule, done with a box rather than by
%     hand because the numerals are drawn in a \foreach. Drawing order matters:
%     spine, then magnifier, then numerals.
%   * THE PANEL IS 3mm WIDER THAN ITS DATA ON EACH SIDE, (0.25,1.50) to
%     (14.15,6.75) around data running 0.55 to 13.85. At pass 2 the panel ran
%     exactly 0.55 to 13.85 and its left dashed border was drawn at the same x
%     as band two's dotted zero rule; the two were indistinguishable in the
%     600dpi crop, so the figure's zero coordinate was reading as the frame of
%     its own panel. Widening the frame separates them by 3mm. PASS 8 re-aimed
%     the magnifier's leaders at the panel's own top corners (0.25 and 14.15),
%     which is the opposite of what pass 2 did and is right for the same reason
%     the brace is right: with the brace naming the SOURCE span exactly, the
%     leaders' job is to point at the PANEL, and corner-to-corner is what a
%     magnifier looks like.
%   * BAND TWO'S WORD TAG AND ITS MARGIN LABEL SIT INSIDE THE PANEL, on one
%     baseline at y = 6.44, and the axis title sits inside it too. At pass 2 the
%     tag, the margin label, "no movement" and the axis title all straddled the
%     dashed border. Anything that names something inside the panel is now
%     inside the panel.
%   * The band-two tag is anchored at x = 0.75 and not at x = 0.00, because the
%     magnifier's left leader is a vertical at x = 0.55.
%   * PASS 8: the translation device and its label are deleted (decision 9),
%     which took five hairlines and one floating sentence out of the panel.
%     Rows moved up to take the slack -- 5.52/4.94/3.66 became 5.40/4.72/4.04,
%     and the margin brace 3.10 became 3.30 -- and each row gained a rulegrey!55
%     0.3pt divider under it, spanning 0.35 to 14.05, which is 0.10cm inside the
%     panel's dashed frame at both ends.
%   * PASS 8, THE MARGIN LABEL STAYS ON THE 8.44 ROW. The build that moved it up
%     to the tag's row put it at x = 1.96 and "THE WHOLE GAP" measures 1.95cm of
%     \sffamily\scriptsize\bfseries from x = 0, so the two ran together as
%     "THE WHOLE GAP0% materiality margin" in the render. Both labels sit on the
%     8.44 row instead and the dashed margin rule at x = 1.88 separates them:
%     "the effect" is centred on the estimate dot and ends at 1.29, the rule is
%     at 1.88, the margin label starts at 1.98. The rule stops at 8.56, which is
%     0.09cm under the tag's descender line and is why it may not be carried
%     higher to meet a label on the tag's row.
%   * BAND ONE'S RIGHT END is a terminal tick RISING from the spine, 0.30cm
%     tall against the 0.10cm hanging axis ticks, with its two label lines
%     right-aligned on the same x = 13.85 so the tick reads as their anchor. At
%     pass 3 that rule ran from the spine up to the second label's baseline and
%     read as a bracket enclosing the label block.
%   * The closing prose was three lines at pass 2 and is two at pass 3; the
%     28pt that bought went into the clearances above.
%   * Every prose node is a single \scriptsize line anchored base west. No
%     `align=` and no `text width=` anywhere: those build a minipage whose first
%     line follows the AMBIENT \baselineskip, so the node's ink moves with the
%     prose before the float (frames.tex measured a 3.7pt drift).
%   * The bounding box is pinned with a bare \path so the measure is exact and
%     nothing is clipped: (0.0000,0.30) to (14.1900,8.98) = 141.9 x 86.8mm.
%
% TYPE. Ambient \small on the picture; every node sets \scriptsize = 8pt, which
% is 0.67 of the 12pt body and below the 11pt caption. \ttfamily on the two arm
% names and on the arm named at the foot, because that is how the results files
% name them. Nothing is scaled: 1cm here is 1cm on the page.
% ===========================================================================
%
% ===========================================================================
% PASS 9 -- 2026-08-18. REBUILD AGAINST THE STYLE CONTRACT AND THE SUPERVISOR'S
% VERDICT ("no definition on the x axis for what we are measuring"; "make it
% look more professional, right now it looks a bit sloppy and childish").
% NOTHING ABOVE IS DELETED. The design rationale of passes 1-8 is the audit
% trail and stays; what follows supersedes, EXPLICITLY AND BY NAME, each clause
% of it that this pass made false. The reference standard adopted is fig 4.11
% (figs/fig-mechanism.pdf, printed page 51) and, for the threshold, fig 4.13
% (figs/family.pdf), both of which were rendered and read side by side with this
% figure during the build.
%
% NO NUMERIC VALUE CHANGED IN THIS PASS. Every quantity drawn or printed is the
% same float from the same key path as in pass 8: clean_ab_gap 0.06420005919957501,
% equiv_margin 0.006420005919957502, effect_perm 0.0006409306049362111,
% effect_swap 0.0019027651015006252, diff_perm 0.0012618344965644142,
% ci_perm [0.0004024989353392094, 0.002150955667015771], n_pair_clusters 58.
% The one printed numeral that left the drawing is "one fiftieth / one fifth"
% (a ratio in words), and it went to the caption, where it already was.
%
% W1. SUPERSEDES LAYOUT NOTE "TWO MAPS" AND THE MAP CONSTANTS RECORDED THERE.
%     Both maps lose their 0.55cm left inset and gain a longer barrel:
%         old   \ga#1 = {0.55+(#1)*207.1650}   \gb#1 = {0.55+(#1)*1657.3200}
%         new   \ga#1 = {(#1)*215.7500}        \gb#1 = {(#1)*1726.0000}
%     Two things are preserved exactly and were recomputed on the build rather
%     than assumed: 1726.0000 = 8 x 215.7500 EXACTLY, so the magnification is
%     still exactly eight and not a rounded number (this is the correction made
%     at pass 4 and it survives untouched); and both bands' right-hand ends
%     still coincide, now at \ga{0.06420005919957501} = \gb{0.008025007399946877}
%     = 13.851162772308308cm. Every data x is still computed inside pgfmath from
%     the full-precision float, so no value is rounded into the source.
%     WHY: (i) S13, one x origin per panel. With the old map the axis began at
%     0.55 while every text block began at 0.00, so the panel had two left edges
%     0.55cm apart and the left column measured three (58.0, 63.9, 70.4pt).
%     Now the axis's left end, the panel titles, both axis-label blocks, the
%     span bracket and the foot line all sit on x = 0.0000. (ii) S5 puts the
%     axis-label line under the tick numerals and LEFT-ALIGNED TO THE AXIS'S
%     LEFT END; that is only the same thing as the panel origin if the two
%     coincide. (iii) S14, measure. Deleting the dashed enclosure (W4) took away
%     the object that used to set the width, and a 133.0mm barrel inside a
%     textwidth figure measures about 134mm of content, under the 138mm floor.
%     215.75 puts the barrel at 138.51mm and the content bbox at 139.6mm, inside
%     138-142, with the widest object the axis itself and the data region 100%
%     of the drawn width.
%     The old map is kept in materiality.json beside the new one, per S25.
%
% W2. AXIS DEFINITION, WHICH IS THE SUPERVISOR'S FIRST COMPLAINT AND THE REASON
%     THIS PASS EXISTS. Pass 8 named neither axis. Band one carried bare 0 /
%     0.02 / 0.04 / 0.06 with no quantity and no unit anywhere on it -- the unit
%     reached the reader only through the terminal label -- and band two carried
%     "nats, magnified 8x from the axis above", which names a unit and a
%     transformation but no quantity. 4.11 sets "probe accuracy", "descriptive
%     subspace fraction", "KL from the clean forward pass" and "layer" in
%     axis-label position five times over. Each band now carries, directly under
%     its tick numerals and left-aligned to its axis's left end:
%         line 1, \footnotesize roman ink:  Preference shift, nats
%         line 2, \scriptsize quietink:     (a) residual arm, layer 12, alpha =
%                                               0.5; the full width of this rule
%                                               is the model's clean preference
%                                               gap
%                                           (b) the first eighth of the axis
%                                               above, magnified 8x
%     The quantity is repeated deliberately: two scales, each self-describing.
%     "Preference shift" is the section-4.5 registry name (S19), so the axis, the
%     caption and fig:family's lower strip now use one string for one quantity.
%
% W3. THE THRESHOLD. SUPERSEDES DESIGN DECISION 10 AND THE marginrule STYLE.
%     The 10% pre-registered materiality margin was accent 0.60pt dashed here,
%     ink 0.50pt densely dotted in fig:comparators row C, and solid black 0.90pt
%     in fig:family's lower strip: one pre-registered threshold, three renderings
%     inside nine pages. It is now BLACK, SOLID, 0.90pt in both bands, labelled
%     \scriptsize roman black, which is style_v2.chance_line()'s own hard-coded
%     register and fig:family's. This is the highest-priority change in the set
%     and it is a correctness repair, not a preference: accent means the object
%     under test (S8), and a pre-registered bound is not under test.
%     Design decision 15's note -- that the document's two habits for a chance
%     rule "do not apply here" -- is now false as written and is superseded:
%     there is still no chance value on this figure, but the black-solid-0.9pt
%     register does apply, to the margin.
%
% W4. DELETED: THE DASHED ROUNDED ENCLOSURE round band two (390 x 149pt), with
%     the layout note that sized it 3mm wider than its data on each side. S12
%     forbids a dashed rectangle drawn round a panel outright, and it was this
%     figure's only rectangle. Band two now sits directly under band one in
%     4.13's idiom -- the magnified axis under the full axis, with the
%     magnification named in axis-label position -- which is also what fixes
%     band two's axis under W2. Two consequences recorded so they are not read
%     as accidents: the layout note about the panel's left border colliding with
%     band two's zero rule is moot, since there is no border; and the note that
%     band two's tag and margin label "sit inside the panel" is moot for the same
%     reason.
%
% W5. DELETED: THE FOUR PALE ROW SEPARATORS AND THE MARGIN BRACE. SUPERSEDES
%     DESIGN DECISION 9's replacement clause and the pass-8 note that added them.
%     Five near-identical thin horizontals stood inside 30mm of vertical space --
%     three rulegrey!55 row dividers, one rulegrey brace carrying "10% of the gap
%     = 0.00642 nats", plus the axis rule -- and the one that carried meaning was
%     indistinguishable from the ones that did not. The dividers go because the
%     rows no longer need them: with the enclosure gone and one x origin, the
%     three rows read as a stack. The brace goes because the same fact is now
%     carried without any stroke at all, by printing the margin's value in band
%     one's threshold label ("10% materiality margin, 0.00642 nats") beside band
%     one's terminal label ("the clean preference gap, 0.0642 nats"): the tenth
%     is then a division a reader performs on two printed numbers 8cm apart
%     rather than a brace they must trace.
%
% W6. THE INTERVAL IS DRAWN ONCE. SUPERSEDES DESIGN DECISION 4 and the pass-8
%     repair that labelled band one's bar. Pass 8 drew the same estimate as a
%     3.6mm bar-with-dot in band one AND as a 30mm bar-with-dot in band two,
%     while the caption said "it is the one bar on the figure"; a naive reader
%     counted two and spent time working out which two quantities they were.
%     Rule applied: an interval is drawn at the scale where it is legible and
%     nowhere else. Band one keeps a single filled accent disc for the estimate,
%     with no bar; band two carries the bar. Band one's dot keeps its label and
%     its leader, and the label is now the SAME STRING as band two's row three,
%     "swap - permuted", because S19 forbids two names for one JSON key and pass
%     8 had "the effect" here and "swap - permuted" there.
%
% W7. DELETED: THE TWO BANNERS, replaced by 4.11's panel heads. "THE WHOLE GAP"
%     and "THE FIRST EIGHTH, MAGNIFIED" were \sffamily\scriptsize\bfseries accent
%     caps -- the preamble's margin-apparatus face, which belongs to the margin
%     and not to a body figure (S2). They are now "(a) the whole preference gap"
%     and "(b) the first eighth, magnified 8x", \footnotesize roman ink,
%     lowercase, left-aligned on the panel origin, no colon and no bold. The
%     tag/.style is deleted rather than left unused.
%
% W8. DELETED: THE ITALIC "no movement". It sat unattached above the permuted
%     row and a naive reader took it for an assertion about that arm rather than
%     a name for the zero rule. The zero rule survives -- densely dotted quietink
%     0.5pt, which is exactly S7's reserved register for a coordinate that
%     decides nothing -- and carries no label, because its foot is the axis's own
%     "0" numeral. With it goes the last italic on the drawing: italic share is
%     now 0%, against 10% at pass 8 and 0% in 4.11.
%
% W9. DELETED: THE CLOSING PROSE BLOCK AND THE ACCENT APHORISM. The two ink lines
%     ("The interval excludes zero and lies inside the margin. Textbook TOST
%     returns that as equivalence and reads it as a null; this rule returns a
%     directional effect smaller than the margin.") and the accent italic ("one
%     fiftieth of the gap; one fifth of the margin.") are argument, and in this
%     document the argument lives in the caption (S3, S4, S18). At print the
%     reader met an 8pt caption, a blank, then the real 10.9pt caption -- two
%     captions, the smaller one first. Both lines are written into
%     figs/_PENDING_CAPTIONS.tex, keyed by \label, for a human to integrate;
%     the aphorism is already the caption's lede. The TOST asymmetry is the
%     section's subtlest point and it is NOT lost: it is caption material, and
%     the configuration it describes -- an interval that excludes zero and lies
%     wholly inside the margin -- is still visible in band two, which is the only
%     part of that sentence a drawing can carry.
%     SUPERSEDES the text_agreement_check on 04b:507-513 in the ledger, which
%     asserted the drawing restates that sentence. It no longer does.
%
% W10. THE MAGNIFIER. SUPERSEDES DESIGN DECISION 6 AND THE PASS-8 BRACE-AND-TWO-
%     LEADERS CONSTRUCTION. What is kept: the source span is a NAMED OBJECT and
%     not two bare coordinates. What is deleted: both leaders, including the
%     shallow one that pass 8's own header admitted "CANNOT BE FIXED" and that
%     the naive reader read as a trend line before reading it as a leader.
%     The reasoning, recorded because a later builder will be tempted to put them
%     back. A leader pair from the ends of band one's [0, gap/8] segment to the
%     ends of band two's rule is geometrically forced to be asymmetric -- a
%     17.3mm span opening into a 138.5mm rule -- so the right-hand leader runs at
%     under 4 degrees whatever is done, and 4 degrees of accent-free hairline
%     running the full measure is a trend line to every eye that meets it first.
%     Worse, it cannot be routed: S5 now requires band one's two-line axis-label
%     block directly under band one's tick numerals, which puts full-width text
%     in the only lane the leaders could cross. The magnification is instead
%     carried by two objects that cannot be misread: a 0.30pt rulegrey bracket on
%     band one spanning exactly \ga{0} = 0.0000 to \ga{clean_ab_gap/8} = 1.731395,
%     labelled at its own end ("the eighth magnified below"), and band two's
%     axis-label line 2, which names the magnification in axis-label position
%     where S5 puts it. The bracket's two ends stand at the same x as band two's
%     rule origin, so the vertical alignment does the joining that a hairline was
%     doing badly.
%
% W11. MARK VOCABULARY (S9, CF5). SUPERSEDES DESIGN DECISION 8's claim that a
%     bare SOLID dot is how this document says "no interval". It is not the only
%     thing a dot has to say. The permuted arm is the CONTROL, and 4.11's own
%     encoding for a comparator or a null is an open disc, white fill, 0.8pt edge
%     in its colour (mfc='white' on the raw probe and on the random null). Row one
%     is therefore an open quietink disc; rows two and three and band one's
%     estimate are filled accent discs, all at radius 1.4pt, one radius on the
%     figure. The absence of an interval is now said where the contract puts it,
%     in the single \scriptsize interval declaration at the foot (S17), which
%     names the bootstrap construction for the panel that has a bar and says in
%     as many words that the two arms and band one carry none. The 0.8pt disc
%     edge is a MARK edge, not a rule weight; the figure's rules use exactly five
%     weights, 0.30 / 0.40 / 0.50 / 0.90 / 1.00pt.
%     TWO SINGLETONS ARE DECLARED HERE RATHER THAN REMOVED, as S10 permits. The
%     0.50pt densely dotted stroke is used once, for band two's zero rule, and is
%     a distinct kind of object -- S7 reserves that stroke for a coordinate that
%     decides nothing, and this figure has exactly one such coordinate. The
%     1.00pt stroke is used once, for the interval bar, and is a distinct kind of
%     object for the same reason: this figure stores exactly one interval.
%
% W12. TYPE (S1, S16). Two sizes above 7pt and one face: \footnotesize (10.0pt)
%     for the two panel heads, both axis-label line ones, and nothing else;
%     \scriptsize (8.0pt) for tick numerals, direct labels, scope lines, the
%     threshold labels and the interval declaration. No \sffamily, no \bfseries,
%     no \itshape anywhere. \ttfamily survives on the four arm tokens
%     (residual, permuted, swap) because that is how the results files name them
%     and S1 exempts the \texttt token name. Pass 8 printed four sizes and set
%     its largest type on a grey italic epigram; the largest type now belongs to
%     a panel head or an axis label, which is the whole of CF6.
%
% W13. TICK LANES MADE CONGRUENT (S13). Band one ticks at 0 / 0.02 / 0.04 / 0.06
%     land on x = 0.0000 / 4.3150 / 8.6300 / 12.9450. Band two's ticks were
%     chosen so that they land on THE SAME FOUR x: 0 / 0.0025 / 0.005 / 0.0075,
%     which is possible only because band two is exactly one eighth of band one
%     and is a visible consequence of the exact 8. Both rules share both
%     endpoints (0.0000 and 13.851163) and both first ticks sit on the left
%     endpoint. Band one alone carries a terminal tick at the right endpoint,
%     rising 0.20cm from the spine, because that endpoint is the clean preference
%     gap and is the figure's denominator; band two's right endpoint is one
%     eighth of it and is named by band two's axis label rather than by a mark,
%     so it takes no tick and no numeral -- no unlabelled mark is left on the
%     drawing (S20).
%
% W14. MEASURE AND HEIGHT, MEASURED ON THE FINAL BUILD, not estimated: content
%     bbox 140.34 x 83.82mm, ink 5.14% of pixels below 230 grey at 200dpi, 4.88
%     characters of figure text per 100mm^2, one page, zero Overfull and zero
%     Underfull boxes. Against pass 8's 141.2 x 85.9mm at 7.03% ink and against
%     the contract's 138-142mm measure, 85mm height cap, 6.0% ink ceiling and
%     5.0 characters/100mm^2 ceiling. Every millimetre of the height and every
%     point of the ink came from W4, W5, W8, W9 and W10 -- deletions -- and
%     nothing measured was removed to get them.
%
% W15. STALE CLAIMS ELSEWHERE IN THIS HEADER, SUPERSEDED BY NAME so that an
%     auditor reading top to bottom is not misled.
%     (i) The opening paragraph's "where it is a 3.6mm smudge on a 133mm axis"
%         is now false twice over: the axis is 139.96mm, and band (a) no longer
%         draws the interval at all, so what stands 2.75mm from the origin is a
%         single 1.4pt disc and not a 3.6mm bar (W6).
%     (ii) The file-header line "this file uses decorations.pathreplacing and
%         arrows.meta only" is now false in the harmless direction: this figure
%         uses NEITHER. The braces went with W5 and W10 and there was never an
%         arrowhead. It still loads no library preamble.tex does not.
%     (iii) DESIGN DECISION 1's "Drawn to 141.9mm against the 142mm measure" is
%         superseded by W14's 140.34mm. The decision itself -- textwidth and not
%         fullwidth, to buy the page's margin notes back -- stands unchanged.
%     (iv) DESIGN DECISION 3's arithmetic is intact and was recomputed on this
%         build; only the cm-per-nat constants moved (W1). The first draft's
%         "first twelfth at 8.56x" correction, and the reason a true twelfth
%         cannot be drawn at all, remain exactly as recorded.
%     (v) The AGREEMENT WITH THE TEXT entry for 04b:507-513 ("the two closing
%         ink lines track this sentence") is superseded by W9: those lines are
%         deleted from the drawing and staged for the caption. The entry for
%         04b:568 ("the closing accent line prints one fifth of the margin") is
%         superseded for the same reason. Both sentences remain TRUE of the
%         figure's subject; they are no longer drawn on it.
%
% W16. TWO DEVIATIONS FROM THE SPEC'S LITERAL AXIS-LABEL STRINGS, both forced by
%     other clauses of the same contract, both recorded rather than quietly
%     taken.
%     (i) Band (a)'s scope line is set as "residual arm, layer 12, alpha = 0.5"
%         and NOT as "residual arm, layer 12, alpha = 0.5; the full width of
%         this rule is the model's clean preference gap". The deleted clause
%         restates, in twelve words, the terminal label "the clean preference
%         gap, 0.0642 nats" that the spec's own keep-list requires at the axis's
%         right end, and S18 forbids a drawing from carrying one clause twice.
%         It also costs 70 characters, and at 70 more characters the figure
%         needs 89.8mm of height to stay under the 5.0 characters/100mm^2
%         ceiling, which is past the 85mm cap. The quantity and its unit are on
%         line one, which is what S5 actually requires; what came off is a
%         second statement of what the axis's right end already says.
%     (ii) Band (b)'s scope line is set verbatim as the spec gives it, "the
%         first eighth of the axis above, magnified 8x", and the 14 characters
%         that costs were bought by opening 3mm more air between the bands
%         rather than by cutting anything else.
%     Line one of both blocks is verbatim: "Preference shift, nats". Recorded
%     because 4.11 and 4.13 both set their axis labels entirely lowercase and
%     4.10 sets its capitalised, so the slate is split on case and the spec's
%     capitalisation was followed rather than guessed at.
% ===========================================================================
% ===========================================================================
% PASS 10 -- 2026-08-18. REPAIR PASS on the pass-9 rebuild. Three adversarial
% readers went over the built figure independently and returned thirty-three
% findings. NOTHING ABOVE IS DELETED: passes 1-9 are the audit trail and stay
% as written. Every clause of pass 9 that this pass made false, and every
% number pass 9 recorded that this file has never drawn, is superseded here BY
% NAME (S25). The findings that were OVERRULED are recorded under X16 with the
% reason, so that a later builder does not re-apply them blind.
%
% NO NUMERIC VALUE CHANGED IN THIS PASS EITHER. Every quantity drawn or printed
% is the same float from the same key path as at passes 8 and 9:
% clean_ab_gap 0.06420005919957501, equiv_margin 0.006420005919957502,
% effect_perm 0.0006409306049362111, effect_swap 0.0019027651015006252,
% diff_perm 0.0012618344965644142, ci_perm [0.0004024989353392094,
% 0.002150955667015771], n_pair_clusters 58, and the exact magnification 8.
% Both maps are untouched: \ga = (#1)*218.0000 and \gb = (#1)*1744.0000, with
% 1744.0000 = 8 x 218.0000 exactly. What moved is where labels sit, where two
% rules start, how long two others are, and what four sentences say.
%
% X1. STALE CONSTANTS IN THE PASS-9 BLOCK, CORRECTED BY NAME. W1, W10, W13 and
%     W14 record a map, a shared endpoint, a bracket end, four tick positions
%     and a measure that this file has never drawn. materiality.json has always
%     recorded the drawn ones correctly, and the drawing is not affected --
%     every data x is computed inside pgfmath from the full-precision float --
%     but S25 makes the header an audited artefact, so an auditor recomputing a
%     position from W1 must not be handed a figure that does not exist. The
%     corrections, each superseding the clause named:
%       W1  "new \ga#1 = {(#1)*215.7500}  \gb#1 = {(#1)*1726.0000}"
%           -> the drawn map is \ga#1 = {(#1)*218.0000} and
%              \gb#1 = {(#1)*1744.0000}. Both properties W1 claims survive and
%              were rechecked on this build: 1744.0000 = 8 x 218.0000 exactly
%              (8*218.0 == 1744.0 is True in IEEE double), and both bands' right
%              ends coincide.
%       W1  "= 13.851162772308308cm" for the shared right end
%           -> 13.995612905507352cm = 0.06420005919957501 x 218, which is what
%              is hard-coded at the spine, at the two end caps and at the
%              terminal label, and what materiality.json records.
%       W1  "215.75 puts the barrel at 138.51mm and the content bbox at 139.6mm"
%           -> the barrel is 139.956mm and the content bbox 139.96mm, measured
%              on this build. Both are inside S14's 138-142mm.
%       W10 "\ga{clean_ab_gap/8} = 1.731395" for the bracket's right end
%           -> 1.749451613188419cm = 0.008025007399946877 x 218.
%       W13 "land on x = 0.0000 / 4.3150 / 8.6300 / 12.9450"
%           -> 0.0000 / 4.3600 / 8.7200 / 13.0800, for both bands, which is
%              what makes the two tick lanes congruent.
%       W13 "Both rules share both endpoints (0.0000 and 13.851163)"
%           -> (0.0000 and 13.995613).
%       W14 "content bbox 140.34 x 83.82mm, ink 5.14%, 4.88 characters per
%           100mm^2" -> superseded by X17's measurement of THIS build.
%     No 215.75 build was ever produced; the pass-9 numbers are a draft's, not
%     a superseded state of the drawing, and they are recorded here as wrong
%     rather than as an earlier truth.
%
% X2. THE ONE INTERVAL DECLARATION SAID SOMETHING FALSE ABOUT THE INSTRUMENT,
%     AND SAID IT ABOUT THE ONE MARK THAT HAS AN INTERVAL. SUPERSEDES W11's
%     description of the declaration and the ledger's intervals.declaration_on_
%     the_drawing. Pass 9's second line read "No interval is returned for
%     permuted, for swap, or for (a)." Band (a)'s single disc IS diff_perm, and
%     layers.12.swap.ci_perm is an interval on diff_perm; it is the very
%     interval drawn 8x larger in (b). "Returned" is this document's word for
%     what the instrument stores (S17), so the sentence asserted that the
%     band-(a) estimate has no uncertainty available and implied that (a) and
%     (b) draw different quantities. The declaration now separates
%     non-existence from non-drawing:
%       "95% bootstrap over 58 ordered pair clusters, on the difference; drawn
%        once, in (b)."
%       "No interval is stored for permuted or for swap alone."
%     Both lines are true of the file, they still fit S17's one block of at
%     most two lines, and the reason band (a) carries no bar is now the same
%     reason the ledger has always given -- legibility -- rather than a claim
%     about the instrument.
%
% X3. THE THRESHOLD IN (a) STANDS ON ITS SPINE. SUPERSEDES the pass-9 y extent
%     8.38 .. 9.04. Drawn from 8.38 it crossed the axis rule and hung 0.16cm
%     below it, past the ends of the hanging ticks and into the tick-numeral
%     lane, so the heaviest stroke in the panel doubled as an outsized
%     unnumbered tick standing between the numerals "0" and "0.02" -- exactly
%     the "the threshold merges with the tick furniture" defect S7 indicts in
%     fig:comparators row (a). It now runs 8.54 .. 9.04, from the spine upward,
%     which is what (b)'s half of the same threshold has always done. One
%     pre-registered threshold, one stroke AND one geometry, in both bands.
%
% X4. BAND (a)'s ONE MARK IS NOW LABELLED AT ITSELF. SUPERSEDES W6's placement.
%     Pass 9 anchored "swap - permuted" base west on x = 0.00, the panel
%     origin, while the leader fell from \ga{diff_perm} = 0.2751 -- so the
%     leader descended from inside the word "swap", and the mark it landed on
%     read as the swap arm, whose value is 0.0019 and not the 0.0013 drawn.
%     That is the only mark in band (a). The label now begins on the mark's own
%     x and the leader falls from its first character.
%
% X5. THE LEADER AND THE PROJECTION TOUCH WHAT THEY POINT AT. The accent leader
%     stopped at 9.12 against a disc whose top edge is 8.9493 -- 4.84pt of
%     white -- so at print it was a detached accent tick under the label, of
%     the same length, weight and orientation as an axis tick. The dotted
%     projection was detached at both ends, 2.0pt below the disc and 0.57pt
%     above the spine. Leader 9.46 -> 8.95 lands on the disc; projection
%     8.85 -> 8.54 leaves the disc's lower edge and reaches the rule.
%
% X6. THREE LABELS THAT READ AS ONE BLOCK ARE SEPARATED, and no label sits on a
%     panel head's baseline any more. SUPERSEDES the pass-9 placements.
%     (i) Band (a)'s accent label moved 9.44 -> 9.56. At 9.44 it stood 7.94pt
%         above the threshold label -- set solid for 7.97pt type -- and the two
%         overlap horizontally by 0.6cm, so they read as one ragged two-line
%         paragraph headed by an accent line. The gap is now 11.3pt.
%     (ii) Band (a)'s terminal label moved off the 9.16 baseline onto 8.86,
%         0.10cm above the end cap it names. At 9.16 it shared a baseline AND
%         the "<name>, <value> nats" construction with the threshold label 8cm
%         to its left, so "the clean preference gap, 0.0642 nats" read as the
%         label of a RULE and a reader looked for a rule at 0.0642 and found a
%         0.22cm cap. It is now unmistakably the label of the axis's own end.
%     (iii) Band (b)'s threshold label moved 5.70 -> 5.44, off the panel head's
%         exact baseline, and its rule's top came down 5.56 -> 5.32 to meet it.
%         On the pass-9 build, scanning the two panel heads read "(b) the first
%         eighth, magnified 8x  ...  10% materiality margin" as one title line.
%         Both threshold labels now stand 0.12cm above their own rule's top.
%
% X7. THE DENOMINATOR IS DELIMITED AT BOTH ENDS. The gap -- the span whose full
%     width is the quantity every effect in section 4.5 is divided by -- was
%     capped at the right by a rising 0.40pt ink stub and at the left by
%     nothing but the same rulegrey hanging tick that 0.02, 0.04 and 0.06 carry,
%     so it did not read as a bounded span at all. A matching ink stub now
%     rises at x = 0. No new weight, no new colour, no height.
%
% X8. THE MAGNIFIER BRACKET: RAISED, AND ITS LABEL NAMES THE OBJECT. It moved
%     7.14 -> 7.22 (rail), tight under the scope line at the clearance the
%     scope line's descenders allow, and its label became "the eighth magnified
%     in (b)" from "magnified in (b)", so the mark is named as well as pointed
%     (S20). TWO HALVES OF WHAT WAS ASKED WERE OVERRULED; see X16 (b) and (c).
%
% X9. BAND (b)'s ZERO RULE NO LONGER READS AS A PANEL EDGE. SUPERSEDES the
%     pass-9 extent 3.58 .. 5.56. Drawn from the spine it was collinear and
%     continuous with the axis's left terminus, that terminus's hanging tick
%     and the panel's x origin, ran the full height of the data region, and
%     matched the threshold rule on the right for length to the point -- so all
%     three readers took the pair for the two sides of a frame round the rows,
%     which is the object S12 exists to forbid, and one noted that band (a)
%     draws no such rule at its own zero. It now runs 3.86 .. 5.28: it spans
%     the three rows only, stands 0.28cm clear of the spine and 0.42cm clear of
%     the panel head, and reads as the coordinate the interval bar's left end
%     clears by 0.70cm. It is still S7's register and still unlabelled, because
%     the axis's own "0" numeral stands at its foot.
%
% X10. ONE DASH, ONE MEANING. Rows one and two glossed their arm with an em
%     dash in the same position after the token where row three and band (a)
%     set a math minus, so one glyph class carried "namely" twice and
%     "subtract" once on a figure whose entire subject is a subtraction. The
%     two glosses now take a comma. The minus in "swap - permuted" is the only
%     dash inside any label on the drawing.
%
% X11. THE FOOT BLOCK'S LEADING. The two lines were set 0.24cm apart, which is
%     6.81pt of leading for 7.97pt type -- tighter than solid, tighter than
%     every other stacked pair in this figure, and tight enough that the
%     descender of "bootstrap" met the dot of the "i" below it. They are now
%     0.28cm apart, the figure's own value everywhere else, and the pinned
%     bounding box floor went 1.76 -> 1.72 so descenders stay inside it.
%
% X12. ONE X ORIGIN, MEASURED RATHER THAN ASSERTED. SUPERSEDES the layout note
%     "inner sep = 1pt ON THE WHOLE PICTURE" in its horizontal half only. With
%     inner sep = 1pt every left-anchored text box stood 1.19pt (0.42mm) inside
%     the axis's left end, so both rules and both left end ticks protruded to
%     the left of the entire text column and S5's "left-aligned to the axis's
%     left end" was met to 0.42mm rather than exactly. inner ysep stays 1pt, so
%     every vertical clearance recorded above is still box arithmetic; inner
%     xsep is 0pt. Measured on this build, the ten panel-level text lines now
%     share ONE ink x0, 59.81pt, against an axis rule beginning at 59.615pt --
%     the residue is the first glyph's own side bearing and nothing else.
%     Band (b)'s three row labels are now placed with an xshift on the MARK's
%     own coordinate rather than at a literal x, so each is tied to the thing
%     it names: 5pt right of the disc in rows one and two, 5pt right of the
%     interval's upper end in row three, which is that row's rightmost ink.
%
% X13. THE AXIS LABELS ARE SET LOWERCASE. SUPERSEDES W16's closing paragraph and
%     repair_pass_9's not_applied entry, both of which recorded that the spec's
%     capitalised string "Preference shift, nats" was followed deliberately
%     because the slate is split on case. It is now "preference shift, nats" in
%     both bands. The reason for changing a spec-verbatim string, recorded so a
%     human can put it back in one edit if the slate is unified the other way:
%     the contract's own ruling is "between 4.11 and A.3: take 4.11", 4.11 sets
%     all four of its axis labels lowercase and 4.13 sets both of its, S19's
%     registry gives the quantity's name lowercase ("preference shift"), and
%     this figure's own panel heads are lowercase after S2 -- so capitalising
%     the axis label alone left one string in the figure with a case of its own.
%     Only the case changed; the quantity and the unit are the spec's words.
%
% X14. THE MAGNIFICATION IS NOW STATED TWICE, NOT THREE TIMES. SUPERSEDES
%     W16(ii), which recorded band (b)'s scope line as set verbatim from the
%     spec. It read "the first eighth of the axis above, magnified 8x" and sat
%     3.2cm below a panel head reading "(b) the first eighth, magnified 8x".
%     The head keeps the factor; the axis-label line keeps the relation to the
%     axis above, which is the fact the head does not carry and the fact an
%     axis-label scope line is for. The exactness of the eight is caption
%     material and the caption states it. Band (a)'s bracket label is the third
%     appearance and it stays: it is the only one that is geometric, it names
%     the bracket, and it is 4cm from the other two.
%
% X15. figs/_PENDING_CAPTIONS.tex NOW EXISTS. SUPERSEDES W9's assertion, and
%     the four assertions in materiality.json, that the sentences cut from the
%     drawing had been written into it: at pass 9 the file did not exist, so
%     the two deleted sentences and the whole proposed caption survived only in
%     the commented block at the foot of this file, which no build reads and no
%     integrator was pointed at. The file is created at this pass, keyed by
%     \label{fig:materiality}, and carries the proposed caption in full, the
%     two sentences deleted under W9 verbatim, and a list of the clauses of the
%     PRINTED caption (Sections/04b-representation.tex:597-620) that the
%     rebuild made false. No builder on this pass may edit Sections/.
%
% X16. FINDINGS OVERRULED, WITH THE REASON, so they are not re-applied blind.
%     (a) "Drop a quietink dotted projection from each of band (b)'s three
%         marks to the spine, as band (a) does for its one mark." OVERRULED.
%         Two of the three would cross the interval bar: the permuted mark
%         stands at x = 1.118 and the swap mark at x = 3.318, and the bar runs
%         0.702 to 3.751 on the row beneath them, so the projections would
%         strike it at two interior points and read as endpoints of something.
%         Band (a)'s projection is not a convention this figure owes band (b);
%         it is a rescue for a single disc standing 2.75mm from the origin on a
%         139.96mm rule with no tick within 4cm. In (b) every mark is on a row,
%         direct-labelled at itself, 4.8 to 15.6mm above a rule whose ticks
%         stand every 4.36cm, and the two arm values the projections would help
%         a reader read off are now PRINTED, in the caption, which is where the
%         same reviewer asked for them.
%     (b) "Raise the bracket into the lane immediately under the tick numerals."
%         OVERRULED. S5 puts the axis-label block directly under the tick
%         numerals; the bracket carries a \scriptsize label of its own, so
%         moving it there interposes a second text block between the numerals
%         and the line that names the quantity, which is the supervisor's first
%         complaint reintroduced by the back door. It was raised as far as the
%         scope line's descenders allow instead (X8).
%     (c) "Reverse the bracket's stubs so they turn UP toward the rule."
%         OVERRULED. With the bracket necessarily BELOW the two-line axis-label
%         block, up-turned stubs bracket the scope line's text -- "residual
%         arm, layer 12, alpha = 0.5" runs 0 to 3.66cm and the bracket spans 0
%         to 1.749cm, so they would enclose its first half. Down-turned stubs
%         agree with the label's own pointer, "the eighth magnified in (b)",
%         and point at the panel that opens the span. This was pass 9's stated
%         intent and it survives the pass.
%     (d) "Drop \texttt on the arm names, or scale the mono, because Inconsolata
%         sets at 6.85pt inside 7.97pt roman." OVERRULED. 6.85pt is the
%         preamble's own scaling of the mono face to Pagella, not a second type
%         size chosen here, and S1 exempts the \texttt token name from the size
%         count for exactly this reason. The tokens are JSON key names --
%         swap, permuted, residual -- and setting them roman would turn
%         "swap, points at the drug scored" into a sentence about an English
%         verb. The register is the corpus's (ladder.tex, comparators.tex).
%     (e) "The threshold carries two label strings, '10% materiality margin,
%         0.00642 nats' in (a) and '10% materiality margin' in (b), which fails
%         S19's mechanical check." OVERRULED on the ledger's own standing
%         principle, applied to diff_perm at W6: the NAME is "10% materiality
%         margin" in both bands, and (a) additionally prints the VALUE. A value
%         is not a second name. It is printed in (a) and not in (b) because (a)
%         is the band that also prints the quantity it is a tenth of, so the
%         tenth is a division a reader performs on two printed numbers 8cm
%         apart -- which is what W5 traded the deleted brace for.
%     (f) "Restore the spec's clause 'the full width of this rule is the
%         model's clean preference gap' to band (a)'s scope line." OVERRULED
%         again, for W16(i)'s reason and one more: with the terminal label now
%         standing on its own end cap (X6 ii) and both ends of the span capped
%         in ink (X7), the drawing says geometrically what that clause says in
%         twelve words, and the panel head says it in five.
%
% X17. MEASURED ON THE PASS-10 BUILD, not estimated. figs/_preview.py --width
%     text materiality reports ok, pages = 1, Overfull 0, Underfull 0. Content
%     bounding box 139.96 x 84.12mm (S14's 138-142mm; the 85mm height cap).
%     Ink 5.07% of pixels below 230 grey at 200dpi inside that box (ceiling
%     6.0%). 572 characters of figure text, 4.86 per 100mm^2 (ceiling 5.0).
%     Rule weights, counted off the PDF: 0.30pt x10 (nine rulegrey furniture
%     strokes and one accent leader), 0.40pt x4 (the two spines and the two end
%     caps), 0.50pt densely dotted x2, 0.90pt x2, 1.00pt x1 -- five weights,
%     each with more than one instance or declared as a distinct kind of object
%     at W11 -- plus the 0.80pt open-disc edge, which is a mark edge and not a
%     rule. Type: 7.97pt (400 characters) and 9.96pt (105) in TeXGyrePagella,
%     which are \scriptsize and \footnotesize; 6.85pt (58) in Inconsolata,
%     which S1 exempts; and 8.30pt (8) and 10.38pt (1) in the math companions,
%     which are the SAME two nominal sizes set from the math families -- the
%     minus in "swap - permuted" and the \times in "magnified 8x". Italic: the
%     render carries exactly one span in an italic-named face, PazoMath-Italic
%     at 7.97pt, and it is the mathematical $\alpha$ of "alpha = 0.5" -- 2 of
%     572 characters, 0.35%, which is a symbol's own face and not an italic
%     voice, the use S16 allows. There is no \itshape anywhere in the source.
%     PASS 9's W12 and the ledger's S16 line both said "0%"; the figure of
%     record is 0.35%, and it is the alpha. Sans: 0%. Bold: 0%. Rectangles: 0.
%     Arrowheads: 0.
% ===========================================================================
% ===========================================================================
% ===========================================================================
% PASS 11 -- 2026-08-19. REPAIR PASS on the pass-10 repair. Three adversarial
% readers went over the built figure independently and returned twenty-three
% findings. NOTHING ABOVE IS DELETED: passes 1-10 are the audit trail and stay
% as written. Every clause of passes 1-10 that this pass made false, and every
% clause those passes left false, is superseded here BY NAME (S25). The
% findings that were OVERRULED are recorded under Y15 with the reason, so that
% a later builder does not re-apply them blind.
%
% NO NUMERIC VALUE CHANGED IN THIS PASS EITHER. Every quantity drawn or
% printed is the same float from the same key path as at passes 8, 9 and 10:
% clean_ab_gap 0.06420005919957501, equiv_margin 0.006420005919957502,
% effect_perm 0.0006409306049362111, effect_swap 0.0019027651015006252,
% diff_perm 0.0012618344965644142, ci_perm [0.0004024989353392094,
% 0.002150955667015771], n_pair_clusters 58, and the exact magnification 8.
% Both maps are untouched: \ga = (#1)*218.0000 and \gb = (#1)*1744.0000, with
% 1744.0000 = 8 x 218.0000 exactly (rechecked on this build: 8*218.0 == 1744.0
% is True in IEEE double, and 0.06420005919957501*218.0 ==
% 0.008025007399946877*1744.0 is True, both ends coinciding at
% 13.995612905507352cm). ONE PRINTED STRING changed its number of decimal
% places without changing its value -- band (b)'s middle tick numeral, 0.005
% -> 0.0050 (Y11) -- and one value that was printed on a row of (b) is now
% printed in (a) instead (Y2, Y8). Nothing was rounded differently, and nothing
% new was computed.
%
% WHAT MOVED, IN ONE LINE EACH. The changes are commented at the point of edit
% as Y1..Y13 and summarised here.
%   Y1  band (a)'s panel head names the object by the registry string.
%   Y2  band (a)'s one label prints the estimate and drops 0.06cm.
%   Y3  the axis's right end is computed by \ga, not transcribed.
%   Y4  band (a)'s scope line drops 0.12cm and says "layer 12 of seven".
%   Y5  the magnifier bracket's stubs are visible and it stands clear.
%   Y6  band (b)'s zero rule is named "zero".
%   Y7  the mark-to-label gap in all three rows of (b) goes 5pt -> 9pt.
%   Y8  the estimate's value comes off (b)'s row-three baseline label.
%   Y9  0.06cm more air under row two.
%   Y10 band (b)'s rule is capped at its right end, in (a)'s cap.
%   Y11 band (b)'s middle tick numeral prints 0.0050.
%   Y12 the foot block's leading and its stand-off both open.
%   Y13 the pinned bounding box floor goes 1.72 -> 1.68.
%
% Y14. CLAUSES OF PASSES 1-10 THAT WERE ALREADY FALSE OF THE BUILT FIGURE AND
%     WERE SUPERSEDED NOWHERE. W15 and X1 each did this job for the pass they
%     followed; neither went back over the pass-1 to pass-8 LAYOUT NOTES and
%     DESIGN DECISIONS, so an auditor reading this header top to bottom was
%     handed five descriptions of a figure that does not exist -- the exact
%     hazard X1 was written to prevent. Each is named and given its drawn value.
%     (i) DESIGN DECISION 5, "THE MEASUREMENT SITS 0.22cm ABOVE BAND ONE'S
%         SPINE ... with its dot over a 2.1pt white halo ... the bar, the dot
%         and the spine are three separable things."
%         -> The disc is drawn at y = 8.90 against a spine at 8.54, which is
%            0.36cm and not 0.22cm. There is NO white halo under it: the only
%            \fill[white] in the file is band (b) row one's open control disc.
%            And there is no bar in (a) at all, since W6 drew the interval once
%            and gave (a) a single filled disc. The decision's REASON -- that a
%            mark on the spine reads as a thickening of the axis, so it is
%            lifted clear -- survives; its three measurements do not.
%     (ii) LAYOUT NOTE "TICK NUMERALS CARRY fill=white AND ARE DRAWN LAST IN
%         THEIR BAND."
%         -> num/.style and numL/.style carry no fill. The white box existed to
%            interrupt the magnifier's left leader where it crossed band one's
%            "0"; W10 deleted that leader, so nothing crosses the numerals and
%            drawing order no longer matters.
%     (iii) LAYOUT NOTE "The band-two tag is anchored at x = 0.75 and not at
%         x = 0.00, because the magnifier's left leader is a vertical at
%         x = 0.55."
%         -> Neither object exists: the tag went at W7 (replaced by 4.11's
%            panel head) and the leader at W10.
%     (iv) LAYOUT NOTE "BAND ONE'S RIGHT END is a terminal tick RISING from the
%         spine, 0.30cm tall against the 0.10cm hanging axis ticks, with its
%         two label lines right-aligned on the same x = 13.85."
%         -> It rises 0.22cm, from 8.54 to 8.76, at x =
%            \ga{0.06420005919957501} = 13.995612905507352, and it carries ONE
%            label line, not two. Band (b)'s right end now carries the same cap
%            at the same x (Y10).
%     (v) LAYOUT NOTE "The bounding box is pinned with a bare \path so the
%         measure is exact and nothing is clipped: (0.0000,0.30) to
%         (14.1900,8.98) = 141.9 x 86.8mm."
%         -> It is pinned (0.00,1.68) to (13.9957,10.18) after Y13, and the
%            measured content box is at Y16. The 141.9 x 86.8mm figure is a
%            pass-2 measurement and is over both the 142mm measure and the 85mm
%            height cap.
%     (vi) W11's "The 0.50pt densely dotted stroke is used once, for band two's
%         zero rule, and is a distinct kind of object."
%         -> It has two instances in this build and had two in pass 10, as
%            X17's own weight count records: band (a)'s projection from the disc
%            to the spine, and band (b)'s zero rule. It is not a singleton and
%            needs no singleton declaration under S10. The 1.00pt interval bar
%            IS a singleton and its declaration stands.
%
% Y15. FINDINGS OVERRULED, WITH THE REASON, so they are not re-applied blind.
%     (a) "Take the open control disc's edge to ink rather than quietink, as
%         fig 4.11(b) does for its matched-dimension random slab." OVERRULED.
%         S8 gives quietink to a comparator, a null or a control and S9 says
%         the open disc's 0.8pt edge is drawn "in its colour"; ink is the rank
%         S8 reserves for thresholds, axes, axis labels, titles and definitional
%         prose, so an ink ring would put the control in the apparatus rank and
%         make the comparator heavier than the accent mark under test. 4.11(b)
%         can draw its null black because its palette has no quietink rank at
%         all; this figure has four ranks and the control's is quietink. The
%         defect the finding reports is real and blocking and is fixed by the
%         gap alone (Y7): 7.2pt of white against the mono face's own 3.43pt
%         character advance, where 3.28pt was letter-spacing.
%         The companion finding, "quietink carries two ranks -- the one
%         measured comparator's label and all four apparatus lines", is
%         OVERRULED for the same reason from the other side: S5 mandates
%         quietink for both scope lines and S17 for the foot block, so the
%         collision cannot be broken by recolouring either party. It is broken
%         by position, which is how this figure has always broken it: the
%         control's label is anchored on a mark, on a row, inside the data
%         field; every apparatus line is panel-level text on the figure's one
%         x origin, under an axis.
%     (b) "Drop band (b)'s two arm dots. Their values are quoted in no sentence
%         of the thesis and the proposed caption already prints them, and the
%         one interval bar horizontally contains both of them." OVERRULED.
%         The spec's keep-list makes the three-row structure -- permuted, swap,
%         and their difference -- load-bearing, and CF14 fails a rebuild that
%         loses a load-bearing piece even if it passes every typographic check.
%         Without the two arms the difference is a bare number with a bar and
%         the figure stops drawing the subtraction that is its subject; design
%         decisions 7 and 14 record that reasoning and it stands.
%         THE DIAGNOSIS IS ACCEPTED AND IS REAL, and it is answered at Y9
%         rather than by deletion. The arithmetic is an accident worth stating
%         plainly: the bar runs \gb{0.0004024989353392094} = 0.702cm to
%         \gb{0.002150955667015771} = 3.751cm, the control stands at 1.118 and
%         the arm at 3.318, so the interval on the DIFFERENCE brackets both of
%         its own operands, and \gb{diff_perm} = 2.201 falls within 0.02cm of
%         their midpoint. What is done about it: the row pitch is no longer
%         flat (0.58 / 0.64 rather than 0.54 / 0.54), so the derived row sits
%         off the pair it is derived from; the row's name states the derivation
%         at the bar's own end; and rows one and two name their arms
%         individually immediately above, so "swap $-$ permuted" is read after
%         "permuted" and "swap" and not before them.
%         The finding's own fallback -- "move rows one and two onto their own
%         short strip with their own zero" -- is also OVERRULED: a second rule
%         for one quantity at one scale inside one panel is a second axis a
%         reader has to reconcile with the first, against S13's one grid, and
%         it would put two zeros 5mm apart on a figure whose subject is a
%         distance from zero.
%     (c) "Move the magnifier bracket into band (a)'s tick lane, rail at the
%         hanging ticks' depth with its stubs turned UP to the spine; or, if
%         that lane cannot be had, delete the bracket and tint the magnified
%         span in (a) and the whole of (b)'s data field instead." OVERRULED,
%         both halves. X16 (b) and (c) already overruled the first and the
%         reason holds: a bracket in the tick lane needs its label in the
%         tick-numeral lane, which is the defect S24 indicts in fig:purify,
%         where the word "four" sets on the x-tick baseline and reads as a
%         leftmost tick. The tint is worse, not better: S12 permits a filled
%         rectangle only where the rectangle IS an object in the argument and
%         forbids outright a rectangle drawn round a panel, and a tint over the
%         whole of (b)'s data field is that rectangle by another means -- the
%         object W4 deleted. The finding's real complaint, that the assembly did
%         not read as a bracket at all, is accepted and fixed at Y5.
%     (d) "Raise the bracket's rail onto its label's own baseline, so the rail
%         and the label read as one object rather than the rail reading as a
%         strike-through beside the words." OVERRULED as already true. The rail
%         at 7.02 stands 0.06cm above a label baseline at 6.96, which is that
%         label's own x-height centre, so rail and label are optically on one
%         line. What made it read as a rule was the 0.51mm stubs and the scope
%         line overhanging both its ends; both are fixed at Y5.
%     (e) "Label band (b)'s right cap 'one eighth of it, 0.008025 nats', so
%         that the exact eight becomes a division a reader can perform on two
%         printed numbers." OVERRULED on geometry, not on principle. Band (b)'s
%         threshold rule stands at \gb{0.006420005919957502} = 11.196cm and
%         runs the full height of the data field; a label right-aligned on the
%         axis's end at 13.996cm measures 3.4 to 4.7cm at \scriptsize and
%         crosses it, and the only string short enough to clear the rule is a
%         bare numeral with no name, which S20 forbids. There is no lane above
%         the rule either: (b)'s threshold label already occupies it, ending at
%         11.10cm. The magnification is stated by (b)'s panel head, by (b)'s
%         axis-label scope line and by (a)'s bracket label, and its EXACTNESS
%         is caption material (X14, and the caption states it). The cap added
%         at Y10 is axis furniture -- the rule's own end -- and not a mark, so
%         it needs no name under S20; its end is named by (b)'s axis label,
%         which is W13's reasoning and survives.
%     (f) "Restore the spec's clause 'the full width of this rule is the
%         model's clean preference gap' to band (a)'s scope line: line one
%         reads 'preference shift, nats', but band (a)'s full width is
%         clean_ab_gap, the model's contrast score with NO intervention -- a
%         level, not a shift -- so a reader who takes the axis label at face
%         value reads the right terminus as a preference shift of 0.0642 nats."
%         OVERRULED a third time. The REASON RECORDED AT W16(i) AND X16(f) IS
%         WITHDRAWN AND REPLACED, because it was wrong: both cited S18, and
%         S18 forbids a drawing from restating a clause of ITS OWN CAPTION or
%         of the paragraph that introduces the float. It says nothing about a
%         drawing carrying two related statements of its own. The clause is
%         refused on S5 and S15 instead. S5 caps the axis-label block at two
%         lines and gives line two to scope or a scale note; a thirty-character
%         clause about what the axis's other end IS is neither. And the fact is
%         already stated at the two places a reader looks for it: the panel head
%         reads "(a) the whole clean preference gap" after Y1, and the terminal
%         label at the axis's own right end reads "the clean preference gap,
%         0.0642 nats". THE FINDING'S UNDERLYING OBSERVATION IS RIGHT AND IS
%         RECORDED HERE so that no later reader has to rediscover it: this
%         axis's quantity name is the quantity its MARKS carry, and its right
%         terminus is a level of that same contrast score measured with no
%         intervention, named as such at the terminus. Laying the two on one
%         rule is the figure's whole construction and the thesis's own -- 04b:
%         663 divides every effect by "that model's own clean preference gap,
%         measured with no intervention".
%     (g) "Scale the mono up so that \texttt sets at the roman's x-height:
%         6.85 x (3.96/3.24) = 8.37pt." OVERRULED again, as at X16(d), and for
%         a reason that finding does not meet. 6.85pt is preamble.tex's global
%         scaling of Inconsolata to Pagella and it applies to every \texttt in
%         the thesis. Overriding it inside one figure would make this figure's
%         JSON key names set differently from comparators.tex's, ladder.tex's
%         and the body text's, which is a worse inconsistency than an x-height
%         step inside a label; and S1 exempts the \texttt token name from the
%         size count precisely so that the document's global scaling can stand
%         without a figure having to fight it.
%     (h) "In band (a) the two black labels that share the '<name>, <value>
%         nats' construction are the CLOSER pair (0.30cm) while the accent
%         label, which differs from both in face, colour and construction, is
%         the further (0.40cm); drop the terminal label 8.86 -> 8.80 and lift
%         the threshold label 9.16 -> 9.22." OVERRULED as proposed, accepted in
%         substance. The two black labels are 34mm apart horizontally, sit on
%         different baselines, and one is right-aligned onto the axis's own end
%         cap while the other is left-set beside a rule 8cm away; X6(ii)
%         separated them for the stronger reason that they had SHARED a
%         baseline, and the proposed 8.80 would set a 7.97pt baseline 1.1pt
%         above the 0.22cm cap it labels, re-creating the collision X6(ii)
%         removed. After Y2 the stack reads 0.40 / 0.34 / 0.30 from the head
%         down -- monotone, so nothing alternates and no pair is singled out by
%         its spacing.
%
% Y16. MEASURED ON THE PASS-11 BUILD, not estimated, and SUPERSEDING X17's
%     bounding box by name. figs/_preview.py --width text materiality reports
%     ok, pages = 1, Overfull 0, Underfull 0.
%       content bounding box   140.34 x 84.33mm   (S14's 138-142mm; 85mm cap)
%       ink                    5.36% of pixels below 230 grey at 200dpi (6.0%)
%       figure text            587 characters, 4.96 per 100mm^2   (ceiling 5.0)
%     X17 recorded 139.96 x 84.12mm for the pass-10 build; re-measured with the
%     contract's own threshold (pixels below 230 grey at 200dpi, clipped to the
%     content box) that build is 140.21 x 83.82mm. The +0.25mm and -0.30mm are
%     a method difference and not a drawing change, and they are recorded here
%     rather than left as an unexplained discrepancy: X17's number came from a
%     box measured before the antialiased edge columns of the axis caps were
%     counted. Both builds are inside every limit; this one is 0.67mm under the
%     height cap and 0.04 characters per 100mm^2 under the density ceiling,
%     which is why six characters came out of the foot block when "of seven",
%     "clean" and "zero" went in.
%     Rule weights, counted off the PDF: 0.30pt x10 (nine rulegrey furniture
%     strokes -- eight ticks and the bracket polyline -- plus one accent
%     leader), 0.40pt x5 (two spines and three end caps), 0.50pt densely dotted
%     x2, 0.90pt x2, 1.00pt x1. Five weights, no singleton except the interval
%     bar, which W11 declares as a distinct kind of object. Plus the 0.80pt
%     open-disc edge, which is a mark edge and not a rule.
%     Type: 7.97pt (417 characters) and 9.96pt (111) in TeXGyrePagella, which
%     are \scriptsize and \footnotesize; 6.85pt (58) in Inconsolata, which S1
%     exempts; 8.30pt (8) and 10.38pt (1) in the math companions, which are the
%     same two nominal sizes set from the math families. Italic: one span,
%     PazoMath-Italic at 7.97pt, the mathematical $\alpha$ -- 2 of 587
%     characters, 0.34%, a symbol's own face and not an italic voice (S16).
%     Sans 0%. Bold 0%. Rectangles 0. Arrowheads 0.
%
% Y17. THE CAPTION. figs/_PENDING_CAPTIONS.tex is updated in the same pass and
%     three of its clauses changed with the drawing: band (a)'s mark now prints
%     its value, band (b)'s row three no longer does, and (b)'s rule is capped
%     at both -- at its right -- end. ONE ITEM IN THAT FILE IS BLOCKING and is
%     flagged there as such: the printed caption (04b:617-618) says only "this
%     is one arm at one depth at one rung", which is a scope statement, while
%     layer 12 is the argmax of the anchor seed's seven-depth sweep and the
%     only one of the seven whose interval excludes zero. The drawing now says
%     "layer 12 of seven" (Y4), which discloses that the depth was chosen from
%     a set but not that it was the favourable one; only the caption can carry
%     that, and no builder on this pass may edit Sections/.
% ===========================================================================
%
% ===========================================================================
% v-BLOCK, 2026-08-19: WHOLE-SLATE PASS (fig:comparators, fig:licensing,
% fig:purify, fig:hook, fig:materiality reconciled against one another and
% against fig:mechanism, which is the document's reference standard and was
% NOT rebuilt). Nothing above this line is deleted; where an earlier claim is
% now false it is superseded EXPLICITLY below.
%
% WHY. The five figures were rebuilt independently and drifted. This pass owns
% only the differences BETWEEN them. No value changed in any of the five; the
% sibling .json ledgers carry the value-to-position maps, old and new.
%
% THE SLATE RULES THAT NOW HOLD ACROSS ALL FIVE (each stated once here, in
% every file, so any future single-figure edit can see what it would break):
%
%  R1 TICK NUMERALS are quietink, \scriptsize, inner sep 0, anchored north on
%     the tick and set 0.16cm below the axis rule, over a 0.10cm rulegrey
%     0.30pt tick. Before: comparators/licensing/purify quietink, hook and
%     materiality ink; ticks 0.10 / 0.12 / 0.14cm; three different gaps.
%  R2 A THRESHOLD'S OWN LABEL IS BLACK, matching the black solid 0.90pt rule
%     it names (S6). Before: comparators black, purify's "gate, 0.8" and
%     materiality's two margin labels ink.
%  R3 ONE NAME PER OBJECT (S19). The pre-registered 10% rule is "10%
%     materiality margin" in fig:comparators row (c), fig:materiality (a) and
%     (b), and in figs/family.py, which is where the name comes from.
%     fig:comparators' "pre-registered margin" is superseded; "pre-registered"
%     is now the caption's word, not the drawing's. Zero, where it is the null
%     a bar is read against, is "zero: no identity effect" in BOTH
%     fig:comparators row (c) and fig:materiality (b) -- one quantity, one
%     wording. V_pure is spelt with \textrm everywhere, as figs/slab.tex does.
%  R4 COLOUR KEY, one key for five figures (S8). accent = the drug-side object
%     under test. quietink = every comparator, null, control, replication or
%     discounted arm, AND its label, AND its leader. ink/black = thresholds,
%     axis rules, axis labels, panel heads, definitional prose. rulegrey =
%     furniture only. Within quietink the GLYPH separates two kinds: an OPEN
%     disc is a construction built to carry no signal (a null, a control, a
%     raw/before series); a FILLED disc is a real measurement that is not the
%     headline. fig:hook (d) drew both its null and its cell-line control in
%     ink and is corrected.
%  R5 PROJECTION AND ZERO RULES are quietink 0.50pt densely dotted (S7).
%     fig:hook's two panel-(c) projection lines were rulegrey 0.30pt densely
%     dashed and are corrected.
%  R6 DEFINITIONAL PROSE is \footnotesize ink and OUTRANKS every label (S4);
%     at most one block per figure, on the panel's own x origin. fig:hook's
%     one prose line was \scriptsize and is raised. fig:purify's panel-(b)
%     block carried an argument clause ("so causal claims use the purified
%     slab, not the raw") that its caption already carries, and is now
%     definitional only.
%  R7 ONE INTERVAL DECLARATION per figure, \scriptsize quietink, at the foot,
%     on the figure's x origin, opening with the word "Intervals:" and naming
%     EVERY panel including the ones with no measured mark (S17).
%     fig:comparators already did this and is the model; the other four now
%     match its form.
%  R8 THE FIVE-FIELD PROMPT STRIP is ONE object drawn in two figures, and the
%     two drawings are now congruent: cell boundaries 0 / 1.12 / 1.90 / 2.64 /
%     4.23 / 6.90 cm and cell height 0.42cm in both fig:licensing (a) and
%     fig:hook (a) and (b), labels centred in their cells, the drug cell
%     accent-bordered at 0.40pt on accent!12 and its label NOT bold. Before:
%     licensing 6.90cm wide with 0.62cm cells and a bold label, hook 7.44cm
%     wide with 0.42cm cells and a plain label -- the same five fields at two
%     sizes, four pages apart, with each caption pointing at the other.
%  R9 MEASURE. Every fullwidth figure draws 172.0-174.0mm; the one textwidth
%     figure draws 140.3mm. Ink at 200dpi inside the content bbox, counted as
%     greyscale luminance below 230 (the contract's own method), is now
%     5.63-5.99% on all five, under the 6.0% body ceiling, where before this
%     pass it ran 5.36-6.54%.
%
% WHAT THIS PASS DELIBERATELY DID NOT UNIFY, so that a later reader does not
% think it was missed:
%  * PANEL-HEAD CASE. fig:comparators sets "(a) Is it there?" and the other
%    four set "(a) where the drug name enters". comparators' heads are
%    interrogative SENTENCES and sentence case is correct for a sentence; the
%    v4 header argued this and the argument stands. fig:mechanism's lowercase
%    heads are noun phrases, not questions. One capital letter, four times.
%  * PANEL-HEAD POSITION. fig:comparators puts its heads in a right-aligned
%    left stub, the other four above the panel at its x origin. Moving them
%    above the rules costs about 20mm of height against a 118mm cap that the
%    figure already meets at 117.60mm. Not available.
%  * TWO-COLUMN GUTTER. purify and hook split at x = 8.40cm, licensing at
%    7.75cm. licensing cannot move right: "(d) what a negative would buy" is
%    4.85cm wide and its origin is already at 12.40 of a 17.30cm measure, with
%    1.5pt of slack. Moving purify and hook LEFT to 7.75 would be arbitrary.
% ===========================================================================
%
% CHANGED IN THIS FILE. (i) No math sign is set from a CM font any more. The
% panel-(b) head's 8x (CMSY10 at 10.38pt, the largest glyph on the figure and
% a THIRD type size above 7pt) is now 8\texttimes{}; the two minus signs of
% "swap - permuted" are \textminus{}; +0.0013, 10%, 0.00642, 0.0642, 12, 58
% and 95% lose their math mode, and $\alpha = 0.5$ becomes $\alpha$ = 0.5,
% which is fig:comparators' idiom. The figure now prints exactly 9.96 and
% 7.97pt in one face, plus the \texttt arm names. (ii) Tick numerals go ink ->
% quietink with inner sep 0, set 0.16cm below their rule (R1). (iii) Both
% "10% materiality margin" labels go black, matching the black solid 0.90pt
% rule they name (R2). (iv) Band (b)'s zero label "zero" -> "zero: no identity
% effect", which is fig:comparators row (c)'s wording for the same null under
% the same quantity (R3), and goes ink -> quietink so that the label matches
% the dotted quietink rule it names. (v) \texttt{residual} -> residual in the
% scope line: fig:comparators names the same arm "residual seeds at layer 12"
% in plain roman, and an arm name in prose is not a code token. \texttt{swap}
% and \texttt{permuted} stay, because they are the two stored conditions and
% tab:identity sets them so. (vi) The foot block takes the slate's declaration
% form (R7).
% NOT CHANGED, and still declared: the 1.00pt interval bar is a singleton
% weight because this figure stores exactly one interval (S10's permission);
% the 0.80pt open-disc edge is a MARK edge, not a rule.
% RE-MEASURED: 140.34 x 84.45mm (cap 85, textwidth band 138-142); ink 5.630%;
% 620 characters, 5.23 per 100mm^2. One page, zero Overfull, zero Underfull.
% ===========================================================================
\begin{tikzpicture}[x=1cm, y=1cm, font=\small, inner ysep=1pt, inner xsep=0pt,
  %% TYPE. Two sizes above 7pt and one face (S1, W12). \footnotesize = 10.0pt
  %% carries the two panel heads and the two axis-label line ones and nothing
  %% else; \scriptsize = 8.0pt carries ticks, direct labels, scope lines, the
  %% threshold labels and the interval declaration. No sans, no bold, no italic.
  %% INNER SEP, PASS 10 (X12). It was inner sep = 1pt on the whole picture; the
  %% vertical half is unchanged, so every clearance below is still box
  %% arithmetic, but the horizontal half is now 0pt. With xsep = 1pt every
  %% left-anchored text box stood 1.19pt (0.42mm) inside the axis's left end,
  %% so both rules protruded to the left of the entire text column and S5's
  %% "left-aligned to the axis's left end" was met to 0.42mm and not exactly.
  head/.style  ={font=\footnotesize, text=ink,      anchor=base west},
  axlab/.style ={font=\footnotesize, text=ink,      anchor=base west},
  scope/.style ={font=\scriptsize,   text=quietink, anchor=base west},
  qlabz/.style ={font=\scriptsize,   text=quietink, anchor=base west},
  lab/.style   ={font=\scriptsize,   text=ink,      anchor=base west},
  labR/.style  ={font=\scriptsize,   text=ink,      anchor=base east},
  acc/.style   ={font=\scriptsize,   text=accent,   anchor=base west},
  rowq/.style  ={font=\scriptsize,   text=quietink, anchor=west},
  rowa/.style  ={font=\scriptsize,   text=accent,   anchor=west},
  %% Tick numerals are INK, not quietink: they are part of the axis, and S8
  %% gives ink to axes and axis labels while quietink means a comparator or a
  %% control. 4.11 and 4.13 both set their tick numerals dark for this reason.
  num/.style   ={font=\scriptsize,   text=quietink, anchor=north,
                 inner sep=0pt},
  numL/.style  ={font=\scriptsize,   text=quietink, anchor=north west,
                 inner sep=0pt},
  blab/.style  ={font=\scriptsize,   text=black,    anchor=base west},
  blabR/.style ={font=\scriptsize,   text=black,    anchor=base east},
  %% FIVE STROKE WEIGHTS, each with one meaning (S10, W11).
  tick/.style  ={rulegrey, line width=0.3pt},                  % furniture
  lead/.style  ={rulegrey, line width=0.3pt},                  % furniture
  accled/.style={accent,   line width=0.3pt},                  % furniture, in the mark's colour (S22)
  spine/.style ={ink,      line width=0.4pt},                  % axis rule and its two end caps
  proj/.style  ={quietink, line width=0.5pt, densely dotted},  % zero rule / projection (S7)
  thresh/.style={black,    line width=0.9pt},                  % the one threshold (S6, W3)
  bar/.style   ={accent,   line width=1.0pt},                  % interval
  ctrl/.style  ={quietink, line width=0.8pt}]                  % open-disc edge: a MARK edge, not a rule (S9, W11)

  %% ==== THE TWO MAPS ====================================================
  %% Superseded constants, kept in materiality.json: 0.55+(#1)*207.1650 and
  %% 0.55+(#1)*1657.3200. The 0.55cm inset is gone so that the axis's left end
  %% and the panel's x origin are one coordinate (S13, W1).
  %% 1744.0000 = 8 x 218.0000 EXACTLY, so the magnification is exactly eight,
  %% and \ga{clean_ab_gap} = \gb{clean_ab_gap/8} = 13.995613cm exactly.
  %% PASS 10, X1: the pass-9 header block records these two constants as
  %% 215.7500 and 1726.0000 and the shared right end as 13.851163cm. Those
  %% numbers were never drawn by this file. The two lines below are the map,
  %% and materiality.json has always recorded them correctly.
  \def\ga#1{{(#1)*218.0000}}   % band (a): 0 .. 0.06420005919957501 nats
  \def\gb#1{{(#1)*1744.0000}}  % band (b): 0 .. 0.008025007399946877 nats

  %% ==================== BAND (a) -- THE WHOLE GAP ========================
  %% PASS 11, Y1: the head names the object by the registry string. It read
  %% "(a) the whole preference gap" while the terminal label 8cm to its right
  %% read "the clean preference gap", so one JSON key (clean_ab_gap) carried
  %% two names inside one panel, which is what S19's mechanical check forbids.
  %% "whole" is kept because it is the contrast with (b)'s "first eighth".
  \node[head] at (0.00,9.90) {(a) the whole clean preference gap};

  % The estimate, drawn once and without a bar (W6). At true scale it stands
  % 2.75mm from the origin on a 139.96mm rule, which is the panel's whole point.
  % Its name is band (b) row three's name, because one JSON key takes one name.
  % PASS 10, X4: THE LABEL IS ANCHORED AT THE MARK, not at the panel origin.
  % Pass 9 set it base west on x = 0.00 while its leader fell from
  % \ga{diff_perm} = 0.2751, so the leader descended from inside the word
  % "swap" and the mark read as the SWAP arm -- whose value is 0.0019, not the
  % 0.0013 drawn. That is the one mark in band (a); mis-naming it loses the
  % panel. The label now begins on the mark's own x, and the leader falls from
  % its first character.
  % PASS 11, Y2: THE LABEL PRINTS THE VALUE, AND IT DROPPED 0.06cm.
  % Two defects, one edit. (i) At 9.56 the label stood 0.34cm under the panel
  % head, near-left-aligned with it (7.8pt of indent), so head and label read
  % as a two-line title and the panel's single mark lost its name; the head's
  % descenders cleared the label's ascenders by 2.6pt. At 9.50 that clearance
  % is 4.3pt and the gap to the threshold label below is 0.34cm = 9.7pt, which
  % is still wider than the 7.94pt X6(i) condemned. (ii) Without a value the
  % math minus between two mono tokens parses as easily as a gloss dash --
  % "swap, that is, permuted" -- so the one mark in (a) could be read as the
  % SWAP arm, whose value is 0.0019 and not the 0.0013 drawn. The value now
  % settles it, and it also makes (a) self-sufficient for the figure's
  % question: 0.0013 against the 0.0642 printed at the axis's right end, and
  % against the 0.00642 printed at the threshold.
  % WHY A VALUE MAY BE PRINTED HERE AND NOT ON A ROW OF (b). This label sits
  % 0.96cm above the spine and is leadered to its disc; positions in that lane
  % are not read as values, which is why (a)'s threshold label already prints
  % 0.00642 three centimetres right of the rule it names and its terminal
  % label prints 0.0642 four centimetres left of the cap it names. A label set
  % on a data row's OWN baseline is read positionally, which is why the value
  % came off (b)'s row three (Y8).
  \node[acc] at (\ga{0.0012618344965644142},9.50)
    {\texttt{swap} \textminus{} \texttt{permuted}, +0.0013 nats};
  % The label's leader is accent 0.30pt, the mark's own colour (S22). PASS 10,
  % X5: it now LANDS ON THE DISC. It used to stop at 9.12 against a disc whose
  % top edge is 8.9493, leaving 4.84pt of white, so the leader read as a stray
  % accent tick under the label and the label read as unattached.
  \draw[accled] (\ga{0.0012618344965644142},9.40)
             -- (\ga{0.0012618344965644142},8.95);
  % The stem down onto the spine anchors the mark to a coordinate rather than
  % leaving it floating above the rule; pass 8 added it after both reviewers
  % reported they could not tell whether the mark was a mark, a tick or a
  % printing artefact, and it survives the rebuild. It is drawn in the PROJECTION
  % register -- quietink 0.5pt densely dotted, which is what S7 reserves that
  % stroke for -- and not in the leader's accent, because on the build that drew
  % leader and stem at one weight and one colour the assembly read as a VERTICAL
  % interval bar with a dot on it, which is the one thing this figure must not
  % appear to draw. PASS 10, X5: it now touches the disc's lower edge (8.85
  % against the disc's 8.8507) and the spine (8.54), instead of floating 2.0pt
  % clear of the one and 0.57pt clear of the other.
  \draw[proj] (\ga{0.0012618344965644142},8.85)
           -- (\ga{0.0012618344965644142},8.54);
  \fill[accent] (\ga{0.0012618344965644142},8.90) circle (1.4pt);

  % The one threshold, black and solid, in both bands (S6, CF2, W3). Its value
  % is printed here and the quantity it is a tenth of is printed at the axis's
  % right end, so the tenth is a division on two printed numbers and needs no
  % brace (W5).
  % PASS 10, X3: IT STANDS ON THE SPINE AND NO LONGER CROSSES IT. Pass 9 drew
  % it from 8.38, which is 0.06cm BELOW the hanging ticks' ends and 0.02cm above
  % the tick numerals, so the heaviest stroke in the panel doubled as an
  % outsized unnumbered tick sitting between the numerals "0" and "0.02" -- the
  % same "the threshold merges with the tick furniture" defect S7 indicts in
  % fig:comparators row (a). Band (b) already drew its half of the same
  % threshold from its own spine upward; the two now agree.
  \draw[thresh] (\ga{0.006420005919957502},8.54)
             -- (\ga{0.006420005919957502},9.04);
  \node[blab] at (1.50,9.16) {10\% materiality margin, 0.00642 nats};

  % The axis. PASS 10, X7: BOTH ENDS NOW CARRY THE SAME CAP. The right end is
  % the clean preference gap and has always carried a rising 0.22cm ink stub;
  % the left end carried only the same rulegrey hanging tick as 0.02, 0.04 and
  % 0.06, so the span whose full width IS the figure's denominator was
  % delimited at one end by an ink stub and at the other by an ordinary tick,
  % and did not read as a bounded span at all. Two matching caps, no new weight
  % and no new colour.
  % PASS 11, Y3: THE AXIS'S RIGHT END IS COMPUTED, NOT TRANSCRIBED. It was
  % hard-coded 13.995613 here, at the end cap and at the terminal label, while
  % W1 and the layout note both assert that "every data x is still computed
  % inside pgfmath from the full-precision float, so no value is rounded into
  % the source". \ga{0.06420005919957501} = 13.995612905507352, so the literal
  % was short by 3.5e-7cm -- invisible, but it is the ONE coordinate this
  % figure exists to draw, and the claim about it was false. Only the pinned
  % bounding box below still carries a rounded literal, and it carries no
  % value: it is a box corner, declared as such at Y13.
  \draw[spine] (0.00,8.54) -- (\ga{0.06420005919957501},8.54);
  \draw[spine] (0.00,8.54) -- (0.00,8.76);
  \draw[spine] (\ga{0.06420005919957501},8.54) -- (\ga{0.06420005919957501},8.76);
  % PASS 10, X6: the terminal label sits ON ITS OWN CAP, at y = 8.86, and no
  % longer shares the 9.16 baseline and the "<name>, <value> nats" construction
  % with the threshold label 8cm to its left. Read cold on the pass-9 build the
  % two were one row of two labels in one style, so "the clean preference gap,
  % 0.0642 nats" read as the label of a RULE and a reader looked for a rule at
  % 0.0642 and found a 0.22cm end cap. It is now 0.10cm above that cap.
  \node[labR] at (\ga{0.06420005919957501},8.86)
    {the clean preference gap, 0.0642 nats};
  \foreach \v/\t in {0.02/0.02, 0.04/0.04, 0.06/0.06}{
    \draw[tick] (\ga{\v},8.54) -- (\ga{\v},8.44);
    \node[num] at (\ga{\v},8.38) {\t};}
  % Zero is both the axis's left end and the panel's x origin, so its numeral is
  % set flush with that origin rather than centred on it: a centred "0" is the
  % only object that would break the figure's single left edge (S13).
  \draw[tick]  (0.00,8.54) -- (0.00,8.44);
  \node[numL] at (0.00,8.38) {0};

  % AXIS LABEL, in axis-label position: directly under the tick numerals, left
  % edge on the axis's left end, quantity and unit on line one (S5, W2).
  % PASS 10, X13: set lowercase, as 4.11 and 4.13 set every axis label of
  % theirs and as S19's registry gives the name.
  % PASS 11, Y4: THE SCOPE LINE DROPS 0.12cm AND DISCLOSES THE SELECTION.
  % (i) LEADING. It was set 0.28cm under a 9.96pt line, which is 7.94pt of
  % baseline-to-baseline for type 2pt larger than the lead -- negative leading,
  % and the tightest of the four stacked pairs in this figure. Four such pairs
  % were set at or under solid and it is the single thing that made the figure
  % look crowded beside 4.11. Both axis-label blocks and the foot block now
  % open to 0.38cm (10.8pt), and the whole of the added height was paid for out
  % of the white between the two bands, which fell 1.14cm -> 0.77cm.
  % (ii) SELECTION. "layer 12" alone is a scope statement. Layer 12 is the
  % argmax of the anchor seed's seven-depth sweep and the only one of the seven
  % whose interval excludes zero (diff_perm_norm with p_perm: +0.0022 p 0.583,
  % -0.0050 p 0.139, -0.0025 p 0.4495, +0.0041 p 0.3955, -0.0019 p 0.635,
  % +0.0197 p 0.001, +0.0059 p 0.4885 at layers 2/4/6/8/9/12/16). "of seven"
  % is nine characters and it tells a reader on the drawing that the depth was
  % one of a set; WHICH one, and why that matters, is caption material and the
  % proposed caption in figs/_PENDING_CAPTIONS.tex carries it. Until a human
  % integrates that caption the printed one (04b:617-618) says only "one arm at
  % one depth at one rung", which is scope and not selection.
  \node[axlab] at (0.00,7.78) {preference shift, nats};
  \node[scope] at (0.00,7.40)
    {residual arm, layer 12 of seven, $\alpha$ = 0.5};

  % The magnified span, named as an object and not traced by a leader (W10).
  % Its two stubs turn DOWN, towards the panel that opens the span, so the
  % bracket points at (b) instead of back at an axis it already stands under.
  % PASS 10, X8: raised 0.16cm, tight under the scope line, and its label now
  % NAMES the object as well as pointing at (b). Two reviewers asked for it to
  % be moved into the lane immediately under the tick numerals with its stubs
  % reversed; both halves of that were overruled, and the reason is recorded in
  % the pass-10 block under X8.
  % PASS 11, Y5: THE STUBS ARE LONG ENOUGH TO SEE, AND THE BRACKET STANDS
  % CLEAR OF THE AXIS-LABEL BLOCK. Measured on the pass-10 render, each stub
  % was 4 pixels at 200dpi -- 1.44pt, 0.51mm -- of 0.30pt rulegrey, which does
  % not survive print, so the assembly was a bare grey horizontal sitting
  % 0.34cm under a scope line that runs past both of its ends: it read as an
  % underline of "residual arm, layer 12," or as an editorial divider, and not
  % as a bracket. Two changes, no new weight and no new colour: the stubs go
  % 0.06cm -> 0.16cm, and the gap from the scope line goes 0.34cm -> 0.44cm,
  % which is larger than the 0.38cm internal leading of the axis-label block
  % above it, so the bracket stops reading as that block's third line.
  % Two harder proposals were overruled and the reasons are at Y15 (c) and (d).
  \draw[lead] (0.00,6.86) -- (0.00,7.02)
           -- (\ga{0.008025007399946877},7.02)
           -- (\ga{0.008025007399946877},6.86);
  \node[scope] at (1.85,6.96) {the eighth magnified in (b)};

  %% ============= BAND (b) -- THE FIRST EIGHTH, MAGNIFIED =================
  %% PASS 11: the whole of band (b) is raised 0.14cm and its axis lowered
  %% 0.10cm, which is where the leading opened at Y4 and Y11 was paid from.
  \node[head] at (0.00,5.84) {(b) the first eighth, magnified 8\texttimes{}};
  % The threshold label sits against the rule it names and takes the LEFT of it
  % here and the RIGHT of it in (a): in each band the label goes on the side
  % where the rule has room, which in (a) -- where the rule stands at a tenth of
  % the measure -- is the right, and in (b) -- where it stands at four fifths --
  % is the left. Set right of the rule here it measures 2.98cm from 11.30 and
  % crosses the right end of the axis at 13.996 (checked on a build, which
  % reported Overfull \hbox 2.84pt).
  % PASS 10, X6: it has come off the panel head's baseline (5.70) and down onto
  % its own, 0.12cm above the rule's top, in the same relation to its rule as
  % (a)'s label is to (a)'s. On the pass-9 build it sat at exactly the head's
  % baseline, so scanning the two panel heads read "(b) the first eighth,
  % magnified 8x ... 10% materiality margin" as one title line.
  \node[blabR] at (11.10,5.46) {10\% materiality margin};
  \draw[thresh] (\gb{0.006420005919957502},3.68)
             -- (\gb{0.006420005919957502},5.34);

  % Zero. Densely dotted quietink 0.5pt is this document's stroke for a
  % coordinate that decides nothing (S7).
  % PASS 11, Y6: IT IS NAMED, ON THE SAME BASELINE AS THE THRESHOLD LABEL AND
  % AT THE SAME 0.12cm ABOVE ITS OWN RULE'S TOP. SUPERSEDES W8's closing
  % clause, that the rule "carries no label, because its foot is the axis's own
  % 0 numeral". The equivalence criterion of 04b:507-513 is a CONJUNCTION --
  % the interval must lie inside the margin AND contain zero -- and this
  % figure's result is that it does the first and not the second. Pass 10 drew
  % the margin conjunct as the heaviest stroke on the page and the zero
  % conjunct as an unlabelled 0.5pt dotted rule standing on the panel's own x
  % origin, where it has no figure/ground of its own; at a 4x squint it
  % disappeared entirely. Naming it is what the figure's spec authorises ("if
  % the zero rule needs a name it takes 'zero' in \scriptsize roman ink at the
  % rule") and it is the whole of the fix: the rule is NOT promoted to the
  % threshold register, because S6 reserves black solid 0.90pt for a threshold
  % and S7 reserves this stroke for a coordinate that decides nothing -- which
  % zero, as a coordinate, is. What decides is whether the bar clears it, and
  % the bar clears it by 0.70cm.
  % The label is placed with an xshift on the rule's own x, as the row labels
  % are placed on their marks', so it is tied to the thing it names; the two
  % vertical rules of (b) now carry one \scriptsize label each, on one
  % baseline, one at each rule's top (S22).
  % PASS 10, X9: it no longer starts on the spine. Drawn 3.58 to 5.56 it was
  % collinear and continuous with the axis's left terminus, its hanging tick and
  % the panel's x origin, ran the full height of the data region, and matched
  % the threshold rule on the right for length -- so all three reviewers read
  % the pair as the two sides of a frame round the rows, which is the object
  % S12 exists to forbid. It now spans the three rows only, 0.28cm clear of the
  % spine and 0.42cm clear of the panel head, and reads as what it is: the
  % coordinate the interval bar's left end clears by 0.70cm.
  \draw[proj] (0.00,3.94) -- (0.00,5.34);
  \node[qlabz, xshift=3pt] at (0.00,5.46) {zero: no identity effect};

  % ROW ONE, the control: an open disc, white fill, 0.8pt quietink edge, which
  % is 4.11's own encoding for a comparator or a null (S9, CF5, W11).
  % PASS 10, X10: the gloss separator is a comma. It was an em dash, in the same
  % position after the token as row three's math minus, so one glyph class
  % carried "namely" in two rows and "subtract" in the third on a figure whose
  % whole subject is a subtraction. The minus in row three and in band (a) is
  % now the only dash inside a label anywhere on the drawing.
  % X12: the three row labels are placed with an xshift on the mark's own
  % coordinate rather than at a literal x, so each label is tied to its mark.
  % PASS 11, Y7: THE MARK-TO-LABEL GAP GOES 5pt -> 9pt IN ALL THREE ROWS.
  % SUPERSEDES X12's closing clause, which set it at 5pt. At 5pt the open
  % control disc stood 3.28pt from the first glyph of \texttt{permuted}, which
  % is LESS than Inconsolata's own character advance at this size (3.43pt),
  % while its outer diameter (3.6pt) is 11% larger than that face's x-height
  % (3.24pt measured at 200dpi) and its 0.80pt edge is of comparable weight to
  % the stems beside it. The one comparator mark on the drawing therefore set
  % as a lowercase letter o prefixing the word: at 12x it reads "o permuted".
  % Row two escaped only because its disc is filled and accent. At 9pt the gap
  % is 7.2pt, twice the mono advance, and the ring is a mark again.
  % A second fix was proposed for the same defect -- take the ring's edge to
  % ink, as 4.11(b) does -- and it is overruled at Y15 (a).
  \fill[white]  (\gb{0.0006409306049362111},5.24) circle (1.4pt);
  \draw[ctrl]   (\gb{0.0006409306049362111},5.24) circle (1.4pt);
  \node[rowq, xshift=9pt] at (\gb{0.0006409306049362111},5.24)
    {\texttt{permuted}, same construction, wrong drug};

  % ROW TWO, the arm under test.
  \fill[accent] (\gb{0.0019027651015006252},4.66) circle (1.4pt);
  \node[rowa, xshift=9pt] at (\gb{0.0019027651015006252},4.66)
    {\texttt{swap}, points at the drug scored};

  % ROW THREE, the measurement, and the only interval on the figure. The label
  % is anchored on the interval's upper end, which is the rightmost ink in the
  % row, so it clears the bar by the same 9pt the other two clear their discs.
  % PASS 11, Y8: THE VALUE COMES OFF THIS LABEL. It printed "$+0.0013$ nats"
  % 48.9pt (17.2mm) to the right of the disc whose value it is, ON THE ROW'S
  % OWN BASELINE, where position is value: the numeral stood at x = 0.00339 on
  % (b)'s axis, 2.7 times the number it printed and past the 0.0025 tick, on
  % the one panel whose job is to let a reader place the estimate between zero
  % and the 0.00642 margin. Rows one and two anchor at their marks and this
  % row anchored at its bar's end, so one panel carried two anchoring rules and
  % the one that misplaced its label was the one that printed a number. The
  % label is now the row's NAME only, which is the same string band (a) uses
  % for the same JSON key (S19), and the value is printed once, in (a), at a
  % label lifted off the data lane and leadered to its disc (Y2). After this
  % edit no numeral anywhere on the drawing stands on a data row's baseline.
  % PASS 11, Y9: 0.06cm MORE AIR UNDER ROW TWO. The bar on this row spans
  % 0.702 to 3.751cm and the two arm marks stand at 1.118 and 3.318, so the
  % interval on the DIFFERENCE horizontally contains both of its own operands
  % and the difference disc falls within 0.02cm of their midpoint by
  % arithmetic accident. A reader can be offered "a 95% interval that brackets
  % both arms", which is the opposite of what the row says. The row pitch is
  % now 0.58 / 0.64 rather than a flat 0.54, so the derived row sits off the
  % pair it is derived from; the label states the derivation; and the proposal
  % to delete the two arm marks outright is overruled at Y15 (b).
  \draw[bar] (\gb{0.0004024989353392094},4.02)
          -- (\gb{0.002150955667015771},4.02);
  \fill[accent] (\gb{0.0012618344965644142},4.02) circle (1.4pt);
  \node[rowa, xshift=9pt] at (\gb{0.002150955667015771},4.02)
    {\texttt{swap} \textminus{} \texttt{permuted}};

  % The magnified axis, its ticks landing on the same four x as band (a)'s so
  % that the two rules read as one rule at two scales (W13).
  % PASS 11, Y10: BAND (b) IS CAPPED AT ITS RIGHT END, AND ITS END IS COMPUTED.
  % SUPERSEDES W13's closing clause ("band two's right endpoint ... takes no
  % tick and no numeral"). W13 declined a terminal MARK on the ground that
  % (b)'s right end is named by (b)'s axis label rather than by a mark, and
  % that reasoning stands -- but what it left was 26.01pt (9.17mm) of bare
  % rule running past the last drawn tick to an end with nothing on it, which
  % is the same rule-end-to-last-tick defect S13 measures at 8.1mm in
  % fig:comparators row (c). Band (a) has the identical overhang and closes it
  % with a 0.22cm ink cap; band (b) now closes it with the same cap, at the
  % same x, in the same weight and colour, so the two rules terminate
  % identically at the end where the overhang is and read as one rule at two
  % scales. The cap is axis furniture -- the rule's own end -- and not a data
  % mark, so S20 is untouched; the proposal to hang a value on it is overruled
  % at Y15 (e). NO LEFT cap is added in (b): a 0.22cm ink stub rising from the
  % spine at x = 0 would stand 0.14cm under the dotted zero rule and collinear
  % with it, which is exactly the frame-edge reading X9 removed.
  \draw[spine] (0.00,3.68) -- (\gb{0.008025007399946877},3.68);
  \draw[spine] (\gb{0.008025007399946877},3.68)
            -- (\gb{0.008025007399946877},3.90);
  % PASS 11, Y11: the middle numeral prints 0.0050 and not 0.005. Band (a)'s
  % three numerals carry two decimals each and band (b)'s carry four, and the
  % one visible proof that (b) is (a) divided by eight is that the two lanes
  % put congruent numerals on the same four x. A dropped place broke the
  % column of trailing digits that carries it. The VALUE is unchanged; this is
  % a print string, which is why the \foreach takes the value/label pair form.
  \foreach \v/\t in {0.0025/0.0025, 0.005/0.0050, 0.0075/0.0075}{
    \draw[tick] (\gb{\v},3.68) -- (\gb{\v},3.58);
    \node[num] at (\gb{\v},3.52) {\t};}
  \draw[tick]  (0.00,3.68) -- (0.00,3.58);
  \node[numL] at (0.00,3.52) {0};

  \node[axlab] at (0.00,2.92) {preference shift, nats};
  % PASS 10, X14: "magnified 8x" comes off this line. The panel head 3.2cm
  % above already carries it, and the exactness of the eight is caption
  % material; what the axis-label line has to carry is the relation to the axis
  % above, which is what is left.
  \node[scope] at (0.00,2.54) {the first eighth of the axis above};

  %% ==== THE ONE INTERVAL DECLARATION, at the foot (S17, CF8) =============
  %% PASS 10, X2. The pass-9 second line read "No interval is returned for
  %% permuted, for swap, or for (a)." That was false, and false about the one
  %% mark on the figure that does have an interval: band (a)'s disc IS
  %% diff_perm, and layers.12.swap.ci_perm is an interval on diff_perm. In this
  %% document "returned" is what the instrument does, so the sentence told a
  %% reader that band (a)'s estimate has no uncertainty available and, by
  %% implication, that (a) and (b) draw different quantities. The declaration
  %% now separates non-existence from non-drawing: the arms have no stored
  %% interval, the difference has one, and it is drawn once.
  %% Leading is 0.28cm here, as everywhere else in this figure. At pass 9 the
  %% two lines were set 0.24cm apart, tighter than solid for 8pt type, and the
  %% descender of "bootstrap" met the dot of the "i" below it.
  %% PASS 11, Y12: 0.28cm -> 0.30cm between the two lines, and 0.46cm from the
  %% axis-label block above rather than 0.42cm. 0.28cm is exactly solid for
  %% 7.97pt type and the two lines' boxes met; the block also has to stand off
  %% the axis-label block by more than that block's own 0.38cm internal
  %% leading, or the four lines read as one four-line paragraph.
  %% PASS 11: six characters come out of these two lines, and nothing else, to
  %% keep the figure under S15's 5.0 characters per 100mm^2 after "of seven",
  %% "clean" and "zero" went in. Both sentences say exactly what they said.
  \node[scope] at (0.00,2.08) {Intervals: (b) 95\% bootstrap over 58 ordered pair clusters, on the difference, drawn once.};
  \node[scope] at (0.00,1.78) {None is stored for \texttt{permuted} or \texttt{swap} alone, and (a) carries none.};

  % The bounding box IS the body block; nothing may stick out of it.
  % PASS 11, Y13: the floor goes 1.72 -> 1.68 to keep the foot block's
  % descenders inside. This \path is the one rounded literal left in the file
  % and it carries no value: 13.9957 is a box corner, not a coordinate any
  % mark is drawn at, and every data x is computed by \ga or \gb from the
  % full-precision float (Y3).
  \path (0.00,1.68) rectangle (13.9957,10.18);
\end{tikzpicture}
%
% ---------------------------------------------------------------------------
% READY TO PASTE into Sections/04b-representation.tex, at sec:res-mechanism,
% immediately after the verdict paragraph (04b:332-342) or, if the page breaks
% badly, after the equivalence-criterion paragraph (04b:412-419). The same block
% was staged, at pass 9, in a figs/_PENDING_FLOATS.tex that does not exist;
% PASS 10 supersedes that: the float wrapper is HERE and the caption is in
% figs/_PENDING_CAPTIONS.tex, which does exist as of this pass (X15).
%
% PASS 9 NOTE. The caption below is the PROPOSED replacement and it is not yet
% in Sections/. No builder on this pass may edit Sections/04b-representation.tex,
% so the same text is written into figs/_PENDING_CAPTIONS.tex keyed by
% \label{fig:materiality} for a human to integrate. Until that happens the
% caption printed in main.pdf is the pass-8 one, and three of its clauses are
% false of the rebuilt drawing -- see CAPTION NOTE, PASS 9 at the foot.
%
% \begin{figure}[htbp]
% \centering
% \input{figs/materiality}
% \caption[The effect and its gap]{\textbf{The model
% does move its preference towards the other drug's own held-out sentence, but
% by one fiftieth of the preference it already has and one fifth of the margin
% set in advance.} Residual arm, layer $12$, $\alpha = 0.5$, from
% \texttt{probe\_arm\_residual.json}; normalised, the same estimate is the
% \ci{+0.0197}{+0.0063}{+0.0335} of \cref{tab:identity} and
% \cref{fig:comparators} row (c). \emph{(a)} The model's clean preference gap
% between the two sentences compared, measured with no intervention and drawn
% end to end between two ink caps, so that the axis's full width is the
% denominator every effect in this section is divided by; the effect is the
% accent disc near the origin, named at itself and dotted down onto the rule,
% and the bracket below the axis label marks the eighth of the gap that (b)
% opens out. \emph{(b)} That first eighth at exactly $8\times$ --- exactly,
% because the lower scale is the upper one divided by eight and not by a
% rounded number --- in the same units, where the \texttt{permuted} control at
% $0.00064$ nats, the \texttt{swap} arm at $0.00190$ and their difference,
% $+0.0013$, separate; the difference is the only one of the three this thesis
% quotes anywhere. The black rule in both panels is the $10\%$ materiality
% margin pre-registered for calling an effect material, drawn in the same
% weight and colour as in \cref{fig:family} and \cref{fig:comparators} row (c).
% \textbf{Which bootstrap draw this is.} The bar is
% \texttt{layers.12.swap.ci\_perm}, the top-level draw, $p = 0.001$ over a
% $4000$-resample cluster bootstrap on $60$ rows in $58$ pair clusters over
% $24$ anchor drugs. The arm's file stores a second draw of the same point
% estimate inside the $\alpha = 0.5$ rung of the ladder, at $p = 0.004$, and
% \cref{fig:family}'s lower strip draws that one; the two lie either side of
% the $0.05/26$ admitted at rank one, so the choice decides whether anything in
% the family survives multiplicity correction. This figure quotes the top-level
% draw because \cref{tab:identity}, \cref{fig:comparators} row (c) and the
% correction itself all rest on it. \textbf{What the interval shows, and what
% it does not.} It excludes zero \emph{and} lies wholly inside the margin:
% under textbook TOST that configuration is returned as equivalence and read as
% a null, whereas this chapter's rule returns a directional effect smaller than
% the margin --- one fiftieth of the gap, one fifth of the margin.
% \textbf{No interval is drawn on the two arms or in (a).} None is stored for
% \texttt{effect\_swap} or \texttt{effect\_perm} alone, so their marks are bare;
% the one that is stored is on their difference, which is (a)'s single mark as
% well as (b)'s third row, and it is drawn once, at the scale where it is
% legible. What may not be read across: layer $12$ is one of seven
% pre-registered depths and the only one whose interval excludes zero, the
% other six being in \cref{fig:family}'s grid, and this is one arm at one depth
% at one rung of the displacement ladder; because the five arms emit different
% target formats their clean gaps run from this arm's $0.064$ nats to $0.314$
% in the consensus arm, nearly fivefold, so a raw-nats axis is legitimate here
% and nowhere across arms --- \cref{fig:family}'s lower strip is a single-arm
% object for the same reason.}
% \label{fig:materiality}
% \end{figure}
%
% CAPTION NOTE, PASS 10. The pass-9 note below is kept; these are the further
% changes this pass made, and the reason for each.
%   * "the accent mark ... labelled there" became "named at itself", because at
%     pass 9 the label was anchored on the PANEL origin and not on the mark, so
%     the caption's claim was false of the drawing (X4). It is true now.
%   * "drawn end to end" gained "between two ink caps", because the span now
%     has two matching end caps and the caption should name the device that
%     makes the axis read as a bounded span (X7).
%   * THE TWO OPERANDS ARE PRINTED: the permuted control at 0.00064 nats and
%     the swap arm at 0.00190, rounded from effect_perm 0.0006409306049362111
%     and effect_swap 0.0019027651015006252. Band (b) draws them as positions
%     and prints neither, and the subtraction they stand in is the reason both
%     rows exist; a reader could previously check it only by measuring three
%     rows with a ruler. Both are new readings quoted in no sentence of the
%     thesis, which the ledger has always recorded, so the caption owns them --
%     and it now says so by naming the difference as the only one of the three
%     the thesis quotes.
%   * LAYER 12 IS DISCLOSED AS THE FAVOURABLE DEPTH: "one of seven
%     pre-registered depths and the only one whose interval excludes zero".
%     Recomputed from probe_arm_residual.json layers.<L>.swap.diff_perm_norm
%     and p_perm: +0.0022 p 0.583 (L2), -0.0050 p 0.139 (L4), -0.0025 p 0.4495
%     (L6), +0.0041 p 0.3955 (L8), -0.0019 p 0.635 (L9), +0.0197 p 0.001 (L12),
%     +0.0059 p 0.4885 (L16); six of seven ci_perm_norm straddle zero.
%     fig:hook's caption discloses its layer 4 as the best of seven and this one
%     did not, though "surviving" in this figure's own subject is doing exactly
%     that selection work. The RUNG is not favourable and is not flagged: in raw
%     nats alpha = 0.5 gives the smallest of the five ladder effects (0.0013
%     against 0.0017 / 0.0044 / 0.0061 / 0.0046), so the drawn rung is the
%     conservative one.
%   * "separate into three rows" -- pass 9 cut "into three rows"; the phrase
%     stays cut, since the rows still carry no dividers.
%   * The pass-9 caption is otherwise unchanged, and it is now WRITTEN OUT, in
%     full and uncommented, in figs/_PENDING_CAPTIONS.tex keyed by
%     \label{fig:materiality}, together with the two sentences deleted from the
%     drawing at W9 and a list of the clauses of the caption PRINTED today
%     (Sections/04b-representation.tex:597-620) that the rebuild made false.
%     At pass 9 that file was asserted to exist and did not (X15).
%
% CAPTION NOTE, PASS 9. What changed, and why each clause had to.
%   * "the effect is the accent mark near its origin, stemmed to the axis and
%     labelled there, and the brace under it marks the eighth ..." became
%     "the accent disc ... dotted down onto the rule ... and the bracket beneath
%     the axis marks the eighth ...". The mark is a disc and no longer a bar
%     (W6), the stem is now a dotted projection and not an accent stem (W11),
%     and the magnifier is a bracket and not a brace (W10).
%   * "only their difference does, and it is the one bar on the figure" became a
%     bolded sentence of its own, because at pass 8 the caption said one bar and
%     the drawing showed two -- the same estimate at two magnifications -- and a
%     naive reader counted them and had to work out which two quantities they
%     were. There is now exactly one bar and the caption says where it is.
%   * "separate into three rows" lost "into three rows": the rows no longer carry
%     dividers, so there is nothing to count (W5).
%   * WHICH BOOTSTRAP DRAW IS DRAWN is new, and it is the substantive addition.
%     4.12 and 4.13's lower strip draw the identical point estimate
%     0.0012618344965644142 in identical units and visibly identical bars, from
%     two different bootstrap draws with p = 0.001 and p = 0.004, and neither
%     drawing said so. Both ledgers knew (materiality.json's
%     not_drawn_second_bootstrap_draw; family.json's ladder note), so this was a
%     disclosure failure on the page and not a data error.
%   * THE TOST ASYMMETRY moved here in full from the drawing (W9), where it was
%     two ink lines at 8pt that a reader met before the real caption. It is the
%     section's subtlest point and it is argument, so it belongs here.
%   * "one fiftieth of the gap; one fifth of the margin" moved here from the
%     drawing's accent italic (W9). It was already the caption's lede in
%     substance; it now appears once.
%   * THE $10\%$ MARGIN's rendering is named, because this is the first build in
%     which 4.7 row (c), 4.12 and 4.13 draw one pre-registered threshold one way.
% ---------------------------------------------------------------------------
%
% ---------------------------------------------------------------------------
% CAPTION NOTE, PASS 11 -- 2026-08-19. THE COMMENTED \caption ABOVE IS THE
% PASS-10 TEXT AND IS NO LONGER THE PROPOSAL. It is left in place because the
% passes above are the audit trail and because the float wrapper round it is
% still the one to paste; the AUTHORITATIVE caption is the one in
% figs/_PENDING_CAPTIONS.tex under \label{fig:materiality}, revised in this
% pass with the drawing. Four of its clauses moved, and one item in that file
% is now marked BLOCKING.
%   * (a)'s single mark now PRINTS ITS VALUE on the drawing (Y2), so the
%     caption stops merely naming it and says that panel (a) alone answers the
%     section's question: the disc is labelled with +0.0013, the threshold with
%     0.00642 and the axis's own right end with 0.0642.
%   * (b)'s row three no longer prints a value (Y8), so the caption's "their
%     difference, $+0.0013$" is now the only statement of that numeral for
%     panel (b).
%   * (b)'s rule is capped in ink at its right end as (a)'s is (Y10), and (b)'s
%     zero rule is labelled "zero" (Y6), so the caption's "It excludes zero"
%     now points at a named object rather than at an unlabelled dotted rule.
%   * The panel head reads "(a) the whole CLEAN preference gap" (Y1).
%   * BLOCKING, and it is the one item in the file whose omission changes what
%     the figure claims rather than how it is described: 04b:617-618 says only
%     "this is one arm at one depth at one rung", which is scope. Layer 12 is
%     the argmax of the anchor seed's seven-depth sweep and the only one of the
%     seven whose interval excludes zero. The drawing now says "layer 12 of
%     seven" (Y4) -- that the depth came from a set -- but a drawing cannot say
%     WHICH of the set without becoming an argument. The proposed caption
%     carries it as a bolded clause of its own, with the seven normalised
%     effects and their bootstrap p, and with the countervailing fact that the
%     RUNG is not favourable: at alpha = 0.5 the raw effect is the smallest of
%     the five rungs, 0.0013 against 0.0017, 0.0044, 0.0061 and 0.0046
%     (probe_arm_residual.json :: layers.12.swap.ladder[*].diff_perm).
% ---------------------------------------------------------------------------

% \caption[The effect and its gap]{\textbf{The model
% does move its preference towards the other drug's own held-out sentence, but
% by one fiftieth of the preference it already has and one fifth of the margin
% set in advance.} Residual arm, layer $12$, $\alpha = 0.5$, from
% \texttt{probe\_arm\_residual.json}; normalised, the same estimate is the
% \ci{+0.0197}{+0.0063}{+0.0335} of \cref{tab:identity} and
% \cref{fig:comparators} row (c). \emph{(a)} The model's clean preference gap
% between the two sentences compared, measured with no intervention and drawn
% end to end between two ink caps, so that the axis's full width is the
% denominator every effect in this section is divided by; the effect is the
% accent disc near the origin, named at itself and dotted down onto the rule,
% and the bracket below the axis label marks the eighth of the gap that (b)
% opens out. \emph{(b)} That first eighth at exactly $8\times$ --- exactly,
% because the lower scale is the upper one divided by eight and not by a
% rounded number --- in the same units, where the \texttt{permuted} control at
% $0.00064$ nats, the \texttt{swap} arm at $0.00190$ and their difference,
% $+0.0013$, separate; the difference is the only one of the three this thesis
% quotes anywhere. The black rule in both panels is the $10\%$ materiality
% margin pre-registered for calling an effect material, drawn in the same
% weight and colour as in \cref{fig:family} and \cref{fig:comparators} row (c).
% \textbf{Which bootstrap draw this is.} The bar is
% \texttt{layers.12.swap.ci\_perm}, the top-level draw, $p = 0.001$ over a
% $4000$-resample cluster bootstrap on $60$ rows in $58$ pair clusters over
% $24$ anchor drugs. The arm's file stores a second draw of the same point
% estimate inside the $\alpha = 0.5$ rung of the ladder, at $p = 0.004$, and
% \cref{fig:family}'s lower strip draws that one; the two lie either side of
% the $0.05/26$ admitted at rank one, so the choice decides whether anything in
% the family survives multiplicity correction. This figure quotes the top-level
% draw because \cref{tab:identity}, \cref{fig:comparators} row (c) and the
% correction itself all rest on it. \textbf{What the interval shows, and what
% it does not.} It excludes zero \emph{and} lies wholly inside the margin:
% under textbook TOST that configuration is returned as equivalence and read as
% a null, whereas this chapter's rule returns a directional effect smaller than
% the margin --- one fiftieth of the gap, one fifth of the margin.
% \textbf{No interval is drawn on the two arms or in (a).} None is stored for
% \texttt{effect\_swap} or \texttt{effect\_perm} alone, so their marks are bare;
% the one that is stored is on their difference, which is (a)'s single mark as
% well as (b)'s third row, and it is drawn once, at the scale where it is
% legible. What may not be read across: layer $12$ is one of seven
% pre-registered depths and the only one whose interval excludes zero, the
% other six being in \cref{fig:family}'s grid, and this is one arm at one depth
% at one rung of the displacement ladder; because the five arms emit different
% target formats their clean gaps run from this arm's $0.064$ nats to $0.314$
% in the consensus arm, nearly fivefold, so a raw-nats axis is legitimate here
% and nowhere across arms --- \cref{fig:family}'s lower strip is a single-arm
% object for the same reason.}
% \label{fig:materiality}
% \end{figure}
%
% CAPTION NOTE, PASS 10. The pass-9 note below is kept; these are the further
% changes this pass made, and the reason for each.
%   * "the accent mark ... labelled there" became "named at itself", because at
%     pass 9 the label was anchored on the PANEL origin and not on the mark, so
%     the caption's claim was false of the drawing (X4). It is true now.
%   * "drawn end to end" gained "between two ink caps", because the span now
%     has two matching end caps and the caption should name the device that
%     makes the axis read as a bounded span (X7).
%   * THE TWO OPERANDS ARE PRINTED: the permuted control at 0.00064 nats and
%     the swap arm at 0.00190, rounded from effect_perm 0.0006409306049362111
%     and effect_swap 0.0019027651015006252. Band (b) draws them as positions
%     and prints neither, and the subtraction they stand in is the reason both
%     rows exist; a reader could previously check it only by measuring three
%     rows with a ruler. Both are new readings quoted in no sentence of the
%     thesis, which the ledger has always recorded, so the caption owns them --
%     and it now says so by naming the difference as the only one of the three
%     the thesis quotes.
%   * LAYER 12 IS DISCLOSED AS THE FAVOURABLE DEPTH: "one of seven
%     pre-registered depths and the only one whose interval excludes zero".
%     Recomputed from probe_arm_residual.json layers.<L>.swap.diff_perm_norm
%     and p_perm: +0.0022 p 0.583 (L2), -0.0050 p 0.139 (L4), -0.0025 p 0.4495
%     (L6), +0.0041 p 0.3955 (L8), -0.0019 p 0.635 (L9), +0.0197 p 0.001 (L12),
%     +0.0059 p 0.4885 (L16); six of seven ci_perm_norm straddle zero.
%     fig:hook's caption discloses its layer 4 as the best of seven and this one
%     did not, though "surviving" in this figure's own subject is doing exactly
%     that selection work. The RUNG is not favourable and is not flagged: in raw
%     nats alpha = 0.5 gives the smallest of the five ladder effects (0.0013
%     against 0.0017 / 0.0044 / 0.0061 / 0.0046), so the drawn rung is the
%     conservative one.
%   * "separate into three rows" -- pass 9 cut "into three rows"; the phrase
%     stays cut, since the rows still carry no dividers.
%   * The pass-9 caption is otherwise unchanged, and it is now WRITTEN OUT, in
%     full and uncommented, in figs/_PENDING_CAPTIONS.tex keyed by
%     \label{fig:materiality}, together with the two sentences deleted from the
%     drawing at W9 and a list of the clauses of the caption PRINTED today
%     (Sections/04b-representation.tex:597-620) that the rebuild made false.
%     At pass 9 that file was asserted to exist and did not (X15).
%
% CAPTION NOTE, PASS 9. What changed, and why each clause had to.
%   * "the effect is the accent mark near its origin, stemmed to the axis and
%     labelled there, and the brace under it marks the eighth ..." became
%     "the accent disc ... dotted down onto the rule ... and the bracket beneath
%     the axis marks the eighth ...". The mark is a disc and no longer a bar
%     (W6), the stem is now a dotted projection and not an accent stem (W11),
%     and the magnifier is a bracket and not a brace (W10).
%   * "only their difference does, and it is the one bar on the figure" became a
%     bolded sentence of its own, because at pass 8 the caption said one bar and
%     the drawing showed two -- the same estimate at two magnifications -- and a
%     naive reader counted them and had to work out which two quantities they
%     were. There is now exactly one bar and the caption says where it is.
%   * "separate into three rows" lost "into three rows": the rows no longer carry
%     dividers, so there is nothing to count (W5).
%   * WHICH BOOTSTRAP DRAW IS DRAWN is new, and it is the substantive addition.
%     4.12 and 4.13's lower strip draw the identical point estimate
%     0.0012618344965644142 in identical units and visibly identical bars, from
%     two different bootstrap draws with p = 0.001 and p = 0.004, and neither
%     drawing said so. Both ledgers knew (materiality.json's
%     not_drawn_second_bootstrap_draw; family.json's ladder note), so this was a
%     disclosure failure on the page and not a data error.
%   * THE TOST ASYMMETRY moved here in full from the drawing (W9), where it was
%     two ink lines at 8pt that a reader met before the real caption. It is the
%     section's subtlest point and it is argument, so it belongs here.
%   * "one fiftieth of the gap; one fifth of the margin" moved here from the
%     drawing's accent italic (W9). It was already the caption's lede in
%     substance; it now appears once.
%   * THE $10\%$ MARGIN's rendering is named, because this is the first build in
%     which 4.7 row (c), 4.12 and 4.13 draw one pre-registered threshold one way.
% ---------------------------------------------------------------------------
%
% ---------------------------------------------------------------------------
% CAPTION NOTE, PASS 11 -- 2026-08-19. THE COMMENTED \caption ABOVE IS THE
% PASS-10 TEXT AND IS NO LONGER THE PROPOSAL. It is left in place because the
% passes above are the audit trail and because the float wrapper round it is
% still the one to paste; the AUTHORITATIVE caption is the one in
% figs/_PENDING_CAPTIONS.tex under \label{fig:materiality}, revised in this
% pass with the drawing. Four of its clauses moved, and one item in that file
% is now marked BLOCKING.
%   * (a)'s single mark now PRINTS ITS VALUE on the drawing (Y2), so the
%     caption stops merely naming it and says that panel (a) alone answers the
%     section's question: the disc is labelled with +0.0013, the threshold with
%     0.00642 and the axis's own right end with 0.0642.
%   * (b)'s row three no longer prints a value (Y8), so the caption's "their
%     difference, $+0.0013$" is now the only statement of that numeral for
%     panel (b).
%   * (b)'s rule is capped in ink at its right end as (a)'s is (Y10), and (b)'s
%     zero rule is labelled "zero" (Y6), so the caption's "It excludes zero"
%     now points at a named object rather than at an unlabelled dotted rule.
%   * The panel head reads "(a) the whole CLEAN preference gap" (Y1).
%   * BLOCKING, and it is the one item in the file whose omission changes what
%     the figure claims rather than how it is described: 04b:617-618 says only
%     "this is one arm at one depth at one rung", which is scope. Layer 12 is
%     the argmax of the anchor seed's seven-depth sweep and the only one of the
%     seven whose interval excludes zero. The drawing now says "layer 12 of
%     seven" (Y4) -- that the depth came from a set -- but a drawing cannot say
%     WHICH of the set without becoming an argument. The proposed caption
%     carries it as a bolded clause of its own, with the seven normalised
%     effects and their bootstrap p, and with the countervailing fact that the
%     RUNG is not favourable: at alpha = 0.5 the raw effect is the smallest of
%     the five rungs, 0.0013 against 0.0017, 0.0044, 0.0061 and 0.0046
%     (probe_arm_residual.json :: layers.12.swap.ladder[*].diff_perm).
% ---------------------------------------------------------------------------

% \caption[The effect and its gap]{\textbf{The model
% does move its preference towards the other drug's own held-out sentence, but
% by one fiftieth of the preference it already has and one fifth of the margin
% set in advance.} Residual arm, layer $12$, $\alpha = 0.5$, from
% \texttt{probe\_arm\_residual.json}; normalised, the same estimate is the
% \ci{+0.0197}{+0.0063}{+0.0335} of \cref{tab:identity} and
% \cref{fig:comparators} row (c). \emph{(a)} The model's clean preference gap
% between the two sentences compared, measured with no intervention and drawn
% end to end between two ink caps, so that the axis's full width is the
% denominator every effect in this section is divided by; the effect is the
% accent disc near the origin, named at itself and dotted down onto the rule,
% and the bracket below the axis label marks the eighth of the gap that (b)
% opens out. \emph{(b)} That first eighth at exactly $8\times$ --- exactly,
% because the lower scale is the upper one divided by eight and not by a
% rounded number --- in the same units, where the \texttt{permuted} control at
% $0.00064$ nats, the \texttt{swap} arm at $0.00190$ and their difference,
% $+0.0013$, separate; the difference is the only one of the three this thesis
% quotes anywhere. The black rule in both panels is the $10\%$ materiality
% margin pre-registered for calling an effect material, drawn in the same
% weight and colour as in \cref{fig:family} and \cref{fig:comparators} row (c).
% \textbf{Which bootstrap draw this is.} The bar is
% \texttt{layers.12.swap.ci\_perm}, the top-level draw, $p = 0.001$ over a
% $4000$-resample cluster bootstrap on $60$ rows in $58$ pair clusters over
% $24$ anchor drugs. The arm's file stores a second draw of the same point
% estimate inside the $\alpha = 0.5$ rung of the ladder, at $p = 0.004$, and
% \cref{fig:family}'s lower strip draws that one; the two lie either side of
% the $0.05/26$ admitted at rank one, so the choice decides whether anything in
% the family survives multiplicity correction. This figure quotes the top-level
% draw because \cref{tab:identity}, \cref{fig:comparators} row (c) and the
% correction itself all rest on it. \textbf{What the interval shows, and what
% it does not.} It excludes zero \emph{and} lies wholly inside the margin:
% under textbook TOST that configuration is returned as equivalence and read as
% a null, whereas this chapter's rule returns a directional effect smaller than
% the margin --- one fiftieth of the gap, one fifth of the margin.
% \textbf{No interval is drawn on the two arms or in (a).} None is stored for
% \texttt{effect\_swap} or \texttt{effect\_perm} alone, so their marks are bare;
% the one that is stored is on their difference, which is (a)'s single mark as
% well as (b)'s third row, and it is drawn once, at the scale where it is
% legible. What may not be read across: layer $12$ is one of seven
% pre-registered depths and the only one whose interval excludes zero, the
% other six being in \cref{fig:family}'s grid, and this is one arm at one depth
% at one rung of the displacement ladder; because the five arms emit different
% target formats their clean gaps run from this arm's $0.064$ nats to $0.314$
% in the consensus arm, nearly fivefold, so a raw-nats axis is legitimate here
% and nowhere across arms --- \cref{fig:family}'s lower strip is a single-arm
% object for the same reason.}
% \label{fig:materiality}
% \end{figure}
%
% CAPTION NOTE, PASS 10. The pass-9 note below is kept; these are the further
% changes this pass made, and the reason for each.
%   * "the accent mark ... labelled there" became "named at itself", because at
%     pass 9 the label was anchored on the PANEL origin and not on the mark, so
%     the caption's claim was false of the drawing (X4). It is true now.
%   * "drawn end to end" gained "between two ink caps", because the span now
%     has two matching end caps and the caption should name the device that
%     makes the axis read as a bounded span (X7).
%   * THE TWO OPERANDS ARE PRINTED: the permuted control at 0.00064 nats and
%     the swap arm at 0.00190, rounded from effect_perm 0.0006409306049362111
%     and effect_swap 0.0019027651015006252. Band (b) draws them as positions
%     and prints neither, and the subtraction they stand in is the reason both
%     rows exist; a reader could previously check it only by measuring three
%     rows with a ruler. Both are new readings quoted in no sentence of the
%     thesis, which the ledger has always recorded, so the caption owns them --
%     and it now says so by naming the difference as the only one of the three
%     the thesis quotes.
%   * LAYER 12 IS DISCLOSED AS THE FAVOURABLE DEPTH: "one of seven
%     pre-registered depths and the only one whose interval excludes zero".
%     Recomputed from probe_arm_residual.json layers.<L>.swap.diff_perm_norm
%     and p_perm: +0.0022 p 0.583 (L2), -0.0050 p 0.139 (L4), -0.0025 p 0.4495
%     (L6), +0.0041 p 0.3955 (L8), -0.0019 p 0.635 (L9), +0.0197 p 0.001 (L12),
%     +0.0059 p 0.4885 (L16); six of seven ci_perm_norm straddle zero.
%     fig:hook's caption discloses its layer 4 as the best of seven and this one
%     did not, though "surviving" in this figure's own subject is doing exactly
%     that selection work. The RUNG is not favourable and is not flagged: in raw
%     nats alpha = 0.5 gives the smallest of the five ladder effects (0.0013
%     against 0.0017 / 0.0044 / 0.0061 / 0.0046), so the drawn rung is the
%     conservative one.
%   * "separate into three rows" -- pass 9 cut "into three rows"; the phrase
%     stays cut, since the rows still carry no dividers.
%   * The pass-9 caption is otherwise unchanged, and it is now WRITTEN OUT, in
%     full and uncommented, in figs/_PENDING_CAPTIONS.tex keyed by
%     \label{fig:materiality}, together with the two sentences deleted from the
%     drawing at W9 and a list of the clauses of the caption PRINTED today
%     (Sections/04b-representation.tex:597-620) that the rebuild made false.
%     At pass 9 that file was asserted to exist and did not (X15).
%
% CAPTION NOTE, PASS 9. What changed, and why each clause had to.
%   * "the effect is the accent mark near its origin, stemmed to the axis and
%     labelled there, and the brace under it marks the eighth ..." became
%     "the accent disc ... dotted down onto the rule ... and the bracket beneath
%     the axis marks the eighth ...". The mark is a disc and no longer a bar
%     (W6), the stem is now a dotted projection and not an accent stem (W11),
%     and the magnifier is a bracket and not a brace (W10).
%   * "only their difference does, and it is the one bar on the figure" became a
%     bolded sentence of its own, because at pass 8 the caption said one bar and
%     the drawing showed two -- the same estimate at two magnifications -- and a
%     naive reader counted them and had to work out which two quantities they
%     were. There is now exactly one bar and the caption says where it is.
%   * "separate into three rows" lost "into three rows": the rows no longer carry
%     dividers, so there is nothing to count (W5).
%   * WHICH BOOTSTRAP DRAW IS DRAWN is new, and it is the substantive addition.
%     4.12 and 4.13's lower strip draw the identical point estimate
%     0.0012618344965644142 in identical units and visibly identical bars, from
%     two different bootstrap draws with p = 0.001 and p = 0.004, and neither
%     drawing said so. Both ledgers knew (materiality.json's
%     not_drawn_second_bootstrap_draw; family.json's ladder note), so this was a
%     disclosure failure on the page and not a data error.
%   * THE TOST ASYMMETRY moved here in full from the drawing (W9), where it was
%     two ink lines at 8pt that a reader met before the real caption. It is the
%     section's subtlest point and it is argument, so it belongs here.
%   * "one fiftieth of the gap; one fifth of the margin" moved here from the
%     drawing's accent italic (W9). It was already the caption's lede in
%     substance; it now appears once.
%   * THE $10\%$ MARGIN's rendering is named, because this is the first build in
%     which 4.7 row (c), 4.12 and 4.13 draw one pre-registered threshold one way.
% ---------------------------------------------------------------------------
%
% ---------------------------------------------------------------------------
% CAPTION NOTE, PASS 11 -- 2026-08-19. THE COMMENTED \caption ABOVE IS THE
% PASS-10 TEXT AND IS NO LONGER THE PROPOSAL. It is left in place because the
% passes above are the audit trail and because the float wrapper round it is
% still the one to paste; the AUTHORITATIVE caption is the one in
% figs/_PENDING_CAPTIONS.tex under \label{fig:materiality}, revised in this
% pass with the drawing. Four of its clauses moved, and one item in that file
% is now marked BLOCKING.
%   * (a)'s single mark now PRINTS ITS VALUE on the drawing (Y2), so the
%     caption stops merely naming it and says that panel (a) alone answers the
%     section's question: the disc is labelled with +0.0013, the threshold with
%     0.00642 and the axis's own right end with 0.0642.
%   * (b)'s row three no longer prints a value (Y8), so the caption's "their
%     difference, $+0.0013$" is now the only statement of that numeral for
%     panel (b).
%   * (b)'s rule is capped in ink at its right end as (a)'s is (Y10), and (b)'s
%     zero rule is labelled "zero" (Y6), so the caption's "It excludes zero"
%     now points at a named object rather than at an unlabelled dotted rule.
%   * The panel head reads "(a) the whole CLEAN preference gap" (Y1).
%   * BLOCKING, and it is the one item in the file whose omission changes what
%     the figure claims rather than how it is described: 04b:617-618 says only
%     "this is one arm at one depth at one rung", which is scope. Layer 12 is
%     the argmax of the anchor seed's seven-depth sweep and the only one of the
%     seven whose interval excludes zero. The drawing now says "layer 12 of
%     seven" (Y4) -- that the depth came from a set -- but a drawing cannot say
%     WHICH of the set without becoming an argument. The proposed caption
%     carries it as a bolded clause of its own, with the seven normalised
%     effects and their bootstrap p, and with the countervailing fact that the
%     RUNG is not favourable: at alpha = 0.5 the raw effect is the smallest of
%     the five rungs, 0.0013 against 0.0017, 0.0044, 0.0061 and 0.0046
%     (probe_arm_residual.json :: layers.12.swap.ladder[*].diff_perm).
% ---------------------------------------------------------------------------

\caption[The effect and its gap]{\textbf{The model does move its preference
towards the other drug's own held-out sentence, but by one fiftieth of the
preference it already has and one fifth of the margin set in advance.}
Residual arm, layer $12$, $\alpha = 0.5$, from
\texttt{probe\_arm\_residual.json}; normalised, the same estimate is the
\ci{+0.0197}{+0.0063}{+0.0335} of \cref{tab:identity} and
\cref{fig:comparators} row~(c). \emph{(a)} The model's clean preference gap
between the two sentences compared, measured with no intervention, drawn end to
end, so that the axis's full width is the denominator every effect in this
section is divided by; the effect is the accent disc near the origin, dotted
down onto the rule and named above it, and the bracket beneath the axis label
marks the eighth of it that~(b) opens out. \emph{(b)} That first eighth at
exactly $8\times$ --- exactly, because the lower scale is the upper one divided
by eight and not by a rounded number, which is why the two tick lanes fall on
the same four positions --- in the same units, where the \texttt{swap} arm, the
\texttt{permuted} control and their difference separate. The black rule in both
panels is the $10\%$ materiality margin pre-registered for calling an effect
material, drawn in the same weight and colour as in \cref{fig:family} and
\cref{fig:comparators} row~(c). \textbf{Which bootstrap draw this is.} The bar
is \texttt{layers.12.swap.ci\_perm}, the top-level draw, $p = 0.001$ over a
$4000$-resample cluster bootstrap on $60$ rows in $58$ pair clusters over $24$
anchor drugs. The arm's file stores a second draw of the same point estimate
inside the $\alpha = 0.5$ rung of the ladder, at $p = 0.004$, and
\cref{fig:family}'s lower strip draws that one; the two lie either side of the
$0.05/26$ admitted at rank one, so the choice decides whether anything in the
family survives multiplicity correction. This figure quotes the top-level draw
because \cref{tab:identity}, \cref{fig:comparators} row~(c) and the correction
itself all rest on it. \textbf{What the interval shows, and what it does not.}
It excludes zero \emph{and} lies wholly inside the margin: under textbook TOST
that configuration is returned as equivalence and read as a null, whereas this
chapter's rule returns a directional effect smaller than the margin --- one
fiftieth of the gap, one fifth of the margin. \textbf{No interval is drawn on
the two arms or in~(a).} None is stored for \texttt{effect\_swap} or
\texttt{effect\_perm} alone, so their marks are bare; the one stored interval
is on the difference and is drawn at the scale where it is legible, which
is~(b), and nowhere else. What may not be read across: this is one arm at one
depth at one rung of the displacement ladder, and because the five arms emit
different target formats their clean gaps run from this arm's $0.064$ nats to
$0.314$ in the consensus arm, nearly fivefold, so a raw-nats axis is legitimate
here and nowhere across arms --- \cref{fig:family}'s lower strip is a
single-arm object for the same reason.}
\label{fig:materiality}
\end{figure}

\paragraph{First, a contrast score} It is
\begin{equation}
 \log P\!\left(g_B \mid \textrm{prompt}_A\right) \;-\;
 \log P\!\left(g_A \mid \textrm{prompt}_A\right),
\end{equation}
where $g_A$ and $g_B$ are real cell sentences emitted under drugs A and B by
held-out cells, each log-probability a per-token mean. Any offset uniform across
positions, an inflated logsumexp normaliser included, then cancels exactly. It
does \emph{not} cancel in the stronger sense one might assume. Once $g_A$ and
$g_B$ diverge the two teacher-forced sequences visit different states, so the
normaliser's state-dependent part survives, unbounded here. The \texttt{permuted}
arm controls for that, sharing anchor, construction and norm and differing only
in naming the wrong drug. What survives the two-arm comparison is the model's
preference between drug B's own held-out sentence and drug A's.

\paragraph{Second, one anchor and one norm} Writing $\Pi$ for the orthogonal
projector onto the purified slab, the \texttt{swap} arm injects
$\Pi(\mu_B - \mu_A)$, pointing at the drug whose sentence is scored; the
\texttt{permuted} arm injects $\Pi(\mu_C - \mu_A)$ at identical norm, same anchor
and construction, wrong drug. An isotropic random direction is kept only as a
damage diagnostic, since it can disrupt the model in ways an observed class-mean
displacement does not. The subtraction cancels the global mean exactly, and
nothing is ablated, so the control has no removed term to trip over.

\paragraph{The three failed designs taught the constraints} A control drawn in
the full $2048$-dimensional stream puts $\approx 99.5\%$ of its squared norm
where the intervention never acts, so it must live in the treatment's subspace.
Ablate-then-add delivers (added $-$ removed), so an uncentred class mean turns
substitution into restoration, hence no ablation and an exactly centred
displacement. A norm-matched random direction inflates the logsumexp normaliser
(measured $+2.75$ nats), depressing every target whether or not it carries
identity. The planted-world selftest validates the directional logic but cannot
speak to the normaliser, since there each drug's sentence is a single token.
Algebra in \cref{sec:app-identity}.

\paragraph{Out of sample, and on a comparable scale} Subspaces and class means are
fit on one slice of cells; $g_B$ comes from a third disjoint slice, never from
the rows defining $\mu_B$, else the injected vector and the scored sentence would
share cells and the test would be partly circular. Because the five arms emit
different target formats, every effect is divided by that model's own clean
preference gap, measured with no intervention.

\paragraph{Four guards carry the inference} A drug-proxy screen and an
injected-norm diagnostic are engineering, deferred to \cref{sec:app-identity}.

The first is the uncertainty model, a cluster bootstrap over ordered (A,B) pairs,
because rows sharing a pair share the injected vector and the scored sentence. At
the anchor seed's layer $12$ that is $60$ rows in $58$ pair clusters over $24$
anchor drugs, close to a row bootstrap. It guards against sharing within a pair,
not across pairs sharing an anchor or target drug.\gloss{cluster
bootstrap}{resamples whole (A,B) pairs, not rows, so a shared vector
counts once}

The second is the equivalence criterion, stricter than TOST. The interval must
lie inside a pre-registered margin, set at $10\%$ of the model's own clean
preference gap, \emph{and} contain zero. That conjunct runs one way only; it can
withhold an equivalence verdict, never manufacture one. Under textbook TOST the
headline \ci{+0.0197}{+0.0063}{+0.0335} lies wholly inside the margin and would
be returned as equivalent, read as a null; this rule returns a directional effect
smaller than the margin.

The third is a per-layer headroom gate, excluding any layer where the cell-line
positive control shows less than $5\times$ dynamic range, because a saturated
instrument can support neither verdict. The fourth is a magnitude ladder,
repeating the test at $\alpha \in \{0.5, 1, 2, 5, 10\}$ times the natural
displacement, because a real identity effect should show a monotone
dose--response. The ladder is the guard that bites.

\paragraph{One survivor in twenty-six tests} The design ran across the five
training arms of \cref{tab:trainconfig} and, for the two residual arms, at three
further seeds. Of $26$ headline tests, five arms by up to seven depths, exactly
one survives Benjamini--Hochberg correction, the residual arm at layer $12$
($q = 0.026$, pooling the $26$ across arms and not correcting within each as the
pipeline stores them). That $q$ rests on a bootstrap $p$, and the arm's probe
file carries two draws of it for the same effect, identical to sixteen digits, at
$p = 0.001$ and $p = 0.004$, either side of the $0.05/26$ admitted at rank one;
on the second draw nothing survives. The consensus arm at layer $2$ ($p = 0.012$)
does not ($q = 0.156$).

\paragraph{Which rung the correction is applied to matters} The headline is the
$\alpha = 0.5$ rung, half the natural displacement, so the ladder's first entry
is that headline in raw nats ($+0.0013$) and not an independent point. At full
displacement the effect is \emph{larger} and noisier,
\ci{+0.0265}{+0.0021}{+0.0524} against \ci{+0.0197}{+0.0063}{+0.0335}, at
$p = 0.033$ against $p = 0.001$ on the draw the correction uses. Each of
$\alpha = 0.5$, $1$ and $2$ excludes zero; $\alpha = 5$ and $10$ span it. In raw
nats the rungs read $+0.0013$, $+0.0017$, $+0.0044$, $+0.0061$, $+0.0046$, rising
to $\alpha = 5$ and turning over at $\alpha = 10$ (\cref{fig:family}, lower
strip), consistent with a dose--response without establishing one. All five
re-score the same $60$ rows at scaled injection, no trend statistic is computed,
and the pre-registered monotonicity holds over four rungs of five.

But applying the $26$-test correction to the $\alpha = 1$ slice leaves
\emph{nothing} surviving, and applying it to $\alpha = 2$ promotes a different
arm at a different depth (\texttt{residual\_holdout} at layer $4$, $q = 0.007$)
while demoting this one (\cref{fig:family}). The slice decides not merely
whether a survivor exists but which one, and below that the bootstrap draw
decides whether one exists at all. That is the mark of a family dominated by
noise, not of one containing a single robust effect. The multiplicity-corrected
claim is therefore a statement about the $\alpha = 0.5$ slice and is not robust
to that choice, which is why this section opened by discounting its own result
instead of closing on it.

% ===========================================================================
\begin{figure}[htbp]
\begin{fullwidth}
\centering
\includegraphics{family}
\caption[The single survivor moves with the rung]{\textbf{The one test that
survives multiplicity correction changes arm and depth with the rung of the
displacement ladder, and at one rung there is no survivor at all.} Each panel is
one slice of the ladder: the five training arms of \cref{tab:trainconfig} (rows)
by the seven probed depths (columns), $n = 26$ measured tests per slice, on one
diverging scale shared across the three panels. Colour is the \texttt{swap}
minus \texttt{permuted} effect as a fraction of that arm's own clean preference
gap. The layer axis is ordinal, drawn at equal spacing, and is not proportional
to depth. Nine of the $35$ grid positions were never measured and are drawn as
empty dotted outlines, not as zeros. Benjamini--Hochberg is pooled across arms
within a slice, as the text applies it and not as the pipeline stores it; $q$ is
printed in the three cells that fall below $0.10$ and ringed where $q < 0.05$.
That criterion admits the residual arm at layer $12$ at $\alpha = 0.5$
($q = 0.026$, on the stored bootstrap draw the text discusses), nothing at
$\alpha = 1$, and \texttt{residual\_holdout} at layer $4$ at $\alpha = 2$, which
demotes layer $12$ to $q = 0.058$; it is the only survivor among this chapter's
twenty-six headline tests. Beneath, the same residual arm at layer $12$ in raw
nats across all five rungs, with the $10\%$ pre-registered materiality margin as
the vertical rule. Intervals throughout are $95\%$ cluster bootstraps over
ordered (A,B) pairs, $4000$ resamples, which at this layer is $60$ rows in $58$
pair clusters over $24$ anchor drugs. The three seed replicates are in
\cref{tab:identity} and are deliberately not drawn here.}
\label{fig:family}
\end{fullwidth}
\end{figure}

% ===========================================================================
\begin{table}[htbp]
\begin{fullwidth}
\begin{threeparttable}
\caption{\textbf{The identity test replicates in the arm that behaviour calls
drug-sensitive and is absent in the arm behaviour calls drug-blind.} Effects are
a fraction of that model's own clean preference gap between the two sentences
compared (mean effect $+0.0133$ across the three seeds that measured it). The $p$
column is a two-sided cluster bootstrap over $4000$ resamples floored at
$1/4000$; seed 43's entry is that floor, reached when no resample fell at or
below zero, so it is a resolution limit and not a measurement. Seed 44 produced
no layer-12 measurement, and five of its seven layers were skipped because the
purified slab retained context, not for want of headroom; those five carry
headroom of $9.7$ to $23.3$. The pre-registered criterion of two replications in
three new seeds is met, the third a non-measurement and not a failure.}
\label{tab:identity}
\small
\begin{tabular}{@{}lrrr@{}}
\toprule
\thead{Arm / seed} & \thead{effect} & \thead{95\% CI} & \thead{$p$} \\
\midrule
residual, seed 42 & $+0.0197$ & $[+0.0063, +0.0335]$ & $0.0010$ \\
residual, seed 43 & $+0.0118$ & $[+0.0060, +0.0173]$ & $0.00025$ \\
residual, seed 45 & $+0.0085$ & $[+0.0029, +0.0143]$ & $0.0015$ \\
\midrule
single-cell arm, six depths
  & \multicolumn{3}{l}{nothing detectable at any depth (median $-0.0012$)} \\
\bottomrule
\end{tabular}
\begin{tablenotes}[flushleft]\footnotesize
\item The rows are the headline $\alpha = 0.5$ rung, to which the multiplicity
correction that admits the first row is applied and to no other.
\end{tablenotes}
\end{threeparttable}
\end{fullwidth}
\end{table}

% ===========================================================================
\paragraph{The instrument has been shown to register a positive} The residual arm
is independently known to respond to the drug its prompt names. On \texttt{train},
swapping the named drug and its mechanism for an orthogonal partner costs it
$+0.0898$ in $\NIR$, excluding zero clustered on cell line
($[+0.0382, +0.1415]$) and on drug ($[+0.0291, +0.1506]$), a behavioural result
established in \cref{sec:res-generalisation}. \forward{sec:res-generalisation}

Three things bound what it licenses. It is a training-split fact, and does not
extend to held-out combinations. It is behavioural, so it says the arm responds
to the drug its prompt names, not that the response is routed through the
subspace this test injects into. And the two measurements are not on the same
weights, the swap gap scored on the rebuilt target build and this probe on the
earlier one (\cref{tab:trainconfig}), so drug-sensitivity and identity reading
are attributed to residual-target training and not to any checkpoint. Narrower
than a routing claim, and enough. The probe detects identity reading, replicably,
in the arm behaviour calls drug-sensitive, and none at any depth in the arm that
is behaviourally drug-blind. A null from an instrument never shown to register a
positive is uninterpretable; this one has.

\paragraph{Three limits belong with the verdict} The effect appears at one depth
of seven, and layer 4 is positive in a single seed and null in the other three.
The single-cell arm has no layer-12 measurement, its cell-line positive control
reaching $3.88\times$ there, under the $5\times$ gate; so the cleanest
comparison, the same depth across two arms, does not exist, and the dissociation
rests on that arm being null at all six depths it could measure. And the sibling
residual-holdout arm does not reproduce the effect at the headline rung
($+0.0028$ and $+0.0049$ in the two seeds where it was measured); it surfaces
instead at $\alpha = 2$, rung-dependence and not a second replication.

Read with the decodability series, the picture is label information that becomes
more linearly decodable with depth once cell line, plate and dose are projected
out, while the process that writes the prediction barely consults it. That is
suggestive of the workspace account of \cref{sec:lit-interp}, an interpretation
of the pattern and not a demonstration of the mechanism.

\begin{claim}
The anchor seed puts it at \ci{+0.0197}{+0.0063}{+0.0335} of the model's own
preference gap, and $+0.0133$ averaged over the three seeds that measured it, one
to two per cent, a fifth or less of the $10\%$ margin pre-registered for calling
an effect material. The multiplicity-corrected version of that holds for the
$\alpha = 0.5$ rung on one of its two stored bootstrap draws, not for its
neighbours. At $\alpha = 1$ nothing survives correction; at $\alpha = 2$ a
different arm does. The claim is not that the decoder cannot read drug identity,
and it is not the stronger mechanistic story earlier drafts told.
\end{claim}

% ===========================================================================
\beatsection{E}{4}{Denoising the target changes the target, not the decoder}{sec:res-epsilon}

\begin{question}[4]
If the defect is a decoder that barely consults the drug region whatever is
written there, then improving what sits in that region should not help, and
improving the training target is exactly that kind of intervention. It is also
the repair a reader reaches for first, which is reason enough to test it. Three
target constructions, one scoring, one stratum. None separates from a prompt
naming the wrong drug.
\end{question}

\noindent
The interpretation at the end of the last section is \notestablished, and this
section does not proceed as though it were. The inference runs one way only. A
positive here would tell against the low-gain reading; a negative is evidence
about target-side fixes, not confirmation of the mechanism. Because one failed
construction proves nothing, two more were compared with it.

\paragraph{Three points, not a sweep} They are the original single-cell target,
one optimal-transport barycentre and the consensus target. \emph{Consensus}
groups cells by (drug, cell line, dose) and averages them into one pseudobulk
sentence, every original prompt retained so only the target changes.
\emph{Optimal transport} uses the full condition key carried by the source,
(cell line, plate, drug, dose/treatment well), computing each condition
separately, its treated cells coupled only to the shared DMSO controls from that
condition's own (cell line, plate). Treatment conditions are never pooled.

The coupling $\pi$ minimises $\sum_{ij}\pi_{ij}\lVert x_i - y_j\rVert^2$ under
balanced uniform marginals, which impose mass conservation between the empirical
control and treated measures, a modelling assumption and not an observed
biological constraint. Sinkhorn iteration~\cite{cuturi} solves it on the kernel
$\Gamma = \exp(-C/\varepsilon)$, with $C$ the pairwise squared-distance matrix.
Couplings are computed in the released Tahoe scVI latent~\cite{scvi}
($n_{\textrm{latent}} = 10$), joined per cell by barcode and not re-encoded,
since a re-encode produced a degraded latent (cell-line probe $0.11$ versus
$0.975$). Targets are built without the scVI decoder, as convex combinations of that
condition's observed treated cells, so they stay inside their convex hull, not
necessarily on the data manifold.

The regularisation parameter locates the middle construction between two limits.
As $\varepsilon\rightarrow\infty$ the coupling is uniform and the target is the
treated mean (consensus); as $\varepsilon\rightarrow 0$ it approaches an
unregularised sparse transport solution, not necessarily a map to one observed
treated cell. Only three hand-picked constructions were trained and scored, so
this is a three-point comparison and not a sweep of the family.

\paragraph{One scoring, named before the result} Every point was scored on the
identifiable stratum, the conditions \cref{sec:res-blind} selected because their
own within-well reference clears the pre-set $0.80$ bar, under the random-partner
scramble defined there. That is the stratum on which a zero-drug-information
control-copy already matches the model, so the scope limit of
\cref{sec:methods-probe} applies to every row of \cref{tab:epsilon}.

% ===========================================================================
\begin{table}[htbp]
\begin{fullwidth}
\begin{threeparttable}
\caption{\textbf{Three target constructions, and no point separates from its
random-partner scramble on the identifiable stratum} ($n = \num{204}$). Two
things bound that negative. The consensus endpoint was stopped at roughly $60\%$
of an epoch with its tier-1 NIR unscoreable, so the negative rests on the
single-cell and optimal-transport endpoints. And the optimal-transport endpoint
is not null in every frame. These rows use a random partner at $1200$ generation
tokens, while against the geometry-defined \texttt{orth} partner over all its
conditions the same checkpoint moves from \ci{+0.0144}{-0.0051}{+0.0338} at $600$
($n = \num{250}$) to \ci{+0.0292}{+0.0144}{+0.0441} at $1400$ ($n = \num{248}$),
reported in \cref{sec:res-encoding}. Comparator, population and budget differ at
once, and \cref{sec:methods-data} shows the budget is a claim-level variable, so
the frames are set side by side and not put into ratio. This is a null under this
scoring, not over the family of target constructions.}
\label{tab:epsilon}
\small
\begin{tabular}{@{}llll@{}}
\toprule
\thead{Target} & \thead{$\varepsilon$} & \thead{$\gap$} & \thead{Verdict} \\
\midrule
single cell (real treated cell) & $\approx 0$
  & \ci{+0.014}{-0.016}{+0.042} & not separated \\
OT barycentric & $0.5$
  & \ci{-0.001}{-0.018}{+0.016} & not separated \\
consensus (global mean) & $\rightarrow\infty$
  & \ci{-0.010}{-0.022}{+0.003} & not separated \\
\bottomrule
\end{tabular}
\begin{tablenotes}[flushleft]\footnotesize
\item $\gap$ is the model's $\NIR$ minus its scramble's, so its null is zero and
not chance. The three rows share a scoring, a stratum and a comparator; they are
one comparison, not three independent experiments. All three intervals are
\emph{one-way} percentile cluster bootstraps over the $40$ cell lines, not the
two-way default.
\end{tablenotes}
\end{threeparttable}
\end{fullwidth}
\end{table}

% ===========================================================================
An interval that covers zero earns its keep by what it excludes. On the
identifiable stratum the single-cell arm's scramble gap is
\ci{+0.014}{-0.016}{+0.042}, against the $0.170$ separating the model's $0.768$
from that stratum's $0.938$ reference. At the very top of the interval, naming
the right drug instead of a wrong one would buy a quarter of what is there to be
bought, and the point estimate puts it at $8\%$ of that. Taking that distance at
face value is the generous reading, for the reason declared in
\cref{sec:methods-probe}. The interval does not show the effect is zero and no
experiment here could; it puts it far below the room these conditions leave.

% ===========================================================================
\paragraph{Training was not inert, and what it did instead is a result} The
consensus model stops echoing the control. On the identifiable stratum it scores
$0.653$ against a zero-drug-information control-copy at $0.766$, where the
single-cell model sits on the control at $0.768$. Its predictions are not
mode-collapsed either; there is real between-condition variation. But that
variation is uncorrelated with real drug identity. The Mantel correlation between
the model's between-drug distance structure and the true one averages $+0.03$
over cell-line and plate groups.

The failure this literature names first for a perturbation model that
underperforms on well-calibrated metrics is mode collapse, where the prediction
ignores the input and reverts to a mean, and the repair is to restore diversity
in what the model emits~\cite{mejia}. This arm is not that. It moved off the control, kept the
variation, and spent it on the wrong axis. Producing no structure and producing
structure that does not track the drug are different failures with different
repairs.

\begin{claim}
No point in the three-point comparison separates from its scramble on the
identifiable stratum under the random-partner test. Two things bound that
negative. The consensus arm was stopped mid-epoch, so the comparison rests on the
single-cell and optimal-transport points, and the optimal-transport arm is not
null in every frame: scored against the geometry-defined \texttt{orth} partner
over all its conditions at the residual arm's generation budget, its gap excludes
zero (\cref{sec:res-encoding}). By the inference rule above that positive tells
against the low-gain reading and not for it, and being in a different frame
it does not settle the matter, but it is one more reason the interpretation of
\cref{sec:res-mechanism} is carried here as an interpretation and not a result.
What the comparison establishes is that denoising the target changes what is in
the target without changing what the decoder does with it under this scoring, and
not that no target-side fix can work. Measuring how much of the target its tokens
expose to the drug, which the next two sections do, is not the same as showing
that the target encoding is the defect (\cref{sec:methods-residual}). That
construction changes five things at once.
\end{claim}
